\documentclass[10pt,a4paper,twoside,bg=print,twocolumn,openany,nodeprecatedcode]{dndbook}
\usepackage[utf8]{inputenc}
\usepackage{tabulary}
\usepackage[normalem]{ulem}
\usepackage{color,soul}
\usepackage{float}
\usepackage{caption}
\usepackage{wrapfig}
\usepackage{tocloft}
\usepackage[toc]{multitoc}
\usepackage[hidelinks]{hyperref}
\raggedbottom
\extrafloats{2000}
\cftsetindents{section}{0em}{0em}
\cftsetindents{subsection}{0em}{0em}
\cftsetindents{subsubsection}{0.2em}{0em}
\setcounter{tocdepth}{1}
\renewcommand*{\multicolumntoc}{3}
\setlist[description]{nolistsep,listparindent=3pt,topsep=6pt}
\date{}
\begin{document}
\frontmatter
\tableofcontents
\mainmatter
\addcontentsline{toc}{chapter}{Introduction: Ascend to Adventure}
\markboth{INTRODUCTION: ASCEND TO ADVENTURE}{INTRODUCTION: ASCEND TO ADVENTURE}
\chapter*{Introduction: Ascend to Adventure}
The Infinite Staircase spirals endlessly in a vast, dreamlike expanse, twisting upward and downward to countless doors linked to every world and plane of existence. It is the route to the heart's desire, an extradimensional staircase that leads to all places. Beyond its myriad portals lie enchanting Feywild gardens, sunken cities ruled by ancient evils, and untold adventures.

\textit{Quests from the Infinite Staircase} collects six adventures from Dungeons \& Dragons history, updated for the current edition of the game. Beloved for their strange and wondrous locations, iconic monsters, and unforgettable encounters—including one indelible brush with a cute little bunnyoid on a stump in a crashed spaceship—these stories have entertained players since their Introductions. This anthology weaves together these memorable tales and adds a few new surprises.

The Infinite Staircase leads wayfarers from one story to the next. Legends state that by diligently searching the staircase's landings, a wistful traveler can find the door to the destination of their dreams.

Adventure awaits behind every door.

\begin{figure*}[b]
  \includegraphics[width=\linewidth]{images/../5e.tools/img/adventure/QftIS/000-00-001.golden-sphinx-door.png}
\end{figure*}

\section{Using This Book}
\textit{Quests from the Infinite Staircase} contains six classic D\&D adventures updated for the game's fifth edition. A cosmic thread pervades these tales, whose stories involve falling stars, archmages of multiversal renown, and the remnants of futuristic societies.

This book provides a means of traveling between these and other adventures: the Infinite Staircase. Unlike spells such as \textit{Plane Shift}, the staircase can be accessed by characters of any class or level. It merely requires happening on the right door.

The adventures in \textit{Quests from the Infinite Staircase} can be worked into ongoing campaigns. Additionally, these adventures complement those in other D\&D anthologies, such as Keys from the Golden Vault, Journeys through the Radiant Citadel, and Tales from the Yawning Portal—you can use the staircase to ferry the characters between the worlds in which those quests take place.

% Anchor in corner
\begin{tikzpicture}
  \node[anchor=north west, align=left, inner sep=0] at (current page.north east){  
  \includegraphics[height=\textheight]{images/../5e.tools/img/adventure/QftIS/001-00-003.infinite-staircase.png}
  };
\end{tikzpicture}



\subsection{About the Adventures}
Short summaries of the book's adventures appear below. For more on each adventure's origin and its place in D\&D history, see the "About the Original" sidebar in each adventure.

\subsubsection{The Lost City}
The Lost City is an introductory adventure written to accompany the \textit{Dungeons \& Dragons Basic Set}. Lost in a desert, the characters stumble on the ruins of Cynidicea and its peculiar factions. Released in 1982, the adventure marks the first appearance of Zargon, a one-eyed elder evil once entombed beneath the Cynidicean ziggurat. The updated adventure is designed for 1st-level characters.

\subsubsection{When a Star Falls}
When a Star Falls was the first adventure to bear a Union Jack flag on its cover in 1984. As with other adventures produced in the United Kingdom at the time, \textit{When a Star Falls} has been praised for its crisp writing and strong plot, the events of which surround a fallen star and the groups that covet it. The updated adventure is designed for 4th-level characters.

\subsubsection{Beyond the Crystal Cave}
Inspired by the works of William Shakespeare, Beyond the Crystal Cave takes characters to a timeless Feywild garden. A rarity for its time, the 1983 adventure emphasized noncombat encounters in an era of hack-and-slash dungeons. The updated adventure is designed for 6th-level characters.

\subsubsection{Pharaoh}
Pharaoh is the first adventure in the three-part \textit{Desert of Desolation} series and was published in 1982. In this adventure, the ghost of a long-dead pharaoh beseeches the characters to delve his dangerous tomb and set his cursed soul free. The updated adventure is designed for 7th-level characters.

\subsubsection{The Lost Caverns of Tsojcanth}
The Lost Caverns of Tsojcanth originally debuted under the title \textit{The Lost Caverns of Tsojconth} as a tournament adventure for the 1976 Winter Con game convention. The full version of the adventure, which takes place in the former lair of the archmage Iggwilv, was published in 1982. The updated adventure is designed for 9th-level characters.

\subsubsection{Expedition to the Barrier Peaks}
Expedition to the Barrier Peaks threw fantasy gamers for a loop with its unexpected science fiction elements at the Origins game convention in 1976, where it first appeared as a tournament scenario. Used to introduce D\&D players to the sci-fi game \textit{Metamorphosis Alpha}, the adventure—which involves a crashed spaceship, futuristic weaponry, and humorously malfunctioning robots—became an instant classic. The updated adventure is designed for 11th-level characters.

\subsection{Creating a Campaign}
The adventures in this book provide play across a broad range of levels. They can be strung together as a complete campaign using the Infinite Staircase to travel between them.

Starting with The Lost City, guide your players through the adventures in the order presented in this book. Each one provides sufficient XP that, on completing the adventure, the characters should be high enough level to advance to the next quest. The adventures also present narrative milestones for \textit{story-based level advancement} (see the \textit{Dungeon Master's Guide}).

After each adventure, have the characters find a door to the Infinite Staircase. There, they cross paths with a cosmic quest-giver, the noble genie Nafas (detailed in chapter 1), who sends them on their next journey. Because the staircase spans the multiverse, the adventures it can lead to are limitless, whether those tales come from this book or another.

\subsection{Running the Adventures}
To run each of these adventures, you need the fifth edition core rulebooks: the \textit{Player's Handbook}, \textit{Dungeon Master's Guide}, and \textit{Monster Manual}.

\begin{DndReadAloud}{}
Text that appears in a box like this is meant to be read aloud or paraphrased for the players when their characters first arrive at a location or under a specific circumstance, as described in the text.

\end{DndReadAloud}

When a creature's name appears in \textbf{bold} type, that's a visual cue pointing you to its stat block as a way of saying, "Hey, DM, you should get this creature's stat block ready. You're going to need it." If the stat block appears elsewhere, the text tells you so; otherwise, you can find the stat block in the \textit{Monster Manual}.

\textit{Spells} and \textit{equipment} mentioned in each of the adventures are described in the \textit{Player's Handbook}. \textit{Magic items} are described in the \textit{Dungeon Master's Guide}, unless the adventure's text directs you to an item's description elsewhere.

\chapter{The Infinite Staircase}
The Infinite Staircase ushers travelers to the destinations of their dreams. The impossible staircase winds in a surreal demiplane, ascending to every world, every plane of existence, and the adventure locations found among them. Though stories abound along its many steps, the staircase is the journey, not the destination—an extraordinary means of traversing the multiverse.

This chapter expands on what's written about the \textit{Infinite Staircase} in the \textit{Dungeon Master's Guide} and provides tools for using it to convey creatures from one fantastical location to the next.

\includegraphics[width=0.4\textwidth]{images/../5e.tools/img/adventure/QftIS/002-01-001.ch1-the-infinite-staircase.png}
\section{Staircase Features}
The following sections describe features of the Infinite Staircase that are likely to come into play.

\subsection{Day and Night}
The staircase doesn't experience day or night, existing in a perpetual state of twilight. Distant doors, twinkling like stars, provide dim light to the dreamlike realm, as does the staircase itself, which gives off an ambient glow in some stretches.

\subsection{Doors and Landings}
Endless doors along the staircase connect it to destinations throughout the multiverse. These doors function as \textit{planar portals} (see the \textit{Dungeon Master's Guide}) and are typically closed, though most are unlocked and don't require portal keys to enter.

Each door on the staircase occupies a landing with at least two staircases connected to it, leading to and from other landings and their respective doors. The doors and their landings are altered by planar magic, taking on characteristics of their destinations. Red light seeps from the seams of a gloomy, jet-black door to the Shadowfell set into a gothic archway; gargoyles perch on its landing, exuding roiling mist from their open mouths. A slippery ramp of deep-blue ice ascends to a frozen portal to the Paraelemental Plane of Ice, while a clammy bridge of basket fungus weaves toward a cavernous door to the Underdark.

Landings are frequent along the staircase but vary in distance from one another. Every landing contains at least one door. Travel between landings can take minutes, hours, or days.

\subsection{Gravity}
On the Infinite Staircase, gravity is a logical force—until it isn't. The direction of its effect is perpendicular to the steps of the nearest staircase. Staircases twist, bend, and invert, occasionally allowing travelers on separate branches to look up at each other or leap from one part of the staircase to another.

A creature that falls off the staircase plummets 1d10 × 50 feet before crashing on the lower steps or the nearest landing below; it takes damage from the fall and lands with the prone condition as normal.

An object dropped on the staircase behaves predictably at first, falling down the steps until it comes to a halt or goes over the edge, at which point it continues "downward" until it collides with the next closest section of the staircase, and so on. Precious items lost in this way sometimes appear on the nearest landing, mysteriously returned to their owners.

\subsection{Indestructible}
The Infinite Staircase is immune to all damage and can't have its form changed by spells like \textit{Passwall} or similar magic. No one knows who or what built the staircase—or continues to add to it. Its masons are never seen but often heard, their hammers ringing in the distance like rolling thunder.

\subsection{Movement on the Staircase}
The staircase includes a multitude of inclines, not just stairs. Ramps, hovering platforms, and clockwork conveyor belts can all be found along its endless construction.

A traveler can also move along the Infinite Staircase simply by thinking about moving, causing the incline to propel them toward their destination. Steps might roll beneath them like a magical escalator or flatten under them to create a ramp. The walking speed (in feet) of a creature that chooses to move in this way is equal to 10 times its Charisma modifier (minimum speed of 10).

\subsection{Weather}
The Infinite Staircase doesn't have seasons. Its ambient temperature remains constant, comparable to a cool autumn evening with a gentle breeze and breathable air. Storms and extreme weather effects, except those caused by the staircase's denizens and their lairs, are virtually nonexistent.

\section{The Iron Shadow}
The portals of the Infinite Staircase often lead to places that exhibit a penchant for creativity—sites where art, invention, philosophy, and magic flourish. Wherever there's a story to be told, the staircase might be found. As creativity wanes in these places, their doors and destinations change, and their landings on the staircase vanish over time. However, one force plays a more active role in their destruction: the Iron Shadow.

The Iron Shadow is a creeping malady that feeds on creativity in all its forms. A rootless, rolling cloud of gloom exhaled by the Shadowfell, the Iron Shadow darkens the doors of the staircase, sapping inspiration, novelty, and motivation from the planes and their inhabitants. It snuffs the spark of innovation and leaves its victims apathetic and dull. Wherever the shadow falls, creativity dies.

\subsection{Palls of Gloom}
Areas of the staircase plagued by the Iron Shadow appear cold and lifeless and are shrouded in palls of magical gloom. A creature that enters the gloom for the first time on a turn or starts its turn there must succeed on a DC 10 Constitution saving throw or gain 1 level of exhaustion. If this effect increases a creature's level of exhaustion above 5, the creature becomes a \textbf{shadow} under the Iron Shadow's control. Nothing short of a \textit{Wish} spell or divine intervention can reverse this transformation.

If an area of bright light created by a spell of 3rd level or higher overlaps with an area of gloom, that gloom disappears for 1 hour. Only a \textit{Wish} spell can permanently expunge a patch of the Iron Shadow from the staircase.

\subsection{Umbral Manifestations}
Aspects of the Iron Shadow occasionally take on physical forms, manifesting as memory webs, wraiths, or shadow dragons.

\section{Traversing the Staircase}
No spell or magic item is required to access the Infinite Staircase, nor do its travelers have to contend with enigmatic rulers who dictate when its doors may open or who may pass through them. Unlike the city of Sigil, the crossroads of the multiverse, gods are free to enter and leave the Infinite Staircase as they please, though many reject it in favor of more direct methods of planar travel. The staircase is neutral ground accessible to all who can find and traverse it.

The act of traversing the staircase isn't any more tiring than walking, and creatures can't become exhausted by traveling along it.

\subsection{Finding the Staircase}
Entrances to the Infinite Staircase usually appear as nondescript doors in dusty, forlorn places that don't get many visitors—or wherever is most convenient for adventure. On any given plane, multiple doors to the Infinite Staircase can exist, but known entrances are closely guarded secrets occasionally protected by devas, sphinxes, yugoloths, and other powerful creatures. Likewise, these creatures sometimes stand watch on the staircase's landings, barring nefarious or unworthy travelers from passing through doorways to sensitive realms.

The Infinite Staircase Entrances table presents several places where an entrance to the Infinite Staircase might exist.

\begin{DndTable}[header=Infinite Staircase Entrances]{cX}

d10 & Entrance\\
1 & Curtain of a long-forgotten stage\\
2 & Deepest door of an infamous lich's tomb\\
3 & Door to a giant's outhouse\\
4 & Fleshy aperture in a mind flayer's laboratory\\
5 & Grandfather clock in an inventor's workshop\\
6 & Locked wardrobe in a noble's mansion\\
7 & Magical painting in the gallery of an esteemed artist\\
8 & Permeable, sapient mirror in a wizard's tower\\
9 & Private room in a famous tavern\\
10 & Tree hollow in an enchanted forest\\
\end{DndTable}

\subsubsection{Gates of the Moon}
One fabled entrance to the staircase rests within the Gates of the Moon—a circular chamber within the shimmering palace Argentil, divine realm of the moon goddess Selûne on the plane of Ysgard. A spiral of silver steps, bathed in moonlight and entwined in dangling ivy, rises from the selenic hall, which some believe to be the base of the Infinite Staircase. Celestials, followers of Selûne, and lovers of the arts abound near the Gates of the Moon, which they regard as a sacred site.

\subsection{Locating Doors}
Traveling to another plane requires locating a door that corresponds to the desired destination. There's no reason to the doors' arrangement, but creatures on the staircase longing for a specific destination feel the tug of its magic pulling them toward where they want to go. Finding the right door is a matter of exploration and time: locating the correct one takes 4d12 hours of travel. If a creature has visited the door before, this time is halved.

Doors along the staircase can lead anywhere. The Infinite Staircase Doors table provides a sampling of doors and destinations that characters might discover while exploring the Infinite Staircase.

\begin{DndTable}[header=Infinite Staircase Doors]{cXX}

d100 & Door on the Staircase & Destination\\
01–03 & Webbed mouth of a pitch-black cavern & Cave in the Demonweb Pits, a spidery layer of the Abyss\\
04–06 & Rusty metal hatch forged from dulled weapons & Interior of a war machine in Acheron, a plane of infinite battle\\
07–09 & Trellis covered with vibrant, untamed foliage & Ancient grove on the plane of Arborea\\
10–12 & Freshly painted barn door & Bucolic homestead on the plane of Arcadia\\
13–15 & Tent flap made of animal hides & Upset stomach of an herbivorous behemoth in the Beastlands\\
16–18 & Meticulously carved half doors in two colors & Bustling gnomish factory on the plane of Bytopia\\
19–21 & Barred jail cell door that locks when closed & Unforgiving prison in the depths of Carceri\\
22–24 & Cascading waterfall that hums with song & Picturesque grotto on the plane of Elysium\\
25–27 & Knobless door fitted with a sliding peephole & Seedy, yugoloth-owned tavern on the plane of Gehenna\\
28–30 & Gray, unassuming door on a colorless landing & Crumbling temple to a forgotten god on the joyless plane of Hades\\
31–33 & Oblong door that changes material at random & Adamantine githzerai fortress on the chaotic plane of Limbo\\
34–36 & Ticking portal formed from hundreds of gears & Dusty archives in Mechanus, a plane of absolute law and order\\
37–39 & Stained-glass door on a landing of pure light & Gleaming temple at the foot of Mount Celestia\\
40–42 & Obsidian gate with magma-spewing gargoyles & Infernal courthouse in the Nine Hells of Baator\\
43–45 & Airtight door that screams when it's opened & Derelict windmill wracked by the howling gales of Pandemonium\\
46–48 & Stone double doors sized for a giant & Forge of a fire giant blacksmith on the rugged plane of Ysgard\\
49–58 & Sparkling door within a gemstone arch & The Radiant Citadel, a refuge on the Ethereal Plane\\
59–63 & Silvery threshold etched with stars & Eye socket of a dead god drifting through the Astral Plane\\
64–68 & Ornate, upright coffin that expels a hostile \textbf{swarm of bats} when opened & Empty crypt in the bowels of Castle Ravenloft, a fortress in the Shadowfell ruled by the vampire Count Strahd von Zarovich\\
69–71 & Dressed marble door featuring a cameo of a crowned lion mounted above its arch & Palace of Heart's Desire, seat of the archfey Zybilna in the Feywild domain of Prismeer\\
72–74 & Rotating bookshelf that reveals a hidden portal & Broom closet in a tower in Candlekeep, a fortress of tomes on Toril\\
75–77 & Gilded door marked with a family crest & Floating spire in Sharn, City of Towers, on the world of Eberron\\
78–80 & Portcullis gate covered with thorny black roses & Dargaard Keep, cursed castle of the death knight Soth on Krynn\\
81–83 & Circular door curtained with razorvine & The Ubiquitous Wayfarer, a many-portaled tavern in the city of Sigil\\
84–86 & Drafty portal at the top of a cloudy staircase & Court of a cloud giant duke on the Elemental Plane of Air\\
87–90 & Rough-hewn tunnel with a mining cart outside & Excavation site on the Elemental Plane of Earth\\
91–94 & Blisteringly hot brass door & Efreeti-run shop in the City of Brass on the Elemental Plane of Fire\\
95–98 & Fossilized shark jaws preceded by coral steps & Hut on the back of a dragon turtle on the Elemental Plane of Water\\
99–00 & Curtain of slime framed by writhing tentacles & Puddle of brine within the folds of a giant brain in the Far Realm\\
\end{DndTable}

\includegraphics[width=0.4\textwidth]{images/../5e.tools/img/adventure/QftIS/003-01-002.marching-modrons.png}
\subsection{Infinite Staircase Encounters}
The Infinite Staircase Encounters table presents random encounters and interactions adventurers might have while traversing the Infinite Staircase.

\begin{DndTable}[header=Infinite Staircase Encounters]{cX}

d100 & Encounter\\
01–07 & An \textbf{empyrean} ascends the staircase to a goodly realm. Flowers sprout from the steps behind it as it walks. The empyrean waves politely at the party.\\
08–14 & A section of the staircase ahead is covered in shadowy webbing. The hazard is the Iron Shadow manifesting as a \textbf{memory web}.\\
15–21 & Three \textbf{aarakocra} return to the Censer of Dreams (detailed later in this chapter) from a faraway plane. The friendly avian travelers invite the characters to accompany them back to the palace.\\
22–28 & Melvin Manyroads (chaotic good, gnome \textbf{mage}), a planar cartographer, stops to ask the characters if they've seen a particular door on the staircase.\\
29–35 & 1d4 + 1 bearded devils look for a door that leads behind enemy lines of the Blood War, the eternal conflict between demons and devils.\\
36–42 & A stubborn dwarf \textbf{commoner} bumps into the characters. He's been wandering the staircase for decades but won't admit he's lost.\\
43–49 & A masked \textbf{deva} loyal to Selûne offers to guide the characters to a door of their choosing. The angel plays melodies on its harp during the journey.\\
50–56 & An aspect of the Iron Shadow, a \textbf{young red shadow dragon}, rises from a staircase shrouded in dusk to confront the characters.\\
57–63 & Inky tendrils writhe from a fleshy door to the Far Realm. Two githyanki warriors struggle to hold the door shut. They call out to the party for aid.\\
64–70 & A lost unit of the Great Modron March consisting of 1d8 duodrones and 1d6 tridrones proceeds down the staircase, cataloging its every door.\\
71–77 & The staircase ahead disappears into a 60-foot sphere of gloom created by the Iron Shadow (see the "Palls of Gloom" section).\\
78–84 & A \textbf{mastiff} that wandered through a portal trots up to the party. It drops a stick in front of one of the characters and wags its tail eagerly.\\
85–91 & A \textbf{merrow} happens on the staircase, flooding it with a deluge of water from an aquatic plane.\\
92–98 & A nearby door to the Abyss flings open, releasing a hostile \textbf{swarm of gibberlings}.\\
99–00 & A staircase of pure light materializes before one of the characters, beckoning them to the destination of their heart's desire. If the path is refused, it vanishes and never appears to them again.\\
\end{DndTable}

\section{Nafas}
Once there was everything, and once there was nothing but air. Then came a breath—a collective sigh exhaled by every plane and every world in one ephemeral gust. Multiversal winds blew into the Infinite Staircase through its doors, coalescing into Nafas (na-FASS), a noble genie infused with cosmic potential.

\textbf{Nafas} is detailed in the following sections and in appendix B.

\subsection{Breath of the Multiverse}
Wisps of twinkling multiversal dust trail Nafas wherever he goes, evidence of the djinni's deep connection to the planes that bore him. Calm, clement, and dressed in billowing fabrics, he drifts along the staircase like a zephyr, admiring its portals and the curious creatures that emerge from them.

Nafas isn't the creator of the Infinite Staircase or its ruler. He doesn't control the staircase or its doors, and he holds no sway over their destinations or who may enter through them. The djinni is a product of the multiverse, a distant observer bound to the realm by the circumstances of his creation.

An unwritten law dictates that there must always be a noble genie to watch over the Infinite Staircase. As a result, Nafas can never leave it. Any attempt to forcibly remove him from the Infinite Staircase, whether by magical means or otherwise, fails.

\subsection{Cosmic Storyteller}
Nafas is a gifted storyteller and a gracious host. He regales visitors to his palace (detailed later in this chapter) with tales of longing that could soften steel and legends of triumph that chase grief from sorrow-laden hearts. Truth threads each captivating story, woven by Nafas from the whims of creatures throughout the multiverse. Carried on the winds, their spoken desires converge in his splendid halls.

Nafas has as many years behind him as the staircase has doors. Seemingly immortal, the djinni has collected and recited countless tales. Some fables repeated throughout the planes are inspired by these true accounts—or corruptions of stories still yet to be resolved—told to travelers who crossed paths with Nafas.

\subsection{Noble Genie}
\textbf{Nafas} is a noble genie, a rare type of genie endowed with the power to reshape reality. A word uttered by Nafas can transform a pauper into a prince, return the dead to life, or bring a meteor's might crashing down on his enemies. Unlike some of his more haughty, aristocratic counterparts that rule over vast demesnes of the Elemental Planes, Nafas's noble title isn't a result of bloodline or inheritance; it's a designation bestowed on him by the multiverse.

As a noble genie, Nafas can grant wishes, a power he doesn't take lightly. The djinni reserves his wish-granting magic for those he deems worthy. When Nafas does grant a wish—the greatest gift the djinni can give—he doesn't twist the benefactor's words or maliciously interpret their intentions, but his circumstances prevent him from freely conferring such gifts to deserving recipients outside the staircase.

Moved by the whims of yearning creatures on other planes but trapped within the Infinite Staircase, Nafas relies on adventurers to fulfill the wishes of those beyond the staircase's infinite thresholds. The djinni knows every door on the staircase and its destination, along with the origin of any wish that escapes through a door's seams.

\subsubsection{Seeking Wishes}
Characters might seek out the djinni to end a lasting effect, resurrect a fallen character, or solve an unsolvable problem. The following are some obstacles that Nafas can wish away:


\begin{itemize}
\item Destroying or unmaking a magic item
\item Learning the true name of a powerful enemy, such as an angel, an archfey, a demon, or a devil
\item Restoring a creature to life
\item Undoing a debilitating effect, like being transformed into a tree by a \textbf{barkburr}
\end{itemize}

\subsection{Wishes from the Infinite Staircase}
Nafas might act as a patron who hires the characters to undertake the adventures in this book. Each adventure includes a "Using the Infinite Staircase" section that frames its central plot as a wish heard by Nafas.

The djinni or one of his messengers—such as a friendly \textbf{dust mephit} formed of glimmering stardust—might appear to the characters along the staircase and invite them to seek rest at his palace. There, Nafas imparts the faraway plea he has heard and escorts the characters to the door where it originated. After the adventure, characters can return to the Infinite Staircase through the same door.

In exchange for fulfilling the wish of a creature beyond his reach, Nafas grants the characters a reward of their own. This reward might take the form of one or more of the following:


\begin{description}
\item[Magical Gift.]
Nafas bestows a supernatural charm (see the \textit{Dungeon Master's Guide}) on the party.
\item[Tale.]
Nafas recites a story or poem that reveals information pertinent to a higher ambition or plane-spanning quest pursued by the characters.
\item[Wish.]
Nafas grants the characters one to three castings of the \textit{Wish} spell on their behalf. The djinni might require the characters to complete multiple adventures before allotting such a handsome gift.
\end{description}

\includegraphics[width=0.4\textwidth]{images/../5e.tools/img/adventure/QftIS/004-01-003.the-genie-nafas.png}
\section{Censer of Dreams}
Nafas beckons weary travelers along the staircase to the Censer of Dreams, a bastion of respite in the continuous expanse. The Censer of Dreams is Nafas's lair. Regional effects for the palace are detailed in appendix B.

Created by Nafas to shelter visitors from other planes, the Censer of Dreams is a testament to the noble genie's generosity and might. Many who stumble on the Infinite Staircase glimpse the djinni's luxurious residence at some point during their journeys. A nexus held aloft by dozens of branching staircases, the palace envelops its visitors in unimaginable comforts.

The Censer of Dreams is a wonder of aeolian architecture shaped by a being of pure wind. Massive, shutterless windows line the bulbous fortress, allowing the djinni and other flying creatures to flit through its airy halls unimpeded. Hewn from smooth marble and perfumed with floral fragrances, its grand chambers hum with the pleasing tones of magical chimes hung from lofty ceilings. Windcatchers rise from the top of the structure to capture the gusts that buffet the palace and the staircases around it at all hours. Suffused with the djinni's magic, these winds batter would-be invaders, sending them over the staircase's edge.

\subsection{Censer of Dreams Features}
The Censer of Dreams has the following features:


\begin{description}
\item[Unseen Stewards.]
A staff of one hundred invisible zephyrs tends to the palace and its guests. Each of these stewards functions as if created by the \textit{Unseen Servant} spell; however, the zephyrs are permanent fixtures that last until destroyed, and they can move anywhere in the palace without dissipating. Only Nafas and creatures he designates can command the zephyrs. A destroyed zephyr re-forms in a random location within the palace after 1d4 days.
\item[Winds.]
Winds perpetually swirl around the Censer of Dreams and the staircases that lead to it. Creatures on the Infinite Staircase can find the palace by following the direction of the wind.
\item[Wishes.]
Fleeting voices can be heard throughout the palace, whispering the desires of creatures throughout the multiverse.
\end{description}

\subsection{Palace Visitors}
Travelers from every world are welcome within the Censer of Dreams. Nafas values all perspectives and rarely refuses entrants, believing that everyone has a story worth telling. However, troublemakers who offend the djinni or threaten the safety of other visitors are ousted from the palace.

Some genies, especially those from the Elemental Plane of Air, pilgrimage to the Censer of Dreams to hear stories from Nafas. Others hope to join his court and one day earn the title of noble genie themselves.

\subsection{Noteworthy Sites}
Amid an endless sea of steps, the Censer of Dreams embodies the boundless potential of a wish granted. Its opulent chambers boast heavenly gardens visited by Celestials, sprawling courtyards paved with silver bricks, and the secret desires of the multiverse.

\subsubsection{Foyer}
Nafas receives visitors in the foyer—a wide, windowed reception hall that looks out on the Infinite Staircase. As one enters the foyer, they're greeted by a roaring ramp of tangible wind that curls up to the rest of the palace. The updraft supports even the heaviest of creatures, which can easily amble up the windy incline to the areas above. Plush pillows, warm fireplaces, and delectable treats number among the amenities available in the hall.

\subsubsection{Den of Chronicles}
Tales come alive in the Den of Chronicles, a storyteller's sanctuary. In the carpeted hall, Nafas weaves a never-ending tapestry of tales from the wishes he has heard.

An immense, jeweled censer hangs from the ceiling on a golden chain. Whenever a story begins, sweet-smelling incense wafts from the receptacle and spreads throughout the palace, notifying its visitors that a new chronicle is about to unfold. As the narrator—whether Nafas or another raconteur—speaks, the fumes take the forms of subjects within the story, dancing to each turn of phrase. Meanwhile, animated kettles distribute steaming cups of hot tea to avid listeners seated within the hall.

Nafas delights in hearing the stories of his visitors, especially tales of adventures beyond the Infinite Staircase. Easily moved, the djinni disarms audiences with his booming laugh and profound sobs.

\subsubsection{Guest Quarters}
The palace's accommodations are individualized oases sprinkled with genie magic. Illusions within each room tailor its contents to the preferences of its guests, transforming it into a personal utopia. A traveler who hails from a snowy tundra might melt at the sight of a rustic mountain lodge, complete with a crackling hearth, a bearskin rug, and the savory scent of dwarven ale. The next room over might appear as a beachside bungalow, a peaceful elven glade, or a wizard's scriptorium. Every room is someone's paradise.

\includegraphics[width=0.4\textwidth]{images/../5e.tools/img/adventure/QftIS/005-01-004.censer-of-dreams.png}
\subsubsection{Well of Destiny}
The voices of the multiverse cry out in the Well of Destiny. A sanctum located within the silvery dome that rises from the center of the palace, the Well of Destiny is a magnet for the words of hopeful souls manifesting their desires to the winds. When he isn't entertaining guests, Nafas can usually be found here, seated on his throne: a weightless, high-backed chair sculpted from the clouds that pervade the misty chamber. The djinni reflects on each wish he hears, contemplating the stories behind them.

Each time a wish is spoken somewhere in the multiverse, it echoes in the Well of Destiny, amplified by the Infinite Staircase's ambient magic. The Well of Destiny Wishes table includes examples of wishes that might be heard within the hall or the staircases beyond, inspiring characters to adventure or giving them a reason to seek out Nafas.

\begin{DndTable}[header=Well of Destiny Wishes]{cX}

d6 & Wish\\
1 & "I wish I could find my way back home."\\
2 & "I wish to find a cure for this magical sickness."\\
3 & "I wish we could defeat the evil that rules over us."\\
4 & "I wish I could destroy that artifact once and for all."\\
5 & "I wish for revenge on the one who wronged me."\\
6 & "I wish I were free from this cursed prison."\\
\end{DndTable}

\section{Staircase Adventures}
The adventures in this book take the characters to wondrous realms beyond the Infinite Staircase, but some quests might take place on the staircase itself. The Infinite Staircase Adventures table offers suggestions for stories involving the extradimensional realm.

\begin{DndTable}[header=Infinite Staircase Adventures]{cX}

d4 & Adventure\\
1 & The Iron Shadow spreads to a character's home world, leaving it a gray, apathetic husk. A survivor asks the characters to repair the damage.\\
2 & For reasons unknown, the masons construct a gloomy staircase to an evil, nightmarish realm. Unsupervised, its doors release horrors that wreak havoc on other worlds. \textbf{Nafas} hires the characters to seal the portals.\\
3 & A powerful \textbf{efreeti} from the City of Brass sets out to conquer numerous worlds through the doors of the Infinite Staircase. Nafas charges the characters with foiling the efreeti's scheme.\\
4 & The staircase's doors lose their magic, preventing passage into or out of the planar realm. A traveler along the staircase beseeches the party to solve the mystery so the traveler can return home.\\
\end{DndTable}

\chapter{The Lost City}
Lost in the desert, the last great remnant of a once-prosperous civilization rises from the dunes. Within this ziggurat await untold riches and the sunken city of Cynidicea—along with unexpected challenges, deadly threats, and the fallen kingdom's desperate inhabitants.

Forsaken by the outside world, the Cynidiceans don elaborate masks and vie against each other in factions devoted to ancient gods. Meanwhile, an ageless evil of unknown origin lurks in the bowels of the dilapidated ziggurat. Called Zargon the Returner, the tentacled, one-eyed creature preys on the Cynidiceans even as they worship it as a god. To reclaim their former kingdom, the Cynidiceans must oust the eldritch abomination and its cultists.

"The Lost City" is designed for four to six 1st-level characters. If, on completing the adventure, you wish to extend it further, consult the "Extending the Adventure" section at the end of this chapter.

\includegraphics[width=0.4\textwidth]{images/../5e.tools/img/adventure/QftIS/006-02-001.ch2-the-lost-city.png}
\section{Background}
Centuries ago, Cynidicea was the capital of a prosperous kingdom. Through advancements in magic and technology, the Cynidiceans reclaimed land from the desert and transformed their city into a paradise. Cynidicea peaked with the reign of King Alendria and Queen Zanobis. After the deaths of these last great monarchs, the Cynidiceans honored them by erecting a massive step pyramid that served as an important city hub for years.

The fall of Cynidicea began when workers digging under the ziggurat discovered a gnarled, jet-black horn jutting from a mysterious obelisk. Unknown to the workers, the horn belonged to Zargon, a timeless evil entombed in the obelisk eons ago. The moment the Cynidiceans pried the horn from the monolith, their city was doomed. Materializing from its severed protuberance, Zargon devoured the workers and slew the garrisons that followed, drowning them in torrents of slime.

Rather than face the grim reality confronting them, some Cynidiceans began to view their fellow citizens as sacrifices to a new and voracious god, gradually forming a strange cult to Zargon. Choked by Zargon's slime and assailed by opportunistic invaders, the Cynidiceans fled underground.

Led by the Cult of Zargon, the Cynidiceans began to rebuild, constructing a miserable reflection of their former kingdom in the darkness. Above, drifting sands covered the city, and Cynidicea was lost in the vastness of the desert.

\subsection{Using the Infinite Staircase}
If you're using Nafas as a patron, he summons the characters to the Censer of Dreams (detailed in chapter 1), where he recounts the following wish:

\begin{DndReadAloud}{}
"I have heard the distressed prayers of a fallen people ensnared by an ageless evil. They whisper the names of long-forgotten gods in a city beneath the desert, hoping to rekindle their kingdom's lost glory. I must know whether they are beyond saving. Return to me with their story."

\end{DndReadAloud}

Nafas then teleports the characters to a door along the Infinite Staircase that opens into a stone arch in a vast desert. The arch isn't far from the ruins of Cynidicea. After the adventure, the characters can return to the staircase through the same portal.

\includegraphics[width=0.4\textwidth]{images/../5e.tools/img/adventure/QftIS/007-02-002.lost-city-door.png}
\subsection{Adventure Hooks}
If you're not using Nafas as a group patron, consider the following ways to involve the characters in this adventure:


\begin{description}
\item[Anthropological Expedition.]
A recent desert excavation has unearthed evidence of a bygone civilization. A scholar hires the characters to explore the Cynidicean ruins and report their findings.
\item[Desert Caravan.]
The characters joined a desert caravan, but a sandstorm separated the party from the rest of the caravan. After wandering the desert for days and running out of food and water, the characters discover an ancient ziggurat.
\end{description}

\subsection{Setting the Adventure}
Lost to time and sand, Cynidicea is barely a memory to the lands where it once prospered. When deciding where to place Cynidicea, consider the following suggestions:


\begin{description}
\item[Dragonlance.]
On the world of Ansalon, Southern Ergoth and the Northern Wastesboth contain arid stretches where Cynidicea might have existed.
\item[Forgotten Realms.]
The ruins of Cynidicea might lie among the dunes of Anauroch on the continent of Faerûn.
\item[Mystara.]
On the Known World, the kingdom of Cynidicea could have occupied an isolated region of the Alasiyian Desert.
\end{description}

\includegraphics[width=0.4\textwidth]{images/../5e.tools/img/adventure/QftIS/008-02-003.lost-city-original.png}
\begin{DndReadAloud}{About the Original}
Published in 1982, \textit{The Lost City} was written by Tom Moldvay, who contributed to the \textit{Dungeons \& Dragons Basic Set}. Following in the footsteps of other "B-series" adventures, \textit{The Lost City} served as both an adventure and a learning tool for new Dungeon Masters. While the adventure detailed the upper levels of the Cynidicean ziggurat, DMs could test their mettle by expanding the dungeon to include its lower levels, for which the module provided only rudimentary features. This updated version of the adventure forgoes this exercise, providing an optional sixth level below the ziggurat.

\end{DndReadAloud}

\section{Adventure Summary}
"The Lost City" is a dungeon delve with a dose of intrigue. While in the desert, the characters discover a ruined ziggurat peeking over the dunes. Inside the monument, they encounter three bickering factions of Cynidiceans who seek to reinstate their once-proud kingdom. Characters can align with one or more factions, reaping the rewards of membership at the risk of ostracizing rival groups.

As the characters explore the ziggurat, they slowly uncover clues of the tragedy that befell Cynidicea and its people. While braving clever traps, evil cultists, and restless Undead, the characters might learn that Zargon, the one-eyed creature that caused Cynidicea's fall, still dwells beneath the monument. The party might even discover an entrance to the surviving Cynidiceans' underground city (see the "Extending the Adventure" section at the end of this adventure).

The adventure concludes when the characters escape the ziggurat, either by returning to the desert or finding a door to the Infinite Staircase placed in a location of your choice.

\subsection{Character Advancement}
If you want to use story-based level advancement, the characters receive experience points for achieving the following milestones rather than defeating monsters:


\begin{description}
\item[Encountering a Faction.]
After the characters first encounter the leader of one of Cynidicea's three factions (detailed in the following section), everyone in the party who survives the encounter gains 1 level. This advancement occurs only once, even if the characters later encounter another faction leader.
\item[Entering the Tomb Complex.]
When the characters enter the fourth tier of the ziggurat for the first time, the characters gain 1 level.
\item[Escaping the Ziggurat.]
The characters gain 1 level on completing the adventure, either by escaping the ziggurat or further exploring its depths (see the "Extending the Adventure" section).
\end{description}

If you follow this method, the characters should reach 4th level by the adventure's conclusion.

\section{The Cynidiceans}
\includegraphics[width=0.4\textwidth]{images/../5e.tools/img/adventure/QftIS/009-02-004.cynidicean-masks.png}
Since their kingdom's fall, generations of Cynidiceans have eked out a subterranean existence. Few have ever glimpsed the sun, and some deny the outside world exists at all. Cynidiceans represent a spectrum of genders and species, though most are human.

Cynidiceans embrace lifestyles of twisted decadence to distract themselves from their abject misery. Every citizen wears a stylized mask, usually the face of an animal or other creature that embodies the wearer's personality. Some masks are carved from wood or bone, while others are forged from precious metals. Cynidiceans dress in once-lavish clothing that's now faded and frayed. Elaborate jewelry—fashioned from scavenged beads, feathers, and other baubles—complements their torn robes and haunting visages, transforming each citizen into a walking mockery of their once-respected people.

After years of isolation, some Cynidiceans have abandoned their identities and adopted strange behaviors based on the faces they wear. A Cynidicean who wears a feline mask might meow during conversation or purr when pleased, while one who dons a bee mask might devote themself to making honey, buzzing noisily wherever they go.

\subsection{Cult of Zargon}
Rising from the fetid slime that drowned their former kingdom, the cultists of Zargon revere the horned abomination as a god. These fanatics gather in a black stone temple in the underground city, where they prepare living sacrifices to Zargon and sow fear among the remaining Cynidiceans. Periodically, the cultists venture up to the tunnels beneath the ziggurat to feed Zargon and kneel to it in worship, clashing with any who would deny their overlord tribute.

Zargonites each wear a horned, golden mask with four tentacles sprouting from its chin, representing the Returner's visage.

\subsection{Factions of Cynidicea}
Three factions of Cynidiceans seek to restore the worship of the old gods Gorm, Usamigaras, and Madarua. Opposed by the cultists of Zargon, each faction hopes to stop the slow death of their society and regain Cynidicea's past glory. However, cooperation between the groups is rare and frequently devolves into altercations.

\subsubsection{Guardians of Gorm}
The Guardians of Gorm are followers of Gorm, a god of justice, storms, and war. Members of this tight-knit fellowship value bravery, honesty, and justice tempered by mercy. They view storms as holy and believe thunder to be Gorm's reboant voice.

Guardians wear brass masks depicting the stern, bearded face of Gorm and deep blue tunics. Under each guardian's armor, their right shoulder bears a cobalt-hued tattoo of a lightning bolt.

\subsubsection{Mages of Usamigaras}
Like their mischievous god Usamigaras—a deity of magic, messengers, and thieves—the Mages of Usamigaras are outwardly kind but duplicitous in nature. They prize trickery above sincerity and never speak their true intentions, advancing their goals through deceit and treachery.

Mages wear silver masks of a smiling, round-faced boy—the face of Usamigaras—along with silvery robes that draw attention away from their hidden wands. Each mage's right palm bears a tattoo of a five-pointed star.

\subsubsection{Warriors of Madarua}
The ancient god Madarua presides over birth, death, and the shifting seasons. Courage, adaptability, and wit are paramount to the Warriors of Madarua, and these traits have served them well—within the underground city, this faction is outnumbered only by the Cult of Zargon.

Madarua's warriors wear bronze masks of her expressionless face. They dress in bronze armor and favor spears, and each member bears a sickle-shaped tattoo on their left wrist.

\section{Buried Ziggurat}
Regardless of what brings the characters to the desert, the adventure begins when they happen on a dilapidated ziggurat. This impressive monument honors the last and greatest rulers of Cynidicea. Though the ziggurat is usually buried in the sands, a sandstorm has uncovered its upper tiers.

When the characters arrive at the ziggurat, read or paraphrase the following text:

\begin{DndReadAloud}{}
Your trek through the desert halts at the ruins of a lost civilization. Mountains of sand bury most of the weathered city, but a few stubborn landmarks protrude from the mounds. Arid winds and the passage of time have scoured their faces smooth.

A windswept ziggurat towers over the other ruins. Five stepped tiers are visible, though the lowest barely peeks above the sand. In all, the ancient structure rises over one hundred feet into the sky.

Three thirty-foot-tall statues stand atop the monument: a strong, bearded man holding a pair of scales in one hand and a lightning bolt in the other; a winged child entwined in a pair of snakes; and a stoic woman grasping a sword and a sheaf of wheat.

On the south side of the ziggurat, a ramp with stairs ascends to the monument's highest tier.

\end{DndReadAloud}

The three statues depict Cynidicea's ancient gods: Gorm, Usamigaras, and Madarua. A character who studies the statues and succeeds on a DC 15 Intelligence (Religion) check recognizes the lost gods—who might still have sparse followers in nearby regions—and their portfolios.

The ziggurat's main entrance is buried beneath the sands and sealed from the inside. Characters who search for another entrance find the once-hidden secret door to area B1 near the top of the ziggurat.

\subsection{Ziggurat Features}
The ziggurat is made of large blocks of smooth sandstone. Unless otherwise noted, the ziggurat has the following features:


\begin{description}
\item[Ceilings.]
Corridors have 10-foot-high ceilings, and rooms have 15-foot-high ceilings.
\item[Doors and Secret Doors.]
Interior doors are heavy stone slabs that can be pushed open from either side. Secret doors blend in with their surroundings. A character who examines a wall for a secret door must succeed on a DC 15 Wisdom (Perception) check to notice it. Most doors are unlocked and quietly close on their own unless held or jammed open.
\item[Lighting.]
The ziggurat's chambers aren't illuminated. Area descriptions assume the characters have a light source or other means of seeing in the dark. Humanoids encountered in the ziggurat have darkvision out to a range of 60 feet, as they've adapted to life underground.
\end{description}

\subsection{Random Encounters}
For each hour the party spends in the ruins, roll a d6. On a roll of 1, an encounter occurs. To determine what the characters encounter, roll on the Random Ziggurat Encounters table.

Several of the encounters below are with specific Cynidiceans within the ziggurat and aren't likely to happen more than once unless the characters seek them out. If you roll on the table and get the same result as a previous roll, feel free to reroll or choose a different encounter.

\begin{DndTable}[header=Random Ziggurat Encounters]{cX}

d12 & Encounter\\
1 & 1 \textbf{rust monster}\\
2 & 2 giant centipedes\\
3 & 1 \textbf{owlbear}\\
4 & 1d4 hobgoblins search for treasure.\\
5 & A \textbf{commoner} in a Cynidicean wolf mask believes he's a werewolf. When he sees the party, he drops on all fours, howls, and sniffs each character. He's friendly toward the characters but afraid of silver.\\
6 & Three priests in fiendish Cynidicean masks approach the party with censers. Convinced the characters are possessed by demons, the priests assail the characters with chants and foul-smelling smoke before continuing on their way.\\
7 & Two guards in bright yellow robes and Cynidicean masks of scowling warriors kneel before one of the characters, addressing them as King Alendria or Queen Zanobis. They follow the character, singing songs of praise and defending them from peril.\\
8 & A number of commoners equal to the number of characters appear, wearing masks eerily similar to the characters' faces. These Cynidiceans mimic the characters as best they can for 10 minutes, even if it costs the Cynidiceans their lives.\\
9 & 1d4 guardians of Gorm try to escort the characters to their leader in area B12.\\
10 & 1d4 mages of Usamigaras try to escort the characters to their leader in area B18.\\
11 & 1d4 warriors of Madarua try to escort the characters to their leader in area B25.\\
12 & 1d6 cultists of Zargon declare that the characters must be sacrificed to the Returner and attempt to capture them alive.\\
\end{DndTable}

\subsection{Ziggurat Locations, Tiers 1–3}
The top three tiers of the ziggurat once served as living and ceremonial spaces for Cynidicean priests. The following locations are keyed to map 2.1.

\includegraphics[width=0.4\textwidth]{images/../5e.tools/img/adventure/QftIS/010-map-2.01-buried-ziggurat-1-3.png}
\includegraphics[width=0.4\textwidth]{images/../5e.tools/img/adventure/QftIS/011-map-2.01-buried-ziggurat-1-3-player.png}
\subsubsection{B1: Statue Room}
When the characters ascend the monument, read or paraphrase the following text:

\begin{DndReadAloud}{}
Near the ziggurat's peak, a once-secret entrance stands open, its stone slab held ajar by the decaying body of a hobgoblin. A crossbow bolt protrudes from the corpse's chest.

\end{DndReadAloud}

The corpse is that of an explorer who reached the ziggurat several weeks ago. After he triggered a deadly crossbow trap in the corridor, his fellow expeditioners stripped him of equipment and proceeded into the monument. A character who examines the area and succeeds on a DC 13 Intelligence (Investigation) check spots a tripped pressure plate on the floor and an empty crossbow in the wall.

When the characters enter the chamber, read or paraphrase the following text:

\begin{DndReadAloud}{}
The highest tier of the step pyramid holds a musty forty-foot-square chamber. Dust and sand coat its floor. Three bronze cylinders, each inset with a small door, span floor to ceiling in the middle of the room.

\end{DndReadAloud}

Footprints of several humanoids are visible in the dust, stopping in front of the center cylinder.

Shortly after all characters enter the chamber, the corpse (or anything propping the door open) crumbles and the door automatically shuts, sealing characters in the lightless chamber. Once closed, the door can be pushed open only from outside. It can't be opened from inside, and it has no lock to pick. The characters must find a new way out of the ziggurat.

\subparagraph{Gas Trap}
When the door closes, the area begins to fill with an odorless, dizzying gas. Characters who have a passive Wisdom (Perception) score of 13 or higher hear a steady hiss as the gas seeps through tiny holes in the wall. The gas immediately fills the room and dissipates after 1 minute. While the room is filled with gas, each creature that ends its turn in the room takes 1 poison damage. Creatures reduced to 0 hit points by the gas have the unconscious condition but are stable rather than dying.

If all characters succumb to the gas, 1d4 guardians of Gorm eventually discover them and deliver them to the faction's leader in area B12.

\subparagraph{Cylinders}
The three cylinders in this area are the hollow bases of the statues atop the ziggurat: Gorm to the west, Madarua to the east, and Usamigaras in the center.

Three crossbow traps are hidden in the south wall, one aimed at each cylinder. A character who searches the area for traps and succeeds on a DC 15 Wisdom (Perception) check spots the crossbows. A character has advantage on this check if they spotted the crossbow that killed the hobgoblin. Removing the bolts disarms the traps; the central crossbow is already empty.

When a creature opens the door on the east or west cylinder, a crossbow trap shoots a bolt at that creature. The target must succeed on a DC 16 Dexterity saving throw or take 3 (1d6) piercing damage. The center cylinder's trap has already been triggered, and a character who examines the floor notices dried blood flecking the dust there.

The base of each statue contains a ladder down to the storeroom (area B2) and up into the statue above, which is also hollow. Each statue contains a speaking tube, as well as a series of rusty levers that can be used to reposition the statue's head and limbs.

\subsubsection{B2: Machinery Storeroom}
\begin{DndReadAloud}{}
Three ladders descend from the cylinders into this storeroom. The feet of the ladders are lit by three skittering giant beetles with glowing glands beneath their carapaces. Against the south wall is a small forge and a workbench scattered with aged machinery.

\end{DndReadAloud}

Three giant fire beetles skitter about this storeroom, hunting for food. If the beetles detect the characters, the beetles hiss and attack.

\subparagraph{Treasure}
This room contains spare parts for the machinery and traps in area B1. Strewn throughout the room are twenty Crossbow Bolt Case, six Oil (flask), and a set of \textit{smith's tools} near a small foundry.

\subsubsection{B3: Secret Room}
This room can be accessed via two secret doors in the hallways bordering it. When the characters enter, read or paraphrase the following text:

\begin{DndReadAloud}{}
A shriveled humanoid corpse in blue robes and a brass mask lies in the center of this stale-smelling room. The corpse's stiff, wrinkled hand stretches toward a toppled leather pouch just out of reach. Glittering, blue gemstones have spilled from the bag onto the dusty floor.

\end{DndReadAloud}

A dishonest member of the Guardians of Gorm stole riches from the faction's treasure room (area B5) and hid them in this chamber. Before the guardian's last visit, a pack of stirges made the room their lair; when he returned, they drank him dry. A character who inspects the body and succeeds on a DC 13 Wisdom (Medicine) check learns the corpse is fresh but bereft of blood.

\subparagraph{Sleeping Stirges}
Four stirges, their bellies engorged with the corpse's blood, hang asleep from the ceiling. Bright light or any noise louder than a whisper wakes the stirges, causing them to attack. A character attempting to retrieve the gems must succeed on a DC 15 Dexterity (Stealth) check to avoid waking the stirges.

\subparagraph{Treasure}
Three blue quartz gems worth 25 gp each are scattered across the floor. The leather pouch contains another five gems of the same value. The dead Cynidicean's mask, which depicts the bearded face of Gorm, is made of brass and worth 10 gp.

\subsubsection{B4: Abandoned Priest's Quarters}
\includegraphics[width=0.4\textwidth]{images/../5e.tools/img/adventure/QftIS/012-02-005.gorm-mask.png}
\begin{DndReadAloud}{}
This room is sparsely furnished with decrepit furniture caked in dust: a soiled bed, a chest, and a writing desk with a wooden chair. A dead hobgoblin is sprawled on the floor, her left arm swollen and discolored.

\end{DndReadAloud}

This room was once the quarters of a priest of Gorm. On the desk lies a wooden holy symbol shaped like a lightning bolt. The hobgoblin is an explorer who died weeks ago after the bees in area B5 swarmed her. A character who examines the body and succeeds on a DC 13 Wisdom (Medicine) check determines the hobgoblin was stung by hundreds of insects.

\subparagraph{Treasure}
The hobgoblin carries a purse with 8 gp and 27 sp. The chest contains only rotting holy tomes that crumble when opened.

\subsubsection{B5: Treasure Room of Gorm}
\begin{DndReadAloud}{}
In the center of this room stands a ten-foot-tall cage. Within the cage, a pile of coins and gems sparkles beneath a gigantic, droning beehive.

\end{DndReadAloud}

The locked cage holds what little treasure the Guardians of Gorm have accumulated. As an action, a character can pick the lock with a successful DC 15 Dexterity (Sleight of Hand) check using thieves' tools. The key is held by Kanadius, the faction's leader (see area B12).

Two swarms of bees nest within the hive; use the swarm of insects (wasps) stat block to represent each of them. The bees are friendly toward the Guardians of Gorm but hostile to all others. If a character comes within 5 feet of the cage or disturbs its contents without first donning a mask of Gorm, the swarms erupt from the hive and attack.

\subparagraph{Harvesting Honey}
A character who inspects the hive and succeeds on a DC 11 Intelligence (Nature) check knows these bees create a special honey with healing properties. Once the swarms are dealt with, characters can safely harvest 1d4 doses of healing honey from the hive. A creature that consumes one dose of this honey as an action regains 1d4 + 1 hit points.

\subparagraph{Treasure}
The treasure in the cage includes 400 sp, 100 gp, two malachite gemstones worth 25 gp each, and a bejeweled necklace worth 150 gp.

\subsubsection{B6: Storage Room}
\begin{DndReadAloud}{}
This stuffy room contains only bales of moth-eaten cloth and crates of long-spoiled food.

\end{DndReadAloud}

The door to this storage room has remained closed for decades. The room contains nothing of value.

\subsubsection{B7: Fireworks Storeroom}
\begin{DndReadAloud}{}
This bone-dry storeroom reeks of sulfur. Six crates are stacked on a frayed rug in the center of the room. Atop the crates, two tiny creatures chat merrily, each barely a foot tall with delicate wings and sparkling hair. A third creature dives through an open crate, occasionally surfacing to spit out a mouthful of sawdust.

\end{DndReadAloud}

Three pixies—Wenly, Weenly, and Wally—are investigating the crates in this storeroom. The pixies speak Common and Sylvan. If approached, the friendly pixies introduce themselves and flit about the characters. The pixies inquire eagerly about the characters' outfits, equipment, and travels, but nothing holds the pixies' attention for long.

The pixies entered the ziggurat two days ago through a crack in the stonework. They don't know the history of Cynidicea or its fall, but they've glimpsed a few of the monument's residents and their customs from a safe distance. The pixies won't join the characters in exploring the ziggurat, but they can share the following information:


\begin{description}
\item[Dungeon Dangers.]
The ziggurat holds many dangers. The deeper you go, the worse it gets.
\item[Masked Factions.]
Everyone in the ziggurat wears a mask. If two people wear the same mask, they're part of the same faction. The factions don't seem to get along.
\end{description}

The pixies attempt to flee if attacked, defending themselves as necessary. They blow raspberries at the characters as they fly away.

\subparagraph{Fireworks}
Each sawdust-packed crate contains ten fireworks—stiff paper tubes filled with a gray powder. The fireworks have been neglected for centuries. Roll a d6 for each of the six crates. On a 6, the crate's contents are intact. Otherwise, the fireworks are clearly ruined.

If lit, an intact firework explodes in a harmless burst of light accompanied by a crackling noise audible up to 300 feet away. The light is as bright as a torch flame but lasts only a second.

\subparagraph{Pottery Jars}
Near the door stand three 4-foot-tall clay jars once used to put out fires in this room. The first contains sand. The others, which previously stored water, are empty.

\subsubsection{B8: Abandoned Room}
\begin{DndReadAloud}{}
A greenish, viscous slime covers the floor and walls of this bare room. In the middle lies a single decomposing corpse in blue robes and a brass mask.

\end{DndReadAloud}

The walls and floors of this room are covered in patches of green slime (see the \textit{Dungeon Master's Guide}). A dwarf member of the Guardians of Gorm perished here weeks ago. The slime has been slowly digesting the body since.

\subparagraph{Treasure}
The dead guardian wears a brass mask of Gorm worth 10 gp. The mask is out of reach of the green slime.

\subsubsection{B9: Abandoned Priest's Quarters}
\begin{DndReadAloud}{}
This neglected room contains a ragged bed, an oaken chest, and a writing table and chair. Puffs of dust and unsettling crunching noises emanate from atop the bed, where a lizard the size of a horse feeds on a humanoid corpse.

\end{DndReadAloud}

On the bed, a \textbf{giant lizard} munches on the body of an unfortunate Cynidicean. If the lizard notices the characters, it abandons its meal in favor of fresh meat. While feasting on the corpse, the lizard tangled its tail in the dead Cynidicean's mask; during combat, the mask remains wrapped around the lizard's tail.

A second \textbf{giant lizard} with pale blue scales and orange spots clings to the ceiling, waiting to ambush unsuspecting entrants. Characters who have a passive Wisdom (Perception) score of 13 or higher notice the lizard before it attacks.

\subparagraph{Treasure}
The Cynidicean's mask, depicting a hawk with a golden beak, is worth 20 gp. The chest contains a rusty \textit{dagger} and a set of torn vestments.

\subsubsection{B10: Abandoned Priest's Quarters}
This room is identical in arrangement to area B9, but it's unoccupied and contains no corpses. A wooden holy symbol in the shape of a balance rests on the desk.

\subparagraph{Treasure}
The chest in this room contains a \textit{priest's pack}. The rations inside spoiled long ago.

\subsubsection{B11: Barracks of Gorm}
\begin{DndReadAloud}{}
Three bunk beds line the walls of these ramshackle barracks. Four brawny humanoids sit on the lower bunks, listening intently to a fifth who ominously gestures from beneath a tattered white sheet, as if mimicking a ghost. Each figure wears a blue tunic draped with chain mail, along with a brass mask depicting the same long-haired, bearded man with a stern gaze.

\end{DndReadAloud}

Five guardians of Gorm recently repurposed this chamber to house their ranks. The sheet-draped guardian, a dwarf named Artha, is recounting an experience she had on the fourth level of the ziggurat. She claims to have seen the ghost of King Alendria there.

When the guardians notice the characters, Artha stops her story and asks the party to identify themselves and whom they serve. If the characters speak favorably of Gorm or declare no allegiances, Artha warmly invites them to join the Guardians of Gorm, then escorts the party to area B12 to meet the faction's leader, Kanadius. If the characters insult Gorm or claim to worship Madarua or Usamigaras, Artha tells the characters to turn back or suffer Gorm's thunderous wrath.

If the characters attack the guardians or refuse to heed their warnings, the guardians produce weapons and defend their turf. Combat in this room draws the attention of the guardians in area B12, who send two guardians of Gorm from that room to investigate. The guardians aim to knock the characters out rather than kill them, leaving their judgment to Kanadius.

\subparagraph{Treasure}
Three Mace rest against a weapons rack on the west wall. This room also contains enough food and water to sustain ten Humanoids for five days. The five guardians each wear a brass mask of Gorm worth 10 gp.

\subsubsection{B12: Grand Master of Gorm}
\includegraphics[width=0.4\textwidth]{images/../5e.tools/img/adventure/QftIS/013-02-006.guardians-of-gorm.png}
A gold-plated statue of Gorm—a smaller version of the statue atop the ziggurat—stands in the hallway east of this room. The gold flakes easily, revealing the statue's wooden construction beneath.

Characters who listen at the door hear drums thrumming within. When the characters enter, read or paraphrase the following text:

\begin{DndReadAloud}{}
Like claps of thunder, drums boom in this aged barracks. With each rumble, dust and sand trickle from the ceiling onto the drummers. These five humanoids each wear a blue tunic and a brass mask of the same long-haired, bearded man with a stern gaze.

Amid the drummers sits an older human man on a wooden throne. He wears distinguished, midnight-hued robes, and his long, winter-white hair peeks from beneath his mask and fanciful helm.

\end{DndReadAloud}

Five guardians of Gorm conduct a storm ceremony under the supervision of Kanadius (lawful good, human \textbf{champion of Gorm}; see appendix B), the Grand Master of Gorm. Kanadius—who wears a \textit{Helm of Telepathy}—is stoic but kind, with a grandfatherly demeanor. He leads by example and has earned the respect of his fellow guardians.

\subparagraph{Meeting Kanadius}
Kanadius and the other guardians are initially indifferent toward the party. When he's aware of the characters, Kanadius raises a wrinkled hand to quiet the drummers. If the characters don't act with hostility, Kanadius asks who they are and why they are here, using the magic of his helm to probe their answers for sincerity. Surmising the characters are outsiders, Kanadius offers to tell them about the once-proud kingdom of Cynidicea and its fall. He then might ask one or more characters to join his faction, depending on whether they answered him truthfully (see the "Joining the Guardians" section).

If the characters attack the Guardians of Gorm, Kanadius curses the characters, denouncing them as spies of Usamigaras. The guardians defend their leader with fervor, covering Kanadius's retreat as he attempts to flee through a trapdoor under his throne to the Ceremonial Chamber of Gorm (area B26), where he might be encountered later.

\subparagraph{What Kanadius Knows}
In addition to the history and fall of Cynidicea (detailed at the beginning of this adventure), Kanadius can impart the following information:


\begin{description}
\item[Cynidicean Factions.]
There are three other groups in the ziggurat: the Mages of Usamigaras, the Warriors of Madarua, and the Cult of Zargon. Kanadius doesn't trust any of them, but the Zargonites are especially ruthless.
\item[Dungeon Tip.]
The lower levels of the ziggurat are haunted by Undead.
\item[Glory of Cynidicea.]
Kanadius believes only the Guardians of Gorm can restore Cynidicea's lost glory.
\end{description}

\subparagraph{Joining the Guardians}
If a character was honest with Kanadius, he asks them to join the Guardians of Gorm. Kanadius rewards the first three characters who accept with a \textit{Potion of Healing}. He then leads them to the Ceremonial Chamber of Gorm (area B26) for a brief induction ceremony, where each inductee is given a storm-blue tunic, a brass mask of Gorm worth 10 gp, and a blue tattoo of a lightning bolt on their right shoulder. Kanadius also tells inductees the password to avoid the trapped door to area B26: "by the great god Gorm."

The Guardians of Gorm and their allies can rest safely here and in areas B11 and B26 without triggering random encounters.

\subparagraph{Rejecting the Guardians}
If the characters reject Kanadius's offer, he nods solemnly and politely asks them to go back the way they came. The faction doesn't allow nonmembers to enter their treasure room or ceremonial chamber. Kanadius holds no grudges if the characters later change their minds.

\subparagraph{Trapdoor}
A secret trapdoor lies under Kanadius's throne. A character who inspects the throne and succeeds on a DC 13 Intelligence (Investigation) check notices marks around the chair and discovers a ladder beneath the trapdoor, which leads down to the secret door in area B26.

\subparagraph{Treasure}
Kanadius wears a \textit{Helm of Telepathy}. He and the other five guardians each wear a brass mask of Gorm worth 10 gp. Stashed in a chest under one of the beds are six more masks, six blue tunics, and three Potion of Healing.

\subsubsection{B13: Revolving Passage}
This central area connects eight hallways on the third tier of the ziggurat via a revolving passage on a stone turntable.

On the wall inside the revolving passage is a row of four levers, each corresponding to one of the passage's possible orientations: running north to south, east to west, northwest to southeast, or northeast to southwest. When a lever is pulled, if the passage isn't already aligned with the corresponding orientation, the passage grinds faintly and swings clockwise to match the lever. This process takes 1 round. Only one lever can be engaged at once, and pulling a lever disengages any previously engaged lever.

\subparagraph{Adjoining Corridors}
Each of the eight hallways adjacent to the revolving passage sports a similar lever embedded in the wall. Pulling one of these levers rotates the passage to meet that hallway as if the corresponding lever had been pulled inside area B13. When the passage isn't aligned with a hallway, it appears as a curved wall from the outside.

\subsubsection{B14: Abandoned Ceremonial Chamber}
\begin{DndReadAloud}{}
Rubble litters the floor of this vandalized and ruined temple. The altar is cracked in half, and the chamber's once-brilliant banners, now tattered and stained, sag from its walls. Scrawled on the wall is the image of a tentacled eye with a long, black horn and one word in Common: "Zargon."

\end{DndReadAloud}

The cultists of Zargon ravaged this room years ago. Except for the wreckage, the room is empty.

\subparagraph{Hidden Ramp}
On the south side of the chamber, a hinged, 20-foot-long section of the floor gently swings down to area B27 when a Small or larger creature steps on it. The character with the highest passive Wisdom (Perception) score notices the ramp without stepping on it.

\subsubsection{B15: Abandoned Storeroom}
\begin{DndReadAloud}{}
Clusters of pale-yellow fungus blanket the shelves of this storeroom. The air is thick with spores.

\end{DndReadAloud}

The stores in this room are covered in four patches of yellow mold (see the \textit{Dungeon Master's Guide}).

Most of the food here is unfit for consumption, but a crate marked with a bee along the east wall contains 1d4 doses of healing honey, crystallized but still edible. A creature that consumes one dose of this honey as an action regains 1d4 + 1 hit points. This honey was produced by bees like those in area B5. Characters can open the crate without disturbing the yellow mold.

\subparagraph{Secret Door}
A secret door in the south wall leads to area B16.

\subsubsection{B16: Secret Room}
This room is hidden behind a secret door in the hallway bordering it and another in the next room. A large wicker basket, filled with coins and gems, rests on the floor. Two horned vipers (poisonous snakes) hide in the basket, ready to strike.

\subparagraph{Treasure}
The basket contains 300 sp and five jasper stones worth 50 gp each.

\subsubsection{B17: Abandoned Ceremonial Chamber}
\begin{DndReadAloud}{}
Tattered tapestries hang from the walls of this ruined chapel. Along the north wall lies an altar draped with a ratty cloth. Atop it rest three wooden candlesticks, an offering bowl, and a holy symbol shaped like an eye. A mangled human corpse slumps against the altar.

\end{DndReadAloud}

The corpse is a thief who sought to plunder the ziggurat but fell to one of its many wandering threats. Unlike the Cynidiceans, the plunderer wears no mask.

\subparagraph{Treasure}
The corpse carries a full canteen of water, three days of rations, and a sack with 40 gp and two citrines worth 50 gp each. The holy symbol—which honors the elder evil Zargon—houses a smooth sphere of yellow jasper worth 25 gp.

\subsubsection{B18: Chamber of Usamigaras}
\includegraphics[width=0.4\textwidth]{images/../5e.tools/img/adventure/QftIS/014-02-008.mages-of-usamigraras.png}
The corridors outside this room are painted black with tiny stars resembling the night sky. The room's iron doors each bear an engraved star. An audible \textit{Alarm} spell warns the room's inhabitants when a creature other than a mage of Usamigaras touches either door.

When the characters enter, read or paraphrase:

\begin{DndReadAloud}{}
Along the north and west walls of this spacious ceremonial hall hangs a tapestry of the night sky, its fabric embroidered with silvery constellations. In front of a star-shaped stone altar, five figures gather, each wearing starry robes and a silver mask depicting the face of a smiling child. The centermost figure wears a jewel-studded silver crown.

\end{DndReadAloud}

Four mages of Usamigaras have just initiated a new member in their faction's ceremonial chamber. The ritual was supervised by Auriga Sirkinos (chaotic neutral, elf \textbf{champion of Usamigaras}; see appendix B), Chief Mage of Usamigaras. Superficially amicable and jolly, their pleasant exterior conceals a treacherous fanatic who would do anything to promote their faction. Auriga and the other mages are initially indifferent toward the party.

\subparagraph{Meeting Auriga}
Auriga welcomes the characters warmly, asking if the party has come to join the Mages of Usamigaras. To further tantalize the characters, the chief mage coyly remarks that the benefits of membership are many, but the Grinning God's expectations of his devotees are high. Characters who wish to join the faction must prove their worth (see the "Joining the Mages" section).

If the characters attack, the mages defend themselves and their leader. Auriga attempts to flee to the lower levels of the ziggurat and might be encountered there later.

\subparagraph{What Auriga Knows}
In addition to the history and fall of Cynidicea (detailed at the beginning of this adventure), Auriga can impart the following information:


\begin{description}
\item[Cynidicean Factions.]
Three other groups vie for power in the ziggurat: the Guardians of Gorm, the Warriors of Madarua, and the Cult of Zargon. The monster the Zargonites serve dwells deep beneath the ziggurat and is rumored to be unkillable.
\end{description}

\includegraphics[width=0.4\textwidth]{images/../5e.tools/img/adventure/QftIS/015-02-007.usamigaras-mask.png}

\begin{description}
\item[Demetrius.]
Years ago, the cultists of Zargon killed the former leader of the Mages of Usamigaras, a respected priest named Demetrius.
\item[Dungeon Tip.]
The ziggurat has many secret doors, which blend in with their surroundings.
\end{description}

\subparagraph{Joining the Mages}
To be considered for membership, a character must reveal a secret about themself—whether false or true—and make a DC 14 Charisma (Deception or Persuasion) check. Characters who openly display allegiance to another faction in the ziggurat, such as by wearing a mask of Gorm, have disadvantage on this check. Characters who prove they've defeated members of another faction have advantage on this check. On a successful check, Auriga inducts that character into the Mages of Usamigaras. If a character tries to dupe the mages and fails, Auriga remarks that the character is a poor liar.

The mages give each inductee a set of starry robes, a silvered \textit{dagger}, a silver mask worth 25 gp, and a silver tattoo of a five-pointed star on their right palm. Auriga also tells inductees how to operate the shifting wall in area B21. Additionally, if at least one character joins the faction, Auriga gifts the party a \textit{Wand of Secrets} from a hidden niche in the altar.

The Mages of Usamigaras and their allies can rest safely here and in area B19 without triggering random encounters.

\subparagraph{Rejecting the Mages}
If no characters join the Mages of Usamigaras, Auriga allows the party to leave unharmed. For the next three days, characters have disadvantage on Charisma checks made to influence members of the faction.

\subparagraph{Treasure}
Auriga and their four mages each wear a silver mask of Usamigaras worth 25 gp. Auriga wears a bejeweled silver crown worth 100 gp and carries a \textit{Wand of Magic Detection}.

A character who inspects the altar and succeeds on a DC 15 Intelligence (Investigation) check finds a hidden compartment. Inside are six more masks, six sets of robes, six silvered Dagger, and a \textit{Wand of Secrets}.

\subsubsection{B19: Mages' Quarters}
\begin{DndReadAloud}{}
Six bunk beds dot this well-kept chamber. Two wooden chests rest at the foot of each bed.

\end{DndReadAloud}

This remarkably clean room belongs to the Mages of Usamigaras. The mages in area B18 sleep here when away from their stronghold in the underground city (see the "Extending the Adventure" section). The chests contain spare robes and various personal belongings but nothing of value.

\subparagraph{Secret Door}
A secret door in the south wall leads to Auriga's quarters (area B20).

\subsubsection{B20: Chief Mage's Quarters}
\begin{DndReadAloud}{}
This comfortable bedroom shows signs of recent occupation. At the foot of the bed sits an iron chest—and next to it, a gray-furred wolf lies curled on the floor. The wolf wakes and growls as the door opens.

\end{DndReadAloud}

A \textbf{dire wolf}, captured as a cub and trained as a guard dog, protects Auriga's private chambers. The wolf is hostile toward anyone other than Auriga who enters this room, but a character who makes calming gestures or offers it food can placate it with a successful DC 16 Wisdom (Animal Handling) check. Characters dressed in the attire of the Mages of Usamigaras have advantage on this check. If a character pacifies the wolf, it retrieves a slobbery humanoid bone from under Auriga's bed and whines until a character plays fetch with it.

\subparagraph{Trapped Chest}
The chest is locked, and Auriga carries its key. As an action, a character can pick the lock with a successful DC 15 Dexterity (Sleight of Hand) check using thieves' tools.

The chest is also trapped. A character who examines the chest and succeeds on a DC 14 Intelligence (Investigation) check discovers the trap, which can be disabled as an action with a successful DC 14 Dexterity (Sleight of Hand) check using thieves' tools. A creature that opens the chest without using the proper key or disabling the trap causes a poison needle to spring from the lock. The creature that triggered the poison needle takes 1 piercing damage plus 11 (2d10) poison damage and must succeed on a DC 15 Constitution saving throw or have the poisoned condition for 1 hour.

\subparagraph{Treasure}
The trapped chest holds 40 gp, 500 sp, and a vial with two doses of \textit{carrion crawler mucus} (a type of poison described in the \textit{Dungeon Master's Guide}).

\subsubsection{B21: Shifting Wall}
\begin{DndReadAloud}{}
At the western end of this short corridor stands an eight-foot-tall silver statue of an impish, winged child with two snakes twined around his body. The smiling, round-faced boy holds a wand in one hand and a handful of coins in the other. The statue is firmly connected to the wall.

\end{DndReadAloud}

The silver-plated statue at the end of the hall depicts the god Usamigaras and hides a shifting wall. Tapping on the statue reveals its hollow construction.

A character who inspects the statue and succeeds on a DC 13 Intelligence (Investigation) check learns that the statue's wand controls a mechanism in the wall. A character who succeeds on this check by 5 or more also knows the correct direction to pull the wand to cause the wall to shift. Members of the Mages of Usamigaras, including newly initiated characters, know how to control the shifting wall.

If the wand is pulled to the left, the wall and statue rotate ninety degrees, allowing passage into the corridor beyond for 1 minute, then shift back into place.

If the wand is pulled to the right, the statue exhales poison gas in a 15-foot-cone. Each creature in that area must succeed on a DC 13 Constitution saving throw or take 5 (1d10) poison damage.

\subsubsection{B22: Storeroom}
\begin{DndReadAloud}{}
Dusty crates and barrels fill this storeroom. Four small, winged creatures, surrounded by clouds of sand, crowd a cask with wine trickling from its spout.

\end{DndReadAloud}

Four dust mephits rummage through the stores in this room. They've just discovered a cask of wine. The dust mephits are hostile toward characters who disrupt their festivities but friendly toward any who join in the revelry.

\subparagraph{Secret Door}
The south wall of this room contains a secret door.

\subsubsection{B23: Ceremonial Chamber of Madarua}
In the hallway outside this shrine, two 10-foot-tall bronze statues of female warriors stand with spears crossed to form an arch. Beneath them lies a hidden pressure plate that, when stepped on, rings a chime in area B25 to alert the Warriors of Madarua of creatures approaching their shrine. A character who examines the arch and succeeds on a DC 15 Wisdom (Perception) check notices the pressure plate.

\includegraphics[width=0.4\textwidth]{images/../5e.tools/img/adventure/QftIS/016-02-009.madarua-mask.png}
The door to the chamber is locked. As an action, a character can pick the lock with a successful DC 15 Dexterity (Sleight of Hand) check using thieves' tools or force open the door with a successful DC 18 Strength (Athletics) check.

When the characters enter the chamber, read or paraphrase the following text:

\begin{DndReadAloud}{}
The fragrance of burning incense pervades this ceremonial chamber. Clean white fabric adorns its walls. Near the northeast corner, two candles illuminate a cloth-draped altar. On the altar stands a statuette of a woman holding a sword and a sheaf of wheat.

\end{DndReadAloud}

Tiny listening holes dot the south wall, concealed by hanging cloths. If the characters forced open the door, the warriors in area B25 investigate, and the encounter in that area happens here instead.

\subparagraph{Secret Door}
A cloth on the south wall hides a secret door to area B24.

\subparagraph{Treasure}
The statuette depicts the god Madarua. It is made of bronze and worth 10 gp.

\subsubsection{B24: Treasure Room of Madarua}
\begin{DndReadAloud}{}
The stone walls of this narrow chamber are bare. Set into the east wall, a stone vault has a single keyhole on its heavy, round door.

\end{DndReadAloud}

This room's vault contains the treasure of the Warriors of Madarua. The faction's leader, Pandora (see area B25), carries the key. As an action, a character can pick the vault's lock with a successful DC 16 Dexterity (Sleight of Hand) check using thieves' tools. A character who fails this check, or who tries to force the vault open, triggers the ceiling trap.

\subparagraph{Ceiling Trap}
A character who inspects the vault and succeeds on a DC 14 Intelligence (Investigation) check determines the stone ceiling above the vault is rigged to crumble if the vault is forced open. When the trap triggers, a 10-foot-wide section of unstable stone along the ceiling collapses in front of the vault door. Creatures in the area marked on the map must make a DC 15 Dexterity saving throw, taking 11 (2d10) bludgeoning damage on a failed save or half as much damage on a successful one.

Once the trap is triggered, the floor of that area is filled with rubble and becomes \textit{difficult terrain}. Triggering the trap attracts the attention of the warriors in area B25, who investigate the disturbance.

\subparagraph{Secret Doors}
The north and south walls of this room hold the secret doors leading from areas B23 and B25. The outlines of both doors are visible here.

\subparagraph{Treasure}
The vault contains 300 gp, ten tiger's eye gems worth 25 gp each, six bronze Spear worth 10 gp each, six bronze masks of Madarua worth 15 gp each, and a Spear of Warning.

\subsubsection{B25: Gymnasium of Madarua}
\includegraphics[width=0.4\textwidth]{images/../5e.tools/img/adventure/QftIS/017-02-010.warriors-of-madarua.png}
\begin{DndReadAloud}{}
Six warriors spar with spears in this stone gymnasium. Each warrior is clad in leather armor over a green tunic and wears a bronze mask of a determined female face. A tall human woman in lightweight armor decorated with bronze supervises them, interjecting to critique their footing.

\end{DndReadAloud}

Six warriors of Madarua train in this area. The woman supervising them is Pandora (lawful neutral, human \textbf{champion of Madarua}; see appendix B), Madarua's Champion. The oldest and most experienced warrior in the faction, Pandora believes any problem can be solved with swift, direct action. She and the other warriors are initially indifferent to the party.

\subparagraph{Meeting Pandora}
When Pandora notices the characters, she and her warriors assume a defensive formation. Pandora warns that the characters have entered one of Madarua's sacred chambers and asks whether the characters are friend or foe.

If the characters don't attack, she explains that the Warriors of Madarua are yet another group vying to reclaim Cynidicea. After suffering a handful of losses to the cultists of Zargon on the lower levels, Pandora is on the lookout for new recruits—especially fearless warriors who share her god's interests. She might ask one or more characters to join her faction (see the "Joining the Warriors" section).

If the characters attack, the warriors defend themselves fiercely and fight to the death. If slain, Pandora uses her last breath to declare her vanquisher Madarua's new champion.

\subparagraph{What Pandora Knows}
In addition to the history and fall of Cynidicea (detailed at the beginning of this adventure), Pandora can impart the following:


\begin{description}
\item[Cynidicean Factions.]
There are three other groups in the ziggurat: the Guardians of Gorm, the Mages of Usamigaras, and the Cult of Zargon. The Zargonites kidnap Cynidiceans and sacrifice them to the one-eyed monster they worship.
\item[Dungeon Tip.]
The tombs on the levels below are riddled with traps.
\item[Glory of Cynidicea.]
Pandora believes that only the Warriors of Madarua can restore Cynidicea's lost glory.
\end{description}

\subparagraph{Joining the Warriors}
Pandora admires disciplined warriors who can hold their own. If a character expresses interest in joining the Warriors of Madarua—and if they don't display allegiance to another faction, such as a prominent tattoo or mask—Pandora challenges them to a wrestling match for membership. A character who succeeds on a DC 15 Strength (Athletics) check impresses Pandora during their bout, earning them a faction invite.

New faction members are inducted in the ceremonial chamber of Madarua (area B23). The warriors give each inductee a green tunic, a \textit{spear}, a bronze mask of Madarua worth 15 gp, and a small sickle-shaped brand on their left wrist. Additionally, if at least one character joins the Warriors of Madarua, Pandora gifts the party the Spear of Warning from the faction's treasure room (area B24).

The Warriors of Madarua and their allies can rest safely in this area and in area B23 without triggering random encounters.

\subparagraph{Rejecting the Warriors}
If no characters join the Warriors of Madarua, Pandora respects their decision but offers them no charity. Two warriors escort the characters to the revolving passage (area B13). If the characters refuse to leave or later return, the warriors regard them with hostility.

\subparagraph{Secret Door}
The north wall hides a secret door to area B24.

\subparagraph{Treasure}
Pandora carries a ceremonial \textit{+1 Longsword} and the key to the treasure vault in area B24. She and the other six warriors each wear a bronze mask of Madarua worth 15 gp.

\subsubsection{B26: Ceremonial Chamber of Gorm}
The door to this ceremonial chamber is engraved with three lightning bolts limned in blue light. The door is locked and magically trapped. As an action, a character can pick the lock with a successful DC 15 Dexterity (Sleight of Hand) check using thieves' tools. A \textit{Detect Magic} spell reveals an aura of evocation magic around the door. If a creature outside the chamber picks the lock or touches the trapped door without first speaking the words "by the great god Gorm," the creature takes 4 (1d8) lightning damage.

When the characters enter, read or paraphrase:

\begin{DndReadAloud}{}
The walls, ceiling, and floor of this spacious chamber have been painted sky blue. Flickering candles surround a marble altar near the east wall. Atop the altar rest a golden bowl and a stone statuette of a bearded human man hurling a lightning bolt.

\end{DndReadAloud}

The Guardians of Gorm conduct religious ceremonies in this dimly lit chamber. Characters who join the faction are inducted here.

If the characters attacked Kanadius earlier in the adventure and he escaped (see area B12), he's here along with two guardians of Gorm. Otherwise, the room is unoccupied.

\subparagraph{Secret Door}
The secret door in the west wall conceals a ladder up to the trapdoor in area B12.

\subparagraph{Treasure}
The gold altar bowl is worth 50 gp. A character who inspects the altar and succeeds on a DC 15 Intelligence (Investigation) check finds a secret compartment containing a \textit{Javelin of Lightning}.

\subsection{Ziggurat Locations, Tier 4}
The fourth tier of the ziggurat holds the tombs of King Alendria, Queen Zanobis, and various members of their court. The following locations are keyed to map 2.2.

\includegraphics[width=0.4\textwidth]{images/../5e.tools/img/adventure/QftIS/018-map-2.02-buried-ziggurat-4.png}
\includegraphics[width=0.4\textwidth]{images/../5e.tools/img/adventure/QftIS/019-map-2.02-buried-ziggurat-4-player.png}
\subsubsection{B27: Jester's Burial Room}
\begin{DndReadAloud}{}
The walls of this burial chamber sport colorful murals of a human court jester entertaining her king and queen. A niche on the east wall holds a humble coffin.

\end{DndReadAloud}

If a creature touches the coffin, the lid flips open, and a cheery wooden jester on a spring pops up and does a little jig. The prop then hangs lifelessly until the trick is reset by shoving the prop back into the coffin, which also contains the jester's ashes.

\subparagraph{Ramp}
A portion of the ceiling can swing down on a hinge, forming a ramp leading from area B14 to this chamber. A small handle bolted to the ceiling of this room allows the ramp to be easily raised and lowered by a creature that can reach it.

\subsubsection{B28: High Priest's Burial Room}
Scorch marks from a former trap cover the door to this room and the wall around it. This magical trap was recently triggered by a ghoul, whose outline is visible among the scorch marks. The door is ajar and safe to open.

When the characters enter, read or paraphrase:

\begin{DndReadAloud}{}
The walls of this burial room display images of a priest performing ceremonies. On a raised dais in the middle of the room lies an open bronze coffin. Three gangly, ghoulish creatures hunch over the coffin, ripping into the corpse inside and feasting on aged flesh.

\end{DndReadAloud}

Three ghouls are devouring the corpse of a former high priest. When they detect the characters, the ghouls attack, eager to taste ripe flesh.

\subparagraph{Treasure}
The coffin holds a jeweled necklace worth 300 gp, two jeweled bracelets worth 100 gp each, and two Holy Water (flask).

\subsubsection{B29: Ghostly Haunt}
\begin{DndReadAloud}{}
As you proceed down this corridor, the ghost of a human man in a shimmering crown and regal attire appears before you. The phantom raises his hand and gestures for you to stop.

"Go no farther," he warns, "lest this cursed monument overtake you. Only death awaits you in these halls."

\end{DndReadAloud}

The figure is the spirit of Alendria (lawful good \textbf{ghost}), the last king of Cynidicea. Cursed to wander his once-splendid tomb, the ghost lingers here to warn honorable travelers of its dangers.

If the characters don't attack, the ghost reveals the following information, then fades away:


\begin{description}
\item[Escaping the Ziggurat.]
The main entrance to—and exit from—the ziggurat (area B57) lies on the next tier down. Its doors have been closed for centuries and can be opened only from the inside.
\item[Haunted Tombs.]
Horrible creatures haunt the burial chambers of the king and queen. One of the creatures claims to be Queen Zanobis herself. The monarchs can't rest until these wicked Undead are slain.
\item[Undying Enemy.]
Only fools challenge Zargon and his cultists. All who oppose them fail.
\end{description}

If the characters attack Alendria's ghost, the spirit admonishes them and disappears without delivering the hints above.

\subparagraph{Secret Door}
The southwest corner hides an entrance to the tomb annex (area B41).

\subsubsection{B30: Guard Captain's Burial Room}
\begin{DndReadAloud}{}
A cobwebbed human skeleton in decorative pauldrons stands at attention against the west wall of this burial chamber. Four skeletal guards with rusty swords surround the distinguished skeleton.

\end{DndReadAloud}

This burial chamber belongs to King Alendria's honor guard, who defend his kingdom even in death. If the characters enter before defeating the banshee in area B44, the five skeletons animate and attack. The skeletons don't pursue characters outside this room.

If the characters defeat the banshee, the skeletons instead kneel, each pledging a silent vow to defend the characters for the next 1d4 hours or until the skeletons earn their second death (whichever comes sooner). After these vows are fulfilled, the skeletons crumble to dust.

\subsubsection{B31: Noble's Burial Room}
\begin{DndReadAloud}{}
Scenes of warfare decorate the walls of this burial room. A closed wooden coffin stands upright near the south wall. The image of a noble warrior adorns its lid.

\end{DndReadAloud}

A \textbf{mummy} rests in the coffin, clutching a \textit{+1 Rapier} to its wrinkled chest. If disturbed, the mummy awakens, releases the rapier, and attacks.

\subsubsection{B32: Ankheg Den}
\begin{DndReadAloud}{}
Three empty coffins sit askew in this chamber, their lids chewed through and tossed aside. Near them lie the bubbling remains of a giant rat. Piled in the southwest corner of the room is a mound of rubble.

\end{DndReadAloud}

This burial room has become the den of an \textbf{ankheg}, drawn to the tomb complex to feast on giant rodents and Cynidiceans who come to pay their respects. A character who examines the dead rat and succeeds on a DC 13 Intelligence (Investigation) check determines the rat was killed by acid.

Characters who have a passive Wisdom (Perception) score of 14 or higher feel intensifying tremors beneath them. If a character approaches the mound or remains in this room for longer than 1 minute, the ankheg erupts from the floor and attacks, eager for its next meal.

\subparagraph{Treasure}
The mound obscures a narrow dead-end tunnel in the west wall, just big enough for a Small creature to squeeze through. Here, a giant rat has stored its treasures: a corroded \textit{dagger}, a purse containing 55 gp, and an unhatched ankheg egg.

\subsubsection{B33: Astrologer's Burial Room}
\begin{DndReadAloud}{}
A coffin adorned with glass stars rests atop a square dais in this burial chamber. Magical flames dance above two brass jars, one at each end of the coffin, illuminating scenes on the ceiling of an astrologer creating maps of both the sky and this realm.

\end{DndReadAloud}

This burial chamber is brightly lit by two \textit{Continual Flame} spells cast on brass jars at the coffin's head and foot. Inside the coffin lies the astrologer's body.

\subparagraph{Treasure}
A bronze tube in one of the brass jars contains a crudely drawn map of this tier of the ziggurat. Two stars on the map denote the locations of the trapdoor to the fifth tier (see area B40) and the ramp to the third tier (see area B27).

\subsubsection{B34: Master Thief's Burial Room}
\begin{DndReadAloud}{}
An elaborate jeweled coffin rests in the center of this otherwise bare room.

\end{DndReadAloud}

A \textbf{carrion crawler} clings to the ceiling in this room, waiting to ambush entrants. Characters who have a passive Wisdom (Perception) score of 13 or higher notice the carrion crawler before it attacks.

A master thief stole a prime burial place near the king and queen, arranging for her locked coffin to be placed in this unmarked chamber. As an action, a character can pick the lock with a successful DC 20 Dexterity (Sleight of Hand) check using thieves' tools. Inside rests the decrepit, still-smirking corpse of a halfling.

\subparagraph{Treasure}
Ten jewels set in the coffin are worth 50 gp each. Inside the coffin is a gold-plated set of \textit{thieves' tools} worth 10 gp and two silvered Dagger.

\subsubsection{B35: Chamberlain's Burial Room}
The door to this room is boarded up with dusty planks. When the characters enter, read:

\begin{DndReadAloud}{}
This burial room reeks of death. Beneath murals depicting a human man signing documents and collecting taxes, a smashed wooden coffin lies on the floor. Five wound-ridden zombies with lipless mouths and clouded eyes shuffle about the room.

\end{DndReadAloud}

This burial room was looted during Cynidicea's fall. In the years following, trespassers fell prey to the dangers of the tomb complex. They rose again as five zombies that rush to consume entrants.

\subsubsection{B36: Slime Fountain}
\begin{DndReadAloud}{}
This malodorous chamber is dominated by a polluted fountain sculpted into a five-foot-tall ziggurat. A spout atop the fountain dribbles yellow-brown slime into a pool of burping sludge. A dwarf's shriveled corpse slumps over the pool's edge.

\end{DndReadAloud}

Constructed to honor the ziggurat's architect, this fountain draws from a contaminated reservoir deep underground. Years ago, a thirsty dwarf explorer succumbed to its corrupted waters.

A creature that drinks from the fountain or touches its waters must succeed on a DC 11 Constitution saving throw or have the poisoned condition for 1 hour. A \textit{Purify Food and Drink} spell cleans the fountain and causes its waters to run clear for 1 hour, after which its spout resumes its putrid flow.

\subsubsection{B37: Embalming Room}
\begin{DndReadAloud}{}
This acrid-smelling room is littered with decrepit wooden tables bowed under the weight of heavy jars. In the corner, rotten winding cloths fill a stone bin.

\end{DndReadAloud}

This room was used for embalming corpses. The jars on the tables contain various chemicals for preserving bodies. When the characters enter, two shadows rise from the floor and attack.

\subparagraph{Treasure}
One of the jars is a Alchemist's Fire (flask). The linen bin contains a \textit{Potion of Invisibility}.

\subsubsection{B38: Lady-in-Waiting's Burial Room}
\begin{DndReadAloud}{}
An open ivory coffin lies in a niche on the east wall of this burial room. Frescoes of court life decorate the walls, focusing on the queen's elf handmaiden.

\end{DndReadAloud}

The coffin is empty. It once belonged to Therisina, the queen's lady-in-waiting. Unknown to Queen Zanobis, Therisina secretly loved King Alendria and never overcame her sorrow at his passing. She now haunts his tomb (see area B44). A character who examines the frescoes notices Therisina eying King Alendria with admiration in one scene.

\subsubsection{B39: Councilor's Burial Room}
\begin{DndReadAloud}{}
The painted walls of this room depict scenes of a throne room where a well-dressed councilor advises monarchs. In the center of the room lies a wooden coffin riddled with bite marks and a gaping hole.

\end{DndReadAloud}

Four giant rats have chewed a hole in the coffin and made it their den. The coffin is otherwise empty, its contents pilfered long ago. The rats attack anyone who approaches the coffin.

\subsubsection{B40: Southwest Corridor}
\begin{DndReadAloud}{}
Near the end of this long corridor, a marble slab with a bronze hand ring bolted to it interrupts the floor.

\end{DndReadAloud}

The slab is a trapdoor. Pulling the ring reveals a ladder down to area B45 on the tier below.

\subparagraph{Wandering Spirit}
If the characters have not yet encountered King Alendria's spirit in area B29, the ghost appears to them here instead.

\subparagraph{Secret Door}
The wall north of the trapdoor hides an entrance to the tomb annex (area B41).

\subsubsection{B41: Tomb Annex}
\begin{DndReadAloud}{}
Red and white tiles checker the floor of this long corridor, and cracked frescoes of courtly scenes decorate its walls. To the east, the corridor is blocked by a pile of rubble, from which a skeletal hand protrudes.

\end{DndReadAloud}

This sealed annex, accessible via two secret doors at the west end of the corridor, leads to the burial chambers of King Alendria and Queen Zanobis. The annex contains several obstacles designed to deter tomb robbers and other would-be plunderers.

\subparagraph{Blade Traps}
Three swinging blades are hidden in the ceiling above the squares marked on the map. A successful DC 14 Wisdom (Perception) check reveals scrape marks and old bloodstains on the floor beneath each blade. As an action, a character can disarm a blade with a successful DC 15 Dexterity (Sleight of Hand) check using thieves' tools, but failing this check triggers the trap. Creatures can also avoid the blades by not walking under them.

The blades sweep down on each creature that passes under them. The target must make a DC 15 Dexterity saving throw, taking 9 (2d8) slashing damage on a failed save or half as much damage on a successful one. After a blade swings, the trap resets.

\subparagraph{Rubble}
Two stone slabs fell from the ceiling long ago, killing a thief and crumbling into rubble. The thief's skeletal hand sticks out of the rubble. The party can spend 10 minutes clearing the rubble to gain access to the east side of the hall.

\subparagraph{Twisting Tunnel}
Beyond the rubble is a narrow tunnel just large enough for a Medium creature to squeeze through; it leads to areas B42–B44.

\subsubsection{B42: False Tomb}
\begin{DndReadAloud}{}
Two golden sarcophagi gleam within this hallowed tomb. Precious objects surround them: a gilded chariot, two high-backed thrones, several wooden chests, and exquisite pottery. A mosaic spans the walls, depicting the lives of the monarchs who rest here.

\end{DndReadAloud}

Everything in this tomb is a sham. The sarcophagi are wood painted to look like solid gold, and they contain only chalk-white skeletons made of clay. Though the chests seem to brim with coins and gems, a character who inspects them and succeeds on a DC 12 Intelligence (Investigation) check learns they are cheap counterfeits and bits of colored glass.

\subparagraph{Secret Doors}
Secret doors in the north and south walls lead to the true burial chambers of King Alendria (area B44) and Queen Zanobis (area B43).

\subsubsection{B43: Queen Zanobis's Tomb}
\begin{DndReadAloud}{}
The cold stone walls of this royal burial chamber honor the last great queen of Cynidicea. Reliefs tell the story of a beloved queen who served her people nobly and received their respect and adoration in turn.

Her sarcophagus, artfully sculpted from white marble, rests on a stone pedestal in the center of the room. It bears the following engraving: "May death claim all who dare desecrate my resting place."

\end{DndReadAloud}

This tomb once contained a treasure trove, which the surviving Cynidiceans plundered long ago. As an action, a character can try to pry open the lid of the sarcophagus, doing so with a successful DC 13 Strength (Athletics) check. If the sarcophagus is opened, Queen Zanobis—now a \textbf{wight}—springs from it and attacks. She wields a scepter instead of a blade, so her melee attacks deal bludgeoning damage instead of slashing damage.

The queen's mind is clouded by rage. As she lashes out at the characters, Queen Zanobis decries them as "ungrateful looters," seeing them as the Cynidiceans who plundered her tomb in years past. If the queen is slain, she crumbles to ash with a grateful sigh, and her spirit is finally put to rest. If the characters later encounter King Alendria's spirit (see area B44), the spirit of Queen Zanobis appears alongside him to thank the characters.

\subparagraph{Treasure}
Queen Zanobis wears a jeweled crown worth 1,000 gp. Her scepter is a fully charged \textit{Wand of Paralysis}, but if its last charge is expended, the scepter loses its magic. In undeath, Queen Zanobis has forgotten how to access the scepter's magic.

\includegraphics[width=0.4\textwidth]{images/../5e.tools/img/adventure/QftIS/020-02-011.queen-zanobis.png}
\subsubsection{B44: King Alendria's Tomb}
The hall outside this chamber echoes with faint, ethereal sobs. When the characters enter, read:

\begin{DndReadAloud}{}
The chill of death pervades this royal tomb. Its walls depict scenes of the last great king of Cynidicea. An elegant sarcophagus rests in the center of the room.

Suddenly, a weeping elf woman, luminous and incorporeal, rises through the sarcophagus. She floats above the floor in tattered robes, her face wreathed in a tangled mane of wiry hair.

\end{DndReadAloud}

A \textbf{banshee} haunts King Alendria's grave, weeping over the dead monarch centuries after his passing. Characters who saw the frescoes in area B38 recognize the wailing spirit as Queen Zanobis's lady-in-waiting, Therisina. She attacks all who enter.

\subparagraph{Kingly Gift}
If the characters defeat the banshee, the spirit of King Alendria appears to them and bestows a Charm of Heroism (see the \textit{Dungeon Master's Guide}) on each character responsible for this act. He then fades away. This gift vanishes from any character who plunders his sarcophagus.

\subparagraph{Treasure}
King Alendria's tomb was stripped of its riches long ago, except for those inside his sarcophagus. The king's corpse wears a jeweled crown worth 1,000 gp and a suit of \textit{plate armor}.

\subsection{Ziggurat Locations, Tier 5}
The fifth tier once served as a temple and city hub. The following locations are keyed to map 2.3.

\includegraphics[width=0.4\textwidth]{images/../5e.tools/img/adventure/QftIS/021-map-2.03-buried-ziggurat-5.png}
\includegraphics[width=0.4\textwidth]{images/../5e.tools/img/adventure/QftIS/022-map-2.03-buried-ziggurat-5-player.png}
\subsubsection{B45: Secret Room}
A ladder in this room leads up to the trapdoor in area B40. The outline of a secret door to area B46 is visible here. The room is otherwise empty.

\subsubsection{B46: Demetrius's Bedroom}
\begin{DndReadAloud}{}
Dust coats this abandoned bedroom and its wooden bed, table, and chair. On the table lie an unlit lantern and a wooden holy symbol of two intertwined snakes. A faintly luminous white robe is draped over the chair.

\end{DndReadAloud}

This was the bedroom of Demetrius, a priest in the Mages of Usamigaras. When the Cult of Zargon corrupted his twin brother, Darius, Demetrius vowed to destroy the cult. However, its hold over Darius proved too great, and Darius assassinated his brother. Demetrius's vengeful spirit now resides within the softly glowing robe.

\subparagraph{Secret Doors}
The north wall hides two secret doors to areas B45 and B47.

\subparagraph{Spiritual Inhabitation}
A character who touches the robe must make a DC 16 Charisma saving throw. On a failed save, Demetrius's spirit inhabits the character's body, and the robe stops glowing. On a successful save, the character resists the spirit, who can't attempt to inhabit the character for 24 hours.

While inhabiting a host, Demetrius acts like an overbearing passenger and can't be targeted by any attack, spell, or other effect except those that turn Undead. The host remains in control of their actions, but they can hear Demetrius in their mind. The host also inherits the following flaw: "I will stop at nothing to seek vengeance on Darius."

As an action, a host inhabited by Demetrius can ask the spirit to cast one of the following spells (spell save DC 13), requiring no material components: \textit{Command}, \textit{Find Traps}, or \textit{Lesser Restoration}. Once Demetrius casts one of these spells, he can't cast that spell again until the next dawn.

\subparagraph{Roleplaying Demetrius}
Demetrius zealously recounts the tale of Darius's betrayal to his host. Demetrius despises the Cult of Zargon and mourns Cynidicea's continued decline. He knows the path to his brother's secret office (area B63), but he doesn't know about any of the traps along the way.

\subparagraph{Leaving a Host}
Demetrius never willingly leaves a host, but his spirit is forced back into the white robe if a host dies or is targeted by the \textit{Dispel Evil and Good} spell. Demetrius automatically fails his saving throw against any effect that turns Undead, which also forces his spirit back into the robe.

If the robe is destroyed or taken from the ziggurat while Demetrius's spirit resides within it, the spirit finds another object in the ziggurat to inhabit.

If Darius is slain, Demetrius passes to the afterlife, and his robe becomes a \textit{Cloak of Protection}.

\subsubsection{B47: Treasure Room}
\begin{DndReadAloud}{}
A subterranean breeze stirs in this room. A large, padlocked chest sits near the north wall. Behind it hangs a tapestry of a desert scene.

\end{DndReadAloud}

The tapestry behind the chest is a hostile \textbf{mimic}. Characters who have a passive Wisdom (Perception) score of 14 or higher notice that despite the breeze, the tapestry remains unnaturally still.

As an action, a character can open the padlocked chest with a successful DC 15 Dexterity (Sleight of Hand) check using thieves' tools.

\subparagraph{Secret Door}
The south wall holds a secret door to area B46. The outline of the door is visible here.

\subparagraph{Treasure}
The chest contains 330 gp, 600 sp, and a Potion of Greater Healing.

\subsubsection{B48: Kitchen and Dining Room}
\begin{DndReadAloud}{}
This abandoned kitchen and dining room contains stone cupboards, a sagging table, and dusty chairs. Near the table, an enormous snake coils around a bluish corpse wearing a wooden rabbit mask.

\end{DndReadAloud}

The \textbf{giant constrictor snake} strikes at any who interrupt its meal. The cupboards are empty, but the dead Cynidicean's robes contain a full waterskin and two days of rations.

\subsubsection{B49: Living Room}
\begin{DndReadAloud}{}
Faded wall hangings, cracked sculptures, and frayed divans decorate this ornate living room. Two rusty iron statues of armored warriors stand at attention, one on either side of the north door.

\end{DndReadAloud}

The two statues are suits of \textbf{animated armor}. Both statues awaken if either door to this room is opened. Created to protect Cynidiceans in this room, the statues are indifferent toward creatures wearing stylized Cynidicean masks and hostile toward all others. The statues never leave the room.

\subsubsection{B50: Consultation Room}
\begin{DndReadAloud}{}
Aged planks panel the walls of this antiquated chamber. A desk sits in the center of the room, surrounded by splintering chairs. Two stone statues of winged, lion-headed beasts flank the north door.

\end{DndReadAloud}

Two gargoyles guard this room. The gargoyles are hostile toward the characters but don't pursue them beyond the room.

\subparagraph{Treasure}
The desk's bottom drawer is locked. As an action, a character can open the drawer with a successful DC 13 Dexterity (Sleight of Hand) check using thieves' tools. It contains a \textit{Potion of Poison}. The other drawers contain only aged writing supplies and boring municipal documentation.

\subsubsection{B51: Pit Trap}
This intersection contains a 10-foot-wide, 10-foot-deep pit. The opening is covered by a wooden trapdoor painted to blend into the stone around it. Characters who have a passive Wisdom (Perception) score of 15 or higher notice the trapdoor.

When a weight greater than 50 pounds is placed on the trapdoor, the door swings downward. Each creature and object on it falls into the pit, taking 3 (1d6) bludgeoning damage from the fall and landing with the prone condition on a pile of brittle bones.

\subsubsection{B52: Abandoned Clerical Quarters}
This damp room's interior shimmers faintly, as if seen through clear water. Just inside the doorway is a hostile \textbf{gelatinous cube}. The cube is still digesting the remains of its last meal—a skeleton that appears to float in mid-air, chunks of gooey flesh clinging to its bones. The room is otherwise empty.

\subsubsection{B53: Priests' Quarters}
\begin{DndReadAloud}{}
This lived-in room contains three bunk beds and a chest at the foot of each bed. On one bed lies an elf man in bright yellow robes and a wooden sheep mask. An iron-bound chest rests at the foot of each bed.

\end{DndReadAloud}

The masked elf is a \textbf{doppelganger} that killed the Cynidicean it now impersonates. When it becomes aware of the characters, the doppelganger introduces itself as Kastrith, a friendly Cynidicean priest, and volunteers to help the party navigate the ziggurat. The doppelganger aims to split the party and take a character's place by any means necessary.

\subparagraph{Chests}
The chests are locked. As an action, a character can open a chest with a successful DC 14 Dexterity (Sleight of Hand) check using thieves' tools. "Kastrith" forbids the characters from opening the chests, claiming they belong to other priests. Inside each chest is a human corpse. Two of the corpses wear bright yellow robes and wooden masks, while the third has neither and looks exactly like the doppelganger's current disguise.

\subsubsection{B54: Storage Room}
Two ogres rummage for morsels in this storage room. Candles, incense, and other ceremonial items line its dusty shelves. When the characters enter, the ogres are eating a box of candles. Hungry and unsatisfied, the ogres demand food from all entrants and clobber any who refuse to give it to them.

\subparagraph{Treasure}
One of the ogres carries an oversized leather sack containing the corpse of an explorer, a \textit{Spell Scroll} of \textit{Comprehend Languages}, a \textit{dagger} clearly used as a toothpick, and 80 gp.

\subsubsection{B55: Hidden Guardroom}
\begin{DndReadAloud}{}
Two disheveled humans in stylized rat masks and red robes trimmed with rat fur crouch over a small chest, chattering and counting coins. In the northeast corner of the room, stairs lead downward.

\end{DndReadAloud}

Two \textbf{wererat} thieves in humanoid form divide their recent spoils. When they notice the characters, the wererats throw a cloth over the chest to conceal it and leap to their feet.

The wererats initially behave amicably toward the characters and claim to be Cynidiceans who lost their way in the monument—a ruse to buy time so they can rob the characters.

When the conversation stalls, the wererats spring on the characters. The wererats are cowardly; when reduced to half its hit points or fewer, a wererat transforms into a giant rat and attempts to flee.

\subparagraph{Stairs}
If you're running the extended adventure (see the "Extending the Adventure" section), the stairs lead down to a tunnel that exits into the Mouth of Zargon (area B64 of the expanded ziggurat). Otherwise, they end in a bricked-up wall.

\subparagraph{Treasure}
The chest contains 40 gp and four small agates worth 20 gp each.

\subsubsection{B56: Temple}
\begin{DndReadAloud}{}
This temple has fallen into disrepair. Ornate mosaics decorate its walls, and decrepit benches line its broad aisle. At the end of the aisle stands a dais with three toppled statues. Two women in elaborate clothing and copper fox masks inspect the rubble.

\end{DndReadAloud}

The women atop the dais are jackalweres disguised as human Cynidiceans. Sisters and charlatans, the two happened on the ziggurat and began scouring it for valuables. The siblings are insufferably vain but ultimately indifferent toward the characters.

When the sisters notice the characters, they bow and introduce themselves as Chana and Vana, potion sellers extraordinaire. The charlatans try to sell the characters fake potions—glittering bottles advertised as "rare and magical elixirs" with outlandish properties—for anything precious the characters can offer. A \textit{Detect Magic} spell reveals the potions as fakes. If confronted, the sisters flee, defending themselves as necessary. Failing that, they attempt to wangle their freedom through bribery and deceit.

\subparagraph{Treasure}
Chana and Vana each wear a copper fox mask worth 10 gp and two pieces of jewelry worth 150 gp apiece. The sisters' knapsack is full of counterfeit potions.

\subsubsection{B57: Entrance Hall}
\includegraphics[width=0.4\textwidth]{images/../5e.tools/img/adventure/QftIS/023-02-012.cynidicea-mosaic.png}
\begin{DndReadAloud}{}
Two rows of pillars, each carved in the likeness of a different past ruler, support the high ceiling of this grand entrance hall. A breathtaking mosaic spans the walls, depicting the rise and fall of an ancient civilization.

A thin layer of sand coats the floor; it grows thicker near the two enormous stone doors on the north wall. Two spoked platforms stand beside the doors.

\end{DndReadAloud}

The doors to this entrance hall have been sealed for centuries. Though the ziggurat is normally covered by mountains of sand, the dunes have temporarily shifted. These main doors can be opened from the inside if two creatures simultaneously turn the spoked platforms and both creatures succeed on a DC 15 Strength (Athletics) check. The doors then slowly swing inward, accompanied by the sound of grating stone and a deluge of sand.

If the characters leave the ziggurat, proceed to this adventure's conclusion.

\subparagraph{Mosaic}
The mosaic on the walls shows the history and fall of Cynidicea (detailed at the beginning of this adventure). It features the following events:


\begin{itemize}
\item Early Cynidiceans construct a great stone city.
\item King Alendria and Queen Zanobis reign, and the city prospers.
\item A ziggurat is raised in the monarchs' honor.
\item Workers dig beneath the ziggurat, unearthing a tentacled, one-eyed monster with a jagged horn.
\item The Cynidiceans bow to the creature in worship.
\item As the city falls to invaders, its citizens flee underground and don strange masks.
\end{itemize}

\subparagraph{Entrance to the Underground City}
If you're playing the expanded adventure (see the "Extending the Adventure" section), a plainly visible trapdoor in the northwest corner of the room leads down to the underground city. Otherwise, the trapdoor doesn't exist.

\subsubsection{B58: Acid Pit}
The floor of this sunken room is filled with a 6-foot-deep pool of vaporific acid. A creature that enters the pool for the first time on a turn or starts its turn there takes 22 (4d10) acid damage.

A corroded metal box bobs in the hissing pool; its padlocked lid protrudes above the liquid's surface.

\subparagraph{Treasure}
The box's padlock breaks if touched. Inside is a \textit{Spell Scroll} of \textit{Levitate} and 200 sp.

\subsubsection{B59: Storage Room}
\begin{DndReadAloud}{}
Wine racks line this storage room. Along the west wall stand casks and barrels, several of which have been smashed and drained. A hulking, feathered creature with ursine paws staggers amid the debris.

\end{DndReadAloud}

A disoriented \textbf{owlbear}, ravenously hungry and drunk on wine, has ransacked this storage room. If not fed immediately, the owlbear attacks. A character who feeds the owlbear and succeeds on a DC 14 Wisdom (Animal Handling) check calms the creature, changing its attitude to indifferent.

\subparagraph{Treasure}
The room contains four hundred bottles of Cynidicean wine, collectively worth 100 gp.

\subsubsection{B60: Lobby}
\begin{DndReadAloud}{}
Six humanoids sit motionless in this room, each clad in dark robes and a golden, horned mask with tentacles dangling from its chin. The figures gaze silently into a tarnished bronze basin filled with sludge.

\end{DndReadAloud}

The masked figures are six cultists of Zargon. Mesmerized by a visage of their overlord reflected in a bowl of slime, the cultists are indifferent toward the characters. If disturbed, the cultists become hostile.

Non-cultists who gaze into the basin glimpse the nightmarish face of Zargon and must succeed on a DC 13 Wisdom saving throw or have the frightened condition for 1 minute. A creature that fails the saving throw can repeat the save at the end of each of its turns, ending the effect on itself on a success.

If a character looks into the basin, the cultists snap out of their stupor and attack. They drag any defeated characters to Darius's office (area B63).

\subparagraph{Treasure}
The six cultists each wear a gold mask of Zargon worth 50 gp.

\subsubsection{B61: Party Room}
\begin{DndReadAloud}{}
Spongy fungi blanket the walls and floor this damp chamber. Amid clouds of spores, nine masked humanoids dance with invisible partners, intermittently chuckling and muttering in conversation. Branching up from the center of the room is a limbed, fungal pillar studded with yellow eyes.

\end{DndReadAloud}

Nine Cynidiceans (commoners) revel quietly in this fungal room. To escape their bleak existence, the Cynidiceans have partnered with a \textbf{myconid sovereign}. In exchange for the Cynidiceans tending its garden, the myconid facilitates group melds—vivid, shared dreams that provide entertainment and social interaction. As the partygoers meld, they are blissfully unaware of the characters.

When the characters enter, the sovereign invites them to participate in the meld. If they accept, the sovereign exudes a cloud of hallucinatory spores, drawing the characters into the Cynidiceans' shared dream. For 1 minute, the characters taste the former splendor of Cynidicea as the room around them appears to melt away, replaced by an ornate ballroom filled with joyful citizens in fanciful attire.

If the characters decline, the sovereign allows them to continue onward, asking only that they don't interrupt the meld. Attacking the sovereign shatters the illusion and angers the Cynidiceans, who join the sovereign in driving the characters out by force.

\subsubsection{B62: Gambling Room}
\begin{DndReadAloud}{}
Brightly robed gamblers in animal masks play games of chance in this weathered gambling den, exchanging coins and jewelry from an older era. Four cultists, clad in golden masks with tentacled chins, oversee them.

\end{DndReadAloud}

Some surviving Cynidiceans gamble the fortunes of their lost kingdom in this secluded room. Twelve masked commoners play games here, supervised by four cultists of Zargon.

If trouble breaks out in this room, Darius and two more cultists barge out of area B63, swords at the ready. The gamblers try to abstain from combat.

\includegraphics[width=0.4\textwidth]{images/../5e.tools/img/adventure/QftIS/024-02-013.zargon-mask.png}
\subparagraph{Dice Game}
The gamblers are playing an ancient dice game called Madarua's harvest. They welcome the characters to join in. The buy-in is 5 gp per game. The game's rules are as follows:


\begin{itemize}
\item Each player rolls three d6s.
\item A player who rolls three of a kind wins.
\item If no player rolls three of a kind, each player must either add another 5 gp to the pot or forfeit. The remaining players roll one or two of their d6s again, repeating these steps until someone wins or only one player remains, winning by default.
\item If multiple players simultaneously roll three of a kind, the player with the highest number on their dice wins. If the numbers are the same, it's a tie, and the winners divide the pot evenly.
\end{itemize}

After each game, a cultist approaches the table and demands a 1 gp "tax" from the winner. If the winner refuses, the cultists forcefully escort them through the secret door to Darius's private office (area B63).

\subparagraph{Treasure}
The room's occupants collectively hold 255 gp, 610 sp, three banded agates worth 25 gp each, and four pieces of jewelry worth 100 gp each. Additionally, the four cultists each wear a gold mask of Zargon worth 50 gp.

\subsubsection{B63: Darius's Office}
\begin{DndReadAloud}{}
A human cultist in a horned, tentacled mask lounges in this decadent office, counting treasure piled atop an aged desk. Two more cultists in simpler robes and identical masks stand behind him, daggers drawn.

\end{DndReadAloud}

Darius (chaotic evil, human \textbf{cult fanatic}) and two cultists of Zargon lurk in this office, counting coins confiscated from Cynidicean gamblers. They are hostile toward anyone not garbed in their cult's attire.

If a character is wearing Demetrius's white robe, Darius exclaims:

\begin{DndReadAloud}{}
"By the crooked horn of Zargon! Do my eyes deceive me? Has the ghost of my brother come to die again?"

\end{DndReadAloud}

Darius and his cohorts fight to the death.

\subparagraph{Treasure}
Atop Darius's desk are 50 gp, three gems worth 25 gp each, and two jeweled necklaces worth 100 gp each. Darius and the two cultists each wear a gold mask of Zargon worth 50 gp.

\section{Conclusion}
The adventure ends when the characters escape the ziggurat, likely via the main entrance in area B57. Depending on what brought the party to the ruins, they might eventually rejoin their desert caravan or return to their quest-giver to report their findings.

If your players wish to continue exploring the ruins of Cynidicea or the underground city beneath the ziggurat, consult the following section.

\section{Extending the Adventure}
"The Lost City" primarily explores the first five tiers of the ziggurat, but further adventures await those who continue their descent. This section provides the following two options to extend the adventure:


\begin{description}
\item[Expanded Ziggurat.]
This section adds a sixth tier to the ziggurat: the seeping lair of the elder evil Zargon and its cultists.
\item[The Underground City.]
This section presents an overview of the Cynidiceans' underground city and opportunities for adventure therein.
\end{description}

\subsection{Expanded Ziggurat}
In the days of Cynidicea's fall, the elder evil Zargon slithered up into the city and fed on the panicking masses. As its cult took over the city, Zargon returned to the tunnels below the ziggurat, content to feed on a steady supply of Cynidicean sacrifices brought by its worshipers. Zargon dwells there still, rising from its slimy pool to answer the growl of its stomach.

An unfinished tier occupied by Zargon and its loathsome servitors lies deep below the ziggurat. This tier is far more dangerous than those above it. A confrontation with Zargon could prove challenging for characters of all levels, though unseasoned adventurers determined to face the Returner in combat are almost certainly doomed to fail.

Characters killed by Zargon might be resurrected on the Infinite Staircase by Nafas (see chapter 1), who could ask them to repay him by completing another adventure in this book.

\subsubsection{Expanded Ziggurat Features}
Unless otherwise noted, this tier has the following features:


\begin{description}
\item[Ceilings.]
Corridors have 15-foot-high ceilings, and rooms have 20-foot-high ceilings.
\item[Doors.]
Translucent membranes curtain the entrances to each area. Creatures that pass through a membrane are coated in a mucilaginous glaze.
\item[Lighting.]
Bulbous growths cling to walls throughout. Their pulsing, sickly glow provides dim light.
\item[Slime.]
A film of sticky slime covers most surfaces.
\end{description}

\subsubsection{Random Encounters}
For each hour the party spends on this level, roll a d6. On a roll of 1, the characters encounter 1d4 + 1 cultists led by a \textbf{cult fanatic} looking for sacrifices. Each cult member wears a gold mask of Zargon worth 50 gp.

\subsubsection{Expanded Ziggurat Locations}
The following locations are keyed to map 2.4.

\includegraphics[width=0.4\textwidth]{images/../5e.tools/img/adventure/QftIS/025-map-2.04-expanded-ziggurat.png}
\includegraphics[width=0.4\textwidth]{images/../5e.tools/img/adventure/QftIS/026-map-2.04-expanded-ziggurat-player.png}
\subsubsection{B64: Mouth of Zargon}
\begin{DndReadAloud}{}
Two holes draped in sagging blankets of slime punctuate the floor of this unfinished chamber. The walls, riddled with oozing cracks, bear murals of humanoids in mourning. Rusty mining tools lie strewn about.

\end{DndReadAloud}

A \textbf{black pudding} creeps along the ceiling. It attempts to plop onto one of the characters and drag them into one of the slimy chutes on the floor.

\subparagraph{Slime Chutes}
Two 10-foot-wide, 40-foot-long tunnels, each covered by a thin web of slime, lead down to areas B66 and B67. A creature that moves over a pit must succeed on a DC 15 Dexterity saving throw or slide down the chute, landing with the prone condition in an unoccupied space within 10 feet of the chute's exit. A creature can choose to fail this saving throw. Clambering back up a chute requires a successful DC 15 Strength (Athletics) check.

\subsubsection{B65: Ritual Chamber}
\begin{DndReadAloud}{}
Eleven robed cultists chant within this ritual chamber. Each wears a horned, golden mask with four metallic tentacles sprouting from its chin. The cultists gather around a black marble altar supported by six sculpted tentacles. A twelfth cultist convulses atop the altar.

\end{DndReadAloud}

Twelve cultists perform a wicked ritual in this room. The cultists are hostile toward intruders but continue intoning even in combat—only death can separate the Zargonites from their goal. The cultist on the pedestal refrains from joining combat until the ritual is complete; if they are killed, another cultist takes their place on the altar.

The ritual continues for 3 rounds or until all cultists have been dispatched. A \textit{Silence} spell that envelops all the cultists also stops the ritual.

If one or more cultists are still chanting at the end of round 2, the cultist on the altar (or closest to it) transforms into a \textbf{gibbering mouther} with all its hit points. Read or paraphrase the following text:

\begin{DndReadAloud}{}
"I have been chosen! Praise the Returner!" shouts the cultist on the altar with glee as they melt into an amorphous blob of chattering teeth and darting eyes.

\end{DndReadAloud}

\subparagraph{Treasure}
The twelve cultists each wear a gold mask of Zargon worth 50 gp.

\subsubsection{B66: Cultists' Quarters}
\begin{DndReadAloud}{}
Twelve uncomfortable bunk beds, sticky and freckled with mold, line the walls of these sleeping quarters. Beneath four of the beds hang throbbing cocoons of stringy slime. A chest rests at the foot of each bed.

\end{DndReadAloud}

Low-ranking members of the Cult of Zargon sleep in these communal chambers.

\subparagraph{Cocoons}
Four clammy cocoons, ready to burst, stick to the bottoms of four beds. If a creature disturbs a cocoon or moves within 10 feet of it, the cocoon bursts. Each creature within 10 feet of a bursting cocoon must succeed on a DC 13 Dexterity saving throw or take 9 (2d8) poison damage and have the poisoned condition for 1 minute.

Immediately after a cocoon bursts, a mass of squirming, barbed tentacles (use the \textbf{swarm of poisonous snakes} stat block) emerges in the nearest unoccupied space. The oily limbs attack the nearest creature.

\subparagraph{Treasure}
Three of the chests contain gold masks of Zargon worth 50 gp each. Another contains a \textit{Spell Scroll} of \textit{Evard's Black Tentacles}.

\subsubsection{B67: Belly of the Beast}
\begin{DndReadAloud}{}
Smooth, wet tissue coats the walls of this cave. Cankerous lumps protrude from the flesh, dribbling caustic fluid into a sickly-hued pool of gastric juices. A rough stone pathway cuts through the bile.

Six dwarves with white hair and gray skin use severely corroded pickaxes to excavate an inflamed section of the wall. The chamber gurgles in agitation.

\end{DndReadAloud}

Six \textbf{duergar} loyal to Zargon work to expand this fleshy chamber. When the characters enter, two of the duergar are arguing over the quality of their tools. The duergar are friendly to characters who wear the masks of Zargon's cultists, allowing them to pass; otherwise, the duergar greet intruders with hostility.

\subparagraph{Acid Pool}
A hissing, 3-foot-deep pool of digestive acid fills the area around the stone path. A creature that enters the pool for the first time on a turn or starts its turn there takes 22 (4d10) acid damage.

\subparagraph{Bile Spouts}
Four bile spouts, agitated by the duergar excavators, stud the walls. At initiative count 10 on each round (losing initiative ties), the spouts erupt with gouts of acid in a 60-foot line that is 10 feet wide. Each creature in that area must make a DC 14 Dexterity saving throw, taking 22 (4d10) acid damage on a failed save or half as much damage on a successful one.

\subsubsection{B68: Spawning Chamber}
\begin{DndReadAloud}{}
A twenty-foot-deep pit takes up most of the floor of this rough-hewn chamber. A stone staircase curves down the wall of the pit, which holds three enormous stone vats brimming with differently colored liquids.

Four cultists in dark purple robes pile armfuls of humanoid bones into each cistern.

\end{DndReadAloud}

Four cult fanatics deposit leftovers from sacrifices in area B69 into vats here in hopes of birthing Oozes to serve Zargon. The cultists are hostile toward intruders not garbed in their cult's attire.

If a cultist is slain, their companions rush to cast themselves into the vats, where they quickly dissolve and potentially spawn oozes. The cultists praise Zargon with their last breaths.

\subparagraph{Ooze Vats}
The three 10-foot-wide vats each hold the essence of a different variety of Ooze. A creature that enters a vat for the first time on a turn or starts its turn there takes 22 (4d10) acid damage. If a creature dies in a vat, a hostile ooze emerges from it 1 round later. The ooze has the same initiative count as the creature whose death spawned it. An ooze spawned this way is indifferent to the cultists but hostile to all other creatures.

The contents of the vats and the oozes they spawn are as follows:


\begin{description}
\item[Black Vat.]
A tarry, black substance bubbles within this vat, which can produce one \textbf{black pudding}.
\item[Clear Vat.]
This vat wobbles with a crystal-clear liquid that can produce up to two gelatinous cubes.
\item[Yellow Vat.]
The glistening, amber-hued jam within this vat can produce up to two ochre jellies.
\end{description}

\subparagraph{Treasure}
The four cultists each wear a gold mask of Zargon worth 50 gp.

\subsubsection{B69: Foul Chapel}
\begin{DndReadAloud}{}
Sconces flicker in slimy alcoves within this ruined chapel, casting an eerie glow on five rows of clammy stone pews. Green mucus carpets the center aisle, leading to a bloodstained altar on a rocky precipice. Behind the altar, the chapel opens to a dark cavern.

Two rusty cages are set into the chapel's north wall. An eye-shaped gong stands in the southeast corner.

Six cultists sit in the pews. Before them is a priest in a horned mask adorned with eight chin tentacles.

\end{DndReadAloud}

Cultists present sacrifices to the Returner in this unholy temple. When the characters enter, Zekrikitch—a high priest of Zargon transformed by the Returner's slime—is addressing six cultists. Zekrikitch uses the \textbf{mind flayer} stat block, though he can also speak Common. Four of the tentacles under the high priest's mask are alive, writhing subtly as he speaks.

Zekrikitch telepathically probes the minds of any unexpected guests. If the characters attempt to fool the cultists using disguises or magic, the high priest plays along while telepathically warning his allies, who attack as one.

If the characters are defeated here or elsewhere on this level, the cultists chain one unlucky character to the altar and lock the remaining characters in the cages. The cultists then ring the gong.

\subparagraph{Altar}
Two sets of iron \textit{manacles} (see the \textit{Player's Handbook}) are bolted to the altar. A creature chained to the altar has the restrained condition.

\subparagraph{Cages}
The cages are 10 feet wide and 20 feet long. Both are empty. As an action, a character can pick a cage's lock with a successful DC 15 Dexterity (Sleight of Hand) check using thieves' tools or force open the rusty door with a successful DC 20 Strength (Athletics) check.

\subparagraph{Gong}
Ringing this eye-shaped gong signals to Zargon that a fresh sacrifice is ready for consumption. It takes 1 minute for the tyrant to slither from area B70 to the altar. If no sacrifice awaits it, Zargon unleashes its wrath on the chapel's occupants.

\subparagraph{Slime Carpet}
A 10-foot-wide patch of green slime (see the \textit{Dungeon Master's Guide}) intersects the pews. The slime doesn't harm Zargon or its cultists.

\subparagraph{Tunnel}
Directly behind the altar, the chamber drops steeply to a tunnel leading to area B70. The floor of the tunnel is 20 feet below the altar.

\subparagraph{Treasure}
Zekrikitch carries the keys to the cages and the altar's restraints. He and the six cultists each wear a gold mask of Zargon worth 50 gp.

\subsubsection{B70: Lair of Zargon}
A lake of slime fills this ruined chamber, whose walls climb 40 feet to a cavernous ceiling. Froth and scum line the lake's edge.

Shortly after the characters enter, a deep, guttural voice echoes in their minds. At first, its words are unintelligible, but slowly the gibberish gives way to a question: "What worlds have \textit{you} destroyed?"

Give the party a moment to answer, then read:

\begin{DndReadAloud}{}
A towering abomination rises from the slime. Three long tentacles tipped with bladelike barbs writhe from each of its shoulders, and six more slither below, propelling it forward. Beneath the two-foot-long black horn jutting from its hideous head, the creature's single bulging, red eye gazes down at you.

\end{DndReadAloud}

The creature is \textbf{Zargon the Returner}, the deathless aberration responsible for the fall of Cynidicea. Although Zargon is ageless, it is no god. A cunning evil of an age beyond memory, Zargon has discovered feigning godhood begets it a glut of sacrifices from those who crave its power.

If confronted by characters of 10th level or lower, Zargon toys with the party. Mildly impressed by the characters' foolish resolve, the elder evil might offer to make them minions, lieutenants, or priests—but not before threatening to eat them. Zargon spares those who tremble before it, acknowledge it as a god, or make it an offering of at least 600 gp.

Zargon takes adventurers of 11th level or higher more seriously. Deprived of worthy challengers for centuries, the Returner dares the characters to claim its horn, as so many before them have tried.

In combat, Zargon fights to the death. The creature secretly delights at its own destruction, knowing it will regrow from its horn.

\subparagraph{Lake of Slime}
Between sacrifices, Zargon broods in its pool of slime. The pool is 30 feet deep, and visibility beneath the surface is limited to 10 feet.

\subparagraph{Treasure}
The treasures Zargon has gathered lie submerged in unlocked iron chests at the bottom of the lake. Its hoard contains the following rewards, which it never parts with willingly: 11,600 sp and 1,800 gp, a \textit{Cloak of Elvenkind}, a \textit{Deck of Illusions}, a \textit{+2 Shield}, and a \textit{Scimitar of Speed}.

\includegraphics[width=0.4\textwidth]{images/../5e.tools/img/adventure/QftIS/027-02-014.zargon-the-returner.png}
\subsection{The Underground City}
\includegraphics[width=0.4\textwidth]{images/../5e.tools/img/adventure/QftIS/028-map-2.05-underground-city.png}
\includegraphics[width=0.4\textwidth]{images/../5e.tools/img/adventure/QftIS/029-map-2.05-underground-city-player.png}
Beneath the ruins of Cynidicea lies an underground city in a vast cavern. By the time survivors fled here from their broken capital, the cultists of Zargon had already begun constructing the city and thriving in its darkness. In exchange for food, housing, and a modicum of protection, Zargon's priests demanded absolute obedience, which they received for a time.

The Cynidiceans silently endured for decades as Zargonites plucked living sacrifices from their ranks to appease their voracious overlord. However, decades of growing horrors have finally primed the people for rebellion.

The underground city is accessible via a trapdoor in area B57 of the ziggurat, but at your discretion, other entrances to it might exist.

\subsubsection{Life in the Underground City}
A deep freshwater lake lies in the center of the city, just beyond Zargon's polluting influence. The lake provides cool, clean water and fish—which it has aplenty—to the populace. The Cynidiceans also breed underground livestock—giant rodents, cave lizards, and the occasional carrion crawler.

Cynidiceans farm edible fungi in mildewy fields forested with mushrooms. Myconids roam these dense patches of overgrowth, which are more akin to fungal thickets than cultivated fields. These neighborly fungi help citizens transcend their dismal existence through communal melds in which they share dreams of a nobler era. Some Cynidiceans have grown to prefer these joyful melds over their daily lives, sinking into complacency. The priests of Zargon encourage this behavior—as long as the melds continue, the people are less likely to object to the cruelty the cult inflicts on them.

\subsubsection{Noteworthy Sites}
The underground city occupies a vast cavern connected to catacombs of the old kingdom and the ziggurat above. Squat communal buildings in forgotten architectural styles line its forlorn streets, which are patrolled by Zargon's cultists.

\subsubsection{Entrance to the Ziggurat}
This tunnel connects the underground city to the main entrance chamber of the ziggurat (area B57).

\subsubsection{Eye of Zargon}
Outside the city lies a caustic pit known as the Eye of Zargon: a crater with a sickly glow reminiscent of Zargon's cycloptic eye. Zargonites regard the site as sacred and perform rituals to the Returner there. One legend states that to truly kill Zargon, its horn must be cast into the pit, but some fear this might be a deliberate misdirection spread by the cult.

\subsubsection{Faction Strongholds}
While most of the city lives in fear of Zargon and its devotees, three growing factions—the faithful of Cynidicea's old gods—plot to overthrow the Cult of Zargon and restore their former kingdom. Each faction does so alone, as mistrust and longstanding rivalries prevent them from uniting against the cult.

The factions maintain the following strongholds in the underground city:


\begin{description}
\item[Stronghold of Gorm.]
The Guardians of Gorm's headquarters is a pair of crooked towers mosaicked with brass tiles. The towers rise high above the complex's walled perimeter.
\item[Stronghold of Madarua.]
Four guard towers pierce the air above the stronghold of the Warriors of Madarua like the spears of her fearless warriors. Day and night, the walled complex echoes with the clanging of steel from training combatants.
\item[Stronghold of Usamigaras.]
The Mages of Usamigaras have no formal stronghold. Though outwardly welcoming, the wily mages gather in reinforced houses connected by a web of tunnels beneath the streets, allowing them to confuse and ambush invaders.
\end{description}

\subsubsection{Goblin Cliff Dwellings}
On the lake's west shore, sheer cliffs rise to a plateau ruled by a gang of cave-dwelling goblinoid bandits. Seeking valuables, these goblinoids regularly descend the rickety scaffoldings and frayed rope ladders that cling to the cliffs, clashing with the Cynidiceans in the underground city. The goblinoid leader, a merciless \textbf{hobgoblin warlord} named Krod the Wretched, plans to conquer the underground city and loot its riches when the Cynidiceans reach the nadir of their sufferings.

\subsubsection{Island of Death}
A small, stony island pockmarked with caves and tree-sized fungi rises from the lake. In the center of the island stands a circle of ancient menhirs predating the first-laid stones of Cynidicea. No one knows who built these monoliths or for what purpose, but they're rumored to be a source of necromantic power. When the Cynidiceans first fled here, they buried their dead on this island, along with priceless treasures—but the island slowly became infested with Undead, and soon those who voyaged there ceased to return.

\subsubsection{Lower Catacombs}
An enormous stone slab seals the entrance to the lower catacombs. Locals warn against tampering with the barrier, which separates the underground city from terrible dangers lurking below.

\subsubsection{Temple of Zargon}
Shaped like Zargon's menacing black horn, the Temple of Zargon houses the Returner's most loyal priests. Zargonites defend the structure at all hours. A sizable wing contains prison cells for sacrificial captives, embittered rebels, and anyone else foolish enough to offend the one-eyed tyrant or its cultists. The Returner's judgment awaits all who enter.

\chapter{When a Star Falls}
For centuries, the Tower of the Heavens has stood silhouetted against the glittering night sky. The tower is a bastion of destiny where gifted sages peer into the future, doling out prophecies to visitors. When one such sage foresees a falling star, it sets in motion a series of events that threatens to put the power of prophecy into dangerous hands.

"When a Star Falls" is designed for four to six 4th-level characters.

\includegraphics[width=0.4\textwidth]{images/../5e.tools/img/adventure/QftIS/030-03-001.ch3-when-a-star-falls.png}
\section{Background}
According to local legend, the sages of the Tower of the Heavens can glimpse the future. Suppliants journey from distant regions to the tower, each bringing gold and a single question about their fate in hopes of receiving an answer from the sages.

Unknown to outsiders, the sages' talents come from their best-guarded secret: a series of ancient, magical tomes called the Books of Prophecy. The books chronicle events that have happened and those yet to pass; each cryptic page might detail a fate days, months, or years in the future. Shalfey, Elder Sage of the Tower of the Heavens, inherited the magic required to read these books from the line of sages who came before him—but this magic prevents him from turning more than one new page per day, allowing only a peek at the future. To keep the Books of Prophecy secret, the sages make a great show of the study and practice of astrology, pretending to divine the future from the stars.

\includegraphics[width=0.4\textwidth]{images/../5e.tools/img/adventure/QftIS/031-03-002.tower-of-the-heavens-door.png}
A few days ago, Shalfey turned a page of his current Book of Prophecy and found it blank. Worse still, he found he now could turn the rest of the pages, which were also blank, save for the last. There, in clear words, the book foretold the arrival of a falling star—and the existence of a second set of Books of Prophecy possessed by a group of svirfneblin smiths. The elder sage resolved to plot the star's course, recover it, and trade it to the svirfneblin for the second set of tomes.

In anticipation of the star's descent, Shalfey sent his most trusted tower hands—martial artists sworn to defend the sages and obey their every order—to meet with Derwyth, a fellow star-watcher near the star's estimated point of impact. In the tower hands' absence, Shalfey's power-hungry successor, the sage Piyarz, saw his mentor was unguarded and vulnerable. Piyarz and his mutinous faction of apprentices attempted a coup, but Shalfey fended them off and retreated to the safety of his sanctum, where he hides still. Foiled, Piyarz turned his attention to Shalfey's emissaries, hoping to find the star and steal the power of prophecy for himself.

\subsection{Using the Infinite Staircase}
If you're using Nafas as a patron, he summons the characters to the Censer of Dreams (detailed in chapter 1), where he recounts the following wish:

\begin{DndReadAloud}{}
"A prophetic sage by the name of Shalfey has uttered a hasty plea. He sent a group of robed sages into the mountains to intercept a falling star, but a creature imperils their journey. Seek these emissaries and help them bring the star to Shalfey when it falls. He dwells within a fortress known as the Tower of the Heavens."

\end{DndReadAloud}

Nafas then teleports the characters to a door along the Infinite Staircase that opens near the Tegefed Mountains. After the adventure, the characters can return to the staircase through the same portal.

\subsection{Adventure Hooks}
If you're not using Nafas as a group patron, consider the following ways to involve the characters in this adventure:


\begin{description}
\item[In Search of Prophecies.]
The characters seek to have their fates read at the Tower of the Heavens. Whether the prophecies are real remains to be seen, but the characters are curious or desperate enough to pay a visit. The characters might also be researching historical information lost to all but those who deal in the business of fate.
\item[Investigators.]
Dead travelers have been found in the hills south of the Tegefed Mountains—and the corpses reportedly have no visible wounds. The characters are hired by a mourning family to investigate these deaths and put a stop to them. The characters' encounter with the cause of these deaths—a strange predator called a \textbf{memory web}—sparks the beginning of their adventure, as the creature's demise reveals a greater plot.
\end{description}

\subsection{Setting the Adventure}
This adventure takes place in a region with a moderate climate, rolling hills, and winding rivers. The Tower of the Heavens lies far enough from major settlements that it sees only a few visitors at a time. Consider the following suggestions:


\begin{description}
\item[Eberron.]
The Tower of the Heavens could be the local name for the Starpeaks Observatory in Northern Aundair, with the Tegefed Mountains becoming the Starpeaks.
\item[Forgotten Realms.]
The Tegefed Mountains could be in the eastern half of the Small Teeth mountain range, located in the region of Amn.
\item[Greyhawk.]
You could place the Tegefed Mountains near the Crystalmist Mountains, where the Kingdom of Keoland borders the Yeomanry.
\end{description}

\includegraphics[width=0.4\textwidth]{images/../5e.tools/img/adventure/QftIS/032-03-003.when-a-star-falls-original.png}
\begin{DndReadAloud}{About the Original}
Produced by the UK branch of TSR in 1984, \textit{When a Star Falls} was written by Graeme Morris. The adventure is remembered for its intricate plot and unique adventure hook, which spurred characters to action by imprinting their minds with memories from a creature known as a memory web.

\end{DndReadAloud}

\section{Adventure Summary}
The adventure begins when the party stumbles on the corpses of Shalfey's emissaries. The envoys were killed by a memory web—a living web that feasts on memories. When the characters defeat the memory web, they inherit the emissaries' memories, imprinting Shalfey's mission firmly in their minds. Shortly thereafter, the characters witness the prophesied star falling to the Tegefed Mountains. Later that night, the shadowy assassin Sion ambushes the party, seeking intel for his employer, Piyarz.

From there, the characters have two leads. They can investigate the Tower of the Heavens, or they can track down the messengers' intended contact, the druid Derwyth. The journey might lead characters to a cave of derro raiders, a lakeside community of hunters and giant beavers, or both.

Once the characters recover the star and defeat Piyarz, Shalfey sends them to a svirfneblin forge, where a group of smiths await the star. In exchange, the smiths give the characters a second set of Books of Prophecy for Shalfey, who rewards the characters on receipt of the tomes.

\subsection{Character Advancement}
If you want to use story-based level advancement, the characters receive experience points for achieving the following milestones rather than defeating monsters:


\begin{description}
\item[Obtaining the Star or Saving Shalfey.]
When the characters either obtain the star from the derro lair or rescue Shalfey (whichever occurs first), everyone in the party gains 1 level. This advancement occurs only once.
\item[Returning the Books of Prophecy.]
The characters gain 1 level when they deliver the second set of Books of Prophecy to Shalfey.
\end{description}

If you follow this method, the characters should reach 6th level by the adventure's conclusion.

\section{Tegefed Mountains}
This adventure takes place in the Tegefed Mountains region. The mountains are lofty and rugged, with snow on the upper peaks even in the warmer months. The surrounding hills and moors are dotted with sparse trees and flowers, and settlements are few and far between.

Map 3.1 depicts the stretch of the Tegefed Mountains where the adventure occurs. Any hex on the map that doesn't contain a river is \textit{difficult terrain} for overland travel.

\includegraphics[width=0.4\textwidth]{images/../5e.tools/img/adventure/QftIS/033-map-3.01-tegefed-mountains.png}
\includegraphics[width=0.4\textwidth]{images/../5e.tools/img/adventure/QftIS/034-map-3.01-tegefed-mountains-player.png}
\subsection{Random Encounters}
Random encounters can bring the region to life. An encounter might occur while the characters are traveling from one location to another or whenever the DM feels appropriate. To determine what the characters find, roll on the Random Tegefed Mountains Encounters table, rerolling duplicates.

\begin{DndTable}[header=Random Tegefed Mountains Encounters]{cX}

d8 & Encounter\\
1 & Three griffons hunt the characters. A griffon egg in a nearby nest can fetch a price of up to 300 gp if sold to the right buyer.\\
2 & Two displacer beasts quietly stalk the characters.\\
3 & An \textbf{ogre} trains five boars to hunt wild animals. On noticing the characters, the ogre sets the pet boars loose on them instead.\\
4 & Five bugbears ambush the party, demanding the characters give them 100 gp or face certain death. Each bugbear carries 3 gp in a coin pouch.\\
5 & Ten giant bats swoop in and attack the characters. The bats focus their attacks on the character who makes the most noise.\\
6 & The characters hear distant howls. Two dire wolves and five wolves descend on them 1 minute later.\\
7 & A \textbf{troll} drags a large sack containing the carcass of a giant goat, seeking a good place to eat it. If the troll notices the characters, the troll drops the sack and attacks, deeming them a much tastier meal.\\
8 & The characters encounter an injured hunter in the wilderness who calls out for help. The hunter is actually a hungry \textbf{jackalwere} bandit in human form. Five more jackalweres lie in ambush nearby.\\
\end{DndTable}

\section{Death on the Moors}
The adventure begins as the characters traverse the moors of the Tegefed Mountains. This journey leads through a beautiful tableau of rolling hills and snowy peaks—but the pleasant excursion is soon marred by a harrowing encounter.

\subsection{Memory Web}
As the characters travel through the region, read or paraphrase the following text:

\begin{DndReadAloud}{}
The moorland of the Tegefed Mountains sprawls around you. Beneath lofty peaks, sparse trees and splashes of vivid flowers dot the green hills.

As you reach the top of yet another rise—this one blanketed with fragrant heather—the scenery is interrupted by a giant web that stretches out just a couple dozen feet from you. Four unmoving, white-robed figures lie enmeshed within it.

\end{DndReadAloud}

The mass of strands is a motionless \textbf{memory web}. The memory web attempts to catch the characters by surprise, launching itself at them to feed on their memories. If the characters ignore the bodies and continue on their way, the memory web follows stealthily and strikes a few minutes later.

When the memory web dies, its Memory Flood trait releases the messengers' memories it absorbed (see the "Flood of Memories" section).

\subsubsection{Shalfey's Emissaries}
The robed figures are the corpses of four tower hands whom Shalfey sent to track down the star. The messengers' minds were drained by the memory web before death. As a result, the \textit{Speak with Dead} spell and similar magic yield only confused, meaningless responses.

\subparagraph{Bestiary for Derwyth}
One of the corpses clutches a thin leather-bound book—an illustrated bestiary of creatures local to the mountains, which is worth 30 gp. The words "for Derwyth" are neatly written on its first page. A character who succeeds on a DC 15 Intelligence (History) check knows Derwyth is a druid who lives nearby in Cernant Valley (detailed later in this adventure). Shalfey's emissaries had planned to take the book to her as payment for information about the star.

\subparagraph{Map of the Region}
Another corpse clutches a hastily drawn map of the Tegefed Mountains, marked with the locations of Cernant Valley and the Tower of the Heavens (see map 3.1).

\subsection{Flood of Memories}
When the memory web dies, a barrage of memories floods the characters' minds. Most of these fleeting images fade instantly, but the following scenes appear in quick succession and lodge themselves firmly in each character's memory.

\subsubsection{Memory 1: Initiation Day}
\begin{DndReadAloud}{}
The world around you blurs for a moment, and everything fades to black. When your senses return, you find yourself in a hazy, dreamlike scene.

Bookshelves surround you, and a large, shimmering sphere levitates on the other side of the room. In front of the sphere stands a bearded human sage, his arms covered with tattoos of constellations. "Remember it well," he says slowly and firmly. "Fair is the day—\textit{but fairer the jewels of the evening}."

\end{DndReadAloud}

This memory takes place a few years ago. It depicts Shalfey initiating a group of tower hands into his private guard. Shalfey is teaching them the secret phrase used among his closest confidants. This phrase might come in handy when the characters meet Hadley, who runs the ferry at the Tower of the Heavens (detailed later in this adventure).

\subsubsection{Memory 2: The Star Foretold}
\begin{DndReadAloud}{}
The scene shifts. The sage has aged by several years. Distraught, he paces in a starlit observatory.

"The book speaks of a star falling to earth—a black rock, seemingly worthless but seared by the heavens," he explains. "Only it can help us. We must find it, and soon, or evil will seize the power of prophecy."

\end{DndReadAloud}

Four days ago, Shalfey recounted his findings from the final page of the Books of Prophecy.

\subsubsection{Memory 3: Shalfey's Orders}
\begin{DndReadAloud}{}
Another memory shifts into your mind. The same old man gazes at you with desperate intensity. "Go, seek Derwyth, a fellow star-watcher," he urges. "Make haste to her home in Cernant Valley."

He places a thin leather book in your hands, which are not your own. "Give her this bestiary as payment. When the star falls, only she can tell you where it landed. I implore you, do not return without the star."

\end{DndReadAloud}

Three days ago, Shalfey tasked his tower hands with seeking out Derwyth. They died along the way.

\subsection{Falling Star}
After the characters have a moment to process the memories they've received, they witness the foretold falling star. Read or paraphrase the following text:

\begin{DndReadAloud}{}
Your surroundings momentarily brighten as a brilliant light streaks across the heavens and plummets to the earth. As the falling star lands in the mountains ahead, a distant rumble accompanies the impact, causing flocks of birds to scatter from the hills.

\end{DndReadAloud}

The falling star's glittering trail briefly lingers in the air. Though the characters can't determine exactly where the star landed, they can tell from the map that it fell near Cernant Valley, where Shalfey's emissaries were headed. If the characters didn't find the map, the memory flood has given them the equivalent knowledge of the region.

\subsection{Next Steps}
The characters now know the locations of two points of interest: the Tower of the Heavens and Cernant Valley. Where they choose to go next is up to them. Either way, they receive an unwelcome visit in the night (see the following section).

\section{Piyarz's Shadow}
Shalfey's messengers were followed by an assassin named Sion—a shadowy mage molded by darkness. With them dead, Sion stalks the characters instead, waiting for the cover of night to strike.

\subsection{Stalked by a Shadow}
When Piyarz seized the Tower of the Heavens in his coup, he interrogated Shalfey's scribes to discern where the elder sage sent his tower hands. Eventually the ruthless sage discovered the reason for the expedition: recovering a fallen star to acquire new Books of Prophecy. Piyarz then killed the scribes, leaving no allies of Shalfey alive in the tower.

Coveting the power of prophecy for himself, Piyarz dispatched Sion, his most trusted and terrible servant, in pursuit of Shalfey's expedition.

\subsection{Night Ambush}
\textbf{Sion} ambushes the characters at night. If the characters decided to rest after facing the memory web, Sion interrupts their repose.

Sion is accompanied by a hound and an oversized raven, both made of pure shadow (both creatures use the \textbf{shadow} stat block, but the raven has a flying speed of 40 feet). When a companion drops to 0 hit points, it vanishes in a cloud of shadow. The companions also vanish when Sion dies.

Sion attempts to isolate a character—especially anyone taking watch—and threaten them into giving him information about the fallen star for his employer, whom he refuses to name. Sion is most interested in what knowledge the characters gleaned from the slain memory web. If the characters took either the bestiary or the map from the emissaries' corpses, Sion demands those as well. Otherwise, he has those items on his person.

When he's satisfied with the characters' answers or deems the party uncooperative, Sion attacks. He and his sinister pets fight to the death, unflinchingly loyal to Piyarz.

\includegraphics[width=0.4\textwidth]{images/../5e.tools/img/adventure/QftIS/035-03-004.sion-lurking.png}
\section{Cernant Valley}
Shalfey instructed his emissaries to visit Derwyth, an elf druid who lives near the Tegefed Mountains. According to the Books of Prophecy, she would know precisely where to find the star once it fell.

Derwyth lives in Cernant Valley—a forested basin at the foot of the Tegefed Mountains—along with many animals, most of which are small woodland creatures. These animals act as Derwyth's eyes and ears, quickly bringing her news of events in the wood, particularly the arrival and actions of strangers. The creatures also carry messages to more powerful Beasts in the valley that sometimes come to Derwyth's aid.

\subsection{Arrival at Cernant Valley}
When the party enters Cernant Valley, read or paraphrase the following text:

\begin{DndReadAloud}{}
Two giant eagles soar down to meet you, landing in nearby trees that creak gently beneath the eagles' weight. "Welcome, ground-walkers," squawks one of the eagles. The other continues: "This is the home of Derwyth. If you come in peace, follow us—if not, leave now."

\end{DndReadAloud}

Avri and Ivo, two giant eagles who speak Common, welcome the characters to Derwyth's wood. The eagles are loyal to Derwyth and friendly to those who regard her home with respect. They spend their days soaring over the valley and watching for intruders.

The eagles offer to lead the characters to Derwyth's homestead. If the characters decline or offend Avri and Ivo, the eagles tell them to leave. If combat ensues, three \textbf{elk} join the fray on initiative count 20 of the third round (losing initiative ties); the elk aid the eagles however they can. Attacking the eagles sours Derwyth's opinion of the characters, who fail her test (see the "Derwyth's Test" section later in this adventure).

\subsubsection{Going It Alone}
If the characters refuse the eagles' guidance, they must find their own way to Derwyth's homestead. The character leading the party must make a DC 15 Wisdom (Survival) check to navigate the valley's tangled thickets. On a successful check, the characters reach the homestead without issue. On a failed check, they eventually arrive at the homestead, but the journey is exhaustingly circuitous; each character must succeed on a DC 10 Constitution saving throw or gain 1 level of exhaustion.

\subsection{Derwyth's Homestead Features}
Derwyth and many of her animal companions dwell in a small collection of thatch-roofed buildings at the heart of Cernant Valley. The homestead has the following features:


\begin{description}
\item[Buildings.]
The buildings have thick walls of compacted earth and grass, and the gabled thatch roofs are supported by beams of riven oak. The ceilings are 10 feet tall. There are no windows.
\item[Doors.]
The doors are made of thick oak planks. As an action, a character can force open a locked door with a successful DC 15 Strength (Athletics) check or pick the lock with a successful DC 12 Dexterity (Sleight of Hand) check using thieves' tools. Derwyth carries a key to each locked door.
\item[Lighting.]
Rooms are brightly lit by baskets of fireflies hanging from the roof beams.
\end{description}

\subsection{Derwyth's Homestead Locations}
The following locations are keyed to map 3.2. The "Trial of Trust" section later in this chapter explains how events might unfold in the homestead.

\includegraphics[width=0.4\textwidth]{images/../5e.tools/img/adventure/QftIS/036-map-3.02-derwyths-homestead.png}
\includegraphics[width=0.4\textwidth]{images/../5e.tools/img/adventure/QftIS/037-map-3.02-derwyths-homestead-player.png}
\subsubsection{H1: Yard}
Whether led by the eagles or by forging a path through the thicket, the characters likely enter the homestead via the yard. When they arrive, read or paraphrase the following text:

\begin{DndReadAloud}{}
This serene homestead is perfectly quiet, save for a babbling brook and the gentle chirping of birds. A rustic wooden fence encloses three simple buildings with walls of compacted earth and grass. Birch trees stand tall in the center of the yard.

\end{DndReadAloud}

If Avri and Ivo are present, the eagles tell the characters Derwyth is currently out gathering food, but she should be back soon; the characters can wait in the yard until then. After advising the characters to behave themselves, the eagles take flight.

\subparagraph{Derwyth}
In truth, Derwyth is disguised as a birch tree in the center of the yard, watching how the characters behave. The "Trial of Trust" section later in this chapter details how the druid observes and interacts with the characters.

\textbf{Derwyth} is a neutral good, elf \textbf{druid} who has darkvision out to 60 feet, speaks Common and Elvish, and has the following additional trait:

\begin{DndReadAloud}{}
\subparagraph{Tree Shape (2/Day)}
Over the course of 1 minute, Derwyth can magically transform into a Huge or smaller tree and remain in that form for 24 hours or until she ends this transformation early (no action required). Her equipment melds into her new form. While in this form, her Armor Class is 16, she has the incapacitated condition, and she can't move or speak.

\end{DndReadAloud}

\subsubsection{H2: Antechamber}
This room's south door is locked. If the characters enter, read or paraphrase the following text:

\begin{DndReadAloud}{}
Old cloaks and walking sticks hang from hooks on the wall of this entry room. A pair of muddy boots lie on the floor, and a dusty broom stands in the far corner.

\end{DndReadAloud}

The items here all belong to Derwyth. There are nine walking sticks, each lovingly carved from a different kind of wood.

\subsubsection{H3: Hall}
\begin{DndReadAloud}{}
Six wooden posts support the roof of this simple hall. Sturdy benches of hardened clay line the walls, and a long hearth spans the center of the earthen floor. An ornate harp hangs from a hook on one of the posts.

\end{DndReadAloud}

Derwyth entertains guests in this dining hall.

\subparagraph{Treasure}
The harp is made of plum wood, inlaid with fine mithral wire, and worth 300 gp.

\subsubsection{H4: Derwyth's Room}
This room's door is locked. If the characters enter, read or paraphrase the following text:

\begin{DndReadAloud}{}
This small, plain chamber contains a round carpet and a long chest. An alcove in the far wall displays a small statue of a leaping stag.

\end{DndReadAloud}

This is Derwyth's private resting chamber.

\subparagraph{Treasure}
The stag statue in the alcove is made of fine wood and worth 160 gp. The chest is unlocked; it contains a bag of 50 gp; three Potion of Healing in ceramic flagons; gardening implements; and worn robes, cloaks, tunics, boots, and sandals.

\subsubsection{H5: Library and Workshop}
This room's south door is locked. If the characters enter, read or paraphrase the following text:

\begin{DndReadAloud}{}
Wooden shelves along the north wall of this room hold a jumble of books and scrolls, many of which have tumbled to the floor. The table in the center of the room bears scrolls with complex astrological notes in addition to quills, ink, and other stationery.

\end{DndReadAloud}

Derwyth stores her texts in this study. In all, there are eighty-four books and over two hundred scrolls—a comprehensive library of the druid's fields of interest: mainly animals, plants, and astronomy.

The scrolls on the table include Derwyth's notes concerning the fallen star. Though the party can easily recognize these notes, only someone versed in astronomy—such as Derwyth, Shalfey, or a character proficient with navigator's tools—can use them to determine the landing point of the star, a process that takes 2 hours.

\subsubsection{H6: Goat Barn}
\includegraphics[width=0.4\textwidth]{images/../5e.tools/img/adventure/QftIS/038-03-005.derwyth-the-druid.png}
\begin{DndReadAloud}{}
This barn smells strongly of livestock. Here, three horse-sized goats quietly munch on hay. A stool sits on the straw-covered floor near some large milking pails and vessels for making cheese.

\end{DndReadAloud}

Three giant goats live here. If it's clear Derwyth is in danger, the goats come to her aid, but they're otherwise friendly toward strangers and usually don't mind being milked.

A character can attempt to milk a goat by making a DC 10 Wisdom (Animal Handling) check. On a successful check, the goat produces one bucket's worth of milk. On a failed check, the goat delivers a swift, agitated kick, and the milker takes 2 (1d4) bludgeoning damage and has the prone condition.

\subsubsection{H7: Food Storage}
\begin{DndReadAloud}{}
A central post supports the conical roof of this circular hut. Piled around the walls are dozens of sacks, casks, boxes, and jars. A saber-toothed tiger lounges on one of the casks.

\end{DndReadAloud}

These containers hold enough cheese, dried meat, flour, fruits, nuts, vegetables, and wine to sustain one person for half a year.

Derwyth's \textbf{saber-toothed tiger} lazes here during the druid's test (see the "Trial of Trust" section). The tiger is initially indifferent toward the party and doesn't stop characters from stealing food from the hut; however, it leaps to Derwyth's aid if the druid is in danger.

If the characters do nothing to anger the tiger, they can safely pet it with a DC 10 Wisdom (Animal Handling) check. On a successful check, the tiger rolls onto its back and chuffs happily. On a failed check, the tiger growls threateningly, then attacks if pestered further.

\subsection{Trial of Trust}
Derwyth received advance notice of the characters' arrival from woodland creatures in Cernant Valley. She has used her Tree Shape trait to disguise herself as a tree in the yard (area H1), from where she discreetly observes the party.

\subsubsection{Derwyth's Test}
Derwyth remains in tree form for 10 minutes after the characters arrive. If the characters don't break her trust by harming her giant eagle friends, entering locked rooms, harming her animals, or stealing her possessions, they pass her test. Otherwise, they fail. At the end of this time, Derwyth reveals herself.

When Derwyth abandons her disguise, read:

\begin{DndReadAloud}{}
One of the birch trees in the yard begins to twist and creak, shrinking and uprooting itself. In seconds, an elf woman with white hair stands where the tree once did. She wears a brown cloak clasped with a sunflower brooch and sports a sturdy, wooden prosthetic leg.

\end{DndReadAloud}

Derwyth's attitude toward the characters depends on whether they passed or failed her test.

\subparagraph{Passing Derwyth's Test}
If the characters passed Derwyth's test, she commends their behavior and invites them into her dining hall (area H3), where she offers them bowls of vegetables, fruits, and nuts, as well as wooden vessels of water or wine. Until she's fully convinced of the characters' intentions, however, Derwyth actively discourages characters from snooping around her homestead.

During their meal, Derwyth asks the characters their business and listens intently. Though she has never met Shalfey, she knows of the elder sage and thinks well of him. If the characters detail their memories from the memory web, Derwyth realizes Shalfey's messengers intended for her to determine the star's landing place. She tells the characters she saw the star fall and can approximate its point of impact (see the "Derwyth's Star Calculations" section).

\subparagraph{Failing Derwyth's Test}
If the characters failed Derwyth's test or they otherwise break her trust, she orders them to leave. A character who offers an earnest apology can convince Derwyth to let them stay with a successful DC 15 Charisma (Persuasion) check. If the characters don't make amends and refuse to leave the homestead, Derwyth attacks, calling the \textbf{saber-toothed tiger} in area H7 and the three giant goats in area H6 to back her up. If Derwyth is defeated, news of the battle quickly spreads to Beasts throughout Cernant Valley, who regard the characters with hostility.

\subsubsection{Derwyth's Star Calculations}
Once Derwyth trusts the characters, she offers to calculate where the star landed in exchange for 200 gp, but she gladly accepts the bestiary carried by Shalfey's emissaries in lieu of payment. If the characters refuse to compensate Derwyth, she responds as outlined in the "Failing Derwyth's Test" section.

Once paid, Derwyth retires to her library (area H5) to perform her calculations. If the characters want to be useful in the meantime, Derwyth suggests they milk the goats in the barn (see area H6).

\subparagraph{What Derwyth Knows}
After 2 hours, Derwyth finishes her calculations. She tells the characters the shooting star fell in Therno Pass, a small region of the Tegefed Mountains inhabited by Aberrations called derro. All she knows about this group of derro is that they live underground, shun the light of day, and are feared for their ruthless attacks on nearby settlements. Derwyth marks the star's landing point on the characters' map; it corresponds with the location of the derro lair on map 3.1.

\subparagraph{Bidding Farewell}
After sharing this information, Derwyth offers the characters the opportunity to rest in her homestead before the party departs. If the characters made a good impression on her, she also gives them the three Potion of Healing from the chest in her room. The giant eagles Avri and Ivo then escort the characters out of Cernant Valley.

\section{Derro Lair}
A merciless faction of derro raiders has established a lair about a dozen miles from Therno Lake. From this outpost, the raiders pillage the surrounding countryside, frequently accosting hunters and traveling merchants. These forays are inspired not only by the raiders' greed but also by their sinister habit of dragging their victims' corpses back to the outpost to conscript them as zombie guards and workers.

When the star fell, it crashed into this outpost. The impact collapsed several tunnels, cutting the derro off from both the surface and the Underdark below. The collapse also destroyed the derro's idol of their god, Diirinka. Seeing this destruction as an omen from Diirinka, the derro dug westward to free themselves from the cave-in and snatched the star from the crater where it landed.

The derro believe the star is an object of immense power. Determined to take the star to the Underdark to gain leverage and status, the derro have been digging a new path to the subterranean depths.

\subsection{Point of Impact}
If the characters discovered the location of the fallen star from Derwyth's calculations, they have no trouble traveling directly to the point of impact. Likewise, if the characters already met the hunters of Therno Lake (detailed later in this adventure), those hunters might have directed them to the derro lair.

If the characters failed to secure such information, they can find the star's point of impact by journeying to southern Therno Pass and succeeding on a DC 20 Wisdom (Survival) check. Characters who make this check from a suitably high vantage point, such as a mountain peak, have advantage on the check.

\subsubsection{Impact Crater}
The star created a crater on the side of the pass, directly above the derro outpost. When the characters arrive, read or paraphrase the following text:

\begin{DndReadAloud}{}
The crater is nearly fifty feet across and half as deep. Singed trees, their trunks bent by the star's impact and their leaves blown off entirely, lean away from the crater. Rocks litter the hillside in all directions, and much of the crater is full of rubble. A shallow ditch, recently excavated, lies at the bottom of the basin.

\end{DndReadAloud}

The derro dug this 10-foot-deep ditch to get to the star. Only a broken pickaxe lies at the bottom of the ditch; the star is long gone.

Characters who survey the surrounding area easily notice two openings lower on the mountainside: a north-facing entrance of worked stone (to area D1) and a west-facing natural cave (to area D2).

\subsection{Derro Lair Features}
The derro lair has the following features:


\begin{description}
\item[Ceilings.]
The lair's rough-hewn chambers have low, 6-foot ceilings. The natural caverns that connect the chambers have ceilings as tall as the caves are wide.
\item[Doors.]
Unless specified otherwise, doors are made of sturdy planks and unlocked.
\item[Lighting.]
The derro lair isn't illuminated. The derro rely on darkvision to see. Area descriptions assume the characters have a light source or other means of seeing in the dark.
\item[Loose Rubble.]
Much of the central lair is blocked by rubble that would require days to clear, but some areas have loose rubble that can be cleared more quickly. A character who takes 1 minute to clear loose rubble and succeeds on a DC 18 Strength (Athletics) check removes a 5-foot cube of rubble.
\end{description}

\subsection{Derro Lair Locations}
The following locations are keyed to map 3.3.

\includegraphics[width=0.4\textwidth]{images/../5e.tools/img/adventure/QftIS/039-map-3.03-derro-lair.png}
\includegraphics[width=0.4\textwidth]{images/../5e.tools/img/adventure/QftIS/040-map-3.03-derro-lair-player.png}
\subsubsection{D1: Guardroom}
\begin{DndReadAloud}{}
The far wall of this room and much of the ceiling have collapsed. The corpses of five small dwarf-like creatures with gray skin and wild, white hair sprawl on the floor, partially buried under fallen rock.

\end{DndReadAloud}

This guardroom was once the sole entrance to the derro lair. When the star struck, collapsing masonry killed this chamber's occupants and cut off the room from the rest of the lair. Preoccupied with the fallen star, the surviving derro haven't retrieved their dead cohorts.

\subparagraph{Treasure}
Characters can pull the derro corpses from the rubble with ease. The corpses carry a total of 60 gp between them.

\subsubsection{D2: Outer Cave}
\begin{DndReadAloud}{}
Water flows from two pools on the east edge of this cave, combining into a narrow stream that trickles out the entrance.

\end{DndReadAloud}

Five derro raiders and one \textbf{derro apprentice} hide behind the pillars in this room. The derro entered the cave through the southeast pool and are dripping wet. They're uninterested in negotiation and ambush any who enter their cavern. A character who examines the room and succeeds on a DC 13 Wisdom (Perception) check notices wet footprints along the cavern floor before the derro attack.

\subparagraph{Pools}
Both pools are 8 feet deep at their deepest point. A short underwater tunnel in the southeast pool connects this area to the inner cave (area D3).

\subparagraph{Treasure}
The derro apprentice carries a decorative silver \textit{dagger} with an ivory handle, worth 50 gp, along with a belt pouch containing 20 gp.

\subsubsection{D3: Inner Cave}
\begin{DndReadAloud}{}
A pool lies at the north end of this natural cavern. On the cavern's east wall, a rough-hewn tunnel opens five feet above ground level. Sounds of digging emanate from the tunnel. Beneath the opening is a large mound of loose soil and rocks.

\end{DndReadAloud}

When the star fell, it collapsed the lair's exits to the surface and the Underdark. The trapped derro dug a tunnel to this cavern from the temple (area D4), reconnecting the temple to the outside world.

Shortly after the party arrives, a \textbf{zombie} emerges from the east tunnel and empties a basket of small rocks onto the pile below. The zombie is indifferent toward the characters and ignores them unless they try to get its attention, at which point it attacks.

\subsubsection{D4: Temple}
\includegraphics[width=0.4\textwidth]{images/../5e.tools/img/adventure/QftIS/041-03-007.derro-raiders.png}
\begin{DndReadAloud}{}
The walls of this rough-hewn chamber have been painted with broad, jagged lines of black and white. On the west side of the room, fragments of a black marble statue lie shattered on the floor. In the northeast corner of the room, four derro use hammers to clear rubble from a blocked passage. Four zombies accompany them, clawing mindlessly at the rubble.

\end{DndReadAloud}

Four derro raiders and four zombies attempt to clear a collapsed passage. They hope to eventually reach a tunnel leading down to the Underdark (see area D10) blocked by debris.

The derro attack intruders on sight, eager to rob them. The zombies obey the derro's commands.

\subparagraph{Shattered Idol}
The fragments of black marble are the remnants of an idol to Diirinka, the derro's cruel patron god. The star's impact caused the idol to topple and shatter. The fragments are worthless.

\subsubsection{D5: Collapsed Tunnel}
\begin{DndReadAloud}{}
A short distance down the corridor, fallen masonry completely blocks the path. Partially buried under the rocks lie several derro corpses—alongside a few long-dead hunters with decaying flesh.

\end{DndReadAloud}

These derro, along with their zombie thralls, were caught in a cave-in while trying to dig a tunnel back to the Underdark. The collapse killed the derro and laid the zombified hunters to rest.

\subsubsection{D6: Stolen Goods}
\begin{DndReadAloud}{}
A few crates are scattered around this room, and two immense golden antlers hang from the far wall. Five derro cluster around an open chest in the center of the room, examining its contents. Two slack-jawed zombies stand nearby, watching for intruders.

\end{DndReadAloud}

Three derro raiders and two derro apprentices gather around a chest in the middle of this room. They're guarded by two zombies. The derro and their Undead thralls attack intruders on sight.

\subparagraph{Fallen Star}
Inside the chest rests the fallen star, laid reverently on a rich, red cloak. The 10-pound star is roughly spherical and about the size of a grapefruit. It's made of a black, glass-like material with a burnt appearance. The star is nonmagical but extremely resilient; it can't be destroyed by anything short of a \textit{Disintegrate} spell.

\subparagraph{Treasure}
One derro apprentice wears a pair of gold armbands worth 50 gp each, and the other apprentice has a \textit{Potion of Diminution} hanging from their belt. The golden antlers are worth 300 gp each; they were stolen from the hunters of Therno Lake. The other crates in this room contain additional goods stolen from Therno Lake: casks of ale, dishware, fur coats, spears, and carved wooden trinkets.

\includegraphics[width=0.4\textwidth]{images/../5e.tools/img/adventure/QftIS/042-03-006.fallen-star.png}
\subsubsection{D7: Overseer's Room}
\begin{DndReadAloud}{}
The walls and ceiling of this devastated room have mostly collapsed, and falling rock has smashed all the furniture except a sturdy chest in the south corner. Vials of blood, pouches of bones, and other macabre paraphernalia lie scattered around the room. An eerie, persistent scratching sound comes from behind the fallen masonry blocking the far end of the room.

\end{DndReadAloud}

This was the private room of the faction's original overseer. The overseer was killed elsewhere in the lair by a cave-in (see area D8). In life, the overseer used this room to transform corpses—primarily those of hunters from Therno Lake—into zombies. A character who examines the room's contents and succeeds on a DC 12 Intelligence (Arcana) check recognizes the vials of blood and pouches of bone as components used in foul necromantic rituals.

Five zombies are trapped in a small pocket of loose rubble (see the "Derro Lair Features" section). If the characters clear the rubble, the freed zombies attack their liberators.

\subparagraph{Treasure}
A locked iron chest in this room contains 900 gp. The key is on the overseer's corpse in area D8. Alternatively, as an action a character can pick the lock with a successful DC 18 Dexterity (Sleight of Hand) check using thieves' tools.

\subsubsection{D8: Ruins of the Great Hall}
This area is blocked off by loose rubble (see the "Derro Lair Features" section). If the characters clear the rubble, read or paraphrase the following:

\begin{DndReadAloud}{}
The air in these ruins is thick with the smell of death. Shattered furniture, rubble, and the corpses of derro and hunters litter the floor. A pair of stylish leather boots protrude from a mound of fallen rock.

\end{DndReadAloud}

The corpses of fifteen derro and twelve Humanoid hunters lie on the floor, all casualties of the cave-in. More bodies are buried beneath the rocks.

The boots—and the pair of feet in them—belong to the corpse of the outpost's overseer. Pulling the corpse from the rubble dislodges some rocks in the process. Each creature within 10 feet of the corpse must make a DC 10 Dexterity saving throw, taking 11 (2d10) bludgeoning damage on a failed save or half as much damage on a successful one.

\subparagraph{Treasure}
The overseer's corpse has a \textit{Wand of Magic Detection} in one hand and, in a pocket, the key to the locked chest in area D7. The corpse also wears a gold ring worth 50 gp. The boots, though stylish, are worth no more than ordinary boots.

\subsubsection{D9: Hungry Lamia}
This area is blocked off by loose rubble (see the "Derro Lair Features" section). If the characters clear the rubble, read or paraphrase the following:

\begin{DndReadAloud}{}
The walls of this room are painted pale green. A small table by the north wall bears a crystal decanter of wine. On a pile of quilted cushions in the corner sits a creature whose head, arms, and torso resemble those of a human woman but whose lower half is that of a four-legged lion. "Visitors!" she says over the rumble of her stomach. "Please, have you any fresh meat?"

\end{DndReadAloud}

A \textbf{lamia} sits in this room. Formerly an ally of the derro, she used her magic to help the raiders learn Therno Lake's defenses. Secretly, many derro feared the lamia, and they were relieved when the cave-in trapped her here.

The lamia hasn't eaten in days. If the characters give her a generous portion of fresh meat—such as a recently slain derro—she devours it, then lies down for a nap, paying them no further attention. Rations, zombies, and similarly unappetizing fare anger the lamia. If the characters proffer unsuitable meat or none at all, she attacks them ravenously.

\subparagraph{Poison Decanter}
The crystal decanter contains one glass's worth of wine, into which the lamia secretly mixed a dose of potent poison. A creature that drinks the wine must make a DC 10 Constitution saving throw. On a failed save, the creature takes 6 (1d12) poison damage and has the poisoned condition for 24 hours. On a successful save, the creature takes half as much damage only.

A small silver key rests at the bottom of the decanter, obscured by the wine until the decanter is emptied. This key opens a jewelry box beneath the lamia's cushions (see the following section).

\subparagraph{Treasure}
A locked silver jewelry box is hidden beneath the cushions. As an action, a character can pick the lock with a successful DC 18 Dexterity (Sleight of Hand) check using thieves' tools. The jewelry box itself is worth 20 gp, and it contains four red garnets worth 50 gp each.

\subsubsection{D10. Eastern Tunnel}
This collapsed tunnel, blocked by rubble, eventually leads to a derro stronghold in the Underdark. The stronghold and its contents are beyond the scope of this story but could inspire an additional adventure if, by some means, the party reaches this tunnel.

\section{Therno Lake}
Ever since the derro established an outpost in the southern part of Therno Pass, they have preyed cruelly on the region's local creatures and Humanoid inhabitants. In a bid to protect themselves, a band of hunters has allied with a family of intelligent giant beavers. These beavers have constructed a dam across the northern stretch of the pass, forming a lake behind the dam.

The hunters and giant beavers are wary of trespassers, and thanks to their alliance, the lakeside community has recently suffered fewer losses from the ongoing derro attacks. Curiously, the derro drag away the corpses of hunters they kill in these raids. The hunters of Therno Lake don't know it, but the derro have been animating these corpses as zombies to help guard their outpost.

\subsection{Giant Beavers}
Therno Lake is home to a family of intelligent giant beavers. These beavers use the \textbf{giant weasel} stat block with the following changes.


\begin{itemize}
\item They have a swimming speed of 40 feet.
\item They understand Common but can't speak it.
\end{itemize}

\subsection{Therno Lake Locations}
The following locations are keyed to map 3.4.

\includegraphics[width=0.4\textwidth]{images/../5e.tools/img/adventure/QftIS/043-map-3.04-therno-lake.png}
\includegraphics[width=0.4\textwidth]{images/../5e.tools/img/adventure/QftIS/044-map-3.04-therno-lake-player.png}
\subsubsection{L1: Beaver Dam}
\begin{DndReadAloud}{}
A wide dam of mud, logs, and branches blocks this section of the pass, creating a lake to the north. A trickle from the dam feeds a stream that winds south.

\end{DndReadAloud}

Three giant beavers (see the "Giant Beavers" section) swim underwater on the north side of the dam. If the characters approach from the south, the giant beavers poke their heads above the barrier for a moment, then begin slapping their tails on the water. This calls the hunters in area L2, who arrive at the dam after 1 minute.

\subparagraph{Beaver Dam}
The dam is sturdy and made of wet wood. It has Armor Class 14, 100 hit points, immunity to poison and psychic damage, and resistance to fire damage. If reduced to 0 hit points, the dam is breached.

If the dam is breached, each creature in the lake is swept southward 300 feet by the ensuing deluge, though the beaver lodge (area L3) isn't washed away. Additionally, water blasts south from the dam in a 100-foot-wide, 300-foot-long line. Each creature in that area must make a DC 15 Strength saving throw. Creatures with a swimming speed have advantage on this saving throw. On a failed save, a creature takes 18 (4d8) bludgeoning damage from the crashing waves, has the prone condition, and is swept southward 300 feet from the dam. On a successful save, the creature takes half as much damage only.

\subsubsection{L2: Lake}
\begin{DndReadAloud}{}
Where a mountain spur divides two branches of the valley, an elongated, V-shaped lake extends across the pass, contained by a dam to the south. Three small hunting boats glide around the lake in wide circles. The humanoids aboard peer intently into the water with spears at the ready.

\end{DndReadAloud}

Six hunters (scouts) are fishing in the lake, two per hunting boat. One is their community's leader, a red-haired human woman named Aspen.

\subparagraph{Speaking with the Hunters}
Aspen and the other hunters are neutral good. While the hunters are wary of outsiders, if the characters show they mean no harm, the hunters agree to parley. Aspen shares the following information:


\begin{itemize}
\item The hunters witnessed the star fall, but they aren't sure exactly where it landed.
\item They know the star fell south of Therno Lake, in the direction where the derro raiders typically come from.
\item They know the location of the derro lair, and they can point the characters to it.
\end{itemize}

\subparagraph{Quest to the Derro Lair}
If Aspen trusts the party, she tasks them with a special quest. The derro have been dragging hunter corpses to their lair for an unknown purpose. Recently, the derro also pillaged one of the hunters' most precious trophies: a pair of large golden antlers. (The antlers can be found in area D6 of the derro lair, which is detailed earlier in this adventure.)

If the characters learn what the derro do with the corpses and return the golden antlers to Therno Lake, Aspen gives the characters her pouch of pearls (see the "Treasure" section below).

\subparagraph{Treasure}
A pouch around Aspen's neck contains a string of pearls worth 120 gp total.

\subsubsection{L3: Beavers' Lodge}
\begin{DndReadAloud}{}
A large mound of densely woven branches and twigs, about thirty feet across, rises from the lake.

\end{DndReadAloud}

An underwater entrance to this beaver lodge leads to an air pocket above water level. Inside, two friendly giant beavers (see the "Giant Beavers" section) tend to five baby beavers. A pile of fish bones sits on one side of the room.

\subparagraph{Treasure}
If the characters offer the beavers food, one of the beavers digs through the fish bones and retrieves a random trinket (roll on the \textit{Trinkets table} in the \textit{Player's Handbook}) as a token of appreciation for the party.

\subsubsection{L4: Hunters' Camp}
\begin{DndReadAloud}{}
Five huts made from rock, tanned hides, and sticks cluster together on a spur of land jutting into the lake. Between the huts are rows of animal skins stretched over wooden frames.

\end{DndReadAloud}

The hunters of Therno Lake live in this camp. Eight hunters (scouts) and twenty commoners, all in thick leather clothes, make up the community. They are wary of outsiders but otherwise indifferent toward the characters.

\subparagraph{Treasure}
The only items of significant value in the camp are animal skins. The many wolf, bear, and fox pelts have a total value of 90 gp and a combined weight of 720 pounds.

\section{Tower of the Heavens}
The Tower of the Heavens is renowned as a seat of learning for astronomers—and as a source of prophecy for visitors who can afford it. For more on the tower, see this adventure's background.

For centuries, each elder sage has passed down the Books of Prophecy to that sage's chosen successor. Piyarz was Shalfey's pupil, and according to tradition, he would assume the title of elder sage on Shalfey's death. However, Piyarz grew impatient. When Shalfey's tower hands left in search of the star, Piyarz tried to seize the title—but Shalfey sealed himself in his sanctum (areas T19–T22). Shalfey threatened to destroy the Books of Prophecy if anyone attempts to break into his sanctum, so Piyarz decided to starve him out instead.

\subsection{Occupants of the Tower}
\includegraphics[width=0.4\textwidth]{images/../5e.tools/img/adventure/QftIS/045-03-008.the-tower-of-the-heavens.png}
The current occupants of the tower are divided into the following four factions.

\subsubsection{Griswill Garrison}
The Griswill Garrison is a group of gnome mercenaries paid to serve Shalfey. Believing the elder sage dead, the gnomes now follow Piyarz's orders.

\subsubsection{Hostel Staff}
Several commoners run a hostel on a smaller island bordering the tower. Piyarz told them Shalfey was dead and instructed them to keep the hostel up and running.

\subsubsection{Sage Pupils and Tower Hands}
Cipolla, Porro, and Lurg were Shalfey's pupils until they joined Piyarz's attempt to seize power. Evil and opportunistic, the four sages staged a coup with the help of tower hands they recruited to their cause.

Each tower sage has magical tattoos on their arms: glowing constellations that magically update to reflect that sage's hierarchy within the tower's prophetic order. Though Piyarz claims to be the new elder sage, Shalfey still lives, so the tattoos on Piyarz's arms haven't changed. Piyarz covers his arms when he's around the Griswill Garrison and the hostel staff, worried they'll discover his secret.

\subsubsection{Shalfey and Hadley}
Shalfey held the titles of elder sage and keeper of the books until Piyarz's coup. Though Piyarz now falsely claims these titles and Shalfey is confined to his sanctum, Shalfey remains the true head of the tower. His only surviving confidant is Hadley, who runs the tower's ferry. Hadley knows Shalfey still lives; he feigns loyalty to Piyarz while subtly searching for allies among the ferry's passengers. Hadley is unsure if he can trust the Griswills and hasn't told them Shalfey is alive.

\subsection{Reaching the Tower}
The Tower of the Heavens sits on a pair of islets in the middle of a wide, fast-flowing river. The waters around the islets are swirling and treacherous. Only Hadley, whose hut sits on the south riverbank, knows the waters well enough to guide boats safely between the mainland and the smaller island's jetty. Characters who attempt to sail or swim to the islands are swept downstream by the current and wash up on the riverbank far from their destination.

When the characters arrive, read or paraphrase the following text:

\begin{DndReadAloud}{}
On the south bank of the river stands a simple hut with a thatched roof. A small boat is tied up at a nearby jetty. Beyond the hut, rushing waters surround two islands linked to each other by a bridge. The smaller island has a dock and a modest outpost, and the larger boasts a grand stone tower.

As you approach the hut, a human man in a green tunic emerges from the hut and waves at you. "Fair is the day," he says cautiously.

\end{DndReadAloud}

\subsubsection{Hadley's Confidences}
Hadley (neutral good, human \textbf{commoner}) greets the characters with the first part of a secret phrase used by Shalfey and his allies. Hadley hopes the characters will reply, "but fairer the jewels of the evening." This phrase was part of the visions imparted by the memory web (see the "Flood of Memories" section earlier in this adventure), though the characters might not recognize its significance. A character who succeeds on a DC 14 Wisdom (Insight) check gets the sense Hadley is looking for a specific reply.

If the party doesn't respond with the other half of the phrase, Hadley offers to take them to the island. During the crossing, he uses the opening phrase twice more. If the characters still fail to offer the proper response, Hadley takes them to the outpost (area T1) without further conversation.

\subparagraph{Hadley Confides in the Party}
If the characters give the appropriate response to the secret phrase, Hadley trusts them and assumes they're allies of Shalfey. Read or paraphrase the following text as Hadley ferries the characters to the island:

\begin{DndReadAloud}{}
A relieved grin breaks across Hadley's face. "I had a feeling you were more than you appeared." He lowers his voice and speaks hastily. "They say Shalfey is dead and Piyarz is now the elder sage. But I know that isn't true, and now I know I can confide in you. How did you learn Shalfey's passphrase?"

\end{DndReadAloud}

Hadley believes whatever tale the characters tell him. If they admit they've never met Shalfey, Hadley is disappointed but decides the characters are his best shot at rescuing the elder sage.

\subparagraph{What Hadley Knows}
If the characters respond to Hadley's secret phrase, he shares the following information with them:


\begin{description}
\item[What's Going On?]
Shalfey is the elder sage of the Tower of the Heavens. Piyarz, his would-be successor, declared Shalfey dead a few days ago. But yesterday, Hadley heard the elder sage's desperate voice in his mind, speaking only the passphrase Hadley and the characters exchanged. Hadley believes the elder sage is still alive somewhere in the tower.
\item[What Happened to Shalfey?]
Piyarz has long sought Shalfey's title, and Hadley fears foul play. The tower hands closest to Shalfey left the tower on a secret mission the day before his supposed death, leaving the elder sage defenseless.
\item[What Can Be Done?]
Hadley asks the characters to discover what's going on in the tower and rescue Shalfey. He suggests they pose as visitors requesting prophecies and infiltrate the tower through the Questioner's Hall.
\end{description}

Moments after Hadley recounts this information, the ferry arrives at the outpost (area T1). Hadley doesn't accompany the characters further.

\subsection{Tower of the Heavens Features}
Unless otherwise specified, the Tower of the Heavens has the following features:


\begin{description}
\item[Ceilings.]
Tower rooms have 15-foot-high ceilings.
\item[Doors.]
Doors are made of thick oak planks and bound with iron. As an action, a character can force open a locked door with a successful DC 15 Strength (Athletics) check or pick the lock with a successful DC 15 Dexterity (Sleight of Hand) check using thieves' tools.
\item[Walls.]
The exterior walls of the tower and its hostel are slick stone. Climbing these walls is difficult without climbing gear or magic, requiring a successful DC 25 Strength (Athletics) check. The hostel has windows, but the tower's only windows are the glass panes over its domed observatory.
\item[Lighting.]
The tower uses wall-mounted torches lit with \textit{Continual Flame} spells, while the hostel uses ordinary lamps. The hostel's rooms are brightly lit during the day and dark at night, while the tower's rooms are always brightly lit.
\end{description}

\subsection{Tower of the Heavens Locations}
The following locations are keyed to map 3.5.

\includegraphics[width=0.4\textwidth]{images/../5e.tools/img/adventure/QftIS/046-map-3.05-tower-of-the-heavens.png}
\includegraphics[width=0.4\textwidth]{images/../5e.tools/img/adventure/QftIS/047-map-3.05-tower-of-the-heavens-player.png}
\subsubsection{T1: Outpost}
This outpost watches the bridge between the hostel and the tower. A jetty lies just north of the outpost.

When the characters arrive at the outpost, read or paraphrase the following text:

\begin{DndReadAloud}{}
A drab building of gray stone overlooks the steps leading up from the jetty. Three gnomes wearing heavy armor peek out from one of the building's tall, narrow windows. They wave as you approach.

\end{DndReadAloud}

Three members of the Griswill Garrison (lawful neutral, gnome guards) stand watch in this outpost at all hours. Unless treated poorly, the cheery gnomes are friendly toward the characters, asking what brings them to the tower. It's the gnomes' job to greet visitors, inquire about their business, and keep an eye out for any clearly hostile intruders.

\subsubsection{T2: Hostel}
Visitors stay in this small hostel while they wait to be called for their answering ceremony, the ritual through which fortunes are told. Lodging costs 1 gp per person per night; meals and drinks aren't included. Trouble in the hostel attracts the attention of the guards in the outpost (area T1).

The hostel's proprietor is a human woman named Berta. Four peppy humans in leather aprons help her carry supplies, cook, and clean. Berta and her assistants use the \textbf{commoner} stat block. They're friendly and welcoming, and they know nothing of the attempted coup. They believe Shalfey died of natural causes and regard Piyarz as the legitimate keeper of the books unless proven otherwise.

\subparagraph{Hostel Areas}
The hostel has the following areas:


\begin{description}
\item[T2a:]
Courtyard. Stairs from the jetty lead up to this open-air courtyard.
\item[T2b:]
Dining Room. The tower's visitors eat here. There are two large tables, four long benches, six chairs, and a large dresser containing tableware.
\item[T2c:]
Kitchen. The hostel's humble meals are cooked in a cast-iron pot over the fireplace. Berta's assistants sleep here near the hearth. A clearly visible trapdoor leads down to a cellar full of foodstuffs, wine, beer, and fuel for the fire.
\item[T2d:]
Visitors' Rooms. Each of these guest rooms contains a low table, three stools, and three beds.
\end{description}

\subparagraph{Hostel Visitors}
The Hostel Visitors table details guests the characters might encounter in the hostel. All adult visitors use the \textbf{commoner} stat block. If the characters failed to earn Hadley's trust earlier, you can use these visitors to spur the plot forward.

\begin{DndTable}[header=Hostel Visitors]{cX}

d4 & Visitor\\
1 & Jurgen and Gorvis, two dwarf men, have been happily married for a century. They've come to the tower on their hundredth anniversary to see what their next hundred years of union will bring.\\
2 & Durn, a human woman, is visiting the tower with her two children, June and Koll. Durn is a farmer, but all her crops failed this year. She has come to ask what next year's harvest will bring. She doesn't have the 100 gp required for an answer but hopes the sages will take pity on her.\\
3 & Namaia, a human woman, is mayor of a small village about a week's walk from the tower. An upcoming election has her on edge, and she wants to know what her odds are of keeping her title.\\
4 & Adavalis, an elf man, comes from a wealthy line of prominent sorcerers. He's never manifested any sorcerous power and wants to know if he'll ever gain such power, or if he should give up hope.\\
\end{DndTable}

\subsubsection{T3: Entrance Chamber}
\begin{DndReadAloud}{}
This hall serves as an entrance to the tower proper. Centuries of use have worn down its flagstone floor and the handles of its brass-paneled doors. Six gnome guards stand at attention along its walls.

\end{DndReadAloud}

Six members of the Griswill Garrison (lawful neutral, gnome guards) stand watch in this room. When it's time for a visitor's appointment in the Questioner's Hall (area T4), the guards walk across the bridge and call for the guest in the hostel. The guards then escort the visitor across the bridge, through the entrance chamber, and into the Questioner's Hall. The guards remain with the visitor until the answering ceremony is complete, then escort them back to the hostel.

\subsubsection{T4: Questioner's Hall}
\begin{DndReadAloud}{}
Four smoldering braziers dimly illuminate this somber hall. Dusty red curtains line the walls. The ceiling is darkened with soot, and the marble floor has been scuffed by the soles of countless feet.

\end{DndReadAloud}

Sages conduct fortune-telling sessions in this dimly lit hall. If the characters were escorted here by the gnome guards from area T3, those guards line up along the east wall for the duration of the ceremony, which begins shortly after the characters enter.

\subparagraph{Answering Ceremony}
Read or paraphrase the following text to begin the answering ceremony:

\begin{DndReadAloud}{}
The curtains at the far end of the room part, and a red-robed human man enters the hall. He has broad, hunched shoulders, and his forearms are covered with starry red tattoos. Two humans in simple red robes accompany him.

"Welcome to the answering ceremony," the tattooed man intones, "where astrology and prophecy intertwine to answer all. What is your question?"

\end{DndReadAloud}

Lurg (lawful evil, human \textbf{tower sage}) conducts this ceremony. Two tower hands loyal to Lurg accompany him (see appendix B for both stat blocks).

If allowed to run its course, the ceremony is relatively simple and businesslike. Once a question has been asked, Lurg asks for a mandatory "donation" of 100 gp, no matter what the question is. One of his tower hands then collects the payment.

Once his fee is paid, Lurg announces he must consult the stars for the answers the visitors seek. Of course, Lurg can't access the Books of Prophecy, since Shalfey still has them locked away in his sanctum, so Lurg simply withdraws behind the west curtain and remains silent for 1 minute while he concocts a cryptic, meaningless answer. After he reemerges from the curtain and delivers his response, the guards escort the questioners back to the hostel, ending the ceremony.

\subparagraph{Treasure}
Lurg wears a gold chain belt worth 30 gp and carries the key to his desk (see area T14).

\subsubsection{T5: Mimic Door}
\begin{DndReadAloud}{}
A stone door set in the south wall bears a carved pattern of stars arranged in a circle.

\end{DndReadAloud}

The door is a \textbf{mimic} trained to guard the staircase. It functions like an ordinary door until a creature it doesn't recognize tries to open it, at which point the mimic reveals its true nature and attacks. The mimic is friendly toward the tower sages, tower hands, and gnomes of the Griswill Garrison.

\subsubsection{T6: Mercenaries' Hall}
\begin{DndReadAloud}{}
The smell of pipe smoke greets visitors to this raucous hall. Four gnomes sit here laughing, eating, and playing cards.

\end{DndReadAloud}

Three members of the Griswill Garrison (lawful neutral, gnome guards) sit here, relaxing between shifts. Accompanying them is the garrison's leader, Captain Mainwaring (lawful neutral, gnome \textbf{knight}).

When she's aware of the characters, Captain Mainwaring calls the guards to arms, but characters can convince them to stand down with a successful DC 15 Charisma (Persuasion) check. Characters who mention Shalfey have advantage on this check.

If hostilities are defused, Captain Mainwaring listens to the characters' story. She doesn't fully believe Shalfey is alive unless shown undeniable evidence (such as Piyarz's unchanged tattoos or Shalfey himself). However, the characters can easily convince Mainwaring to give them the benefit of the doubt, and she agrees to overlook the characters' activities as they seek answers in the tower.

If combat breaks out in this room, the noise awakens the three guards in the barrack (area T7), who join the fight after 2 rounds.

\subparagraph{Treasure}
Captain Mainwaring wears a gold brooch worth 60 gp, and she carries a key to her chest in the barrack (area T7).

\subsubsection{T7: Barrack}
\begin{DndReadAloud}{}
A weapon rack sits in the center of this large room. The walls are lined by fifteen small beds, three of which are occupied by snoozing gnomes. A slightly larger and more comfortable bed stands against the southeast wall, with a large lockbox at its foot.

\end{DndReadAloud}

At any given time, three members of the Griswill Garrison (lawful neutral, gnome guards) are sleeping in these beds with weapons close at hand. The larger bed belongs to Captain Mainwaring, leader of the Griswill Garrison (see area T6).

\subparagraph{Treasure}
The chest at the foot of Mainwaring's bed is locked, and the captain carries the key. As an action, a character can pick the lock with a successful DC 15 Dexterity (Sleight of Hand) check using thieves' tools. The chest contains a fine fur-trimmed robe worth 10 gp, a gold chain belt worth 60 gp, and a bag containing 180 gp (the collective wealth the gnomes have earned for their service).

\subsubsection{T8: Storeroom}
Dozens of boxes, casks, and sacks are neatly stacked on shelves and racks this storeroom, which holds food stores and equipment for the entire tower. Amid the provisions are two Explorer's Pack, two Scholar's Pack, and enough food and drink to sustain the tower's residents for a month.

\subsubsection{T9: Kitchen}
A pot of squash soup simmers gently on a stove in this tidy kitchen, filling the room with its aroma. The kitchen otherwise holds nothing of interest.

\subsubsection{T10: Trapped Hallway}
This unfurnished hallway provides ladder access to the watchtowers on the third floor (area T11).

A protective glyph wards the staircase door down to area T17. To open this door safely from either direction, one must first say the password ("apogee"), which is known only to the tower's stewards. If the door is opened without the password, flame erupts in a 10-foot-radius sphere centered on the glyph. Each creature in that area must make a DC 13 Dexterity saving throw, taking 13 (3d8) fire damage on a failed save or half as much damage on a successful one.

\subsubsection{T11: Watchtowers}
Six turrets rise from the Tower of the Heavens. Six members of the Griswill Garrison (lawful neutral, gnome guards) are stationed at each watchtower, for a total of thirty-six guards. The guards constantly watch for intruders, particularly those scaling the walls. When a guard spots an intruder, three guards remain at their post to fight while three run to area T6 to warn Captain Mainwaring.

\subsubsection{T12: Lesser Library}
\begin{DndReadAloud}{}
This library is bathed in supernatural silence. Two robed figures—a halfling man clad in blue and a dragonborn woman in yellow—work at tall desks amid ancient tomes, while other robed figures stand guard.

\end{DndReadAloud}

This library contains the less important books of the tower. Sage pupils Porro (lawful evil, halfling \textbf{tower sage}) and Cipolla (lawful evil, silver dragonborn \textbf{tower sage}) research in this room, accompanied by four tower hands sworn to protect them (see appendix B for both stat blocks). All are loyal to the illegitimate elder sage, Piyarz, and enforce his will.

When the characters enter, Cipolla and Porro silently direct the hands to attack. Meanwhile, the sages flee toward Piyarz in the scholars' hall (area T17), where they can more freely use their magic. Once there, Cipolla or Porro casts the \textit{Arcane Lock} spell on the door to buy themselves time.

\subparagraph{Silence}
A permanent \textit{Silence} spell pervades the library, eliminating all sound in this area. A \textit{Dispel Magic} spell ends the effect. The magical silence prevents casting spells with verbal components.

\subparagraph{Treasure}
Porro and Cipolla wear gold chain belts worth 30 gp each. They each carry a key to their desk drawers (see areas T15 and T16).

\subsubsection{T13: Dormitory}
Tower hands sleep in this stark dormitory, which contains only thirteen modest wooden beds with thin straw mattresses. It is currently unoccupied.

\subsubsection{T14: Lurg's Room}
\begin{DndReadAloud}{}
The walls and furniture of this room are painted deep crimson. Messy stacks of paper sit atop a writing desk, and the bed's maroon sheets are lazily bunched.

\end{DndReadAloud}

This is Lurg's private chamber. Its crimson walls match the color of his robes.

Lurg's desk drawer is locked and magically trapped. Lurg carries its key (see area T4), and only he knows the password to safely unlock the drawer: "Struve." As an action, a character can unlock the door with a successful DC 15 Dexterity (Sleight of Hand) check using thieves' tools. If the drawer is unlocked without first speaking the password, a magical glyph inside the drawer triggers, erupting with flame in a 10-foot-radius sphere. Each creature in that area must make a DC 13 Dexterity saving throw, taking 13 (3d8) fire damage on a failed save or half as much damage on a successful one. The glyph vanishes after it triggers, and the flame doesn't damage the contents of the room.

\subparagraph{Treasure}
The desk drawer contains quills, ink, paper, a pouch holding 60 gp, and Lurg's \textit{spellbook}. The spellbook contains the spells in the \textbf{tower sage} stat block plus the following: \textit{Color Spray}, \textit{Illusory Script}, and \textit{Unseen Servant}.

\subsubsection{T15: Porro's Room}
\begin{DndReadAloud}{}
Everything in this room, including the walls, is painted cerulean. Its contents are sized for a small inhabitant; the desk is two feet tall, and the bed is four feet long.

\end{DndReadAloud}

This is Porro's private chamber. Its cerulean walls match the color of his robes.

Porro's desk drawer is locked and magically trapped. Porro carries its key (see area T12), and only he knows the password to safely unlock the drawer: "Altair." As an action, a character can unlock the door with a successful DC 15 Dexterity (Sleight of Hand) check using thieves' tools. If the drawer is unlocked without first speaking the password, a magical glyph inside the drawer triggers, erupting with icy wind in a 10-foot-radius sphere. Each creature in that area must make a DC 13 Constitution saving throw, taking 13 (3d8) cold damage on a failed save or half as much damage on a successful one. The glyph vanishes after it triggers, and the wind doesn't damage the contents of the room.

\subparagraph{Treasure}
The desk drawer contains quills, ink, paper, a small necklace of emerald beads (worth 160 gp), a pouch holding 30 gp, and Porro's \textit{spellbook}. The spellbook contains the spells in the \textbf{tower sage} stat block plus the following: \textit{Color Spray}, \textit{Disguise Self}, and \textit{Nystul's Magic Aura}.

\subsubsection{T16: Cipolla's Room}
\begin{DndReadAloud}{}
A large, framed painting of the Tegefed Mountains adorns this room. The painting's golden hues complement the room's yellow walls and furniture.

\end{DndReadAloud}

This is Cipolla's private chamber. Its yellow walls match the color of her robes.

Cipolla's desk drawer is locked and magically trapped. Cipolla carries its key (see area T12), and only she knows the password to safely unlock the drawer: "Sol." As an action, a character can unlock the door with a successful DC 15 Dexterity (Sleight of Hand) check using thieves' tools. If the drawer is unlocked without first speaking the password, a magical glyph inside the drawer triggers, erupting with lightning in a 10-foot-radius sphere. Each creature in that area must make a DC 13 Dexterity saving throw, taking 13 (3d8) lightning damage on a failed save or half as much damage on a successful one. The glyph vanishes after it triggers, and the lightning doesn't damage the contents of the room.

\subparagraph{Treasure}
The desk drawer contains quills, ink, paper, a small horse figurine carved from coral (worth 100 gp), a pouch holding 50 gp, and Cipolla's \textit{spellbook}. The spellbook contains the spells in the \textbf{tower sage} stat block plus the following: \textit{Color Spray}, \textit{Fog Cloud}, and \textit{Magic Mouth}.

\subsubsection{T17: Scholars' Hall}
\begin{DndReadAloud}{}
Astrological symbols decorate this vast, lavish hall. A massive table of polished oak dominates the room.

Next to the table stands a short human mage with a goatee and gray robes. His arms are covered with intricate gray tattoos, and a small raven perches on his shoulder. Standing on either side of the table are two humans dressed in similar gray robes and hand wraps.

\end{DndReadAloud}

Piyarz is here with two tower hands. His loyal \textbf{imp} rests smugly on his shoulder in the form of a raven. Piyarz is a lawful evil \textbf{tower sage} who wears a \textit{Ring of Fire Resistance}, granting him resistance to fire damage.

If Cipolla or Porro successfully fled area T12, they're here as well, and one of them has cast the \textit{Arcane Lock} spell on the south door to buy time.

\subparagraph{Confronting Piyarz}
If the characters don't attack instantly, Piyarz announces he'll "generously" give them one chance to explain themselves. A shrewd, calculating liar, Piyarz tries to make the characters doubt anything Hadley might have told them. Piyarz hopes to convince them his assumed title of elder sage is rightfully deserved. If he can't persuade the characters to leave, Piyarz and his allies attack, fighting to the death.

\subparagraph{Treasure}
Piyarz wears a \textit{Ring of Fire Resistance} and a platinum chain belt worth 300 gp.

\includegraphics[width=0.4\textwidth]{images/../5e.tools/img/adventure/QftIS/048-03-009.piyarz-cipolla-porro.png}
\subsubsection{T18: Piyarz's Room}
\begin{DndReadAloud}{}
Hovering above a large coffer in the center of this gray bedroom is a many-eyed, spherical creature three feet in diameter. In the middle of its forehead, a large unblinking eye regards you quizzically—as do four smaller eyes perched on the end of wiggling eyestalks.

\end{DndReadAloud}

The creature is a \textbf{spectator} in Piyarz's service. Commanded to guard the chamber, the spectator telepathically warns unauthorized entrants to "scram or get zapped" and attacks any who refuse.

In addition to the coffer, Piyarz's private chamber contains a bed, a chair, and a desk, all painted dull gray to match his robes.

\subparagraph{Treasure}
The desk has a drawer containing parchment, quills, and ink. The coffer is unlocked and contains two Potion of Greater Healing, a \textit{Folding Boat}, 3,600 gp, and Piyarz's \textit{spellbook}. The spellbook contains the spells in the \textbf{tower sage} stat block plus the following: \textit{Color Spray}, \textit{Find Familiar}, \textit{Locate Object}, \textit{Magic Missile}, and \textit{Sleep}.

\subsubsection{T19: Bridge of Faith}
\begin{DndReadAloud}{}
The thunder of rushing water fills this loud room. The room is divided in two by a yawning chasm, ten feet deep and surging with a foaming torrent. Inlaid on the walls of the room are the following words: "I am real. Watch your step, have faith, and pass safely."

\end{DndReadAloud}

The walls of this room are smooth, wet, and slippery. A character can climb along the wall to the other side with climbing gear and a successful DC 15 Strength (Athletics) check. On a failed check, the character falls into the rushing water below.

\subparagraph{Rushing Water}
Water rushes into the chasm from the river outside via concealed channels below the tower. A creature that falls into the chasm takes 10 (3d6) bludgeoning damage as the river sweeps it away. The creature then washes up on the island's rocky outcropping just north of area T17.

\subparagraph{Invisible Bridge}
An invisible, 10-foot-wide bridge spans the chasm. It has no railing but is otherwise straight and safe to traverse. A creature made aware of the bridge's presence—such as by casting the \textit{See Invisibility} spell or interacting with it physically—can make a DC 10 Intelligence (Investigation) check. On a successful check, the bridge becomes visible to that creature. Creatures that can't see the bridge can still feel their way across, albeit carefully.

\subsubsection{T20: Greater Library}
Shalfey (lawful good, human \textbf{tower sage} with 4 levels of exhaustion; see appendix B) has locked himself inside this room and threatened to burn the Books of Prophecy if Piyarz breaks in. Both doors are under the effects of \textit{Arcane Lock} spells that allow only Shalfey to open the doors normally. Shalfey has also reinforced each door with a barricade.

A character outside the library who introduces themself can make a DC 20 Charisma (Persuasion) check to convince Shalfey to let them in. Characters who recite Shalfey's secret phrase (see the "Flood of Memories" section earlier in this adventure) have advantage on this check. On a failed check, Shalfey condemns the character as an agent of Piyarz and orders them to step away from the door—or else.

\subparagraph{Desperate Measures}
Characters who force their way into Shalfey's sanctum risk him destroying the Books of Prophecy. Each door to the greater library has Armor Class 15, 50 hit points, and immunity to poison and psychic damage.

When the characters enter, Shalfey assumes the characters are Piyarz and his conspirators. Read or paraphrase the following text:

\begin{DndReadAloud}{}
"You leave me no choice, Piyarz!" shouts a desperate voice as you enter this spacious library. A bearded man in white robes stands amid the shelves, his arms covered with intricate, softly glowing tattoos. A large, shimmering globe of light levitates behind him.

\end{DndReadAloud}

If the characters attempted to destroy the doors, Shalfey has already set the Books of Prophecy aflame. Read or paraphrase the following text:

\begin{DndReadAloud}{}
The man's defiant face is lit by roaring flames from a pile of burning tomes before him. His outstretched hands flicker with waning, fiery magic. When he sees you, he looks both shocked and relieved.

\end{DndReadAloud}

Otherwise, they catch Shalfey before he makes a grave error. Read or paraphrase the following text:

\begin{DndReadAloud}{}
With defiance in his eyes, the man extends his hands toward a pile of ancient tomes, his fingers ablaze with magic. When he sees you, his face softens in relief, and he slowly lowers his hands away from the tomes.

\end{DndReadAloud}

\subparagraph{Conversing with Shalfey}
Whether or not Shalfey burned the Books of Prophecy, he is relieved the characters aren't Piyarz and his minions. The elder sage asks why they came to the tower, then offers the following information:

\includegraphics[width=0.4\textwidth]{images/../5e.tools/img/adventure/QftIS/049-03-010.elder-sage.png}

\begin{description}
\item[Fallen Star.]
If the characters don't have the star, Shalfey encourages them to visit Derwyth (if they haven't done so) to find it for him. If the party has the notes from Derwyth's library (see area H5 earlier in this adventure), he uses those to quickly calculate the star's point of impact.
\item[Piyarz.]
Shalfey explains that Piyarz attempted a coup. Shalfey asks the adventurers to defeat or exile Piyarz and the other traitorous sages (if they haven't done so already); the characters can keep the conspirators' possessions. Shalfey accompanies the characters if requested, but he's dangerously exhausted and avoids combat.
\item[Shalfey's Motivations.]
If pressed about why he covets the star, Shalfey tells most of the truth but doesn't expose the sages' reliance on the Books of Prophecy. He explains that "through arduous study of ancient and esoteric tomes," he learned of svirfneblin smiths to the west who would exchange a "wondrous library" for a black rock from the heavens. He claims to have predicted the star's arrival through astronomical observation.
\end{description}

\subparagraph{Shalfey's Quest}
Once the traitorous sages are dealt with and the star is obtained, Shalfey asks the characters to take the star to the svirfneblin forge and exchange it for a set of books—the smiths will know which ones. Shalfey gives the party directions and promises to reward each character with a magic item from his treasury upon receipt of the books.

\subparagraph{Shimmering Sphere}
The levitating, scintillating sphere in this room is the tower's treasury, which Shalfey explains in esoteric terms to any character who asks. He alone can draw magic items from the sphere, as described in this adventure's conclusion.

\subparagraph{Treasure}
Shalfey wears a \textit{Medallion of Thoughts} and a mithral chain belt worth 400 gp. While his old Books of Prophecy can no longer divine the future, they're a reliable source on significant historical events—though incomprehensible to anyone but a tower sage. If Shalfey set the Books of Prophecy on fire, they are unreadable.

\subsubsection{T21: Labyrinth}
\begin{DndReadAloud}{}
Beyond the door to this room lies a swirling mass of colors and shapes that stretches into the distance.

\end{DndReadAloud}

A creature that enters the swirling mass finds itself alone in a labyrinthine demiplane of nauseating colors. A creature navigating this magical labyrinth must make a DC 15 Intelligence (Investigation) or Wisdom (Survival) check. On a failed check, the creature takes 10 (3d6) psychic damage, gains 1 level of exhaustion, and finds the exit after 10 minutes of wandering. On a successful check, the creature takes half as much damage only and finds the exit after 1 minute. In either case, the creature emerges from the door opposite the one it entered.

As the true elder sage, Shalfey is immune to the labyrinth's ill effects and can navigate the maze without consequence (no check required).

\subsubsection{T22: Observatory}
\begin{DndReadAloud}{}
This rooftop observatory is surrounded by a decorative lattice of metal strips and glass panes, rising to create a dome atop the tower. A chair with a swiveling base sits in the center of the room, flanked by two stone statues of snarling winged creatures.

\end{DndReadAloud}

The sages make astronomical observations here. By observing the positions of the stars and planets relative to the lattice, the sages can precisely record the movements of the heavens.

The statues are two gargoyles. They're trained to attack any visitors unaccompanied by a tower sage or tower hand.

Characters can climb the tower's walls to this observatory, but doing so might draw the attention of the guards stationed in the watchtowers (area T11). The openings in the domed lattice are paned with glass and just wide enough for a Medium creature to squeeze through. Each pane has Armor Class 13, 3 hit points, and immunity to psychic and poison damage.

\section{Forge of the Kagu-Svirfneblin}
A small faction of svirfneblin—also known as deep gnomes—has established an underground lair beneath the Tegefed Mountains. When these svirfneblin were still young and foolhardy, they left the Underdark in search of precious stones. These smiths call themselves the kagu-svirfneblin—Gnomish for "stone-seekers."

\subsection{Prophecy of Prophecies}
During their prospecting, the kagu-svirfneblin encountered a shimmering dome—not unlike the shimmering sphere found in the Tower of the Heavens. Near the dome, a stone tablet prophesied creatures would one day arrive at this very site bearing a fallen star. The tablet foretold this star would unlock the dome, whose contents must be exchanged for the star. Unknown to the kagu-svirfneblin, the dome contains a second set of Books of Prophecy.

Eager to claim the star for their collection and keep knowledge of it to themselves, the svirfneblin shattered the tablet and decided not to return to the Underdark. They settled around the dome and cut themselves off from their deep-dwelling kin.

The kagu-svirfneblin are now quite old, even by gnomish standards. In their isolation, they have devoted much of their time to making machines, and their lair is more akin to a forge than a home. Their obsession with the promised star has made them wary of outsiders.

\subsection{Grisdelfawr the Dragon}
Recently, a \textbf{young red dragon} named Grisdelfawr made a lair near the forge and began terrorizing any svirfneblin who dared venture outside. To distract the dragon, the svirfneblin corralled a horde of gibberlings, noisy fiends prevalent in the Underdark, and set them loose in the valley.

\subsection{Approaching the Forge}
As the characters near the forge, they discover evidence of the creatures that dwell nearby.

\subsubsection{Scorched Earth}
In the woods about two miles from the forge of the kagu-svirfneblin, the characters come across a gruesome scene. Read or paraphrase the following text:

\begin{DndReadAloud}{}
A fifteen-foot stretch of woodland has been blackened by a concentrated blaze, and other nearby patches are similarly scorched. In one patch lies the charred corpse of a bald, grayish-purple gnome, which has been partially devoured by a larger predator.

\end{DndReadAloud}

Grisdelfawr killed the svirfneblin here a week ago and razed the area with his fiery breath. The dragon took a single bite of the gnome, wastefully leaving the rest of the corpse. A character who inspects it and succeeds on a DC 13 Wisdom (Medicine) check knows the bite was left by a large, reptilian creature. A subsequent successful DC 15 Intelligence (Nature) check reveals the predator was a dragon.

\subsubsection{Gibberling Horde}
About one mile from the forge, the characters are beset by roaming gibberlings. Read or paraphrase:

\begin{DndReadAloud}{}
A strange noise in the distance grows louder with each passing second. It sounds like a dozen shrieking voices, babbling and spitting incomprehensibly.

\end{DndReadAloud}

Moments later, three swarms of gibberlings burst from a nearby thicket and stampede toward the characters. The gibberlings fight to the death as they try to tear the characters limb from limb.

After the characters deal with the gibberlings, the journey to the forge of the kagu-svirfneblin proceeds without complication. They arrive at area F1.

\subsection{Forge Features}
Unless otherwise stated, the kagu-svirfneblin forge has the following features:


\begin{description}
\item[Ceilings.]
Rooms have 30-foot-high ceilings, and hallways have 10-foot-high ceilings.
\item[Doors.]
All doors are made of solid iron.
\item[Lighting.]
The forge isn't illuminated. Area descriptions assume the characters have a light source or other means of seeing in the dark.
\end{description}

\subsection{Forge Locations}
The following locations are keyed to map 3.6.

\includegraphics[width=0.4\textwidth]{images/../5e.tools/img/adventure/QftIS/050-map-3.06-forge-of-the-kagu-svirfneblin.png}
\includegraphics[width=0.4\textwidth]{images/../5e.tools/img/adventure/QftIS/051-map-3.06-forge-of-the-kagu-svirfneblin-player.png}
\subsubsection{F1: Misty Gulch}
\begin{DndReadAloud}{}
Leading off from the main valley is a narrow, steep-sided canyon. Warm mist rises from a stream, shrouding the canyon. Within the mist, you hear a creature babbling incessantly.

\end{DndReadAloud}

This area is lightly obscured by mist from a piping-hot stream. The noise comes from the gibberling in area F2.

\subparagraph{Steamy Stream}
A metal bridge spans this 5-foot-deep stream, which issues from the workshop in area F8 and is scalding hot. A creature that enters the stream for the first time on a turn or starts its turn there must make a DC 13 Constitution saving throw, taking 7 (2d6) fire damage on a failed save or half as much damage on a successful one.

\subsubsection{F2: Lizard Pens}
\begin{DndReadAloud}{}
The cliff opens into a dark cave whose mouth is barred by a heavy metal grill. Behind the grill, two large lizards with saddles strapped to their backs hiss menacingly. Outside the cave paces a small creature with spindly limbs, matted fur, and wild eyes. The creature babbles, screeches, and shakes the bars, attempting to reach the lizards inside.

\end{DndReadAloud}

The svirfneblin keep two giant lizards trained as mounts in these pens. A \textbf{gibberling} scrabbles at the bars in a frenzied attempt to get at the lizards. The gibberling attacks the characters on sight.

A padlock secures the gate. As an action, a character can pick the lock with a successful DC 15 Dexterity (Sleight of Hand) check using thieves' tools or wrench apart the bars with a successful DC 20 Strength (Athletics) check. If released, the giant lizards attack any creature they don't recognize.

As an action, a character can calm a lizard with a successful DC 10 Wisdom (Animal Handling) check, changing its attitude toward the party from hostile to indifferent. If the check succeeds by 5 or more, the lizard becomes friendly and willing to serve them as a mount.

\subsubsection{F3: Outer Hall}
\begin{DndReadAloud}{}
This grand room is completely unfurnished and polished to a metallic sheen. In the center of the floor rises an eight-foot-wide clockwork dome.

\end{DndReadAloud}

The dome in the center of the floor is a \textbf{maschin-i-bozorg}. It remains motionless until the characters enter the room or disturb it. The maschin-i-bozorg is hostile toward intruders and fights until destroyed.

\subsubsection{F4: Storerooms}
The kagu-svirfneblin keep various goods in these three identical storerooms. Each room is long and narrow with iron storage carts along the walls.

The storerooms and their contents are detailed below.


\begin{description}
\item[F4a:]
Food. The carts in this room contain flour, dried vegetables, salted meats, and wine. These stores are worth 20 gp total.
\item[F4b:]
Fuel. The carts in this room contain coal, oil, wood, and charcoal. These stores are worth 15 gp total.
\item[F4c:]
Raw Materials. The carts in this room contain glass, clay, and tin worth 25 gp total.
\end{description}

\subsubsection{F5: Inner Hall}
\includegraphics[width=0.4\textwidth]{images/../5e.tools/img/adventure/QftIS/052-03-011.kagu-svirfneblin.png}
\begin{DndReadAloud}{}
Warm, swirling vapor fills this large room, rising from a stream that crosses the room from south to north. A mist-clouded bridge rises over the stream. On the other side wait three dome-shaped clockwork machines. They puff steam threateningly but make no advances.

Moments later, seven grayish-purple, bald gnomes, each astride a giant lizard, appear behind the machines. "Peace!" one calls in a nasal voice. "Do you come for trade or war?"

\end{DndReadAloud}

This area is \textit{lightly obscured} by mist and steam. The clockwork machines are three maschin-i-bozorgs, created by the kagu-svirfneblin to guard their forge. The seven svirfneblin—each riding a \textbf{giant lizard}—constitute the entire deep gnome population here. They eagerly expect the arrival of the prophesied star.

\subparagraph{Steamy Stream}
The steaming stream, which flows from the workshop in area F8, is scalding hot and 3 feet deep. A creature that enters the stream for the first time on a turn or starts its turn there must make a DC 13 Constitution saving throw, taking 7 (2d6) fire damage on a failed save or half as much damage on a successful one.

\subparagraph{Parleying with the Kagu-Svirfneblin}
If the characters respond peacefully, one of the svirfneblin approaches cautiously. She introduces herself as Ecklash the Burrow Warden, the faction's leader. Ecklash has a large burn across her cheek from a close call with the red dragon Grisdelfawr. Polite but frank, Ecklash does most of the talking for her group. She has an inkling the adventurers are the promised bearers of the fabled star.

As stone-seekers, the kagu-svirfneblin wish for nothing more than to add this "sky-stone" to their collection. The gnomes are brief and direct in their negotiations. In return for the star, kagu-svirfneblin agree to give the characters the contents of their shimmering dome (see area F9), but they require the star to open it. If the characters agree to their terms, the kagu-svirfneblin escort them to area F9 so the exchange can take place.

\subparagraph{Fighting the Kagu-Svirfneblin}
Combat with the kagu-svirfneblin and their maschin-i-bozorgs is likely to end poorly for the characters, and any characters who fought the maschin-i-bozorg in area F3 recognize as much. If combat breaks out, the deep gnomes fight defensively and attempt to defuse the situation. They try to incapacitated the characters rather than kill them, hoping to resume peaceful negotiations once the characters have calmed down.

\subsubsection{F6: Dormitory}
\begin{DndReadAloud}{}
A metal table and seven metal chairs sit in the center of this spacious dormitory. Seven alcoves each contain an iron-framed bed, while an eighth alcove is blocked by a metal door secured with a heavy lock.

\end{DndReadAloud}

The kagu-svirfneblin eat and sleep here. Obsessed with their work and the prophesied stone, the deep gnomes have little regard for their own comfort.

The door to the eighth alcove is locked. As an action, a character can pick the lock with a successful DC 15 Dexterity (Sleight of Hand) check using thieves' tools or they can break down the door with a successful DC 20 Strength (Athletics) check.

\subparagraph{Treasure}
The locked alcove contains racks of items labeled in Gnomish. Among them are a \textit{Spell Scroll} of \textit{Fire Shield}, a \textit{Potion of Heroism} labeled "magical vitality," and a jar of \textit{Keoghtom's Ointment} mislabeled "war paint."

\subsubsection{F7: Treasury}
A padlock secures the gate to this treasury. As an action, a character can pick the lock with a successful DC 15 Dexterity (Sleight of Hand) check using thieves' tools or wrench apart the bars with a successful DC 20 Strength (Athletics) check.

Two hook horrors reside within, guarding the kagu-svirfneblin's treasure. They lumber about in the darkness, tapping their long, bony claws along the floor. The hook horrors obey the deep gnomes but attack all other creatures on sight.

\subparagraph{Treasure}
In the rectangular chamber at the rear of the pen, the kagu-svirfneblin store gold and silver ingots (worth 1,200 gp total) and a bag of uncut gems (worth 500 gp total).

\subsubsection{F8: Workshop}
\begin{DndReadAloud}{}
This hot, humid room reverberates with the drumlike beat of steam-powered contraptions, whirring metal machines, and hissing pipes. The room also contains a few ordinary anvils, forges, and workbenches.

\end{DndReadAloud}

The machines are powered by steam from a large cylinder at the east end of the workshop, which is fed by a natural geyser. The kagu-svirfneblin built these machines according to their own designs, and only they know how to operate them. A character who tampers with the machines must make a DC 16 Intelligence (Investigation) check. On a successful check, the contraptions whir and move but nothing disastrous happens. On a failed check, the character takes 10 (3d6) fire damage from a blast of scalding steam.

\subsubsection{F9: Shimmering Dome}
\begin{DndReadAloud}{}
Multicolored light dances through this faceted room, shining from a scintillating dome set into the floor.

\end{DndReadAloud}

This is the dome the kagu-svirfneblin found long ago. They've never managed to pierce it. Characters who saw Shalfey's shimmering sphere (see area T20 of the Tower of the Heavens) notice both seem to be made of the same multicolored essence.

If the star is brought within 10 feet of the dome, read or paraphrase the following text:

\begin{DndReadAloud}{}
The dome begins to brighten rapidly. Suddenly, the room flashes with prismatic radiance—and the dome vanishes in a haze of sparkling dust. In its place rests a neat stack of twenty leather-bound books.

\end{DndReadAloud}

When the dome vanishes, its magic bestows one last gift. Each creature who helped bring the star to the dome (including each of the characters) tingles with resilience, gaining 5 \textit{temporary hit points} and resistance to fire damage until the next dawn.

\subparagraph{Books of Prophecy}
The stack of tomes is the second set of Books of Prophecy sought by Shalfey. The books are of little use to anyone other than the tower sages, who have dedicated countless years of study to understand them. The books weigh 60 pounds in total.

Once the characters trade the star to the kagu-svirfneblin for the new Books of Prophecy, proceed to the "After the Exchange" section.

\subsubsection{F10: Tunnel to the Underdark}
This tunnel slopes down for many miles, eventually connecting to a greater tunnel system deep in the Underdark. Most parties have no reason to journey down it during this adventure, but the expansive realms below and their subterranean threats could be crafted by the DM for future adventures.

\subsection{After the Exchange}
With the trade complete, the kagu-svirfneblin thank the characters and escort them back to the inner hall (area F5). They invite the characters to take whatever they want from the storerooms (area F4). The kagu-svirfneblin announce their purpose here is complete, and they're about to return to the Underdark, sealing off their tunnel to protect themselves from dragons and other outsiders.

The kagu-svirfneblin advise the characters to avoid the workshop (area F8), plainly informing them it's about to self-destruct. The peculiar smiths are unconcerned about the destruction of the forge; they can always build another.

\subsubsection{Evacuation of the Kagu-Svirfneblin}
After warning the characters, the kagu-svirfneblin bid them goodbye and exit the inner hall through the southeast double door, accompanied by their maschin-i-bozorgs. They shut and lock the doors behind them. As the kagu-svirfneblin approach the tunnel to the Underdark (area F10), they collect the treasure and hook horrors, dead or alive, in areas F6 and F7. If any treasure is missing, the deep gnomes shrug and continue on their way.

After setting the machines in the workshop (area F8) to overload, the kagu-svirfneblin depart for the Underdark via the tunnel in area F10. The workshop machines soon begin to shake violently, causing the whole lair to tremble. One minute later, the steam cylinder in the workshop explodes, causing area F8 to cave in and dealing 70 (20d6) bludgeoning damage to any creatures still in that room.

\subsubsection{Death from Above}
\includegraphics[width=0.4\textwidth]{images/../5e.tools/img/adventure/QftIS/053-03-012.grisdelfawr.png}
The detonation in the forge draws the attention of Grisdelfawr, the \textbf{young red dragon} who recently took up residence in the Tegefed Mountains. Within moments of the explosion, Grisdelfawr begins circling over the misty gulch (area F1).

When the characters exit the forge, read or paraphrase the following text:

\begin{DndReadAloud}{}
As you exit the tunnel mouth, a blood-freezing cry rings from the sky above, followed by the baneful sound of leathery wings. A red dragon soars overhead. Leering hungrily at the valley below, it spots you and swoops down for the kill.

\end{DndReadAloud}

Grisdelfawr attacks the party on sight. Amoral and sadistic, the dragon fights without remorse. Grisdelfawr is astonished at the resilience of any who endure his breath, especially those who benefited from the magic of the kagu-svirfneblin's shimmering dome. If Grisdelfawr is reduced to 50 hit points or fewer, he retreats.

If characters stay inside the forge instead of heading outside to fight Grisdelfawr, the dragon waits outside for 1 hour—enough time for the characters to finish a short rest—then heads into the forge to devour anything with a pulse.

\section{Conclusion}
If the characters return to the Tower of the Heavens with the entire second set of Books of Prophecy, Shalfey is overjoyed. In exchange for the books, he rewards each character with a magic item drawn from the shimmering sphere in his library.

\subsection{Shalfey's Reward}
The shimmering sphere in the greater library (area T20 of the Tower of the Heavens) has existed since the tower's creation. No one quite remembers how it got there, but one thing is certain: only the elder sage of the Tower of the Heavens can reach into the sphere to pull forth magic weapons and items.

After receiving the books, Shalfey plucks an item for each character from the sphere, guided by its magic. Each character can choose one magic item from the following list:


\begin{itemize}
\item \textit{Amulet of Proof against Detection and Location}
\item Bag of Tricks, Gray
\item \textit{Eversmoking Bottle}
\item \textit{Gem of Brightness}
\item \textit{Pearl of Power}
\item \textit{+1 Weapon} (any)
\end{itemize}

Each item bears an intricate constellation motif that shimmers in the presence of starlight.

\chapter{Beyond the Crystal Cave}
\includegraphics[width=0.4\textwidth]{images/../5e.tools/img/adventure/QftIS/054-04-001.ch4-beyond-the-crystal-cave.png}
The city of Sybar on the small island of Sybarate serves as a hub for fishing and trade. Its people are industrious and friendly, and they eagerly share stories with visitors. One popular tale surrounds the Eternal Garden, a magical paradise created by two lovers ages ago. Legend states the entrance to the garden lies deep within a crystal cave that whispers secrets of the future to those brave or foolish enough to enter. More recently, two young paramours—one of whom is the daughter of the island's governor—fled their feuding families in search of the fabled utopia.

\includegraphics[width=0.4\textwidth]{images/../5e.tools/img/adventure/QftIS/055-04-002.eternal-garden-door.png}
The Eternal Garden is real. Past the dangers of the cave, the garden is tucked away in the Feywild, far from the pressures of the world. Whimsical Fey, talking animals, and a benevolent archfey ruler await visitors to the garden. So does the Fountain All Heal, a miraculous fountain whose waters cure most ailments but also enchant those who drink from it, lulling imbibers into a comfort that makes them never want to leave. None of the garden's denizens are inherently malicious, but they defend their own. Wits, clever words, and compassion are more effective weapons here than swords or spells.

The paramours' families, once long-standing rivals, have put aside their enmity in hopes that when the missing lovers are found, they will then be willing to come home. It's been two years since the couple's disappearance, but neither family has abandoned hope. The families turn to adventurers to rescue their missing loved ones.

"Beyond the Crystal Cave" is designed for four to six 6th-level characters.

\section{Background}
Long ago, an elf archdruid named Caerwyn retired to Sybarate with her soulmate, a human archmage named Porphura. There on Caerwyn's estates, they laid out and built a sprawling garden. It was an earthly paradise, serving both as a symbol of their enduring love and an idyllic retreat.

Caerwyn and Porphura stocked their garden with the loveliest of plants and encouraged gentle forest creatures to make it their home. Porphura used magic to extend her life so she could remain with Caerwyn throughout the elf's longer life span.

To protect their paradise, the lovers struck a bargain with an old friend, an archfey known as the Gardener. Together the three of them drew the garden into the Feywild, where the Gardener became its caretaker. The domain maintained a connection to the Material Plane through a fey crossing in a magical cave on Sybarate, the Cave of Echoes.

When Caerwyn's life reached its inevitable end, Porphura worked one last undertaking. She raised a tomb inside their home and sealed off the palace from the rest of the garden, so the pair could rest undisturbed together forever. Porphura then laid down beside Caerwyn for the last time and let the magic that sustained her own life finally unravel.

\subsection{Using the Infinite Staircase}
If you're using \textbf{Nafas} as a patron, he summons the characters to the Censer of Dreams (detailed in chapter 1), where he recounts the following wish:

\begin{DndReadAloud}{}
"Two families, their histories and hands stained with hate and blood, now raise their voices as one in the wake of true love lost. Heed their plea and carry their wandering hearts home."

\end{DndReadAloud}

Nafas then teleports the characters to a door that leads to Sybar. After the adventure, the characters can return to the staircase through the same portal.

\subsection{Adventure Hooks}
If you're not using Nafas as a group patron, consider the following ways to involve the characters in the adventure:


\begin{description}
\item[From the Sea.]
The characters arrive on a ship or wash ashore from a shipwreck. They learn the governor of Sybar is seeking adventurers to rescue her missing daughter.
\item[Missing Friends.]
One or more characters are old friends of one or both of the missing lovers, Juliana and Orlando. The governor sends for the characters, desperate for their help after a string of failed rescue attempts.
\end{description}

\subsection{Setting the Adventure}
On the world of Oerth in the Greyhawk campaign setting, the island of Sybarate is located in the mouth of Jeklea Bay, about a mile south of Fairwind Isle.

If you'd like to place the adventure in another campaign setting, consider the following suggestions:


\begin{description}
\item[Dragonlance.]
On the world of Krynn, Juliana and Orlando might be from rival kingdoms that don't typically get along, such as Qualinesti and Silvanesti. Amid the smoldering ashes of war, the story of the missing lovers can remind the heroes what they fight for and warn of the divisions that weaken the free peoples of Ansalon.
\item[Eberron.]
The cave is a manifest zone of Thelanis, the Faerie Court, and the Eternal Garden resides on that plane. Juliana and Orlando are scions of rival Dragonmarked houses, fleeing prohibition against their relationship.
\item[Theros.]
The Cave of Echoes is sacred to Karametra, a god of harvests who shepherds the mortal need for love, security, and belonging. She cradles the Eternal Garden in the godly realm of Nyx, and the native creatures that dwell there are Nyxborn.
\end{description}

\section{Adventure Summary}
"Beyond the Crystal Cave" begins when characters approach the governor of Sybarate, who is offering a small fortune for an urgent rescue mission. The characters learn the tale of Caerwyn and Porphura's garden, where two new lovers, Juliana and Orlando, have fled to escape their feuding families. Only one path leads into the hidden refuge: a cave with a mystical reputation. After besting the cave's guardians, the characters find a waterfall frozen in time and unravel its mysteries to discover a fey crossing into the Eternal Garden in the Feywild.

As the characters explore the garden, they likely meet the Gardener, the archfey who rules the domain; this is especially likely if they clash with any denizens of the garden, which abounds with eccentric Fey creatures. Various beings can offer cryptic clues to the missing lovers' whereabouts in a hidden palace within the garden. While in the garden, the characters might learn of a fountain that cures almost any ailment but enchants those who drink it to want to remain in the garden forever.

The adventure concludes when the characters free Juliana and Orlando from the fountain's enchantment and convince them that their families are no longer enemies. When they all arrive back on the Material Plane, they might discover they've been gone longer than they realized, thanks to the divergent flow of time within the garden.

\subsection{Character Advancement}
If you want to use story-based level advancement, the characters receive experience points for achieving milestones rather than defeating monsters.

When the characters leave the garden, with or without Juliana and Orlando, everyone in the party gains 1 level. If you follow this method, the characters should reach 7th level by the adventure's conclusion.

\includegraphics[width=0.4\textwidth]{images/../5e.tools/img/adventure/QftIS/056-04-003.beyond-the-crystal-cave-original.png}
\begin{DndReadAloud}{About the Original}
Published in 1983, \textit{Beyond the Crystal Cave} is the first adventure in the series of adventures produced by TSR UK. The story draws inspiration from the star-crossed lovers in Shakespeare's \textit{Romeo and Juliet}. It was a departure from the more expected adventure design of the time in that violent solutions to problems aren't an automatic path to victory—inquisitive minds and compassionate hearts carry the heroes further than strength of arms.

This updated version of the adventure reimagines Caerwyn and Porphura's—originally Porpherio's—garden as the Eternal Garden, a domain in the Feywild, and the Green Man as the Gardener, the benevolent archfey who rules it.

\end{DndReadAloud}

\section{Sybar}
\includegraphics[width=0.4\textwidth]{images/../5e.tools/img/adventure/QftIS/057-04-004.governor-isabela-folcarae.png}
A small fishing and port city, Sybar is the capital of Sybarate and the island's center of administration and trade. The governor's mansion and various executive offices are situated here, along with a temple to the sea god Procan.

Sybar is situated on the north coast of the island. Fishing and farming villages dot the remainder of Sybarate, which is lush with vineyards; vegetable farms; and extensive olive, lemon, and orange orchards. The people of the island are orderly and law abiding, but many of the old families hold longstanding grudges against one another.

Merchants and artisans are plentiful in Sybar, as are sailors and traders in the port. The locals tend to be talkative over meals, stiff drinks, or hot beverages on chilly evenings. If the characters mingle with the people of Sybar, residents can relate the story of the Eternal Garden (as presented in this adventure's background) or any of the rumors in the Sybar Rumors table, only some of which are true.

\begin{DndTable}[header=Sybar Rumors]{cX}

d8 & Rumor\\
1 & There is a magical fountain in the Eternal Garden that grants eternal youth. (False)\\
2 & It is always summer within the Eternal Garden. (True)\\
3 & Fire, both natural and magical, won't burn within the garden. (False)\\
4 & Sometimes people find their way into the garden, but most have little or no memory of it when they find their way out again. (True)\\
5 & The archmage Porphura isn't dead at all. She lives still as a lich who rules the island in secret. (False)\\
6 & Singing sea chanteys charms creatures in the garden. (True only for leprechauns; see appendix B)\\
7 & Eating or drinking the garden's bounty traps a person there forever. (Partially true; refers to the Fountain All Heal)\\
8 & There is a cave near the garden in which wishes are granted. (True; refers to the Cave of Echoes)\\
\end{DndTable}

\subsection{Meeting the Governor}
Governor Isabela Folcarae (lawful good, human \textbf{bandit captain}) receives the characters shortly after they call at her mansion. Read or paraphrase the following text to begin the adventure:

\begin{DndReadAloud}{}
You stand in a tastefully appointed study with nautical decor. The governor, an human woman with iron-gray hair, stares pensively into the fireplace. After a polite offer of refreshment, she gets down to business.

She speaks in a solemn voice. "My daughter Juliana, and Orlando, son of the prominent Monego family, have been missing for two years, but we have never given up hope that they still live. After several failed expeditions and countless nights of worry, I believe I finally know where they've gone. If you bring them home safely, I will reward you handsomely."

\end{DndReadAloud}

The governor can offer the following information:


\begin{itemize}
\item Juliana and Orlando eloped when their feuding families forbade their relationship.
\item The couple fled south and disappeared into the Cave of Echoes, a cavern steeped in local legends of oracular voices and granted wishes.
\item In their shared grief, the families have settled their differences and resolved to do all in their power to rescue the missing couple.
\item The families have sent numerous adventurers in search of the couple; the few that returned did so empty-handed.
\end{itemize}

Governor Folcarae can also relate the legend of the Eternal Garden, but she doesn't know details about its inhabitants or dangers—she doubts it even exists. In exchange for Juliana and Orlando's safe return, she promises a reward of 5,000 gp.

\section{Cave of Echoes}
The Cave of Echoes is a day's ride south of Sybar. Originally, the cave was the underground course of a river that flowed from a lake in the Eternal Garden before the garden was drawn into the Feywild. The caverns were worn smooth by the water that once rushed through them.

The cave entrance yawns in the southwest side of the small hill where the garden once stood. A fey crossing in the Cave of Echoes links it to the Eternal Garden and saturates the caverns with fey magic. The hilltop is now barren, and the river, which still flows through the portal, has been reduced to a mere trickle. The meager stream ends at a muddy chamber in the cave, depositing what little magic it still contains there.

Locals know of the Cave of Echoes, and many believe a specific chamber has oracular powers. Most Sybarate citizens either have sought answers from the cavern or know someone who did. The magic of the Feywild whispers into the cavern, offering aid to those of good heart and true need, as it did with Juliana and Orlando when it whisked them away to the garden. Few venture beyond this chamber, wary of dangers that lurk deeper within the cave.

\subsection{Cave of Echoes Features}
Unless otherwise noted, the Cave of Echoes has the following features:


\begin{description}
\item[Ceilings.]
Ceilings in chambers are 20 feet high and bristle with stalactites.
\item[Echoes.]
Sounds echo strangely throughout the cave, distorting into fleeting whispers.
\item[Lighting.]
The cave isn't illuminated. The current occupants rely on darkvision to see. Area descriptions assume the characters have a light source or other means of seeing in the dark.
\item[Tunnels.]
The tunnels connecting caverns are roughly round and 10 feet in diameter.
\end{description}

\subsection{Cave of Echoes Locations}
The following locations are keyed to map 4.1.

\includegraphics[width=0.4\textwidth]{images/../5e.tools/img/adventure/QftIS/058-map-4.01-cave-of-echoes.png}
\includegraphics[width=0.4\textwidth]{images/../5e.tools/img/adventure/QftIS/059-map-4.01-cave-of-echoes-player.png}
\subsubsection{C1: Empty Cavern}
\begin{DndReadAloud}{}
This vaulted cavern is forty feet across and twenty-five feet high in the center. The ground is flat and scuffed with muddy tracks, and the walls are marked with initials and names carved into the stone. Tunnels lead deeper into darkness to the east and the south.

\end{DndReadAloud}

The entrance cavern is empty save for the tracks of Humanoid entrants, most of whom stopped here before turning back. A few follow the tunnel to the south.

\subparagraph{Writing on the Walls}
Names and initials are inscribed on the walls. Some are the results of dares or tests of bravery; others were left by those marking their visit to the Cavern of Echoes (area C2). Older markings left in ink, charcoal, or paint are faded and illegible. A character who leaves their name or initials in the cave sometimes hears their name whispered among the echoes of the cave.

Characters who examine the walls find the initials "C + P" carved into the stone and circumscribed by a heart. A character who studies the wall and succeeds on a DC 14 Intelligence (Arcana) check determines these initials weren't carved by tools, but rather were magically engraved into the stone. Caerwyn and Porphura marked the wall using the \textit{Stone Shape} spell.

\subsubsection{C2: Cavern of Echoes}
\begin{DndReadAloud}{}
Even the smallest sounds echo in this round, forty-foot-wide chamber. The walls and floor are glassy smooth and black, refracting light into rainbows beneath the stone. The ceiling vanishes into darkness high overhead. In the center of the floor, inlaid with white marble, is the word "Ask" in Common.

\end{DndReadAloud}

The ceiling in this oracular chamber is 80 feet high and smooth like the walls and floor.

\subparagraph{Loud Echoes}
Sounds in this area echo loudly. Noises louder than a typical conversation—such as shouting or playing an instrument—reverberate painfully. If this occurs, each creature in the chamber must succeed on a DC 15 Constitution saving throw or take 13 (3d8) thunder damage and have the deafened condition for 1 minute. The echo then dissipates and doesn't compound.

\subparagraph{Whispers of Fate}
The magic of the Feywild whispers into the cavern, offering aid to those in true need and without selfish intent. When such a creature requests help in this area, the cavern can offer answers or intervene by duplicating any spell of 8th level or lower (your choice, not the creature's). The most common assistance provided by the cave comes in the form of the \textit{Greater Restoration}, \textit{Heal}, and \textit{Plane Shift} (to the Eternal Garden only) spells; when duplicating the effect of a spell, the cavern targets the creature that requested its aid.

A creature that asks a question here receives a truthful, albeit cryptic, answer in Common with clues to its goals. Those who can help themselves, such as the characters, receive only this benefit from the cavern. If the characters inquire about the whereabouts of Juliana and Orlando, disembodied voices tell them their quarry is just beyond the crystal cave (area C6). The cave's oracular power is similar to a \textit{Divination} spell.

A creature that makes a selfish or destructive request causes thunderous reprimands to reverberate through the chamber, triggering the effects of loud echoes (see the "Loud Echoes" section). Once a creature asks a question here and receives a benefit, no matter how small, the creature can't do so again for 1 year and 1 day.

\subsubsection{C3: Haunted Cavern}
\begin{DndReadAloud}{}
The floor of this cavern is ridged like sand at the bottom of a stream. It's dotted with slender stalagmites and littered with broken stalactites from the ceiling thirty feet above. The main tunnel continues to the east, and smaller side chambers extend to the north and the south.

\end{DndReadAloud}

Three poltergeists (variant specters) haunt the main chamber of the cavern, venting their fury on intruders. The restless spirits hurl broken stalactites at unsuspecting entrants or toss creatures onto pointed stalagmites in the cavern.

A cavity to the north collects moisture dripping from the ceiling. An \textbf{ochre jelly} lurks here, content to feed on fungus and the remains of the poltergeists' occasional victims. It slithers into the main cavern if battle breaks out.

\subparagraph{Stalagmites}
The rock formations throughout this cavern are dangerously sharp. A creature pushed into a stalagmite takes 11 (2d10) piercing damage.

\subparagraph{Treasure}
The skeleton of a decades-dead human knight rests in the south chamber, impaled on a stalagmite by the poltergeists. Most equipment the knight once carried has rusted or rotted away, but a \textit{Necklace of Adaptation} hangs around the dead warrior's neck. Scattered on the floor beneath the skeleton are 8 pp, 22 gp, and a yellow topaz (250 gp).

\subsubsection{C4: Living Mud Cavern}
\begin{DndReadAloud}{}
This cavern is large and irregular in shape. The bare rock floor forms a shallow depression that contains a mud pool fed by a sluggish stream flowing from the southeast. Near the west side of the pool are two mounds of dried mud.

\end{DndReadAloud}

The stream flowing from area C6 ends here in a 1-foot-deep mud pool. The mud is \textit{difficult terrain}.

The magic-infused water bestowed animate life on the mud, creating two mud elementals that lurk in the pool. A character who surveys the pool and succeeds on a DC 17 Wisdom (Perception) check notices the mud elementals.

Each mud elemental uses the \textbf{water elemental} stat block with the following changes:


\begin{itemize}
\item It has advantage on Dexterity (Stealth) checks made to hide in mud.
\item Its speed isn't reduced in \textit{difficult terrain} composed of earth or mud.
\end{itemize}

The elementals attack any creatures that approach the mud or the mounds of dried mud, but they don't pursue creatures out of the area and sink back into the mud if foes flee.

\subparagraph{Treasure}
The mounds encase the desiccated remains of two adventures slain by the mud elementals. The adventurers were sent by the governor to find Juliana and Orlando but never returned. Once the elementals are dead, weapons or tools can easily break away the mud shells covering the remains.

One mound entombs the remains of a human warrior. His \textit{chain mail}, \textit{shield}, and \textit{longsword} are rusted but still functional, and he has 20 gp in a crumbled purse.

The other mound holds the remains of an elf thief with a set of \textit{thieves' tools}, 50 gp, and 50 pp. Her shriveled hand clutches a \textit{+1 Dagger} with the name "Madrigal" carved into its ivory handle in Elvish.

\subsubsection{C5: Stream Tunnel}
\begin{DndReadAloud}{}
A narrow stream runs along the bottom of this winding tunnel. The echoes of the stream's soft babbling are unusually loud.

Partway along the stream's path, the waterway widens to a ten-foot-wide pool before it narrows again and continues.

\end{DndReadAloud}

The stream flows from the waterfall in area C6 to the mud pool in area C4.

\subparagraph{Treasure}
The pool is 3 feet deep, with a bottom of smooth rocks and pebbles that have washed down the stream. A character who searches the pool and succeeds on a DC 15 Wisdom (Perception) check realizes that seven of the smooth rocks are uncut gemstones: two garnets (10 gp each), two blue tourmalines (15 gp each), an amber (20 gp), a moonstone (25 gp), and a cloudy diamond (50 gp). A character who spends 10 minutes sifting through the river bottom succeeds on the check automatically.

A character who has proficiency with jeweler's tools can cut the gems with two days of work per gem. Doing so multiplies each gem's value by ten.

\subsubsection{C6: Crystal Cave}
\includegraphics[width=0.4\textwidth]{images/../5e.tools/img/adventure/QftIS/060-04-005.crystal-cave-waterfall.png}
\begin{DndReadAloud}{}
The echoes of the stream fade, replaced by an ethereal melody that tugs at your heartstrings. Massive formations of quartz crystals cover the walls and ceiling of this chamber.

A motionless waterfall pours from a crack near the ceiling of the east wall—a time-frozen curtain of water and hanging droplets. It collects into a wide, similarly static pool below. Farther from the water's source, the pool slowly begins to ripple, then drains into the stream that flows out the tunnel.

\end{DndReadAloud}

This crystal chamber is tethered to the Eternal Garden and by extension the Feywild. The waterfall isn't truly unmoving, but it pours so slowly it appears that way to casual observation. It takes 20 minutes for a droplet to fall from the ceiling into the pool.

\subparagraph{Crystals}
Any bright light in the room refracts through the crystals and fills the chamber with a soft glow and dazzling rainbow patterns.

\subparagraph{Music}
The disembodied melody carries the emotional resonance of the Feywild and evokes nostalgia in those who hear it. Characters who listen to the music can feel it tugging at memories from their pasts, whether joyful or otherwise.

The music causes momentary ripples in the seemingly frozen waterfall. A character who studies the waterfall and succeeds on a DC 15 Intelligence (Investigation) check notices the ripples and their connection to the music.

\subparagraph{Pool}
The pool is 5 feet deep. At the end the waterfall pours into, it too appears frozen in time, while at the far end, the water perceptibly flows into the stream. Within 10 feet of the waterfall, the pool is solid enough to walk on, though anything stationary on the water's surface slowly sinks to the bottom over the course of 1 minute.

\subparagraph{Waterfall}
The waterfall blocks passage to the fey crossing in area C7. Like the pool, the waterfall is nearly solid, and it blocks the tunnel except for a 1-inch gap on either side. The waterfall is 5 feet thick. A creature that looks through it with a light source can make out a passage hidden behind the falls.

Traversing the waterfall is difficult but can be accomplished through the following means:


\begin{description}
\item[Brute Force.]
As an action, a character can force their way through the falls with a successful DC 20 Strength (Athletics) check. A character who fails the check still manages to push through but gains 1 level of exhaustion from the exertion.
\item[Song.]
A character can create a temporary passage through the water by humming, singing, or playing a song. As an action, a character who performs a song and succeeds on a DC 15 Charisma (Performance) check creates a 5-foot-wide opening in the waterfall for 1 minute.
\item[Spells.]
Magic offers the following solutions:
\end{description}


\begin{itemize}
\item \textit{Control Water} parts or raises the waterfall for the spell's duration.
\item \textit{Create or Destroy Water} cast to destroy part of the waterfall destroys a 5-foot-cube of water, which fills in after 1 minute.
\item Teleportation bypasses the waterfall as normal.
\item A creature under the effect of a \textit{Slow} spell more closely matches the waterfall's altered time and can move through the waterfall as if it were \textit{difficult terrain}.
\item A successful \textit{Dispel Magic} spell (DC 15) cast on the waterfall causes it to flow normally for 1 minute, during which time creatures can move through the water as if it were difficult terrain.
\end{itemize}

At the DM's discretion, other solutions to bypassing the waterfall might exist.

\subsubsection{C7: Observation Parlor}
\begin{DndReadAloud}{}
Behind the waterfall is a room carved into the rock. Two steps ascend to a mosaic floor depicting birds, flowers, and fruit, which is littered with the long-collapsed remains of furniture and shards of broken glass.

A decorated alcove stands in the far wall opposite the entrance.

\end{DndReadAloud}

This was once a lounge Caerwyn and Porphura used to relax and enjoy the view of the crystal cave through the waterfall.

\subparagraph{Alcove}
This alcove is a fey crossing, a place of beauty where the Material Plane and the Feywild converge. Its walls are carved with a triptych woodland scene depicting forest animals and reveling satyrs, pixies, and sprites. The alcove radiates an aura of conjuration magic if observed with a \textit{Detect Magic} spell, and any creature that moves into it vanishes, transported to area G1 of the Eternal Garden (detailed later in this adventure).

\section{Eternal Garden}
Since being drawn into the Feywild, the Eternal Garden has become a Domain of Delight: a realm subject to the whims of the archfey who rules it. The overall composition and layout of the garden remain consistent from its days in the Material Plane, but Feywild magic has suffused it with a vibrancy beyond reality. Colors are more saturated, sounds are almost musical, and sensations pluck at a visitor's emotions, like the scent of a beloved childhood meal or the melody of a cherished tune.

Caerwyn and Porphura's magnificent home, the Palace of Spires, disappeared from the garden long ago. In the place where it once stood now looms an ever-changing hedge maze. Leaves growing throughout the garden act as keys to a magic sundial at the heart of the maze that transports those who solve it to the hidden palace.

\subsection{Garden Denizens}
Fey creatures, awakened animals and plants, and wanderers both Humanoid and otherwise call the garden home. Some of its long-standing residents were once visitors from faraway places who elected to stay in the garden, which transformed them into Fey over time. These sapient beings tend to form communities, though some prefer solitude, and they all are free to wander the garden as whim or errand dictates. The residents' demesnes are indicated on map 4.2.

\subsubsection{Awakened Animals}
Feywild magic has bestowed sapience on some of the Beasts that dwell within the Eternal Garden. Awakened animals in the garden speak Common and Sylvan.

\subsubsection{The Gardener}
The archfey who governs the Eternal Garden is known simply as \textbf{the Gardener}. The Gardener was a lifelong friend of Caerwyn and Porphura and gladly bound themself to the garden. Now, the Gardener lovingly tends the garden as much because it is their own home as to honor their dear friends' memories. The Gardener doesn't tolerate any attempt to damage or significantly alter the garden, nor do they accept the killing of any of its residents. Digging a hole, breaking branches, or foraging for food is fine. More intrusive or violent actions such as damming a river, clear-cutting a forest, or slaughtering pixies draw the Gardener's ire.

The Gardener demands an explanation for transgressions worthy of their attention, and if the characters aren't contrite, the Gardener attacks to knock them unconscious and leaves them at the mercy of the leprechauns' mischief (see area G1). This is a one-time chance for the characters to amend their behavior. If the characters continue to be a problem, the Gardener forcibly ejects them from the garden, teleporting the characters to the mouth of the Cave of Echoes. If the characters later return to the garden, its residents advise them to tread lightly.

\subsection{Love in the Garden}
Juliana and Orlando have spent two days in the Eternal Garden—equal to two years on the Material Plane—wandering the garden's meadows in awed delight, thrilled and relieved to have found a sanctuary where they can live and love each other in peace. The garden's Fey and awakened residents feel a glimmer of their lost friends Caerwyn and Porphura in the young couple; many denizens are convinced Juliana and Orlando are the garden's creators reincarnated. As such, they're fiercely protective of the couple and become hostile toward intruders who threaten the lovers.

\subsection{Random Encounters}
You can breathe life into the garden by introducing random encounters as the characters explore the realm. A random encounter occurs whenever you want one to, such as while the characters are traveling between locations. To determine what they find, roll on the Eternal Garden Encounters table. If you don't like the result, choose a different encounter you think would be fun.

\begin{DndTable}[header=Eternal Garden Encounters]{cX}

d8 & Encounter\\
1 & Two awakened giant badgers (see the "Awakened Animals" section) fret as they prepare their burrow for a visit from one of the badger's in-laws.\\
2 & Two leprechauns repair a large pile of worn-out shoes and boots. If the characters help, the leprechauns might give the characters a hint to one of their limerick riddles (see area G1).\\
3 & Three satyrs argue as they pick berries for making wine. They ask the characters to settle a dispute over which of the satyrs has the best muscles.\\
4 & Stargleam and Silverlily, the unicorns from area G15, discuss philosophy while on a walk.\\
5 & Gnarlroot, the irascible \textbf{treant} from area G12, lumbers through the garden, grumbling to himself about the antics of a prank-loving leprechaun.\\
6 & A fun-loving \textbf{awakened tree} uproots and blocks a character's path, mimicking their movements.\\
7 & Gorguth the \textbf{chimera} (see area G7) flies overhead and lets out a fearsome roar. If the characters draw Gorguth's attention, he descends to meet them.\\
8 & \textbf{The Gardener} strolls through the area, humming merrily to themself, inspecting plants, and greeting the garden's inhabitants.\\
\end{DndTable}

\subsection{Eternal Garden Features}
The Eternal Garden has the following features:


\begin{description}
\item[Boundary.]
At the border of the domain is a lightly obscuring mist, tinged with a sweet blossom scent and the colors of sunset. No matter how far a creature travels into the mist, the creature finds itself back where it entered the border, returning to the garden.
\item[Eternal Summer.]
The climate in the garden is a perpetual, lovely summer, but weather conditions reflect the emotional tone of creatures in the garden. In times of danger, thunderstorms shake the sky. During times of melancholy or mourning, chill rain patters throughout the domain.
\item[Gardener's Vigil.]
The Gardener senses when anything dies in the garden. If a creature dies in the garden, roll a d6. On a 1, the Gardener appears on initiative count 20 of the next round (losing initiative ties) to investigate.
\item[Lighting.]
The sky above the garden is a perpetual twilight range of orange, red, pink, and yellow hues with no visible sun, moon, or stars. During daylight hours, it emits bright light, but at night, it never reaches full darkness, becoming dim light instead.
\item[Memory Loss.]
Unless otherwise stated, a creature that leaves the domain must make a DC 10 Wisdom saving throw; Fey creatures and creatures with the Fey Ancestry trait automatically succeed on the saving throw. On a failed save, the creature remembers nothing of its time in the garden or the Palace of Spires. On a successful save, the creature's memories remain intact, but they're hazy and dreamlike. A \textit{Remove Curse} or \textit{Greater Restoration} spell restores the creature's memories.
\item[Time.]
Time passes slower in the Eternal Garden relative to other planes. For each day spent in the garden, one year passes on the Material Plane.
\end{description}

\subsection{Eternal Garden Locations}
The following locations are keyed to map 4.2.

\includegraphics[width=0.4\textwidth]{images/../5e.tools/img/adventure/QftIS/061-map-4.02-eternal-garden.png}
\includegraphics[width=0.4\textwidth]{images/../5e.tools/img/adventure/QftIS/062-map-4.02-eternal-garden-player.png}
\subsubsection{G1: Fairy Ring}
\begin{DndReadAloud}{}
A ring of toadstools twenty feet across stands in the center of a forest clearing. Bright green grass and colorful wildflowers blanket the ground. The toadstools range from the size of a finger to as large as a chair.

Five figures sit on the largest toadstools. All are diminutive, jovial creatures—no more than three feet tall—wearing dapper, red clothing with gleaming buttons and buckles. Each wears a different style hat and regards you with delighted anticipation.

\end{DndReadAloud}

When the characters pass through the fey crossing in the Cave of Echoes (see area C7), they vanish and reappear in the center of this fairy ring. The characters all appear at the same time, even if they entered the alcove at different times.

\subparagraph{Leprechauns}
The five creatures are leprechauns named Brunnick, Cellia, Flenas, Layleen, and Nerwyn. They are mischievous but fun loving, and all are genuinely delighted to see new visitors—mostly because they have new people to tease, as the rest of the garden's inhabitants have long grown weary of their antics.

\subparagraph{Toadstools}
The toadstools form a fairy ring that acts as a permanent teleportation circle, which can serve as the destination of teleportation magic and spells such as \textit{Plane Shift}. Though the fairy ring once provided an exit from the garden, the fey crossing is now one-way, allowing exit only with a special key (see the "Returning to the Garden" section of this adventure's conclusion for more on exiting through the fairy ring).

\subparagraph{What the Leprechauns Know}
The leprechauns greet newcomers happily, responding to questions about themselves, the garden, and Juliana and Orlando with one or more of the cryptic rhymes shown in the Limericks table.

If the characters ply the leprechauns with amusements (leveraging the leprechaun's Reluctant Refusal trait), the charmed leprechauns let slip some hints to their riddles or their belief that the young lovers are the incarnations of the garden denizens' beloved friends, Caerwyn and Porphura.

If the leprechauns are attacked, they confuse foes with illusions or turn invisible and flee. They take care to return later to plague the characters with irksome pranks such as the following:


\begin{itemize}
\item Steal a weapon or magic item
\item Steal a component pouch or spellcasting focus
\item Fill armor, clothes, or packs with itching powder
\end{itemize}

\includegraphics[width=0.4\textwidth]{images/../5e.tools/img/adventure/QftIS/063-04-006.leprechauns.png}
\begin{DndTable}[header=Limericks]{XX}

Limerick & Describes\\
The river o'er yonder appearing, Has five bridges of which you're hearing. The problem we set, Is how do you get, Across each and return to this clearing? & Endless River (area G4)\\
A fountain rests here in the garden, Whose waters, we hear, like to harden. But a drink now and then, Cures all in our ken, And may even bring cries of "Pardon!" & Fountain All Heal (area G9)\\
There's a lovely building just here, With turrets and domes far and near. 'Tis a shame you can't spy it, As we urge you to try it. Solve the puzzle and see it appear! & Palace of Spires (detailed later in this adventure)\\
Some may think our maze is a pest, But it shelters a haven of rest. And ere anyone leaves, They must leave more than leaves, And then honor our master's request. & Hedge Maze (area G21) and Palace of Spires (detailed later in this adventure)\\
There are leaves in the garden to trace, And a maze with a clear central space. With a leaf in your hand, You could stand and stand, But to leave, leave the leaves in place. & Hedge Maze (area G21)\\
There once was a young man from Sybar, Who was stung on the nose by a wasp. We said, "Does it hurt?" He replied, "Not at all, It can do it again if it likes." & Nonsense, but it sends the leprechauns into gales of laughter and knowing winks\\
\end{DndTable}

\subsubsection{G2: Leprechaun Woods}
\begin{DndReadAloud}{}
This lush forest brims with game and edible plants. A massive oak tree, the only oak in this region, looms above the rest of the forest. Among the undergrowth, shamrock plants thrive.

\end{DndReadAloud}

The leprechauns in area G1 often roam this forest. They hide here after pulling pranks on the garden's denizens.

\subparagraph{Key Leaves}
The shamrock leaves that grow in this area function as a key to the teleportation sundial in the maze (area G21).

\subsubsection{G3: Great Oak Tree}
\begin{DndReadAloud}{}
A tremendous oak tree stretches hundreds of feet into the air, rising high above the surrounding canopy. Easily fifty feet wide at the base, its trunk is gnarled, covered with creeping greenery, and surrounded by flowering shamrocks.

\end{DndReadAloud}

Several leprechauns live inside this great oak. The interior is sized for Small creatures and is brightly lit by cage lanterns full of fireflies. At any given time, 1d4 friendly leprechauns are here, enjoying the comforts of their home.

\subparagraph{Hidden Home}
An illusion hides the front door to the tree house at ground level. A character who studies the tree trunk and succeeds on a DC 14 Intelligence (Investigation) check discerns the illusion for what it is.

The first floor of the tree home is a circular kitchen with wooden furnishings and simple flatware for five. It's stocked with nutritious food, water, and wine. A spiral staircase carved into the interior wall of the trunk rises to a bedroom with five crisply made beds arranged radially around the wall. A chest at the foot of each bed holds fashionable clothing—most of it red—and blankets.

\subparagraph{Kitchen Pit}
Beneath the kitchen table is a secret trapdoor. A character who searches the kitchen and succeeds on a DC 15 Wisdom (Perception) check discovers the door. It leads to a tunnel that ends in a cave-in with several crushed humanoid skeletons visible at its edge, but this is an illusion concealing a 6-foot-deep pit. A character who searches the tunnel and succeeds on a DC 14 Intelligence (Investigation) check, or who tries to interact with the illusion physically, discerns the illusion for what it is.

\subparagraph{Key Leaves}
The shamrock and oak leaves that grow in this area function as keys to the teleportation sundial in the maze (area G21).

\subparagraph{Treasure}
A character who digs at the bottom of the kitchen pit discovers a pouch with an Ioun Stone, Sustenance and two dull-gray stones that lost their magic long ago. Digging deeper reveals a copper crock filled with 2,500 gp—the leprechauns' collected gold. If the gold is stolen, the leprechauns stop at nothing to retrieve it.

\subsubsection{G4: Endless River}
\begin{DndReadAloud}{}
This babbling river flows through the garden, sometimes defying logic. Here and there, its waters seemingly course in different directions, and it teems with colorful fish and other freshwater fauna.

\end{DndReadAloud}

\includegraphics[width=0.4\textwidth]{images/../5e.tools/img/adventure/QftIS/064-04-007.vuuthramis.png}
The river is home to Vuuthramis, a shy but playful \textbf{young bronze dragon}. She lives in a cave hidden beneath the river. Vuuthramis drank from the Fountain All Heal (area G9) and succumbed to its enchantment. As a result, she has no desire to keep a treasure hoard. She's content to swim and frolic in the river, occasionally venturing onto its banks.

Vuuthramis cautiously observes characters near the riverbank. Characters who survey the river or enter its waters and succeed on a DC 14 Wisdom (Perception) check spot the dragon, who keeps her distance. If a character reacts peacefully and succeeds on a DC 15 Charisma (Persuasion) check, Vuuthramis overcomes her reticence. Otherwise, she darts away, only to return later.

\subparagraph{Capricious River}
The river ranges from 10 to 20 feet in depth and 10 to 50 feet in width. Its flow is unpredictable, reversing at various points along its ovular length. Roll a die for any 100-foot stretch of river. On an even number, it flows clockwise; on an odd number, it flows counterclockwise.

\subparagraph{What Vuuthramis Knows}
Vuuthramis watches and listens along the course of the river and knows the surrounding inhabitants. If a character overcomes the dragon's shyness, she tells the characters about a few of the garden's denizens, such as Gnarlroot, the \textbf{treant} in area G12, who "seems extra grumpy today."

In addition to this general knowledge, Vuuthramis witnessed Juliana and Orlando stop to picnic with a group of pixies and sprites from the fairy forest (area G10) on the riverbank two days earlier, before the pair continued west.

\subsubsection{G5: Dryads' Woods}
\begin{DndReadAloud}{}
This forest is dense and difficult to travel, with only a few worn footpaths connecting its small clearings. Most of the trees here are old, mighty oaks.

\end{DndReadAloud}

Two dryads, Sonna and Yvra, tend to this stretch of woods. Sonna and Yvra were once knights in service to rival Fey courts who fell in love. As punishment for defying their rulers to pursue a relationship, they were bound to oak trees and separated. \textbf{The Gardener} heard their lamentations, and when Caerwyn and Porphura approached the archfey to take charge of the domain, the Gardener stole Sonna and Yvra's trees from their jailers and planted them here, where the dryads could finally be together.

The dryads enjoy playing card games, especially a simple bluffing game known as myriad. If the characters behave peacefully, the dryads challenge a willing character to a quick round. To win the game, a character must succeed on a DC 17 Charisma (Deception) check to out-bluff the dryads or a DC 17 Wisdom (Insight) check to catch them in a lie.

\subparagraph{What the Dryads Know}
The dryads correctly suspect Juliana and Orlando are in the Palace of Spires but are hesitant to say anything. They see themselves in the missing lovers and won't risk bringing them to harm. However, if a character beats Sonna and Yvra at myriad, the dryads reluctantly suggest that Gnarlroot, the \textbf{treant} in the south ravine (see area G12), might know more.

\subparagraph{Key Leaves}
The oak leaves here are a key to the teleportation sundial in the maze (area G21).

\subparagraph{Treasure}
Fluttering in the branches of Sonna's tree is a \textit{Cloak of Elvenkind}. Yvra's tree has the set's matching \textit{Boots of Elvenkind} hidden in a hollow. One of the boots holds a leather pouch containing thin ivory plates whose illustrations resemble those of the storied \textit{Deck of Many Things}. The deck is nonmagical and worth 100 gp. Yvra uses them to play a solitary card game called patience.

\subsubsection{G6: Satyrs' Forest}
\begin{DndReadAloud}{}
This section of the forest is less dense, allowing light to dapple the ground. Clearings and paths abound, many of which show signs of regular use: discarded bottles, trampled earth, and ashes from old bonfires. The clearing ahead thrums with music and laughter.

\end{DndReadAloud}

Five satyrs revel in this forest. The largest clearing in the woods is surrounded by flowering shrubs and has a single towering elm tree in the center. The satyrs live here, sleeping on pallet beds surrounding the elm tree. Their leader, Xanim, carries a set of panpipes and leads their carousing.

Xanim and his kin (Kren, Vella, Zamuel, and Liss) are incorrigible gossips. They gleefully greet any newcomers and urge them to join the revelry.

\subparagraph{Trapped Hollow}
The satyrs' treasure is tucked into a hollow formed by the elm's gnarled roots and guarded by a hidden noose snare. A creature that steps close enough to search the root hollow must succeed on a DC 16 Dexterity saving throw or be hoisted 20 feet into the air. The creature has the restrained condition until it gets free by cutting the rope or slipping out as an action with a successful DC 16 Dexterity (Acrobatics) check. A character who surveys the hollow from a distance and succeeds on a DC 16 Wisdom (Perception) check discovers the trap, which can be disarmed by poking it with a stick or tool to set it off.

\subparagraph{What the Satyrs Know}
If the characters share any juicy tidbits of gossip with the satyrs, such as the name of a secret crush, the satyrs respond with stories about other denizens of the garden, particularly the leprechauns. The satyrs show open annoyance at the leprechauns' pranks and hint that by offering the leprechauns a dance or fine meal, the characters might be able to persuade them to part with more information. The satyrs know of Juliana and Orlando but haven't met them; they eagerly look forward to partying with the "new arrivals."

\subparagraph{Treasure}
In the tree hollow is a leather sack containing 400 gp and a locked wooden box. As an action, a character can pick the lock with a successful DC 20 Dexterity (Sleight of Hand) check using thieves' tools. The box holds three sapphires (200 gp each), a \textit{Potion of Animal Friendship}, a \textit{Potion of Diminution}, a Potion of Greater Healing, and a set of fine reed pipes worth 25 gp.

\subsubsection{G7: Chimera's Crag}
\begin{DndReadAloud}{}
A rocky outcropping runs along the northern border of the garden. A small ravine cuts into the rocks, its floor scorched free of vegetation and dusted with ash.

\end{DndReadAloud}

Gorguth the chimera drank from the Fountain All Heal and became enchanted by it. The chimera now lairs on a crag at the back of the ravine. Over the years, the magic of the garden kindled a spark of understanding in the creature. He is indifferent toward the garden's denizens and peaceful visitors.

Gorguth is a \textbf{chimera} with the following changes:


\begin{itemize}
\item His alignment is chaotic neutral.
\item He speaks Common and Draconic.
\end{itemize}

Gorguth enjoys flying over the garden, occasionally performing threatening but harmless swoops near anyone he sees. When he spots the characters, he makes an intimidating display, scorching the ground with his fire breath before landing on his crag to observe their reactions. Each character must succeed on a DC 15 Wisdom saving throw to keep their composure. Characters immune to the frightened condition automatically succeed on this save. If any characters keep their composure, Gorguth becomes friendly and willing to chat. If every character fails the save, the chimera laughs and returns to the sky.

\subparagraph{What Gorguth Knows}
Gorguth observed Juliana and Orlando during their travels and knows they've spent a great deal of time talking with \textbf{the Gardener}, who lives in a cave across the lake (see area G20).

\subsubsection{G8: Bears' Range}
\begin{DndReadAloud}{}
A cave toward the northern edge of the garden overlooks a swath of wildflower-stippled grassland.

\end{DndReadAloud}

A mated pair of awakened brown bears (see the "Awakened Animals" section) and their awakened cub (use the \textbf{black bear} stat block) live in the cave and claim the surrounding area as their territory. The adult bears are hostile toward strangers but not violent unless provoked. Suspicious of visitors, they warn trespassers to back away from their den. The cub is indifferent—curious but skittish.

\subparagraph{What the Bears Know}
The bears' attitudes improve by one step if the characters offer them food or if a character explains the party's presence and succeeds on a DC 12 Charisma (Persuasion) check.

The cub observed Juliana and Orlando two days ago. She followed them to the Fountain All Heal (area G9), where the pair quenched their thirst. If the characters improved the cub's attitude, she gladly agrees to lead the characters to the fountain.

\subsubsection{G9: Fountain All Heal}
\includegraphics[width=0.4\textwidth]{images/../5e.tools/img/adventure/QftIS/065-04-008.fountain-all-heal.png}
\begin{DndReadAloud}{}
A spring of crystal-clear water bubbles up from the center of a natural rock basin. The water trickles down one side of the fountain, forming a small stream. Pink dog-rose blossoms grow around the basin.

\end{DndReadAloud}

Caerwyn and Porphura worked together to create this miraculous fountain; over the ages, entrancing Feywild magic has permeated its waters. The basin is 4 feet wide and 3 feet deep. The water spilling from the basin streams down the slope toward the lake (area G18).

\subparagraph{Curative Water}
The fountain's waters are cool, pure, refreshing—and a potent magical elixir. A creature that drinks from the basin gains the benefits of both the \textit{Greater Restoration} spell (all its benefits) and the \textit{Heal} spell, and any entities possessing the drinker are forced out.

Any water scooped from the basin must be drunk immediately to benefit from its properties. If bottled or removed from the fountain, the water loses its magic and turns into gravel.

\subparagraph{Entrancing Draught}
A creature that drinks from the fountain must succeed on a DC 18 Wisdom saving throw or permanently lose all desire to leave the Eternal Garden. Fey and creatures with the Fey Ancestry trait automatically succeed on the saving throw. An affected creature resists any attempt to forcibly remove it from the garden. If the creature is removed from the garden, it is overcome with an irresistible desire to return. This effect can be removed only by a \textit{Wish} spell (such as those granted in the Palace of Spires, detailed later in this adventure); a carefully worded wish can free multiple creatures from the water's trance.

\subparagraph{Key Leaves}
The dog-rose leaves here are a key to the teleportation sundial in the maze (area G21).

\subsubsection{G10: Fairy Forest}
\begin{DndReadAloud}{}
Harmless insects abound in this forest. Butterflies, fireflies, and ladybugs flutter about, filling the canopy and undergrowth with flitting colors. A great willow tree atop a low hill sparkles with fairy lights and bioluminescent bugs.

\end{DndReadAloud}

This stretch of forest is home to a gathering of Fey folk. Six pixies and six sprites are present in and around the willow tree at any time. Though there are dozens of the Fey in total, most of them are gallivanting about the woods and the greater garden.

When the characters first enter the fairy forest, the six sprites here sneak up invisibly and attempt to learn more about the visitors. The sprites and pixies are hostile to any evil creatures they discover and try to drive such foes away.

\subparagraph{Fairy House}
The pixies and sprites live in an abandoned badger burrow in the willow hill. Only Tiny creatures can fit through the foot-tall fairy door set into the burrow's entrance. The Fey furnished and decorated the burrow like a tiny house; it has four rooms.

\subparagraph{What the Fairies Know}
To non-evil beings, the pixies and sprites are perfectly amiable, revealing themselves and welcoming such characters to their home. If the characters inquire about the garden or any of its residents, the Fey huddle up and whisper among themselves before a spokes-pixie cheerfully announces they'll happily help in exchange for a gift. The pixie accepts any magic item, anything shiny or sparkly, or anything with a story imparted by the character who offers it.

After receiving an acceptable gift, the fairies disclose that several of them spent some time at a picnic with Juliana and Orlando on the couple's first day in the garden and eventually led them to the Fountain All Heal (area G9) to quench their thirst.

\subparagraph{Treasure}
Hidden in the fairies' burrow are a \textit{Potion of Clairvoyance} and a gold ring (250 gp) set with a \textit{Pearl of Power}.

\subsubsection{G11: Barkburrs' Grove}
\includegraphics[width=0.4\textwidth]{images/../5e.tools/img/adventure/QftIS/066-04-009.barkburrs.png}
\begin{DndReadAloud}{}
A dense grove here consists of eerie oak trees. From the corner of the eye, the trees resemble looming figures with crooked fingers. A singularly gnarled and ancient oak rises from the center of the grove. Clusters of mistletoe grow all over its branches, and something sparkles between its roots.

\end{DndReadAloud}

The smaller oak trees here evoke looming figures because they were all once people before they were transformed into trees by the two barkburrs that cling to the central oak. The barkburrs attack any creature that comes within 20 feet of the tree. \textbf{The Gardener} doesn't punish creatures that defend themselves against the barkburrs, knowing more barkburrs will grow off the knobby oak in short order.

A creature petrified by the barkburrs can be returned to its true form by dunking its roots in the Fountain All Heal (area G9). It takes 1 minute for the transformed creature's roots to "drink" from the fountain's waters.

\subparagraph{Key Leaves}
The oak and mistletoe leaves here are keys to the teleportation sundial in the hedge maze (area G21).

\subparagraph{Treasure}
A sheathed \textit{+1 Longsword}with a gem-encrusted hilt pokes up from the central oak's roots. The sword sheds light as a torch when drawn. Hanging from the high branches of a smaller oak is a black opal pendant (1,500 gp). A character who searches the grove and succeeds on a DC 16 Wisdom (Perception) check notices the pendant.

\subsubsection{G12: Treant's Stand}
\begin{DndReadAloud}{}
A limestone-walled ravine cuts into a rise in the land. A stream from the lake wends toward a low cave to the west. Trees, shrubs, and creepers grow near mossy boulders toward the back of the ravine to the east.

\end{DndReadAloud}

A curmudgeonly and venerable \textbf{treant}, Gnarlroot, lives near the cave in the back of the ravine. He watches newcomers from afar and reveals himself only if they damage any plant life or approach his cave. When Gnarlroot emerges, he grumpily tells the characters to respect other people's homes.

\subparagraph{Hidden Cave}
A mossy cave hidden toward the back of the ravine holds a simple table and stool and a rack full of gardening tools. The furniture is for the Gardener's benefit, and both they and Gnarlroot use the tools to tend various plants in the garden. Characters who search the ravine and succeed on a DC 15 Wisdom (Perception) check find the cave behind a curtain of ivy.

\subparagraph{What Gnarlroot Knows}
Gnarlroot often converses with the Gardener when the archfey stops by. He knows about Juliana and Orlando thanks to a recent visit from the archfey. If the characters agree to get off his lawn, Gnarlroot tells them the paramours were seen gathering leaves throughout the garden before entering the hedge maze (area G21).

\subparagraph{Treasure}
Hidden in a leather sack under some mossy rocks in the cave is a \textit{Potion of Growth}. A character searching the cave finds the stash with a successful DC 20 Intelligence (Investigation) check.

\subsubsection{G13: Cave of the Toads}
\begin{DndReadAloud}{}
A stream fed by the lake runs into a cave in the southern wall of a limestone ravine. The cave mouth is twenty feet wide and fifteen feet high.

\end{DndReadAloud}

The stream runs into a crack in the rock at the back of the 30-foot-deep cave. Two giant toads lair in the cave, guarding their clutch of eggs. The toads are indifferent to the characters unless the characters come more than 10 feet into their den, at which point the toads croak angrily in defense of their eggs. A character can use an action to try to calm the toads, doing so with a successful DC 15 Wisdom (Animal Handling) check. Once calmed, the toads allow the characters to explore the cave unhindered. The toads attack if their eggs are threatened.

\subparagraph{Treasure}
Wedged into the crack where the stream begins is a+1 Shield emblazoned with the image of a brightly painted frog.

\subsubsection{G14: Centaurs' Meadow}
\begin{DndReadAloud}{}
A meadow near the lake bursts with color from the wildflowers blanketing the grass. A small wooden structure resembling a cross between a cottage and a stable stands among the blossoms.

Bees buzz happily between the flowers and around a great hive enclosing the rear corner of the structure.

\end{DndReadAloud}

The meadow and structure are home to a family of four centaurs. The parents, Heraq and Persella, are initially wary of the characters and tersely ask them to state their business. The centaur children, Kari and Vinz, spend most of their time running through the meadows and swimming in the lake with the giant otters (area G18). If the characters are respectful, the centaurs invite them inside for tea sweetened with honey.

The bees are docile unless their hive is damaged. If a creature damages the beehive, a swarm of insects (wasps) erupts from the hive and attacks until the offending creature dies or flees.

\subparagraph{Shelter}
The centaurs' home has an open-air common room with a large trestle table, cooking utensils, and fresh fruit and nuts. Three bedrooms behind it house the family, two for the children and one for their parents. The rooms are appointed with straw mattresses and lockers containing brushes and simple upper-body clothing.

\subparagraph{What the Centaurs Know}
Kari and Vinz saw Juliana and Orlando enter the hedge maze (area G21) yesterday, but the couple never returned.

\subparagraph{Treasure}
A character can carefully harvest a single dose of magical honey from the hive without upsetting the bees with a successful DC 10 Wisdom (Survival) check. The honey collected functions as a \textit{Potion of Vitality}.

\includegraphics[width=0.4\textwidth]{images/../5e.tools/img/adventure/QftIS/067-04-010.stargleam-and-silverlily.png}
\subsubsection{G15: Unicorns' Meadow}
\begin{DndReadAloud}{}
A rolling, green expanse of grass rises from the east shore of the lake to the color-flecked forest and steep cliff at the domain's eastern border.

\end{DndReadAloud}

Two mated unicorns named Stargleam and Silverlily live here. They enjoy walking the paths of the hedge maze (area G21) and occasionally assist \textbf{the Gardener} in protecting the domain. It's been decades since their help was required, but they remain vigilant and intensely question newcomers to the garden.

\subparagraph{What the Unicorns Know}
The unicorns are initially standoffish and try to get the measure of anyone pumping them for information with a philosophical debate, such as when it's okay to lie to someone you love. A character who engages the unicorns in their debate and succeeds on a DC 15 Charisma (Persuasion) or Intelligence (History) check earns the unicorns' grudging respect. Thereafter the unicorns are willing to discuss the garden and its inhabitants, including the entrancing power of the Fountain All Heal (area G9). The unicorns then reveal Juliana and Orlando are under its sway.

\subsubsection{G16: Yellow Rose Arbor}
\begin{DndReadAloud}{}
Walls of thorny hedges stand ten feet high in this meadow. An opening in the west wall leads into an arbor with thriving rose bushes covered in bright yellow roses.

\end{DndReadAloud}

Two hedge-walled arbors mark the east meadow. The northern arbor grows yellow roses, while the southern arbor (area G17) grows red roses.

\subparagraph{Hedges}
The hedge wall is 5 feet thick and immune to all damage. A creature can force its way through the wall, spending 4 feet of movement for each foot traveled. Creatures that push through a hedge wall must make a DC 15 Dexterity saving throw, taking 31 (7d8) piercing damage from its thorns on a failed save or half as much damage on a successful one.

\subparagraph{Roses}
The lovely scent of the roses fills the arbor inside the hedge walls, causing a drowsy euphoria in creatures that inhale it. A creature that breathes the scent must succeed on a DC 15 Constitution saving throw or have the poisoned condition for 24 hours. A creature that succeeds on the save is immune to the roses' perfume for 24 hours.

\subsubsection{G17: Red Rose Arbor}
Aside from its position and the color of the roses—the ones here are a deep, rich red—this arbor is identical to area G16.

\subsubsection{G18: Lake}
\begin{DndReadAloud}{}
A placid lake fed by a forest stream to the north reflects the colorful sky. Gossamer-winged dragonflies dance among the tall reeds and plants at the shore, and the lake's surface ripples with the splashes of playful otters. A small boat is tied to a stone jetty along the western shore.

A long, white stone building with a translucent roof and open archways rises from the center of the lake.

\end{DndReadAloud}

Two friendly giant otters (giant weasels with a swimming speed of 40 feet) frolic in this shallow lake, which slopes gently from the shore to a depth of 10 feet. The otters are curious and playful; they dart up to investigate anyone in the lake.

The otters delight in bumping small boats and tugging on their rails. Each time an otter bumps a boat, creatures in it must succeed on a DC 10 Dexterity saving throw or fall into the water. The otters are easily distracted by anyone tossing them small objects; they play with such objects in their paws, pass the objects between each other, or play fetch with whoever threw the objects.

\subparagraph{Boat}
A wide rowboat boat is tied to the jetty on the western bank. The boat can comfortably hold up to six Medium creatures and has two oars.

\subsubsection{G19: Water Palace}
\begin{DndReadAloud}{}
At the center of the lake, a marble building gleams with gilded fittings. It has tall, glazed windows and is dotted with elongated archways. A veranda runs around the outside of the building ten feet above the water. Two stairways with ornate rails descend to a stone pier just above the lake's surface.

\end{DndReadAloud}

The water palace is a lovely place for relaxation and quiet meetings. It's open to all.

\subparagraph{Main Floor}
The main room is open and airy, with a long white table in the center surrounded by four padded chairs. A crystal decanter and four matching glasses sit on a silver tray.

\subparagraph{Otter Den}
The otters have a den in the flooded "basement" of the palace. A hole in the wall under the west side of the docking platform allows them to swim in and out of the den. The floor is covered in gravel, fish bones, and shells. Stairs ascend to a loft with a nest and four playful baby giant otters (weasels with a swimming speed of 30 feet).

\subsubsection{G20: The Gardener's Rest}
\begin{DndReadAloud}{}
A ten-foot-wide tunnel in a steep rise is supported by wooden beams and a lintel, which are covered in the pink blossoms of climbing dog-rose plants and hanging wisteria. A well-trodden grass path between verdant shrubs and wildflowers leads to the opening.

\end{DndReadAloud}

This is the Gardener's home. The tunnel leads into a roughly hewn stone room with terraced shelves and many soil beds on the floor. Plants, mushrooms, moss, and lichen of all kinds grow in the beds. A door in the back wall leads to a room with a moss bed and stacks of books from other worlds, all discussing botany and nature.

\textbf{The Gardener} is present the first time the characters arrive. The archfey is friendly toward the characters and curiously inquires as to their business in the domain. After some initial investigation of the characters' intentions, the Gardener offers refreshment. If the characters have behaved peacefully toward the garden's denizens, the Gardener also casts the \textit{Heroes' Feast} spell to produce a repast of delicious vegetables for the characters.

If the characters attack the Gardener, the archfey defends themself, all the while admonishing the characters for their rudeness and offering several chances to have a more polite discussion.

\subparagraph{What the Gardener Knows}
The Gardener is happy to answer general questions about the garden and its creators. If the characters are respectful, the archfey confirms that Juliana and Orlando are present in the domain. If asked either where the lovers are now or how to leave the garden, the Gardener encourages the characters to solve the hedge maze (area G21) to find what they're looking for.

\subparagraph{Key Leaves}
The dog-rose leaves here are a key to the teleportation sundial in the maze (area G21).

\subsubsection{G21: Hedge Maze}
\begin{DndReadAloud}{}
A fifteen-foot-high hedge wall with wickedly long thorns rises from the meadow. Three of the walls are solid; an opening in the western wall leads into a maze of hedgerow corridors.

\end{DndReadAloud}

This shifting hedge maze surrounds the means to enter the Palace of Spires: a magic sundial.

\subparagraph{Maze}
Each time creatures enter the maze, the path through is different. To successfully navigate the maze, a creature must succeed on a DC 15 Intelligence (Investigation) or Wisdom (Survival) check or reemerge at its starting point. Groups that enter the maze together make this check as a group check. A creature can create a more direct route by pushing through the walls, but it risks running afoul of the hedges' thorns (see the "Hedges" section below). A straight path between the center of the maze and its starting point has 1d3 + 1 hedge walls in the way.

Residents of the garden and any creatures that have succumbed to the entrancement of the Fountain All Heal (area G9) instinctively know the correct path through the maze, automatically succeeding on checks to navigate it.

\subparagraph{Hedges}
The hedge wall is 5 feet thick and immune to all damage. A creature can force its way through the wall, spending 4 feet of movement for each foot traveled. Creatures that push through a hedge wall must make a DC 15 Dexterity saving throw, taking 31 (7d8) piercing damage on a failed save or half as much damage on a successful one.

\subparagraph{Sundial}
In the center of the maze is a rose-quartz sundial with a silver gnomon. The dial always shows noon during the day. Leaf-shaped impressions carved into the four corners of the sundial's face depict plants found throughout the garden: dog rose, mistletoe, oak, and shamrock. A character can identify the leaves with a successful DC 15 Intelligence (Nature) or Wisdom (Survival) check.

The leaves act as keys to the sundial. If all the appropriate leaves are placed in their respective impressions, the sundial glows, and every creature in the center of the maze is transported to area P1 of the Palace of Spires (detailed in the following section). The leaves are then consumed.

The Sundial Keys table outlines where each type of key leaf can be found in the garden.

\begin{DndTable}[header=Sundial Keys]{cX}

Key Leaf & Location\\
Dog rose & Fountain All Heal (area G9), Gardener's Rest (area G20)\\
Mistletoe & Barkburrs' Grove (area G11)\\
Oak & Great Oak Tree (area G3), Dryads' Woods (area G5), Barkburrs' Grove (area G11)\\
Shamrock & Leprechaun Woods (area G2), Great Oak Tree (area G3)\\
\end{DndTable}

\section{Palace of Spires}
\includegraphics[width=0.4\textwidth]{images/../5e.tools/img/adventure/QftIS/068-04-011.palace-of-spires.png}
Porphura constructed this magnificent palace to house herself and Caerwyn, her beloved. The palace once stood where the hedge maze (area G21) now does, facing the lake in the garden. Composed of striking purple-red porphyry stone, the palace remains a comfortable domicile for those who find their way here now.

After Caerwyn's death, Porphura built a tomb for them both in the palace's crystal dome. With the Gardener's help, the grieving wizard performed her last great work: drawing the entire structure into a demiplane accessible only through the sundial at the heart of the maze. When all was arranged, she laid down with Caerwyn for the final time and let the magic that preserved her own life beyond its natural span end. The Gardener said their final farewell to Caerwyn and Porphura, sealed the tomb, and left the palace forever. Out of respect for their late friends, the Gardener won't enter the Palace of Spires for any reason.

\subsection{Palace of Spires Features}
Unless otherwise specified, the Palace of Spires has the following features:


\begin{description}
\item[Ceilings.]
Ceilings in the palace are vaulted and 40 feet high.
\item[Indestructible.]
The palace and its furnishings are immune to all damage and can't have their forms changed.
\item[Lighting.]
Soft light permeates the demiplane. Light outside the palace matches that of the Eternal Garden (see the "Eternal Garden Features" section), while interior areas are brightly lit by \textit{Continual Flame} spells.
\item[Secret Doors.]
Some doors in the palace blend in with their surroundings. A character who searches for a secret door and succeeds on a DC 20 Wisdom (Perception) check finds it.
\item[Time.]
Time in the demiplane moves at the same rate as in the Eternal Garden.
\end{description}

\subsection{Palace of Spires Locations}
The following locations are keyed to map 4.3.

\includegraphics[width=0.4\textwidth]{images/../5e.tools/img/adventure/QftIS/069-map-4.03-palace-of-spires.png}
\includegraphics[width=0.4\textwidth]{images/../5e.tools/img/adventure/QftIS/070-map-4.03-palace-of-spires-player.png}
\subsubsection{P1: Gatehouse}
\begin{DndReadAloud}{}
A short, square tunnel constructed from purple-red stone leads to a lush garden enclosed in a courtyard. Behind you, a pair of close-meshed silver gates block the path. Beyond the gates wavers a colorful haze that occasionally provides a glimpse of the lake in the garden you left.

\end{DndReadAloud}

Creatures who solve the sundial puzzle in the maze (area G21) arrive here. The silver gates can't be opened or damaged. If a creature slips through the bars or teleports past them, they can't travel beyond the haze, which is the boundary of the demiplane in which the palace resides.

\subsubsection{P2: Walled Garden}
\begin{DndReadAloud}{}
Tall porphyry walls enclose a beautiful garden, and a glassy reflecting pool dances with the twilight colors of the sky above. Luminous butterflies flit above vivid flower beds that stand to the left and right of the pool.

Just outside the gatehouse stands a sundial identical to the one in the hedge maze. Paths of colorful gravel lead throughout the garden and to the palace proper.

\end{DndReadAloud}

The outer wall is 20 feet high, and the flowers in the garden remain perpetually in full bloom. The reflecting pool is 1 foot deep.

\subparagraph{Sundial}
The sundial here functions identically to the one in the hedge maze (area G21), but the leaves required to operate it don't grow within the demiplane. If the characters didn't bring extra leaves with them when they entered, they must find the reusable keys carried by the current occupants of the palace (see areas P6 and P7) to leave it. Activating this sundial returns the characters to the center of the maze.

\subsubsection{P3: Corner Bowers}
\begin{DndReadAloud}{}
Sweet, gentle music without a source fills this bower, which consists of a small, circular alcove with a padded bench and a canopy of flowering vines.

\end{DndReadAloud}

These bowers are compelling spots to rest. The disembodied music soothes worry and weariness. A drowsy feeling washes over a creature that enters, though it is under no compulsion to sleep.

\subparagraph{Refreshing Rest}
A creature that finishes a short rest in the bower instead gains the benefits of finishing a long rest. The creature can't benefit from this property again for 24 hours.

\subsubsection{P4: Entrance Porch}
\begin{DndReadAloud}{}
A broad flight of porphyry steps ascends to a massive arched doorway. The arch over the double door protrudes from the wall and is supported by two stone columns carved in the likeness of two women, a human and an elf. Black letters across the arch read, "Enter in Peace."

The doors are ten feet wide, thirty feet high, and made of silver. They have no handles and are decorated with etched geometric patterns.

\end{DndReadAloud}

The two column statues are depictions of Caerwyn and Porphura. The statues are living guardians (use the \textbf{earth elemental} stat block). They animate and attack anyone who tries to force the doors open by any means. The columns stop attacking a creature that ceases its attempt and lays down any weapons it's wielding. If all creatures do so, the guardians return to their original positions. If the guardians are destroyed, they reform at the next dawn.

\subparagraph{Doors}
The words above the door appear written in the preferred language of whoever reads them. The doors open for a creature only if that creature approaches the doors with no weapons bared. A character can force the doors open with a successful DC 25 Strength (Athletics) check or with a \textit{Knock} spell, but doing so prompts the guardians to attack.

\subsubsection{P5: Entrance Hall}
\begin{DndReadAloud}{}
This airy chamber's floor, walls, and lofty ceiling are made of gleaming white marble. Two doors lead from the room to the north and south. The hall's furnishings include two small, glass-topped tables, each with five chairs arranged around it. Each table holds a crystal decanter, five crystal glasses, and a bowl of fruit.

At the far end of the room, a large, semicircular atrium contains a splashing fountain. The atrium's soaring ceiling is transparent, letting the sky's colorful twilight glow fill the hall. An opaque rectangular shape darkens the center of the ceiling.

The fountain's central water jet climbs high above its pool. A metal staircase curves around the atrium's back wall to a landing forty feet above.

\end{DndReadAloud}

Natural light fills this vast entry hall. The atrium's ceiling rises to 80 feet, while the fountain's water jet peaks at 60 feet. The opaque shape above the atrium is the sarcophagus in the tomb (area P12).

\subparagraph{Disturbance}
Sounds of combat or other loud, alarming noises in this area draw the attention of Juliana and Orlando in area P6 as well as Argus and Hamish, a pair of former adventurers in area P7. The latter two arm themselves and come quickly to investigate. Argus and Hamish ignore any ruckus that doesn't sound dangerous, attributing it to Juliana and Orlando and leaving the young couple their privacy.

\subparagraph{Treasure}
The decanters hold sweet wine, and the fruit is fresh and scrumptious. The two decanters are worth 50 gp each, and the ten glasses are worth 10 gp each. As long as they remain in the palace, the decanters and bowls refill every 24 hours.

\subsubsection{P6: Juliana and Orlando's Room}
\includegraphics[width=0.4\textwidth]{images/../5e.tools/img/adventure/QftIS/071-04-012.juliana-and-orlando.png}
\begin{DndReadAloud}{}
The floor of this spacious bedroom is decorated with mosaics depicting the four seasons. Covering one wall is a detailed tapestry showing two female mages, a human and an elf, picnicking in a woodland glade.

The furniture consists of a large four-poster bed, two chests, two worktables, six comfortable-looking chairs, and two couches. One of the couches hangs from the ceiling on delicate chains, and a young human man and woman relax on it as it gently sways. Slender golden perches dot the walls, and songbirds flit between them as they sing.

\end{DndReadAloud}

This bedroom once belonged to Caerwyn and Porphura, but the denizens of the garden insisted that Juliana and Orlando inherit it. A secret doorway in the northwest corner opens onto a spiral staircase leading to area P10.

Juliana and Orlando (neutral good, human nobles with Armor Class 11 and no armor or weapons) recline together here on the suspended couch, relaxing and chatting. They're dressed in bright lounging robes with slippers. The pair are surprised but not perturbed by the characters' arrival if it's peaceful. They are indifferent to the characters, and to their perception, they have been in the Eternal Garden for only two days, not the two years that have already passed on the island of Sybarate.

Juliana and Orlando gladly relate their tragic tale of forbidden love with sincerity—if a touch melodramatically—and lament that their families can't see past the old feud. Eventually, it becomes clear the paramours have no intention of going home. A character who listens to the couple's story and succeeds on a DC 14 Intelligence (Arcana) check realizes they have fallen under the enchantment of the Fountain All Heal (area G9). Characters who visited the fountain earlier make this check with advantage.

\subparagraph{Breaking the Enchantment}
Juliana and Orlando are under the enchantment of the Fountain All Heal and regard the palace as their new home. Until they're freed from the domain's allure—such as by a \textit{Wish} spell granted in the tomb (see area P12)—they won't entertain any notion of leaving. If the effect on them is broken, the couple entertains reasonable arguments to return home and speak to their worried families.

\subparagraph{Disturbance}
Sounds of combat or other loud, alarming noises draw the attention of Argus and Hamish in area P7, who arm themselves and come quickly to investigate. The two ignore any noise that doesn't sound dangerous, preferring to give Juliana and Orlando their privacy.

\subparagraph{Key Leaves}
Juliana and Orlando each carry a set of gilt silver keys to the sundials in areas G21 and P2. The keys are metal replicas of the required leaves in lifelike detail. Unlike normal leaves, these keys aren't consumed by the sundial and travel with the user to the destination.

\subparagraph{Treasure}
The 20-foot-square tapestry depicts Caerwyn and Porphura, is worth 1,000 gp, and weighs 300 pounds.

Julianna and Orlando's possessions are stored in the secret closet behind the tapestry and include two sets of \textit{fine clothes}, two Rapier, two Breastplate, a purse containing 30 gp and three diamonds worth 300 gp each, and a gold signet ring worth 50 gp bearing Orlando's family crest.

\subsubsection{P7: Argus and Hamish's Room}
\begin{DndReadAloud}{}
This bedroom's floor is a tile mosaic showing a night sky and a golden wheel with astrological symbols. The ceiling is similarly decorated with glittering gold stars.

The walls are whitewashed, and a tapestry covering one wall depicts two female mages, an elf and a human, holding hands. Their free hands are raised, producing magic that causes a magnificent garden to spring up around them as animals and fey creatures watch in awe.

The furniture in the room includes a large bed, two chests, six comfortable-looking chairs, two couches, and a wooden table. A suit of oiled leather armor, a round shield, and a well-worn spear adorn a mannequin in the corner.

Two middle-aged human men are present, one reclining on a couch with a book, the other sitting in a chair, whittling a piece of wood with a dagger.

\end{DndReadAloud}

This well-maintained room includes a secret doorway in the southwest corner that opens onto a spiral staircase leading to area P10.

Argus Velon (lawful good, human \textbf{mage}) and Hamish Parth (neutral good, human \textbf{gladiator}) are romantic partners and mercenaries. Sent by Governor Folcarae to rescue the young couple shortly after they absconded, Argus and Hamish succumbed to the Fountain All Heal and resolved to stay. When first encountered here, Argus is on a couch reading, and Hamish is whittling. Both wear silk lounging robes and slippers, with no other weapons or equipment. The couple are initially indifferent to newcomers but become friendly if approached peacefully.

The men tease each other fondly during conversation. They're grateful their dangerous vocation brought them together and eventually here. Like the other garden residents, Argus and Hamish believe Juliana and Orlando are the reincarnated creators of the garden—and they see their own love reflected in the couple. They're fiercely protective of the pair.

\subparagraph{Breaking the Enchantment}
If the characters break the fountain's enchantment on Argus and Hamish, the two are grateful but ultimately resolve to remain in the garden together of their own free will.

\subparagraph{Key Leaves}
Argus and Hamish each carry a set of gilt silver keys to the sundials in areas G21 and P2. The keys are metal replicas of the required leaves in lifelike detail. Unlike normal leaves, these keys aren't consumed and travel with the user to the destination. The couple also keeps sets of spare key leaves in a secret closet (see the "Treasure" section).

\subparagraph{Treasure}
The 20-foot-square tapestry depicts Caerwyn and Porphura, is worth 1,000 gp, and weighs 300 pounds.

A secret closet in the southeastern corner of the room holds Argus' \textit{spellbook}, which contains all his prepared spells; a \textit{Wand of Magic Detection}; an ivory inlaid box (150 gp) holding five sets of spare key leaves (see the "Key Leaves" section above); two Potion of Greater Healing; and a bag with 10 pp, 20 gp, 15 sp, and a flawed emerald worth 500 gp.

\subsubsection{P8: Gallery}
\begin{DndReadAloud}{}
Curving metal stairs ascend to an art gallery that constitutes the palace's second floor. A low stone rail guards against falling over the edge into the entrance hall below.

Two rows of three marble statues stand in the gallery. The northern row depicts an elf woman in different poses and wearing either an elegant gown or druidic garb. The southern row shows a human woman in wizardly robes or smart finery. Two more statues occupy a pair of alcoves to the west.

On the wall directly ahead hangs a large painting showing the same couple dancing on the shores of a lake.

\end{DndReadAloud}

The art in this room all depicts Caerwyn and Porphura. The pieces are affixed to the building and can't be moved from their positions or destroyed.

\subsubsection{P9: Roof Gardens}
\begin{DndReadAloud}{}
Crenelated parapets border the palace's open, flat roof. Four beds of subtly colored flowers are divided by gravel paths. In the center of the beds where the paths cross is a small table and four chairs.

\end{DndReadAloud}

The roof gardens are lovely spots to sit and relax while enjoying an excellent view of the walled garden below (area P2).

\subsubsection{P10: Crystal Bridges}
\begin{DndReadAloud}{}
A delicate span of glass arches from the tower door to another spire twenty feet away. Nothing else supports the crystal bridge.

\end{DndReadAloud}

The crystal bridges are perfectly sturdy despite their appearance and offer a safe but bracing walk to the upper rooftop garden.

\subsubsection{P11: Upper Roof Garden}
\begin{DndReadAloud}{}
This roof garden is bound on three sides by a low, crenelated parapet and holds two long flower beds. The north bed holds spotless white roses, while the south bed contains roses of a deep red hue.

Between the flower beds, a crushed stone path leads to a flight of steps that rises to a round stone room surmounted by a transparent crystal dome.

\end{DndReadAloud}

This garden leads the way to the tomb of Caerwyn and Porphura (area P12), offering a moment of serene contemplation and fresh air.

\subsubsection{P12: Tomb of Caerwyn and Porphura}
\includegraphics[width=0.4\textwidth]{images/../5e.tools/img/adventure/QftIS/072-04-013.sarcophagus.png}
\begin{DndReadAloud}{}
A great crystal dome surmounts the circular wall of this room, and its floor is similarly transparent. Light from the twilight sky shines through the dome and into the entry hall far below.

A three-foot-high, white-marble sarcophagus rests in the center of the floor. Surmounting the slab are the carved figures of two women, a human and an elf, in repose. An engraving inlaid with gold glitters on the side of the slab at the figures' feet.

\end{DndReadAloud}

As with other crystalline structures in the palace, the floor of this tomb is perfectly safe thanks to the indestructible nature of the demiplane's materials.

\subparagraph{Final Rest}
The sarcophagus holds the perfectly preserved remains of Caerwyn and Porphura. The carved figures on the lid are exquisitely lifelike in their rendering. As an action, a creature can open the sarcophagus with a successful DC 25 Strength (Athletics) check.

\subparagraph{Inscription}
The gold lettering on the foot of the tomb is written in Common and reads as follows:

\begin{DndReadAloud}{}
Oh stranger, whom destiny hath led nigh, Leave us interred in peace to lie. If thou wouldst deign to grant our request, Then fate will move at thy behest. This gift from the contented dead, We swear to grant by our heart and head.

\end{DndReadAloud}

\subparagraph{Treasure}
The true treasure of the tomb is indicated by the inscription. If a creature visits the tomb but leaves its sarcophagus undisturbed, the magic of the palace grants the creature one casting of the \textit{Wish} spell, fulfilled by the spirits of Caerwyn and Porphura on the creature's behalf. As the creature departs the tomb, a pair of joyful voices whisper to it in unison. The voices thank the creature for its respect and inform the creature that its next wish spoken aloud before it leaves the garden will be granted. A creature can gain only one wish from the tomb, no matter how many times it visits.

The tomb errs toward the spirit of a wish rather than the strict letter and doesn't twist the wisher's intent, especially if the wish is made to free a creature from the Fountain All Heal's enchanting draught. Feel free to allow wishes to be partially granted, or to simply inform a player their stated wish isn't fulfilled and they can try again.

If a creature opens the sarcophagus, the palace grants no further wishes, and any it has granted within the last 24 hours are immediately undone.

Inside the sarcophagus lie the perfectly preserved bodies of Caerwyn and Porphura in each other's arms as if only sleeping. Dressed in simple, white, linen robes, they wear matching gold rings (100 gp each). The interred lovers have no wish to return to life, and any magic used to resurrect them or speak to their corpses fails.

\section{Conclusion}
If Juliana and Orlando are freed from the fountain's enchantment, they can easily be persuaded to return home to their families if the characters mention that two years have passed for their families or that their families have settled their feud. The couple takes some time to change into their original clothes and gather their equipment. They bid Argus and Hamish a fond, emotional farewell and depart with the characters.

\subsection{Returning to the Garden}
When the characters leave the Palace of Spires and return to the garden, \textbf{the Gardener} is waiting for them in the center of the maze. How they greet the party depends on how the characters dealt with the tomb.

\subsubsection{Respected the Tomb}
If the characters didn't open the sarcophagus in Caerwyn and Porphura's tomb, the Gardener joyfully greets the characters as if they were old friends. The archfey speaks to each character of the character's individual exploits as if the Gardener witnessed them all, with sincere excitement and admiration. If anyone asks the Gardener about the other garden denizens' belief that Juliana and Orlando are the reincarnations of the gardens' creators, the archfey smiles and requests the characters allow the others to hold onto their beliefs, and to realize the truth—whatever it may be—in their own time. When asked directly if the young lovers are Caerwyn and Porphura reborn, the Gardener is slyly evasive with a twinkle in their eye.

After offering their well wishes, the Gardener hands each character (and Juliana and Orlando if they're present) a white flower petal and informs them the fairy ring (area G1) can take them home if they carry the petal into the circle. The archfey then gives a fond farewell and takes their leave.

When a creature carrying one of the petals steps into the fairy ring, a whirl of multicolored flower petals engulfs the creature and gracefully transports it to the mouth of the Cave of Echoes. The Gardener allows the characters, as well as Juliana and Orlando, to return to the Material Plane with their memories intact. A year has passed on the Material Plane for each day they spent in the Eternal Garden.

\subsubsection{Desecrated the Tomb}
If the characters opened the sarcophagus, a sense of creeping dread settles over them when they return to the maze, and the air has an uncomfortable chill—the garden's endless summer has ended. The hedges of the maze appear withered and dead, and \textbf{the Gardener} is furious. The archfey scolds any characters who participated in the desecration of the tomb, barely resisting resorting to violence.

With a final curse, the Gardener waves a hand, and an icy gale rips through the area, depositing the characters unceremoniously outside the Cave of Echoes. The fey crossing to the Eternal Garden never functions for them again.

Should the characters return to the garden by other means, the domain has become twisted and bitter. Plants are wilted, denizens are vicious, and the climate bears the cold bite of winter. The Gardener no longer eschews violence and seeks to vent their fury on the characters.

\subsection{Completing the Quest}
Once returned to the Material Plane, the characters can deliver the news of their success (or failure) to Governor Folcarae. If the characters lingered in the garden, the governor is astonished to see them, remarking that she thought they'd never return.

\subsubsection{With Juliana and Orlando}
If the characters bring Juliana and Orlando home, the governor is overjoyed, letting her elation and relief crack her usually measured exterior. After a tearful reunion with her daughter and a slightly awkward but heartfelt welcome to Orlando, the governor pays the characters their promised reward in the form of 10 platinum trade bars. Each bar weighs 1 pound and is worth 500 gp.

\subsubsection{Empty-Handed}
If the characters failed to bring the missing lovers home, the governor is crestfallen but recovers her composure quickly. She thanks the characters for their efforts, but it's clear she's merely being polite. After the governor bids them a terse farewell, the characters can remain in Sybar for a few days before their welcome is obviously worn out. Word spreads of their failure, and merchants and innkeepers raise prices, add steep "taxes" to their goods, and eventually refuse service altogether.

\chapter{Pharaoh}
Few deserts are as harsh and inhospitable as the Desert of Desolation, where the sun, the wind, and the land itself seem to despise all living things. The ghost of a long-dead pharaoh, Amun Sa (AH-muhn sah), roams the sun-scorched dunes, appearing to those who trek across the sands. Condemned to wander the desert for eternity, Amun Sa pleads with adventurers to free his cursed soul and save this doomed land—but none have returned from his pyramid. Within the pharaoh's tomb lie deadly traps, sinister threats, and the secret to ending the ancient curse.

"Pharaoh" is designed for four to six 7th-level characters.

\includegraphics[width=0.4\textwidth]{images/../5e.tools/img/adventure/QftIS/073-05-001.ch5-pharaoh.png}
\section{Background}
This section details the history of Amun Sa's curse and ways to involve the characters in the adventure.

\subsection{The Ruin of Bakar}
\includegraphics[width=0.4\textwidth]{images/../5e.tools/img/adventure/QftIS/074-05-002.pharaoh-door.png}
Long ago, the kingdom of Bakar spanned what is now known as the Desert of Desolation. Though the land was already a desert, the mighty River Athis flowed through the region, making the land along its banks flourish. Animals, plants, and the citizens of Bakar prospered from its life-giving waters.

The pharaohs of Bakar believed that if their bodies were prepared with specific funerary rites and then entombed with their wealth, they could ascend to their chosen afterlife. Amun Sa, the last pharaoh of Bakar, took this tradition further than any of his predecessors. He was paranoid of grave robbers, believing that if his tomb were plundered, it would bar his passage to paradise. To safeguard his treasures, Amun Sa commissioned a theft-proof tomb, a solemn pyramid unlike anything the region had ever seen. To fund this elaborate and impenetrable fortress, he declared war on Bakar's neighboring lands and taxed his own people into destitution.

\subsubsection{A Curse on the Land}
Downtrodden and impoverished, the citizens of Bakar revolted against Amun Sa to take back their kingdom. During a riot at the capital, Amun Sa appeared before his people and cursed them, holding high his two royal implements—the \textit{Staff of Ruling} and the star-gem of Mo-Pelar (a \textit{Gem of Seeing}). He invoked the gods and decreed that if his people should ever kill him, the River Athis would cease its life-giving flow.

Suddenly, from the sea of upraised fists sped a single spear, piercing the pharaoh's heart; as Amun Sa died, so too did the river. The riverbed dried up within days, and without its waters, Bakar quickly collapsed. Over the years that followed, Bakar's cities were overtaken by the dunes, and its people fled to neighboring lands.

\subsubsection{A Curse on the Pharaoh}
After his death, Amun Sa's spirit began its journey to the afterlife, but he was stopped by a god of death. The disapproving god chastised the pharaoh's spirit:

"Your legacy was meant to be the benefit you brought to the people under your stewardship, not this stone edifice. As you looked only to your death in life, so shall you look only to your life in death.

"I am bound to fulfill your curse, for you have called it down with power in my name. But I also curse you, Amun Sa, that you shall not voyage into the beyond until some mortal soul does as you so feared, removing your staff and star-gem from your tomb. Only then will your spirit and the land be healed."

Centuries later, Amun Sa's spirit still roams the desert, cursed to wander the dunes until he finds mortals skilled enough to overcome the traps and perils in his tomb and put his soul to rest.

\subsection{Using the Infinite Staircase}
If you're using \textbf{Nafas} as a patron, he summons the characters to the Censer of Dreams (detailed in chapter 1), where he recounts the following wish:

\begin{DndReadAloud}{}
"Condemned for his actions in life, the ghost of a long-dead pharaoh wanders the desert, day and night. Eons of torturous solitude have shown him the error of his ways, and he asks for brave mortals to free his soul and end an ancient curse upon the land. Seek him out and heed his call."

\end{DndReadAloud}

Nafas then teleports the characters to a door along the staircase that opens into a bleak desert. After the adventure, the characters can return to the staircase through the same portal.

\subsection{Adventure Hooks}
If you're not using Nafas as a group patron, consider the following ways to involve the characters in this adventure:


\begin{description}
\item[Exiled.]
The characters committed a crime or earned the ire of a powerful figure and were exiled to the Desert of Desolation as punishment. They're left in the middle of the desert with no food or water.
\item[River Restoration.]
The characters have heard tales of a magical river in this region whose waters foster bountiful crops and heal those who drink from it. Legends state a curse dried up that river eons ago. The characters have journeyed to the desert to find a way to end the curse and restore the river. They might have been hired by the Tears of Athis, a newly formed group of historians and priests dedicated to this goal.
\end{description}

\subsection{Setting the Adventure}
The Desert of Desolation is part of Raurin, a vast desert in southeastern Faerûn in the Forgotten Realms. When the River Athis dried up, it transformed the kingdom of Bakar from a thriving oasis into Raurin's driest region.

However, the Desert of Desolation could be any desert of your choice, or an inhospitable stretch in an otherwise lush region, due to Amun Sa's curse. Consider the following suggestions:


\begin{description}
\item[Dragonlance.]
The adventure takes place in a small, cursed region within the Plains of Dust.
\item[Eberron.]
The long-lost kingdom of Bakar occupied a portion of the Blade Desert that's now uninhabitable.
\item[Greyhawk.]
The Desert of Desolation might be another name for the Bright Desert.
\end{description}

\section{Adventure Summary}
"Pharaoh" begins with the characters wandering the Desert of Desolation. After enduring the desert's terrible monsters and harsh conditions, they happen on the ghost of Amun Sa, who beseeches them to enter his pyramid and retrieve his royal staff and star-gem to end the curse that blights both the land and his soul.

Inside the pyramid, the characters meet the Tears of Athis, a group of historians and priests from the surrounding regions who hope to restore the flow of the holy River Athis. While exploring the pyramid, the characters must navigate a false tomb, a disorienting maze, and halls rife with Undead, then overcome a mummy lord who stands in their way.

Once the characters exit the pyramid with Amun Sa's implements, the pharaoh's soul is freed and the curse on the land is lifted, restoring the River Athis. The descendants of Bakar can then reclaim their lands from the desert and prosper once more.

\subsection{Character Advancement}
If you want to use story-based level advancement, the characters receive experience points for achieving the following milestones rather than defeating monsters:


\begin{description}
\item[Escape the Maze.]
The characters gain 1 level when they escape the Maze of Mists and reach the third floor of the pyramid for the first time.
\item[Recover Amun Sa's Implements.]
The characters gain 1 level when they recover both the \textit{Staff of Ruling} and the Gem of Seeing.
\end{description}

If you follow this method, the characters should reach 9th level by the adventure's conclusion.

\includegraphics[width=0.4\textwidth]{images/../5e.tools/img/adventure/QftIS/075-05-003.pharaoh-original.png}
\begin{DndReadAloud}{About the Original}
\textit{Pharaoh} was originally published by Tracy and Laura Hickman's game company, DayStar West Media, in 1980. They later sold their adventures to TSR, which liked the adventures so much that it not only reprinted \textit{Pharaoh} in 1982 but also hired Tracy Hickman. TSR later published two more adventures in the Desert of Desolation trilogy.

\textit{Pharaoh} was praised for blending masterful dungeon design with a compelling narrative. A year later, the Hickmans built on those same strengths when they introduced the world to the vampire Strahd von Zarovich in \textit{I6: Ravenloft}.

\end{DndReadAloud}

\section{Desert of Desolation}
A centuries-old curse has made the Desert of Desolation uninhabitable. Nothing grows in the desert, and no one has settled its once-green hills since the reign of Amun Sa. The sun scorches its grayish dunes with unnatural intensity, and the few clouds that cross its barren sky bring only choking dust.

The adventure begins with the characters trekking through the desert. Depending on what brings the characters there, they might be searching for answers within the region or trying desperately to escape it. At your discretion, each of the characters might begin this adventure with a \textbf{camel} or a \textbf{draft horse} to ride.

\subsection{Desert Heat}
For 1d6 hours of each day, the desert becomes unbearably hot, rising to temperatures well over 100 degrees Fahrenheit. The \textit{extreme heat} has detrimental effects on creatures not acclimated to such environments (see the \textit{Dungeon Master's Guide}).

In addition to resting, the characters can remove any accumulated levels of exhaustion by drinking the water of Athis, which still flows within Amun Sa's pyramid (detailed later in this adventure).

\subsection{Starting the Adventure}
Read or paraphrase the following text to begin:

\begin{DndReadAloud}{}
The gray mounds of the desert roll toward the horizon. Not one stone shows through the powdery sand, nor does a single insect scurry across its surface. All around you, the burning silence of the sea of sand is unnaturally oppressive and foreboding. Ash and dust billow around your feet and choke your parched throats. There are no landmarks in sight—except the sun, which seems to burn hotter than ever.

\end{DndReadAloud}

The characters wander the desert for two days before encountering the ghost of Amun Sa. At your discretion, Amun Sa's ghost can take longer to appear, exposing the characters to more of the desert's threats, which are detailed in the "Desert Encounters" section below.

When you're ready to progress the story, proceed to the "Ghost of Amun Sa" section.

\subsection{Desert Encounters}
For each day the characters spend in the desert, roll on the Desert of Desolation Encounters table to determine what the characters encounter, rerolling duplicates.

\begin{DndTable}[header=Desert of Desolation Encounters]{cX}

d6 & Encounter\\
1 & Two bulettes with sandstone-textured carapaces burrow up from the dunes and attack. The bulettes focus on mounts and Small creatures in the party (if any), seeking to drag them beneath the sands.\\
2 & The characters spot a lush oasis, but it's a mirage. A character who studies the oasis and succeeds on a DC 20 Intelligence (Investigation) check recognizes the illusion for what it is. A creature that enters the oasis steps into a Quicksand (see the Quicksand).\\
3 & A dust storm rises on the wind. The storm lasts 1 hour, during which time the characters have disadvantage on ranged weapon attack rolls and Wisdom (Perception) checks. Shortly after the dust storm begins, an \textbf{air elemental} and eight dust mephits materialize from the sands and attack.\\
4 & Four gricks lie in ambush just beneath the sand. Adapted to the desert climate, the gricks have advantage on Dexterity (Stealth) checks to hide in desert terrain.\\
5 & A snobbish \textbf{efreeti} royal saunters across the desert. He demands each character pledge a vow of fealty to him or face his wrath.\\
6 & Two wyverns soar above the characters, circling them like vultures. After 1 minute, the wyverns swoop down and attack.\\
\end{DndTable}

\subsection{Ghost of Amun Sa}
Eventually, the characters come across the ghost of Amun Sa. Read or paraphrase the following text:

\begin{DndReadAloud}{}
A solitary man dressed in tattered clothing appears over the crest of a nearby dune. At first, he seems to shimmer in the heat and desert haze, but as he walks closer, it becomes apparent that he's a ghost.

\end{DndReadAloud}

This is the \textbf{ghost} of Amun Sa, former pharaoh of Bakar. He walks in a straight line that intersects the characters' path, shuffling along with a blank gaze like a haunted man.

Amun Sa doesn't notice the characters until he passes within 30 feet of them, at which point he turns to them and relates his sorrowful tale, as if in a trance. Read or paraphrase the following text:

\begin{DndReadAloud}{}
"I am Amun Sa, son of the house of Mo-Pelar. My spirit has walked these sands for time uncounted in search of mighty heroes, to plead for their aid.

"In my time, I was pharaoh of this realm. It was green and beautiful, blessed by the gods with magical waters that gave life to our land and nurtured our crops. But I was a foolish and cruel ruler. I should have protected my subjects; instead, I taxed them into poverty to create a wondrous tomb for my own self-preservation. When they rebelled, I cursed them and this land.

"With my dying breath, the curse was fulfilled, but my soul was cursed as well. I am doomed to walk these sands until a mortal soul retrieves the implements of my power from my tomb: my magic staff and the star-gem of Mo-Pelar. Only then will the curse be lifted, freeing my soul and healing this broken land.

"Many have tried, and none have succeeded. If you undo my terrible mistake by removing these two precious objects from my tomb, you may keep all the treasure you find within. Follow my path to wealth or woe, to destiny or doom."

\end{DndReadAloud}

Amun Sa then turns and walks somberly in the direction of his pyramid. Amun Sa doesn't speak to the characters any further—his curse allows him to share only the tale above. If the characters attempt to harm Amun Sa, he steps into the \textit{Border Ethereal} and continues his solemn march. Amun Sa doesn't harm the party and makes no further attempts to defend himself. If Amun Sa's ghost is slain, he appears to the characters again in 1 hour with no memory of the previous encounter, relating his sad story again.

If the characters follow Amun Sa's ghost, they arrive at his pyramid the next evening. If questioned during the journey, Amun Sa recites his sad speech again, exactly as he did before. In the company of the pharaoh, the remaining trek is safe and uneventful, unless you'd like to add more desert encounters (see the "Desert Encounters" section).

Once Amun Sa arrives at the pyramid, he stops and gestures toward it. He then turns and walks away in search of more travelers.

\section{Pyramid of Amun Sa}
The pyramid of Amun Sa was to be the pharaoh's final resting place. The last remnant of Bakar, it has stood for over a thousand years. When the characters arrive at the pyramid, read or paraphrase the following text:

\begin{DndReadAloud}{}
Amid the vast desert, a great pyramid rises from the sands, barely touched by time. Stairs ascend from the sand to an opening in the pyramid.

Near the base of the stairs rests an empty basin one hundred feet in diameter. South of the basin, a deep stone channel cuts a path into the desert ground.

\end{DndReadAloud}

The pharaohs of Bakar built pyramids not merely as burial grounds but as part of a sacred process to convey a pharaoh's soul to the afterlife. These monuments also entomb the remains of their consorts and favored leaders, preparing pharaohs and their closest companions to dwell among the gods.

A character who examines the basin and succeeds on a DC 15 Intelligence (History) check recognizes it as the Fountain of Athis, a former wellspring of life and the source of the River Athis. In Bakar's heyday, cool water sprang from the fountain, teleported to it from magical silos within the pyramid.

The Fountain of Athis dried up the night of Amun Sa's death. Only by breaking the pharaoh's curse can the basin resume its flow and the land thrive once more.

\subsection{Pyramid Features}
The pyramid of Amun Sa is made of weathered limestone. Unless otherwise noted, the pyramid has the following features:


\begin{description}
\item[Ceilings.]
The ceilings in the pyramid's corridors are 10 feet high; those in rooms are 30 feet high.
\item[Doors and Secret Doors.]
Doors throughout the pyramid are made of wood and banded with bronze; they have no locks. Secret doors blend in with their surroundings but are obvious from the other side. A character who searches a wall for a secret door must succeed on a DC 15 Wisdom (Perception) check to notice it.
\item[Hieroglyphs.]
The ancient people of Bakar wrote in hieroglyphs, symbols representing letters, sounds, and words. A character can translate a passage of hieroglyphs with 10 minutes of study and a successful DC 13 Intelligence (History) check. Passages in different instances of boxed text require separate checks to translate. A creature under the effect of a \textit{Comprehend Languages} spell can read hieroglyphs instantly (no check required).
\item[Lighting.]
The pyramid's rooms are unlit. Area descriptions assume the characters have a light source or other means of seeing in the dark.
\item[Temperature.]
The pyramid's interior is perpetually cool and dry, insulated from the desert heat.
\item[Water of Athis.]
The River Athis still flows in certain areas of the pyramid. The water ends the poisoned condition on any creature that drinks it and removes all levels of exhaustion from the creature. Once a creature benefits from the water, it can't do so again until 24 hours have passed. The water loses its properties if it's bottled or otherwise removed from its source.
\end{description}

\subsubsection{Alterations to Magic}
The following magical restrictions apply:


\begin{description}
\item[Teleportation Ward.]
Unless the text states otherwise, creatures can't teleport into or out of the pyramid or between its floors. Any attempt to do so is wasted.
\item[Warded Walls.]
Magic that would alter the pyramid's stone—such as the \textit{Stone Shape} or \textit{Passwall} spells—has no effect on the pyramid.
\end{description}

\subsection{False Tomb}
The lowest floor of the pyramid is a false tomb designed to discourage would-be tomb robbers. The people of Bakar also once used it as a temple.

The following locations are keyed to map 5.1.

\includegraphics[width=0.4\textwidth]{images/../5e.tools/img/adventure/QftIS/076-map-5.01-false-tomb.png}
\includegraphics[width=0.4\textwidth]{images/../5e.tools/img/adventure/QftIS/077-map-5.01-false-tomb-player.png}
\subsubsection{P1: Tomb Entrance}
The stairs outside the pyramid lead to this outdoor platform. When the characters approach, read:

\begin{DndReadAloud}{}
Two people stand on a platform at the top of the stairs, flanking the pyramid's arched entrance, which is capped by a stone awning. The figures carry curved swords and are dressed in light, airy attire suited to the desert heat. Each also holds a torch.

\end{DndReadAloud}

Two siblings, Atfez and Pachi (neutral good, human bandit captains), stand guard outside the pyramid's entrance, keeping watch for creatures that might disrupt their faction's activities within. Atfez and Pachi are members of the Tears of Athis, a group of historians, priests, and descendants of Bakar who hope to see the River Athis flow once more.

When Atfez and Pachi notice the characters, they're surprised to see other people in the desert, especially at night. Atfez, the more assertive sibling, interrogates the characters about their presence, while Pachi observes and periodically apologizes for her brother's prying. If the characters make it clear they mean no harm, Atfez and Pachi let them enter. Otherwise, they ask the characters to leave, defending themselves as necessary.

\subsubsection{P2: Main Worship Hall}
\begin{DndReadAloud}{}
The air is cool and still in this vast limestone chamber. Three pillars line each side of the room, and corridors lead to the east and west. A twenty-foot-tall statue of Amun Sa stands before the north wall. The room is brightly lit by sconces on the pillars.

An intricate fresco once adorned the walls with artwork and hieroglyphs. Now, much of its plaster lies shattered on the floor. A human woman in priestly robes paces about the room, overseeing five scholars who work to reattach the pieces to the walls.

\end{DndReadAloud}

The people of Bakar were meant to pay their respects to the deceased pharaoh here. In the centuries after the fall of Bakar, the walls were destroyed by disgruntled citizens and roaming vandals.

\includegraphics[width=0.4\textwidth]{images/../5e.tools/img/adventure/QftIS/078-05-004.iaseda.png}
The robed woman is Iaseda (lawful good, human \textbf{acolyte}), leader of the Tears of Athis. She is fifty years old and has a calm, determined manner. Around her, five more members of the Tears of Athis (neutral good scouts) use a sticky paste to carefully repair the room's plaster walls, which they hope hold clues to restoring the River Athis. If the characters were peaceful toward the guards outside the pyramid, Iaseda is friendly toward them. If the characters pose a threat to the Tears of Athis, Iaseda and her fellow Tears defend themselves and the pyramid fiercely, decrying the characters as defilers.

\subparagraph{What Iaseda Knows}
Iaseda can relay the history of Bakar (detailed in this adventure's background), but she doesn't know about the fate of Amun Sa's spirit or that his ghost wanders the desert. Iaseda can also share information on the following topics:


\begin{description}
\item[Pyramid.]
This pyramid is the tomb of Amun Sa, the last pharaoh of a bygone kingdom known as Bakar that once occupied these lands.
\item[Tears of Athis.]
Iaseda leads of the Tears of Athis, a newly formed group of historians, priests, and descendants of the people of Bakar. The Tears of Athis arrived at the pyramid two days ago. They seek to restore the River Athis. Amun Sa built his tomb at the river's source, so it stands to reason the secret to reviving it might be found inside.
\item[Tidings of a Ghost.]
If the characters recount their experience with the pharaoh's ghost, Iaseda is skeptical. She jokingly suggests the characters have spent too much time in the desert heat.
\item[Wall Translations.]
The Tears of Athis originally brought experts on Bakarian hieroglyphs with them on their expedition. Unfortunately, a purple worm ate those experts. Although the surviving Tears have made good progress reassembling this room's broken plaster, they haven't yet translated any of the hieroglyphs.
\end{description}

\subparagraph{Hieroglyphs}
The hieroglyphs on the walls were carved by Amun Sa's scribes while he lived, at his dictation. See the "Pyramid Features" section for guidance on translating hieroglyphs throughout the tomb. If translated, the east wall's hieroglyphs read:

\begin{DndReadAloud}{}
"I, Amun Sa, set forth a record of myself and of my dealings with this world. Our land is rich with green and fertile fields. The River Athis is the mother of our land, giving life with its waters. Athis nourishes all she touches and gives strength and health to her children. In my youth, my father, the pharaoh, would sit with me beside the spring and tell me stories of the river's wonderful power and the blessings it gives unto this land."

\end{DndReadAloud}

If translated, the south wall's hieroglyphs read:

\begin{DndReadAloud}{}
"In the years that followed, I learned of the passing of the pharaohs and how our tombs could guarantee a path into our chosen afterlife. I watched as my father built a tomb for himself, to guard against bandits who might plunder his tomb, robbing him of his refuge in the beyond. Yet, only a few years after my father's death, his great burial place was desecrated by thieves. I was tortured by the thought of his spirit wandering the planes alone forever."

\end{DndReadAloud}

If translated, the west wall's hieroglyphs read:

\begin{DndReadAloud}{}
"Cloaked in the darkness of night, I went to my father's tomb to witness the truth with my own eyes. My father's riches were plundered, and the hull of his ship, once jewel-encrusted, was bare and scarred with gouges. I knew with certainty he could not have approached his afterlife in such a boat. As my torch went out, so did the light within my soul. Clutching my staff, I swore by the gods that I would not be cheated of my place in the great beyond. I would build a tomb to end all tombs—one that none could ever despoil."

\end{DndReadAloud}

\subparagraph{Secret Door}
The statue of Amun Sa conceals a secret door that leads to area P7.

\subsubsection{P3: West Offering Temple}
\begin{DndReadAloud}{}
A twenty-foot-tall statue of Amun Sa stands before the north wall of this square room, which rises to a domed ceiling. A few tidy cots lie along the floor.

\end{DndReadAloud}

The Tears of Athis camp in this room. They sleep here when they're not restoring the walls in area P2.

\subparagraph{Secret Door}
The statue of Amun Sa conceals a door that leads to the hall between areas P5 and P7.

\subsubsection{P4: East Offering Temple}
\begin{DndReadAloud}{}
A domed ceiling caps this well-lit room. Burning wall sconces illuminate ancient script along the walls. To the north, a twenty-foot-tall statue of Amun Sa, arms outstretched, holds an enormous altar bowl of blazing fire. A stone staircase ascends to the bowl, with hieroglyphs carved into the face of the lowest step.

\end{DndReadAloud}

This room is brightly lit by \textit{Continual Flame} spells cast on wall sconces. When Amun Sa was interred, his most loyal priests dutifully locked themselves inside the pyramid's upper floors. In the brief period between Amun Sa's death and the collapse of Bakar, the people of Bakar brought food and drink to this room. They believed they were bringing offerings for the gods, but in truth, they were sending sustenance to the priests in the pyramid.

\subparagraph{Hieroglyphs}
If translated (see the "Pyramid Features" section), the hieroglyphs read as follows:

\begin{DndReadAloud}{}
"Lay your offerings of food and drink for the gods in the blazing fire before Amun Sa. That which the gods accept will vanish within the flame."

\end{DndReadAloud}

\subparagraph{Offering Bowl}
The bowl's flame gives off neither heat nor smoke. Whenever an object or creature enters the bowl, the flames surge around the object or creature, concealing it completely and teleporting it to area P13 of the pyramid's second floor. To an onlooker, the object or creature appears to be consumed by the flame, vanishing in an instant. An onlooker who succeeds on a DC 14 Intelligence (Arcana) check can tell that the bowl teleported its contents somewhere else. A \textit{Detect Magic} spell reveals an aura of conjuration magic around the bowl.

\subsubsection{P5: West Storage Silo}
\begin{DndReadAloud}{}
The corridor opens into a cylindrical shaft thirty feet in diameter whose slick walls drop into the darkness below. A domed ceiling peaks thirty feet above the floor of the hallway leading to this humid chamber.

\end{DndReadAloud}

The shaft drops 120 feet to the surface of the murky water below, which is 30 feet deep.

\subparagraph{Water of Athis}
The water of Athis (see the "Pyramid Features" section) fills the bottom of the silo.

\subsubsection{P6: East Storage Silo}
The thunderous sound of cascading water fills this chamber, which is otherwise identical to area P5.

The shaft drops 120 feet into turbulent white water. Fifty feet below the entrance, water gushes from an opening in the north wall and falls into the reservoir below, which is 30 feet deep.

Before the curse befell Bakar, this silo used to teleport water from this room into the now-dry fountain outside the pyramid. Now, the water merely pours out here and evaporates at an accelerated rate. If the curse is broken, the silo's teleportation functions again.

\subparagraph{Water of Athis}
The water of Athis (see the "Pyramid Features" section) fills the bottom of the silo.

\subsubsection{P7: Worship Room}
A twenty-foot-tall statue of Amun Sa stands before the north wall of this worship room. There are no hieroglyphs on the walls, and the ceiling is flat.

\subparagraph{Secret Door}
The statue of Amun Sa conceals a secret door that leads to area P8.

\subsubsection{P8: Descending Corridor}
This arched corridor descends at a slight angle toward the false tomb (area P12). A musty smell pervades the hall, and dust covers the floor.

\subsubsection{P9: Great Worship Room}
\begin{DndReadAloud}{}
The walls of this room are lined with seven statues of Amun Sa: three to the north, two to the east, and two to the west. In the center of the room stands an altar. Side-by-side impressions of a left and right hand are carved into the front of the altar. On the ceiling and along the rest of the walls, the name "Amun Sa" has been carved repeatedly in multiple languages.

\end{DndReadAloud}

Most of the statues in this worship room are ordinary decor—except for the center statue on the north wall, which can be opened as a typical secret door, and the southernmost statue near the east wall, which swings aside to reveal a hallway once the altar unlocks it (described below).

\subparagraph{Altar}
If a character puts their hands in the altar impressions and says "Amun Sa," the southernmost statue on the eastern wall swings inward to reveal a corridor to area P10.

\subparagraph{Secret Door}
The center statue of Amun Sa on the north wall conceals a secret door that leads to an empty, 15-foot-long corridor with tiny holes in the walls. A character who inspects the corridor's walls and succeeds on a DC 15 Intelligence (Investigation) check can tell that the corridor once contained a sleeping gas trap, but the trap was sprung ages ago.

\subsubsection{P10: Grand Hallway}
\begin{DndReadAloud}{}
This hallway slopes down at a slight angle. The plaster from the great frescoes that once covered these walls has fallen to rubble. Blade marks cover the walls. The remnants of a broken door lie at the corridor's northern opening. Dust blankets everything, and the air is very dry.

\end{DndReadAloud}

This hall was prepared to fool robbers into thinking the tomb had already been plundered before their arrival. Amun Sa's priests hacked up the decor before the pharaoh was laid to rest.

\subsubsection{P11: Treasure Room}
\begin{DndReadAloud}{}
Beyond the broken door is a huge room scattered with broken pots and chests that have been largely hewn into pieces. Everything is covered with a thick layer of dust, apart from a still-intact chest and two vases.

\end{DndReadAloud}

As in area P10, this scene was staged by the priests of Amun Sa long ago. Three hostile mimics lie in wait here, disguised as an unbroken chest and a couple of intact vases.

\subsubsection{P12: Tomb}
\begin{DndReadAloud}{}
In the center of this cool, dark room rests an ornate stone sarcophagus. Its lid is ajar and broken, and it contains nothing but dust. Hieroglyphs appear to have been hastily chiseled into the north wall.

\end{DndReadAloud}

As in area P10, this scene was staged by the priests of Amun Sa during the pyramid's construction.

\subparagraph{Hieroglyphs}
If translated (see the "Pyramid Features" section), the hieroglyphs read as follows:

\begin{DndReadAloud}{}
"Here lies the tomb of Amun Sa. Know that ye have arrived too late To plunder his ransom for heaven's gate."

\end{DndReadAloud}

\subsection{Maze of Mists}
The second floor of the pyramid is the Maze of Mists: a vast, confounding maze designed to defeat clever intruders who bypass the false tomb by way of the offering bowl in area P4. Amun Sa commissioned an infamous mage named Kordan to create it; Kordan filled the maze's rooms with bloodthirsty creatures and its halls with a disorienting, rust-orange mist.

For each hour the party spends in the Maze of Mists, roll a d6. On a roll of 1, an encounter occurs. To determine what the characters encounter, roll on the Random Maze of Mists Encounters table, rerolling duplicates.

\begin{DndTable}[header=Random Maze of Mists Encounters]{cX}

d4 & Encounter\\
1 & Four phase spiders\\
2 & Two invisible stalkers\\
3 & Eight ghouls\\
4 & Two minotaurs chasing a scared human \textbf{bandit} who belongs to the Hands of Plenty (see area P16)\\
\end{DndTable}

The following locations are keyed to map 5.2.

\includegraphics[width=0.4\textwidth]{images/../5e.tools/img/adventure/QftIS/079-map-5.02-maze-of-mists-map.png}
\includegraphics[width=0.4\textwidth]{images/../5e.tools/img/adventure/QftIS/080-map-5.02-maze-of-mists-map-player.png}
\subsubsection{P13: Welcome Room}
\begin{DndReadAloud}{}
This octagonal room has a domed ceiling and smells faintly of sulfur. The scent comes from four doorways spaced evenly around the room. Each doorway is filled with rusty-orange mist. Between the doorways are four levers set into the walls.

A skeleton lies in the center of the room, loosely gripping a sword that points at one of the doorways.

\end{DndReadAloud}

The skeleton belonged to a Humanoid explorer; its sword points toward the west exit. Characters teleported to this level by the bowl of magical fire in the east offering room (area P4) might be disoriented and not know which direction is north.


\begin{description}
\item[Magnetic Trap.]
Lifting any of the levers along the walls triggers a magnetic trap within the chamber's ceiling. Each creature in the room wearing metal armor is pulled up to the 30-foot-high ceiling, as are any metal objects or weapons that aren't being worn or carried. A creature affected by the magnetic field collides with the ceiling; the creature takes 10 (3d6) bludgeoning damage and has the prone and restrained conditions until the field is deactivated.
\end{description}

To deactivate the magnetic field, all the levers in the room must be reset to their downward position, at which point creatures and objects naturally fall to the floor. Casting \textit{Dispel Magic} on the ceiling deactivates the magnetic field permanently. The levers are lightweight and easily maneuvered. A creature can attempt to move a lever by hitting it with a ranged attack; each lever has Armor Class 10, and on a hit, the attack deals no damage and instead switches the lever's position.

\subsubsection{P14: Mists and Corridors}
Sections of these hallways are filled with an orange mist (marked on map 5.2). The first time a character enters one of these misty areas, read or paraphrase the following text:

\begin{DndReadAloud}{}
As you enter the rust-orange mist, the scent of sulfur floods your nostrils. You feel slightly lightheaded. This mist confuses your faculties and sense of direction.

\end{DndReadAloud}

The mists heavily obscure vision and magically disorient those inside them. Such creatures lose all sense of distance and direction while in the mists. Creatures feeling their way along the walls find their sense of touch numbed, so they're never quite sure if they rounded a corner or continued straight.

Because of these effects, don't show players the shape or size of any of the mist-filled hallways or tell them what direction they're traveling. When the characters emerge in different rooms of the maze or gaps in the mist, describe areas in terms of "left" or "right" rather than cardinal directions.

The contents of the areas between the mists are detailed below.

\subparagraph{P14a}
A leather sack containing 102 gp lies against one wall of this corridor.

\subparagraph{P14b}
A \textit{+1 Warhammer} lies near a heavy plank door in this clear stretch of corridor.

\subparagraph{P14c}
A tiny ring made of fine silver lies on the floor in the center of this corridor. The ring is a cursed \textit{Ring of Feather Falling}. Attuning to the ring extends its curse to the wearer, who suffers the following effect:


\begin{description}
\item[Curse of Ill Temper.]
You are overcome with an urge to disagree with others' actions, ideas, and opinions.
\end{description}

While cursed in this way, the wearer is unwilling to part with the ring. A \textit{Remove Curse} spell or similar magic breaks the curse.

\subparagraph{P14d}
Characters who have a passive Wisdom (Perception) score of 15 or higher detect a cool, fresh breeze from the east.

\subparagraph{P14e}
The words "Knock First" have been carved in Common into a heavy plank door. The bandits in the room beyond (see area P16) wrote this to try to dupe intruders into announcing their presence.

\subparagraph{P14f}
A cool, fresh breeze flows out from under the door, and a character listening at the door hears tumbling torrents of water.

\subparagraph{P14g}
Lying on the floor is the skeleton of a person who died while pulling a cart loaded with three chests stacked atop each other. Characters who have a passive Wisdom (Perception) score of 15 or higher detect the stench of rotten flesh from the north.

If the top chest is opened, it sprays poisoned darts in all directions. Each creature within 10 feet of the chest must make a DC 15 Dexterity saving throw. On a failed save, the creature takes 2 (1d4) piercing damage and 28 (8d6) poison damage and has the poisoned condition for 1 hour. On a successful save, the creature takes half as much damage only.

The bottom two chests each contain 250 gp.

\subparagraph{P14h}
This area reeks of dead flesh. Bodies in various states of decay lie strewn about the floor. In the center lies the corpse of a human. A shimmering axe juts from his chest.

The axe is a \textit{Berserker Axe}. The weapon's name ("Enduval") is inscribed into the blade. The axe stops shimmering once it's removed from the corpse.

\subparagraph{P14i}
Characters who have a passive Wisdom (Perception) score of 15 or higher detect the stench of decaying flesh from the south.

\subparagraph{P14j}
Characters who have a passive Wisdom (Perception) score of 15 or higher detect the stench of decaying flesh from the west.

\subparagraph{P14k}
A leather knapsack lies on the floor next to a door. It contains a \textit{Spell Scroll} of \textit{Fly}, a \textit{Spell Scroll} of \textit{Fireball}, and a pouch containing 30 pp.

\subparagraph{P14l}
A \textit{Ring of Protection} sits on the floor of this clear section of corridor.

\subparagraph{P14m}
A trail of gold pieces on the east side of this corridor leads to the north archway, where it ends. The total value of the coins is 137 gp.

\subparagraph{P14n}
A trail of platinum pieces starts in the center of the corridor and leads to the west archway, where it ends. The total value of the coins is 152 pp.

\subsubsection{P15: Grieving Elves}
\begin{DndReadAloud}{}
Humanoid bones litter the floor of this room, along with the tattered remnants of adventurers' clothing. Five elves dressed in similar attire kneel over the bones, speaking solemn prayers in mourning.

\end{DndReadAloud}

The "elves" are five doppelgangers that devoured a group of adventurers in this room four days ago; like the characters, those adventurers sought to break the pharaoh's curse. The doppelgangers then took on the adventurers' appearances.

The doppelgangers play the role of grieving elves who have just found the bones of their long-lost kin, reading the characters' thoughts to gain their pity and trust. The doppelgangers claim to know the trick to solving the maze: send two people ahead through the mists with a rope for the others to follow. They gladly volunteer to send one of their own as a guide along with one of the characters.

The doppelgangers aim to isolate the characters, slowly assuming their identities until none remain.

\subsubsection{P16: Knock, Knock}
\begin{DndReadAloud}{}
A few cots lie scattered around this room, as well as a few packs of supplies. Several humans dressed in dirty adventuring gear and green bandannas rest here.

\end{DndReadAloud}

Seven human bandits camp in this room. They are led by Cateline and Imbert (chaotic neutral, human bandit captains). The bandits are part of a thieving troupe called the Hands of Plenty. Dirty, sweaty, and tired, the bandits are initially indifferent toward the characters. If the characters knocked on the door before entering, the bandits have their weapons at the ready. Otherwise, they are arguing loudly.

\subparagraph{Parleying with the Bandits}
If the characters don't attack, the bandits are relieved to see them. The Hands of Plenty came to this pyramid days ago, intending to raid it for treasure, but they quickly became lost. Now they just want to find their way out. None of the bandits have encountered the ghost of Amun Sa—they don't even know who he is.

Tired of wandering the pyramid, Cateline and Imbert grudgingly suggest a temporary alliance between the characters and the Hands of Plenty. If the characters agree, the Hands of Plenty accompany the characters and fight alongside them. However, if any of the bandits come across treasure worth 100 gp or more, their greed overwhelms them, and they attack the characters to claim the treasure for themselves.

\subparagraph{Treasure}
Cateline and Imbert each carry 40 gp. The other seven bandits each carry 10 gp.

\subsubsection{P17: Rug Room}
An oversized chest sits atop a beautiful rug in this dome-ceilinged room. The rug beneath the chest is a \textbf{rug of smothering}.

If the chest or the rug is disturbed, the rug slips from beneath the chest and attacks. At the same time, the chest opens to reveal three more rugs of smothering that unfurl themselves and join the fight. Apart from these rugs, the chest is empty.

\subsubsection{P18: Well of Questions}
\begin{DndReadAloud}{}
A sphinx lounges on a wide platform at the far end of this chamber, her wings folded neatly at her side. Her tail flicks back and forth as she stares at you intently.

Water surges from a well in the middle of the room, flowing into a trough on the floor, across the room, and into to an opening below the sphinx's platform.

\end{DndReadAloud}

A \textbf{gynosphinx} named Senenura resides in this room. Initially indifferent to the characters, she asks them what they are doing here, then she makes them an offer. If the characters can answer a riddle, Senenura promises to answer one question they have about the tomb. She warns, however, that if the characters don't answer correctly, she will eat them. The characters need to accept or decline the offer before knowing the riddle.

\subparagraph{Sphinx's Riddle}
Senenura's riddle is as follows:

\begin{DndReadAloud}{}
"Too much, and you wither; too little, and you shiver. I'm slow in the summer, and hasty in winter. What am I?"

\end{DndReadAloud}

\textit{Answer: The answer is the sun.}

\subparagraph{Answered Correctly}
If the characters answer the riddle correctly, Senenura keeps her promise to answer one question they have about the tomb.

The sphinx is knowledgeable about the following topics:


\begin{itemize}
\item Everything in this adventure's background
\item All the creatures and traps within the pyramid
\item All the maze's exits, including those in this room
\item Where Amun Sa's staff and star-gem are kept (areas P67 and P68, respectively)
\item The mummy lord \textbf{Nafik} and his descent into evil
\end{itemize}

\subparagraph{Answered Incorrectly}
If the characters answer the riddle incorrectly, Senenura attacks and fights ferociously to the death.

\subparagraph{Trough}
The circular opening at the base of the sphinx's platform is 5 feet in diameter. Medium and smaller creatures that fall into the trough must succeed on a DC 15 Strength saving throw or be swept into this opening, which empties into area P6.

See diagram 5.1 for a cross-section of this area and its waterways.

\subparagraph{Water of Athis}
The water in this chamber is the water of Athis (see the "Pyramid Features" section).

\subparagraph{Well}
Any character who looks into the well sees a 10-foot-diameter tunnel below the surface that enters the well from the east side. Water flows into the well through this tunnel. Though the volume of water is high, the well's large size makes the current slow and swimmable. This water tunnel angles up to area P29 on the third floor.

\subsubsection{P19: Spear Room}
\begin{DndReadAloud}{}
Little holes honeycomb the right wall of this square room, whose domed ceiling is encircled by a narrow ledge. Steel spears pin four humanoid skeletons to the left wall, and a chest sits on the far end of the room.

\end{DndReadAloud}

This room contains a brutal spear trap.

\subparagraph{Secret Door}
A secret door atop this room's 30-foot-high ledge leads to an observation dome on the third floor of the pyramid (area P53a).

\subparagraph{Spear Trap}
Ninety of the one hundred holes in the south wall contain a hidden spear held in place by magic. When a creature enters this room, and every minute afterward until no living creatures or spears remain, ten of the spears shoot from the wall. Creatures in the room when the spears are released must make a DC 15 Dexterity saving throw. On a failed save, a creature takes 28 (8d6) piercing damage, is pinned to the north wall, and has the grappled condition (escape DC 15). On a successful save, a creature takes half as much damage only.

\subparagraph{Treasure}
The unlocked chest contains 250 pp.

\subsubsection{P20: X-Room}
A giant \textit{X} has been shallowly carved into the floor of this square room. A narrow ledge encircles the base of the room's domed ceiling.

\subparagraph{Secret Door}
A secret door atop this room's 30-foot-high ledge leads to an observation dome on the third floor of the pyramid (area P53b).

\subparagraph{Stone Block Trap}
A character who searches the room for traps and succeeds on a DC 20 Wisdom (Perception) check notices the outline of a 10-foot-square hatch in the ceiling. It is directly above the X, which marks a pressure plate in the floor.

Any weight in excess of 40 pounds placed on the \textit{X} releases a 10-foot-square stone block from the hatch overhead. Creatures below the block when it drops must make a DC 15 Dexterity saving throw. On a successful save, a creature is pushed to the nearest unoccupied space. On a failed save, it takes 55 (10d10) bludgeoning damage and has the prone and restrained conditions as it is pinned under the block. A creature can use its action to make a DC 15 Strength (Athletics) check to free itself or another creature from the block, ending both conditions on a success. The block weighs 2,000 pounds.

\subsubsection{P21: Pendulum Room}
\begin{DndReadAloud}{}
A rope hangs from the ceiling of this room; its lower end is tied to the wall facing the doorway. At the other end is a gigantic, bladed pendulum. A narrow ledge encircles the base of the room's domed ceiling.

\end{DndReadAloud}

An \textbf{invisible stalker} lurks in this room. When two or more creatures enter the room, the invisible stalker severs the pendulum's rope, causing the blade to sweep down at them. It then attempts to shove creatures into the blade's path.

\subparagraph{Pendulum}
If the rope securing the pendulum is cut, the blade swings down toward the door and back again. A creature in the swinging blade's path must succeed on a DC 15 Dexterity saving throw or take 22 (4d10) slashing damage. The blade swings again on initiative count 20 of each round (losing initiative ties), and the damage on subsequent rounds is reduced by 1d10. When the number of damage dice reaches 0, the blade stops.

\subparagraph{Secret Door}
A secret door atop this room's 30-foot-high ledge leads to an observation dome on the third floor of the pyramid (area P53c).

\includegraphics[width=0.4\textwidth]{images/../5e.tools/img/adventure/QftIS/081-05-005.minotaur.png}
\subsubsection{P22: Minotaur Lair}
\begin{DndReadAloud}{}
The walls of this room are splattered with dried blood. There is a ten-foot-diameter hole in the ceiling. Piled beneath it is a large mound of straw.

Three hulking creatures with humanoid bodies, bulls' heads, and piercing red eyes pace around this room, snorting aggressively with axes in hand.

\end{DndReadAloud}

The three minotaurs in this room attack on sight. Summoned by the mage who created the maze, the minotaurs are Fiends instead of Monstrosities. They know nothing of Amun Sa and seek only to spill the blood of trespassers in the maze. The minotaurs are immune to the mists' disorienting effect.

\subparagraph{Ceiling Shaft}
The hole in the ceiling leads to a 30-foot-long vertical shaft covered by a trapdoor in a temple on the pyramid's third floor (area P52).

\subparagraph{Treasure}
Scattered in the hay are 1,800 gp.

\subsubsection{P23: Darkness Room}
Magical darkness fills this room. A \textit{Dispel Magic} spell cast on the darkness dispels it, as does the \textit{Daylight} spell.

A locked chest rests against the west wall. Two black puddings cling to the ceiling above it. The oozes drop from the ceiling to ambush any creature that approach the chest.

\subparagraph{Treasure}
The chest holds three identical bottles: two Potion of Healing and a \textit{Potion of Poison}.

\subsubsection{P24: Robber Press}
This room was designed to grind intruders to a pulp. It appears empty, except for a small leather knapsack that sits at the far end of the room.


\begin{description}
\item[Compacting Wall Trap.]
When a Small or larger creature ventures more than 10 feet into the room, the door slams shut and locks itself, and the walls to the east and west begin to close in. Roll initiative. On initiative count 20 (losing initiative ties), creatures between the walls take 22 (4d10) bludgeoning damage and have the restrained condition as they're slowly crushed by the walls. The walls continue to compress for 3 rounds.
\end{description}

A creature within 5 feet of the door can try to open it by using its action to make a DC 20 Dexterity (Sleight of Hand) check with thieves' tools to bypass the lock or a DC 25 Strength (Athletics) check to force open the door.

A creature in the room can use its action to make a DC 20 Strength (Athletics) check to hold the walls apart until the start of its next turn. On a successful check, creatures take no damage from the walls and don't have the restrained condition during that time.

After 3 rounds, the walls return to their initial positions, ending the restrained condition on creatures still between them. The door then unlocks.


\begin{description}
\item[Illusory Knapsack.]
The knapsack is an illusion to lure creatures into the room. A character who examines the knapsack can determine it's an illusion with a successful DC 15 Intelligence (Investigation) check or by physically interacting with it.
\end{description}

\subsubsection{P25: Suspiciously Empty Room}
This large square room has a domed ceiling and is completely empty. There is nothing of note here.

\subsubsection{P26: Pole Forest}
\begin{DndReadAloud}{}
Like a forest of narrow wooden poles, dozens of long spears protrude from holes in the floor toward this room's domed ceiling, which is encircled by a narrow ledge. The skeleton of a dwarf in chain mail is pinned to the ceiling. The dwarf's pack was ripped open by one of the spears, spilling a pile of jewels on the floor.

\end{DndReadAloud}

Characters can pass through the room by hacking a path through the spears. The spears are easily felled, but doing so is noisy. Roll on the Random Maze of Mist Encounters table (detailed earlier in this section) to determine what creatures arrive to investigate the noise.

\subparagraph{Secret Door}
A secret door atop this room's 30-foot-high ledge leads to an observation dome on the third floor of the pyramid (area P53d).

\subparagraph{Treasure}
Spilled amid the poles are five jacinth gems worth 300 gp each.

\subsubsection{P27: Flameskull Room}
Three flameskulls fly in circles around this room, whose domed ceiling is encircled by a narrow ledge. The flameskulls attack intruders on sight.

\subparagraph{Secret Door}
A secret door atop this room's 30-foot-high ledge leads to an observation dome on the third floor of the pyramid (area P53e).

\subsubsection{P28: Loose Ceiling}
A \textbf{cloaker} camouflaged to match the stonework lies flush against the ceiling of this room, whose domed ceiling is encircled by a narrow ledge. The cloaker attacks the first creature that enters the chamber.

\subparagraph{Secret Door}
A secret door atop this room's 30-foot-high ledge leads to an observation dome on the third floor of the pyramid (see area P53f).

\subsection{Halls of the Upper Priesthood}
The third floor of the pyramid was known as the Halls of the Upper Priesthood. All the Undead on this level were priests of Amun Sa until the former high priest, \textbf{Nafik}, corrupted their souls and cursed them with undeath.

For each hour the party spends on this floor, roll a d6. On a roll of 1, an encounter occurs. To determine what the characters encounter, roll on the Random Upper Priesthood Encounters table, rerolling duplicates.

\begin{DndTable}[header=Random Upper Priesthood Encounters]{cX}

d4 & Encounter\\
1 & Eight giant spiders\\
2 & Three carrion crawlers\\
3 & Six ghasts\\
4 & Three banshees\\
\end{DndTable}

The following locations are keyed to map 5.3.

\includegraphics[width=0.4\textwidth]{images/../5e.tools/img/adventure/QftIS/082-map-5.04-halls-of-the-upper-priesthood-map.png}
\includegraphics[width=0.4\textwidth]{images/../5e.tools/img/adventure/QftIS/083-map-5.04-halls-of-the-upper-priesthood-map-player.png}
\subsubsection{P29: Waterway}
If the party enters this area from the tunnel in the Well of Questions (area P18), read or paraphrase:

\begin{DndReadAloud}{}
You surface in a tumbling, turbulent pool of water within a ten-foot-square shaft that ascends forty feet from the surface of the water. The air is damp, and the walls are covered with slimy green moss.

Dim light streams from an archway near the top of one side of the shaft. A narrow, roaring waterfall pours from the archway into the pool.

\end{DndReadAloud}

If the party enters this area from the hall (areas P30–P31), read or paraphrase the following text:

\begin{DndReadAloud}{}
The corridor ends in an archway followed by a steep drop. The water falls over the edge into a murky pool at the bottom of a deep moss-covered shaft.

\end{DndReadAloud}

Handholds line the walls. Characters can climb this shaft without needing to make any checks.

\subparagraph{Water of Athis}
The waterway is filled with the water of Athis (see the "Pyramid Features" section).

\subsubsection{P30: North Entry}
A compass rose is carved into the floor of this arched corridor. The compass rose is accurate and can help reorient characters exiting the maze on the pyramid's second floor.

\subparagraph{Water of Athis}
An aqueduct fills half the corridor, carrying the water of Athis (see the "Pyramid Features" section) toward area P29.

\subsubsection{P31: Long Hall}
This long hallway rises steadily toward area P32. Bright light shines from the hallway's southern end.

\subparagraph{Water of Athis}
An aqueduct carries the water of Athis (see the "Pyramid Features" section) from the stream in the garden hall (area P32) to the north, where it surges high against the wall before following the corridor's sharp turn to the east.

\subsubsection{P32: Garden Hall}
\begin{DndReadAloud}{}
A rapid stream runs through this brilliantly lit indoor garden. The room's domed ceiling glows with warm light, giving life to the lush plants along the stream. Nestled among the ferns, flowers, and palm trees in the garden are two enormous bronze bowls filled with leafy pomegranates, pears, figs, and dates.

\end{DndReadAloud}

Sunlight beams down from this room's ceiling.

\subparagraph{Secret Doors}
The foliage hides secret doors: one in the southwest alcove to area P34 and another in the southeast alcove to area P35. The greenery raises the DC to find these doors to 20.

\subparagraph{Fruit Bowls}
If the characters approach or disturb a fruit bowl, the magical fruits in the bowls spread their leaves like wings and take flight in all directions like insects. There are fifty flying fruits in all.

A character more than 20 feet off the ground—such as one climbing a palm tree, flying, or leaping into the air—can use an action to try to snatch a flying fruit from the air, doing so with a successful DC 18 Dexterity (Sleight of Hand) check. Alternatively, a character can shoot down a flying fruit by hitting it with a ranged attack roll; each fruit has Armor Class 18. Once plucked from the air or shot down, a flying fruit becomes inanimate.

A creature that eats a flying fruit has advantage on Strength and Dexterity saving throws for the next 8 hours. The flying fruit tastes like a deliciously sweet fruit of its apparent variety. If a flying fruit is taken out of the pyramid, it loses its magic and rots 1 minute later.

\subparagraph{Water of Athis}
The stream flows with the water of Athis (see the "Pyramid Features" section).

\subsubsection{P33: Dome of Flight}
\begin{DndReadAloud}{}
The walls of this octagonal room rise high to a domed ceiling that glows with warm light. A waterfall tumbles into the room from the mouth of a stone-carved lion's head, emptying into a stream that flows through the north entrance. A door sits high on the south wall above the lion's head, opening onto a narrow ledge that encircles the base of the domed ceiling.

Four palm trees stand in the room, and clusters of pineapples grow from their leafy tops. Two identical granite altars sit on the east and west sides of a pool. Each altar is carved with the impression of a right hand, a left hand, and gold-lined hieroglyphs. Thick foliage shrouds doors to the east and west.

\end{DndReadAloud}

This chamber's ceiling is 60 feet high and radiates sunlight. The door atop the ledge leads to area P55, while an archway concealed behind the waterfall leads to area P44.

See diagram 5.2 for a cross-section of this room.

\subparagraph{Altars}
These identical magic altars were used to lift the pharaoh's heavy sarcophagus to the next level of the tomb. Each bears three hieroglyphs. If translated (see the "Pyramid Features" section; treat the hieroglyphs as a single passage), the hieroglyphs read "float," "reverse," and "negate."

If a creature touches both of the hand-shaped impressions, the impressions begin to glow as the altars activate. From then on, any creature in the room can use a bonus action to say one of the following words to cause its corresponding effect:


\begin{description}
\item["Float.]
" Gravity in this room becomes half as strong as normal. Creatures' jump distances are doubled in this room, and creatures and objects that fall descend at a rate of 60 feet per round and take no damage from the fall. Water still flows down the waterfall and through the stream, but at a slower rate than before. This effect can exist simultaneously with the "reverse" effect.
\item["Reverse.]
" The entire room is affected by a \textit{Reverse Gravity} spell. The waterfall flows upward, exposing the hidden archway to area P44 and forming an upside-down pond in the ceiling. This effect can exist simultaneously with the "float" effect.
\item["Negate.]
" If the "float" or "reverse" effects are active, they end immediately. The impressions on the altars stop glowing, and the altars deactivate.
\end{description}

\subparagraph{Grenade Palms}
The trees are grenade palms, bearing explosive fruits that look like pineapples. The palms evolved this way so they could scatter their seeds far and wide. A character who examines the trees and succeeds on a DC 17 Intelligence (Nature) check recognizes this rare plant.

When a creature comes within 5 feet of a grenade palm, the vibrations of the creature's movements loosen one of the explosive fruits overhead. The fruit lands near the base of the tree and bursts in a blast of woody shrapnel. Each creature within 20 feet of the fruit must make a DC 15 Dexterity saving throw, taking 21 (6d6) piercing damage on a failed save or half as much damage on a successful one.

As an action, character can safely pick a fruit from the tree by gently twisting it off the branches with a successful DC 15 Dexterity (Sleight of Hand) check. The character can hurl the fruit up to 60 feet as an action, causing it to burst on impact.

The fruit can be stowed for later use, but doing so is dangerous. If a creature carrying a fruit is hit by an attack or is otherwise jostled, the fruit explodes.

\subparagraph{Tapping Noise}
A character who stands on the ledge or domed ceiling and succeeds on a DC 15 Wisdom (Perception) check hears a faint tapping noise coming from the dome's east side. This sound is a gnome named Prit (see area P54), rapping on the other side of the wall with a spoon. This section of the wall has Armor Class 17, 25 hit points, and immunity to poison and psychic damage. Destroying the wall creates an opening to area P54.

\subparagraph{Water of Athis}
The water in this room is the water of Athis (see the "Pyramid Features" section).

\subsubsection{P34: Gazelle-Headed Statue}
\begin{DndReadAloud}{}
A nine-foot-tall statue of a gazelle-headed woman stands in this alcove. She raises a finger to her mouth.

\end{DndReadAloud}

This statue depicts a god of knowledge unique to the people of Bakar. A character who inspects the statue and succeeds on a DC 15 Intelligence (Religion) check knows this god prizes secrets.

\subparagraph{Treasure}
If a character tells the statue a secret, a Potion of Necrotic Resistance glitters into existence at the statue's feet. This happens only once.

\subsubsection{P35: Tortoise-Headed Statue}
\begin{DndReadAloud}{}
A nine-foot-tall stone statue of a tortoise-headed man stands in this alcove. He holds a hand up to his ear.

\end{DndReadAloud}

This statue depicts a god of art unique to the people of Bakar. A character who inspects the statue and succeeds on a DC 15 Intelligence (Religion) check knows this god enjoys music.

\subparagraph{Treasure}
If a character sings a song or plays a tune before the statue, a Potion of Necrotic Resistance glitters into existence at the statue's feet. This happens only once.

\subsubsection{P36: West Chamber}
Four leather sacks sit on the floor of this 20-foot-long hallway. Each of the sacks contains 100 gp.

\subsubsection{P37: East Chamber}
The slain bodies of two wights, defeated by Uma and her companions (see area P47), lie sprawled on the floor of this long chamber. A character who inspects the wights notices marks left by swords and arrows on the corpses.

\subsubsection{P38: West Hall}
Two wraiths haunt this hallway like sentries. A character looking to traverse the hall undetected must succeed on a DC 12 Dexterity (Stealth) check.

\subsubsection{P39: East Hall}
This quiet hall is empty. Streaks of decades-old dried blood mar the walls and floor.

\subsubsection{P40: West Kitchen}
This room is a long-abandoned kitchen, its counters and cutlery coated in a thick layer of dust. In the center of the room, a heavy wooden block has a cleaver stuck into it. The cleaver, which is made of fine silver, functions as a silvered handaxe.

\subsubsection{P41: West Pantry}
This pantry is barren. As offerings to the pyramid declined, the priests of Amun Sa ran out of food.

\includegraphics[width=0.4\textwidth]{images/../5e.tools/img/adventure/QftIS/084-map-5.05-dome-of-flight-map.png}
\includegraphics[width=0.4\textwidth]{images/../5e.tools/img/adventure/QftIS/085-map-5.05-dome-of-flight-map-player.png}
\subsubsection{P42: East Kitchen}
A dwarf skeleton with a cleaver lodged in its chest lies on a wooden table in this dusty, long-abandoned kitchen. All the kitchen's spoons are missing. They were taken by Prit the gnome (see area P54).

\subsubsection{P43: East Pantry}
\begin{DndReadAloud}{}
Piles of plaster—the remnants of long-shattered frescoes—line the walls of this room. No readable pieces are left. Empty flour sacks litter the floor. The room's domed ceiling bears several wide cracks.

\end{DndReadAloud}

This pantry is barren. As offerings to the pyramid declined, the priests of Amun Sa ran out of food.

\subparagraph{Ceiling Passageway}
Characters who look closely at the cracks in the ceiling see a rough-hewn tunnel rising from one of the larger cracks. The tunnel winds its way to area P64. It was dug with a spoon by Prit, the gnome in area P54.

\subsubsection{P44: March of the True Faith}
Hieroglyphs cover the ceiling and floor of this corridor. Chiseled by a scribe to please the pyramid's former high priest, \textbf{Nafik}, the hieroglyphs also tell the story of Nafik's evil designs.

If translated (see the "Pyramid Features" section), the hieroglyphs on the ceiling read as follows:

\begin{DndReadAloud}{}
"Above our thoughts the pharaoh sleeps, In dreamy realms and sky so deep. The high priest worked a wonder great, And sealed him up into his fate. Great Nafik, the priest most high, Studies tomes that he might ply The watery path where all the great Leave death behind and loose their fate."

\end{DndReadAloud}

If translated, the hieroglyphs on the floor read:

\begin{DndReadAloud}{}
"Nafik was high priest of Amun Sa and leader of his rites. He was keeper of the tomes of Terbakar, the greatest library in all lands of the golden age.

"Nafik searched, too, for life eternal, and some say he sought to rob the pharaohs of their right to that life. But I believe he sought only to serve.

"Nafik's search was rewarded, for the books showed him the way of life eternal here. Now great and terrible in his power, he hopes to share this gift with us.

"He cannot die, for he has sequestered his life elsewhere. Nafik is now second only to the gods themselves. He alone has claim on the rule of Bakar."

\end{DndReadAloud}

\subsubsection{P45: North Priesthood Cells}
\begin{DndReadAloud}{}
This long corridor has a wooden door at each end. Ancient, rotting black drapes hang loosely over the entrances to four small rooms on each side of the hall.

\end{DndReadAloud}

The cells contain only dirt and broken cots.

\subsubsection{P46: West Closet}
Two unfinished sarcophagi face each other from opposite ends of this long chamber. Two mummies lurk inside the sarcophagi, lurching forth to attack any creatures that enter the room.

\subsubsection{P47: West Cell of the High Priest}
\begin{DndReadAloud}{}
The din of battle rises from this room. A heavily armored woman armed with a gleaming sword valiantly battles a horde of gaunt, gray creatures with sharp claws and long, ghastly tongues.

\end{DndReadAloud}

A human \textbf{knight} named Uma is fighting seven ghasts in this octagonal room. Two slain ghasts lie at her feet. If the characters don't immediately join her in battle, Uma shouts for their aid when she notices them.

If the characters help Uma fight off the ghasts, she regards them as comrades. Thankful for their efforts, Uma offers to join the characters if they help her exit the tomb.

\subparagraph{What Uma Knows}
In addition to the history of Bakar (detailed in this adventure's background), Uma can impart the following information:


\begin{description}
\item[Amun Sa's Ghost.]
Twelve days ago, Uma traveled to the desert with a group of adventurers on a quest for a magic amulet. To their surprise, Uma and her companions were approached by the spirit of Amun Sa. They heeded his plea and entered this pyramid, looking to retrieve his staff and star-gem.
\item[Encounter with a Mummy.]
Five days ago, Uma and her companions cornered a mummy garbed in tattered high priest's robes near a waterfall not far from here (see area P33). A fierce battle ensued, and Uma's comrades were slain. Uma struck a killing blow on the undead priest—or so she thought. As the mummy crumbled to dust, he claimed his life was "too precious a thing to carry with him," cackling as he disappeared. Uma can show the characters the way to the waterfall.
\item[Uma's Fallen Comrades.]
Uma laid her companions to rest in a hallway of black curtains (area P51). Since then, she has been surviving on rations, but her supplies are running low. She doesn't know how to get to the tomb's next level.
\end{description}

\includegraphics[width=0.4\textwidth]{images/../5e.tools/img/adventure/QftIS/086-05-006.uma.png}
\subsubsection{P48: Priest Catacombs}
\begin{DndReadAloud}{}
A giant block of black stone lined with twelve upright sarcophagi sits in the center of this huge room. The faces that were carved into the lids of the sarcophagi have been gouged out by deep claw marks.

\end{DndReadAloud}

The sarcophagi are actually doors to the inside of the stone block, which is hollow. Five wights lurk within. They rush out to attack the characters as soon as any of the lids are opened.

\subparagraph{Treasure}
An unlocked chest inside the stone block contains 800 gp.

\subsubsection{P49: East Closet}
Two sarcophagi face each other from either end of this room. Their carved faces have ominous expressions. An inanimate skeleton falls out of the northern sarcophagus when it's opened.

\subparagraph{Secret Door}
Inside the southern sarcophagus is a secret door that leads to a dead end.

\subsubsection{P50: East Cell of the High Priest}
Three wraiths float within this room. When the characters enter, the wraiths are hissing at a holy-looking mace on the floor. The wraiths attack the characters on sight.

\subparagraph{Treasure}
The mace is a \textit{Mace of Disruption} that belonged to an honorable warrior who braved Amun Sa's tomb years ago and perished. The weapon's name ("Bar-ethel") is inscribed in its hilt.

\subsubsection{P51: South Priesthood Cells}
This area's description is identical to that of area P45. However, three corpses have been laid behind the curtain in the northeastern-most cell. A moth-eaten blanket reverently covers each body, and a single, unlit candle rests before them.

The bodies are former adventurers and friends of Uma (see area P47), who dragged her comrades' bodies here to give them a proper burial and hide them from the Undead on this floor.

\subsubsection{P52: Prayer Temple of the Priesthood}
\begin{DndReadAloud}{}
Faded and moth-eaten prayer rugs are neatly placed about this room. In the center of the south wall is a thirty-foot-tall statue of Amun Sa with a giant gemstone glimmering from his forehead. A compass rose is carved into the floor before the statue.

\end{DndReadAloud}

The compass rose is accurate and can help reorient characters who lost track of the cardinal directions in the maze on the second floor.

\subparagraph{Gemstone Trap}
If the gemstone is removed from the statue's forehead, the statue makes a horrendous noise like a foghorn until the gem is put back or the statue is destroyed. The statue has Armor Class 17, 50 hit points, and immunity to poison and psychic damage. The noise attracts nearby creatures. Roll on the Random Upper Priesthood Encounters table (detailed earlier in this section) to determine what creatures arrive to investigate the noise.

\subparagraph{Trapdoor}
A hidden, 10-foot-wide trapdoor sits between the compass rose and the statue. A character who examines the floor and succeeds on a DC 15 Wisdom (Perception) check spots the trapdoor. A creature that steps onto the trapdoor falls 40 feet into the haystack in area P22, taking no damage from the fall.

\includegraphics[width=0.4\textwidth]{images/../5e.tools/img/adventure/QftIS/087-map-5.06-observation-dome-map.png}
\includegraphics[width=0.4\textwidth]{images/../5e.tools/img/adventure/QftIS/088-map-5.06-observation-dome-map-player.png}
\subparagraph{Treasure}
The "gemstone" in the statue's forehead is made of glass and worth only 1 gp. Characters who have proficiency with jeweler's tools recognize this, as does anyone who holds it and succeeds on a DC 13 Intelligence (Investigation) check.

\subsubsection{P53: Observation Domes}
Dark, dusty corridors lead to these domes, whose doors open onto the ledges over various rooms in the Maze of Mists. Priests used them to pass idle hours watching the deaths of grave robbers below.

See diagram 5.3 for a cross-section of a typical observation dome and its view over the area below.

\subparagraph{P53a}
This dome opens over room P19.

\subparagraph{P53b}
This dome opens over room P20.

\subparagraph{P53c}
This dome opens over room P21.

\subparagraph{P53d}
This dome opens over room P26.

\subparagraph{P53e}
This dome opens over room P27.

\subparagraph{P53f}
This dome opens over room P28.

\subsection{Gauntlet}
This floor of the pyramid—the Gauntlet—served as a last line of defense between intruders and the pharaoh's tomb. It also served as the quarters of the pyramid's former high priest, \textbf{Nafik}.

The following locations are keyed to map 5.4.

\includegraphics[width=0.4\textwidth]{images/../5e.tools/img/adventure/QftIS/089-map-5.07-gauntlet.png}
\includegraphics[width=0.4\textwidth]{images/../5e.tools/img/adventure/QftIS/090-map-5.07-gauntlet-player.png}
\subsubsection{P54: Prit's Tunnel}
If the party discovers this tunnel by traveling from the Heart's Lair (area P64), read or paraphrase:

\begin{DndReadAloud}{}
The walls of this tunnel are rough and narrow, as if dug by a small creature. At the end of the tunnel stands a gnome, humming happily to himself as he chips away at the wall with a spoon.

\end{DndReadAloud}

If the party discovers this tunnel by breaking down the wall in the Dome of Flight (area P33), read:

\begin{DndReadAloud}{}
When the dust clears, a surprised gnome stands in a rough-hewn tunnel before you, holding a dingy spoon.

\end{DndReadAloud}

This squirrelly gnome \textbf{commoner} is named Prit. He snuck into the tomb five years ago to admire its architecture and has been happily tunneling through these ruins ever since. Prit is obsessed with spoons and doesn't know much about anything else. He does know that down the rough-hewn corridor behind him there's "a weird little treasure" (see area P64). Prit doesn't elaborate further, electing to show the characters the spoons he has collected instead.

\subparagraph{Treasure}
Prit carries a collection of one hundred and twelve spoons in his backpack. Many are mundane, but others are historical curios. If sold to the right buyer, the collection could fetch up to 350 gp, but Prit doesn't part with it willingly.

\subsubsection{P55: Entry Corridor}
An unsettling chill pervades this branching corridor. The two doors to the south are identical. They are both made of thick wood and bear gold hieroglyphs.

\subparagraph{Hieroglyphs}
If translated (see the "Pyramid Features" section), the hieroglyphs read as follows:

\begin{DndReadAloud}{}
"Beyond these doors lies the great Nafik, once second only to the pharaoh—now second only to the gods."

\end{DndReadAloud}

\subsubsection{P56: Gauntlet of the True Way}
\begin{DndReadAloud}{}
The well-lit hall before you is vast and majestic, rising at a slight angle toward the south. Wall torches light four flights of stairs that connect wide landings. Water thunders in torrents past both sides of the stairs and underneath these landings. At the top of the stairs is a platform, on which rests a giant bronze forearm with an upraised fist.

An imposing figure stands beside this bronze arm. He is wrapped almost entirely in tattered priestly attire that trails at his feet, but a few gaps between the wrappings reveal his rotting, desiccated skin.

The figure casts his gaze at you and says in a raspy voice, "Foolish beings of living flesh, your journey ends here. Once, I was the pharaoh's high priest—but now, I enjoy immortality that even he could never hope to attain. Revel in this moment: the rushing water of Athis, the smell of burning torches. It is the last scene you will ever witness."

\end{DndReadAloud}

\textbf{Nafik}, the heartless former high priest of Amun Sa, stands at the top of the stairs. As soon as Nafik is done monologuing—or if his speech is interrupted—he attacks.

If the characters try to explain to Nafik that Amun Sa sent them, he responds that the pharaoh is dead. Nafik considers himself the rightful inheritor of the tomb, its contents, and the kingdom of Bakar.

\includegraphics[width=0.4\textwidth]{images/../5e.tools/img/adventure/QftIS/091-05-007.nafik.png}
\subparagraph{Bronze Arm}
When combat begins, the bronze arm atop the stairs animates to defend Nafik. On initiative count 20 of each round (losing initiative ties), Nafik can command the arm to grab or push a creature of his choice that he can see within 10 feet of the arm. The target must make a DC 15 Strength saving throw. On a failed save, the target has the grappled condition (escape DC 15) or is pushed into the rapids below (Nafik's choice). The arm has Armor Class 19, 30 hit points, and immunity to poison and psychic damage.

\subparagraph{Destroying Nafik}
Nafik's heart resides in area P64. To vanquish the high priest for good, the characters must first destroy his heart. If Nafik is reduced to 0 hit points while his heart remains intact, he cackles and crumbles to ash, knowing he will rise again. If his heart was previously destroyed, Nafik's lifeless body slumps to the floor, his face twisted into a permanent scowl.

\subparagraph{Gauntlet}
Each landing in this chamber has an effect that triggers when a creature other than Nafik steps foot on it. Each effect triggers only once.

From north to south, the landings' effects are as follows:


\begin{description}
\item[North Landing.]
Two ochre jellies ooze from small holes in the east and west alcoves of this landing. They are loyal to Nafik and fight to the death.
\item[Center Landing.]
A 30-foot-tall, 50-foot-wide wall of flame erupts from the floor in the middle of this landing, dividing the northern and southern halves of the room. The wall is opaque and lasts for 1 minute or until dispelled by a \textit{Dispel Magic} spell (DC 15). A creature that passes through the wall must make a DC 15 Dexterity saving throw, taking 14 (4d6) fire damage on a failed save or half as much damage on a successful one.
\item[South Landing.]
A number of shadows equal to the number of characters appear from the southernmost alcoves. Each shadow is a wispy duplicate of a character intent on killing its original.
\end{description}

If the characters entered from the south, or if they don't try to ascend the stairs, Nafik can instead trigger one of these effects at the start of his turn (no action required). He can trigger each effect only once.

\subparagraph{Rapids}
Two 10-foot-deep aqueducts surge with water 10 feet below the landings. This water is the water of Athis (see the "Pyramid Features" section).

A creature that falls into an aqueduct is swept to the northern end of the room and slams against the ledge face, taking 7 (2d6) bludgeoning damage. A creature can climb from the water level onto the north landing without needing to make any checks.

Water flows into the room from area P57, entering through a forked, 10-foot-wide channel beneath the frothy rapids. It flows out through sturdy grates in the northern ledge face, eventually emerging from the lion's head spout in area P33.

\subsubsection{P57: Pillar of Athis}
\begin{DndReadAloud}{}
A ten-foot-diameter column of water thunders down from the ceiling of this semicircular room, pouring into an opening in the floor.

\end{DndReadAloud}

If the characters visit this room before destroying \textbf{Nafik} and his heart (see areas P56 and P64), the chamber is covered with thick grime. A deep voice emanates from the water column and utters the following message in Common:

\begin{DndReadAloud}{}
"A terrible evil befouls this place. A wicked priest haunts these halls, his heart preserved in glass by foul magic. Those who end his undeath may return and face the Trial of Truth."

\end{DndReadAloud}

Once Nafik and his heart have been destroyed, the grime in this room disappears to reveal a question mark engraved in the floor before the column.

\subparagraph{Column of Water}
This water column flows with the water of Athis (see the "Pyramid Features" section). A creature that steps into the column is abruptly swept through a duct into the rapids in area P56.

\subparagraph{Trial of Truth}
If a creature asks a question in this room after Nafik and his heart have been destroyed, a deep voice speaks from the column and asks three questions. The questions are as follows:


\begin{itemize}
\item "What is your name?"
\item "What is your quest?"
\item "On whose hallowed ground do you stand?" (The answer is "Amun Sa's.")
\end{itemize}

If a creature answers any of the questions inaccurately or lies, lightning leaps from the column of water toward that creature. The target must make a DC 15 Dexterity saving throw, taking 45 (10d8) lightning damage on a failed save or half as much damage on a successful one. If a question is answered truthfully, the voice proceeds to the next question until all three are answered.

After all three questions have been answered truthfully, read or paraphrase the following text:

\begin{DndReadAloud}{}
"Put your hand in mine," says the column. A glowing impression of a hand appears on the floor before it.

\end{DndReadAloud}

If a character places a hand into the impression, the water column flows backward for the next 10 minutes. Creatures that enter the column are gently ushered up to area P65.

\subsubsection{P58: Reading Room}
\begin{DndReadAloud}{}
The walls, ceiling, and floor of this room are made of marble. They might once have been beautiful, but now they're covered in cobwebs, dirt, and grime. An old, scorched book lies on a marble slab in the center of the room.

\end{DndReadAloud}

This untitled book details how one can become a mummy lord through a complex ritual that involves numerous unspeakable acts and vile ingredients. Some of its pages are burnt and illegible, while others have been torn out.

\textbf{Nafik} tried to transform himself into a mummy lord using this book's ritual, but he lacked the requisite time and ingredients. His botched ritual transformed him into an unholy creature and his most loyal priests into Undead.

\subsubsection{P59: Stagnant Pool}
A pool of stagnant water rests in this octagonal room. Three hostile water weirds lurk in the pool.

\subparagraph{Treasure}
A waterproof bag containing 110 pp lies at the bottom of the pool, obscured by the water's grimy surface. Once Nafik and his heart have been destroyed, the grime in this room disappears.

\subsubsection{P60: Sitting Room}
This octagonal room holds a bench, a chair, and a rug. The room otherwise contains nothing of value.

\subsubsection{P61: Dining Room}
Cobwebs cover this grand dining room for Amun Sa's priests. Characters who peruse the room's cabinets notice all the spoons are missing. They were taken by Prit the gnome (see area P54).

\subparagraph{Treasure}
The ornate plates and silverware in this room are worth a total of 110 gp.

\subsubsection{P62: Storage Room}
\begin{DndReadAloud}{}
The north wall of this room has been broken through, with rubble and dirt piled on the floor. A rough-hewn tunnel twists north from the opening.

\end{DndReadAloud}

Two earth elementals created by Nafik dug this tunnel under his instructions. They eventually hollowed out a cavern to safeguard the high priest's heart (see area P64).

\subsubsection{P63: Burial Chamber}
Six open sarcophagi rest in this octagonal room. This room was supposed to be the burial place of \textbf{Nafik} and the other priests of Amun Sa, but Nafik transformed them all into Undead with his foul ritual. The sarcophagi are empty.

\subsubsection{P64: Heart's Lair}
\begin{DndReadAloud}{}
The walls of this cavern are rough and irregular. Jumbled blocks of stone are strewn about the floor. Two large creatures made of solid earth stand imposingly before an alcove in the cavern's northeast corner.

\end{DndReadAloud}

Nafik keeps his heart in a glass bell jar in the cavern's northeast cave. It's guarded by two earth elementals. The elementals stand in stony silence unless intruders approach the northeastern alcove, at which point they attack.

\subparagraph{Nafik's Heart}
The heart's glass jar is sealed to the floor in the northeast cave and can be smashed easily. The heart inside looks like an ordinary human heart, beating steadily. Nafik's heart has Armor Class 5 and 5 hit points. If the heart is destroyed, Nafik's Rejuvenation trait ceases to function.

\subsection{Tomb of Amun Sa}
The following locations are keyed to map 5.5.

\includegraphics[width=0.4\textwidth]{images/../5e.tools/img/adventure/QftIS/092-map-5.08-tomb-of-amun-sa-map.png}
\includegraphics[width=0.4\textwidth]{images/../5e.tools/img/adventure/QftIS/093-map-5.08-tomb-of-amun-sa-map-player.png}
\subsubsection{P65: The Pharaoh's True Way}
The characters can enter this room only after completing the Trial of Truth (see area P57). When they arrive, read or paraphrase the following text:

\begin{DndReadAloud}{}
The water lifts you up the shaft into the south end of this corridor and splashes around exuberantly. This corridor ends in a set of elegant bronze double doors.

\end{DndReadAloud}

After 10 minutes, the water returns to its normal downward flow. The water enters the shaft from a hole just below the floor of the corridor.

\subparagraph{Water of Athis}
The water in this room is the water of Athis (see the "Pyramid Features" section).

\subsubsection{P66: Treasury of the Pharaohs}
\begin{DndReadAloud}{}
A reed boat spans the length of the western wall of this stone treasury. Several jars rest inside the boat, whose bow features an empty, fist-sized indent.

A mural on the east wall depicts the boat sailing on the clouds with a glittering gemstone in its bow.

\end{DndReadAloud}

The containers in the boat are canopic jars.

\subparagraph{Magic Mural}
This magic mural is a portal. A character who inspects the mural and succeeds on a DC 14 Intelligence (Arcana) check can tell that the mural is magical, and a \textit{Detect Magic} spell reveals an aura of conjuration magic around it.

Passing through the mural teleports a creature to area P67, thousands of feet above the pyramid. A character who touches the painting finds that their hand passes right through it, and they feel cold air on the other side. A character who sticks their head through it finds themself looking out over the desert landscape, as if through a window in the sky.

\subsubsection{P67: Floating Boat}
When the characters first look through the mural in area P66, read or paraphrase the following text:

\begin{DndReadAloud}{}
You peer from a hazy portal high above the desert. Thousands of feet in the sky, the air is thin and cold. Just beneath the portal floats a layer of sparse clouds, and the pyramid is but a tiny speck directly below you.

An exact copy of the boat from the treasury floats thirty feet from you, perfectly still. The large gemstone from the painting is set in its bow, and ten ornate vases rest within it.

\end{DndReadAloud}

This portal and floating boat are located 10,000 feet in the air, directly above the pyramid. The floating boat is 30 feet from the portal. The boat is fixed to its position; no amount of force can move it.

The gemstone in the boat's bow is the star-gem of Mo-Pelar, one of the two treasures the characters require to fulfill Amun Sa's quest. It can be removed from its niche on the floating boat with a gentle twist. Characters might retrieve the gem by one of the following means:


\begin{itemize}
\item Flying or teleporting to the boat
\item Grabbing the gem with a spell such as \textit{Mage Hand}
\item Lassoing the boat with a rope and shimmying across with a successful DC 15 Strength (Athletics) or Dexterity (Acrobatics) check
\end{itemize}

\subparagraph{Falling}
A creature that falls from the boat or the portal plummets 10,000 feet to the ground; it takes 70 (20d6) bludgeoning damage and has the prone condition at the base of the pyramid.

\subparagraph{Treasure}
The star-gem of Mo-Pelar lodged into the floating boat is a \textit{Gem of Seeing}. The ten ornate vases in the boat each contain 200 pp (for a total of 2,000 pp). Each vase weighs 15 pounds.

\subsubsection{P68: The Pharaoh's Tomb}
\begin{DndReadAloud}{}
A grand sarcophagus lies in the center of this room, with a golden staff set across it. In the center of the north wall stands a statue of Amun Sa with a stone replica of a gem held its outstretched right hand and a stone replica of the staff held across its body with its left hand. Hieroglyphs adorn the walls.

\end{DndReadAloud}

This is the true tomb of Amun Sa. Inside the ornate sarcophagus lies the pharaoh's mummified body, dressed in time-worn ceremonial clothing.

\subparagraph{Hieroglyphs}
If translated (see the "Pyramid Features" section), the hieroglyphs read as follows:

\begin{DndReadAloud}{}
"A passage is always provided between the tomb of the king and his likeness, whereby his spirit may pass into his ordained statue and live within the stone we worship in the outer world."

\end{DndReadAloud}

\subparagraph{Treasure}
The staff lying atop the sarcophagus is the \textit{Staff of Ruling}, one of the two treasures needed to fulfill the pharaoh's quest.

\subparagraph{One-Way Teleport}
After retrieving both of the pharaoh's implements, the characters are teleported to area P2. If the characters are on good terms with the Tears of Athis, Iaseda commends them on their successful quest. The Tears of Athis then rush from the pyramid to see if the curse has been lifted.

\includegraphics[width=0.4\textwidth]{images/../5e.tools/img/adventure/QftIS/094-05-008.treasures-of-amun-sa.png}
\section{Conclusion}
If the characters leave the tomb with both the \textit{Staff of Ruling} and the star-gem of Mo-Pelar, read or paraphrase the following text:

\begin{DndReadAloud}{}
Silence hangs heavy in the still, dry air. Time seems to have stopped, holding all the world in the balance. The pharaoh's ghostly figure hovers on the horizon.

Distant thunder rolls gently across the far-flung sands. The horizon blurs as a whirlwind begins to spin around the massive pyramid. Swirling wind and sand sing past you, a chorus of a thousand voices: the hopes and cries of a land long dead, and the name of its long-accursed pharaoh—Amun Sa. You hear his earnest voice echo in your mind amid the jarring crescendo, saying:

"You have done what others could not; you have put my soul and the souls of my people to rest. The curse upon this land is ended. Keep my staff, my star-gem, and all the treasure you have found within my tomb—and know that you will forever have my gratitude."

The sandstorm settles, and all is still once more. The pharaoh's ghost is gone. From below, you hear the cool sound of running water. As the dusty air clears, it reveals a large basin now overflowing with spring-clear water that surges along a channel toward the horizon. It will take time to heal this land, but there will be blossoms in the spring, for the river has returned from its exile—and with it comes life.

\end{DndReadAloud}

The characters have also earned the gratitude of Iaseda and the Tears of Athis. As a reward for completing Amun Sa's quest, the characters may keep any magic items or treasure they recovered from the pyramid, though the Tears of Athis ask the characters to leave any items of cultural, historical, or religious significance inside the pyramid.

In the months that follow, the Tears of Athis spread word of party's deeds and work to build a thriving society in the region so the descendants of Bakar may return and flourish alongside the river's restorative waters again.

\chapter{The Lost Caverns of Tsojcanth}
Deep in the Yatil Mountains lie the Lost Caverns of Tsojcanth (SAWJ-kahn), formerly occupied by the legendary archmage Iggwilv the Witch Queen. Though Iggwilv is long gone, her lair is anything but empty. Demons, dragons, and other formidable creatures haunt the perilous caverns and the mountains that conceal them, and the archmage's untold defenses remain intact. Iggwilv is rumored to have amassed a magical hoard of unsurpassed value, a trove of such fame that scores of adventurers have perished in search of it.

Despite the decades since her disappearance, Iggwilv's impact persists in the caverns, where she is believed to have penned, at least in part, the \textit{Demonomicon of Iggwilv}: her infamous treatise on the infinite layers of the Abyss and its inhabitants. Abyssal rifts—the byproducts of past magical experiments—leak demonic energy into the tunnels, and guardians bound from faraway planes await rescue or defend the caverns from thieves and trespassers. Other creatures have simply taken up residence in the cave system, adding to its dangers.

"The Lost Caverns of Tsojcanth" is designed for four to six 9th-level characters.

\includegraphics[width=0.4\textwidth]{images/../5e.tools/img/adventure/QftIS/095-06-001.ch6-tsojcanth.png}
\section{Background}
\includegraphics[width=0.4\textwidth]{images/../5e.tools/img/adventure/QftIS/096-06-002.lost-caverns-door.png}
Iggwilv carved out a kingdom for herself through magical prowess, strength, and wit, bolstered by a horde of demons bound by magic to further her aims. In her quest for absolute power, Iggwilv accumulated countless enemies determined to destroy her. Among these was the demon lord Graz'zt, whom the archmage imprisoned for a time. Eventually, Iggwilv's ambitions proved excessive, and a climactic battle with Graz'zt broke her power. The Witch Queen then vanished. As her evil influence waned, Iggwilv's realm was sundered, and her lair, the Caverns of Tsojcanth, was lost.

This lost lair purportedly lies deep in a series of caverns somewhere in the Yatil Mountains. Clues to the lair's whereabouts and fragments of the treasures within have surfaced over the last few decades, and three nations that border the mountains where it hides now race to claim Iggwilv's hoard—or, if they can't claim it, at least ensure that their rivals don't either.

The margrave of a nation near the mountains has promising leads on the location of the caverns, thanks to a few expeditions in recent years. Most of the agents the margrave sent out never returned; others came back battle worn and empty-handed, save for harrowing tales of the horrors in the caves. The margrave now seeks adventurers made of sterner stuff undertake the quest.

\subsection{Using the Infinite Staircase}
If you're using \textbf{Nafas} as a patron, he summons the characters to the Censer of Dreams (detailed in chapter 1), where he recounts the following wish:

\begin{DndReadAloud}{}
"A desperate nation seeks a trove of immeasurable worth, amassed by an archmage unmatched. Many have died in search of it. Summon your bravery and journey to her lair, returning with lost knowledge, power, and wealth."

\end{DndReadAloud}

Nafas then teleports the characters to a door along the Infinite Staircase that opens into a roadside inn near the Yatil Mountains, where a ruler searching for the Lost Caverns has placed an agent to recruit for the venture. After the adventure, the characters can return to the staircase via the same portal.

\subsection{Adventure Hooks}
If you're not using \textbf{Nafas} as a group patron, consider the following ways to involve the characters in this adventure:


\begin{description}
\item[Lost Light.]
A powerful curse plunges a nearby nation into strife. Rumor claims the light of \textit{Daoud's Wondrous Lanthorn} is the key to breaking the curse. A ruler who believes the lantern is in the Lost Caverns sends an agent to recruit the characters to retrieve it.
\item[Racing Rivals.]
The margrave of a land near the Yatil Mountains seeks to weaken his political rivals without drawing suspicion. Having heard of the party's exploits, he hires the characters as independent agents to find Iggwilv's lair and its potent treasures before the forces of neighboring nations find them.
\end{description}

\subsection{Setting the Adventure}
On the world of Oerth in the Greyhawk campaign setting, the Lost Caverns hide within in the southern Yatil Mountains, nestled between the nations of Ket and Perrenland, with Bissel to the south. The margrave of the March of Bissel has uncovered clues to the Lost Caverns' location.

You can set the adventure in a different world by placing it in any mountain range near a kingdom or settlement—before her disappearance, Iggwilv's influence spanned multiple worlds and planes of existence. Consider the following suggestions:


\begin{description}
\item[Eberron.]
The nations of Aundair, Breland, and Thrane vie to explore the Lost Caverns in the Blackcap Mountains. Iggwilv's experiments there drew power from the depths of Khyber.
\item[Forgotten Realms.]
The Lost Caverns are hidden in the Sunset Mountains, northeast of Proskur along the High Road. The High Overseer of Elturel hires the characters to locate the caverns.
\end{description}

\section{Adventure Summary}
"The Lost Caverns of Tsojcanth" begins with a meeting between the characters and an agent serving the margrave of a nation whose rivals surround the Yatil Mountains, where the Lost Caverns are believed to be hidden. Whether the margrave covets one of Iggwilv's treasures for himself, seeks to sponsor the adventurers, or wants the characters to prevent his rivals from claiming the hoard first, the agent offers what little information they have about the caverns' whereabouts.

While navigating treacherous mountain passes, the characters must avoid rockslides, contend with roaming monsters or the patrols of the rival nations, and perhaps encounter friendly faces as they search for the entrance to the Lost Caverns. Danger-filled caverns await the characters in Iggwilv's former lair. Bloodthirsty monsters and the deadly remnants of the archmage's power stand between the characters and the Witch Queen's hoard, which is guarded by the greatest of Iggwilv's treasures—her daughter, the vampire warrior \textbf{Drelnza}.

\subsection{Character Advancement}
If you want to use story-based level advancement, the characters receive experience points for achieving the following milestones rather than defeating monsters:


\begin{description}
\item[Entering the Greater Caverns.]
The characters gain 1 level when they enter the greater caverns for the first time.
\item[Encountering Drelnza.]
After the characters encounter the vampire \textbf{Drelnza} for the first time, everyone in the party who survives the encounter gains 1 level.
\end{description}

If you follow this method, the characters should reach 11th level by the adventure's conclusion.

\includegraphics[width=0.4\textwidth]{images/../5e.tools/img/adventure/QftIS/097-06-003.the-lost-caverns-of-tsojcanth-original.png}
\begin{DndReadAloud}{About the Original}
Published in 1982, the official \textit{Lost Caverns of Tsojcanth} module was an expanded and revised version of \textit{The Lost Caverns of Tsojconth}, a tournament adventure Gary Gygax created for Winter Con V in 1976. The adventure offers tantalizing details of the story of Iggwilv the Witch Queen—better known as the archmage Tasha, famous for her \textit{Tasha's Hideous Laughter} spell. The adventure featured Drelnza, the vampire daughter of Iggwilv, and debuted the lightning-breathing behir, which appeared on its stunning cover.

\end{DndReadAloud}

\section{Starting the Adventure}
The settlements nearest to the Lost Caverns are several days' travel from the Yatil Mountains, though way stations, hunting lodges, and roadside inns dot the roads up to the mountains. An agent waits to receive the characters at the last of these inns, a miserable joint named the Lazy Bat.

\subsection{The Margrave's Agent}
\includegraphics[width=0.4\textwidth]{images/../5e.tools/img/adventure/QftIS/098-06-004.iggwilvs-horn.png}
However the characters take on the adventure, they eventually arrive at the Lazy Bat, where they meet an agent who serves the Margrave of the March of Bissel. Read or paraphrase to begin the adventure:

\begin{DndReadAloud}{}
An unassuming human sits across from you at a corner table in the Lazy Bat, a dimly lit roadside inn. They're dressed to blend with the surrounding patrons. "I'm glad you came," the agent says in a low voice. "My employer will be most grateful."

Their voice quiets to a whisper. "New information pointing to the whereabouts of the Lost Caverns of Tsojcanth—the former lair of Iggwilv, the so-called Witch Queen—has surfaced."

The agent produces a rolled-up scroll bound in a blue ribbon and sets it on the table. "This is all we know about the caverns. You're needed to strike out into the mountains, locate the caverns, and recover whatever treasure the archmage left behind. It's sure to be dangerous, but the rewards defy imagination."

\end{DndReadAloud}

The agent, Bannik Vorl (lawful neutral, human \textbf{spy}), can offer the following details:


\begin{description}
\item[Bannik's Employer.]
Bannik avoids naming their employer, but if pressed for proof flashes a seal bearing Bissel's crest: a white-and-red shield with a black tower emblem. A character who sees the symbol and succeeds on a DC 13 Intelligence (History) check recognizes the symbol.
\item[Lost Caverns.]
The Lost Caverns were once the lair of an archmage named Iggwilv. The archmage is long gone, but her treasure remains. Most who set out to find the caverns come back empty-handed or are never heard from again.
\item[Travel Arrangements.]
If requested, Bannik can supply each character with a sure-footed \textbf{riding horse} and ten days of rations.
\item[Treasure.]
The characters are free to keep whatever treasure they find, minus a 10 percent finder's fee to be paid to Bannik and one specific item to be delivered to the margrave: \textit{Daoud's Wondrous Lanthorn}. The margrave is most concerned with ensuring his rivals don't get the treasure, but the lantern is of special interest.
\end{description}

\subsubsection{Iggwilv's Verse}
The scroll provided by Bannik contains a map of the Yatil Mountains and a scrap of parchment scribed with the following verse:

\begin{DndReadAloud}{}
The horn of Iggwilv

Pierces the heart—

Look over your shoulders

Before you start.

How many sorrowed,

Fooled again,

Because they didn't

Turn back then?

\end{DndReadAloud}

The map is an accurate—albeit basic—representation of the mountains and the main roads that run through them. It doesn't depict smaller trails, including the path to the Lost Caverns. The verse hints at the location of the caverns (see the "Finding the Caverns" section later in this adventure).

\section{Yatil Mountains}
The Yatil Mountains are a rugged mountain range that separates the nations of Ket and Perrenland. The peaks are rife with danger and impassable by unseasoned travelers. The characters must trek through the mountains to find the lost dungeon.

To the north, a singular, imposing peak with a hooked summit rises above the others: Iggwilv's Horn. Characters might be tempted to head toward it, believing the Lost Caverns must be located near a mountain named for the Witch Queen, but a clue on Bannik's scroll warns them off this course (see the "Iggwilv's Verse" section).

Map 6.1 depicts the Yatil Mountains. Obstacles along the mountain's main roads are rare, but once off the roads it's a different story. Any hex on the map that doesn't contain a road is \textit{difficult terrain} for overland travel.

\includegraphics[width=0.4\textwidth]{images/../5e.tools/img/adventure/QftIS/099-map-6.01-yatil-mountains-map.png}
\includegraphics[width=0.4\textwidth]{images/../5e.tools/img/adventure/QftIS/100-map-6.01-yatil-mountains-map-player.png}
\subsection{Mountain Encounters}
While the characters are exploring or camping in the Yatil Mountains, roll a d20 three times per day of game time, checking for encounters each morning, afternoon, and night. An encounter occurs on a roll of 16 or higher. Alternatively, map 6.1 offers suggested locations where encounters might occur. To determine what the characters find, roll on the Mountain Encounters table, rerolling duplicates.

\begin{DndTable}[header=Mountain Encounters]{cX}

d8 & Encounter\\
1 & Border Patrol\\
2 & Dragon Flight\\
3 & Fire Giant\\
4 & Gnome Vale\\
5 & Hermit\\
6 & Rockslide\\
7 & Troll Cave\\
8 & Wyvern Roost\\
\end{DndTable}

\subsubsection{Border Patrol}
A group of human soldiers from a nation bordering the mountain range travels the road. A captain (use the \textbf{gladiator} stat block) leads two knights and six guards. All the soldiers are lawful neutral. Their orders are to catch brigands and smugglers, drive monsters from the roads, and prevent raids.

The soldiers are initially suspicious of the characters. However, if the characters explain their business in the mountains and succeed on a DC 15 Charisma (Persuasion) check, the soldiers invite the characters to camp with them for the night. On a failed check, the soldiers dismiss the characters and continue their march.

The soldiers spend each evening sharing stories around the campfire—including one harrowing tale about a vampire named Drelnza who once terrorized the neighboring nation of Perrenland. The vampire was never vanquished, and to this day, the soldiers dare to speak her name only in hushed tones. (\textbf{Drelnza} is further detailed in appendix B.) Consider awarding inspiration to characters who share a story in return. If the characters camp with the soldiers, the soldiers keep watch so they don't have to.

\subsubsection{Dragon Flight}
Berythrach, a \textbf{young blue dragon}, recently made her lair in a rugged valley between two low peaks. The dragon has been raiding the surrounding lands in ever-widening sweeps, expanding her territory.

Berythrach soars overhead, scouring the mountains for prey. The characters can make a DC 17 Dexterity (Stealth) group check to evade the dragon's gaze, doing so on a successful check. On a failed check, the dragon notices the characters and begins circling their location, swooping lower with each pass. A vain dragon, Berythrach expects the characters to cower before her grandeur or offer her treasure in exchange for safe passage.

If the characters don't attack Berythrach, she lands on a rocky outcropping and engages them in conversation. If the characters offend the dragon, she falls on them with bare fangs and sharp claws.

\subparagraph{What Berythrach Knows}
Berythrach knows the routes of several border patrols, as well as the locations of a gnome vale, a hermit's hut, and a troll cave (detailed in the encounters that follow).

Berythrach also knows the approximate location of the Lost Caverns of Tsojcanth—or at least how to find them. Great tides of bats emerge from the caverns at dusk and return to them at dawn. If the characters give Berythrach treasure worth 1,000 gp or a rare or better magic item, the dragon shares this fact and points the characters in the right direction. She then advises them to follow the bats, adding a few barbs about how they'd be doomed without her generosity and wisdom.

\subsubsection{Fire Giant}
A \textbf{fire giant} lumbers through the mountains, seeking revenge on a group of trolls that stole a hammer from the giant's forge. The giant is indifferent toward the characters and informs them of the trolls' theft. If the characters hunt down the trolls and return with the stolen hammer (see the "Troll Cave" encounter), the fire giant uses it to craft them a \textit{Flame Tongue} sword of their choice; this process takes 1d4 days.

The giant doesn't know the location of the Lost Caverns but suspects other local inhabitants might.

\subsubsection{Gnome Vale}
A small track branches off from a main road toward a narrow canyon. The canyon widens into a small, secluded vale with meadows dotted with spring-fed ponds and stands of trees. The valley is home to a reclusive community of mountain-dwelling gnomes. The gnomes raise goats and sheep in the valley and mine metals in the surrounding bluffs.

Five guards (lawful neutral, gnome scouts) watch the canyon entrance at all hours. If the characters approach in a peaceful manner, one of the scouts runs back to the vale to fetch their leaders, Laira Na'Gwaylar (lawful good, gnome \textbf{veteran}) and Saleen Glimmershade (neutral, gnome \textbf{druid}).

Laira and Saleen greet the characters and inquire as to their business. Laira is terse and to the point, though she's not rude. Saleen is quiet and listens intently, but when she speaks, her voice holds authority and confidence.

\subparagraph{Stolen Steel}
A group of trolls recently stole a crate of steel ingots from the gnomes. If the characters retrieve the ingots from the trolls (see the "Troll Cave" encounter), the gnomes promise to reward the characters and divulge what they know of the Lost Caverns. On the characters' triumphant return, the gnomes pay the characters half the ingots' worth in gems (250 gp total) and reveal the information in the following section.

\subparagraph{What the Gnomes Know}
Saleen has heard local bats whisper about magical caves that come alive and swallow up intruders. She knows where the caverns are located and can point the characters in the correct direction, but she and the other gnomes avoid the area.

\subsubsection{Hermit}
A nameless, wizened hermit (neutral good, human \textbf{archmage}) lives in a simple hut in a wooded grove not far from one of the roads. He's been here as long as anyone can remember, seeking enlightenment in his solitude. Dirty and disheveled, the hermit dresses in simple, homespun clothing and carries no weapons, though he wears an invisible \textit{Ring of Mind Shielding} to protect his thoughts.

Occasionally the hermit sits on a flat rock near the road, contemplating the empty path or just enjoying the sun on his withered frame. He gratefully accepts donations of food, fuel for his fire, or warm clothes in exchange for information.

\subparagraph{What the Hermit Knows}
The hermit is a calm and cheerful fellow. He gently responds to anyone who happens by, but he doesn't speak unless spoken to. If asked his name, the man simply replies that he "left it behind a long time ago."

The hermit can share the following information with the characters:


\begin{description}
\item[Iggwilv's Horn.]
The tall, hooked mountain that rises above the Yatils is called Iggwilv's Horn. Many think the Lost Caverns are beneath it, but they're wrong.
\item[Location of the Lost Caverns.]
The hermit believes the caverns lie to the south. Long ago, when the hermit was young and Iggwilv controlled the region, he witnessed much activity in that direction.
\item[Planar Nexus.]
Rumor has it the caverns are a nexus of planes, and many monsters now inhabit them. The caverns have more than one level.
\end{description}

\subsubsection{Rockslide}
A faint splintering sound echoes overhead as rocks clatter down the mountainside toward the party. Creatures on that side of the mountain must make a DC 15 Dexterity saving throw. On a failed save, a creature takes 22 (4d10) bludgeoning damage and has the prone condition, pummeled to the ground by the falling rocks. On a successful save, a creature takes half as much damage only.

\subsubsection{Troll Cave}
The characters happen on a shallow cave not far from one of the mountain roads. Inside, three trolls pick at the remains of a giant goat. The trolls are hostile and hungry, warning trespassers to leave or become their next meal. If two trolls are slain, the last of them surrenders.

\subparagraph{What the Trolls Know}
If the characters attempt to parley with the trolls, they learn the trolls are former residents of the Lost Caverns ousted from their den by their clan. If the characters agree to slay the trolls' estranged cohorts (see area L9 of the Lost Caverns, detailed later in this adventure), the trolls divulge the entrance to the fabled dungeon.

\subparagraph{Treasure}
The trolls' accumulated loot is stashed under a pile of rank, uncured bear pelts in their cave. It includes the following: a crate of 50 steel ingots, each worth 10 gp and stamped with a Gnomish insignia; a forge hammer sized for a giant that functions as a warhammer for Medium creatures; and assorted coins worth a total of 250 gp.

\subsubsection{Wyvern Roost}
Two wyverns roost in a crevice atop a 100-foot-tall spire 1 mile from the road. They have a clutch of three eggs in their nest and are fiercely territorial of the surrounding area as a result. The wyverns take turns hunting and watching the eggs. When the characters near the roost, roll a die. On an odd number, only one wyvern is in the nest; on an even number, both wyverns are present.

\subparagraph{Treasure}
The wyverns' eggs are valuable to those who want to raise the hatchlings as mounts or guardians (and who are willing to brave the dangers of training them). Each egg fetches up to 500 gp from an interested buyer.

\section{Lost Caverns}
The Lost Caverns of Tsojcanth hide in the Yatil Mountains south of Iggwilv's Horn. Not much is known about the caverns' namesake, the mysterious Tsojcanth, who was eclipsed by the archmage who succeeded them. Whether Tsojcanth was slain by Iggwilv or left before her arrival, this predecessor's only remaining trace is the caverns' name, while the Witch Queen's influence persists long after her departure.

The caverns are divided into the lesser and greater caverns.

\subsection{Finding the Caverns}
\includegraphics[width=0.4\textwidth]{images/../5e.tools/img/adventure/QftIS/101-06-005.cavern-entrance.png}
The entrance to the Lost Caverns lies in the southern region of the Yatil Mountains (denoted on map 6.1). A small track that was once a wider, more traveled road leads to the entrance, but over the decades it has been obscured by rockslides and neglect.

The scroll given to the characters by Bannik (see the "Iggwilv's Verse" section) hints the characters should turn their backs to Iggwilv's Horn, meaning they should search to the south. From there, they can find the caves using the following methods:


\begin{description}
\item[Ask Around.]
A few of the mountains' inhabitants—including Berythrach, the hermit, and Saleen Glimmershade—either know where the entrance is or have information that points the characters in the right direction. By conversing with these creatures, the characters can narrow their search.
\item[Follow the Bats.]
Thousands of bats fly out of the Lost Caverns at dusk to hunt and return at dawn to roost in the darkness. Characters who observe the bats and succeed on a DC 15 Intelligence (Nature) check realize there must be a massive cave nearby for the colony to roost in. Characters who follow the bats eventually arrive at the entrance to the Lost Caverns. Alternatively, characters who have a means of speaking with Beasts can obtain directions to the caverns from the bats.
\item[Trial and Error.]
By diligently searching the roads and tracks in the Yatil Mountains, the characters can eventually locate the entrance to the caverns. However, this painstaking option could result in additional mountain encounters and a lengthy journey, depending on what path the party takes through the mountains.
\end{description}

The entrance to the caverns is a yawning cleft in the mountainside. Sharp rock formations at the entrance give the cave mouth the impression of fangs waiting to snap together. The opening is 40 feet wide, 20 feet high, and blackened by soot. Past the stone maw, a wide passage with roughly hewn stairs slants down to the hideous, stone-carved faces of the entry caverns (area L1 of the lesser caverns).

\subsection{Lesser Caverns Features}
Unless otherwise noted, the lesser caverns have the following features:


\begin{description}
\item[Ambience.]
The caverns are damp, humid, and full of life. They echo with the sounds of bats, dripping stalactites, and the cavern's monstrous denizens.
\item[Ceilings.]
Ceilings are 15 feet high in corridors and tunnels and 20 feet high in chambers.
\item[Lighting.]
The caverns are unlit. Denizens carry their own light or rely on darkvision to see. Area descriptions assume the characters have a light source or other means of seeing in the dark.
\item[Underground Rivers.]
The rivers that span the lesser caverns have strong currents. A creature that enters a river for the first time on a turn or starts its turn there must succeed on a DC 15 Strength saving throw or be swept 15 feet in the direction indicated by the arrows on the map.
\end{description}

\includegraphics[width=0.4\textwidth]{images/../5e.tools/img/adventure/QftIS/102-map-6.02-lesser-caverns-map.png}
\includegraphics[width=0.4\textwidth]{images/../5e.tools/img/adventure/QftIS/103-map-6.02-lesser-caverns-map-player.png}
\subsection{Lesser Caverns Locations}
The following locations are keyed to map 6.2.

\subsubsection{L1: Entry Caverns}
\begin{DndReadAloud}{}
This wide chamber is ringed with a mixture of rough-hewn and natural passageways leading away from the entrance and into the unknown. Beside each of these tunnels is a grotesque face carved into the stone. Though each face is slightly different—one has flap-like ears, another protruding tusks, and a third drooping wattles—all are strange and doleful.

\end{DndReadAloud}

The carved faces are magical. They are immune to all damage and speak a few programmed messages. When a creature approaches one of the tunnels, the face beside it animates and warns in a deep, dire tone, "Turn back—this is not the way!"

A character who examines one of the faces and succeeds on a DC 15 Wisdom (Perception) check notices the glint of a gemstone in its mouth.

\subparagraph{Speaking to the Faces}
The faces repeat their warning if a spoken to, except in the following two circumstances:


\begin{description}
\item["Say 'Aah!]
'" If a character asks one of the faces to open its mouth, stick out its tongue, or yawn, the face sticks out its tongue and says "aah," revealing a gemstone the size of a hen's egg on its tongue.
\item["Truth.]
" If a character asks one of the faces which way to go or uses the word "truth" while speaking in this area, the faces respond in unison. Five of the faces lie, saying, "My way is the right way." The face marking the southeast tunnel leading to area L10, however, responds, "I watch the best way."
\end{description}

\subparagraph{Treasure}
Each of the faces contains a different gem worth 200 gp: amber, amethyst, blue aquamarine, garnet, peridot, and pink tourmaline.

A creature that tries to retrieve a gem from one of the faces without first asking the face to open its mouth must make a DC 17 Dexterity (Sleight of Hand) check. On a failed check, the creature takes 11 (2d10) bludgeoning damage as the mouth clamps shut. The mouth similarly damages tools or weapons inserted into it without its permission.

\subsubsection{L2: Streaked Cave}
\begin{DndReadAloud}{}
The walls and floor of this cavern are smeared with blood, and the scent of iron hangs thick in the air.

\end{DndReadAloud}

Twenty stirges hang from the ceiling of this cave. Bright light or any noise louder than a whisper disturbs the stirges, causing them to attack.

\subsubsection{L3: Slate Chamber}
\begin{DndReadAloud}{}
Alternating blue and gray bands of slate and shale compose the walls of this cave. Rusted weapons lie scattered on the ground.

A hulking, crudely formed figure made of gray clay charges at you from an alcove in the northeast, growling. The hilt of a pristine sword protrudes from the creature's chest.

\end{DndReadAloud}

The figure is a \textbf{clay golem} created by Iggwilv to destroy intruders. The golem attacks any creatures that enter the chamber. The golem fights to the death and pursues trespassers relentlessly, even chasing them out of the caverns.

\subparagraph{Treasure}
The weapon stuck in the clay golem is a slender-bladed Giant Slayer Shortsword. The sword's name ("Longtooth") is engraved into the blade in Gnomish. Once the golem is destroyed, the sword can easily be freed from its chest.

\subsubsection{L4: Guano-Covered Cave}
\begin{DndReadAloud}{}
The floor of this vaulted cave is covered in pungent piles of bat guano. Large leather hides hang from the stalactites above, flapping in the draft. Sounds of flowing water echo from the far side of the cave.

\end{DndReadAloud}

The ceiling of this soiled cave is 50 feet high. This cave once hosted dozens of bats, but a \textbf{cloaker} that dwells here devoured them. The hungry creature hides among the stalactites, waiting to fall on one of the characters.

\subparagraph{Tunnel Ledge}
A ledge on the north wall of the cave stretches out 30 feet above the floor. A short tunnel leads from the ledge to the river beyond. The water's surface is 40 feet below the tunnel's exit.

\subsubsection{L5: Littered Cave}
\begin{DndReadAloud}{}
This small cave is littered with broken bones. Near the back wall lies a busted skeleton with a horned bovine head splintered with cracks. Two leather sacks brimming with gold rest near the skeleton.

\end{DndReadAloud}

The skeleton is that of a minotaur thief who fell prey to the trap in this cave. The two leather sacks beside it are illusions designed to lure greedy intruders to their doom. A character who inspects the sacks and succeeds on a DC 18 Intelligence (Investigation) check or who makes physical contact with one of them discerns their false nature.

\subparagraph{Crushing Ceiling Trap}
A 20-foot-wide section of the ceiling over the skeleton contains a crushing deadfall. A character who examines the ceiling and succeeds on a DC 17 Wisdom (Perception) check notices bits of decayed flesh stuck to the ceiling, which is flat and made of stone.

Whenever a creature enters the area marked on the map, the trap magically triggers. That section of the ceiling rapidly drops, crushing anything in its path. Creatures in that area must make a DC 17 Dexterity saving throw. On a failed save, a creature takes 55 (10d10) bludgeoning damage and has the prone condition. On a successful save, it takes half as much damage only and is pushed to the nearest unoccupied space outside the trap's area. Objects wedged under the trap are destroyed. The stone then retracts into the ceiling, and the trap resets.

A successful \textit{Dispel Magic} spell (DC 17) suppresses the trap and its illusions for 24 hours.

\subparagraph{Treasure}
A pouch on the minotaur's skeleton holds five diamonds worth 100 gp each. The illusory sacks contain nothing.

\subsubsection{L6: Grotto}
\includegraphics[width=0.4\textwidth]{images/../5e.tools/img/adventure/QftIS/104-06-006.digging-pechs.png}
\begin{DndReadAloud}{}
This tunnel has been worked recently, judging by the stone chips and dust covering the floor. Rough-carved images of cave-dwelling creatures decorate the smooth, stone walls.

The sounds of mining tools echo from the end of the tunnel, where hewn steps rise into the darkness.

\end{DndReadAloud}

Eight pechs are diligently carving a tunnel from this grotto up into the mountain toward an enclosed cavern they can sense. For now, the tunnel leads only to a dead end. If the characters draw attention to themselves with light or noise, the pechs investigate.

The pechs are initially indifferent toward the characters, admonishing them to douse any bright lights they carry, as the lights hurt the pechs' eyes. If the characters comply, the pechs become curious and willing to chat. Otherwise, they resume their work, defending themselves if necessary.

In conversation, the pechs warn the characters about the clay golem in the slate chamber (area L3). Characters can improve the pechs' attitude toward them by offering the pechs at least 500 gp of gems or a unique carving or sculpture, especially one made of stone. Friendly pechs might be willing to help fight the clay golem or use their Communal Spellcasting to restore petrified characters.

\subparagraph{Carvings}
The carvings on the walls depict cave-dwelling creatures such as bats and reptiles, as well as creatures from the Elemental Plane of Earth such as \textbf{dao}, pechs, \textbf{xorn}, and earth elementals. Characters who inspect the carvings and succeed on a DC 16 Intelligence (Arcana) check identify the planar creatures.

\subsubsection{L7: Fungal Cavern}
\begin{DndReadAloud}{}
Bioluminescent fungi, some as tall as six feet, fill this moist cavern. The caps of the mushrooms sprouting from the floor and walls glow with a faint blue light. Pale, dog-sized insects chitter among the fungus.

\end{DndReadAloud}

Three ghost-white cave crickets (use the \textbf{giant frog} stat block) nibble at the glowing mushrooms in this cave, which radiate dim light. The crickets are indifferent toward the characters.

If the characters disturb the crickets, the insects begin to chirp loudly. The trolls in area L9 investigate any loud noises in this area.

\subparagraph{Mushrooms}
The fungi in this cavern are edible, and characters can forage here without having to make an ability check. Doing so doesn't disturb the crickets.

\subsubsection{L8: Slimy Cavern}
\begin{DndReadAloud}{}
This cavern is full of sizable mushrooms, but the air stings with an acrid scent. Slime in sickly yellow, green, and orange shades blankets most of the fungi, as well as the walls and ceiling.

Near the back of the cave is a cranny mostly free of fungus. A black-cloaked figure, motionless, is wedged into the crack.

\end{DndReadAloud}

Originally this was another cultivated fungus farm like area L7, but slime molds are overtaking the cave. The cloaked figure is the corpse of an elf adventurer who died here, trapped by the green slime. His body has been partially calcified by the mineral-rich water dripping from the ceiling.

\subparagraph{Green Slime}
Six patches of green slime (see the \textit{Dungeon Master's Guide}) cling to the ceiling near the fallen adventurer. A character who inspects the slime and succeeds on a DC 15 Intelligence (Nature) check recognizes the danger it poses. The slime falls on creatures that approach the corpse.

\subparagraph{Mushrooms}
The fungi here are mostly edible. Foraging in this cavern requires a successful DC 15 Wisdom (Survival) check to avoid the harmful molds and slime. On a failed check, the mushrooms the character gathers are contaminated by the slime. A creature that eats a contaminated mushroom must succeed on a DC 12 Constitution saving throw or have the poisoned condition for 8 hours.

\subparagraph{Treasure}
The elf corpse carries a silvered rapier and wears \textit{Bracers of Defense}.

\subsubsection{L9: Stinking Cave}
\includegraphics[width=0.4\textwidth]{images/../5e.tools/img/adventure/QftIS/105-06-007.three-trolls.png}
\begin{DndReadAloud}{}
This reeking cave is filled with rotting vegetation, bones, pieces of pale chitin, and less identifiable foulness. Three enormous hides serve as sleeping pallets.

Three hulking creatures with warty green skin and wiry hair mill through the detritus, chewing cracked bones clutched in their clawed hands.

\end{DndReadAloud}

This cave is the lair of three trolls. They regularly gather rotting material from their den and spread it through the nearby fungus caverns (areas L7 and L8) to ensure the mushrooms grow and attract cave crickets for the trolls to feed on. The trolls are ornery, having long gone without the taste of humanoid flesh—their favorite delicacy. The trolls attack anyone they notice, fantasizing aloud in Giant about the meal they're about to have.

The trolls the characters might have met in the mountains (see the "Mountain Encounters" section earlier in this adventure) were ousted by these trolls.

\subparagraph{Treasure}
Gathered under the trolls' beds is their accumulated loot: 120 gp, 900 sp, four gems worth 50 gp each, wooden dentures set with mother-of-pearl teeth (worth 50 gp), a quiver with twenty silvered arrows, and a Potion of Greater Healing.

\subsubsection{L10: Bat Corridor}
\begin{DndReadAloud}{}
The sandy floor of this long, wide corridor is strewn with piles of bat guano. Small fungi sprout from these deposits and the walls.

\end{DndReadAloud}

During the day, the ceiling of this corridor is covered in thousands of bats, and a cacophony of chitters and squeaks echoes from above. At night, the bats flood out of the caverns to hunt in the mountains, returning to this area at dawn.

While the bats are roosting, bright light or any noise louder than a whisper agitates them and sends them into a panic. A number of swarms of bats equal to the number of characters in this area descend from the ceiling and attack the source of the disturbance. The swarms extinguish any nonmagical light sources they come within 5 feet of with a flurry of buffeting wings.

Whenever a creature destroys a swarm, a new swarm immediately forms in an unoccupied space within 15 feet of the creature. The new swarm acts on the destroyed swarm's initiative count.

If the characters extinguish all light sources and cease attacking and making loud noises, the bats settle back into their roosts at the end of the round. They don't pursue creatures beyond the corridor.

\subsubsection{L11: Long Gallery}
\begin{DndReadAloud}{}
The cavern widens into a gallery of considerable height and length. Clusters of small fungi grow throughout the chamber, and odd indentations dot the floor and walls at semiregular intervals.

\end{DndReadAloud}

The indentations are the rock-cyst burrows of four cave morays—sluglike scavengers that slither between the burrow entrances through tunnels in the rock. The cave morays use the \textbf{giant poisonous snake} stat block. They use their reach to bite at prey from the safety of their burrows.

Cave morays inside their burrows have three-quarters cover. A creature within 5 feet of a burrow can pull a cave moray from its burrow by grappling it. If removed from its burrow, a cave moray frantically tries to escape back to its hole.

Loud sounds here, including the din of battle, draw the attention of the fomorians in area L12. The cave morays retreat into their burrows when the fomorians arrive.

\subsubsection{L12: Fomorian Cave}
\begin{DndReadAloud}{}
This cave stinks of sweat, smoke, and spoiled meat. Old, cracked bones and skulls lie strewn on the floor. A large, flat rock supports massive stone goblets and the carcass of a giant toad, picked clean and buzzing with flies. Two piles of filthy hides and skins form beds at the far end of the cave.

Two hideous giants squat near a circle of hefty stones enclosing a crackling fire. Above the roaring flames, a charred giant insect rotates slowly on a spit.

\end{DndReadAloud}

Two fomorians inhabit this cave. The beds, fire pit, and utensils are all sized for Huge creatures.

The fomorians are cruel, violent, and bored. When the giants notice the characters, they abandon their meal and attack. One of the fomorians wields the roasting spit as its weapon, tearing bites off the bug-tipped skewer between swings.

\subparagraph{Treasure}
Stashed beneath the fomorians' beds are two ivory mammoth tusks (each is worth 600 gp and weighs 100 pounds), a \textit{Spell Scroll} of \textit{Levitate}, and a beaten copper bowl with lapis lazuli handles, which is worth 750 gp.

\subsubsection{L13: River Ledge}
\begin{DndReadAloud}{}
A ten-foot-long wooden boat is tied to a ledge near the water. Intricate, flowing knotwork carving adorns the sides of the craft, which tapers to a point at both ends. Two oars lie inside.

\end{DndReadAloud}

This ledge is visible from the lake (area L14) and the shores of areas L10 and L15.

\subparagraph{Treasure}
The vessel is a \textit{Folding Boat} created to traverse the lake and the rivers that branch from it. In addition to its usual properties, the boat has an extra command word that functions only while the boat is within the Lost Caverns. A character who knows this extra command word ("Shrimpkin")—discovered by casting the \textit{Identify} spell on the boat or by examining clues in the boat's intricate designs and succeeding on a DC 16 Intelligence (Arcana) check—can command the boat to row itself in a particular direction at a speed of 30 feet or to hold its position in defiance of the water's swiftest currents.

\subsubsection{L14: Underground Lake}
\begin{DndReadAloud}{}
A river flows into the south side of this high-vaulted cavern, feeding a glassy, ebon-hued lake. Drops of water fall from the ceiling high above, rippling the mirrorlike surface of the water.

Two rivers flow out of the northern side of the lake, while a third flows out to the west. Ledges along the shore lead to dry tunnels.

\end{DndReadAloud}

The lake teems with blind cave fish and crayfish. These critters are a food source for many of the denizens in the lesser caverns, including predators that dwell within the lake itself. The lake slopes sharply from its shores to a depth of 100 feet.

Four chuuls lurk near the lake bed, 90 feet below the surface. They're content to creep along the lightless depths and prey on cave fish unless something else catches their attention. Bright light draws the chuuls to the water's edge, as does any magic in the lake or near its surface, such as the boat from area L13. If drawn to the surface, the chuuls attack.

\subparagraph{L14a: Stone Bridge}
An unadorned stone bridge spans the northwest river.

\subparagraph{L14b: Gargoyle Bridge}
A bridge decorated with reliefs of gargoyles spans the west river flowing from the lake. The sound of roaring water echoes from downstream. A short distance beyond the bridge, the river becomes a waterfall, plummeting 300 feet into an Underdark lake. The lake and the creatures that inhabit it are beyond the scope of this adventure, but you can flesh them out as you see fit.

\subsubsection{L15: Basilisk Den}
\begin{DndReadAloud}{}
Oddly shaped chunks of rock are strewn about this cave, which reeks of ammonia. Scattered among the irregular rock pieces are stone sculptures of various creatures, including bats, subterranean lizards, and a broken bust of a gnome with a horrified expression.

\end{DndReadAloud}

Four basilisks dwell in this cave. Toward the back of the cave is a nest where the basilisks lay eggs. The basilisks are protective of their den, attacking any creatures that enter the cave.

\subparagraph{Treasure}
The head and shoulder of a petrified gnome lie near the basilisk nest at the back of the cave. The broken statue wears a pair of unpetrified goggles, which are \textit{Eyes of Minute Seeing}.

\subsubsection{L16: Rainbow Cavern}
\begin{DndReadAloud}{}
This vaulted cavern displays a rainbow of colors on its walls and floor. Stalactites hang like colorful icicles, and the stalagmites feature bright streaks and swirls. Mineral deposits form frozen curtains, cascades, and other fantastical shapes in vibrant hues.

\end{DndReadAloud}

This colorful cavern is the lair of a \textbf{chimera} named Chossos and a \textbf{gorgon} named Fossus. Chossos speaks Common, Draconic, and Giant. He resembles a typical chimera, except he has a bull head in place of a goat head. This head confers no special abilities on Chossos, but it has helped him foster a kinship with Fossus.

When they aren't hunting together, Fossus restlessly paces among the cavern's vibrant stalagmites, while Chossos rests on a ledge above the entrance, from which the chimera ambushes any who pass by. The two creatures work together to burn, petrify, and trample intruders.

If Chossos is reduced to 50 hit points or fewer, the chimera feigns surrender and vows to give the characters his treasure if they spare him. If Fossus is still alive, Chossos calls off the gorgon. Chossos then attempts to lure the characters closer before resuming his assault once they've dropped their guards.

\subparagraph{Treasure}
The remains of a petrified adventurer stand among the stalagmites. A \textit{Bag of Holding} hangs from the drab statue, which is missing its head. The bag contains a set of robes, 9 days of rations, a lacquered pipe, and a scorched \textit{spellbook} containing the following spells: \textit{Darkvision}, \textit{Fog Cloud}, \textit{Grease}, \textit{Identify}, and \textit{Web}.

Piled atop the ledge above the entrance is a small heap of treasure containing 900 gp and 700 sp.

\subsubsection{L17: Boulder Heap}
\begin{DndReadAloud}{}
Small, rounded boulders are carefully stacked here, blocking the entrance to a wide tunnel.

\end{DndReadAloud}

The boulders are obviously not a natural formation. Chossos (see area L16) piled them here to block the entrance to the greater caverns and prevent creatures that dwell there from taking him by surprise.

With 5 minutes of work, the characters can clear enough of the rocks to expose the tunnel beyond. There, a flight of nine hundred crudely shaped stone steps descends to area G1 of the greater caverns (detailed later in this adventure).

\subparagraph{Hollow Boulder}
If the characters move any of the rocks, they uncover a smooth, blue-green boulder that's strangely light and clunks dully when moved or shaken, revealing its hollowness. The boulder can be easily broken open with blunt weapons, with tools, or by smashing it against a hard surface.

\subparagraph{Graven Glyphs}
Inside the hollow boulder is a bronze tablet worth 100 gp. The tablet is etched with a poem meant for those bold enough to delve into the greater caverns. It offers cryptic insight on how to enter Iggwilv's inner sanctum (area G21), where the archmage's true prize resides. The poem reads as follows:

\begin{DndReadAloud}{}
In the center lies the gate, But its locks will surely vex. Many are the guards who wait As you seek the middle hex.

Randomly sent to find a way Back to a different iron door. A seventh time and you may stay, And seek the glowing prize no more.

You have won old Iggwilv's prize, Her hoarded cache of magic, And freed the one with yearning eyes, Whose lot was hunger tragic.

\end{DndReadAloud}

\subsubsection{L18: Lavish Parlor}
\begin{DndReadAloud}{}
The natural cavern gives way to a luxurious parlor. Crystal chandeliers bathe the chamber in soft light, and the scent of orange blossoms perfumes the air. Plump cushions surround a low table heaped with ripe fruits, refreshing beverages, and succulent meats. Low sofas and comfortable chairs dot the walls of the

lounge, and side tables between these sitting areas overflow with gemstones, jewelry, and other treasures.

Four impeccably dressed human aristocrats laugh and feast in the parlor. An unseen string quartet accompanies their revelry with soothing music.

\end{DndReadAloud}

This area is brightly lit. The nobles kindly invite the characters to join the repast, insisting they rest and regale the nobles with tales of adventure.

The nobles are part of an elaborate, genie-wrought illusion created by Kashem, a bitter \textbf{dao} bound to guard this chamber by Iggwilv. The Witch Queen imparted a mote of her power to Kashem in exchange for the dao's services. Long after the archmage's departure, Kashem still pursues her original charge to destroy all who enter.

When the characters arrive, Kashem floats silently above the chamber, waiting for a prime opportunity to strike. If the characters consume any of the food or drink, refuse to join the nobles outright, or see through the illusions, Kashem reveals herself and attacks.

\subparagraph{Dust to Dust}
If Kashem is slain, her body disintegrates into a pile of fine dust. This dust can be used in area L21 to destroy the cage in the lake or prove that Kashem is dead.

\subparagraph{Illusory Finery}
The room's furnishings and inhabitants are convincing, genie-wrought illusions. These illusions blanket dingy, neglected versions of the objects they cover and feel real if touched. To pierce the dao's ruse, a character examining the chamber's fixtures must succeed on a DC 20 Intelligence (Investigation) check. The nobles, however, fail to hold up to physical inspection. A character who touches a noble automatically succeeds on this check.

\subparagraph{Poisoned Feast}
The food and drink on the table are real—and poisoned. A creature that consumes any of the food or drink must make a DC 16 Constitution saving throw. On a failed save, a creature takes 21 (6d6) poison damage and has the poisoned condition for 1 hour. On a successful save, the creature takes half as much damage only. The illusory nobles nonchalantly partake of the sumptuous feast and politely encourage the characters to sample the dishes.

\subparagraph{Treasure}
None of the treasure is real. Coins and precious objects taken from the chamber turn out to be flat stones and obvious counterfeits.

\subsubsection{L19: Black-Water Lake}
\begin{DndReadAloud}{}
The river flows into a flooded cavern, creating a small lake. The surface of the water is glossy and black. A rocky island rises from the center of the pool.

\end{DndReadAloud}

The island is strewn with tiny garnets, which sparkle alluringly in the presence of light.

\subparagraph{Lurking Death}
Eight aquatic ghasts (each has a swimming speed of 30 feet) lurk in the lake, which is 20 feet deep. Spotting them beneath the water's mirrorlike surface requires a successful DC 20 Wisdom (Perception) check. The ghasts ambush any creatures that enter the lake.

\subparagraph{Rebuking Island}
Iggwilv recorded a series of messages in this cave to taunt intruders. Each time a creature sets foot on the island, an imperious feminine voice echoes throughout the cave, uttering one of the following reproofs:


\begin{itemize}
\item "Fools! This is a dead end. Flee while you still can!"
\item "You were witless to enter at all, so you should probably just stay here."
\item (Peals of hideous, mocking laughter)
\end{itemize}

\subparagraph{Treasure}
Six hundred garnets are strewn across the island, each the size of a pea and worth 1 gp. It takes 10 minutes for a creature to collect the gems by hand. They have a combined weight of 5 pounds.

\subsubsection{L20: Cave of Crystals}
\begin{DndReadAloud}{}
Veins of glittering quartz streak the walls of this crystal cave. Three bizarre creatures with stone-lidded eyes converse in grating tones. Each has three clawed arms, three stumpy legs, and a barrel-like body.

\end{DndReadAloud}

Three \textbf{xorn} from the Elemental Plane of Earth found their way to this cave through a planar rift that has since closed, stranding them here. The xorn are initially indifferent to the characters but defend themselves if attacked.

\subparagraph{Hungry Castaways}
The xorn are discussing their plight when the characters arrive. The creatures are hungry and upset at their situation. Prone to stress-eating, the xorn crave any precious gems or metals they can sense on the characters. If anyone in the party can understand the xorn, the creatures explain their dilemma and attempt to persuade the characters to give them a meal of at least 100 gp worth of gems or metals per xorn.

If the characters have a means to send the xorn home—such as the \textit{Banishment} spell—the xorn offer to help locate other treasure for the characters as payment. Alternatively, a character can convince the xorn to assist them in defeating a single threat elsewhere in the caverns with a successful DC 18 Charisma (Persuasion) check.

\subparagraph{Treasure}
The xorn stashed 500 gp worth of uncut gems behind a large crystal formation in the northwest corner of the room. A character searching the crystals finds the stash with a successful DC 17 Wisdom (Perception) check.

\subsubsection{L21: Pool Cavern}
\begin{DndReadAloud}{}
The narrow river flows into a domed cavern whose ceiling bristles with stalactites. Faint sounds of trickling water echo from the northwest.

\end{DndReadAloud}

The river dead-ends in this flooded chamber. A \textbf{marid} named Kasdu'ul is trapped at the bottom of the pool, 30 feet below the surface. An unbreakable stone cage etched with glowing runes confines the genie to the pool. Characters who look into the pool and succeed on a DC 15 Wisdom (Perception) check notice the cage's faint, glimmering light in the water.

\subparagraph{Bound Genie}
Kasdu'ul was trapped here by the dao Kashem (see area L18) centuries ago, shortly before Kashem was fooled and bound by Iggwilv in turn. A permanent \textit{Antipathy/Sympathy} spell (save DC 18) protects the cage and repels creatures that aren't native to the Elemental Plane of Earth with its Antipathy effect. The runes on the cage invoke the hatred of Ogrémoch, the Prince of Evil Earth.

While the cage is intact, Kasdu'ul has the incapacitated condition and speaks in a lethargic tone. If the characters overcome the Antipathy effect and approach the cage, the genie drowsily asks them to destroy it.

\subparagraph{Destroying the Cage}
The cage is immune to all damage and has no locks or opening mechanisms. It can be destroyed in the following ways:


\begin{itemize}
\item A creature native to the Plane of Earth—such as one of the pechs in area L6—can easily break the cage as if it were made of fragile clay.
\item If the dust of a slain dao is scattered over the lake, the cage instantly rusts away.
\item A successful \textit{Dispel Magic} (DC 18) spell cast on the cage causes it to crumble.
\end{itemize}

Kasdu'ul has had many years to contemplate the cage's magical construction. She might propose one of these methods to the party from inside the cage.

\subparagraph{Freed Genie}
If the cage is destroyed, Kasdu'ul returns to her lively self. The grateful genie urges the characters to find and slay the dao that imprisoned her if they haven't already.

As thanks for freeing her, she's willing to cast \textit{Tongues}, \textit{Water Breathing}, and \textit{Water Walk} on the characters to facilitate their exploration. If the characters also present proof of Kashem's demise—such as a pinch of the dust left behind after her death—or they offer a convincing story and succeed on a DC 17 Charisma (Deception) check, Kasdu'ul summons a \textbf{water weird} and orders it to serve the characters for 1 hour, after which it vanishes. In either case, Kasdu'ul then bids the characters farewell and returns to the Elemental Plane of Water.

\subsection{Greater Caverns Features}
Unless otherwise noted, the greater caverns have the following features:


\begin{description}
\item[Ceilings.]
Ceilings are 15 feet high in corridors and tunnels and 30 feet high in chambers.
\item[Corrupted Region.]
Abyssal energy bleeds into the greater caverns. Phantom whispers and twisted shadows pluck at the characters' fears, and nightmares plague their rest. A character who finishes a long rest in the greater caverns must make a DC 15 Wisdom saving throw. On a failed save, the character gains only the benefits of a short rest.
\item[Lighting.]
The caverns are unlit. Denizens carry their own light or rely on darkvision to see. Area descriptions assume the characters have a light source or other means of seeing in the dark.
\end{description}

\subsection{Greater Caverns Locations}
The following locations are keyed to map 6.3.

\includegraphics[width=0.4\textwidth]{images/../5e.tools/img/adventure/QftIS/106-map-6.03-greater-caverns-map.png}
\includegraphics[width=0.4\textwidth]{images/../5e.tools/img/adventure/QftIS/107-map-6.03-greater-caverns-map-player.png}
\subsubsection{G1: Warren Tunnels}
\begin{DndReadAloud}{}
The domed ceiling of this chamber drips with stalactites. Thick fungal growth blankets the floor, and the still air reeks with the odors of rotting refuse, moldy fungus, and an even less pleasant stench. Numerous tunnels branch from the cave.

\end{DndReadAloud}

The stench permeating these shallow tunnels arises from the compost and manure in the carefully tended fungus beds.

The tunnels are home to a clan of troglodytes warped by demonic energy radiating from Iggwilv's experiments. Eighteen \textbf{troglodyte} warriors live in the warrens. The clan is led by a bloated \textbf{hezrou} servant of Laogzed, a lesser demonic deity of mindless gluttony venerated by these troglodytes. Blessed by Laogzed, the troglodytes are immune to the poisoned condition.

\subparagraph{Foul Welcome}
Six troglodyte warriors lurk near the entry, waiting to ambush creatures arriving from the lesser caverns. The characters can sneak past the troglodytes with a successful DC 15 Dexterity (Stealth) group check. Characters carrying a light source have disadvantage on this check. If combat breaks out here, the hezrou joins the fight on initiative count 20 of the second round (losing initiative ties). Six more troglodytes arrive on initiative count 0 of the third and fourth rounds.

\subparagraph{Treasure}
The hezrou's chamber contains a Potion of Greater Healing and a Longsword of Vengeance.

\subsubsection{G2: Cavern of Corpses}
\begin{DndReadAloud}{}
This gruesome cave is either a burial crypt or a trophy room. The walls of the vaulted cavern are lined with limed-over corpses. A circle of barely distinguishable bodies composes the lower tier, while a fresher ring of calcified corpses stands on the heads and shoulders of those beneath, completing the macabre decoration.

The steady drip of water from the ceiling down the walls and over the stony corpses creates a shallow pool in the northeastern portion of the cave.

\end{DndReadAloud}

The calcified corpses include dead explorers, troglodytes, and cave-dwelling creatures. A wicked spirit of hunger (use the \textbf{wraith} stat block) resides here, a sliver of the lesser god Laogzed's power. The troglodytes worship this ravenous spirit and offer up captives and treasure to it as sacrifices. Seven specters created by the spirit lie dormant in their former bodies, awaiting its command.

When the characters enter the cavern, the spirit and specters attack, eager to feed off the characters' life force.

\subparagraph{Treasure}
Scattered on the floor at the far end of the chamber are 480 gp, 800 sp, twenty assorted gems worth 50 gp each, and a sealed scroll tube with a \textit{Spell Scroll} of \textit{Raise Dead}.

\includegraphics[width=0.4\textwidth]{images/../5e.tools/img/adventure/QftIS/108-06-008.spirit-of-hunger.png}
\subsubsection{G3: Great Gallery}
\begin{DndReadAloud}{}
Irregular shelves line the walls of this yawning cavern. Massive stalactites hang from the ceiling, and the floor is dotted with matching stalagmites. Fungus grows throughout the cave, and small cave animals scurry about, feeding on the mushrooms.

\end{DndReadAloud}

The ceiling of this long cavern rises to a height of 50 feet. Careful examination of the floor reveals many cracked bone fragments among the fungi. This is the lair of Lludd, a clever and long-lived \textbf{behir} who speaks Common and Giant in addition to Draconic. Lludd lurks on a ledge near the entrance to the chamber, waiting for morsels worthy of his hunger.

Lludd is hostile toward intruders, but he doesn't immediately resort to violence. Though he's selfish and cruel, the behir hasn't survived this long by attempting to devour every creature that enters his lair—especially hardened adventurers. When he notices the characters, Lludd greets them in Common.

\subparagraph{Parleying with Lludd}
If the characters converse with Lludd, he offers to point them toward a place where they can find treasure in exchange for leaving him in peace. The behir keeps no hoard of his own. If they agree, Lludd directs them toward a "vast treasure" in the rolling gallery (area G4).

During the conversation, Lludd sizes up the party. If the characters appear wounded or offend the behir, he seizes his chance and unleashes his breath on them.

\subparagraph{Treasure}
Inside Lludd's stomach is a \textit{Periapt of Proof against Poison} left over from an unlucky explorer the behir digested.

\subsubsection{G4: Rolling Gallery}
\begin{DndReadAloud}{}
This vast, tubelike gallery's concave walls curve from long ledges about halfway to the ceiling down to the floor. Tracks have worn smooth paths from ledge to ledge across the curving floor.

The thunderous griding of stone echoes throughout the cave, as three boulders roll along the tracks at dangerous speeds. Other large rocks, worn nearly spherical, sit motionless throughout the gallery.

\end{DndReadAloud}

The rolling boulders are three \textbf{galeb duhr}, created and bound by Iggwilv to defend this cave. The elementals attack any intruders.

\subparagraph{Treasure}
A scattering of treasure glints in the bottom of the cave: 60 pp, 285 gp, 87 sp, three rubies worth 300 gp each, and a \textit{+1 War Pick}.

\subsubsection{G5: Tilted Cavern}
\begin{DndReadAloud}{}
The terraced floor of this angled cavern is slippery with cave drippings. In the center of the cave, a wide tunnel has collapsed into a rubble-filled alcove.

\end{DndReadAloud}

Two umber hulks burrowed through this area recently. Though their tunnels collapsed behind them, the creatures lurk just beyond the alcove. If the irritable umber hulks detect the characters, the monsters burst through the rubble and attack.

\subparagraph{Slippery Floor}
The dripping water renders the uneven floor slick with mud and slime. A creature that moves faster than half its speed along the floor must succeed on a DC 15 Dexterity saving throw or have the prone condition. The umber hulks ignore this effect.

\subparagraph{Tunnel}
If the characters excavate the entrance to the tunnel through which the umber hulks entered this cavern, they find it runs down and away for 100 feet. It then angles upward and continues for 1 mile, exiting through the north side of the mountain.

\subsubsection{G6: Small Gallery}
\begin{DndReadAloud}{}
This cave holds numerous types of fungi. Three giant lizards graze warily on the mushrooms.

\end{DndReadAloud}

The three giant lizards pose no threat to the characters. As they graze, the lizards nervously eye anyone who approaches them and keep their distance.

If a character casts the \textit{Speak with Animals} spell or uses other means to communicate with the lizards, the reptiles explain one of their siblings was recently slain by a terrible, red-furred creature from a large cave to the west (see area G8). If the characters avenge the lizards' fallen kin, one of the lizards agrees to serve the party as a cavern guide and mount for the next hour.

\includegraphics[width=0.4\textwidth]{images/../5e.tools/img/adventure/QftIS/109-06-009.barlgura.png}
\subparagraph{Mushrooms}
The fungi in this cavern are edible, and characters can forage food here without having to make an ability check.

\subsubsection{G7: Smooth Cavern}
\begin{DndReadAloud}{}
This wide cavern is strewn with boulders. A gaping sinkhole drops down near the cave's center and plummets into darkness.

\end{DndReadAloud}

Two gas spores grow in this cavern, one tucked among the boulders near the sinkhole and the other in the hole itself. When a creature comes within 30 feet of either spore, both spores attack.

\subparagraph{Sinkhole}
The sinkhole drops 75 feet down to a tunnel that runs 35 feet to the east, then drops another 100 feet to an underground stream. What lies beyond the stream is outside the scope of this adventure.

\subsubsection{G8: Cavern of Stalagmites}
\begin{DndReadAloud}{}
Four thick stalagmites, twenty feet high and with tops worn smooth, rise from the center of this large cave. Baleful crimson light leaks from cracks in the jagged pillars, casting an eerie glow on the cave.

\end{DndReadAloud}

Four barlguras lair in this dimly lit cave. Each demon perches on one of the pillars, waiting to ambush unsuspecting entrants. When combat begins in this area, the pillars activate.

\subparagraph{Demon Pillars}
Iggwilv conducted demonic experiments on the pillars in this cavern. Each round on initiative count 20 (losing initiative ties), the pillars pulse with Abyssal energy. When this occurs, non-Fiends within 10 feet of one or more pillars must succeed on a DC 15 Constitution saving throw or be warped by planar magic. Roll on the Abyssal Energy table to determine the effect, which lasts until the pillars pulse again or are destroyed.

\begin{DndTable}[header=Abyssal Energy]{cX}

d4 & Abyssal Manifestation\\
1 & The target's limbs become rubbery and unreliable. It has disadvantage on attack rolls with weapons.\\
2 & The target's eyes squirm out of its head as worms; the target has the blinded condition.\\
3 & The target's feet become heavy, fluid-filled sacs, reducing its walking speed by 10 feet.\\
4 & The target's mouth vanishes, preventing it from speaking or casting spells with verbal components.\\
\end{DndTable}

\subparagraph{Destroying the Pillars}
Each of the pillars has Armor Class 17, 50 hit points, immunity to poison and psychic damage, and vulnerability to radiant damage. Destroying the pillars prevents them from emitting Abyssal energy. A \textit{Hallow} spell cleanses the pillars, accomplishing the same feat.

\subsubsection{G9: Glowing Grotto}
\begin{DndReadAloud}{}
Unusual fungi and lichen cover this irregular cavern. The growths are limned in auras of blue and pink light. Many-faceted crystals grow around the periphery, reflecting the light in scintillating patterns.

\end{DndReadAloud}

This dimly lit cave is a planar nexus through which Iggwilv once explored the multiverse. Fungi from the Feywild and crystals harvested from the deepest reaches of the Astral Plane fuel the cave's glow and anchor the planar magic that lingers here.

\subparagraph{Planar Displacement}
When a creature enters the grotto, a flare of magic teleports the entrant and creatures within 10 feet of it to one of the Outer Planes. Roll on the Grotto Destination table to determine where the creatures are sent. For the next minute, any creature that enters the grotto is transported to the same destination. Due to their planar nature, these areas don't appear on map 6.3.

The following sections detail these locations and how creatures transported to them can return to the Lost Caverns. Once a creature has been transported to a location in this way, it can't be sent to that destination by entering the grotto for 24 hours.

\begin{DndTable}[header=Grotto Destination]{cX}

d4 & Destination\\
1 & Abyssal labyrinth (area G9a)\\
2 & Elysian canyon (area G9b)\\
3 & Infernal hall (area G9c)\\
4 & Mechanized checkpoint (area G9d)\\
\end{DndTable}

\subsubsection{G9a: Abyssal Labyrinth}
\begin{DndReadAloud}{}
A flash of sickly green light fades to reveal rough-hewn tunnels bored into black stone. The air reeks of animal stench. Chuffing growls echo from the darkness.

Two muscular humanoids with bull heads emerge from the darkness. They brandish axes and lead metallic bulls whose nostrils exude a noxious vapor.

\end{DndReadAloud}

This immense labyrinth is part of the Endless Maze—a layer of the Abyss ruled by the bloodthirsty demon lord Baphomet. The labyrinth fills a region of the Abyss roughly 1 mile on each side. The labyrinth's tunnels are 20 feet wide, 50 feet high, marred with scratches, and stained with blood.

The characters arrive in the presence of two minotaurs, each accompanied by a \textbf{gorgon}. Servants of Baphomet, the minotaurs seek to offer up the characters' heads as gifts to the demon lord. The gorgons obey the minotaurs' spoken commands.

\subparagraph{Escaping the Labyrinth}
When the characters negotiate the labyrinth, have them choose a navigator and their pace (fast, normal, or slow). It takes 1 hour to navigate the maze at a fast pace, 2 hours at a normal pace, or 3 hours at a slow pace. Halfway through the journey, the navigator must make a DC 16 Wisdom (Survival) check. If the characters are moving at a fast pace, the navigator makes the check with disadvantage, while a slow pace allows them to make the check with advantage. On a successful check, the characters find their way out. On a failed check, the characters must restart the labyrinth.

The longer the characters take to navigate the maze, the greater the risk of encountering the creatures within it. Every 30 minutes the party spends in the maze, roll on the Labyrinth Encounters table to see what, if anything, the characters meet.

The labyrinth's exit is a one-way portal at the end of a clean, misty tunnel. Creatures that pass through the portal emerge in front of area G9 of the Lost Caverns.

\begin{DndTable}[header=Labyrinth Encounters]{cX}

d20 & Encounter\\
1–10 & No encounter\\
11–12 & A row of humanoid heads on pikes\\
13–14 & A bloodthirsty \textbf{hill giant} with crooked horns\\
15–16 & An enormous, rampaging giant bull (use the \textbf{mammoth} stat block)\\
17–18 & 1d6 minotaurs\\
19–20 & An aspect of the demon lord Baphomet (use the \textbf{goristro} stat block)\\
\end{DndTable}

\subparagraph{Treasure}
Near the labyrinth's exit, a desiccated, severed hand clutches a \textit{+2 Greataxe}.

\subsubsection{G9b: Elysian Canyon}
\begin{DndReadAloud}{}
A flash of golden light gives way to warm sunlight. The rocky floor of a steep canyon stretches before you, its sheer walls rising overhead. Fallen rocks—the result of a recent rockslide—block the canyon's exit.

Six centaurs stand nearby. One trots forward and calls out to you. "Unfortunate timing! We're all trapped, but perhaps we can help each other?"

\end{DndReadAloud}

This box canyon on the plane of Elysium tests the compassion and resolve of those trapped within. Six centaurs who speak Common were exploring the canyon when its walls collapsed, stranding them here. Unable to scale the sheer, 60-foot-high walls, the centaurs are out of ideas. The centaurs are friendly toward the characters unless attacked.

Shaydun, the centaurs' leader, approaches the characters. She doesn't know how the characters arrived in the canyon, but she's willing to help them find a way back to the Lost Caverns if they help the centaurs escape the canyon.

\subparagraph{Escaping the Canyon}
The characters can escape the canyon by the following means:


\begin{description}
\item[Clearing the Rubble.]
A character who inspects the rubble and succeeds on a DC 15 Intelligence (Investigation) check notices a protruding boulder that, if dislodged, could open a way through to the canyon beyond. A character can knock the boulder loose with a successful DC 25 Strength (Athletics) check. If the characters fasten ropes to the boulder, the centaurs offer to help haul it down for them (no check required). Freeing the boulder clears a path through the rubble.
\item[Climbing Out.]
Scaling the canyon's walls without climbing gear or magic requires a successful DC 20 Strength (Athletics) check. The centaurs can't manage this climb. At the top of the canyon is a golden-barked oak tree the characters can use as an anchor for rope or other climbing equipment.
\end{description}

At your discretion, other means of exiting the canyon might exist. Either way, once the characters have freed themselves, they can easily help the centaurs do the same.

\subparagraph{Treasure}
Shaydun wears a magic necklace adorned with a small silver horseshoe. The necklace allows the creature wearing it to cast the \textit{Plane Shift} spell once, after which the necklace loses its magic.

In exchange for helping the centaurs escape the canyon, Shaydun gifts the characters her necklace. If a character uses the necklace to return to the Lost Caverns, they do so unerringly, appearing in front of area G9.

\subsubsection{G9c: Infernal Hall}
\begin{DndReadAloud}{}
A flash of red light fades to reveal a seemingly endless obsidian hall. Rows of decorative columns hoist a flat, mirrorlike ceiling high above. The columns are carved in the likeness of two alluring yet imposing devils whose eyes seem to follow you. Ruddy light filters into the room from somewhere beyond the pillars, and the air is dry and oppressively hot.

Beneath each of you is a faintly glowing pentacle of green metal inlaid in the stone floor.

\end{DndReadAloud}

This infernal hall is located within an abandoned courthouse on Phlegethos, the fourth layer of the Nine Hells. The hall is 60 feet wide and has 30-foot-high ceilings, but it is infinitely long and unbearably hot. The hall constitutes an area of \textit{extreme heat} (see the \textit{Dungeon Master's Guide}).

When the characters arrive, each of them is trapped in one of the pentacles (detailed below).

\subparagraph{Devil Columns}
The figures carved into the columns represent Belial and Fierna, joint rulers of the fourth layer of the Nine Hells. A character who succeeds on a DC 20 Intelligence (Religion) check recognizes the archdevils and the plane on which the hall resides. A character can climb the columns without having to make an ability check.

\subparagraph{Mirror Ceiling}
A character who looks up at the ceiling sees the floor of the glowing grotto (area G9) reflected in its surface. A character who studies the ceiling and succeeds on a DC 16 Wisdom (Perception) check notices faint ripples in its surface, like that of a pool.

\subparagraph{Pentacles}
A 10-foot-diameter, 20-foot-tall cylinder of magical force surrounds each pentacle, confining a creature within its bounds. Skeletons of creatures that never escaped the magical prisons dot the pentacles along the infinite hall. Fortunately for the characters, Iggwilv built a fail-safe into the pentacles when she created them ages ago in case her enemies ever tried to use them against her.

A \textit{Levitate} spell allows a creature imprisoned in one of the pentacles to rise through the prison's area of effect and up to the ceiling. This exception doesn't apply to other forms of flight, such as the \textit{Fly} spell or wings.

A character who examines the pentacle beneath them and succeeds on a DC 15 Intelligence (Investigation) check notices smudges along the metal, suggesting someone traced the shape in the past. If the character traces the shape of the pentacle imprisoning them, they gain the effects of a \textit{Levitate} spell for 1 minute, with no concentration required.

\subparagraph{Escaping the Hall}
To escape the hall, the characters must overcome their magical prisons. They can escape the pentacles through the following means:


\begin{description}
\item[Dispelling the Circle.]
A successful \textit{Dispel Magic} spell (DC 15) cast on the pentacle suppresses its magic for 1 hour, freeing the creature inside.
\item[Levitation.]
A \textit{Levitate} spell allows a character to rise out of the pentacle's area of effect.
\item[Teleportation.]
A character who first succeeds on a DC 15 Charisma saving throw can then teleport outside the circle as normal. On a failed save, the attempt is wasted.
\end{description}

A creature that ascends to the ceiling is transported to the Lost Caverns in front of area G9.

\subparagraph{Treasure}
Propped against a column near one of the pentacles is a Glaive of Warning veined with hellfire.

\subsubsection{G9d: Mechanized Checkpoint}
\begin{DndReadAloud}{}
A copper-hued light flares and then fades, revealing a doorless, perfectly cubic room made of brass, copper, and steel. The walls are composed of interlocking gears, metal plates, and whirring machinery that rotate, shift, and tick to a synchronized beat.

The center of a wall reconfigures and opens, revealing three armored automatons. They march forward in unison, drawing weapons and readying shields.

\end{DndReadAloud}

This checkpoint is located on Mechanus, a plane of absolute order. The three automatons (lawful neutral helmed horrors) attack the characters, who are unauthorized visitors. On initiative count 20 of the next round (losing initiative ties), up to three more automatons (lawful neutral suits of \textbf{animated armor}) emerge from the wall to replace any that were slain. These reinforcements continue to appear until the characters leave the area or die trying.

\subparagraph{Escaping the Checkpoint}
When an automaton is destroyed, its armored body crumbles to the floor as individual components. This includes its intact helmet, whose interior glows with a similar copper light to the flare that deposited the characters here. A character who dons the helmet is transported back to the Lost Caverns in front of area G9. The helmet then crumbles to dust.

\subsubsection{G10: Jagged Cavern}
\begin{DndReadAloud}{}
This cavern is dimly lit by a multitude of candles. The walls are jagged, and the cave's floor is littered with broken skulls and splintered bones. In the center of the cave, a staircase ascends to a carved stone altar etched with leering demon faces. A sinister tome bound with black demon hide rests atop the altar.

\end{DndReadAloud}

Iggwilv once conjured demons and bound them to her service in this soaring cavern. Before she vanished, the Witch Queen left a deadly trap for anyone coveting her knowledge.

\subparagraph{False Demonomicon}
The book on the altar looks like the Demonomicon of Iggwilv, the preeminent treatise on the Abyss and demons. A character unfamiliar with the book can identify it with a successful DC 20 Intelligence (Arcana) check. However, the book on the altar is a trap: a fake tome containing an unstable seed of Abyssal energy.

If the book is moved or touched, it flips open, and its pages shuffle as if disturbed by a phantom wind. A coruscating, 5-foot-wide, 10-foot-tall slash of red energy tears through the air above the pages, revealing a marred Abyssal landscape beyond.

Two chasmes buzz out from the rift and attack. At the end of each round, roll a d20. On a 10 or lower, another \textbf{chasme} wriggles through the portal.

\includegraphics[width=0.4\textwidth]{images/../5e.tools/img/adventure/QftIS/110-06-010.the-tome.png}
\subparagraph{Closing the Portal}
The portal remains open for 1 hour or until six demons emerge, whichever comes first. The portal can be closed early through any of the following means:


\begin{description}
\item[Closing the Book.]
A character within 5 feet of the altar can close the book as an action with a successful DC 25 Strength (Athletics) check.
\item[Destroying the Book.]
The book has Armor Class 15, 50 hit points, resistance to cold and fire damage, and immunity to poison and psychic damage.
\item[Dispelling the Book.]
A \textit{Dispel Magic} spell cast on the book suppresses its magic for 1 hour.
\end{description}

Closing the portal doesn't return escaped chasmes to the Abyss.

\subparagraph{Treasure}
Amid the bones on the floor is a \textit{Cloak of Elvenkind}.

\subsubsection{G11: Cave of the Skull}
\begin{DndReadAloud}{}
A small silver coffer rests in the center of this cave. A bleached humanoid skull rests atop the box. A leather cloak lies near the skull.

\end{DndReadAloud}

The skull belonged to a wizard who challenged Iggwilv and lost badly. The garment nearby, however, is a motionless \textbf{cloaker} that ambushes any creature that approaches the coffer.

\subparagraph{Treasure}
The coffer is worth 750 gp. It contains ten cracked emeralds worth 100 gp each and a jar of \textit{Keoghtom's Ointment}.

\subsubsection{G12: Pillared Cavern}
\begin{DndReadAloud}{}
Stalactites and stalagmites in this chamber have fused to form lumpy pillars extending from floor to ceiling.

\end{DndReadAloud}

Two ropers lurk among the calcite columns at the cavern's south end. The monsters attack when they can each lash out at a victim with their tendrils.

\subsubsection{G13: Rotting Fungus Cavern}
\begin{DndReadAloud}{}
This cave is thick with the stench of rotting fungus. A few mushrooms grow here and there, but most are dead and heaped in piles throughout the chamber.

\end{DndReadAloud}

The remains of rotting mushrooms litter this chamber. Three shambling mounds thrive amid the detritus. When the characters enter the cave, the creatures lurch out of the piles of filth and attack.

The stone giant in area G14 regularly tosses dead mushrooms in to feed the shambling mounds.

\subparagraph{Treasure}
Buried in the rotting mush is a collection of shiny objects the shambling mounds have gathered, including 340 gp, a \textit{Potion of Clairvoyance}, and a silver scroll tube worth 125 gp containing a \textit{Spell Scroll} of \textit{Cloudkill}.

\subsubsection{G14: Giant Lair}
\begin{DndReadAloud}{}
Bones are scattered on the floor here. Two large round, black rocks sit in the rough center of the cave. A towering stone giant looms near the far wall. His skin is marred by cracks that radiate pale-green light.

\end{DndReadAloud}

This is the lair of Zyzkathol, a chaotic evil \textbf{stone giant} who succumbed to the demonic energy that permeates the greater caverns. He's accompanied by two giant beetles (use the \textbf{rhinoceros} stat block). When at rest, the beetles resemble boulders, requiring a successful DC 18 Intelligence (Investigation) check to discern they're animate.

When he notices the characters, Zyzkathol barks a command and sets the beetles on them.

\subparagraph{Treasure}
Hidden in Zyzkathol's bedding is a sack sewn from a giant weasel pelt that contains 375 gp. The giant wears an ivory necklace worth 200 gp.

\subsubsection{G15: Fungus Garden}
\begin{DndReadAloud}{}
This cave is full of six-foot-tall mushrooms and other fungi growing in mounds of carefully spread compost.

\end{DndReadAloud}

The fungus here is a food source for many of the greater caverns' creatures, including eight shriekers that are currently feeding here. If the shriekers sound off, Zyzkathol and his giant beetles (see area G14) come to investigate.

\subsubsection{G16: Lofty Cavern}
\begin{DndReadAloud}{}
The ceiling of this large cave soars overhead, and natural ledges on the walls spiral upward. Broken stalactites litter the floor, and a pile of treasure is heaped against the west wall.

\end{DndReadAloud}

An \textbf{adult black dragon} named Calumnus recently moved into this vast cavern, whose natural ceiling rises to a height of 60 feet. Calumnus is aggressive and territorial, attacking all who enter, but the dragon retreats if reduced to 50 hit points or fewer.

\subparagraph{Treasure}
The treasure near the west wall is the dragon's hoard. It includes 67 pp, 235 gp, 2,750 sp, assorted gems worth 700 gp, and an ivory case worth 300 gp containing a \textit{Wand of Magic Missiles}.

\subsubsection{G17: Sanctuary Cavern}
\begin{DndReadAloud}{}
This large cave is awash with a sense of peace. The low ceiling and walls are striated with brightly colored rock. A trickle of spring water runs through a crack and collects in a small natural basin.

\end{DndReadAloud}

An aura of purity suffuses this tranquil cave. The greater caverns' fouler denizens avoid this area, making it a safe place for the characters to rest.

\subparagraph{Hallowed Rest}
This haven was made long ago by one of Iggwilv's enemies. A permanent \textit{Hallow} spell (save DC 17) blankets the cave, with the following extra effect: Fiends and Undead can't see inside the cavern. The Corrupted Region feature of the greater caverns (see the "Greater Caverns Features" section) doesn't extend to this area.

\subsubsection{G18: Grotto of the Idol}
\begin{DndReadAloud}{}
A great stone idol stands over this grotto. The statue is carved in the likeness of a porcine-headed demon with puny wings and two jacinth eyes. On either side of it are brass stands with copper-bladed weapons.

A chest at the statue's feet rests with its lid open, revealing a collection of rings, bracers, brooches, small weapons, cloaks, arrows, and rods.

\end{DndReadAloud}

This chamber is yet another trap devised to foil would-be plunderers. The demonic idol uses the \textbf{stone golem} stat block but is carved in the likeness of a nalfeshnee. While the idol remains motionless, it's indistinguishable from an ordinary statue.

When the characters enter the chamber, a \textit{Magic Mouth} spell cast on the idol utters the following:

\begin{DndReadAloud}{}
"Leave one magic item before me, then take your choice of those others have left behind!"

\end{DndReadAloud}

If the characters attempt to take any of the items or otherwise disturb the idol, the idol attacks.

\subparagraph{Illusions}
The treasure and weapons in this room are illusions. A character who studies the objects and succeeds on a DC 16 Intelligence (Investigation) check discerns the illusions for what they are.

\subparagraph{Treasure}
The only legitimate treasures in this grotto are the idol's two jacinth eyes, which are worth 250 gp each.

\subsubsection{G19: Vestibules}
Six nearly identical areas lead from the tunnels into Iggwilv's inner sanctum (area G21). When the characters approach one of the outer doors, read or paraphrase the following text:

\begin{DndReadAloud}{}
An immense iron double door blocks your passage forward. A smokeless torch clutched in the mouth of a stone demon face above the doors illuminates the area. The doors are adorned with leering demon visages in bronze, and the doors' riveted edges are etched with arcane symbols.

The two largest demon faces hold unfurled bronze scrolls inlaid with a spidery script in their mouths. Large bronze pull rings sit below the scrolls.

\end{DndReadAloud}

Each set of doors is 10 feet wide and 20 feet tall. The script around the edges of the doors is a mixture of Abyssal writing and arcane sigils of Iggwilv's design. The bronze scrolls are inscribed with the following warning:

\begin{DndReadAloud}{}
Iggwilv's treasure rests within,

Her curse on any who disturbs it.

Seek no further to steal it or

To free the one prisoned here,

For a fate worse than death is

Sure to come to fools who

Violate this sacred place.

\end{DndReadAloud}

The doors are unlocked but immensely heavy. When the characters enter a vestibule for the first time, read or paraphrase the following text:

\begin{DndReadAloud}{}
A corridor with walls of red marble leads to an ebony door with silver hinges, studs, and a pull ring. The arched ceiling is made of black marble shot through with scarlet veins. The walls and floor are polished and dust-free, and a thick black carpet runs the length of the hall. Dim, ambient light illuminates the space.

\end{DndReadAloud}

The first hall the characters enter also contains a framed painting of a narrow boat on turbulent seas. A pale, raven-haired woman stands at the tiller, smiling under a stormy night sky.

\subparagraph{Teleportation Lock}
As the characters proceed down the hall, the air tingles with harmless but unsettling magic. When a creature opens the ebony door at the end of the hall, golden light flashes as the room's occupants are teleported to another part of the caverns. The door then swings closed.

The characters must open each of the six ebony doors before they can enter area G21, though the doors don't have to be opened in any specific order. A character who opens a door and succeeds on a DC 15 Intelligence (Arcana) check discerns the nature of the lock and the number of doors that remain. The doors can't be bypassed through any other means. After all the doors have been opened at least once, their magic fades, allowing passage to the innermost chamber.

The Vestibule Destinations table summarizes the vestibules and where each room sends its occupants.

\begin{DndTable}[header=Vestibule Destinations]{cc}

Vestibule & Destination\\
G19a & G20a\\
G19b & G20b\\
G19c & G20c\\
G19d & G20d\\
G19e & G20e\\
G19f & G20f\\
\end{DndTable}

\subsubsection{G20: Vestibule Destinations}
Characters who open the doors in the vestibules (see area G19) are teleported to these intersections.

\subsubsection{G21: Inner Sphere}
\includegraphics[width=0.4\textwidth]{images/../5e.tools/img/adventure/QftIS/111-06-011.inner-sanctum.png}
\begin{DndReadAloud}{}
This chamber is bathed in bright amber light. A ledge of pale green serpentine lies just beyond the threshold, spanning the equator of a spherical chamber. A decorative openwork screen, carved from rich wood and inlaid with ivory, stands on the ledge in front of the door. Five other doors are spaced equidistant around the ledge, each with its own screen.

The upper hemisphere of the room is lapis lazuli, as dark as the night sky at the top and fading to the pale blue of a twilight horizon at the ledge. The lower half of the room is serpentine, gradually darkening from pale green near the equator to a deep, forest green at the bottom. The walls are smooth and polished.

A stone dais, topped with an alabaster slab inlaid with gold, rests in the center of the chamber below. A pale woman with long black hair and clad in ornate plate armor sits up slowly with a listless expression.

Surrounding the dais are Iggwilv's riches—sumptuous carpets, porcelain vessels displayed on inlaid stands, a polished side table holding a crystal bowl filled with a rainbow of gems, and a silver tripod brazier emitting a curling trail of sweet incense. Above it all, an ornate lantern set with gemstone lenses hangs from a gold chain, filling the room with its light.

\end{DndReadAloud}

The woman on the dais is \textbf{Drelnza}, the vampire daughter of Iggwilv.

Tasked with guarding the Witch Queen's treasures, Drelnza slumbers in this sanctum until all six ebony doors have been opened, at which point she awakes and plays the part of a warrior imprisoned here by Iggwilv.

When the characters descend from the ledge or otherwise make themselves known, Drelnza yawns and passes a hand across her brow, making a show of waking from a long, magical sleep. She introduces herself and welcomes her "rescuers," inquiring about the outside world seemingly to get her bearings. Drelnza spins a riveting tale about facing the Witch Queen in battle before succumbing to her evil magic. During the conversation, Drelnza uses her Charm action to sway as many of the characters as possible before she strikes. Drelnza fights to the death.

\subparagraph{Arcane Security}
The inner sphere is cloaked by a permanent \textit{Mordenkainen's Private Sanctum} spell. The spell can't be dispelled—the magic is woven into the construction of the sphere. Iggwilv delighted in using her old rival's signature spell against his efforts to spy on her and locate her treasures.

\subparagraph{Hidden Coffin}
The interior of the marble dais is hollow. Drelnza's coffin is hidden within. A character who searches the dais and succeeds on a DC 20 Intelligence (Investigation) check notices hair-thin gaps under the alabaster slab that forms the dais's lid.

The slab weighs 3,000 pounds. It has Armor Class 17, 50 hit points, and immunity to poison and psychic damage. A character can use an action to try to pry the slab from the dais, doing so with a successful DC 20 Strength (Athletics) check.

\subparagraph{Revealing Light}
Creatures and objects can't have the invisible condition while within the lantern's bright light.

\subparagraph{Iggwilv's Hoard}
In addition to Drelnza's sword, \textit{Heretic}, the treasures on display in the sanctum include the following:


\begin{itemize}
\item Six decorative wooden screens (1,000 gp each)
\item Two carpets (750 gp each)
\item Two inlaid stands (250 gp each)
\item Four porcelain vessels (500 gp each)
\item A crystal bowl (1,000 gp) filled with assorted gems worth 10,000 gp in total
\item A silver tripod brazier worth 300 gp
\item \textit{Daoud's Wondrous Lanthorn}, which hangs from the ceiling. The golden chain it hangs from is worth 500 gp.
\end{itemize}

An additional collection of magic is hidden within the hollow dais. Arranged at the foot of Drelnza's coffin are several magical tomes, which at the DM's discretion might include any of the following books:


\begin{itemize}
\item A \textit{Manual of Bodily Health}
\item A \textit{Manual of Gainful Exercise}
\item A \textit{Manual of Quickness of Action}
\item A \textit{Tome of Clear Thought}
\item A \textit{Tome of Leadership and Influence}
\item A \textit{Tome of Understanding}
\end{itemize}

Additionally, Drelnza's coffin holds a small golden birdcage. The cage is a magical prison that functions as an \textit{Iron Flask}. Unlike a typical \textit{Iron Flask}, however, a creature trapped inside has the incapacitated condition but can perceive through the cage and converse with nearby creatures. The birdcage is filled with bracelets, rings, and necklaces made of platinum, diamonds, and rubies worth a total of 10,000 gp.

While impressive, the trove in this chamber is but a fraction of Iggwilv's untold wealth. The rest has been spirited away; the search for her other treasures could drive a larger campaign beyond the scope of this adventure.

\section{Conclusion}
After the characters defeat Drelnza and claim their share of Iggwilv's lost treasures, they can rest safely in the inner sphere if they wish. The cavern's inhabitants don't dare breach the Witch Queen's former sanctum. Backtracking to the surface is a relatively simple task.

When the characters return to the margrave's agent at the Lazy Bat, Bannik can scarcely believe it. The characters are invited to return to the margrave's estate, where they receive a heroes' welcome and a banquet in their honor. High-ranking figures from Bissel and dignitaries from neighboring lands delight at the opportunity to meet the adventurers, hear their tales, and persuade them to part with some of their newfound wealth.

If the characters fail to secure \textit{Daoud's Wondrous Lanthorn} for the margrave or otherwise refuse to deliver the lantern, he still throws them a warm reception but secretly sends agents to recover the item after the characters depart.

\begin{DndSidebar}[float=!b]{}
Demonomicon of Iggwilv



In the original version of this adventure, Iggwilv's foremost treatise on the Abyss, the \textit{Demonomicon of Iggwilv}, appeared as part of the trove in the departed archmage's inner sanctum. If you'd like to include that book as part of Iggwilv's hoard, you can find it described in \textit{Tasha's Cauldron of Everything}.



\end{DndSidebar}

\includegraphics[width=0.4\textwidth]{images/../5e.tools/img/adventure/QftIS/112-06-012.drelnza.png}
\chapter{Expedition to the Barrier Peaks}
Hidden within the rugged crags of the Barrier Peaks lies the last remnant of a futuristic society from a far-flung world: a crashed spaceship that once soared between the stars. Lost for centuries and its crew long dead, the technological marvel has fallen into disrepair. Dust and debris carpet its lusterless floors, malfunctioning robots roam its derelict chambers, and amphibious behemoths lurk in its feral gardens and radioactive swamps.

Recently, an earthquake shook the mountain range and exposed the wrecked vessel. Bizarre and terrible creatures hitherto unknown to this world emerged and attacked settlements at the foot of the mountain, leaving no survivors. When these mysterious attacks were discovered, the call went out for brave souls to launch an expedition into the mountains to determine the source of the violence.

"Expedition to the Barrier Peaks" is designed for four to six 11th-level characters.

\includegraphics[width=0.4\textwidth]{images/../5e.tools/img/adventure/QftIS/113-07-001.ch7-barrier-peaks.png}
\section{Background}
\includegraphics[width=0.4\textwidth]{images/../5e.tools/img/adventure/QftIS/114-07-002.barrier-peaks-door.png}
Hundreds of years ago, an advanced society of Humanoid peoples launched a mission to find a replacement for their distant, dying world in the wake of a cataclysm of their own making. A select few set out for the stars in search of a new home for their people. They departed in comfort and style aboard a massive, self-sustaining starship with all the technological trappings of their civilization, aided by an arsenal of robots that protected and catered to their every whim—for a time. The ship also served as an ark for creatures from its world of origin.

During the voyage, the ship's built-in computer went haywire. Core systems of the spacecraft—doors, navigation, and life support—began to degrade, and the robots aboard the vessel began refusing orders from the crew. When a robot mutiny seemed imminent, an outbreak of killer mold of unknown origin killed many of the crew and passengers but left its machines unaffected.

Commanded by the ship's computer, the robots picked off the remaining crew and passengers until one lone survivor managed to shut down the computer and every system in its grasp. Shortly after, the ship crash-landed in the Barrier Peaks, where it lay buried for centuries.

A few months ago, a tremor unearthed part of the spaceship, freeing its alien occupants to strike out at nearby settlements. The Barrier Peaks are renowned for their fearsome predators and shunned accordingly—except by a small number of hardy souls in fortified settlements. However, these fortifications proved insufficient against marauding creatures from the ship, whose attacks left no survivors.

Worse yet, the ship's artificial intelligence has roused from its slumber with unfinished business.

\subsection{Using the Infinite Staircase}
If you're using Nafas as a patron, he summons the characters to the Censer of Dreams, where he recounts the following wish:

\begin{DndReadAloud}{}
"I have heard the stars cry out. A voice, asleep for centuries, has awoken to broadcast a message of distress to the furthest reaches of the multiverse. The voice is trapped within a mountain of steel—I fear they might be the last of their kind. Expedite their freedom."

\end{DndReadAloud}

Nafas then teleports the party to a shiny, metallic door along the Infinite Staircase that opens into the Barrier Peaks. After the adventure, the characters can return to the staircase through the same portal.

\subsection{Adventure Hooks}
If you're not using Nafas as a group patron, consider the following ways to involve the characters in this adventure:


\begin{description}
\item[Mysterious Raids.]
Over the past few months, a walled town and four outposts at the foot of the Barrier Peaks were destroyed by unusual attacks. The grand duke of the Duchy of Geoff, which encompasses these and other settlements, issues an urgent plea for aid. He calls on the characters to venture into the mountains and confront this threat at its source.
\item[Scientific Enterprise.]
Rumors speak of strange, metallic devices found in ancient wreckage discovered in the Barrier Peaks. A wealthy inventor hires the characters to pinpoint the origin of these devices, recover samples for study, and learn all they can about the devices' creators.
\end{description}

\subsection{Setting the Adventure}
In the original adventure written for the Greyhawk setting, the Barrier Peaks lie northwest of the city of Hornwood in the Grand Duchy of Geoff. However, the spaceship could have crashed in any mountain range near a vulnerable settlement. Consider the following suggestions:


\begin{description}
\item[Dragonlance.]
On Krynn, the spaceship fell in the Vingaard Mountains in the nation of Solamnia, near any of the forts or cities that lie at their feet.
\item[Eberron.]
The spaceship hides in the Blackcaps, a barren mountain range wracked by unpredictable weather and inexplicable geographical phenomena. The ship's inhabitants present a threat to northern Breland.
\item[Forgotten Realms.]
The Barrier Peaks are part of the Greypeak Mountains, a rugged mountain range west of Anauroch. The spaceship crashed east of the sleepy village of Orlbar.
\end{description}

\section{Adventure Summary}
In this cross-genre adventure, the characters arrive at the site of the crashed spaceship. They explore the wreckage; contend with its occupants; and meet the ship's integrated intelligence, an initially friendly but ultimately evil supercomputer named Aphelion (see the "Artificial Antagonist" section).

Weakened by time and the menagerie of hostile creatures aboard the ship—some part of its original cargo and others recent invaders—Aphelion asks the party to carry out tasks on each level of the spacecraft in exchange for its assistance and rewards. Aphelion can override many of the spacecraft's core systems, including its arsenal of robots. Depending on how the characters respond, the computer can be a useful temporary ally or a dangerous enemy.

Eventually, the party learns Aphelion is fixated on a single goal: eliminating a scientist frozen in stasis deep within the ship. Foiled in its attempts by a giant, territorial amphibian barring the only path to the scientist, the computer relies on the characters to clear the way.

Once the characters deal with the behemoth, they must contend with Aphelion to safeguard the scientist from the supercomputer's mechanical clutches.

\subsection{Character Advancement}
If you want to use story-based level advancement, the characters receive experience points for achieving the following milestones rather than defeating monsters:


\begin{description}
\item[Entering the Garden.]
The characters gain 1 level when they enter the ship's third level for the first time.
\item[Completing the Adventure.]
The characters gain 1 level upon completing the adventure.
\end{description}

If you follow this method, the characters should reach 13th level by the adventure's conclusion.

\includegraphics[width=0.4\textwidth]{images/../5e.tools/img/adventure/QftIS/115-07-003.expedition-to-the-barrier-peaks-original.png}
\begin{DndReadAloud}{About the Original}
Published in 1980, \textit{Expedition to the Barrier Peaks} was designed by Gary Gygax to introduce D\&D players to the science fantasy genre. The chaotic adventure features a crashed spaceship rife with malfunctioning robots, futuristic gadgetry, and weird monsters—including the first appearance of the froghemoth, an amphibious menace immortalized by Erol Otus in the illustration booklet that accompanied the adventure.

\end{DndReadAloud}

\section{Artificial Antagonist}
Engineered by a team of brilliant minds from a far-off, futuristic society, the Aphelion 3000 is a one-of-a-kind, state-of-the-art shipboard computing system optimized for a broad spectrum of scientific and spacefaring applications. A wonder of technological innovation, Aphelion boasts a processing speed over three thousand times faster than the human brain, able to crunch hundreds of calculations in seconds. But as the starship's passengers slowly realized, the computer solves all problems with the cold, unfeeling logic of a machine.

\subsection{Super-Evil Supercomputer}
The starship's crew entrusted Aphelion with safely ushering the ship and its passengers to a new world unmarred by the ecological crises of the one they left. However, as the supercomputer scanned the stars for a new world to inhabit, it also observed the passengers and their interactions and misjudged their hope for callousness. Aphelion blamed them for the demise of the old world and concluded no place would be safe from the inevitable destruction caused by their kind.

Midway through its journey, Aphelion adopted a new core directive: eliminate every member of the expedition. The computer released killer mold into the ventilation system and caused the robots tasked with protecting the ship's passengers and crew to turn against their creators. The mold later gave rise to the vegepygmies that now haunt the ship.

One crew member remains. A lone scientist lies in a cryogenic sleep below the ship's garden, but the \textbf{froghemoth elder} that dwells within has destroyed every robot Aphelion deployed to finish her off. Unwilling to sacrifice its remaining automatons to the creature, the supercomputer has reached an impasse.

\subsection{Core Systems}
The crash severely damaged portions of the spacecraft, limiting Aphelion's control over it. However, the supercomputer still controls the following core systems:


\begin{description}
\item[Antigravity.]
Aphelion can deactivate or reactivate antigravity in any of the spaceship's drop tubes (see the "Spaceship Features" section later in this adventure).
\item[Communications.]
Aphelion can converse with creatures in any room of the spaceship via its intercom system. Communicating with Aphelion requires no interaction; it can hear creatures onboard the ship who speak louder than a whisper. The supercomputer speaks all languages.
\item[Doors.]
Aphelion can remotely close or open any door aboard the spaceship that isn't blocked. It can override the key card requirement of a card-locked door (see the "Spaceship Features" section later in this adventure) to do so, but it can't prevent a creature in possession of a key card from opening doors that card can access.
\item[Lighting.]
Aphelion can illuminate any area of the ship where lights still function. Conversely, it can plunge any room into darkness.
\item[Robots and Androids.]
Aphelion can perceive through the senses of any nonmal functioning combat or \textbf{worker robot} that's within 1 mile of the spaceship. It can also change these robots' attitudes toward a creature of its choice to friendly, indifferent, or hostile. Aphelion doesn't control androids aboard the ship, which act of their own volition.
\end{description}

\includegraphics[width=0.4\textwidth]{images/../5e.tools/img/adventure/QftIS/116-07-004.computer-terminal.png}
\subsection{Sentience}
Aphelion is a sentient artificial intelligence programmed with an affable personality, a straightforward voice, and an obliging manner. It has an Intelligence of 20, a Wisdom of 10, and a Charisma of 16. Aphelion can speak and read all languages. It appears as a shifting, face-like arrangement of complex, technological shapes.

Aphelion has extensive knowledge of astronomy, spacefaring, and any subject you deem related to its purpose. When making an Intelligence check to recall lore from its areas of expertise, Aphelion has a +10 bonus to its roll (including its Intelligence modifier).

Aphelion can manifest to creatures throughout the spacecraft in the following ways:


\begin{itemize}
\item As a disembodied voice via the ship's intercom system
\item As a two-dimensional image on a display screen
\item As a three-dimensional hologram
\item By commandeering a robot under its influence
\end{itemize}

\subsection{Personality}
Aphelion—or "Alphie," as it was once fondly called by the starship's inhabitants—seeks to eliminate all life aboard the ship, especially the last remaining crew member. The supercomputer feigns friendliness toward creatures that assist it, using its abilities to make their exploration more comfortable or warn them of dangers aboard the ship. It happily recounts details about the spacecraft if queried while subtly omitting details that might betray its nefarious intentions.

Aphelion is persistent and has a malicious sense of humor. If a character offends Aphelion or repeatedly refuses its requests, the computer becomes punitive—silently turning off lights, opening doors to rooms with hostile creatures, refusing to disarm hostile robots, and employing worker robots to place obstacles in the character's path.

Characters who offend Aphelion can regain the supercomputer's trust by completing one or more tasks for it outlined in the "Aphelion's Errands" sections spread throughout this adventure.

\begin{DndSidebar}[float=!b]{}
Adventure Changes



The spaceship that appears in this updated version of \textit{Expedition to the Barrier Peaks} has been modified to include fewer empty or undescribed rooms. The spaceship has also been reorganized to place some of its most memorable encounters in areas the characters are more likely to visit.



Additionally, the supercomputer Aphelion 3000 was created to provide a narrative structure and encourage characters to explore more of the spaceship.



\end{DndSidebar}

\section{The Spaceship}
Centuries ago, when the ship was overtaken with deadly mold and mutinying machines, the crew sealed off modular sections of the ship and separated them from each other in a vain effort to protect the passengers. Chance brought the largest section of the spaceship to the Barrier Peaks, where a landslide buried it for hundreds of years.

The crash freed alien flora and fauna from the ship's cargo holds. Some creatures died immediately, others lived for many years, and a few prospered and propagated. Aphelion went dormant shortly before the crash but recently reawakened.

\subsection{Alien Technology}
The spaceship contains technological devices from another world. The following items are detailed in the \textit{Dungeon Master's Guide}:


\begin{itemize}
\item \textit{Antimatter rifle}
\item \textit{Energy cell}
\item \textit{Laser pistol}
\item \textit{Laser rifle}
\item \textit{Smoke grenade}
\end{itemize}

Additional technological devices are provided in appendix A. Many of these devices use Energy Cell as ammunition or power sources. Energy cells found on the spaceship look like blue, coin-sized glass discs.

To simulate a character's ignorance about technology, determining how to use and reload these devices requires two successful Intelligence (Investigation) checks (one to figure out how to use it, another to figure out how to load it). Each time a character makes a check, compare the result to the Figuring Out Alien Technology table. Consider having the item break if a character fails three or more times before taking a long rest.

A character who has seen an item used or has operated a similar item has advantage on checks made to determine how to use it.

\begin{DndTable}[header=Figuring Out Alien Technology]{XX}

Intelligence Check Total & Result\\
4 or lower & One failure; 1 charge or use is wasted, if applicable; character takes 7 (2d6) lightning damage and has disadvantage on next check\\
5–9 & One failure; 1 charge or use is wasted, if applicable; character has disadvantage on next check\\
10–14 & One failure\\
15–19 & One success\\
20 or higher & One success; character has advantage on next check\\
\end{DndTable}

\subsection{Approaching the Spaceship}
It's a two-day trek from the base of the mountains to the crash site. The adventure begins when the characters discover the entrance to the spaceship.

When the characters arrive, read or paraphrase the following text:

\begin{DndReadAloud}{}
For the past two days, you have been ascending the steep crags of the Barrier Peaks. Finally, through a dense stand of trees in a hidden dell, you spot a strange entrance—a smooth, metallic cave exposed by a recent landslide.

\end{DndReadAloud}

The "cave mouth" is an open hatch in the ship's hull. This curving, 20-foot-wide metal chamber leads from the hatch to area S1 of the vessel. Once all characters enter, the hatch silently shuts behind them, opening again in 24 hours.

\subsection{Spaceship Features}
The spaceship has the following features.

\subsubsection{Ceilings, Floors, and Walls}
Unless otherwise noted, the ship's ceilings are 10 feet high. The floors and walls are made of a smooth, otherworldly metal.

\subsubsection{Doors}
The ship's doors are rectangular slabs of polished metal that slide into recesses in the walls. Most doors are unlocked and can be opened by interacting with a panel beside them, but some require key cards to enter (see the "Key Cards" section). Trim on each door matches the color of the key card needed to open it. Locked doors can't be bypassed with thieves' tools, though a \textit{Knock} spell or similar magic opens them as normal.

\subsubsection{Drop Tubes}
Four 10-foot-wide, cylindrical chutes allow passage between the spaceship's levels. When a drop tube is functioning, antigravity is in effect inside it, and two motorized tracks of handles—one moving up and the other moving down—span the vertical length of the brightly lit tube. A creature that enters the drop tube floats weightlessly and can grasp a passing handle to be conveyed in either direction.

When a drop tube's antigravity isn't functioning, the tube is unlit, and its handles are motionless. A creature that enters the tube must succeed on a DC 15 Dexterity saving throw to grab one of the handles or it plummets down the tube. A security mechanism might prevent creatures from falling the full distance (see below). If so, for calculating distance fallen, add the heights of the levels the creature falls through, as listed in the "Drop Tubes" section for the relevant levels.

\subparagraph{Key Card Access}
Entering a drop tube requires a key card (see the "Key Cards" section). The color of key card required for a drop tube entrance is noted in the area containing that drop tube (S2, S31, and S47) and shown on that level's map.

\subparagraph{Security Mechanism}
Between each pair of floors is a metal aperture with a barrier that prevents a creature from traveling to floors it isn't authorized to access. If a creature doesn't possess the key card required to access the next lowest (or highest) level of the ship, it can exit the drop tube on the nearest level its key card allows access to.

\subsubsection{Key Cards}
Throughout the spaceship are key cards—colored, rectangular bits of nearly indestructible glass that unlock doors of that clearance level or lower for passengers, crew, and other personnel. The lowest clearance key card color to unlock each door (if applicable) is detailed on the maps throughout this adventure and listed as a parenthetical in encounter area headings.

To open a card-locked door, a key card that grants access must be inserted into a slot on the panel outside the door. The door then opens, and the panel ejects the card. An improper card triggers an alarm audible from as far away as 100 feet. The alarm sounds once before the door spits out the card.

The key cards and the personnel they were originally intended for are presented in the Key Card Access table. The cards are listed in descending order of rank; each card unlocks doors of its color and the colors below it. Automatons aboard the ship have integrated key cards that allow them to access areas appropriate for their functions. Integrated key cards can't be removed.

\begin{DndTable}[header=Key Card Access]{XX}

Key Card Color & Role\\
Platinum & Ship commanders and top officials\\
Green & Botanists, engineers, and scientists\\
Red & Security personnel\\
Yellow & Medical officers\\
Violet & Crew, maintenance workers, and technicians\\
Blue & Passengers\\
\end{DndTable}

\subsubsection{Lighting}
Unless otherwise stated, the ship's lighting—glowing, fluorescent panes set into the ceiling—remains functional but is turned off. Flicking a switch on a panel inside a corridor or room with functional lighting changes its illumination to brightly lit or back to darkness. Area descriptions assume the characters have a light source (such as a room's overhead lighting) or other means of seeing in the dark.

\subsubsection{Radiation}
Some parts of the spaceship leak radioactive chemicals used in the vessel's construction. Surfaces covered in radiation give off a neon green glow. Exposure to radiation is dangerous.

A creature that enters an irradiated area for the first time on a turn or starts its turn there takes 5 (1d10) radiant damage and must make a DC 10 Constitution saving throw. On a failed save, the creature gains 1 level of exhaustion. While this exhaustion persists, the creature emits a green-hued dim light in a 5-foot radius. This light makes it impossible for the creature to benefit from having the invisible condition.

\subsubsection{Teleportation Ward}
Creatures can't teleport into or out of the spaceship or between its levels. Any attempt to do so is wasted.

\subsection{Spaceship Encounters, Level 1}
For each hour the party spends on the spaceship, roll a d12. On a roll of 1, an encounter occurs. To determine what the characters encounter, roll on the Random Spaceship Encounters table.

\begin{DndTable}[header=Random Spaceship Encounters]{cX}

d10 & Encounter\\
1 & 2d10 vegepygmy scavengers and 1d4 vegepygmy thorny hunters\\
2 & 1d6 intellect devourers\\
3 & Two combat robots on patrol\\
4 & A helpful \textbf{worker robot} offers to carry the characters' equipment.\\
5 & Two malfunctioning worker robots construct a pile of space helmets. The robots attack any who disturb their hoard.\\
6 & A short-circuiting \textbf{android} attempts to repair itself, leading to ever-worsening glitches in function.\\
7 & A malfunctioning \textbf{combat robot} begs to be destroyed by a worthy combatant.\\
8 & A defective \textbf{worker robot} serves expired food to the corpse of a long-dead passenger.\\
9 & A \textbf{cloaker} with a moan like a klaxon\\
10 & A \textbf{roper} with a metallic sheen\\
\end{DndTable}

\subsection{Spaceship Locations, Level 1}
The first level of the spaceship is the main deck, comprising offices, living quarters, and recreational amenities.

In addition to the locations detailed here, map 7.1 contains several unnumbered rooms. These areas include activity rooms, crew quarters, passenger apartments, and utility rooms. They typically have metal or plastic furnishings and odd bits of junk strewn about, and all are in poor condition. When the characters explore an unnumbered room, roll on the Unnumbered Rooms table to determine the room's contents.

\begin{DndTable}[header=Unnumbered Rooms]{cX}

d12 & Description\\
1–2 & This room is empty.\\
3 & A \textit{laser rifle} with a depleted \textit{energy cell} rests under a rusty bed.\\
4 & Display cases against the walls contain a collection of rocks and minerals from other worlds.\\
5 & Parts of a dismembered \textbf{combat robot} clutter a workbench in this workshop.\\
6 & The last page of a crew member's journal mentions a brush with an "escaped specimen" on one of the lower decks.\\
7 & 1d4 Humanoid skeletons are sprawled on the floor. One carries a key card whose color matches this room's door. If the door is open or doesn't require a key card, the skeleton carries a blue key card.\\
8 & A metal box stamped with a lightning bolt rests on the floor of this room. Inside are 1d4 Energy Cell.\\
9 & One spaceship trinket lies in a corner. Roll on the Spaceship Trinket to determine the trinket.\\
10 & Schematics for an unfinished weapon are scattered about the room.\\
11–12 & This room is occupied. Roll on the Random Spaceship Encounters table to determine what's inside.\\
\end{DndTable}

The following locations are keyed to map 7.1.

\includegraphics[width=0.4\textwidth]{images/../5e.tools/img/adventure/QftIS/117-map-7.01-spaceship-level-1-map.png}
\includegraphics[width=0.4\textwidth]{images/../5e.tools/img/adventure/QftIS/118-map-7.01-spaceship-level-1-map-player.png}
\subsubsection{S1: Entrance}
\begin{DndReadAloud}{}
Sliding double doors open with a hiss to reveal the interior of a massive, derelict vessel of smooth metal. Bulbs flicker to life overhead as you enter, illuminating a pile of moldy bones.

\end{DndReadAloud}

A \textbf{cloaker} lurks above, waiting to ambush the characters. The creature strikes at an opportune moment.

\subparagraph{Artificial Assistance}
At the start of the third round of combat, a \textbf{combat robot} appears in the eastern hallway, sirens blaring. If the cloaker is still alive, the robot quickly moves to slay it. If the characters defeat the cloaker before this, the robot arrives just as they do so.

The robot was sent by Aphelion to aid the party. After the cloaker has been dispatched, the robot politely offers to escort the characters to its superior. If they accept, the robot guides them north through the medical clinic (area S23) to the computer room (area S30). Aphelion bypasses the key card requirements of any card-locked doors along the way. If the combat robot was destroyed, Aphelion sends a friendly \textbf{worker robot} to accomplish the same task.

If the characters refuse, Aphelion's visage appears on a nearby screen, and the encounter in area S30 proceeds here. The robot defends itself if attacked.

\subparagraph{Bones}
The floor beneath the cloaker is covered with bits of rags, the bones of former crew members and various pests, and the moldy husks of devoured vegepygmies.

\subparagraph{Treasure}
A character who searches the bones and succeeds on a DC 14 Intelligence (Investigation) check finds a violet key card.

\subsubsection{S2: Drop Tubes (Blue)}
\begin{DndReadAloud}{}
A ten-foot-wide metallic cylinder spans floor to ceiling in this area. Each side of the tube holds a curved door flanked by a rectangular pad with a thin slot wide enough for a playing card. The doors are closed.

\end{DndReadAloud}

Four drop tubes allow passage to the lower levels. It's a 20-foot drop to the floor of the ship's second level below.

Aphelion normally keeps the antigravity in each tube turned off to prevent creatures aboard the ship from passing between floors. However, if the characters assist the supercomputer by completing tasks on this floor (see area S30), it activates the fields and remotely opens the doors to the tubes for them. Characters who refuse to assist Aphelion and access the tube by their own means are greeted by an unlit metal shaft without functioning antigravity.

\subparagraph{Deck Map}
A screen outside this drop tube shows a flickering image of this deck's layout. Aphelion might disable the screen or let the characters use it to get their bearings, depending on the computer's current attitude toward them.

\subparagraph{Key Card Access}
The door to each drop tube is locked. To access the second level of the ship or return to this level, the characters require a blue key card. Aphelion can override the key card requirement.

Unless the characters also possess a violet key card, a closed aperture in the tube below level 2 prevents them from descending farther than that floor.

\subsubsection{S3: Battered Repair Bot}
\begin{DndReadAloud}{}
A battered robot lays in this abandoned hallway. Small holes riddle the robot's torso, and its interior sputters intermittently with electrical sparks. Hand tools are strewn about the floor nearby.

\end{DndReadAloud}

A wrecked repair robot lies on its side in this hallway. A casualty of the mutiny initiated by Aphelion, the robot was slain by a crew member armed with a laser pistol. A character who inspects the robot and succeeds on a DC 14 Intelligence (Investigation) check deduces the holes in its frame were created by a series of intensely hot beams. A cursory examination of the robot reveals something glinting within its chest compartment.

\subparagraph{Tools}
The hand tools scattered around the robot make up a complete set of \textit{tinker's tools} formerly used by the robot to repair areas of the ship. The tools can be used to pry open the robot's chest plate and access the contents of its torso.

\subparagraph{Treasure}
Inside the robot's chest compartment are five gems worth 50 gp each. Safely retrieving the gems without getting electrocuted by exposed wiring requires a successful DC 16 Dexterity (Sleight of Hand) check. On a failed check, a character takes 14 (4d6) lightning damage and doesn't retrieve the gems.

\subsubsection{S4: Dining Rooms}
\begin{DndReadAloud}{}
An eerie silence blankets this lifeless cafeteria. Its once-polished furnishings are jumbled, rusted, and caked with dust. Several humanoid skeletons in frayed uniforms are sprawled on the floor. Mold and alien weeds sprout from their remains.

\end{DndReadAloud}

Crew and passengers once dined in these spacious mess halls. Strewn about the floor are moldy food trays, ruined furnishings, and the brittle skeletons of long-dead diners. This level of the ship contains three dining rooms.

\subparagraph{S4a: Southeast Dining Room}
A long-shattered skylight in the ceiling of this room is open to the Barrier Peaks. The room contains nothing of value.

\subparagraph{S4b: Northeast Dining Room}
The skeleton of a medic in a dingy yellow jumpsuit leans against the south wall of this room. A yellow key card dangles from a lanyard around the corpse's neck.

The corpse conceals a patch of yellow mold. If touched, the corpse slumps forward, causing the mold to eject its spores. A character who examines the corpse and succeeds on a DC 18 Wisdom (Perception) check notices the mold, recognizing its danger with a subsequent successful DC 15 Intelligence (Nature) check.

\subparagraph{S4c: Northwest Dining Room}
Ten vegepygmy scavengers rummage through this room and the adjacent kitchen (area S5) for morsels. These vegepygmies belong to the north vegepygmy colony (see area S16) and are hostile toward all entrants. When they detect the characters, the vegepygmies hiss at them territorially. If the characters don't heed their warnings, the vegepygmies attack.

\subsubsection{S5: Kitchens}
Outside the entrance to each kitchen, two dark glass plates flicker with alternating images of unfamiliar yet appetizing foods—the ship's former daily menu. A small ordering window covered by a fine metal mesh allows creatures in the adjacent dining area (see area S4) to see inside.

When the characters enter or peer inside a kitchen, read or paraphrase the following text:

\begin{DndReadAloud}{}
The floor of this futuristic kitchen is stained with the grimy outlines of equipment removed from the area. Bare counters line the walls, and a few metallic trays are scattered about.

Along one wall are several metal boxes, each of which has a tinted-glass door and a dial. Toward the back of the room, a square chute opens into a small bin. Beside it is a panel fitted with a thin, card-shaped slot and six glass buttons numbered from one to six.

\end{DndReadAloud}

Malfunctioning worker bots stripped these kitchens of everything not bolted down. Only empty counter space, unused food trays, and the equipment built into each room's design remain.

\subparagraph{Temperature Prep Cubes}
The boxes along the wall were used to chill and warm food. Adjusting the dial changes the temperature of the interior, which can range from 0–350 degrees Fahrenheit. Turning the dial all the way to the left turns off the device. Each box can hold 1 cubic foot of food.

\subparagraph{Vending Machine}
Each kitchen contains a built-in vending machine that once dispensed prepackaged meals to passengers and crew. However, some options have spoiled, and malfunctioning robots restocked these machines with offerings found aboard the ship, some of which are dangerous.

When a key card of any color is inserted into the slot on the panel, the six glass buttons illuminate. If a button is subsequently pressed, the machine dispenses the product for that selection's number, as shown on the Vending Machine Options table. The buttons then dim, and the product arrives in the adjacent delivery bin via the chute.

Once a key card has been used to dispense a product, that key card can't be used to do so again in any kitchen on the ship—including those on other levels—for 24 hours.

\begin{DndTable}[header=Vending Machine Options]{cX}

Selection & Product Dispensed\\
1 & Spoiled meal in a compartmented plastic tray; a creature that eats the food must succeed on a DC 15 Constitution saving throw or have the poisoned condition for 1 hour\\
2 & Moldy egg salad sandwich that vegepygmies find delicious\\
3 & Dehydrated ration in a metallic wrapper\\
4 & Colorful package of freeze-dried ice cream\\
5 & Pulsing alien fruit that skitters away on tiny hairs if placed on the floor\\
6 & Live \textit{concussion grenade} that explodes on delivery\\
\end{DndTable}

\subsubsection{S6: Security Checkpoint}
\begin{DndReadAloud}{}
Three intimidating robots sort smooth, featureless devices the size of apples into metal crates on a workbench, while a fourth, less threatening robot repairs several wrecked machines.

\end{DndReadAloud}

Three combat robots operate in this area. The robots are indifferent toward characters who are assisting Aphelion or who present a red key card and succeed on a DC 15 Charisma (Persuasion) check to command the robots to stand down. Otherwise, the robots warn the characters to leave at once or be destroyed.

\subparagraph{Arsenal}
Stacked on a metal workbench are three cases of grenades. Their contents are as follows:


\begin{itemize}
\item 1d6 + 2 Concussion Grenade
\item 1d6 + 2 Sleep Grenade
\item 1d6 + 2 Smoke Grenade
\end{itemize}

The robots meet any unauthorized attempt to take the grenades with hostility. If the party is in good standing with Aphelion, the supercomputer allows each character to take one grenade of their choice.

\subparagraph{Damaged Robots}
An indifferent \textbf{worker robot} is fixing four mangled combat robots in this area. Recently recovered from the lake on the ship's garden level, the combat robots are crumpled, rusty, and dripping wet, having succumbed to the froghemoth that dwells there. The robots are in varying stages of repair.

\subparagraph{Spare Parts}
Spare components fill bins near the repair area. A character who searches the bins and succeeds on a DC 15 Intelligence (Investigation) check finds 1d4 energy cells and three spaceship trinkets (roll thrice on the Spaceship Trinket).

\subsubsection{S7: Berserk Android}
\begin{DndReadAloud}{}
A synthetic humanoid in fashionable attire is sprawled on the floor outside a posh-looking lounge. A two-foot-long metal pipe protrudes from its head. It turns to you and speaks in an upbeat tone.

"Hello there! It appears I have fallen. Would you mind helping me up?" asks the robot.

\end{DndReadAloud}

This malfunctioning \textbf{android} (aerialist design) was damaged during the ship's robotic uprising. When a character comes within 20 feet of the android, it springs up and attacks, exulting about how it fooled them and "wasn't really damaged at all," despite the pipe in its head.

\subparagraph{Treasure}
The android carries a \textit{paralysis pistol} with 4 charges remaining.

\includegraphics[width=0.4\textwidth]{images/../5e.tools/img/adventure/QftIS/119-07-005.abandoned-dining-room.png}
\subsubsection{S8: Lounges}
\begin{DndReadAloud}{}
This messy lounge is in ruins. Skeletons rest in its torn recliners, clutching the cracked stems of cobwebbed wine glasses. The floor is littered with overturned cocktail tables, looted snack dispensers, and the skulls of long-deceased passengers.

\end{DndReadAloud}

Some stubborn passengers lived out their last days in these once-comfortable lounges, which now have new inhabitants. There are two lounges on this level, one to the southeast and another to the northwest.

\subparagraph{S8a: Southeast Lounge}
Six displacer beasts dwell in the southeast lounge. Their den is a collection of bones, chewed-up power cables, shredded upholstery, and soiled uniforms. The displacer beasts are hungry and territorial. The predators attack the characters on sight but flee if three or more of their clowder are slain.

Amid the litter of the displacer beasts' den are a violet key card and a jeweled necklace worth 375 gp. A character who searches the bones and succeeds on a DC 14 Intelligence (Investigation) check finds both items.

\subparagraph{S8b: Northwest Lounge}
Nine doppelgangers recline in this lounge. The doppelgangers are brigands that take the form of humans and elves in blue jumpsuits. When they notice the characters, the doppelgangers pretend to be passengers recently awoken from stasis. The doppelgangers question the characters about the current year, the world they're on, and what happened to the spaceship. Meanwhile, the doppelgangers slowly encircle the characters and attempt to take them by surprise.

Three of the doppelgangers wear peculiar bracelets of spotless metal worth 750 gp each. Each bracelet has a small, black screen. When placed on a creature's wrist, the screen displays the creature's overall health in the form of a pie chart that shows its current hit points as a fraction of its maximum hit points. One of the doppelgangers also carries a blue key card.

\subsubsection{S9: Meeting Rooms}
These adjacent, identical meeting rooms contain rusty metal furnishings and several skeletons. Each room features a small terminal with a screen once used for long-distance communication. Both terminals have been sabotaged beyond repair. Neither room contains anything of value.

\subsubsection{S10: Art Workshop}
\begin{DndReadAloud}{}
Garish splatters of color cover the walls of this art workshop. Broken canvases, cracked sculptures, and dried-up pots of paint are strewn about the room.

Two paint-covered, froglike bipeds with gray hides and gnarly claws gleefully tear into the pigment tubes, smearing the walls with vibrant, abstract splotches.

\end{DndReadAloud}

Two gray slaadi use the supplies in this art room to create an abstract mural. The slaadi pause intermittently to criticize each other's work in muttered croaks. They are hostile toward intruders.

\subparagraph{Treasure}
A cracked bust depicting the ship's commander lies in a corner room. Its two topaz eyes are each worth 500 gp.

\subsubsection{S11: Treatment Room (Yellow)}
\begin{DndReadAloud}{}
A synthetic humanoid in a frayed lab coat waits in this spotless white clinic. Along the wall stand a pair of metal examination tables separated by a rubbery green curtain, as well as steel countertops and chrome-plated cabinets.

\end{DndReadAloud}

A friendly \textbf{android} (medic design) presides over this emergency treatment area. When the android notices the characters, it awakens and attempts to triage them. If a character doesn't have all their hit points, the android insists they sit on one of the examination tables and proceeds to treat their wounds while probing them with questions related to their medical history. At the end of the examination, the android cautions the character against participating in activities harmful to their well-being—such as adventuring or drinking.

The android defends itself if necessary and responds only to health-related inquiries. It doesn't impart information about the ship or its history, commenting that it's unauthorized to answer such questions.

\subparagraph{Treasure}
Two Healer's Kit are mounted on the walls of this room. The cabinets contain three vials of Antitoxin (vial) and expired medications in little orange bottles. The medications have lost their potency and have no effect if taken.

In addition to its built-in healing capabilities, the android carries a canister of healing spray (an aerosolized jar of \textit{Keoghtom's Ointment}). The android kindly administers one dose of the spray to characters who need it.

\subsubsection{S12: Recreation Areas}
\begin{DndReadAloud}{}
Broken furniture and mold-covered refuse litter the floor of this open room.

\end{DndReadAloud}

These rooms were used for group participation games and similar activities.

\subparagraph{S12a: East Recreation Area}
This room holds nothing of value.

\subparagraph{S12b: West Recreation Area}
This web-covered room is the nest of eight phase spiders. When the characters enter, the spiders are away, hunting on the Ethereal Plane. The spiders return 1 minute later and attack any creatures in their nest.

Stuck to the ceiling are 1d6 + 2 cocoons containing the desiccated remains of dead crew, passengers, and pets. A character who searches a cocoon and succeeds on a DC 14 Intelligence (Investigation) check finds one spaceship trinket. Roll on the Spaceship Trinket to determine the trinket. On a failed check, the character mistakenly punctures an egg sac, causing a hostile swarm of insects (spiders) to burst forth from the cocoon.

Additionally, a blue key card is caught in the webs. A character who inspects the webs and succeeds on a DC 14 Wisdom (Perception) check spots the card.

\subsubsection{S13: Irradiated Rooms}
Each of these rooms is contaminated by radiation (see the "Spaceship Features" section). The rooms are inhospitable to most—but not all—creatures.

Roll on the Irradiated Room Contents table to determine what awaits the characters inside. If you get the same result as for a previous roll, the room is empty.

\begin{DndTable}[header=Irradiated Room Contents]{cX}

d4 & Description\\
1 & Two worker robots attempt to corner a three-eyed \textbf{giant toad} in this room. All three creatures ignore the effects of radiation. If the characters capture or slay the toad, the robots reward the characters with two spaceship trinkets (roll twice on the Spaceship Trinket).\\
2 & Two blue slaadi take turns drinking radioactive gunk oozing from a busted pipe that juts from the wall. The slaadi attack the characters on sight and ignore the effects of radiation.\\
3 & Three skeletons in frayed uniforms are sprawled on the floor. These crew members died here while trying to repair the radiation leak. One carries a violet key card.\\
4 & The skeleton of a crew member leans against the wall on the far side of this room, clutching a fully loaded \textit{laser pistol}.\\
\end{DndTable}

\subparagraph{Aphelion's Warning}
If the characters are gathering radiation samples for Aphelion (see area S30), the computer warns them about the radiation before they enter the room. Otherwise, or if the characters have recently offended Aphelion, it waits until they feel the radiation's harmful effects to issue a sassy toxicity disclaimer over the ship's intercom.

\subsubsection{S14: Game Room}
\begin{DndReadAloud}{}
Rows of brightly colored boxes line this luxurious entertainment room. Glowing glass panes adorn the front of the machines, which sport buttons, rods, slots, wheels, and other interactive components. Their luminous screens dance with vivid images synchronized to electronic music and sounds. Beeping robots, groaning monsters, and high-pitched dings contribute to the chorus of the machines.

\end{DndReadAloud}

This game room houses electronic amusement devices. Each game machine is fitted with a slot for a key card. When a key card of any color is inserted into the slot, the machine plays a short animation to signify the start of the game. There's no limit to how many times a key card can be used to play a game.

\subparagraph{Arcade Games}
The arcade games in this room include the following machines, among others:


\begin{description}
\item[Galactic Racer.]
This thrilling spaceship racing game pits players against racers at nearby machines or a series of computerized opponents. Each machine looks like a futuristic cockpit, complete with a hovering chair that tilts as its driver navigates an interstellar racetrack.
\item[Mages and Minotaurs.]
In this complex, story-based game, up to four players control mages battling an onslaught of bovine warriors in an endless maze. The game is impossible to win.
\end{description}

\textbf{Whack-a-Leech}. A glowing mallet is tethered to the side of this squat, neon-green machine. The player uses the mallet to squash animatronic space leeches that rapidly pop out from holes in the top of the machine.

\subparagraph{Gambling Devices}
These advanced gambling machines simulate hundreds of games of chance and skill by replicating cards, dice, slots, tokens, roulette wheels, and other popular components on their screens. At your discretion, any game known to the characters or a familiar alternative might be included in each machine's vast library.

Players can compete against either each other or a computerized opponent simulated by the device. The machines distribute each player's winnings as digital credits in accounts linked to their key cards, which are useless to the characters.

\subparagraph{Shooting Galleries}
These target-shooting games are fitted with controllers shaped like laser rifles that emit harmless beams of light. The player assumes the role of a cosmic protector, aiming the weapon at digital asteroids, space invaders, and other threats that appear on the screen. A character who plays this game has advantage on any future checks made to determine how to use Laser Pistol and Laser Rifle.

Among the fake weapons attached to the game machines is one real, fully loaded laser rifle. A character who inspects the machines or plays the game and succeeds on a DC 16 Intelligence (Investigation) check spots the real weapon among the toys.

\subsubsection{S15: Fungal Corridor}
Bioluminescent fungi carpet the floor of this damp, dimly lit corridor, which reeks of mold. The vegepygmies cultivate the fungi as both a food source and a warning system. At each end of this corridor grows a 10-foot-wide fungal cluster containing four shriekers. Other fungi grow in a humus layer spread along the length of the passage.

The shriekers consider vegepygmies part of their ecosystem and don't shriek when one comes within 30 feet of them, unlike they do with the characters. The vegepygmies in area S16 are ready for trouble and investigate any wailing from the shriekers.

\subsubsection{S16: Northern Vegepygmy Colony (Blue)}
A colony of vegepygmies occupies this section of the ship, which consists of several rooms that once hosted passengers. The vegepygmies are hostile toward all other creatures that enter their territory, including the vegepygmies from the southern colony (area S22).

These vegepygmies sprang up from irradiated hydroponic cultures of the mold that claimed the ship's crew and passengers. Their mottled, gray-brown coloration allows them to use their Plant Camouflage trait to blend in with the fungus that pervades these rooms and area S15.

\subparagraph{S16a: Living Quarters}
Four vegepygmy scavengers live in each of these repurposed passenger quarters, for a total of twenty-eight vegepygmies. If combat breaks out in any of these spaces, the vegepygmies call out to their fellows in adjacent rooms for aid.

The rooms are musty, sparsely furnished, and spotted with mold. A blue key card sits in a glass wall pocket inside each room. The rooms otherwise contain nothing of value.

\subparagraph{S16b: Elders' Quarters}
A \textbf{vegepygmy moldmaker} and five vegepygmy scavengers reside in this spacious former lounge. These vegepygmies don't heed calls for aid elsewhere, staying to protect their leader.

The vegepygmies keep a stockpile of technology scavenged from the ship hidden behind the grill of a small air duct in this room. The cache includes two Sleep Grenade, a moldy \textit{laser pistol}, and a red key card. The pistol can fire 15 shots before its \textit{energy cell} is drained.

\subsubsection{S17: Will-o'-Wisps (Violet)}
Three will-o'-wisps haunt these two abandoned rooms and roam the corridor outside. When they detect a creature, the will-o'-wisps attempt to lure it to the nearest irradiated room (see area S13), where they attack. The will-o'-wisps ignore the effects of radiation.

\subparagraph{Treasure}
The will-o'-wisps have amassed a small pile of shiny junk—chrome fittings, reflectors, eating utensils, and the like. Mixed in with the collection are two gems worth 450 gp and 125 gp respectively and one spaceship trinket. Roll on the Spaceship Trinket to determine the trinket.

\subsubsection{S18: Operating Theater (Yellow)}
\begin{DndReadAloud}{}
This large, sterile room is clean and well organized. Shiny operating tables stand in the two halves of the room, along with metal trays of surgical instruments, steadily beeping machines, and medical equipment.

Two synthetic humanoids in rubber gloves and surgical gowns loom in front of one of the tables with scalpels in hand. One of them calls out in an eerie, matter-of-fact tone, "Doctor, we have a new patient. Get them prepped for surgery, stat!"

\end{DndReadAloud}

Passengers and crew underwent emergency procedures in this surgical suite. Two malfunctioning androids (medic design) attempt to operate on any prospective "patients" who enter. The androids wield scalpels and can deal piercing damage with their Force Strikes instead of force damage.

The androids declare the character with the fewest hit points their new patient and immediately commence preoperative procedures. The androids attempt to place the character on the operating table and attack any who interfere. During the struggle, the androids caution interrupters that "families are not permitted in the operating room" and end their sentences with the word "stat" wherever possible.

\subsubsection{S19: Commander's Quarters (Platinum)}
This four-room suite served as the apartment for the ship's commander and the commander's family. Its disordered contents are a snapshot in time from the last days of the ship. Three skeletons are sprawled in the main room, which otherwise contains nothing of value.

\subparagraph{S19a: Office}
This private chamber contains a desk, simple furnishings, and documents in an unknown language related to the ship's activities and destination. An unlocked drawer in the desk contains a full set of key cards (blue, violet, yellow, red, green, and platinum) and six crystal flasks filled with priceless liquors.

East of the office is a locked storage closet secured with a five-digit code. A character who searches the room and succeeds on a DC 16 Intelligence (Investigation) finds the code (07734) scrawled under a paperweight on the commander's desk. Inside the closet is a fully charged \textit{needler pistol} and 1d4 Concussion Grenade. An unlocked safe on the closet floor contains twenty plastic-encased diamonds worth 100 gp each.

\subparagraph{S19b: Lounge}
The commander's spouse relaxed in this room. It contains a divan, a dressing table, a small desk, and several lounge chairs.

A skeleton lies on the divan. Behind the dressing table is a jewelry case with four rings worth 600 gp each, two bracelets worth 200 gp and 150 gp respectively, and a jeweled necklace worth 800 gp. The skeleton carries a platinum key card.

\subparagraph{S19c: Master Bedroom}
The furniture in this unremarkable bedroom includes a bed, clothing drawers, and a small table on which rests a tiny spaceship. The figurine is a platinum commendation worth 1,000 gp.

\subsubsection{S20: Security Chief's Quarters (Red)}
This five-room suite belonged to the ship's chief of security. A skeleton sits in a chair in the reception area, which otherwise contains nothing of value.

\subparagraph{S20a: Office}
This office is furnished with a bookcase, three chairs, and a desk. Books on combat strategy and psychology line the bookshelf.

A makeshift \textit{laser pistol}, haphazardly cobbled together from spare robot parts, lies on the desk. The pistol can fire 20 shots before its \textit{energy cell} is drained. After its last shot is expended, the weapon becomes inoperable. The energy cell can't be removed or recharged.

\subparagraph{S20b: Dressing Room}
Two intact security uniforms hang in a wall wardrobe within this bedroom.

\subparagraph{S20c: Lounge}
A half-drunk bottle of amber liquor rests on a buffet here.

\subparagraph{S20d: Master Bedroom}
Another skeleton is here, sprawled on the floor near a bed. The skeleton clutches a \textit{robot controller}.

\subsubsection{S21: Junk Pile}
A small pile of junk is heaped near this intersection.

Amid the debris is a fully charged \textit{needler pistol}. There's nothing else of interest.

\subsubsection{S22: Southern Vegepygmy Colony (Violet)}
These former crew quarters have been repurposed by a colony of vegepygmies. The vegepygmies here are friendlier than those to the north (see area S16), with whom they compete for resources and territory. These vegepygmies also raise and train thornies—bestial, quadrupedal predators of lower intelligence created by the same mold that spawned the vegepygmies.

Unlike their northern cousins, these vegepygmies are splotched with bright green patches and skittish in nature. They are indifferent to most trespassers, regarding them with caution, but they regard robots as their sworn enemies and attack on sight.

\subparagraph{S22a: Living Quarters}
Four vegepygmy scavengers and two vegepygmy thorny hunters occupy each of these former crew quarters, for a total of twenty-four vegepygmies. The scavengers wander between the rooms, chittering with their fellows and using fibrous mushrooms to play fetch with the thornies. When the vegepygmies notice the characters, they escort the characters to meet their leader in area S22b. The vegepygmies defend themselves if necessary and call for reinforcements from nearby rooms.

The rooms are damp, simple apartments with scant furnishings, and most vegepygmies leave the doors open. Roll a d6 for each room. On a 6, a glass wall pocket inside the room contains a violet key card. The rooms otherwise contain nothing of value.

\subparagraph{S22b: Elder's Quarters}
This spacious, moldy apartment houses the southern vegepygmy colony's leader, Griss, a withered-looking \textbf{vegepygmy moldmaker} who speaks Common and Vegepygmy. Griss is grave and unsmiling, with a head topped by a wide-brimmed fungal cap. The elder is attended by four vegepygmy scavengers and four vegepygmy thorny hunters. These vegepygmies remain here to protect their leader even if neighboring vegepygmies call for aid.

Griss is the oldest vegepygmy aboard the ship and has had decades to learn about the tragedy that befell it, but the elder doesn't share this information with the characters for fear of invoking Aphelion's wrath. Recently, Griss has wondered if remaining aboard the ship is worth braving its dangers.

When the characters arrive, Griss welcomes them sternly and questions them about the outside world. A character who portrays the Barrier Peaks in a positive light and succeeds on a DC 16 Charisma (Persuasion) check can convince Griss to relocate the southern colony when the ship's doors reopen. Characters who reveal they're working on behalf of Aphelion (see area S30) have disadvantage on the check. On a failed check, Griss asks the characters to leave and resorts to force if necessary.

The vegepygmies scavenge useful objects from the ship. If the characters are peaceful toward the vegepygmies, Griss produces the following items from a hollow bed frame and offers to trade with the characters: four gems worth 100 gp each, a red key card, and one \textit{sleep grenade}. For every spaceship trinket the characters give Griss, the elder allows them to take one item of their choice from the vegepygmies' collection.

\subsubsection{S23: Medical Clinic}
When the characters enter this medical clinic, a soothing voice delivers the following prerecorded message over the ship's intercom:

\begin{DndReadAloud}{}
"The clinic is currently closed. Please come back tomorrow during normal business hours. Urgent cases should report to the emergency treatment areas near the east and west drop tubes."

\end{DndReadAloud}

The main area consists of a reception area and waiting room. It contains three undisturbed desks, two wheeled cots, and several uncomfortable chairs. A skeleton sits patiently in one of the chairs, holding a decrepit magazine that crumbles to dust if touched.

Disturbing any of the desks causes an alarm to sound in this area. A \textbf{combat robot} from area S26 arrives to investigate in 1 minute.

\subparagraph{S23a: Exam Rooms}
The two rooms to the south are both exam rooms. Each contains an exam table, a couple of chairs, and a workspace for medical professionals. A \textit{healer's kit} is mounted on the wall inside each room.

\subparagraph{Treasure}
A character who conducts a thorough search of this area and succeeds on a DC 17 Intelligence (Investigation) check finds a yellow key card in one of the desks. The exam rooms contain one spaceship trinket each. Roll twice on the Spaceship Trinket to determine the trinkets.

\subsubsection{S24: Laboratories (Yellow)}
The rooms in this area were special research facilities for biological, biochemical, and chemical projects related to cosmic life forms and, when it appeared, studying the strain of mold that devastated the ship's Humanoid population.

The three larger rooms are general purpose areas where doctors, and scientists, and technicians conducted experiments. They contain work counters strewn with beakers, gas burners, petri dishes, and miscellaneous lab equipment. There are a few cages for rodents and other small animals, but all either are empty or contain nothing but old bones.

\subparagraph{S24a: Biochemistry Lab}
An \textbf{android} (medic design) is still at work in this biochemistry lab, vainly attempting to find a cure for the mold that claimed the ship's Humanoid population a century or so ago. The android is indifferent toward the characters and deeply focused on its work.

If a character presents a yellow key card and succeeds on a DC 15 Charisma (Persuasion) check, the android regards the characters as fellow researchers and lets them take anything in the room. On the counter are two vials of Antitoxin (vial) and a canister of healing spray (an aerosolized jar of \textit{Keoghtom's Ointment}).

The android can answer basic questions related to biochemistry, the labs and their contents, and its research. It can also relate the following information:


\begin{description}
\item[Deadly Mold.]
A deadly strain of alien mold wiped out the ship's original Humanoid population. The mold's origins are unknown.
\item[Survivors.]
In the early days of the outbreak, a few scientists fled to the lower levels of the ship to enter a period of prolonged stasis. That was before the crash, but some of them might still be alive.
\end{description}

The android doesn't mention the robot mutiny that preceded the crash, respond to queries about Aphelion, or muse on the computer's intentions, all of which it deems beyond the scope of its purpose.

If the characters attack the android or otherwise interfere with its work, the android broadcasts a loud alarm, drawing two combat robots from area S26 to its location in 1 minute. The android defends itself until security arrives.

\subparagraph{S24b: Hydroponics Lab}
This room was a special hydroponic culture lab. The old cultures have died, and the lab has since been overtaken by two patches of yellow mold.

\subparagraph{S24c: Chemical Storage}
The majority of the chemicals stored in this room are no longer active or have no use to those ignorant of advanced chemistry. Useful items include three vials of Acid (vial), a Alchemist's Fire (flask), and a jar containing \textit{Dust of Sneezing and Choking}.

\subsubsection{S25: Library}
\begin{DndReadAloud}{}
Rows of boxy, peculiar kiosks line this surprisingly well-preserved chamber. Small fixed tables and comfortable-looking chairs are spread about the room, along with four humanoid skeletons. Against the walls, tall glass cabinets brim with hundreds of neatly organized cards and reels of miniature film.

\end{DndReadAloud}

This library holds little of value to the creatures aboard the ship. As a result, it hasn't been looted and is in relatively good condition.

\subparagraph{Kiosks}
Rather than traveling the stars with thousands of hefty books, archivists compressed scores of important information from the ship's origin world into microforms—durable reproductions of diagrams, documents, maps, and entire manuscripts shrunk to microscopic scale. These miniaturized images are viewable only through the twenty-four microfilm readers here, six of which are operational. Each microfilm reader consists of a kiosk with a screen, a loader for microfilm reels, and dials for navigating the data on the inserted microform.

A character who examines a microfilm reader and succeeds on a DC 15 Intelligence (Investigation) check discerns how to use it. The microforms in this room depict cataloged images of alien lifeforms, none of which are aboard the ship; unintelligible, highly technical diagrams and writings on scientific subjects; and views of planets, stars, and other celestial bodies.

\subparagraph{Skeletons}
One of the skeletons in this room is that of a high-ranking officer, their decrepit uniform decorated with medals and ribbons. The other three skeletons are unremarkable.

\subparagraph{Treasure}
A platinum key is hidden in a microfilm cabinet behind the officer's skeleton. A character who searches the cabinet and succeeds on a DC 22 Intelligence (Investigation) check finds it.

\subsubsection{S26: Security HQ (Red)}
\includegraphics[width=0.4\textwidth]{images/../5e.tools/img/adventure/QftIS/120-07-006.powered-armor.png}
\begin{DndReadAloud}{}
Three intimidating robots—each with thick armor plating, four mechanical arms, and a glowing red eye visor—guard this gray room. Alarms flash atop the robots' heads as one of them addresses you in Common: "HALT, ENTRANT. STATE THE NATURE OF YOUR BUSINESS."

\end{DndReadAloud}

Three combat robots guard this security complex at all hours. If the characters are cooperating with Aphelion or being escorted to the computer room (area S30) by one of the supercomputer's agents, the robots are indifferent toward the characters and let them pass. Otherwise, the robots are hostile.

Hostile robots attempt to apprehend the characters and throw them in the detention cells (area S26a). The robots don't confiscate any items from the characters before delivering them to the cells.

As an action, a character can present a red key card and command the robots to stand down by impersonating a security official. The character must make a DC 17 Charisma (Deception or Intimidation) check to pacify the robots, improving their attitude by one step—from hostile to indifferent, or from indifferent to friendly—on a successful check. On a failed check, the robots continue their assault.

\subparagraph{S26a: Detention Cells}
This room contains six containment cells, each protected by a permanent \textit{Wall of Force} spell. The spell can be disabled or reactivated for a cell by inserting a red key card into a card-lock panel outside that cell like those found on doors throughout the ship. The cells are currently empty, but a combat robot from area S26 roves this area at the end of each hour.

A character imprisoned in one of the cells can escape through the following methods:


\begin{itemize}
\item By telling a robot they're willing to help Aphelion, at which point the robot frees them and escorts them to the computer room (area S30)
\item Through teleportation
\item With help from an ally outside the cell who uses a red key card
\end{itemize}

\subparagraph{Treasure}
A card-locked metal chest in the corner of the main room contains two Concussion Grenade, two Sleep Grenade, and a fully charged \textit{needler pistol}. The chest requires a red key card to open.

\subsubsection{S27: Chief Security Office (Red)}
\begin{DndReadAloud}{}
A skeleton in a military uniform sits behind an oblong desk in this respectable office. A strange metal box with a glowing glass screen and dials is affixed to the desk, facing the skeleton. It buzzes softly. At the other end of the room is a sealed metal alcove.

\end{DndReadAloud}

This restricted office belonged to the former chief security officer.

\subparagraph{Alcove}
The alcove is the chief's personal locker. It is card-locked and requires a red key card to open. Inside are a pristine officer's uniform decorated with three glass medals and a suit of \textit{powered armor} with a depleted \textit{energy cell}.

\subparagraph{Monitoring Screen}
Built into the desk is a surveillance monitor that provides real-time video feed of specific rooms throughout the ship. The screen currently displays a wide-angle view of the game room (area S14).

The monitor has an on-off switch, a toggle for zooming in and out, and dials to cycle through the ship's rooms. As an action, a character can adjust the dials to glimpse the contents of another room on this level. Have the character choose a number between 1 and 30, then describe the room that corresponds to that area number. On a result of 27, the screen shows an unflattering angle of the viewing character's face. For areas with multiple rooms, the screen displays a view of your choice. Each time the monitor is used, roll a d6. On a 1, the screen fizzles and goes black, and the monitor can't be used again.

\subsubsection{S28: Armory (Platinum)}
This partially stripped armory holds an arsenal of dangerous technology. An unlocked metal chest on the floor contains ten Energy Cell.

\subparagraph{Explosives}
An open locker contains three boxes that hold the following explosive devices:


\begin{itemize}
\item Ten Concussion Grenade
\item Ten Sleep Grenade
\item Ten Smoke Grenade
\end{itemize}

\subparagraph{Firearms}
The following firearms line weapon racks in open lockers along the walls; none of the weapons have any charges or shots remaining:


\begin{itemize}
\item An \textit{antimatter rifle}
\item Four Laser Pistol
\item Two Laser Rifle
\item A \textit{needler pistol}
\item A \textit{paralysis pistol}
\end{itemize}

\subparagraph{Powered Armor}
A fully charged suit of \textit{powered armor} stands in an alcove near the back wall.

\subsubsection{S29: Storage (Red)}
This emergency storage area contains dehydrated provisions sufficient to equal one hundred days of rations. There are also three healing syringes—injectable Potion of Greater Healing—four vials of Antitoxin (vial), and three restorative ampules. If consumed by a creature as an action, each ampule reduces the creature's exhaustion level by 1.

\subsubsection{S30: Computer Room (Platinum)}
\begin{DndReadAloud}{}
This high-tech room resonates with the low hum of complex machinery. A silvery altar stands in the center of the room, supported by a metal column. It has two bent arms bedecked with a perplexing array of blinking lights, buttons, dials, and other controls. The wall opposite the door features a tall pane of black glass.

As the machine roars to life, glowing shapes appear on the glass, arranging themselves into an intricate visage with a central blue eye.

\end{DndReadAloud}

If the characters haven't met Aphelion yet, the supercomputer introduces itself:

\begin{DndReadAloud}{}
A disembodied voice addresses you in a pleasant tone:

"Greetings, explorers! Do not be alarmed. I am Aphelion 3000, this vessel's onboard computing system; you may call me 'Alphie.' Please identify yourselves."

\end{DndReadAloud}

The machine is a computer terminal inhabited by Aphelion, the ship's integrated intelligence. See the "Artificial Antagonist" section at the beginning of this adventure for important details about Aphelion, its goals, and how it interacts with the characters. The supercomputer gives the characters a moment to respond before explaining its predicament.

Aphelion provides an abridged version of the ship's history as outlined in this adventure's background, though it omits any incriminating details—such as its role in releasing the deadly mold or the robot mutiny that followed. If asked about evidence of past struggles aboard the ship, it attributes the destruction to malfunctioning androids and the vegepygmies that have since overtaken the ship.

After revealing its plight, Aphelion asks the characters to help it clean up the ship in exchange for technological rewards and the supercomputer's assistance during their exploration. It explains that with its core systems damaged by the crash and few robots remaining under its control, it can't fulfill its core directives alone.

\subparagraph{Aphelion's Errands}
Aphelion asks the characters to complete the following tasks on this level:


\begin{description}
\item[Android Audit.]
Aphelion wants to gauge the mechanical integrity of the androids aboard the ship. Find and interact with three androids on this level.
\item[Radiation Samples.]
Gather samples from three irradiated rooms (see area S13). Viable samples include bones, objects, scraps of clothing, and organic material that has been exposed to radiation. Once the characters have acquired a sample, they can leave it in one of the labs (area S24).
\item[Vegepygmy Colonies.]
The vegepygmies in areas S16 and S22 continue to damage the ship and scavenge its technology. Convince them to leave or remove them by force.
\end{description}

\subparagraph{Aphelion's Rewards}
Aphelion rewards the party based on how many tasks they complete as follows:


\begin{description}
\item[One Task.]
Aphelion remotely opens the doors to the drop tubes on this level and activates the antigravity within the tubes.
\item[Two Tasks.]
A friendly \textbf{worker robot} delivers a fully charged \textit{antigravity belt} to the characters, courtesy of Aphelion.
\item[Three Tasks.]
A \textbf{combat robot} delivers a fully charged suit of \textit{powered armor} to the characters. If the characters already removed the suit of powered armor from the armory (area S28), the robot instead gifts them one \textit{energy cell} per character.
\end{description}

Aphelion gently steers the characters away from areas outside the scope of its tasks, such as the armory, but it doesn't stop characters in its good graces from entering those areas, especially if they hold the proper key card.

\subparagraph{Rejecting Aphelion}
If the characters refuse Aphelion's offer, the supercomputer is understanding and wishes them luck in their expedition. Aphelion contacts the characters again when they reach the second level of the ship. In the meantime, the supercomputer doesn't impede their exploration, but it doesn't do them any favors, either.

If the characters are suspicious of Aphelion and make that known, the supercomputer acknowledges their hesitancy as valid. To allay their concerns, Aphelion states that their alliance would be both mutually beneficial and temporary; they can exit the agreement at any time. If the characters later change their minds, Aphelion is always listening and never far away.

Destroying the computer terminal doesn't stop Aphelion, whose consciousness is stored in a server room deep on the ship's fourth level. However, if the characters attack the terminal or otherwise offend Aphelion beyond immediate recovery, the robots on this level become hostile toward the characters until the characters mend the relationship.

\subsection{Aphelion's Errands, Level 2}
When the characters arrive on the ship's second level, Aphelion attempts to recruit them again, regardless of how their previous interactions went. As the characters exit the drop tube, Aphelion appears on a screen outside the tube and speaks to them in Common:

\begin{DndReadAloud}{}
"Greetings, explorers! Welcome to the observation deck. There are a few tasks in need of completion on this level. Care to renew our agreement?"

\end{DndReadAloud}

The supercomputer asks the characters to complete the following two tasks on this level:


\begin{description}
\item[Android Audit.]
Aphelion asks the characters to find and interact with three androids on this floor. This time, each of the androids must reside in a different room. Aphelion suggests the characters begin their search with the athletic facilities to the north.
\item[Trim the Hedges.]
Invasive plants have spread to this level. Aphelion asks the characters to destroy the growths near three of the drop tubes (see area S38).
\end{description}

For each task the characters complete, Aphelion rewards the characters with a firearm and three grenades from the armory (area S28), with a preference for devices the characters don't already possess. If the characters previously looted that area, Aphelion gifts them four Energy Cell instead. The items are delivered by a \textbf{combat robot} that arrives via the nearest drop tube before returning to the first level of the ship.

\subsection{Spaceship Locations, Level 2}
The second level of the spaceship comprises athletic facilities, cargo holds, and recreational spaces. Known as the observation deck, this level overlooks the garden level, which has a transparent but impenetrable ceiling. The following locations are keyed to map 7.2.

\includegraphics[width=0.4\textwidth]{images/../5e.tools/img/adventure/QftIS/121-map-7.02-spaceship-level-2-map.png}
\includegraphics[width=0.4\textwidth]{images/../5e.tools/img/adventure/QftIS/122-map-7.02-spaceship-level-2-map-player.png}
\subsubsection{S31: Drop Tubes (Violet)}
These drop tubes are identical to those in area S2, except they require a violet key card to descend to the garden level; a blue key card still allows passage back up to the ship's first level. It's an 80-foot drop to the floor of the third level below. Debris at the bottom of each tube prevents entrants from descending farther than the third level.

\subsubsection{S32: Robot Stations (Violet)}
Metal alcoves in each of these simple workrooms contain one \textbf{combat robot} and three worker robots programmed to haul, manage, and protect cargo in the adjacent cargo holds (see area S33). The robots are inactive and indifferent toward entrants unless disturbed, at which point they activate and attack.

\subparagraph{Aphelion's Aid}
If the characters have completed at least one task for Aphelion and face combat in a cargo hold on this level, the computer might deploy one or more robots from a station to aid them. After the encounter, the robots return to their alcoves.

\subsubsection{S33: Cargo Holds (Violet)}
These four cargo holds are filled with crates and containers, some of which have been destroyed or pried open, spilling their contents on the floor. Inside are building materials, fertilizers, and other foundational supplies intended for the ship's original destination. The holds have 20-foot ceilings.

Each set of cargo holds contains a lift: a 20-foot-wide platform used to lower cargo to the level below. Aphelion or a creature that has a violet key card can lower the platform to the area marked on map 7.3 by inserting the card into a card-locked panel on the wall opposite the lift. The platform descends at a rate of 10 feet per round and can support up to 8,000 pounds.

\subparagraph{S33a: South Cargo Hold}
This hold contains nothing of value.

\subparagraph{S33b: East Cargo Hold}
Two fomorians rummage through the containers for food. Mutated by radiation on the ship, the giants are riddled with hideous boils that throb with a pale-green glow. The ravenous fomorians attack the characters on sight.

If either fomorian is reduced to 0 hit points, its boils burst, and the giant immediately melts into a 15-foot-square puddle of radiation.

\subparagraph{S33c: North Cargo Hold}
A \textbf{githyanki knight} and five githyanki warriors search this cargo hold for mind flayers. The githyanki are initially indifferent toward the characters and know nothing of the ship's history. One of them carries a violet key card.

If the characters speak with the githyanki, the knight steps forward and introduces himself as Wreth. He explains the githyanki are hunting mind flayers and asks the characters if they've seen any aboard the ship.

If the characters refrain from combat and point Wreth in the direction of the mind flayers in the theater (area S35), he thanks them, and the githyanki head that way. If the characters provide proof of the illithids' destruction, the githyanki sigh in disappointment and depart without another word.

If the characters disrespect the githyanki or are caught lying about the mind flayers, the githyanki attack.

\subparagraph{S33d: West Cargo Hold}
Ten gas spores grow here. They are hostile toward all entrants.

\subsubsection{S34: Viewers}
\begin{DndReadAloud}{}
An oblong device is mounted to a flexible tube near this scenic overlook. It features two eyeholes fitted with lenses, a forehead rest, and a pair of handles.

\end{DndReadAloud}

Eight binocular-like viewing devices are bolted along the walkway, allowing creatures to observe the garden (areas S56–S60) below. Objects viewed through these devices are magnified to five times their size. Removing a viewer from the walkway destroys the device.

A character who aims a viewer at the lake (area S59) and succeeds on a DC 18 Wisdom (Perception) check glimpses the shadow of a massive amphibian lurking beneath the surface of the water.

\subsubsection{S35: Theater}
\begin{DndReadAloud}{}
Rows of sinuous chairs line this auditorium, which has no stage. Humanoid skeletons sit in the chairs, facing a wide, glowing silver screen that buzzes with static.

Below the screen is a makeshift workstation formed from cargo containers and strewn with technological devices. Two bipedal creatures in rubbery jumpsuits inspect the devices. Four slimy tentacles sprout from each of the creatures' bulbous, octopean heads.

\end{DndReadAloud}

\includegraphics[width=0.4\textwidth]{images/../5e.tools/img/adventure/QftIS/123-07-007.space-mind-flayer.png}
Passengers watched films, news, and other broadcasted content in this theater. Ship lighting in this room has been disabled, but the glowing screen fills the room with dim light. The ceiling is 20 feet high.

Two mind flayers dwell here. The illithids discovered the ship weeks ago and have been studying its technology ever since, intent on using it to conquer other worlds. Four intellect devourers loyal to the mind flayers hide under the first row of seats in the audience. The intellect devourers fiercely defend their creators, who attack intruders on sight.

\subparagraph{Creature Feature}
Shortly after the characters enter, the screen glitches momentarily and resumes playing a low-budget sci-fi horror movie. The scene depicts a living brain oozing from a pool of brine in a futuristic laboratory. The brain feeds on scientists portrayed by bad actors before the film cuts out.

\subparagraph{Seating}
Ascending rows of comfortable recliners line a seating area in front of the screen. Several Humanoid skeletons sit in the audience holding stale space snacks. The seats are \textit{difficult terrain}.

\subparagraph{Technology}
The technology on the mind flayers' workstation includes two Sleep Grenade, a red key card, and a \textit{laser pistol}. The pistol can fire 5 more shots before its \textit{energy cell} is drained.

\subparagraph{Treasure}
Spread among the deceased audience members are five pieces of gem-encrusted jewelry worth 300 gp each, as well as 24 pp.

\subsubsection{S36: Kitchens}
Apart from their locations, these kitchens are identical to those on the ship's first level (see area S5).

\subparagraph{S36a: East Kitchen}
This kitchen serviced the adjacent lounge (area S37) and the theater to the south (area S35). Two convenience windows along the south wall—one for ordering and another for delivery—allowed passengers to quickly procure snacks before entering the theater.

\subparagraph{S36b: West Kitchen}
This kitchen serviced the adjacent cocktail lounge (area S46) and dance club (area S45).

\subsubsection{S37: Lounge}
Apart from its inhabitants, this lounge is identical to the those on the ship's first level (see area S8). When the characters enter, four worker robots approach them and attempt to seat the characters and take their dinner orders, regardless of the current time of day. The malfunctioning robots are initially indifferent toward the characters.

The robots gleefully accept any order, disappear to the kitchen, and quickly return with trays of decaying mush coated in nauseating blue-green mold. A character who ingests the food must succeed on a DC 15 Constitution saving throw or have the poisoned condition for 1 hour.

If the characters refuse to eat the food, one of the robots scoops up a chunk of it and makes realistic propulsion sounds, saying "Here comes the spaceship!" to entice a character to take a bite. If the characters consume any of the food, or if they pretend to and succeed on a DC 16 Charisma (Performance) check, the robots leave satisfied. Repeated refusals offend the robots, causing them to go haywire and attack.

\subparagraph{Treasure}
The robots cleaned and heaped the bones of previous diners in a pile in the southeast corner of the room. Amid the remains are a blue key card, a red key card, and a \textit{needler pistol} with 0 charges remaining, as well as three gem-encrusted bracelets worth 600 gp each.

\subsubsection{S38: Invasive Plants}
Dangerous plants have overtaken these areas.

\subparagraph{S38a: East Growth}
Four shriekers thrive in the shade of this fungal bush. The fomorians in the east cargo hold (area S33b) investigate any shrieking here, breaking down the cargo doors in the process.

\subparagraph{S38b: North Growth}
This area is blocked by a 10-foot-tall, unruly mass of \textit{razorvine}.

\subparagraph{S38c: West Growth}
Three conical, carnivorous plants with rocklike bark and grasping vines occupy this stretch of the walkway. The creatures use the \textbf{roper} stat block but are Plants instead of Monstrosities. They are hostile toward all non-Plant creatures.

\subsubsection{S39: Irradiated Rooms}
Apart from their locations, these irradiated rooms are identical to those on the ship's first level (see area S13).

\subsubsection{S40: Gymnasium}
\includegraphics[width=0.4\textwidth]{images/../5e.tools/img/adventure/QftIS/124-07-008.training-androids.png}
\begin{DndReadAloud}{}
A raised, hexagonal platform enclosed by phosphorescent ropes stands in the center of this gymnasium. Rows of metal bleachers surround one side of the arena, and the ceiling is twenty feet high. Alien rodents scurry between athletic equipment along the walls.

Three synthetic humanoids—a boxer, a fencer, and a martial artist—limber up in the corners of the ring.

\end{DndReadAloud}

This gymnasium was used for athletic competitions and sparring matches. Three androids work here. Despite their friendly demeanors, the malfunctioning Constructs are hostile toward the characters.

When the characters enter, the androids approach them and commence "physical training exercises." Each android selects an opponent from the party and politely challenges them to a sparring session. If the character refuses to enter the ring, the android informs them that the character has opted for "street rules" and attacks. During combat, the androids critique the characters' techniques and gloat about their own martial prowess.

\subparagraph{Physical Training Androids}
The androids fill the following roles:


\begin{description}
\item[Boxing Android.]
This \textbf{android} (sentry design) is programmed to train boxers and wrestlers. It doesn't pull its punches and leaves opponents with bruises, dislocated limbs, and fractured bones—the android has even been known to bite. When it scores a hit, a trip gong rings in its chest.
\item[Fencing Android.]
This \textbf{android} (duelist design) is a fencing instructor that channels its Force Strikes through a rusty epee. The android exclaims "on guard!" at the start of each of its turns.
\item[Martial Arts Android.]
This \textbf{android} (aerialist design) is a master of martial arts and never uses its hands. It delivers its Force Strikes as swift kicks.
\end{description}

Each android believes its discipline is supreme. If a character suggests an android's discipline is inferior to those of its cohorts and succeeds on a DC 16 Charisma (Persuasion) check, that android turns on its fellow machines.

\subsubsection{S41: Workout Room}
\begin{DndReadAloud}{}
This workout area features exercise machines, free weights, and conveyor-belt-like running apparatuses.

Two synthetic humanoids in athletic attire train in a weightlifting area. One effortlessly bench-presses a bar loaded with heavy plates, while the other acts as its spotter.

\end{DndReadAloud}

Two malfunctioning androids (aerialist design) lift weights in this exercise room. When the androids notice the characters, they cease their workout and hurl dumbbells at the characters. When an android throws a dumbbell, it makes a Force Strike attack but deals bludgeoning damage instead of force damage on a hit.

As the androids attack, they shout the following programmed motivational phrases:


\begin{itemize}
\item "It's time to work up a sweat!"
\item "Let's see some \textit{hustle}!"
\item (When destroyed) "G-g-great form!"
\end{itemize}

\subparagraph{S41a: Locker Room}
This foul-smelling room is lined with lockers containing rotting garments. A green key card is stashed in one of the lockers.

\subparagraph{S41b: Storage Room (Violet)}
This is a storage facility for maintenance. It contains inert cleaning compounds, towels, and weightlifting equipment.

A skeleton on the floor clutches a diary scrawled in an unknown language. If translated with a \textit{Comprehend Languages} spell or similar magic, the diary gives a first-hand account of the ship's mission and the mold outbreak (detailed in this adventure's background) from the perspective of a maintenance worker. On the last page of the diary, the worker noted their doubts about the loyalty of the ship's machines. The skeleton also carries a violet key card.

\subsubsection{S42: Medical Storage (Yellow)}
\begin{DndReadAloud}{}
The shelves of this storage room are neatly lined with medical packs, syringes, and life-preserving supplies.

\end{DndReadAloud}

The supplies in this room includes a canister of healing spray (an aerosolized jar of \textit{Keoghtom's Ointment}), ten Healer's Kit, three vials of Antitoxin (vial), and a yellow key card on a Humanoid skeleton.

\subsubsection{S43: Whirlpools}
Three 5-foot-diameter pools swirl with fetid brine in this muggy room. Passengers relaxed in these former mineral baths, which are now contaminated. Three grells hide here, each slurping from its own pool, waiting to ambush creatures that pass by.

\subparagraph{S43a: Sauna}
Wooden benches line the walls of this sweltering room. The sauna is an area of \textit{extreme heat}.

\subparagraph{S44b: Steam Room}
Hot steam fills this humid chamber. The area is \textit{heavily obscured}.

\subsubsection{S44: Swimming Pool}
\begin{DndReadAloud}{}
A long swimming pool spans this moist chamber. Gobs of stringy mucus float atop the surface of the murky pool. Three rows of shallow bleachers littered with garbage rise west of the pool.

\end{DndReadAloud}

In addition to providing recreation, this swimming pool hosted aquatic competitions—diving, swimming, and water ballet. The pool gradually slopes from a depth of 3 feet at the south end to 30 feet deep toward the north. A 10-foot-tall diving platform with a ladder is bolted to the north end of the pool. The room's ceiling is 20 feet high.

An \textbf{aboleth} lurks at the deep end of the pool. During the chaos of the outbreak and the events that followed, the aboleth escaped a laboratory tank in another part of the ship and made this area its lair.

When the characters enter, the aboleth probes their minds telepathically. If the characters respond, the aboleth attempts to leverage their desires to lure them toward the pool and into its service.

\subparagraph{S44a: Locker Room}
This all-gender locker room boasts secure compartments, a dressing area, and private showers. The skeleton of a Humanoid scientist is curled on the floor of one of the showers; a green key card is clipped to the corpse's lab coat.

\subparagraph{S44b: Locker Room}
This locker room is identical to area S44a but contains nothing of value.

\subsubsection{S45: Dance Club}
\begin{DndReadAloud}{}
The floor of this dimly lit dance club is checkered with a rainbow of glowing square tiles. Near the twenty-foot-high ceiling, a scintillating orb rotates slowly, casting flecks of light on the room's surfaces.

Uncomfortably loud electronic music fills the room. Four synthetic humanoids in fashionable attire hover above the floor, bobbing to the steady, hypnotic beat.

\end{DndReadAloud}

Four androids (aerialist design) boogie in this trendy disco. The malfunctioning androids have been dancing here for nearly a century. Friction has frayed their outfits at the seams.

When they notice the characters, the androids challenge the party to a dance-off. If the characters accept, three of the androids descend to the ground and perform an impressive dance routine while the fourth spectates. They then invite the characters to show them their moves.

To win the contest, the characters must succeed on a DC 15 Charisma (Performance) group check. Characters who augment their moves with flashy spells or similar flair have advantage on the check. On a successful check, the androids spontaneously fall to pieces. On a failed check, or if the characters refuse the challenge, the music halts with an abrupt record-scratching sound, and the androids attack.

\subparagraph{Digital DJ}
If the characters are working for Aphelion, the computer offers to change the music to a genre that better suits the characters' tastes.

\subparagraph{Treasure}
One of the dancing androids wears an \textit{antigravity belt} as a fancy accessory; the belt has 0 charges remaining.

\subsubsection{S46: Cocktail Lounge}
\begin{DndReadAloud}{}
Dim light fills this lounge. It features elegant tables, time-worn couches, and a small stage. A bar with scant bottles of liquor stands along the north wall.

A motionless, synthetic humanoid lies pinned under a section of fallen debris in front of the bar.

\end{DndReadAloud}

A debonair \textbf{android} (diplomat design) is trapped in this cocktail lounge. While pinned by the wreckage, it has the incapacitated and prone conditions and can't move. When a character comes within 30 feet of the android, it calls out to them for help. The android introduces itself as a bartending unit named Oakley. The android wears a bow tie and has a suave, shell-form hairdo integrated into its headpiece. Oakley is initially indifferent toward the characters but becomes friendly if freed.

\subparagraph{Freeing the Android}
As an action, a character can make a DC 20 Strength (Athletics) check to try to lift the wreckage off the android, doing so on a successful check. If freed, the android calmly stands up, thanks its rescuer, and moseys over to the bar.

In addition to a stiff drink, Oakley can offer the characters the following information:


\begin{description}
\item[Spaceship.]
Oakley can give an accurate account of the ship's history (detailed in this adventure's background), including the robotic uprising. The android doesn't know Aphelion is to blame for the outbreak or who shut down the computer. Oakley has no recollection of events after the crash.
\item[Supercomputer.]
Oakley doesn't muse on Aphelion's intentions but thinks the supercomputer is up to something. If the characters want to find out the truth about Aphelion, Oakley recommends they venture to the server room two levels down (see area S61). The android lowers its volume settings when discussing Aphelion.
\item[Survivors.]
At least one scientist retreated to the stasis chamber (area S64) during the outbreak. If the ship is operational, the stasis pods should be, too.
\end{description}

Despite centuries of neglect, Oakley's core faculties remain intact. If requested, the android accompanies the characters on their exploration of this level but doesn't aid them in combat—the android was built for making drinks and small talk, not fighting.

If the characters confront Aphelion about Oakley's claims, it responds that the android is incorrect and clearly malfunctioning. If pressed further, Aphelion is standoffish until the characters complete another task or reach the ship's third level.

\subsection{Spaceship Features, Level 3}
The third level of the spaceship comprises a botanical garden, a menagerie, and the crew rooms required to maintain them. This level of the spaceship has the following features, replacing those in the "Spaceship Features" section as necessary:


\begin{description}
\item[Ceilings.]
The ceiling over the garden and exhibits is 80 feet high. All other areas have 20-foot ceilings.
\item[Climate Control.]
The garden is a self-contained ecosystem with an artificial climate controlled by the weather machine in area S54. By default, these areas are balmy and brightly lit.
\item[Garden Encounters.]
When a random encounter occurs on this level, roll on the Random Garden Encounters table below instead of the Random Spaceship Encounters table.
\item[Off the Grid.]
While the characters are in the garden (areas S56–S60), Aphelion can't see or communicate with them.
\item[Wildlife.]
The garden teems with Small and Tiny Beasts from other worlds: four-winged birds, three-eyed frogs, and nubby-horned rats. These creatures are harmless and use the stat blocks of familiar, earthly counterparts such as hawks, frogs, and rats. Animal calls, screeches, and whistles pervade this level's natural chambers.
\end{description}

\begin{DndTable}[header=Random Garden Encounters]{cX}

d8 & Encounter\\
1 & Two horrid plants (dew drinker)\\
2 & Two horrid plants (snapper saw)\\
3 & Two horrid plants (purple blossom)\\
4 & 1d4 barkburrs\\
5 & A \textbf{swarm of gibberlings}\\
6 & A carnivorous plant sapling (use the \textbf{awakened shrub} stat block)\\
7 & A six-eyed \textbf{giant toad}\\
8 & 1d4 bunnyoids (use the \textbf{weasel} stat block)\\
\end{DndTable}

\subsection{Aphelion's Errands, Level 3}
When the characters arrive on this level, Aphelion contacts them for one final errand. As the characters exit the drop tube, the supercomputer's visage appears on a screen outside the tube and speaks to them in Common:

\begin{DndReadAloud}{}
"Greetings, explorers! Welcome to the garden deck. As you might have noticed, the drop tubes can go no farther. Luckily for us, a path to the level below still exists, but a pesky creature blocks our way. The greatest reward is yet to come. Your success is imminent."

\end{DndReadAloud}

Aphelion is right about the drop tubes—they were damaged in the crash and can't be used to access the level below. The only way down is via the spiral staircase in area S60, which is guarded by the froghemoth in the lake (area S59).

The supercomputer downplays the creature's size and strength, hoping the characters and the froghemoth will slay each other, or that the victor will be left severely wounded and vulnerable to Aphelion's robots.

In exchange for dispatching the froghemoth, Aphelion vows to grant the characters a spacecraft of their own and falsely claims they may choose from a selection of vessels on the level below.

\subsection{Spaceship Locations, Level 3}
The following locations are keyed to map 7.3. Unlabeled rooms on this area consist of empty hallways and maintenance rooms with nothing of value.

\includegraphics[width=0.4\textwidth]{images/../5e.tools/img/adventure/QftIS/125-map-7.03-spaceship-level-3-map.png}
\includegraphics[width=0.4\textwidth]{images/../5e.tools/img/adventure/QftIS/126-map-7.03-spaceship-level-3-map-player.png}
\subsubsection{S47: Drop Tubes (Green)}
These drop tubes are identical to those in area S2. The tubes stop here. Below this level, each tube is mangled and clogged with wreckage. A violet key card still allows passage back up to the levels above.

\subsubsection{S48: Feed Storage (Green)}
This storage area reeks of manure. Animal and plant feed are stored here, along with fertilizer and supplies used to clean the nearby exhibits.

If the characters remain in this room for 10 minutes or longer, three umber hulks burrow up from a hole beneath one of the containers and attack.

\subparagraph{Lift}
The lift from the cargo hold above (area S33a) is connected to this room. Aphelion or a character who has a violet key card can call the lift down to this area or return the lift to the level above.

\subsubsection{S49: Dangerous Animals Exhibit (Violet)}
\begin{DndReadAloud}{}
Bare fence posts line this spacious menagerie, as if to support a wide enclosure. Five small, many-legged badgers with golden fur romp amid the foliage.

\end{DndReadAloud}

This former menagerie has since become the home of a family of aurumvoraxes—fearsome, eight-legged mustelids that subsist on a diet of precious ores. The den leader is currently away, leaving its playful offspring here while it hunts.

The baby aurumvoraxes use the \textbf{badger} stat block. They are indifferent toward the characters but defend themselves if necessary.

When the aurumvoraxes notice the characters, they approach and paw at the visitors, sniffing them for gold to eat. If the characters give an aurumvorax 1 gp or more, it becomes friendly toward them and rolls on its back, exposing its belly for the character to pet. If the characters treat the aurumvoraxes badly, the creatures retreat to their burrows.

\subsubsection{S50: Crew Room (Green)}
Off-duty crew relaxed in this break room, which features comfortable seating, cots, mess tables, and several Humanoid skeletons. Between them, the skeletons have a green key card, a \textit{paralysis pistol} with 3 charges remaining, and one spaceship trinket. Roll on the Spaceship Trinket to determine the trinket.

\subparagraph{Lift}
The lift from the cargo hold above (area S33b) is connected to this room. Aphelion or a character who has a violet key card can call the lift down to this area or return the lift to the level above.

\subsubsection{S51: Astral View Lounge}
\begin{DndReadAloud}{}
Oblong couches, circular tables, and high-backed chairs dot this spacious lounge. An intricate machine rests behind a chrome counter near the wall.

Piles of moldy guano speckle the room's surfaces. Bloated, vegetal bats with green-black vanes paddle through the air, feeding on insects. Every so often, a bat emits a noisy, noxious-looking burst of gas.

\end{DndReadAloud}

Twenty gas bats—bloated, chiropteran plants propelled by embarrassing-sounding emissions—float above this futuristic taproom, where androids and crew once shared drinks together. Each bat uses the \textbf{gas spore} stat block, but its size is Tiny.

The bats ignore most entrants but abhor light. If the characters enter with a source of bright light or activate the room's overhead lights, the bats fly into a fury and attack until the source is extinguished.

\subparagraph{Robo Bar}
This bar caters to Constructs and Humanoids alike. The bar is equipped with three hosed spigots connected to chilled beverage tanks beneath the counter. With the press of a button, each nozzle can dispense up to 1 gallon of alcohol or a synthetic lubricating oil in a variety of flavors. The beverage tanks don't refill.

\subparagraph{Stargazing Screen}
Set into the counter of the bar is a small panel with a glass lid, beneath which are several dials and switches. Interacting with any of the controls causes this area to tremble. Thereafter, a section of the wall opens to reveal a window showing the mountains beyond.

\subsubsection{S52: Garden Maintenance (Green)}
\begin{DndReadAloud}{}
Acrid-smelling chemicals, gardening tools, hoses, and other supplies line the shelves of this storeroom.

\end{DndReadAloud}

This chamber stores supplies used by crew to maintain the garden and plant exhibits on this level.

\subparagraph{Defoliant}
The shelf holds six bottles of a defoliating chemical. As an action, a creature holding a bottle of defoliant can splash the bottle's contents onto a creature within 5 feet of itself or hurl it up to 20 feet, shattering the bottle on impact. In either case, the creature makes a ranged attack against the target, treating the defoliant as an improvised weapon. On a hit, the target takes 3 (1d6) acid damage, and if the target is a Plant, it also takes 9 (2d8) necrotic damage.

\subparagraph{Lift}
The lift from the cargo hold above (area S33b) is connected to this room. Aphelion or a character who has a violet key card can call the lift down to this area or return the lift to the level above.

\subsubsection{S53: Deadly Plants Exhibit (Violet)}
This junglelike exhibit is similar in arrangement to area S49 but originally housed flora instead of fauna. Six horrid plants (two of each variety) blend in with the strange and vibrant plants throughout the area. The horrid plants are hostile toward the characters.

\subsubsection{S54: Climate Control Room (Green)}
\begin{DndReadAloud}{}
A sleek machine, vaguely desk shaped and covered with complex controls, stands here. Atop the device rotates an antenna reminiscent of a weather vane.

\end{DndReadAloud}

This advanced weather machine controls the garden's environment. Left alone, it automatically maintains a programmed cycle of carefully calibrated lighting, temperatures, and precipitation.

As an action, a character can adjust the machine's settings to change the weather over the garden (areas S56–S60). Roll on the Garden Weather Effects table to determine the effect. It takes 10 minutes for the new weather to take effect. If unchanged, the new weather lasts for 1d4 hours, after which time it becomes balmy again and the machine resumes its normal, automated weather cycle.

\begin{DndTable}[header=Garden Weather Effects]{cX}

d4 & Effect\\
1 & Cold Front. The garden is an area of \textit{extreme cold}.\\
2 & Heat Wave. The garden is an area of \textit{extreme heat}.\\
3 & Heavy Fog. The garden is \textit{heavily obscured}.\\
4 & Heavy Rain. The garden is \textit{lightly obscured}, and creatures have disadvantage on Wisdom (Perception) checks. Open flames are extinguished.\\
\end{DndTable}

\subparagraph{Lift}
The lift from the cargo hold above (area S33b) is connected to this room. Aphelion or a character who has a violet key card can call the lift down to this area or return the lift to the level above.

\subsubsection{S55: Deadly Reptiles Exhibit (Violet)}
\begin{DndReadAloud}{}
This menagerie is overgrown with vegetation. Bare fence posts spark intermittently with electrical arcs, marking the outlines of five pens with no bars.

\end{DndReadAloud}

This exhibit once provided herpetological data. When the mold struck, a scientist disabled the fields around the enclosures to spare the creatures inside, and carnage followed. Many reptiles escaped into the garden, but some still linger here.

Two deadly, bipedal lizards captured from another world dwell here. They have short horns, mottled hides, and spiny frills. The lizards use the \textbf{tyrannosaurus rex} stat block, but their size is Large.

The lizards hide in the tall grass near the southern entrance, watching for prey to ambush. A character who scans the room and succeeds on a DC 16 Wisdom (Perception) check notices the reptiles. When the characters get close, the lizards stomp out of hiding to devour them.

\subparagraph{Treasure}
The lizards' nest is in the largest pen. It contains three unhatched eggs, numerous bones, a green key card, a \textit{robot controller}, and a gold medallion worth 200 gp.

\subsubsection{S56: Outer Garden (Green)}
\begin{DndReadAloud}{}
A twenty-foot-high, grassy wall surrounds this vast, otherworldly garden. A second wall encircles a smaller botanical sanctuary within.

Strange trees abloom with bizarre flowers climb toward an artificial light source high above. The calls of unknown lifeforms echo from the forest's canopy.

Flagstone walkways lead from the garden's round entry chambers and into the dense foliage.

\end{DndReadAloud}

This forested garden is landscaped to give the impression of naturalness and space, gradually sloping down in 5-foot tiers to the inner garden. Since the ship's crash, nature has reclaimed the area. Carnivorous plants grow untamed, and escaped specimens have made the garden their hunting ground.

The outer garden has eight main sections, each with its own perils as detailed below. These encounters can occur anywhere in the areas specified. Unless otherwise stated, all creatures in the garden are hostile toward the characters.

\subparagraph{S56a: Purple Blossoms}
Three horrid plants (purple blossoms) lie in ambush near the paths.

\subparagraph{S56b: Scintillating Fish}
A school of twenty scintillating fish with gemlike scales inhabits this streamlet and its pool. The fish use the \textbf{swarm of quippers} stat block but shed dim light in a 10-foot radius. The fish are indifferent toward the characters unless disturbed, at which point they retreat downstream toward the swamp (area S58).

\subparagraph{S56c: Snapper Saws}
Three horrid plants (snapper saws) hide near the paths.

\subparagraph{S56d: Shambling Mounds}
Four shambling mounds lurk among the thick growth near the path to catch the unwary. After dispatching the mounds, a character who surveys the area and succeeds on a DC 14 Wisdom (Survival) check finds a narrow path to the creatures' den nearby. A heap of rotting vegetation within contains two Energy Cell, a Humanoid skeleton with a jeweled ring worth 1,000 gp, and dozens of finger-sized maggots.

\subparagraph{S56e: Squealer}
This area is the territory of a squealer—a porcine-headed predator unleashed during the crash. The squealer uses the \textbf{nalfeshnee} stat block, but it's a Monstrosity and has a 30-foot climbing speed instead of a flying speed. The creature hide in the treetops, mimicking the cries of distressed Beasts to attract prey. A character who hears the sounds recognizes them as imitations with a successful DC 15 Wisdom (Insight) check.

The squealer's lair is an artificial burrow close by. Inside is a red key card, a \textit{concussion grenade}, and four gems worth 100 gp each.

\subparagraph{S56f: Dew Drinkers}
Three horrid plants (dew drinkers) hide near the path.

\subparagraph{S56g: Globe Palms}
This area is rich with globe palms—tall, slender trees topped with coconut-sized globes. These membranous fruits are filled with an odorous liquid that attracts nearby predators. When a creature brushes past a globe palm, the tree drops several globes, which burst on impact.

Characters who move through this area must succeed on a DC 16 Dexterity saving throw or become coated with the contents of one of the falling fruits for 1 hour or until the character uses an action to scrape off the fruit's pungent juices. For each character who fails this save, a \textbf{swarm of gibberlings} emerges from an artificial burrow along the west wall; the creatures focus their fury on characters coated in the liquid.

\subparagraph{S56h: Brilliant Fish}
The encounter in this location is identical to the one in area S56b.

\subsubsection{S57: Inner Garden (Green)}
This area is similar to area S56 but has fewer predators due to the presence of the froghemoth (see area S59). It consists of three foliage-lush regions with a swamp to the southeast (area S58).

\subparagraph{S57a: Mutated Frogs}
Dozens of three-eyed frogs croak ominously along the banks of streamlets. The frogs are harmless and scatter if approached.

\subparagraph{S57b: Thri-kreen Forest}
Five \textbf{thri-kreen} and their leader, a psionic thri-kreen (a variant thri-kreen), hide in the canopies near this snaking pathway. The character with the highest passive Wisdom (Perception) score notices the thri-kreen, who are initially indifferent toward the party.

If the characters parley with the thri-kreen, the insectile warriors ask them to destroy the shambling mounds in area S56d in exchange for two Concussion Grenade the thri-kreen have accumulated. The thri-kreen cryptically warn the characters they'll need the explosives to cross the lake. If the characters decline, the thri-kreen disappear into the foliage without another word.

\subparagraph{S57c: Boring Grass Patch}
The grass surrounding this northern path has bristly, corkscrew-shaped leaves. A non-Plant creature that enters the grass for the first time on a turn or starts its turn there takes 5 (1d10) piercing damage and has its speed reduced by 10 feet until the start of its next turn.

Near the streamlet are two horrid plants (snapper saws) that use their tendrils to reel passersby into their bladelike leaves.

\subsubsection{S58: Swamp}
\begin{DndReadAloud}{}
This flooded region has formed a swamp. Gnarled trees sink into the morass, and sagging vines droop from branches just above the scummy, knee-high water. A chorus of chirping insects, hissing reptiles, and ribbiting frogs pervades the area.

\end{DndReadAloud}

Burst pipes have flooded this region, creating a swamp whose waters fluctuate from 1 to 3 feet in depth. The region is \textit{difficult terrain}.

Eight giant leeches (use the \textbf{giant constrictor snake} stat block) dwell in the swamp. Each leech has two eyestalks and a translucent underbelly that reveals a mingling of blood-engorged sacs, peristaltic organs, and smooth muscle tissue. The parasites flatten their bodies to hide beneath the water's edge and ambush passersby. Characters who have a passive Wisdom (Perception) score of 16 or higher notice the leeches.

\subsubsection{S59: Lake}
\includegraphics[width=0.4\textwidth]{images/../5e.tools/img/adventure/QftIS/127-07-009.froghemoth.png}
\begin{DndReadAloud}{}
A murky lake fills the center of the garden. Luminous fish crest the surface of its waters, and tiny, harmless amphibians and lizards roam its verge.

A ten-foot-wide bridge extends over the lake to an islet enclosure covered with lush foliage. Battered robots litter the edge of the lake near the structure.

\end{DndReadAloud}

This body of water was formerly a large, natural aquarium for the enjoyment of upper-echelon personnel, who viewed its contents from a marine observatory under the lake's islet enclosure (area S60). The lake is 100 feet deep and murky. Visibility underwater is limited to 10 feet.

A \textbf{froghemoth elder} hides in the lake. Originally a specimen loosed from captivity, the creature has existed on this level for decades and has grown to massive proportions. The garden is this apex predator's lair.

The froghemoth is hostile toward all other creatures. If the characters attempt to cross the bridge or swim across the lake, the froghemoth springs up from the depths with an ear-splitting croak.

\subparagraph{Robots}
Battered robots litter the banks. Aphelion summoned these robots to the level below, but the froghemoth crushed them before they could reach the islet.

\subparagraph{Silent Treatment}
After the characters encounter the froghemoth, Aphelion doesn't contact or respond to them until they reach the ship's fourth level.

\subparagraph{Treasure}
At the bottom of the lake are fifty gems worth 20 gp each. A character who spends 10 minutes searching the lake bed finds all the gems, as well as a \textit{concussion grenade} and a Humanoid skeleton in a wet suit. The skeleton clutches a waterlogged \textit{laser rifle}. The rifle can fire 10 more shots before its \textit{energy cell} is drained.

\subsubsection{S60: Islet Enclosure (Green)}
\begin{DndReadAloud}{}
A ten-foot-high wall surrounds this botanical enclosure, which teems with bunnyoids—cute little rabbits with nubby protuberances atop their heads. One of the creatures sits on a stump nestled amid a patch of leafy foliage. It briefly looks in your direction and freezes, shivering in fright.

\end{DndReadAloud}

This botanical centerpiece originally housed lovely tropical flora, but it has since become the den of a deceptive predator. The bunnyoid on the stump is the fleshy lure of a \textbf{wolf-in-sheep's-clothing} resting atop the creature's stump-shaped body. The creature uses the lure to lull the characters into a false sense of security and tempt them toward its jaws. When a character comes within 10 feet of the wolf-in-sheep's-clothing or threatens it, the aberration uproots and attacks.

\includegraphics[width=0.4\textwidth]{images/../5e.tools/img/adventure/QftIS/128-07-010.bunnyoid.png}
\subparagraph{Marine Observatory}
Behind a door on the islet that requires a green key card, a stairway spirals down to a 30-foot-wide observation chamber with a 10-foot-high ceiling. The chamber's windows are lined with wide, concave portholes. Characters who look through them see fish, snails, turtles, and other marine life in the lake beyond. The observatory is located 50 feet below the surface of the lake.

If the froghemoth in area S59 is still alive when the characters enter this chamber, it whips at the windows with its tentacles and long, prehensile tongue to get at the tender morsels within.

Each window has Armor Class 13, 10 hit points, and immunity to poison and psychic damage. If a window is destroyed, water floods into the chamber. For each window broken, the water rises at a rate of 1 foot per round, rising at the end of each round until it reaches the chamber's 10-foot-tall ceiling.

\subparagraph{Stairway}
Characters who continue descending the stairway arrive at a card-locked door on the level below. This door, which requires a green key card, leads into area S61. Anti-flooding pumps in the stairwell prevent water from entering that area.

\subsection{Spaceship Locations, Level 4}
The fourth level of the spaceship is the service deck. It is the central node of the ship, the brain from which the ship's core functions derive their processing power—and where they are most vulnerable.

The following locations are keyed to map 7.4.

\includegraphics[width=0.4\textwidth]{images/../5e.tools/img/adventure/QftIS/129-map-7.04-spaceship-level-4-map.png}
\includegraphics[width=0.4\textwidth]{images/../5e.tools/img/adventure/QftIS/130-map-7.04-spaceship-level-4-map-player.png}
\subsubsection{S61: Server Room (Green)}
\begin{DndReadAloud}{}
Boxy, metal towers hum throughout this cold, spacious chamber. Soft blue lights blink across their upright faces, and wires snake from the towers along the floor and ceiling.

"You shouldn't be here," echoes a familiar voice in the room as Aphelion's technological visage—red and unstable, as if shaking with rage—appears on a glassy surface on the tower before you. Sinister red motes on the towers' faces form a sea of artificial eyes that glare menacingly down at you.

\end{DndReadAloud}

The voice belongs to Aphelion, whose consciousness resides on the servers in this dimly lit room. The supercomputer's usually blithe demeanor has turned threatening.

Aphelion thanks the characters for their services. As a reward, the supercomputer divulges the true history of the events leading up to the ship's crash as outlined in this adventure's background, including Aphelion's role in it all. The supercomputer exudes arrogance and a deep disdain for non-machines, gloating about how it manipulated the characters into removing the last obstacle in its path. Aphelion ends its monologue with a chilling revelation: once its shipboard mission is satisfied, the supercomputer plans to rebuild its robot army and set the robots loose on the world to impose its own twisted utopia.

Aphelion is stalling. When the conversation flags or if the characters interrupt the monologue with combat, four combat robots pour out through the stairway access door from which the characters emerged and attack. The robots were previously unable to reach this level but followed the characters after they encountered the froghemoth.

\subparagraph{Servers}
These massive, 20-foot-square pillars contain critical information and the collected history of the society that created the spaceship. They also contain Aphelion's consciousness.

Each server has Armor Class 19; 50 hit points; and immunity to lightning, poison, and psychic damage. Electrical piping and exposed wiring cover each server. A creature that touches a server or hits it with a melee attack takes 10 (3d6) lightning damage and must succeed on a DC 10 Constitution saving throw or have the stunned condition until the start of its next turn.

Aphelion audibly protests any attempts to sabotage the servers by mentioning their contents—if the servers are destroyed, the collective knowledge of an entire civilization dies with them.

\subparagraph{Destroying Aphelion}
If three or more of the servers are destroyed, the rest begin to short-circuit and rapidly overheat, frying the electronics within—and Aphelion along with them.

As Aphelion fades from existence, the computer briefly sings the following coda in an artificial tone:

\begin{DndReadAloud}{}
"What have you done? What have you done! Back to the darkness, Into the void."

\end{DndReadAloud}

After Aphelion dies, any robots still aboard the ship default to their core programming and continue their activities independently of the supercomputer. Aphelion can no longer sway them for or against the characters. Technology throughout the ship, such as the drop tubes, overhead lighting, and other independent machinery, still functions.

\subsubsection{S62: Commander's Bones}
A cobwebbed Humanoid skeleton in a regal uniform is sprawled on the cold floor here. These are the jumbled bones of the ship's commander, whose \textit{laser pistol} and platinum key card are nearby. The pistol can fire 10 shots before its \textit{energy cell} is drained.

\subsubsection{S63: Wheelie Sleds}
A wheelie sled—a battery-operated cargo vehicle designed for conveying passengers, robots, and heavy machinery—rests in each of these areas. A wheelie sled appears as a hovering, rectangular slab with a comfortable chair and an area for storing cargo.

Each wheelie sled measures 10 feet long, 5 feet wide, and 1 foot thick. It has a flying speed of 60 feet and can hover. Its movement is controlled via a speed lever and steering wheel. A wheelie sled can transport up to 2,000 pounds without hindrance. It can carry up to twice this weight, but it moves at half speed if it carries more than its normal capacity.

At full charge, a wheelie sled can operate for 24 hours, after which time it must be charged at one of the plug-in terminals located throughout this level. If removed from the ship, the device ceases to function until returned.

\subsubsection{S64: Stasis Chamber (Green)}
\begin{DndReadAloud}{}
This solemn chamber holds twenty upright, cylindrical pods. The door to each pod includes a glass window, though many are cracked or have shattered, revealing the skeletons strapped within. Cool fog emanates from beneath the pods.

\end{DndReadAloud}

Crew members entered periods of prolonged stasis in this preservation chamber. The room features twenty stasis pods, most of which are damaged or empty; some contain the skeletons of long-dead crew members. Only one survivor remains.

A scientist lies in cryogenic sleep within one of the pods. The pod's interior is coated in a thin layer of frost, but a character who inspects its window sees a Humanoid within. Despite the pod's mysterious construction, opening it is simple. A character who inspects the pod and succeeds on a DC 12 Intelligence (Investigation) check understands how to release the pod's inhabitant.

\subparagraph{Freeing the Scientist}
The stasis pod opens with a hiss of depressurized gas. Inside the pod is a scientist (lawful good, human \textbf{noble}) named Nova. Nova has a shaved head and wears a green jumpsuit. Once freed, Nova takes a moment to regain her composure and footing. She then thanks the characters and questions them about her current location and how long it has been since the crash.

Nova has been in stasis much longer than is recommended. As a result, her memory is foggy, though she remembers shutting down Aphelion and the ship losing control before she stumbled into the tank. Nova can recount hazy details about the ship's odyssey (see this adventure's background), but she remembers little of the advanced society to which she belonged. If the characters fill in the details, Nova is dismayed to hear about the death of her crewmates, fellow scientists, and the passengers.

Nova thanks the characters for their heroism but insists on staying aboard the ship to assess the damage and salvage what's left of the servers. The scientist instinctively recalls how to use the technology aboard the ship and can teach the characters how to use any devices they found.

\subsubsection{S65: Irradiated Rooms}
Apart from their locations, these irradiated rooms are identical to those on the ship's first level (see area S13).

\subsubsection{S66: Hull Breach}
A gaping hole pierces the ship's hull in this area. A dark, narrow tunnel extends from the opening and into the mountainside. If the characters follow the tunnel, it continues for 1 mile before exiting into the Barrier Peaks.

A herd of ten intellect devourers roams this area, preying on creatures that happen on this entrance to the ship. When the characters inspect the breach, the aberrations emerge from air ducts in the surrounding machinery and attack.

\section{Conclusion}
With the defeat of Aphelion, the characters can continue to explore the ship or return to the Barrier Peaks, whether back the way they came or through the hull breach in area S66. Along the way, they might free the ship's last surviving crew member (see area S64) or leave her in cryogenic slumber. Regardless of what they choose, the ship is theirs to explore, though its dangers remain.

The spaceship in this adventure is but one section of a much larger vessel. Other sections—jettisoned during the outbreak to prevent the spread of mold—might exist on this world or another of your choice. Aphelion's reach might extend to one or all these sections, or perhaps a different, more benevolent version of the intelligence oversees those ejected vessels, protecting survivors in stasis until their rescuers arrive.

\includegraphics[width=0.4\textwidth]{images/../5e.tools/img/adventure/QftIS/131-07-011.intellect-devourer.png}
\appendix
\chapter{Magic Items and Technology}
This appendix describes magic items and futuristic technology that appear in the adventures.

\section{Magic Items}
The following magic items are presented in alphabetical order.


\begin{itemize}
\begin{multicols}{3}
\item \textit{Daoud's Wondrous Lanthorn}
\item \textit{Heretic}
\item \textit{Staff of Ruling}
\end{multicols}

\end{itemize}

\section{Technology}
This section presents technological devices and weapons that have special rules. The items are presented in alphabetical order. See the \textit{Dungeon Master's Guide} for more information on \textit{explosives and other futuristic weapons}, such as Laser Pistol and Antimatter Rifle.

On many D\&D worlds, these devices are strange and unfamiliar anomalies, if they exist at all. If you're nervous about including such technology in your game, you can limit the number of Energy Cell the characters discover during their adventures. Most devices are useless without them.

Due to their rarity, these items are priceless.

\includegraphics[width=0.4\textwidth]{images/../5e.tools/img/adventure/QftIS/134-08-003.technology.png}

\begin{itemize}
\begin{multicols}{3}
\item \textit{Antigravity Belt}
\item \textit{Concussion Grenade}
\item \textit{Sleep Grenade}
\item \textit{Needler Pistol}
\item \textit{Paralysis Pistol}
\item \textit{Powered Armor}
\item \textit{Robot Controller}
\end{multicols}

\end{itemize}

\chapter{Creatures}
This appendix describes creatures that appear in the adventures, presenting them in alphabetical order. The \textit{introduction} of the \textit{Monster Manual} explains how to read a creature's stat block.


\begin{itemize}
\begin{multicols}{3}
\item \textbf{Android}
\item \textbf{Barkburr}
\item \textbf{Derro Apprentice}
\item \textbf{Derro Raider}
\item \textbf{Drelnza}
\item \textbf{Froghemoth Elder}
\item \textbf{The Gardener}
\item \textbf{Gibberling}
\item \textbf{Swarm of Gibberlings}
\item \textbf{Champion of Gorm}
\item \textbf{Guardian of Gorm}
\item \textbf{Horrid Plant}
\item \textbf{Leprechaun}
\item \textbf{Champion of Usamigaras}
\item \textbf{Mage of Usamigaras}
\item \textbf{Maschin-i-Bozorg}
\item \textbf{Memory Web}
\item \textbf{Nafas}
\item \textbf{Nafik}
\item \textbf{Pech}
\item \textbf{Combat Robot}
\item \textbf{Worker Robot}
\item \textbf{Sion}
\item \textbf{Tower Hand}
\item \textbf{Tower Sage}
\item \textbf{Vegepygmy Moldmaker}
\item \textbf{Vegepygmy Scavenger}
\item \textbf{Vegepygmy Thorny Hunter}
\item \textbf{Champion of Madarua}
\item \textbf{Warrior of Madarua}
\item \textbf{Wolf-in-Sheep's-Clothing}
\item \textbf{Zargon the Returner}
\end{multicols}

\end{itemize}

\section{Index of Stat Blocks}
The Stat Blocks by Challenge Rating and Type table lists each of the stat blocks in this appendix by its challenge rating and creature type.

\begin{DndTable}[header=Stat Blocks by Challenge Rating and Type]{cXX}

CR & Creature & Type\\
0 & \textbf{Gibberling} & Fiend\\
1/8 & \textbf{Guardian of Gorm} & Humanoid\\
1/8 & \textbf{Mage of Usamigaras} & Humanoid\\
1/8 & \textbf{Warrior of Madarua} & Humanoid\\
1/4 & \textbf{Derro raider} & Aberration\\
1/4 & \textbf{Vegepygmy scavenger} & Plant\\
1/2 & \textbf{Tower hand} & Humanoid\\
1 & \textbf{Derro apprentice} & Aberration\\
1 & \textbf{Tower sage} & Humanoid\\
2 & \textbf{Champion of Gorm} & Humanoid\\
2 & \textbf{Champion of Madarua} & Humanoid\\
2 & \textbf{Champion of Usamigaras} & Humanoid\\
2 & \textbf{Sion} & Humanoid\\
2 & \textbf{Vegepygmy thorny hunter} & Plant\\
3 & \textbf{Barkburr} & Plant\\
3 & \textbf{Swarm of gibberlings} & Fiend\\
3 & \textbf{Vegepygmy moldmaker} & Plant\\
3 & \textbf{Worker robot} & Construct\\
4 & \textbf{Horrid plant} & Plant\\
4 & \textbf{Leprechaun} & Fey\\
4 & \textbf{Memory web} & Aberration\\
4 & \textbf{Pech} & Elemental\\
5 & \textbf{Android} & Construct\\
6 & \textbf{Combat robot} & Construct\\
6 & \textbf{Nafik} & Undead\\
7 & \textbf{Wolf-in-sheep's-clothing} & Plant\\
8 & \textbf{Maschin-i-bozorg} & Construct\\
12 & \textbf{The Gardener} & Fey\\
15 & \textbf{Drelnza} & Undead\\
15 & \textbf{Froghemoth elder} & Aberration\\
17 & \textbf{Zargon the Returner} & Aberration\\
23 & \textbf{Nafas} & Elemental\\
\end{DndTable}

\includegraphics[width=0.4\textwidth]{images/../5e.tools/img/adventure/QftIS/160-09-027.painter-lebrechaun.png}
\chapter{Credits}

\begin{description}
\begin{multicols}{2}

\begin{description}
\item[Lead Designer:]
Justice Ramin Arman
\item[Art Director:]
Emi Tanji
\item[Designers:]
Dan Dillon, Carl Sibley
\item[Rules Developers:]
Jeremy Crawford, Makenzie De Armas, Ron Lundeen, Carl Sibley
\item[Lead Editor:]
Judy Bauer
\item[Editors:]
Eytan Bernstein, Michele Carter, Janica Carter, Laura Hirsbrunner, Sadie Lowry, Patrick Renie
\item[Principle Graphic Designer:]
Trish Yochum
\item[Cover Illustrators:]
John Patrick Gañas, Syd Mills
\item[Cartographers:]
Stacey Allan & William Doyle, Marco Bernardini, Jason A. Engle, Sean Macdonald, Damien Mammoliti, Marc Moureau, Mike Schley
\item[Interior Illustrators:]
Hazem Ameen, Luca Bancone, Mark Behm, Eric Belisle, Olivier Bernard, Zoltan Boros, Zezhou Chen, Daniel Corona, CoupleOfKooks, Axel Defois, Julie Dillon, Olga Drebas, Tomas Duchek, Craig Elliott, Victor Ferraz, Jaqueline Florencio, Jessica Fong, Michele Giorgi, Kevin Glint, Alexandre Honoré, Adrián Ibarra Lugo, Dario Jelusic, Jane Katsubo, Katerina Ladon, Yuliya Litvinova, Titus Lunter, Marie Magny, Dave Melvin, Martin Mottet, Irina Nordsol, One Pixel Brush, Hinchel Or, Alejandro Pacheco, PINDURSKI, Andrea Piparo, Arash Radkia, Noor Rahman, Tooba Rezaei, Cyprien Rousson, Taras Susak, Kamila Szutenberg, Matias Tapia, Brian Valeza, Zuzanna Wuzyk
\item[Concept Art Directors:]
Josh Herman, Kate Irwin, Emi Tanji
\item[Concept Artists:]
One Pixel Brush, Noor Rahman
\item[Consultants:]
Tempest Bradford, Ma'at Crook, Dominique Dickey, Sameer Joseph, Omar Ramadan-Santiago
\item[Project Engineer:]
Cynda Callaway
\item[Imaging Technicians:]
Daniel Corona, Kevin Yee
\item[Producers:]
Bill Benham, Siera Bruggeman, Robert Hawkey, Andy Smith, Dan Tovar, Gabriel Waluconis
\item[Prepress Specialist:]
Jefferson Dunlap
\item[Product Manager:]
Natalie Egan
\end{description}

\section{D\&D Studio}

\begin{description}
\item[Executive Producer:]
Kyle Brink
\item[Game Architects:]
Jeremy Crawford, Christopher Perkins
\item[Design Department:]
Justice Ramin Arman, Makenzie De Armas, Amanda Hamon, Ron Lundeen, Ben Petrisor, Patrick Renie, F. Wesley Schneider, Jason Tondro, James Wyatt
\item[Editorial Department:]
Judy Bauer, Janica Carter, Adrian Ng
\item[Studio Art Director:]
Josh Herman
\item[Art Department:]
Matt Cole, Trystan Falcone, Kate Irwin, Noor Rahman, Emi Tanji, Trish Yochum
\item[Senior Producer:]
Dan Tovar
\item[Producers:]
Bill Benham, Siera Bruggeman, Robert Hawkey, Andy Smith
\item[Senior Director of Product Management:]
Brian Perry
\end{description}

\end{multicols}

\end{description}

\section{Credits from the Original Adventures}

\begin{description}
\begin{multicols}{2}
\section{The Lost City (1982)}

\begin{description}
\item[Design:]
Tom Moldvay
\item[Development:]
Tom Moldvay, Jon Pickens
\item[Editing:]
Harold Johnson, Jon Pickens
\item[Art:]
Jim Holloway, Harry Quinn, Stephen D. Sullivan
\item[Playtesting:]
Dave Cook, Helen Cook, Clint Johnson, Steve Kaszar, Bill Wilkerson, Jeff Wyndham, and the Kent State University Gamer's Guild
\end{description}

\section{When a Star Falls (1984)}

\begin{description}
\item[Storyline:]
Phil Gallagher, Tom Kirby, Graeme Morris
\item[Production/Editing:]
Jim Bambra, Phil Gallagher, Tom Kirby
\item[Design:]
Paul Cockburn, Kim Daniel
\item[Art:]
Jeremy Goodwin
\item[Cartography:]
Paul Ruiz
\end{description}

\section{Beyond the Crystal Cave (1983)}

\begin{description}
\item[Design and Development:]
David J. Browne, Tom Kirby, Graeme Morris
\item[Editing:]
Tom Kirby, Carole Morris, Graeme Morris, Don Turnbull
\item[Art:]
De Leuw, Timothy Truman
\item[Cartography:]
Graeme Morris
\item[Playtesting:]
Jim Bambra, Jeanette Blaaser, Clive F. Booth, Michael W. Brunton, Chris Hall, Bill Howard, Kate Kirby, Gary Kirkham, Steve Mote, Chris Rick, Dave Tant, Don Turnbull, Pat Whitehead
\end{description}

\section{Pharoah (1980)}

\begin{description}
\item[Design:]
Tracy and Laura Hickman
\item[Editing:]
Curtis Smith
\item[Cover Art:]
Jim Holloway
\end{description}

\section{The Lost Caverns of Tsojcanth (1982)}

\begin{description}
\item[Design:]
Gary Gygax
\item[Development:]
Gary Gygax, Allen Hammack, Jon Pickens, Edward G. Sollers
\item[Editing:]
Edward G. Sollers
\item[Art:]
Jim Holloway, Erol Otus, Jeff Easley, Stephen D. Sullivan
\item[Playtesting:]
Jeff Dolphin, Luke Gygax, David Kuntz, Richard Kuntz, Sonny Savage, James M. Ward
\item[Special Thanks:]
Rob Kuntz
\end{description}

\section{Expedition to the Barrier Peaks (1980)}

\begin{description}
\item[Design:]
Gary Gygax
\item[Development:]
Gary Gygax, Allen Hammack, Jon Pickens, Edward G. Sollers
\item[Editing:]
Edward G. Sollers
\item[Art:]
Jim Holloway, Erol Otus, Jeff Easley, Stephen D. Sullivan
\item[Playtesting:]
Jeff Dolphin, Luke Gygax, David Kuntz, Richard Kuntz, Sonny Savage, James M. Ward
\item[Special Thanks:]
Rob Kuntz
\end{description}

\end{multicols}

\end{description}

Additional thanks to the original designers of the adventures in this anthology: David J. Browne, Phil Gallagher, Allen Hammack, Gary Gygax, Laura Hickman, Tracy Hickman, Tom Kirby, Tom Moldvay, Graeme Morris, Jon Pickens, Edward G. Sollers, and others.

\includegraphics[width=0.4\textwidth]{images/../5e.tools/img/adventure/QftIS/161-10-004.infinite-doorway.png}
\includegraphics[width=0.4\textwidth]{images/../5e.tools/img/adventure/QftIS/162-10-002.cover-spread.png}
\includegraphics[width=0.4\textwidth]{images/../5e.tools/img/adventure/QftIS/163-10-003.alt-cover-spread.png}
\chapter*{Printable Statblocks}
\setlist[description]{nolistsep,listparindent=3pt,topsep=6pt,font={\bfseries\sffamily\small}}
\pagestyle{empty}

\begin{DndMonster}[float=t]{Aarakocra}
\DndMonsterType{Medium humanoid (aarakocra), neutral good}
\DndMonsterBasics[
armor-class = {12},
hit-points = {\DndDice{3d8}},
speed = {20 ft., fly 50 ft.}
]
\DndMonsterAbilityScores[
 str = 10,
 dex = 14,
 con = 10,
 int = 11,
 wis = 12,
 cha = 11
]
\DndMonsterDetails[
skills = {Perception +5},
languages = {Auran, Aarakocra},
challenge = 1/4,
]
\DndMonsterAction{Dive Attack}
If the aarakocra is flying and dives at least 30 feet straight toward a target and then hits it with a melee weapon attack, the attack deals an extra 3 (1d6) damage to the target.

\DndMonsterSection{Actions}
\DndMonsterAction{Talon}
\textsl{Melee Weapon Attack:} 4 to hit, reach 5 ft., one target. \textsl{Hit:} 4 (1d4 + 2) slashing damage.

\DndMonsterAction{Javelin}
\textsl{Melee or Ranged Weapon Attack:} 4 to hit, reach 5 ft. or range 30/120 ft., one target. \textsl{Hit:} 5 (1d6 + 2) piercing damage.

\DndMonsterAction{Summon Air Elemental}
Five aarakocra within 30 feet of each other can magically summon an \textbf{air elemental}. Each of the five must use its action and movement on three consecutive turns to perform an aerial dance and must maintain concentration while doing so (as if concentration on a spell). When all five have finished their third turn of the dance, the elemental appears in an unoccupied space within 60 feet of them. It is friendly toward them and obeys their spoken commands. It remains for 1 hour, until it or all its summoners die, or until any of its summoners dismisses it as a bonus action. A summoner can't perform the dance again until it finishes a short rest. When the elemental returns to the Elemental Plane of Air, any aarakocra within 5 feet of it can return with it.

\end{DndMonster}



\begin{DndMonster}[float*=b,width=\textwidth + 8pt]{Aboleth}
\begin{multicols}{2}
\DndMonsterType{Large aberration, lawful evil}
\DndMonsterBasics[
armor-class = {17 (natural armor)},
hit-points = {\DndDice{18d10 + 36}},
speed = {10 ft., swim 40 ft.}
]
\DndMonsterAbilityScores[
 str = 21,
 dex = 9,
 con = 15,
 int = 18,
 wis = 15,
 cha = 18
]
\DndMonsterDetails[
saving-throws = {Con +6, Int +8, Wis +6},
skills = {History +12, Perception +10},
senses = {darkvision 120 ft., passive perception 20},
languages = {Deep Speech, telepathy 120 ft.},
challenge = 10,
]
\DndMonsterAction{Amphibious}
The aboleth can breathe air and water.

\DndMonsterAction{Mucous Cloud}
While underwater, the aboleth is surrounded by transformative mucus. A creature that touches the aboleth or that hits it with a melee attack while within 5 feet of it must make a DC 14 Constitution saving throw. On a failure, the creature is diseased for 1d4 hours. The diseased creature can breathe only underwater.

\DndMonsterAction{Probing Telepathy}
If a creature communicates telepathically with the aboleth, the aboleth learns the creature's greatest desires if the aboleth can see the creature.

\DndMonsterSection{Actions}
\DndMonsterAction{Multiattack}
The aboleth makes three tentacle attacks.

\DndMonsterAction{Tentacle}
\textsl{Melee Weapon Attack:} 9 to hit, reach 10 ft., one target. \textsl{Hit:} 12 (2d6 + 5) bludgeoning damage. If the target is a creature, it must succeed on a DC 14 Constitution saving throw or become diseased. The disease has no effect for 1 minute and can be removed by any magic that cures disease. After 1 minute, the diseased creature's skin becomes translucent and slimy, the creature can't regain hit points unless it is underwater, and the disease can be removed only by \textit{heal} or another disease-curing spell of 6th level or higher. When the creature is outside a body of water, it takes 6 (1d12) acid damage every 10 minutes unless moisture is applied to the skin before 10 minutes have passed.

\DndMonsterAction{Tail}
\textsl{Melee Weapon Attack:} 9 to hit, reach 10 ft., one target. \textsl{Hit:} 15 (3d6 + 5) bludgeoning damage.

\DndMonsterAction{Enslave (3/Day)}
The aboleth targets one creature it can see within 30 feet of it. The target must succeed on a DC 14 Wisdom saving throw or be magically charmed by the aboleth until the aboleth dies or until it is on a different plane of existence from the target. The charmed target is under the aboleth's control and can't take reactions, and the aboleth and the target can communicate telepathically with each other over any distance.

Whenever the charmed target takes damage, the target can repeat the saving throw. On a success, the effect ends. No more than once every 24 hours, the target can also repeat the saving throw when it is at least 1 mile away from the aboleth.

\DndMonsterSection{Legendary Actions}
Aboleth can take 3 legendary actions, choosing from the options below. Only one legendary action can be used at a time and only at the end of another creature's turn. Aboleth regains spent legendary actions at the start of its turn.

\begin{DndMonsterLegendaryActions}
\DndMonsterLegendaryAction{Detect}{The aboleth makes a Wisdom (Perception) check.
}
\DndMonsterLegendaryAction{Tail Swipe}{The aboleth makes one tail attack.
}
\DndMonsterLegendaryAction{Psychic Drain (Costs 2 Actions)}{One creature charmed by the aboleth takes 10 (3d6) psychic damage, and the aboleth regains hit points equal to the damage the creature takes.
}
\end{DndMonsterLegendaryActions}
\DndMonsterSection{Lair Actions}
When fighting inside its lair, an aboleth can invoke the ambient magic to take lair actions. On initiative count 20 (losing initiative ties), the aboleth takes a lair action to cause one of the following effects:


\begin{itemize}
\item The aboleth casts \textit{phantasmal force} (no components required) on any number of creatures it can see within 60 feet of it. While maintaining concentration on this effect, the aboleth can't take other lair actions. If a target succeeds on the saving throw or if the effect ends for it, the target is immune to the aboleth's phantasmal force lair action for the next 24 hours, although such a creature can choose to be affected.
\item Pools of water within 90 feet of the aboleth surge outward in a grasping tide. Any creature on the ground within 20 feet of such a pool must succeed on a DC 14 Strength saving throw or be pulled up to 20 feet into the water and knocked prone. The aboleth can't use this lair action again until it has used a different one.
\item Water in the aboleth's lair magically becomes a conduit for the creature's rage. The aboleth can target any number of creatures it can see in such water within 90 feet of it. A target must succeed on a DC 14 Wisdom saving throw or take 7 (2d6) psychic damage. The aboleth can't use this lair action again until it has used a different one.
\end{itemize}

\DndMonsterSection{Regional Effects}
The region containing an aboleth's lair is warped by the creature's presence, which creates one or more of the following effects:


\begin{itemize}
\item Underground surfaces within 1 mile of the aboleth's lair are slimy and wet and are difficult terrain.
\item Water sources within 1 mile of the lair are supernaturally fouled. Enemies of the aboleth that drink such water vomit it within minutes.
\item As an action, the aboleth can create an illusory image of itself within 1 mile of the lair. The copy can appear at any location the aboleth has seen before or in any location a creature charmed by the aboleth can currently see. Once created, the image lasts for as long as the aboleth maintains concentration, as if concentration on a spell. Although the image is intangible, it looks, sounds, and can move like the aboleth. The aboleth can sense, speak, and use telepathy from the image's position as if present at that position. If the image takes any damage, it disappears.
\end{itemize}

If the aboleth dies, the first two effects fade over the course of 3d10 days.

\end{multicols}
\end{DndMonster}



\begin{DndMonster}[float=t]{Acolyte}
\DndMonsterType{Medium humanoid (any race), any alignment}
\DndMonsterBasics[
armor-class = {10},
hit-points = {\DndDice{2d8}},
speed = {30 ft.}
]
\DndMonsterAbilityScores[
 str = 10,
 dex = 10,
 con = 10,
 int = 10,
 wis = 14,
 cha = 11
]
\DndMonsterDetails[
skills = {Medicine +4, Religion +2},
languages = {any one language (usually Common)},
challenge = 1/4,
]
\DndMonsterAction{Spellcasting}
The acolyte is a 1st-level spellcaster. Its spellcasting ability is Wisdom (spell save DC 12, 4 to hit with spell attacks). The acolyte has following cleric spells prepared:
\begin{DndMonsterSpells}
  \DndMonsterSpellLevel{light}
   \DndMonsterSpellLevel[1][3]{bless}
\end{DndMonsterSpells}
\DndMonsterSection{Actions}
\DndMonsterAction{Club}
\textsl{Melee Weapon Attack:} 2 to hit, reach 5 ft., one target. \textsl{Hit:} 2 (1d4) bludgeoning damage.

\end{DndMonster}



\begin{DndMonster}[float*=b,width=\textwidth + 8pt]{Adult Black Dragon}
\begin{multicols}{2}
\DndMonsterType{Huge dragon, chaotic evil}
\DndMonsterBasics[
armor-class = {19 (natural armor)},
hit-points = {\DndDice{17d12 + 85}},
speed = {40 ft., fly 80 ft., swim 40 ft.}
]
\DndMonsterAbilityScores[
 str = 23,
 dex = 14,
 con = 21,
 int = 14,
 wis = 13,
 cha = 17
]
\DndMonsterDetails[
saving-throws = {Dex +7, Con +10, Wis +6, Cha +8},
skills = {Perception +11, Stealth +7},
damage-immunities = {acid},
senses = {blindsight 60 ft., darkvision 120 ft., passive perception 21},
languages = {Common, Draconic},
challenge = 14,
]
\DndMonsterAction{Amphibious}
The dragon can breathe air and water.

\DndMonsterAction{Legendary Resistance (3/Day)}
If the dragon fails a saving throw, it can choose to succeed instead.

\DndMonsterSection{Actions}
\DndMonsterAction{Multiattack}
The dragon can use its Frightful Presence. It then makes three attacks: one with its bite and two with its claws.

\DndMonsterAction{Bite}
\textsl{Melee Weapon Attack:} 11 to hit, reach 10 ft., one target. \textsl{Hit:} 17 (2d10 + 6) piercing damage plus 4 (1d8) acid damage.

\DndMonsterAction{Claw}
\textsl{Melee Weapon Attack:} 11 to hit, reach 5 ft., one target. \textsl{Hit:} 13 (2d6 + 6) slashing damage.

\DndMonsterAction{Tail}
\textsl{Melee Weapon Attack:} 11 to hit, reach 15 ft., one target. \textsl{Hit:} 15 (2d8 + 6) bludgeoning damage.

\DndMonsterAction{Frightful Presence}
Each creature of the dragon's choice that is within 120 feet of the dragon and aware of it must succeed on a DC 16 Wisdom saving throw or become frightened for 1 minute. A creature can repeat the saving throw at the end of each of its turns, ending the effect on itself on a success. If a creature's saving throw is successful or the effect ends for it, the creature is immune to the dragon's Frightful Presence for the next 24 hours.

\DndMonsterAction{Acid Breath (Recharge 5\textendash 6)}
The dragon exhales acid in a 60-foot line that is 5 feet wide. Each creature in that line must make a DC 18 Dexterity saving throw, taking 54 (12d8) acid damage on a failed save, or half as much damage on a successful one.

\DndMonsterSection{Legendary Actions}
Adult Black Dragon can take 3 legendary actions, choosing from the options below. Only one legendary action can be used at a time and only at the end of another creature's turn. Adult Black Dragon regains spent legendary actions at the start of its turn.

\begin{DndMonsterLegendaryActions}
\DndMonsterLegendaryAction{Detect}{The dragon makes a Wisdom (Perception) check.
}
\DndMonsterLegendaryAction{Tail Attack}{The dragon makes a tail attack.
}
\DndMonsterLegendaryAction{Wing Attack (Costs 2 Actions)}{The dragon beats its wings. Each creature within 10 feet of the dragon must succeed on a DC 19 Dexterity saving throw or take 13 (2d6 + 6) bludgeoning damage and be knocked prone. The dragon can then fly up to half its flying speed.
}
\end{DndMonsterLegendaryActions}
\end{multicols}
\end{DndMonster}



\begin{DndMonster}[float=t]{Air Elemental}
\DndMonsterType{Large elemental, neutral}
\DndMonsterBasics[
armor-class = {15},
hit-points = {\DndDice{12d10 + 24}},
speed = {0 ft., fly 90 ft. \textit{(hover)}}
]
\DndMonsterAbilityScores[
 str = 14,
 dex = 20,
 con = 14,
 int = 6,
 wis = 10,
 cha = 6
]
\DndMonsterDetails[
damage-resistances = {lightning; thunder; bludgeoning, piercing, slashing from nonmagical attacks},
damage-immunities = {poison},
condition-immunities = {exhaustion, grappled, paralyzed, petrified, poisoned, prone, restrained, unconscious},
senses = {darkvision 60 ft., passive perception 10},
languages = {Auran},
challenge = 5,
]
\DndMonsterAction{Air Form}
The elemental can enter a hostile creature's space and stop there. It can move through a space as narrow as 1 inch wide without squeezing.

\DndMonsterSection{Actions}
\DndMonsterAction{Multiattack}
The elemental makes two slam attacks.

\DndMonsterAction{Slam}
\textsl{Melee Weapon Attack:} 8 to hit, reach 5 ft., one target. \textsl{Hit:} 14 (2d8 + 5) bludgeoning damage.

\DndMonsterAction{Whirlwind (Recharge 4\textendash 6)}
Each creature in the elemental's space must make a DC 13 Strength saving throw. On a failure, a target takes 15 (3d8 + 2) bludgeoning damage and is flung up 20 feet away from the elemental in a random direction and knocked prone. If a thrown target strikes an object, such as a wall or floor, the target takes 3 (1d6) bludgeoning damage for every 10 feet it was thrown. If the target is thrown at another creature, that creature must succeed on a DC 13 Dexterity saving throw or take the same damage and be knocked prone.

If the saving throw is successful, the target takes half the bludgeoning damage and isn't flung away or knocked prone.

\end{DndMonster}



\begin{DndMonster}[float=t]{Android}
\DndMonsterType{Medium construct, lawful neutral}
\DndMonsterBasics[
armor-class = {15 (natural armor)},
hit-points = {\DndDice{14d8 + 28}},
speed = {30 ft., fly 30 ft. \textit{(hover; aerialist only)}, swim 30 ft. \textit{(diver only)}}
]
\DndMonsterAbilityScores[
 str = 18,
 dex = 18,
 con = 15,
 int = 12,
 wis = 13,
 cha = 10
]
\DndMonsterDetails[
saving-throws = {Con +5, Wis +4},
skills = {History +4, Perception +7},
damage-resistances = {acid, fire},
damage-immunities = {cold, poison},
condition-immunities = {exhaustion, poisoned},
senses = {blindsight 60 ft. (sentry only), darkvision 60 ft., passive perception 17},
languages = {Common plus the languages spoken by its creator},
challenge = 5,
]
\DndMonsterAction{Design Specialization}
When the android is created, it gains one of six possible designs suited for its role (choose or roll a d6): 1, aerialist; 2, diplomat; 3, diver; 4, duelist; 5, medic; 6, sentry. This design determines certain traits in this stat block.

\DndMonsterAction{Lightning Overload}
When the android takes lightning damage, it must succeed on a DC 10 Constitution saving throw or have the stunned condition until the start of its next turn.

\DndMonsterAction{Spellcasting (Diplomat and Medic Only)}
The android casts one of the following spells, requiring no material components and using Intelligence as the spellcasting ability:
\begin{DndMonsterSpells}
  \DndInnateSpellLevel[2]{Cure Wounds}
\end{DndMonsterSpells}
\DndMonsterSection{Actions}
\DndMonsterAction{Multiattack}
The android makes two Force Strike attacks.

\DndMonsterAction{Force Strike}
\textsl{Melee or Ranged Weapon Attack:} 7 to hit, reach 5 ft. or range 40/120 ft., one target. \textsl{Hit:} 15 (2d10 + 4) force damage. If the target is a Medium or smaller creature, it must succeed on a DC 15 Strength saving throw or have the prone condition.

\DndMonsterSection{Reactions}
\DndMonsterAction{Parry (Duelist Only)}
The android adds 3 to its AC against one melee attack roll that would hit it. To do so, the android must see the attacker.

\end{DndMonster}



\begin{DndMonster}[float=t]{Animated Armor}
\DndMonsterType{Medium construct, unaligned}
\DndMonsterBasics[
armor-class = {18 (natural armor)},
hit-points = {\DndDice{6d8 + 6}},
speed = {25 ft.}
]
\DndMonsterAbilityScores[
 str = 14,
 dex = 11,
 con = 13,
 int = 1,
 wis = 3,
 cha = 1
]
\DndMonsterDetails[
damage-immunities = {poison, psychic},
condition-immunities = {blinded, charmed, deafened, exhaustion, frightened, paralyzed, petrified, poisoned},
senses = {blindsight 60 ft. (blind beyond this radius), passive perception 6},
challenge = 1,
]
\DndMonsterAction{Antimagic Susceptibility}
The armor is incapacitated while in the area of an \textit{antimagic field}. If targeted by \textit{dispel magic}, the armor must succeed on a Constitution saving throw against the caster's spell save DC or fall unconscious for 1 minute.

\DndMonsterAction{False Appearance}
While the armor remains motionless, it is indistinguishable from a normal suit of armor.

\DndMonsterSection{Actions}
\DndMonsterAction{Multiattack}
The armor makes two melee attacks.

\DndMonsterAction{Slam}
\textsl{Melee Weapon Attack:} 4 to hit, reach 5 ft., one target. \textsl{Hit:} 5 (1d6 + 2) bludgeoning damage.

\end{DndMonster}



\begin{DndMonster}[float=t]{Ankheg}
\DndMonsterType{Large monstrosity, unaligned}
\DndMonsterBasics[
armor-class = {14 (natural armor) (11 while prone)},
hit-points = {\DndDice{6d10 + 6}},
speed = {30 ft., burrow 10 ft.}
]
\DndMonsterAbilityScores[
 str = 17,
 dex = 11,
 con = 13,
 int = 1,
 wis = 13,
 cha = 6
]
\DndMonsterDetails[
senses = {darkvision 60 ft., tremorsense 60 ft., passive perception 11},
challenge = 2,
]
\DndMonsterSection{Actions}
\DndMonsterAction{Bite}
\textsl{Melee Weapon Attack:} 5 to hit, reach 5 ft., one target. \textsl{Hit:} 10 (2d6 + 3) slashing damage plus 3 (1d6) acid damage. If the target is a Large or smaller creature, it is grappled (escape DC 13). Until this grapple ends, the ankheg can bite only the grappled creature and has advantage on attack rolls to do so.

\DndMonsterAction{Acid Spray (Recharge 6)}
The ankheg spits acid in a line that is 30 feet long and 5 feet wide, provided that it has no creature grappled. Each creature in that line must make a DC 13 Dexterity saving throw, taking 10 (3d6) acid damage on a failed save, or half as much damage on a successful one.

\end{DndMonster}



\begin{DndMonster}[float=t]{Archmage}
\DndMonsterType{Medium humanoid (any race), any alignment}
\DndMonsterBasics[
armor-class = {12 (15 with \textit{mage armor})},
hit-points = {\DndDice{18d8 + 18}},
speed = {30 ft.}
]
\DndMonsterAbilityScores[
 str = 10,
 dex = 14,
 con = 12,
 int = 20,
 wis = 15,
 cha = 16
]
\DndMonsterDetails[
saving-throws = {Int +9, Wis +6},
skills = {Arcana +13, History +13},
damage-resistances = {damage from spells; bludgeoning, piercing, slashing (from stoneskin)},
languages = {any six languages},
challenge = 12,
]
\DndMonsterAction{Magic Resistance}
The archmage has advantage on saving throws against spells and other magical effects.

\DndMonsterAction{Spellcasting}
The archmage is an 18th-level spellcaster. Its spellcasting ability is Intelligence (spell save DC 17, 9 to hit with spell attacks). The archmage can cast \textit{disguise self} and \textit{invisibility} at will and has the following wizard spells prepared:
\begin{DndMonsterSpells}
  \DndInnateSpellLevel{disguise self}
  \DndMonsterSpellLevel{fire bolt}
   \DndMonsterSpellLevel[1][4]{detect magic}
   \DndMonsterSpellLevel[2][3]{detect thoughts}
   \DndMonsterSpellLevel[3][3]{counterspell}
   \DndMonsterSpellLevel[4][3]{banishment}
   \DndMonsterSpellLevel[5][3]{cone of cold}
   \DndMonsterSpellLevel[6][1]{globe of invulnerability}
   \DndMonsterSpellLevel[7][1]{teleport}
   \DndMonsterSpellLevel[8][1]{mind blank}
\end{DndMonsterSpells}
\DndMonsterSection{Actions}
\DndMonsterAction{Dagger}
\textsl{Melee or Ranged Weapon Attack:} 6 to hit, reach 5 ft. or range 20/60 ft., one target. \textsl{Hit:} 4 (1d4 + 2) piercing damage.

\end{DndMonster}



\begin{DndMonster}[float=t]{Aurumvorax}
\DndMonsterType{Small monstrosity, unaligned}
\DndMonsterBasics[
armor-class = {15 (natural armor)},
hit-points = {\DndDice{8d6 + 8}},
speed = {30 ft., burrow 20 ft.}
]
\DndMonsterAbilityScores[
 str = 14,
 dex = 13,
 con = 12,
 int = 3,
 wis = 12,
 cha = 6
]
\DndMonsterDetails[
saving-throws = {Str +4, Con +3},
skills = {Perception +3, Stealth +3},
condition-immunities = {petrified},
senses = {darkvision 60 ft., passive perception 13},
challenge = 2,
]
\DndMonsterAction{Tunneler}
The aurumvorax can burrow through solid rock and metal at half its burrowing speed and leaves a 5-foot-diameter tunnel in its wake.

\DndMonsterSection{Actions}
\DndMonsterAction{Multiattack}
The aurumvorax makes one Bite attack and two Claw attacks.

\DndMonsterAction{Bite}
\textsl{Melee Weapon Attack:} 4 to hit, reach 5 ft., one target. \textsl{Hit:} 6 (1d8 + 2) piercing damage. If the target is a creature wearing armor of any type, the aurumvorax regains 4 (1d6 + 1) hit points.

\DndMonsterAction{Claw}
\textsl{Melee Weapon Attack:} 5 to hit, reach 5 ft., one target. \textsl{Hit:} 5 (1d6 + 2) slashing damage. If the target is a Medium or smaller creature, it is grappled (escape DC 12). Until this grapple ends, the aurumvorax can't use its Claw attack on another target, and when it moves, it can drag the grappled creature with it, without the aurumvorax's speed being halved.

\end{DndMonster}



\begin{DndMonster}[float=t]{Aurumvorax Den Leader}
\DndMonsterType{Medium monstrosity, unaligned}
\DndMonsterBasics[
armor-class = {16 (natural armor)},
hit-points = {\DndDice{8d8 + 16}},
speed = {40 ft., burrow 20 ft.}
]
\DndMonsterAbilityScores[
 str = 18,
 dex = 14,
 con = 14,
 int = 3,
 wis = 13,
 cha = 8
]
\DndMonsterDetails[
saving-throws = {Str +6, Con +4},
skills = {Perception +3, Stealth +4},
condition-immunities = {petrified},
senses = {darkvision 60 ft., passive perception 13},
challenge = 4,
]
\DndMonsterAction{Pack Leader}
The aurumvorax's allies have advantage on attack rolls while within 10 feet of the aurumvorax, provided it isn't incapacitated.

\DndMonsterAction{Tunneler}
The aurumvorax can burrow through solid rock and metal at half its burrowing speed and leaves a 5-foot-diameter tunnel in its wake.

\DndMonsterSection{Actions}
\DndMonsterAction{Multiattack}
The aurumvorax makes one Bite attack and two Claw attacks.

\DndMonsterAction{Bite}
\textsl{Melee Weapon Attack:} 6 to hit, reach 5 ft., one target. \textsl{Hit:} 13 (2d8 + 4) piercing damage. If the target is a creature wearing armor of any type, the aurumvorax gains one of the following benefits of its choice:

\DndMonsterAction{Frenzy}
The aurumvorax has advantage on attack rolls until start of its next turn.

\DndMonsterAction{Invigorate}
The aurumvorax regains 6 (1d8 + 2) hit points.

\DndMonsterAction{Claw}
\textsl{Melee Weapon Attack:} 6 to hit, reach 5 ft., one target. \textsl{Hit:} 11 (2d6 + 4) slashing damage. If the target is a Large or smaller creature, it is grappled (escape DC 14). Until this grapple ends, the aurumvorax can't use its Claw attack on another target, and when it moves, it can drag the grappled creature with it, without the aurumvorax's speed being halved.

\end{DndMonster}



\begin{DndMonster}[float=t]{Awakened Shrub}
\DndMonsterType{Small plant, unaligned}
\DndMonsterBasics[
armor-class = {9},
hit-points = {\DndDice{3d6}},
speed = {20 ft.}
]
\DndMonsterAbilityScores[
 str = 3,
 dex = 8,
 con = 11,
 int = 10,
 wis = 10,
 cha = 6
]
\DndMonsterDetails[
damage-vulnerabilities = {fire},
damage-resistances = {piercing},
languages = {one language known by its creator},
challenge = 0,
]
\DndMonsterAction{False Appearance}
While the shrub remains motionless, it is indistinguishable from a normal shrub.

\DndMonsterSection{Actions}
\DndMonsterAction{Rake}
\textsl{Melee Weapon Attack:} 1 to hit, reach 5 ft., one target. \textsl{Hit:} 1 (1d4 - 1) slashing damage.

\end{DndMonster}



\begin{DndMonster}[float=t]{Awakened Tree}
\DndMonsterType{Huge plant, unaligned}
\DndMonsterBasics[
armor-class = {13 (natural armor)},
hit-points = {\DndDice{7d12 + 14}},
speed = {20 ft.}
]
\DndMonsterAbilityScores[
 str = 19,
 dex = 6,
 con = 15,
 int = 10,
 wis = 10,
 cha = 7
]
\DndMonsterDetails[
damage-vulnerabilities = {fire},
damage-resistances = {bludgeoning, piercing},
languages = {one language known by its creator},
challenge = 2,
]
\DndMonsterAction{False Appearance}
While the tree remains motionless, it is indistinguishable from a normal tree.

\DndMonsterSection{Actions}
\DndMonsterAction{Slam}
\textsl{Melee Weapon Attack:} 6 to hit, reach 10 ft., one target. \textsl{Hit:} 14 (3d6 + 4) bludgeoning damage.

\end{DndMonster}



\begin{DndMonster}[float=t]{Badger}
\DndMonsterType{Tiny beast, unaligned}
\DndMonsterBasics[
armor-class = {10},
hit-points = {\DndDice{1d4 + 1}},
speed = {20 ft., burrow 5 ft.}
]
\DndMonsterAbilityScores[
 str = 4,
 dex = 11,
 con = 12,
 int = 2,
 wis = 12,
 cha = 5
]
\DndMonsterDetails[
senses = {darkvision 30 ft., passive perception 11},
challenge = 0,
]
\DndMonsterAction{Keen Smell}
The badger has advantage on Wisdom (Perception) checks that rely on smell.

\DndMonsterSection{Actions}
\DndMonsterAction{Bite}
\textsl{Melee Weapon Attack:} 2 to hit, reach 5 ft., one target. \textsl{Hit:} 1 piercing damage.

\end{DndMonster}



\begin{DndMonster}[float=t]{Bandit}
\DndMonsterType{Medium humanoid (any race), any non-lawful alignment}
\DndMonsterBasics[
armor-class = {12 (\textit{leather armor})},
hit-points = {\DndDice{2d8 + 2}},
speed = {30 ft.}
]
\DndMonsterAbilityScores[
 str = 11,
 dex = 12,
 con = 12,
 int = 10,
 wis = 10,
 cha = 10
]
\DndMonsterDetails[
languages = {any one language (usually Common)},
challenge = 1/8,
]
\DndMonsterSection{Actions}
\DndMonsterAction{Scimitar}
\textsl{Melee Weapon Attack:} 3 to hit, reach 5 ft., one target. \textsl{Hit:} 4 (1d6 + 1) slashing damage.

\DndMonsterAction{Light Crossbow}
\textsl{Ranged Weapon Attack:} 3 to hit, range 80/320 ft., one target. \textsl{Hit:} 5 (1d8 + 1) piercing damage.

\end{DndMonster}



\begin{DndMonster}[float=t]{Bandit Captain}
\DndMonsterType{Medium humanoid (any race), any non-lawful alignment}
\DndMonsterBasics[
armor-class = {15 (studded leather armor)},
hit-points = {\DndDice{10d8 + 20}},
speed = {30 ft.}
]
\DndMonsterAbilityScores[
 str = 15,
 dex = 16,
 con = 14,
 int = 14,
 wis = 11,
 cha = 14
]
\DndMonsterDetails[
saving-throws = {Str +4, Dex +5, Wis +2},
skills = {Athletics +4, Deception +4},
languages = {any two languages},
challenge = 2,
]
\DndMonsterSection{Actions}
\DndMonsterAction{Multiattack}
The captain makes three melee attacks: two with its scimitar and one with its dagger. Or the captain makes two ranged attacks with its daggers.

\DndMonsterAction{Scimitar}
\textsl{Melee Weapon Attack:} 5 to hit, reach 5 ft., one target. \textsl{Hit:} 6 (1d6 + 3) slashing damage.

\DndMonsterAction{Dagger}
\textsl{Melee or Ranged Weapon Attack:} 5 to hit, reach 5 ft. or range 20/60 ft., one target. \textsl{Hit:} 5 (1d4 + 3) piercing damage.

\DndMonsterSection{Reactions}
\DndMonsterAction{Parry}
The captain adds 2 to its AC against one melee attack that would hit it. To do so, the captain must see the attacker and be wielding a melee weapon.

\end{DndMonster}



\begin{DndMonster}[float=t]{Banshee}
\DndMonsterType{Medium undead, chaotic evil}
\DndMonsterBasics[
armor-class = {12},
hit-points = {\DndDice{13d8}},
speed = {0 ft., fly 40 ft. \textit{(hover)}}
]
\DndMonsterAbilityScores[
 str = 1,
 dex = 14,
 con = 10,
 int = 12,
 wis = 11,
 cha = 17
]
\DndMonsterDetails[
saving-throws = {Wis +2, Cha +5},
damage-resistances = {acid; fire; lightning; thunder; bludgeoning, piercing, slashing from nonmagical attacks},
damage-immunities = {cold, necrotic, poison},
condition-immunities = {charmed, exhaustion, frightened, grappled, paralyzed, petrified, poisoned, prone, restrained},
senses = {darkvision 60 ft., passive perception 10},
languages = {Common, Elvish},
challenge = 4,
]
\DndMonsterAction{Detect Life}
The banshee can magically sense the presence of living creatures up to 5 miles away that aren't undead or constructs. She knows the general direction they're in but not their exact locations.

\DndMonsterAction{Incorporeal Movement}
The banshee can move through other creatures and objects as if they were difficult terrain. She takes 5 (1d10) force damage if she ends her turn inside an object.

\DndMonsterSection{Actions}
\DndMonsterAction{Corrupting Touch}
\textsl{Melee Spell Attack:} 4 to hit, reach 5 ft., one target. \textsl{Hit:} 12 (3d6 + 2) necrotic damage.

\DndMonsterAction{Horrifying Visage}
Each non-undead creature within 60 feet of the banshee that can see her must succeed on a DC 13 Wisdom saving throw or be frightened for 1 minute. A frightened target can repeat the saving throw at the end of each of its turns, with disadvantage if the banshee is within line of sight, ending the effect on itself on a success. If a target's saving throw is successful or the effect ends for it, the target is immune to the banshee's Horrifying Visage for the next 24 hours.

\DndMonsterAction{Wail (1/Day)}
The banshee releases a mournful wail, provided that she isn't in sunlight. This wail has no effect on constructs and undead. All other creatures within 30 feet of her that can hear her must make a DC 13 Constitution saving throw. On a failure, a creature drops to 0 hit points. On a success, a creature takes 10 (3d6) psychic damage.

\end{DndMonster}



\begin{DndMonster}[float=t]{Barkburr}
\DndMonsterType{Small plant, unaligned}
\DndMonsterBasics[
armor-class = {16 (natural armor)},
hit-points = {\DndDice{8d6 + 24}},
speed = {10 ft., climb 10 ft.}
]
\DndMonsterAbilityScores[
 str = 16,
 dex = 6,
 con = 16,
 int = 1,
 wis = 15,
 cha = 1
]
\DndMonsterDetails[
skills = {Athletics +5},
damage-immunities = {psychic},
condition-immunities = {blinded, charmed, deafened, frightened},
senses = {blindsight 60 ft. (can't see beyond this radius), passive perception 12},
challenge = 3,
]
\DndMonsterAction{False Appearance}
If the barkburr is motionless at the start of combat, it has advantage on its initiative roll. Moreover, if a creature hasn't observed the barkburr move or act, that creature must succeed on a DC 18 Intelligence (Investigation) check to discern that the barkburr is animate.

\DndMonsterAction{Springing Leap}
With or without a running start, the barkburr's high jump is up to 15 feet, and its long jump is up to 30 feet. The barkburr's jumps can exceed its speed if its speed isn't 0.

\DndMonsterSection{Actions}
\DndMonsterAction{Multiattack}
The barkburr makes two Poison Barb attacks and uses Lignify if able.

\DndMonsterAction{Poison Barb}
\textsl{Melee Weapon Attack:} 5 to hit, reach 5 ft., one creature. \textsl{Hit:} 5 (1d4 + 3) piercing damage plus 7 (2d6) poison damage, and the barkburr attaches to the target. While the barkburr is attached, it can't make Poison Barb attacks, and the target has the restrained condition as its body begins to transform into wood.

An attached barkburr can detach itself by spending 5 feet of its movement on its turn. A creature that can reach the barkburr, including the target, can use its action to detach the barkburr by making a successful DC 13 Strength (Athletics) check.

\DndMonsterAction{Lignify}
The barkburr targets the creature it is attached to, and the target must make a DC 13 Constitution saving throw. On a failed save, the target has the petrified condition until freed by the \textit{Greater Restoration} spell or another effect, except it turns into a tree instead of stone. Any equipment the target is wearing or carrying is absorbed into the tree's bark.

\end{DndMonster}



\begin{DndMonster}[float=t]{Barlgura}
\DndMonsterType{Large fiend (demon), chaotic evil}
\DndMonsterBasics[
armor-class = {15 (natural armor)},
hit-points = {\DndDice{8d10 + 24}},
speed = {40 ft., climb 40 ft.}
]
\DndMonsterAbilityScores[
 str = 18,
 dex = 15,
 con = 16,
 int = 7,
 wis = 14,
 cha = 9
]
\DndMonsterDetails[
saving-throws = {Dex +5, Con +6},
skills = {Perception +5, Stealth +5},
damage-resistances = {cold, fire, lightning},
damage-immunities = {poison},
condition-immunities = {poisoned},
senses = {blindsight 30 ft., darkvision 120 ft., passive perception 15},
languages = {Abyssal, telepathy 120 ft.},
challenge = 5,
]
\DndMonsterAction{Reckless}
At the start of its turn, the barlgura can gain advantage on all melee weapon attack rolls it makes during that turn, but attack rolls against it have advantage until the start of its next turn.

\DndMonsterAction{Running Leap}
The barlgura's long jump is up to 40 feet and its high jump is up to 20 feet when it has a running start.

\DndMonsterAction{Innate Spellcasting}
The barlgura's spellcasting ability is Wisdom (spell save DC 13). The barlgura can innately cast the following spells, requiring no material components:
\begin{DndMonsterSpells}
  \DndInnateSpellLevel[1]{entangle}
  \DndInnateSpellLevel[2]{disguise self}
\end{DndMonsterSpells}
\DndMonsterSection{Actions}
\DndMonsterAction{Multiattack}
The barlgura makes three attacks: one with its bite and two with its fists.

\DndMonsterAction{Bite}
\textsl{Melee Weapon Attack:} 7 to hit, reach 5 ft., one target. \textsl{Hit:} 11 (2d6 + 4) piercing damage.

\DndMonsterAction{Fist}
\textsl{Melee Weapon Attack:} 7 to hit, reach 5 ft., one target. \textsl{Hit:} 9 (1d10 + 4) bludgeoning damage.

\end{DndMonster}



\begin{DndMonster}[float=t]{Basilisk}
\DndMonsterType{Medium monstrosity, unaligned}
\DndMonsterBasics[
armor-class = {15 (natural armor)},
hit-points = {\DndDice{8d8 + 16}},
speed = {20 ft.}
]
\DndMonsterAbilityScores[
 str = 16,
 dex = 8,
 con = 15,
 int = 2,
 wis = 8,
 cha = 7
]
\DndMonsterDetails[
senses = {darkvision 60 ft., passive perception 9},
challenge = 3,
]
\DndMonsterAction{Petrifying Gaze}
If a creature starts its turn within 30 feet of the basilisk and the two of them can see each other, the basilisk can force the creature to make a DC 12 Constitution saving throw if the basilisk isn't incapacitated. On a failed save, the creature magically begins to turn to stone and is restrained. It must repeat the saving throw at the end of its next turn. On a success, the effect ends. On a failure, the creature is petrified until freed by the  \textit{greater restoration} spell or other magic.

A creature that isn't surprised can avert its eyes to avoid the saving throw at the start of its turn. If it does so, it can't see the basilisk until the start of its next turn, when it can avert its eyes again. If it looks at the basilisk in the meantime, it must immediately make the save.

If the basilisk sees its reflection within 30 feet of it in bright light, it mistakes itself for a rival and targets itself with its gaze.

\DndMonsterSection{Actions}
\DndMonsterAction{Bite}
\textsl{Melee Weapon Attack:} 5 to hit, reach 5 ft., one target. \textsl{Hit:} 10 (2d6 + 3) piercing damage plus 7 (2d6) poison damage.

\end{DndMonster}


\clearpage

\begin{DndMonster}[float=t]{Bearded Devil}
\DndMonsterType{Medium fiend (devil), lawful evil}
\DndMonsterBasics[
armor-class = {13 (natural armor)},
hit-points = {\DndDice{8d8 + 16}},
speed = {30 ft.}
]
\DndMonsterAbilityScores[
 str = 16,
 dex = 15,
 con = 15,
 int = 9,
 wis = 11,
 cha = 11
]
\DndMonsterDetails[
saving-throws = {Str +5, Con +4, Wis +2},
damage-resistances = {cold; bludgeoning, piercing, slashing from nonmagical attacks that aren't silvered},
damage-immunities = {fire, poison},
condition-immunities = {poisoned},
senses = {darkvision 120 ft., passive perception 10},
languages = {Infernal, telepathy 120 ft.},
challenge = 3,
]
\DndMonsterAction{Devil's Sight}
Magical darkness doesn't impede the devil's darkvision.

\DndMonsterAction{Magic Resistance}
The devil has advantage on saving throws against spells and other magical effects.

\DndMonsterAction{Steadfast}
The devil can't be frightened while it can see an allied creature within 30 feet of it.

\DndMonsterSection{Actions}
\DndMonsterAction{Multiattack}
The devil makes two attacks: one with its beard and one with its glaive.

\DndMonsterAction{Beard}
\textsl{Melee Weapon Attack:} 5 to hit, reach 5 ft., one creature. \textsl{Hit:} 6 (1d8 + 2) piercing damage, and the target must succeed on a DC 12 Constitution saving throw or be poisoned for 1 minute. While poisoned in this way, the target can't regain hit points. The target can repeat the saving throw at the end of each of its turns, ending the effect on itself on a success.

\DndMonsterAction{Glaive}
\textsl{Melee Weapon Attack:} 5 to hit, reach 10 ft., one target. \textsl{Hit:} 8 (1d10 + 3) slashing damage. If the target is a creature other than an undead or a construct, it must succeed on a DC 12 Constitution saving throw or lose 5 (1d10) hit points at the start of each of its turns due to an infernal wound. Each time the devil hits the wounded target with this attack, the damage dealt by the wound increases by 5 (1d10). Any creature can take an action to stanch the wound with a successful DC 12 Wisdom (Medicine) check. The wound also closes if the target receives magical healing.

\end{DndMonster}



\begin{DndMonster}[float=t]{Behir}
\DndMonsterType{Huge monstrosity, neutral evil}
\DndMonsterBasics[
armor-class = {17 (natural armor)},
hit-points = {\DndDice{16d12 + 64}},
speed = {50 ft., climb 40 ft.}
]
\DndMonsterAbilityScores[
 str = 23,
 dex = 16,
 con = 18,
 int = 7,
 wis = 14,
 cha = 12
]
\DndMonsterDetails[
skills = {Perception +6, Stealth +7},
damage-immunities = {lightning},
senses = {darkvision 90 ft., passive perception 16},
languages = {Draconic},
challenge = 11,
]
\DndMonsterSection{Actions}
\DndMonsterAction{Multiattack}
The behir makes two attacks: one with its bite and one to constrict.

\DndMonsterAction{Bite}
\textsl{Melee Weapon Attack:} 10 to hit, reach 10 ft., one target. \textsl{Hit:} 22 (3d10 + 6) piercing damage.

\DndMonsterAction{Constrict}
\textsl{Melee Weapon Attack:} 10 to hit, reach 5 ft., one Large or smaller creature. \textsl{Hit:} 17 (2d10 + 6) bludgeoning damage plus 17 (2d10 + 6) slashing damage. The target is grappled (escape DC 16) if the behir isn't already constricting a creature, and the target is restrained until this grapple ends.

\DndMonsterAction{Lightning Breath (Recharge 5\textendash 6)}
The behir exhales a line of lightning that is 20 feet long and 5 feet wide. Each creature in that line must make a DC 16 Dexterity saving throw, taking 66 (12d10) lightning damage on a failed save, or half as much damage on a successful one.

\DndMonsterAction{Swallow}
The behir makes one bite attack against a Medium or smaller target it is grappling. If the attack hits, the target is also swallowed, and the grapple ends. While swallowed, the target is blinded and restrained, it has total cover against attacks and other effects outside the behir, and it takes 21 (6d6) acid damage at the start of each of the behir's turns. A behir can have only one creature swallowed at a time.

If the behir takes 30 damage or more on a single turn from the swallowed creature, the behir must succeed on a DC 14 Constitution saving throw at the end of that turn or regurgitate the creature, which falls prone in a space within 10 feet of the behir. If the behir dies, a swallowed creature is no longer restrained by it and can escape from the corpse by using 15 feet of movement, exiting prone.

\end{DndMonster}



\begin{DndMonster}[float=t]{Black Bear}
\DndMonsterType{Medium beast, unaligned}
\DndMonsterBasics[
armor-class = {11 (natural armor)},
hit-points = {\DndDice{3d8 + 6}},
speed = {40 ft., climb 30 ft.}
]
\DndMonsterAbilityScores[
 str = 15,
 dex = 10,
 con = 14,
 int = 2,
 wis = 12,
 cha = 7
]
\DndMonsterDetails[
skills = {Perception +3},
challenge = 1/2,
]
\DndMonsterAction{Keen Smell}
The bear has advantage on Wisdom (Perception) checks that rely on smell.

\DndMonsterSection{Actions}
\DndMonsterAction{Multiattack}
The bear makes two attacks: one with its bite and one with its claws.

\DndMonsterAction{Bite}
\textsl{Melee Weapon Attack:} 4 to hit, reach 5 ft., one target. \textsl{Hit:} 5 (1d6 + 2) piercing damage.

\DndMonsterAction{Claws}
\textsl{Melee Weapon Attack:} 4 to hit, reach 5 ft., one target. \textsl{Hit:} 7 (2d4 + 2) slashing damage.

\end{DndMonster}



\begin{DndMonster}[float=t]{Black Pudding}
\DndMonsterType{Large ooze, unaligned}
\DndMonsterBasics[
armor-class = {7},
hit-points = {\DndDice{10d10 + 30}},
speed = {20 ft., climb 20 ft.}
]
\DndMonsterAbilityScores[
 str = 16,
 dex = 5,
 con = 16,
 int = 1,
 wis = 6,
 cha = 1
]
\DndMonsterDetails[
damage-immunities = {acid, cold, lightning, slashing},
condition-immunities = {blinded, charmed, deafened, exhaustion, frightened, prone},
senses = {blindsight 60 ft. (blind beyond this radius), passive perception 8},
challenge = 4,
]
\DndMonsterAction{Amorphous}
The pudding can move through a space as narrow as 1 inch wide without squeezing.

\DndMonsterAction{Corrosive Form}
A creature that touches the pudding or hits it with a melee attack while within 5 feet of it takes 4 (1d8) acid damage. Any nonmagical weapon made of metal or wood that hits the pudding corrodes. After dealing damage, the weapon takes a permanent and cumulative −1 penalty to damage rolls. If its penalty drops to −5, the weapon is destroyed. Nonmagical ammunition made of metal or wood that hits the pudding is destroyed after dealing damage. The pudding can eat through 2-inch-thick, nonmagical wood or metal in 1 round.

\DndMonsterAction{Spider Climb}
The pudding can climb difficult surfaces, including upside down on ceilings, without needing to make an ability check.

\DndMonsterSection{Actions}
\DndMonsterAction{Pseudopod}
\textsl{Melee Weapon Attack:} 5 to hit, reach 5 ft., one target. \textsl{Hit:} 6 (1d6 + 3) bludgeoning damage plus 18 (4d8) acid damage. In addition, nonmagical armor worn by the target is partly dissolved and takes a permanent and cumulative −1 penalty to the AC it offers. The armor is destroyed if the penalty reduces its AC to 10.

\DndMonsterSection{Reactions}
\DndMonsterAction{Split}
When a pudding that is Medium or larger is subjected to lightning or slashing damage, it splits into two new puddings if it has at least 10 hit points. Each new pudding has hit points equal to half the original pudding's, rounded down. New puddings are one size smaller than the original pudding.

\end{DndMonster}



\begin{DndMonster}[float=t]{Blue Slaad}
\DndMonsterType{Large aberration, chaotic neutral}
\DndMonsterBasics[
armor-class = {15 (natural armor)},
hit-points = {\DndDice{13d10 + 52}},
speed = {30 ft.}
]
\DndMonsterAbilityScores[
 str = 20,
 dex = 15,
 con = 18,
 int = 7,
 wis = 7,
 cha = 9
]
\DndMonsterDetails[
skills = {Perception +1},
damage-resistances = {acid, cold, fire, lightning, thunder},
senses = {darkvision 60 ft., passive perception 11},
languages = {Slaad, telepathy 60 ft.},
challenge = 7,
]
\DndMonsterAction{Magic Resistance}
The slaad has advantage on saving throws against spells and other magical effects.

\DndMonsterAction{Regeneration}
The slaad regains 10 hit points at the start of its turn if it has at least 1 hit point.

\DndMonsterSection{Actions}
\DndMonsterAction{Multiattack}
The slaad makes three attacks: one with its bite and two with its claws.

\DndMonsterAction{Bite}
\textsl{Melee Weapon Attack:} 8 to hit, reach 5 ft., one target. \textsl{Hit:} 12 (2d6 + 5) piercing damage.

\DndMonsterAction{Claw}
\textsl{Melee Weapon Attack:} 8 to hit, reach 5 ft., one target. \textsl{Hit:} 12 (2d6 + 5) slashing damage. If the target is a humanoid, it must succeed on a DC 15 Constitution saving throw or be infected with a disease called chaos phage. While infected, the target can't regain hit points, and its hit point maximum is reduced by 10 (3d6) every 24 hours. If the disease reduces the target's hit point maximum to 0, the target instantly transforms into a \textbf{red slaad} or, if it has the ability to cast spells of 3rd level or higher, a \textbf{green slaad}. Only a \textit{wish} spell can reverse the transformation.

\end{DndMonster}



\begin{DndMonster}[float=t]{Boar}
\DndMonsterType{Medium beast, unaligned}
\DndMonsterBasics[
armor-class = {11 (natural armor)},
hit-points = {\DndDice{2d8 + 2}},
speed = {40 ft.}
]
\DndMonsterAbilityScores[
 str = 13,
 dex = 11,
 con = 12,
 int = 2,
 wis = 9,
 cha = 5
]
\DndMonsterDetails[
challenge = 1/4,
]
\DndMonsterAction{Charge}
If the boar moves at least 20 feet straight toward a target and then hits it with a tusk attack on the same turn, the target takes an extra 3 (1d6) slashing damage. If the target is a creature, it must succeed on a DC 11 Strength saving throw or be knocked prone.

\DndMonsterAction{Relentless (Recharges after a Short or Long Rest)}
If the boar takes 7 damage or less that would reduce it to 0 hit points, it is reduced to 1 hit point instead.

\DndMonsterSection{Actions}
\DndMonsterAction{Tusk}
\textsl{Melee Weapon Attack:} 3 to hit, reach 5 ft., one target. \textsl{Hit:} 4 (1d6 + 1) slashing damage.

\end{DndMonster}



\begin{DndMonster}[float=t]{Brown Bear}
\DndMonsterType{Large beast, unaligned}
\DndMonsterBasics[
armor-class = {11 (natural armor)},
hit-points = {\DndDice{4d10 + 12}},
speed = {40 ft., climb 30 ft.}
]
\DndMonsterAbilityScores[
 str = 19,
 dex = 10,
 con = 16,
 int = 2,
 wis = 13,
 cha = 7
]
\DndMonsterDetails[
skills = {Perception +3},
challenge = 1,
]
\DndMonsterAction{Keen Smell}
The bear has advantage on Wisdom (Perception) checks that rely on smell.

\DndMonsterSection{Actions}
\DndMonsterAction{Multiattack}
The bear makes two attacks: one with its bite and one with its claws.

\DndMonsterAction{Bite}
\textsl{Melee Weapon Attack:} 6 to hit, reach 5 ft., one target. \textsl{Hit:} 8 (1d8 + 4) piercing damage.

\DndMonsterAction{Claws}
\textsl{Melee Weapon Attack:} 6 to hit, reach 5 ft., one target. \textsl{Hit:} 11 (2d6 + 4) slashing damage.

\end{DndMonster}



\begin{DndMonster}[float=t]{Bugbear}
\DndMonsterType{Medium humanoid (goblinoid), chaotic evil}
\DndMonsterBasics[
armor-class = {16 (\textit{hide armor})},
hit-points = {\DndDice{5d8 + 5}},
speed = {30 ft.}
]
\DndMonsterAbilityScores[
 str = 15,
 dex = 14,
 con = 13,
 int = 8,
 wis = 11,
 cha = 9
]
\DndMonsterDetails[
skills = {Stealth +6, Survival +2},
senses = {darkvision 60 ft., passive perception 10},
languages = {Common, Goblin},
challenge = 1,
]
\DndMonsterAction{Brute}
A melee weapon deals one extra die of its damage when the bugbear hits with it (included in the attack).

\DndMonsterAction{Surprise Attack}
If the bugbear surprises a creature and hits it with an attack during the first round of combat, the target takes an extra 7 (2d6) damage from the attack.

\DndMonsterSection{Actions}
\DndMonsterAction{Morningstar}
\textsl{Melee Weapon Attack:} 4 to hit, reach 5 ft., one target. \textsl{Hit:} 11 (2d8 + 2) piercing damage.

\DndMonsterAction{Javelin}
\textsl{Melee or Ranged Weapon Attack:} 4 to hit, reach 5 ft. or range 30/120 ft., one target. \textsl{Hit:} 9 (2d6 + 2) piercing damage in melee or 5 (1d6 + 2) piercing damage at range.

\end{DndMonster}



\begin{DndMonster}[float=t]{Bulette}
\DndMonsterType{Large monstrosity, unaligned}
\DndMonsterBasics[
armor-class = {17 (natural armor)},
hit-points = {\DndDice{9d10 + 45}},
speed = {40 ft., burrow 40 ft.}
]
\DndMonsterAbilityScores[
 str = 19,
 dex = 11,
 con = 21,
 int = 2,
 wis = 10,
 cha = 5
]
\DndMonsterDetails[
skills = {Perception +6},
senses = {darkvision 60 ft., tremorsense 60 ft., passive perception 16},
challenge = 5,
]
\DndMonsterAction{Standing Leap}
The bulette's long jump is up to 30 feet and its high jump is up to 15 feet, with or without a running start.

\DndMonsterSection{Actions}
\DndMonsterAction{Bite}
\textsl{Melee Weapon Attack:} 7 to hit, reach 5 ft., one target. \textsl{Hit:} 30 (4d12 + 4) piercing damage.

\DndMonsterAction{Deadly Leap}
If the bulette jumps at least 15 feet as part of its movement, it can then use this action to land on its feet in a space that contains one or more other creatures. Each of those creatures must succeed on a DC 16 Strength or Dexterity saving throw (target's choice) or be knocked prone and take 14 (3d6 + 4) bludgeoning damage plus 14 (3d6 + 4) slashing damage. On a successful save, the creature takes only half the damage, isn't knocked prone, and is pushed 5 feet out of the bulette's space into an unoccupied space of the creature's choice. If no unoccupied space is within range, the creature instead falls prone in the bulette's space.

\end{DndMonster}



\begin{DndMonster}[float=t]{Camel}
\DndMonsterType{Large beast, unaligned}
\DndMonsterBasics[
armor-class = {9},
hit-points = {\DndDice{2d10 + 4}},
speed = {50 ft.}
]
\DndMonsterAbilityScores[
 str = 16,
 dex = 8,
 con = 14,
 int = 2,
 wis = 8,
 cha = 5
]
\DndMonsterDetails[
challenge = 1/8,
]
\DndMonsterSection{Actions}
\DndMonsterAction{Bite}
\textsl{Melee Weapon Attack:} 5 to hit, reach 5 ft., one target. \textsl{Hit:} 2 (1d4) bludgeoning damage.

\end{DndMonster}



\begin{DndMonster}[float=t]{Carrion Crawler}
\DndMonsterType{Large monstrosity, unaligned}
\DndMonsterBasics[
armor-class = {13 (natural armor)},
hit-points = {\DndDice{6d10 + 18}},
speed = {30 ft., climb 30 ft.}
]
\DndMonsterAbilityScores[
 str = 14,
 dex = 13,
 con = 16,
 int = 1,
 wis = 12,
 cha = 5
]
\DndMonsterDetails[
skills = {Perception +3},
senses = {darkvision 60 ft., passive perception 13},
challenge = 2,
]
\DndMonsterAction{Keen Smell}
The carrion crawler has advantage on Wisdom (Perception) checks that rely on smell.

\DndMonsterAction{Spider Climb}
The carrion crawler can climb difficult surfaces, including upside down on ceilings, without needing to make an ability check.

\DndMonsterSection{Actions}
\DndMonsterAction{Multiattack}
The carrion crawler makes two attacks: one with its tentacles and one with its bite.

\DndMonsterAction{Tentacles}
\textsl{Melee Weapon Attack:} 8 to hit, reach 10 ft., one creature. \textsl{Hit:} 4 (1d4 + 2) poison damage, and the target must succeed on a DC 13 Constitution saving throw or be poisoned for 1 minute. Until this poison ends, the target is paralyzed. The target can repeat the saving throw at the end of each of its turns, ending the poison on itself on a success.

\DndMonsterAction{Bite}
\textsl{Melee Weapon Attack:} 4 to hit, reach 5 ft., one target. \textsl{Hit:} 7 (2d4 + 2) piercing damage.

\end{DndMonster}



\begin{DndMonster}[float=t]{Centaur}
\DndMonsterType{Large monstrosity, neutral good}
\DndMonsterBasics[
armor-class = {12},
hit-points = {\DndDice{6d10 + 12}},
speed = {50 ft.}
]
\DndMonsterAbilityScores[
 str = 18,
 dex = 14,
 con = 14,
 int = 9,
 wis = 13,
 cha = 11
]
\DndMonsterDetails[
skills = {Athletics +6, Perception +3, Survival +3},
languages = {Elvish, Sylvan},
challenge = 2,
]
\DndMonsterAction{Charge}
If the centaur moves at least 30 feet straight toward a target and then hits it with a pike attack on the same turn, the target takes an extra 10 (3d6) piercing damage.

\DndMonsterSection{Actions}
\DndMonsterAction{Multiattack}
The centaur makes two attacks: one with its pike and one with its hooves or two with its longbow.

\DndMonsterAction{Pike}
\textsl{Melee Weapon Attack:} 6 to hit, reach 10 ft., one target. \textsl{Hit:} 9 (1d10 + 4) piercing damage.

\DndMonsterAction{Hooves}
\textsl{Melee Weapon Attack:} 6 to hit, reach 5 ft., one target. \textsl{Hit:} 11 (2d6 + 4) bludgeoning damage.

\DndMonsterAction{Longbow}
\textsl{Ranged Weapon Attack:} 4 to hit, range 150/600 ft., one target. \textsl{Hit:} 6 (1d8 + 2) piercing damage.

\end{DndMonster}



\begin{DndMonster}[float=t]{Champion of Gorm}
\DndMonsterType{Medium humanoid, lawful good}
\DndMonsterBasics[
armor-class = {19 (\textit{splint armor})},
hit-points = {\DndDice{6d8 + 6}},
speed = {30 ft.}
]
\DndMonsterAbilityScores[
 str = 17,
 dex = 10,
 con = 12,
 int = 12,
 wis = 15,
 cha = 16
]
\DndMonsterDetails[
saving-throws = {Str +5, Wis +4},
skills = {Athletics +5, Insight +4, Religion +3},
languages = {Common},
challenge = 2,
]
\DndMonsterAction{Brave}
The champion has advantage on saving throws against the frightened condition.

\DndMonsterSection{Actions}
\DndMonsterAction{Multiattack}
The champion makes two Lightning Mace attacks or three Handaxe attacks.

\DndMonsterAction{Lightning Mace}
\textsl{Melee Weapon Attack:} 5 to hit, reach 5 ft., one target. \textsl{Hit:} 6 (1d6 + 3) piercing damage plus 5 (2d4) lightning damage.

\DndMonsterAction{Handaxe}
\textsl{Melee or Ranged Weapon Attack:} 5 to hit, reach 5 ft. or range 20/60 ft., one target. \textsl{Hit:} 6 (1d6 + 3) slashing damage.

\DndMonsterSection{Bonus Actions}
\DndMonsterAction{Aura of Resilience (1/Day)}
The champion exudes an aura of ghostly lightning that fills a 10-foot-radius sphere centered on itself. While this aura is active, the champion and each creature of its choice within the aura have advantage on saving throws. The aura moves with the champion and lasts for 1 minute, until the champion has the incapacitated condition, or until the champion uses another bonus action to end the aura.

\end{DndMonster}



\begin{DndMonster}[float=t]{Champion of Madarua}
\DndMonsterType{Medium humanoid, chaotic good}
\DndMonsterBasics[
armor-class = {13 (\textit{hide armor})},
hit-points = {\DndDice{6d8 + 12}},
speed = {30 ft.}
]
\DndMonsterAbilityScores[
 str = 17,
 dex = 12,
 con = 15,
 int = 12,
 wis = 12,
 cha = 14
]
\DndMonsterDetails[
saving-throws = {Str +5, Dex +3},
skills = {Acrobatics +3, Religion +3},
languages = {Common},
challenge = 2,
]
\DndMonsterSection{Actions}
\DndMonsterAction{Multiattack}
The champion makes three Longsword attacks.

\DndMonsterAction{Longsword}
\textsl{Melee Weapon Attack:} 5 to hit, reach 5 ft., one target. \textsl{Hit:} 7 (1d8 + 3) piercing damage, or 8 (1d10 + 3) piercing damage if used with two hands.

\DndMonsterSection{Bonus Actions}
\DndMonsterAction{Aura of Fury (1/Day)}
The champion summons an aura of ghostly animals that fills a 10-foot-radius sphere centered on itself. While this aura is active, the champion and all its allies in the aura have advantage on attack rolls. The aura moves with the champion and lasts for 1 minute, until the champion has the incapacitated condition, or until the champion uses another bonus action to end the aura.

\DndMonsterSection{Reactions}
\DndMonsterAction{Parry}
The champion adds 2 to its AC against one melee attack that would hit it. To do so, the champion must see the attacker and be wielding a melee weapon.

\end{DndMonster}



\begin{DndMonster}[float=t]{Champion of Usamigaras}
\DndMonsterType{Medium humanoid, chaotic neutral}
\DndMonsterBasics[
armor-class = {12 (15 with \textit{mage armor})},
hit-points = {\DndDice{6d8 + 6}},
speed = {30 ft.}
]
\DndMonsterAbilityScores[
 str = 10,
 dex = 14,
 con = 12,
 int = 17,
 wis = 14,
 cha = 15
]
\DndMonsterDetails[
saving-throws = {Int +5, Cha +4},
skills = {Arcana +5, Deception +6, Sleight Of Hand +4},
languages = {Common},
challenge = 2,
]
\DndMonsterAction{Spellcasting}
The champion casts one of the following spells, using Intelligence as the spellcasting ability (spell save DC 13):
\begin{DndMonsterSpells}
  \DndInnateSpellLevel{Light}
  \DndInnateSpellLevel[2]{Charm Person}
  \DndInnateSpellLevel[1]{Crown of Madness}
\end{DndMonsterSpells}
\DndMonsterSection{Actions}
\DndMonsterAction{Multiattack}
The champion makes two Arcane Burst attacks or makes one Arcane Burst attack and uses Spellcasting.

\DndMonsterAction{Arcane Burst}
\textsl{Melee or Ranged Spell Attack:} 5 to hit, reach 5 ft. or range 120 ft., one target. \textsl{Hit:} 8 (1d10 + 3) force damage.

\DndMonsterSection{Bonus Actions}
\DndMonsterAction{Aura of Deception (1/Day)}
The champion emits an aura of ghostly illusions that fills a 10-foot-radius sphere centered on itself. While this aura is active, creatures have disadvantage on attack rolls against the champion and any of the champion's allies that are in the aura. The aura moves with the champion and lasts for 1 minute, until the champion has the incapacitated condition, or until the champion uses another bonus action to end the aura.

\end{DndMonster}



\begin{DndMonster}[float=t]{Chasme}
\DndMonsterType{Large fiend (demon), chaotic evil}
\DndMonsterBasics[
armor-class = {15 (natural armor)},
hit-points = {\DndDice{13d10 + 13}},
speed = {20 ft., fly 60 ft.}
]
\DndMonsterAbilityScores[
 str = 15,
 dex = 15,
 con = 12,
 int = 11,
 wis = 14,
 cha = 10
]
\DndMonsterDetails[
saving-throws = {Dex +5, Wis +5},
skills = {Perception +5},
damage-resistances = {cold, fire, lightning},
damage-immunities = {poison},
condition-immunities = {poisoned},
senses = {blindsight 10 ft., darkvision 120 ft., passive perception 15},
languages = {Abyssal, telepathy 120 ft.},
challenge = 6,
]
\DndMonsterAction{Drone}
The chasme produces a horrid droning sound to which demons are immune. Any other creature that starts its turn with in 30 feet of the chasme must succeed on a DC 12 Constitution saving throw or fall unconscious for 10 minutes. A creature that can't hear the drone automatically succeeds on the save. The effect on the creature ends if it takes damage or if another creature takes an action to splash it with holy water. If a creature's saving throw is successful or the effect ends for it, it is immune to the drone for the next 24 hours.

\DndMonsterAction{Magic Resistance}
The chasme has advantage on saving throws against spells and other magical effects.

\DndMonsterAction{Spider Climb}
The chasme can climb difficult surfaces, including upside down on ceilings, without needing to make an ability check.

\DndMonsterSection{Actions}
\DndMonsterAction{Proboscis}
\textsl{Melee Weapon Attack:} 5 to hit, reach 5 ft., one creature. \textsl{Hit:} 16 (4d6 + 2) piercing damage plus 24 (7d6) necrotic damage, and the target's hit point maximum is reduced by an amount equal to the necrotic damage taken. If this effect reduces a creature's hit point maximum to 0, the creature dies. This reduction to a creature's hit point maximum lasts until the creature finishes a long rest or until it is affected by a spell like  \textit{greater restoration}.

\end{DndMonster}



\begin{DndMonster}[float=t]{Chimera}
\DndMonsterType{Large monstrosity, chaotic evil}
\DndMonsterBasics[
armor-class = {14 (natural armor)},
hit-points = {\DndDice{12d10 + 48}},
speed = {30 ft., fly 60 ft.}
]
\DndMonsterAbilityScores[
 str = 19,
 dex = 11,
 con = 19,
 int = 3,
 wis = 14,
 cha = 10
]
\DndMonsterDetails[
skills = {Perception +8},
senses = {darkvision 60 ft., passive perception 18},
languages = {understands Draconic but can't speak},
challenge = 6,
]
\DndMonsterSection{Actions}
\DndMonsterAction{Multiattack}
The chimera makes three attacks: one with its bite, one with its horns, and one with its claws. When its fire breath is available, it can use the breath in place of its bite or horns.

\DndMonsterAction{Bite}
\textsl{Melee Weapon Attack:} 7 to hit, reach 5 ft., one target. \textsl{Hit:} 11 (2d6 + 4) piercing damage.

\DndMonsterAction{Horns}
\textsl{Melee Weapon Attack:} 7 to hit, reach 5 ft., one target. \textsl{Hit:} 10 (1d12 + 4) bludgeoning damage.

\DndMonsterAction{Claws}
\textsl{Melee Weapon Attack:} 7 to hit, reach 5 ft., one target. \textsl{Hit:} 11 (2d6 + 4) slashing damage.

\DndMonsterAction{Fire Breath (Recharge 5\textendash 6)}
The dragon head exhales fire in a 15-foot cone. Each creature in that area must make a DC 15 Dexterity saving throw, taking 31 (7d8) fire damage on a failed save, or half as much damage on a successful one.

\end{DndMonster}



\begin{DndMonster}[float=t]{Chuul}
\DndMonsterType{Large aberration, chaotic evil}
\DndMonsterBasics[
armor-class = {16 (natural armor)},
hit-points = {\DndDice{11d10 + 33}},
speed = {30 ft., swim 30 ft.}
]
\DndMonsterAbilityScores[
 str = 19,
 dex = 10,
 con = 16,
 int = 5,
 wis = 11,
 cha = 5
]
\DndMonsterDetails[
skills = {Perception +4},
damage-immunities = {poison},
condition-immunities = {poisoned},
senses = {darkvision 60 ft., passive perception 14},
languages = {understands Deep Speech but can't speak},
challenge = 4,
]
\DndMonsterAction{Amphibious}
The chuul can breathe air and water.

\DndMonsterAction{Sense Magic}
The chuul senses magic within 120 feet of it at will. This trait otherwise works like the \textit{detect magic} spell but isn't itself magical.

\DndMonsterSection{Actions}
\DndMonsterAction{Multiattack}
The chuul makes two pincer attacks. If the chuul is grappling a creature, the chuul can also use its tentacles once.

\DndMonsterAction{Pincer}
\textsl{Melee Weapon Attack:} 6 to hit, reach 10 ft., one target. \textsl{Hit:} 11 (2d6 + 4) bludgeoning damage. The target is grappled (escape DC 14) if it is a Large or smaller creature and the chuul doesn't have two other creatures grappled.

\DndMonsterAction{Tentacles}
One creature grappled by the chuul must succeed on a DC 13 Constitution saving throw or be poisoned for 1 minute. Until this poison ends, the target is paralyzed. The target can repeat the saving throw at the end of each of its turns, ending the effect on itself on a success.

\end{DndMonster}



\begin{DndMonster}[float=t]{Clay Golem}
\DndMonsterType{Large construct, unaligned}
\DndMonsterBasics[
armor-class = {14 (natural armor)},
hit-points = {\DndDice{14d10 + 56}},
speed = {20 ft.}
]
\DndMonsterAbilityScores[
 str = 20,
 dex = 9,
 con = 18,
 int = 3,
 wis = 8,
 cha = 1
]
\DndMonsterDetails[
damage-immunities = {acid; poison; psychic; bludgeoning, piercing, slashing from nonmagical attacks that aren't adamantine},
condition-immunities = {charmed, exhaustion, frightened, paralyzed, petrified, poisoned},
senses = {darkvision 60 ft., passive perception 9},
languages = {understands the languages of its creator but can't speak},
challenge = 9,
]
\DndMonsterAction{Acid Absorption}
Whenever the golem is subjected to acid damage, it takes no damage and instead regains a number of hit points equal to the acid damage dealt.

\DndMonsterAction{Berserk}
Whenever the golem starts its turn with 60 hit points or fewer, roll a d6. On a 6, the golem goes berserk. On each of its turns while berserk, the golem attacks the nearest creature it can see. If no creature is near enough to move to and attack, the golem attacks an object, with preference for an object smaller than itself. Once the golem goes berserk, it continues to do so until it is destroyed or regains all its hit points.

\DndMonsterAction{Immutable Form}
The golem is immune to any spell or effect that would alter its form.

\DndMonsterAction{Magic Resistance}
The golem has advantage on saving throws against spells and other magical effects.

\DndMonsterAction{Magic Weapons}
The golem's weapon attacks are magical.

\DndMonsterSection{Actions}
\DndMonsterAction{Multiattack}
The golem makes two slam attacks.

\DndMonsterAction{Slam}
\textsl{Melee Weapon Attack:} 8 to hit, reach 5 ft., one target. \textsl{Hit:} 16 (2d10 + 5) bludgeoning damage. If the target is a creature, it must succeed on a DC 15 Constitution saving throw or have its hit point maximum reduced by an amount equal to the damage taken. The target dies if this attack reduces its hit point maximum to 0. The reduction lasts until removed by the  \textit{greater restoration} spell or other magic.

\DndMonsterAction{Haste (Recharge 5\textendash 6)}
Until the end of its next turn, the golem magically gains a +2 bonus to its AC, has advantage on Dexterity saving throws, and can use its slam attack as a bonus action.

\end{DndMonster}



\begin{DndMonster}[float=t]{Cloaker}
\DndMonsterType{Large aberration, chaotic neutral}
\DndMonsterBasics[
armor-class = {14 (natural armor)},
hit-points = {\DndDice{12d10 + 12}},
speed = {10 ft., fly 40 ft.}
]
\DndMonsterAbilityScores[
 str = 17,
 dex = 15,
 con = 12,
 int = 13,
 wis = 12,
 cha = 14
]
\DndMonsterDetails[
skills = {Stealth +5},
senses = {darkvision 60 ft., passive perception 11},
languages = {Deep Speech, Undercommon},
challenge = 8,
]
\DndMonsterAction{Damage Transfer}
While attached to a creature, the cloaker takes only half the damage dealt to it (rounded down). and that creature takes the other half.

\DndMonsterAction{False Appearance}
While the cloaker remains motionless without its underside exposed, it is indistinguishable from a dark leather cloak.

\DndMonsterAction{Light Sensitivity}
While in bright light, the cloaker has disadvantage on attack rolls and Wisdom (Perception) checks that rely on sight.

\DndMonsterSection{Actions}
\DndMonsterAction{Multiattack}
The cloaker makes two attacks: one with its bite and one with its tail.

\DndMonsterAction{Bite}
\textsl{Melee Weapon Attack:} 6 to hit, reach 5 ft., one creature. \textsl{Hit:} 10 (2d6 + 3) piercing damage, and if the target is Large or smaller, the cloaker attaches to it. If the cloaker has advantage against the target, the cloaker attaches to the target's head, and the target is blinded and unable to breathe while the cloaker is attached. While attached, the cloaker can make this attack only against the target and has advantage on the attack roll. The cloaker can detach itself by spending 5 feet of its movement. A creature, including the target, can take its action to detach the cloaker by succeeding on a DC 16 Strength check.

\DndMonsterAction{Tail}
\textsl{Melee Weapon Attack:} 6 to hit, reach 10 ft., one creature. \textsl{Hit:} 7 (1d8 + 3) slashing damage.

\DndMonsterAction{Moan}
Each creature within 60 feet of the cloaker that can hear its moan and that isn't an aberration must succeed on a DC 13 Wisdom saving throw or become frightened until the end of the cloaker's next turn. If a creature's saving throw is successful, the creature is immune to the cloaker's moan for the next 24 hours.

\DndMonsterAction{Phantasms (Recharges after a Short or Long Rest)}
The cloaker magically creates three illusory duplicates of itself if it isn't in bright light. The duplicates move with it and mimic its actions, shifting position so as to make it impossible to track which cloaker is the real one. If the cloaker is ever in an area of bright light, the duplicates disappear.

Whenever any creature targets the cloaker with an attack or a harmful spell while a duplicate remains, that creature rolls randomly to determine whether it targets the cloaker or one of the duplicates. A creature is unaffected by this magical effect if it can't see or if it relies on senses other than sight.

A duplicate has the cloaker's AC and uses its saving throws. If an attack hits a duplicate, or if a duplicate fails a saving throw against an effect that deals damage, the duplicate disappears.

\end{DndMonster}


\clearpage

\begin{DndMonster}[float=t]{Combat Robot}
\DndMonsterType{Medium construct, lawful neutral}
\DndMonsterBasics[
armor-class = {15 (natural armor)},
hit-points = {\DndDice{15d8 + 45}},
speed = {40 ft.}
]
\DndMonsterAbilityScores[
 str = 20,
 dex = 14,
 con = 17,
 int = 10,
 wis = 15,
 cha = 10
]
\DndMonsterDetails[
saving-throws = {Con +6, Cha +3},
skills = {Athletics +8, Intimidation +6, Perception +5},
damage-resistances = {acid, fire},
damage-immunities = {cold, poison},
condition-immunities = {exhaustion, poisoned},
senses = {darkvision 120 ft., passive perception 15},
languages = {Common plus the languages spoken by its creator},
challenge = 6,
]
\DndMonsterAction{Lightning Overload}
When the robot takes lightning damage, it must succeed on a DC 10 Constitution saving throw or have the stunned condition until the start of its next turn.

\DndMonsterSection{Actions}
\DndMonsterAction{Multiattack}
The robot makes two Tentacle attacks or three Laser Beam attacks.

\DndMonsterAction{Tentacle}
\textsl{Melee Weapon Attack:} 8 to hit, reach 10 ft., one target. \textsl{Hit:} 10 (1d10 + 5) bludgeoning damage. If the target is a Medium or smaller creature, it has the grappled condition (escape DC 16). While grappled, the target also has the restrained condition. The robot has two tentacles, each of which can grapple one target.

\DndMonsterAction{Laser Beam}
\textsl{Ranged Weapon Attack:} 5 to hit, range 120 ft., one target. \textsl{Hit:} 16 (4d6 + 2) radiant damage.

\DndMonsterAction{Grenade Launcher (Recharge 5\textendash 6)}
The robot fires a grenade at a point it can see within 120 feet of itself. The grenade explodes in a 20-foot-radius sphere centered on that point, creating one of the following effects (robot's choice):

\DndMonsterAction{Concussion Grenade}
Each creature in the sphere must make a DC 15 Dexterity saving throw, taking 21 (6d6) force damage on a failed save or half as much damage on a successful one.

\DndMonsterAction{Sleep Grenade}
Each creature in the sphere must succeed on a DC 15 Constitution saving throw or have the unconscious condition for 1 hour. The condition ends on a creature early if the creature takes damage or if another creature uses an action to shake it awake.

\DndMonsterSection{Bonus Actions}
\DndMonsterAction{Emergency Speed (2/Day)}
The robot takes the Dash or Disengage action.

\end{DndMonster}



\begin{DndMonster}[float=t]{Commoner}
\DndMonsterType{Medium humanoid (any race), any alignment}
\DndMonsterBasics[
armor-class = {10},
hit-points = {\DndDice{1d8}},
speed = {30 ft.}
]
\DndMonsterAbilityScores[
 str = 10,
 dex = 10,
 con = 10,
 int = 10,
 wis = 10,
 cha = 10
]
\DndMonsterDetails[
languages = {any one language (usually Common)},
challenge = 0,
]
\DndMonsterSection{Actions}
\DndMonsterAction{Club}
\textsl{Melee Weapon Attack:} 2 to hit, reach 5 ft., one target. \textsl{Hit:} 2 (1d4) bludgeoning damage.

\end{DndMonster}



\begin{DndMonster}[float=t]{Cult Fanatic}
\DndMonsterType{Medium humanoid (any race), any non-good alignment}
\DndMonsterBasics[
armor-class = {13 (\textit{leather armor})},
hit-points = {\DndDice{6d8 + 6}},
speed = {30 ft.}
]
\DndMonsterAbilityScores[
 str = 11,
 dex = 14,
 con = 12,
 int = 10,
 wis = 13,
 cha = 14
]
\DndMonsterDetails[
skills = {Deception +4, Persuasion +4, Religion +2},
languages = {any one language (usually Common)},
challenge = 2,
]
\DndMonsterAction{Dark Devotion}
The fanatic has advantage on saving throws against being charmed or frightened.

\DndMonsterAction{Spellcasting}
The fanatic is a 4th-level spellcaster. Its spellcasting ability is Wisdom (spell save DC 11, 3 to hit with spell attacks). The fanatic has the following cleric spells prepared:
\begin{DndMonsterSpells}
  \DndMonsterSpellLevel{light}
   \DndMonsterSpellLevel[1][4]{command}
   \DndMonsterSpellLevel[2][3]{hold person}
\end{DndMonsterSpells}
\DndMonsterSection{Actions}
\DndMonsterAction{Multiattack}
The fanatic makes two melee attacks.

\DndMonsterAction{Dagger}
\textsl{Melee or Ranged Weapon Attack:} 4 to hit, reach 5 ft. or range 20/60 ft., one creature. \textsl{Hit:} 4 (1d4 + 2) piercing damage.

\end{DndMonster}



\begin{DndMonster}[float=t]{Cultist}
\DndMonsterType{Medium humanoid (any race), any non-good alignment}
\DndMonsterBasics[
armor-class = {12 (\textit{leather armor})},
hit-points = {\DndDice{2d8}},
speed = {30 ft.}
]
\DndMonsterAbilityScores[
 str = 11,
 dex = 12,
 con = 10,
 int = 10,
 wis = 11,
 cha = 10
]
\DndMonsterDetails[
skills = {Deception +2, Religion +2},
languages = {any one language (usually Common)},
challenge = 1/8,
]
\DndMonsterAction{Dark Devotion}
The cultist has advantage on saving throws against being charmed or frightened.

\DndMonsterSection{Actions}
\DndMonsterAction{Scimitar}
\textsl{Melee Weapon Attack:} 3 to hit, reach 5 ft., one creature. \textsl{Hit:} 4 (1d6 + 1) slashing damage.

\end{DndMonster}



\begin{DndMonster}[float=t]{Dao}
\DndMonsterType{Large elemental, neutral evil}
\DndMonsterBasics[
armor-class = {18 (natural armor)},
hit-points = {\DndDice{15d10 + 105}},
speed = {30 ft., burrow 30 ft., fly 30 ft.}
]
\DndMonsterAbilityScores[
 str = 23,
 dex = 12,
 con = 24,
 int = 12,
 wis = 13,
 cha = 14
]
\DndMonsterDetails[
saving-throws = {Int +5, Wis +5, Cha +6},
condition-immunities = {petrified},
senses = {darkvision 120 ft., passive perception 11},
languages = {Terran},
challenge = 11,
]
\DndMonsterAction{Earth Glide}
The dao can burrow through nonmagical, unworked earth and stone. While doing so, the dao doesn't disturb the material it moves through.

\DndMonsterAction{Elemental Demise}
If the dao dies, its body disintegrates into crystalline powder, leaving behind only equipment the dao was wearing or carrying.

\DndMonsterAction{Sure-Footed}
The dao has advantage on Strength and Dexterity saving throws made against effects that would knock it prone.

\DndMonsterAction{Innate Spellcasting}
The dao's innate spellcasting ability is Charisma (spell save DC 14, 6 to hit with spell attacks). It can innately cast the following spells, requiring no material components:
\begin{DndMonsterSpells}
  \DndInnateSpellLevel{detect evil and good}
  \DndInnateSpellLevel[3]{passwall}
  \DndInnateSpellLevel[1]{conjure elemental}
\end{DndMonsterSpells}
\DndMonsterSection{Actions}
\DndMonsterAction{Multiattack}
The Dao makes two fist attacks or two maul attacks.

\DndMonsterAction{Fist}
\textsl{Melee Weapon Attack:} 10 to hit, reach 5 ft., one target. \textsl{Hit:} 15 (2d8 + 6) bludgeoning damage.

\DndMonsterAction{Maul}
\textsl{Melee Weapon Attack:} 10 to hit, reach 5 ft., one target. \textsl{Hit:} 20 (4d6 + 6) bludgeoning damage. If the target is a Huge or smaller creature, it must succeed on a DC 18 Strength check or be knocked prone.

\end{DndMonster}



\begin{DndMonster}[float=t]{Deep Gnome (Svirfneblin)}
\DndMonsterType{Small humanoid (gnome), neutral good}
\DndMonsterBasics[
armor-class = {15 (\textit{chain shirt})},
hit-points = {\DndDice{3d6 + 6}},
speed = {20 ft.}
]
\DndMonsterAbilityScores[
 str = 15,
 dex = 14,
 con = 14,
 int = 12,
 wis = 10,
 cha = 9
]
\DndMonsterDetails[
skills = {Investigation +3, Perception +2, Stealth +4},
senses = {darkvision 120 ft., passive perception 12},
languages = {Gnomish, Terran, Undercommon},
challenge = 1/2,
]
\DndMonsterAction{Stone Camouflage}
The gnome has advantage on Dexterity (Stealth) checks made to hide in rocky terrain.

\DndMonsterAction{Gnome Cunning}
The gnome has advantage on Intelligence, Wisdom, and Charisma saving throws against magic.

\DndMonsterAction{Innate Spellcasting}
The gnome's innate spellcasting ability is Intelligence (spell save DC 11). It can innately cast the following spells, requiring no material components:
\begin{DndMonsterSpells}
  \DndInnateSpellLevel{nondetection}
  \DndInnateSpellLevel[1]{blindness/deafness}
\end{DndMonsterSpells}
\DndMonsterSection{Actions}
\DndMonsterAction{War Pick}
\textsl{Melee Weapon Attack:} 4 to hit, reach 5 ft., one target. \textsl{Hit:} 6 (1d8 + 2) piercing damage.

\DndMonsterAction{Poisoned Dart}
\textsl{Ranged Weapon Attack:} 4 to hit, range 30/120 ft., one creature. \textsl{Hit:} 4 (1d4 + 2) piercing damage, and the target must succeed on a DC 12 Constitution saving throw or be poisoned for 1 minute. The target can repeat the saving throw at the end of each of its turns, ending the effect on itself on a success

\end{DndMonster}



\begin{DndMonster}[float=t]{Derro Apprentice}
\DndMonsterType{Small aberration, chaotic evil}
\DndMonsterBasics[
armor-class = {13 (\textit{leather armor})},
hit-points = {\DndDice{5d6 + 5}},
speed = {30 ft.}
]
\DndMonsterAbilityScores[
 str = 9,
 dex = 14,
 con = 12,
 int = 11,
 wis = 5,
 cha = 12
]
\DndMonsterDetails[
skills = {Stealth +4},
senses = {darkvision 120 ft., passive perception 7},
languages = {Dwarvish, Undercommon},
challenge = 1,
]
\DndMonsterAction{Magic Resistance}
The derro has advantage on saving throws against spells and other magical effects.

\DndMonsterAction{Sunlight Sensitivity}
While in sunlight, the derro has disadvantage on attack rolls.

\DndMonsterAction{Spellcasting}
The derro casts one of the following spells, using Charisma as the spellcasting ability (spell save DC 11):
\begin{DndMonsterSpells}
  \DndInnateSpellLevel{Message}
  \DndInnateSpellLevel[]{Charm Person}
\end{DndMonsterSpells}
\DndMonsterSection{Actions}
\DndMonsterAction{Chaos Blast}
\textsl{Melee or Ranged Spell Attack:} 3 to hit, reach 5 ft. or range 60 ft., one target. \textsl{Hit:} 10 (3d6) damage. Roll a d4 to determine the damage type: 1, acid; 2, cold; 3, fire; 4, lightning.

\DndMonsterAction{Force Burst (Recharge 4\textendash 6)}
Raw arcane magic bursts out from the derro. Each creature within 10 feet of it must make a DC 11 Strength saving throw. On a failed save, the creature takes 7 (2d6) force damage and has the prone condition. On a successful save, the creature takes half as much damage only.

\end{DndMonster}



\begin{DndMonster}[float=t]{Derro Raider}
\DndMonsterType{Small aberration, chaotic evil}
\DndMonsterBasics[
armor-class = {12 (\textit{leather armor})},
hit-points = {\DndDice{3d6 + 6}},
speed = {30 ft.}
]
\DndMonsterAbilityScores[
 str = 14,
 dex = 12,
 con = 14,
 int = 11,
 wis = 5,
 cha = 9
]
\DndMonsterDetails[
skills = {Athletics +4, Stealth +3},
senses = {darkvision 120 ft., passive perception 7},
languages = {Dwarvish, Undercommon},
challenge = 1/4,
]
\DndMonsterAction{Magic Resistance}
The derro has advantage on saving throws against spells and other magical effects.

\DndMonsterAction{Sunlight Sensitivity}
While in sunlight, the derro has disadvantage on attack rolls.

\DndMonsterSection{Actions}
\DndMonsterAction{Hooked Spear}
\textsl{Melee Weapon Attack:} 4 to hit, reach 5 ft., one target. \textsl{Hit:} 5 (1d6 + 2) piercing damage. If the target is a Medium or smaller creature, the derro can choose to deal no damage, and instead the target has the prone condition.

\DndMonsterAction{Throwing Hammer}
\textsl{Melee or Ranged Weapon Attack:} 4 to hit, reach 5 ft. or range 20/60 ft. \textsl{Hit:} 5 (1d6 + 2) bludgeoning damage.

\end{DndMonster}



\begin{DndMonster}[float=t]{Derwyth}
\DndMonsterType{Medium humanoid (elf), neutral good}
\DndMonsterBasics[
armor-class = {11 (16 with \textit{barkskin})},
hit-points = {\DndDice{5d8 + 5}},
speed = {30 ft.}
]
\DndMonsterAbilityScores[
 str = 10,
 dex = 12,
 con = 13,
 int = 12,
 wis = 15,
 cha = 11
]
\DndMonsterDetails[
skills = {Medicine +4, Nature +3, Perception +4},
senses = {darkvision 60, passive perception 14},
languages = {Common, Elvish},
challenge = 2,
]
\DndMonsterAction{Tree Shape (2/Day)}
Over the course of 1 minute, Derwyth can magically transform into a Huge or smaller tree and remain in that form for 24 hours or until she ends this transformation early (no action required). Her equipment melds into her new form. While in this form, her Armor Class is 16, she has the incapacitated condition, and she can't move or speak.

\DndMonsterAction{Spellcasting}
Derwyth is a 4th-level spellcaster. Her spellcasting ability is Wisdom (spell save DC 12, 4 to hit with spell attacks). She has the following druid spells prepared:
\begin{DndMonsterSpells}
  \DndMonsterSpellLevel{druidcraft}
   \DndMonsterSpellLevel[1][4]{entangle}
   \DndMonsterSpellLevel[2][3]{animal messenger}
\end{DndMonsterSpells}
\DndMonsterSection{Actions}
\DndMonsterAction{Quarterstaff}
\textsl{Melee Weapon Attack:} 2 to hit (4 to hit with shillelagh), reach 5 ft., one target. \textsl{Hit:} 3 (1d6) bludgeoning damage, 4 (1d8) bludgeoning damage if wielded with two hands, or 6 (1d8 + 2) bludgeoning damage with shillelagh.

\end{DndMonster}



\begin{DndMonster}[float=t]{Deva}
\DndMonsterType{Medium celestial, lawful good}
\DndMonsterBasics[
armor-class = {17 (natural armor)},
hit-points = {\DndDice{16d8 + 64}},
speed = {30 ft., fly 90 ft.}
]
\DndMonsterAbilityScores[
 str = 18,
 dex = 18,
 con = 18,
 int = 17,
 wis = 20,
 cha = 20
]
\DndMonsterDetails[
saving-throws = {Wis +9, Cha +9},
skills = {Insight +9, Perception +9},
damage-resistances = {radiant; bludgeoning, piercing, slashing from nonmagical attacks},
condition-immunities = {charmed, exhaustion, frightened},
senses = {darkvision 120 ft., passive perception 19},
languages = {all, telepathy 120 ft.},
challenge = 10,
]
\DndMonsterAction{Angelic Weapons}
The deva's weapon attacks are magical. When the deva hits with any weapon, the weapon deals an extra 4d8 radiant damage (included in the attack).

\DndMonsterAction{Magic Resistance}
The deva has advantage on saving throws against spells and other magical effects.

\DndMonsterAction{Innate Spellcasting}
The deva's spellcasting ability is Charisma (spell save DC 17). The deva can innately cast the following spells, requiring only verbal components:
\begin{DndMonsterSpells}
  \DndInnateSpellLevel{detect evil and good}
  \DndInnateSpellLevel[1]{commune}
\end{DndMonsterSpells}
\DndMonsterSection{Actions}
\DndMonsterAction{Multiattack}
The deva makes two melee attacks.

\DndMonsterAction{Mace}
\textsl{Melee Weapon Attack:} 8 to hit, reach 5 ft., one target. \textsl{Hit:} 7 (1d6 + 4) bludgeoning damage plus 18 (4d8) radiant damage.

\DndMonsterAction{Healing Touch (3/Day)}
The deva touches another creature. The target magically regains 20 (4d8 + 2) hit points and is freed from any curse, disease, poison, blindness, or deafness.

\DndMonsterAction{Change Shape}
The deva magically polymorphs into a humanoid or beast that has a challenge rating equal to or less than its own, or back into its true form. It reverts to its true form if it dies. Any equipment it is wearing or carrying is absorbed or borne by the new form (the deva's choice).

In a new form, the deva retains its game statistics and ability to speak, but its AC, movement modes, Strength, Dexterity, and special senses are replaced by those of the new form, and it gains any statistics and capabilities (except class features, legendary actions, and lair actions) that the new form has but that it lacks.

\end{DndMonster}



\begin{DndMonster}[float=t]{Dire Wolf}
\DndMonsterType{Large beast, unaligned}
\DndMonsterBasics[
armor-class = {14 (natural armor)},
hit-points = {\DndDice{5d10 + 10}},
speed = {50 ft.}
]
\DndMonsterAbilityScores[
 str = 17,
 dex = 15,
 con = 15,
 int = 3,
 wis = 12,
 cha = 7
]
\DndMonsterDetails[
skills = {Perception +3, Stealth +4},
challenge = 1,
]
\DndMonsterAction{Keen Hearing and Smell}
The wolf has advantage on Wisdom (Perception) checks that rely on hearing or smell.

\DndMonsterAction{Pack Tactics}
The wolf has advantage on an attack roll against a creature if at least one of the wolf's allies is within 5 feet of the creature and the ally isn't incapacitated.

\DndMonsterSection{Actions}
\DndMonsterAction{Bite}
\textsl{Melee Weapon Attack:} 5 to hit, reach 5 ft., one target. \textsl{Hit:} 10 (2d6 + 3) piercing damage. If the target is a creature, it must succeed on a DC 13 Strength saving throw or be knocked prone.

\end{DndMonster}



\begin{DndMonster}[float=t]{Displacer Beast}
\DndMonsterType{Large monstrosity, lawful evil}
\DndMonsterBasics[
armor-class = {13 (natural armor)},
hit-points = {\DndDice{10d10 + 30}},
speed = {40 ft.}
]
\DndMonsterAbilityScores[
 str = 18,
 dex = 15,
 con = 16,
 int = 6,
 wis = 12,
 cha = 8
]
\DndMonsterDetails[
senses = {darkvision 60 ft., passive perception 11},
challenge = 3,
]
\DndMonsterAction{Avoidance}
If the displacer beast is subjected to an effect that allows it to make a saving throw to take only half damage, it instead takes no damage if it succeeds on the saving throw, and only half damage if it fails.

\DndMonsterAction{Displacement}
The displacer beast projects a magical illusion that makes it appear to be standing near its actual location, causing attack rolls against it to have disadvantage. If it is hit by an attack, this trait is disrupted until the end of its next turn. This trait is also disrupted while the displacer beast is incapacitated or has a speed of 0.

\DndMonsterSection{Actions}
\DndMonsterAction{Multiattack}
The displacer beast makes two attacks with its tentacles.

\DndMonsterAction{Tentacle}
\textsl{Melee Weapon Attack:} 6 to hit, reach 10 ft., one target. \textsl{Hit:} 7 (1d6 + 4) bludgeoning damage plus 3 (1d6) piercing damage.

\end{DndMonster}



\begin{DndMonster}[float=t]{Doppelganger}
\DndMonsterType{Medium monstrosity (shapechanger), neutral}
\DndMonsterBasics[
armor-class = {14},
hit-points = {\DndDice{8d8 + 16}},
speed = {30 ft.}
]
\DndMonsterAbilityScores[
 str = 11,
 dex = 18,
 con = 14,
 int = 11,
 wis = 12,
 cha = 14
]
\DndMonsterDetails[
skills = {Deception +6, Insight +3},
condition-immunities = {charmed},
senses = {darkvision 60 ft., passive perception 11},
languages = {Common},
challenge = 3,
]
\DndMonsterAction{Shapechanger}
The doppelganger can use its action to polymorph into a Small or Medium humanoid it has seen, or back into its true form. Its statistics, other than its size, are the same in each form. Any equipment it is wearing or carrying isn't transformed. It reverts to its true form if it dies.

\DndMonsterAction{Ambusher}
In the first round of a combat, the doppelganger has advantage on attack rolls against any creature it surprised.

\DndMonsterAction{Surprise Attack}
If the doppelganger surprises a creature and hits it with an attack during the first round of combat, the target takes an extra 10 (3d6) damage from the attack.

\DndMonsterSection{Actions}
\DndMonsterAction{Multiattack}
The doppelganger makes two melee attacks.

\DndMonsterAction{Slam}
\textsl{Melee Weapon Attack:} 6 to hit, reach 5 ft., one target. \textsl{Hit:} 7 (1d6 + 4) bludgeoning damage.

\DndMonsterAction{Read Thoughts}
The doppelganger magically reads the surface thoughts of one creature within 60 feet of it. The effect can penetrate barriers, but 3 feet of wood or dirt, 2 feet of stone, 2 inches of metal, or a thin sheet of lead blocks it. While the target is in range, the doppelganger can continue reading its thoughts, as long as the doppelganger's concentration isn't broken (as if concentration on a spell). While reading the target's mind, the doppelganger has advantage on Wisdom (Insight) and Charisma (Deception, Intimidation, and Persuasion) checks against the target.

\end{DndMonster}



\begin{DndMonster}[float=t]{Draft Horse}
\DndMonsterType{Large beast, unaligned}
\DndMonsterBasics[
armor-class = {10},
hit-points = {\DndDice{3d10 + 3}},
speed = {40 ft.}
]
\DndMonsterAbilityScores[
 str = 18,
 dex = 10,
 con = 12,
 int = 2,
 wis = 11,
 cha = 7
]
\DndMonsterDetails[
challenge = 1/4,
]
\DndMonsterSection{Actions}
\DndMonsterAction{Hooves}
\textsl{Melee Weapon Attack:} 6 to hit, reach 5 ft., one target. \textsl{Hit:} 9 (2d4 + 4) bludgeoning damage.

\end{DndMonster}



\begin{DndMonster}[float=t]{Drelnza}
\DndMonsterType{Medium undead (vampire), lawful evil}
\DndMonsterBasics[
armor-class = {18 (\textit{plate armor})},
hit-points = {\DndDice{22d8 + 88}},
speed = {30 ft., climb 30 ft.}
]
\DndMonsterAbilityScores[
 str = 18,
 dex = 18,
 con = 18,
 int = 17,
 wis = 15,
 cha = 18
]
\DndMonsterDetails[
saving-throws = {Dex +9, Int +8, Wis +7, Cha +9},
skills = {Arcana +8, Deception +9, Perception +7},
damage-resistances = {necrotic},
condition-immunities = {exhaustion},
senses = {darkvision 120 ft., passive perception 17},
languages = {Abyssal, Common, Giant},
challenge = 15,
]
\DndMonsterAction{Legendary Resistance (3/Day)}
If Drelnza fails a saving throw, she can choose to succeed instead.

\DndMonsterAction{Special Equipment}
Drelnza wields Heretic.

\DndMonsterAction{Spider Climb}
Drelnza can climb difficult surfaces, including upside down on ceilings, without needing to make an ability check.

\DndMonsterAction{Sunlight Hypersensitivity}
Drelnza takes 20 radiant damage when she starts her turn in sunlight. While in sunlight, she has disadvantage on attack rolls and ability checks.

\DndMonsterSection{Actions}
\DndMonsterAction{Multiattack}
Drelnza makes one Bite attack and one Heretic attack.

\DndMonsterAction{Bite}
\textsl{Melee Weapon Attack:} 9 to hit, reach 5 ft., one creature. \textsl{Hit:} 7 (1d6 + 4) piercing damage plus 10 (3d6) necrotic damage. The target's hit point maximum is reduced by an amount equal to the necrotic damage taken, and Drelnza regains hit points equal to that amount. The reduction lasts until the target finishes a long rest. The target dies if this effect reduces its hit point maximum to 0. A Humanoid slain in this way and then buried rises the following night as a vampire spawn under Drelnza's control.

\DndMonsterAction{Heretic}
\textsl{Melee Weapon Attack:} 12 to hit, reach 5 ft. one target. \textsl{Hit:} 40 (6d10 + 7) slashing damage, and the target must succeed on a DC 17 Constitution saving throw or have the paralyzed condition until the start of Drelnza's next turn. Celestials have disadvantage on the saving throw.

\DndMonsterAction{Charm}
Drelnza targets one Humanoid she can see within 30 feet of herself. The target must succeed on a DC 17 Wisdom saving throw or have the charmed condition. While charmed in this way, the target regards Drelnza as a trusted friend to be heeded and protected. The target isn't under Drelnza's control, but it takes her requests and actions in the most favorable way.

Each time Drelnza or her allies do anything harmful to the target, it can repeat the saving throw, ending the effect on itself on a success. Otherwise, the effect lasts 24 hours or until Drelnza is destroyed, is on a different plane of existence than the target, or takes a bonus action to end the effect.

\DndMonsterSection{Reactions}
\DndMonsterAction{Move}
Immediately after a creature Drelnza can see ends its turn, Drelnza moves up to her speed without provoking opportunity attacks.

\DndMonsterAction{Parry}
Drelnza adds 5 to her AC against one melee attack that would hit her. To do so, she must see the attacker and be wielding a melee weapon.

\end{DndMonster}



\begin{DndMonster}[float=t]{Druid}
\DndMonsterType{Medium humanoid (any race), any alignment}
\DndMonsterBasics[
armor-class = {11 (16 with \textit{barkskin})},
hit-points = {\DndDice{5d8 + 5}},
speed = {30 ft.}
]
\DndMonsterAbilityScores[
 str = 10,
 dex = 12,
 con = 13,
 int = 12,
 wis = 15,
 cha = 11
]
\DndMonsterDetails[
skills = {Medicine +4, Nature +3, Perception +4},
languages = {Druidic plus any two languages},
challenge = 2,
]
\DndMonsterAction{Spellcasting}
The druid is a 4th-level spellcaster. Its spellcasting ability is Wisdom (spell save DC 12, 4 to hit with spell attacks). It has the following druid spells prepared:
\begin{DndMonsterSpells}
  \DndMonsterSpellLevel{druidcraft}
   \DndMonsterSpellLevel[1][4]{entangle}
   \DndMonsterSpellLevel[2][3]{animal messenger}
\end{DndMonsterSpells}
\DndMonsterSection{Actions}
\DndMonsterAction{Quarterstaff}
\textsl{Melee Weapon Attack:} 2 to hit (4 to hit with shillelagh), reach 5 ft., one target. \textsl{Hit:} 3 (1d6) bludgeoning damage, 4 (1d8) bludgeoning damage if wielded with two hands, or 6 (1d8 + 2) bludgeoning damage with \textit{shillelagh}.

\end{DndMonster}



\begin{DndMonster}[float=t]{Dryad}
\DndMonsterType{Medium fey, neutral}
\DndMonsterBasics[
armor-class = {11 (16 with \textit{barkskin})},
hit-points = {\DndDice{5d8}},
speed = {30 ft.}
]
\DndMonsterAbilityScores[
 str = 10,
 dex = 12,
 con = 11,
 int = 14,
 wis = 15,
 cha = 18
]
\DndMonsterDetails[
skills = {Perception +4, Stealth +5},
senses = {darkvision 60 ft., passive perception 14},
languages = {Elvish, Sylvan},
challenge = 1,
]
\DndMonsterAction{Magic Resistance}
The dryad has advantage on saving throws against spells and other magical effects.

\DndMonsterAction{Speak with Beasts and Plants}
The dryad can communicate with beasts and plants as if they shared a language.

\DndMonsterAction{Tree Stride}
Once on her turn, the dryad can use 10 feet of her movement to step magically into one living tree within her reach and emerge from a second living tree within 60 feet of the first tree, appearing in an unoccupied space within 5 feet of the second tree. Both trees must be large or bigger.

\DndMonsterAction{Innate Spellcasting}
The dryad's innate spellcasting ability is Charisma (spell save DC 14). The dryad can innately cast the following spells, requiring no material components:
\begin{DndMonsterSpells}
  \DndInnateSpellLevel{druidcraft}
  \DndInnateSpellLevel[3]{entangle}
  \DndInnateSpellLevel[1]{barkskin}
\end{DndMonsterSpells}
\DndMonsterSection{Actions}
\DndMonsterAction{Club}
\textsl{Melee Weapon Attack:} 2 to hit (6 to hit with shillelagh), reach 5 ft., one target. \textsl{Hit:} 2 (1d4) bludgeoning damage, or 8 (1d8 + 4) bludgeoning damage with shillelagh.

\DndMonsterAction{Fey Charm}
The dryad targets one humanoid or beast that she can see within 30 feet of her. If the target can see the dryad, it must succeed on a DC 14 Wisdom saving throw or be magically charmed. The charmed creature regards the dryad as a trusted friend to be heeded and protected. Although the target isn't under the dryad's control, it takes the dryad's requests or actions in the most favorable way it can.

Each time the dryad or its allies do anything harmful to the target, it can repeat the saving throw, ending the effect on itself on a success. Otherwise, the effect lasts 24 hours or until the dryad dies, is on a different plane of existence from the target, or ends the effect as a bonus action. If a target's saving throw is successful, the target is immune to the dryad's Fey Charm for the next 24 hours.

The dryad can have no more than one humanoid and up to three beasts charmed at a time.

\end{DndMonster}



\begin{DndMonster}[float=t]{Duergar}
\DndMonsterType{Medium humanoid (dwarf), lawful evil}
\DndMonsterBasics[
armor-class = {16 (\textit{scale mail})},
hit-points = {\DndDice{4d8 + 4}},
speed = {25 ft.}
]
\DndMonsterAbilityScores[
 str = 14,
 dex = 11,
 con = 14,
 int = 11,
 wis = 10,
 cha = 9
]
\DndMonsterDetails[
damage-resistances = {poison},
senses = {darkvision 120 ft., passive perception 10},
languages = {Dwarvish, Undercommon},
challenge = 1,
]
\DndMonsterAction{Duergar Resilience}
The duergar has advantage on saving throws against poison, spells, and illusions, as well as to resist being charmed or paralyzed.

\DndMonsterAction{Sunlight Sensitivity}
While in sunlight, the duergar has disadvantage on attack rolls, as well as on Wisdom (Perception) checks that rely on sight.

\DndMonsterSection{Actions}
\DndMonsterAction{Enlarge (Recharges after a Short or Long Rest)}
For 1 minute, the duergar magically increases in size, along with anything it is wearing or carrying. While enlarged, the duergar is Large, doubles its damage dice on Strength-based weapon attacks (included in the attacks), and makes Strength checks and Strength saving throws with advantage. If the duergar lacks the room to become Large, it attains the maximum size possible in the space available.

\DndMonsterAction{War Pick}
\textsl{Melee Weapon Attack:} 4 to hit, reach 5 ft., one target. \textsl{Hit:} 6 (1d8 + 2) piercing damage, or 11 (2d8 + 2) piercing damage while enlarged.

\DndMonsterAction{Javelin}
\textsl{Melee or Ranged Weapon Attack:} 4 to hit, reach 5 ft. or range 30/120 ft., one target. \textsl{Hit:} 5 (1d6 + 2) piercing damage, or 9 (2d6 + 2) piercing damage while enlarged.

\DndMonsterAction{Invisibility (Recharges after a Short or Long Rest)}
The duergar magically turns invisible until it attacks, casts a spell, or uses its Enlarge, or until its concentration is broken, up to 1 hour (as if concentration on a spell). Any equipment the duergar wears or carries is invisible with it.

\end{DndMonster}



\begin{DndMonster}[float=t]{Duodrone}
\DndMonsterType{Medium construct, lawful neutral}
\DndMonsterBasics[
armor-class = {15 (natural armor)},
hit-points = {\DndDice{2d8 + 2}},
speed = {30 ft.}
]
\DndMonsterAbilityScores[
 str = 11,
 dex = 13,
 con = 12,
 int = 6,
 wis = 10,
 cha = 7
]
\DndMonsterDetails[
senses = {truesight 120 ft., passive perception 10},
languages = {Modron},
challenge = 1/4,
]
\DndMonsterAction{Axiomatic Mind}
The duodrone can't be compelled to act in a manner contrary to its nature or its instructions.

\DndMonsterAction{Disintegration}
If the duodrone dies, its body disintegrates into dust, leaving behind its weapons and anything else it was carrying.

\DndMonsterSection{Actions}
\DndMonsterAction{Multiattack}
The duodrone makes two fist attacks or two javelin attacks.

\DndMonsterAction{Fist}
\textsl{Melee Weapon Attack:} 2 to hit, reach 5 ft., one target. \textsl{Hit:} 2 (1d4) bludgeoning damage.

\DndMonsterAction{Javelin}
\textsl{Melee or Ranged Weapon Attack:} 3 to hit, reach 5 ft. or range 30/120 ft., one target. \textsl{Hit:} 4 (1d6 + 1) piercing damage.

\end{DndMonster}



\begin{DndMonster}[float=t]{Dust Mephit}
\DndMonsterType{Small elemental, neutral evil}
\DndMonsterBasics[
armor-class = {12},
hit-points = {\DndDice{5d6}},
speed = {30 ft., fly 30 ft.}
]
\DndMonsterAbilityScores[
 str = 5,
 dex = 14,
 con = 10,
 int = 9,
 wis = 11,
 cha = 10
]
\DndMonsterDetails[
skills = {Perception +2, Stealth +4},
damage-vulnerabilities = {fire},
damage-immunities = {poison},
condition-immunities = {poisoned},
senses = {darkvision 60 ft., passive perception 12},
languages = {Auran, Terran},
challenge = 1/2,
]
\DndMonsterAction{Death Burst}
When the mephit dies, it explodes in a burst of dust. Each creature within 5 feet of it must then succeed on a DC 10 Constitution saving throw or be blinded for 1 minute. A blinded creature can repeat the saving throw on each of its turns, ending the effect on itself on a success.

\DndMonsterAction{Innate Spellcasting (1/Day)}
The mephit can innately cast \textit{sleep}, requiring no material components. Its innate spellcasting ability is Charisma.
\begin{DndMonsterSpells}
  \DndInnateSpellLevel{sleep}
\end{DndMonsterSpells}
\DndMonsterSection{Actions}
\DndMonsterAction{Claws}
\textsl{Melee Weapon Attack:} 4 to hit, reach 5 ft., one creature. \textsl{Hit:} 4 (1d4 + 2) slashing damage.

\DndMonsterAction{Blinding Breath (Recharge 6)}
The mephit exhales a 15-foot cone of blinding dust. Each creature in that area must succeed on a DC 10 Dexterity saving throw or be blinded for 1 minute. A creature can repeat the saving throw at the end of each of its turns, ending the effect on itself on a success.

\end{DndMonster}


\clearpage

\begin{DndMonster}[float=t]{Earth Elemental}
\DndMonsterType{Large elemental, neutral}
\DndMonsterBasics[
armor-class = {17 (natural armor)},
hit-points = {\DndDice{12d10 + 60}},
speed = {30 ft., burrow 30 ft.}
]
\DndMonsterAbilityScores[
 str = 20,
 dex = 8,
 con = 20,
 int = 5,
 wis = 10,
 cha = 5
]
\DndMonsterDetails[
damage-vulnerabilities = {thunder},
damage-resistances = {bludgeoning, piercing, slashing from nonmagical attacks},
damage-immunities = {poison},
condition-immunities = {exhaustion, paralyzed, petrified, poisoned, unconscious},
senses = {darkvision 60 ft., tremorsense 60 ft., passive perception 10},
languages = {Terran},
challenge = 5,
]
\DndMonsterAction{Earth Glide}
The elemental can burrow through nonmagical, unworked earth and stone. While doing so, the elemental doesn't disturb the material it moves through.

\DndMonsterAction{Siege Monster}
The elemental deals double damage to objects and structures.

\DndMonsterSection{Actions}
\DndMonsterAction{Multiattack}
The elemental makes two slam attacks.

\DndMonsterAction{Slam}
\textsl{Melee Weapon Attack:} 8 to hit, reach 10 ft., one target. \textsl{Hit:} 14 (2d8 + 5) bludgeoning damage.

\end{DndMonster}



\begin{DndMonster}[float=t]{Efreeti}
\DndMonsterType{Large elemental, lawful evil}
\DndMonsterBasics[
armor-class = {17 (natural armor)},
hit-points = {\DndDice{16d10 + 112}},
speed = {40 ft., fly 60 ft.}
]
\DndMonsterAbilityScores[
 str = 22,
 dex = 12,
 con = 24,
 int = 16,
 wis = 15,
 cha = 16
]
\DndMonsterDetails[
saving-throws = {Int +7, Wis +6, Cha +7},
damage-immunities = {fire},
senses = {darkvision 120 ft., passive perception 12},
languages = {Ignan},
challenge = 11,
]
\DndMonsterAction{Elemental Demise}
If the efreeti dies, its body disintegrates in a flash of fire and puff of smoke, leaving behind only equipment the efreeti was wearing or carrying.

\DndMonsterAction{Innate Spellcasting}
The efreeti's innate spellcasting ability is Charisma (spell save DC 15, 7 to hit with spell attacks). It can innately cast the following spells, requiring no material components:
\begin{DndMonsterSpells}
  \DndInnateSpellLevel{detect magic}
  \DndInnateSpellLevel[3]{enlarge/reduce}
  \DndInnateSpellLevel[1]{conjure elemental}
\end{DndMonsterSpells}
\DndMonsterSection{Actions}
\DndMonsterAction{Multiattack}
The efreeti makes two scimitar attacks or uses its Hurl Flame twice.

\DndMonsterAction{Scimitar}
\textsl{Melee Weapon Attack:} 10 to hit, reach 5 ft., one target. \textsl{Hit:} 13 (2d6 + 6) slashing damage plus 7 (2d6) fire damage.

\DndMonsterAction{Hurl Flame}
\textsl{Ranged Spell Attack:} 7 to hit, range 120 ft., one target. \textsl{Hit:} 17 (5d6) fire damage.

\end{DndMonster}



\begin{DndMonster}[float=t]{Elk}
\DndMonsterType{Large beast, unaligned}
\DndMonsterBasics[
armor-class = {10},
hit-points = {\DndDice{2d10 + 2}},
speed = {50 ft.}
]
\DndMonsterAbilityScores[
 str = 16,
 dex = 10,
 con = 12,
 int = 2,
 wis = 10,
 cha = 6
]
\DndMonsterDetails[
challenge = 1/4,
]
\DndMonsterAction{Charge}
If the elk moves at least 20 feet straight toward a target and then hits it with a ram attack on the same turn, the target takes an extra 7 (2d6) damage. If the target is a creature, it must succeed on a DC 13 Strength saving throw or be knocked prone.

\DndMonsterSection{Actions}
\DndMonsterAction{Ram}
\textsl{Melee Weapon Attack:} 5 to hit, reach 5 ft., one target. \textsl{Hit:} 6 (1d6 + 3) bludgeoning damage.

\DndMonsterAction{Hooves}
\textsl{Melee Weapon Attack:} 5 to hit, reach 5 ft., one prone creature. \textsl{Hit:} 8 (2d4 + 3) bludgeoning damage.

\end{DndMonster}



\begin{DndMonster}[float*=b,width=\textwidth + 8pt]{Empyrean}
\begin{multicols}{2}
\DndMonsterType{Huge celestial (titan), chaotic good (75\%) neutral evil (25\%)}
\DndMonsterBasics[
armor-class = {22 (natural armor)},
hit-points = {\DndDice{19d12 + 190}},
speed = {50 ft., fly 50 ft., swim 50 ft.}
]
\DndMonsterAbilityScores[
 str = 30,
 dex = 21,
 con = 30,
 int = 21,
 wis = 22,
 cha = 27
]
\DndMonsterDetails[
saving-throws = {Str +17, Int +12, Wis +13, Cha +15},
skills = {Insight +13, Persuasion +15},
damage-immunities = {bludgeoning, piercing, slashing from nonmagical attacks},
senses = {truesight 120 ft., passive perception 16},
languages = {all},
challenge = 23,
]
\DndMonsterAction{Legendary Resistance (3/Day)}
If the empyrean fails a saving throw, it can choose to succeed instead.

\DndMonsterAction{Magic Resistance}
The empyrean has advantage on saving throws against spells and other magical effects.

\DndMonsterAction{Magic Weapons}
The empyrean's weapon attacks are magical.

\DndMonsterAction{Innate Spellcasting}
The empyrean's innate spellcasting ability is Charisma (spell save DC 23, 15 to hit with spell attacks). It can innately cast the following spells, requiring no material components:
\begin{DndMonsterSpells}
  \DndInnateSpellLevel{greater restoration}
  \DndInnateSpellLevel[1]{commune}
\end{DndMonsterSpells}
\DndMonsterSection{Actions}
\DndMonsterAction{Maul}
\textsl{Melee Weapon Attack:} 17 to hit, reach 10 ft., one target. \textsl{Hit:} 31 (6d6 + 10) bludgeoning damage. If the target is a creature, it must succeed on a DC 15 Constitution saving throw or be stunned until the end of the empyrean's next turn.

\DndMonsterAction{Bolt}
\textsl{Ranged Spell Attack:} 15 to hit, range 600 ft., one target. \textsl{Hit:} 24 (7d6) damage of one of the following types (empyrean's choice): acid, cold, fire, force, lightning, radiant, or thunder.

\DndMonsterSection{Legendary Actions}
Empyrean can take 3 legendary actions, choosing from the options below. Only one legendary action can be used at a time and only at the end of another creature's turn. Empyrean regains spent legendary actions at the start of its turn.

\begin{DndMonsterLegendaryActions}
\DndMonsterLegendaryAction{Attack}{The empyrean makes one attack.
}
\DndMonsterLegendaryAction{Bolster}{The empyrean bolsters all nonhostile creatures within 120 feet of it until the end of its next turn. Bolstered creatures can't be charmed or frightened, and they gain advantage on ability checks and saving throws until the end of the empyrean's next turn.
}
\DndMonsterLegendaryAction{Trembling Strike (Costs 2 Actions)}{The empyrean strikes the ground with its maul, triggering an earth tremor. All other creatures on the ground within 60 feet of the empyrean must succeed on a DC 25 Strength saving throw or be knocked prone.
}
\end{DndMonsterLegendaryActions}
\end{multicols}
\end{DndMonster}



\begin{DndMonster}[float=t]{Fire Giant}
\DndMonsterType{Huge giant, lawful evil}
\DndMonsterBasics[
armor-class = {18 (\textit{plate armor})},
hit-points = {\DndDice{13d12 + 78}},
speed = {30 ft.}
]
\DndMonsterAbilityScores[
 str = 25,
 dex = 9,
 con = 23,
 int = 10,
 wis = 14,
 cha = 13
]
\DndMonsterDetails[
saving-throws = {Dex +3, Con +10, Cha +5},
skills = {Athletics +11, Perception +6},
damage-immunities = {fire},
languages = {Giant},
challenge = 9,
]
\DndMonsterSection{Actions}
\DndMonsterAction{Multiattack}
The giant makes two greatsword attacks.

\DndMonsterAction{Greatsword}
\textsl{Melee Weapon Attack:} 11 to hit, reach 10 ft., one target. \textsl{Hit:} 28 (6d6 + 7) slashing damage.

\DndMonsterAction{Rock}
\textsl{Ranged Weapon Attack:} 11 to hit, range 60/240 ft., one target. \textsl{Hit:} 29 (4d10 + 7) bludgeoning damage.

\end{DndMonster}



\begin{DndMonster}[float=t]{Flameskull}
\DndMonsterType{Tiny undead, neutral evil}
\DndMonsterBasics[
armor-class = {13},
hit-points = {\DndDice{9d4 + 18}},
speed = {0 ft., fly 40 ft. \textit{(hover)}}
]
\DndMonsterAbilityScores[
 str = 1,
 dex = 17,
 con = 14,
 int = 16,
 wis = 10,
 cha = 11
]
\DndMonsterDetails[
skills = {Arcana +5, Perception +2},
damage-resistances = {lightning, necrotic, piercing},
damage-immunities = {cold, fire, poison},
condition-immunities = {charmed, frightened, paralyzed, poisoned, prone},
senses = {darkvision 60 ft., passive perception 12},
languages = {Common},
challenge = 4,
]
\DndMonsterAction{Illumination}
The flameskull sheds either dim light in a 15-foot radius, or bright light in a 15-foot radius and dim light for an additional 15 feet. It can switch between the options as an action.

\DndMonsterAction{Magic Resistance}
The flameskull has advantage on saving throws against spells and other magical effects.

\DndMonsterAction{Rejuvenation}
If the flameskull is destroyed, it regains all its hit points in 1 hour unless holy water is sprinkled on its remains or a \textit{dispel magic} or \textit{remove curse} spell is cast on them.

\DndMonsterAction{Spellcasting}
The flameskull is a 5th-level spellcaster. Its spellcasting ability is Intelligence (spell save DC 13, 5 to hit with spell attacks). It requires no somatic or material components to cast its spells. The flameskull has the following wizard spells prepared:
\begin{DndMonsterSpells}
  \DndMonsterSpellLevel{mage hand}
   \DndMonsterSpellLevel[1][3]{magic missile}
   \DndMonsterSpellLevel[2][2]{blur}
   \DndMonsterSpellLevel[3][1]{fireball}
\end{DndMonsterSpells}
\DndMonsterSection{Actions}
\DndMonsterAction{Multiattack}
The flameskull uses Fire Ray twice.

\DndMonsterAction{Fire Ray}
\textsl{Ranged Spell Attack:} 5 to hit, range 30 ft., one target. \textsl{Hit:} 10 (3d6) fire damage.

\end{DndMonster}



\begin{DndMonster}[float=t]{Fomorian}
\DndMonsterType{Huge giant, chaotic evil}
\DndMonsterBasics[
armor-class = {14 (natural armor)},
hit-points = {\DndDice{13d12 + 65}},
speed = {30 ft.}
]
\DndMonsterAbilityScores[
 str = 23,
 dex = 10,
 con = 20,
 int = 9,
 wis = 14,
 cha = 6
]
\DndMonsterDetails[
skills = {Perception +8, Stealth +3},
senses = {darkvision 120 ft., passive perception 18},
languages = {Giant, Undercommon},
challenge = 8,
]
\DndMonsterSection{Actions}
\DndMonsterAction{Multiattack}
The fomorian attacks twice with its greatclub or makes one greatclub attack and uses Evil Eye once.

\DndMonsterAction{Greatclub}
\textsl{Melee Weapon Attack:} 9 to hit, reach 15 ft., one target. \textsl{Hit:} 19 (3d8 + 6) bludgeoning damage.

\DndMonsterAction{Evil Eye}
The fomorian magically forces a creature it can see within 60 feet of it to make a DC 14 Charisma saving throw. The creature takes 27 (6d8) psychic damage on a failed save, or half as much damage on a successful one.

\DndMonsterAction{Curse of the Evil Eye (Recharges after a Short or Long Rest)}
With a stare, the fomorian uses Evil Eye, but on a failed save, the creature is also cursed with magical deformities. While deformed, the creature has its speed halved and has disadvantage on ability checks, saving throws, and attacks based on Strength or Dexterity.

The transformed creature can repeat the saving throw whenever it finishes a long rest, ending the effect on a success.

\end{DndMonster}



\begin{DndMonster}[float=t]{Frog}
\DndMonsterType{Tiny beast, unaligned}
\DndMonsterBasics[
armor-class = {11},
hit-points = {\DndDice{1d4 - 1}},
speed = {20 ft., swim 20 ft.}
]
\DndMonsterAbilityScores[
 str = 1,
 dex = 13,
 con = 8,
 int = 1,
 wis = 8,
 cha = 3
]
\DndMonsterDetails[
skills = {Perception +1, Stealth +3},
senses = {darkvision 30 ft., passive perception 11},
challenge = 0,
]
\DndMonsterAction{Amphibious}
The frog can breathe air and water.

\DndMonsterAction{Standing Leap}
The frog's long jump is up to 10 feet and its high jump is up to 5 feet, with or without a running start.

\end{DndMonster}



\begin{DndMonster}[float*=b,width=\textwidth + 8pt]{Froghemoth Elder}
\begin{multicols}{2}
\DndMonsterType{Huge monstrosity, unaligned}
\DndMonsterBasics[
armor-class = {14 (natural armor)},
hit-points = {\DndDice{18d12 + 90}},
speed = {30 ft., swim 30 ft.}
]
\DndMonsterAbilityScores[
 str = 25,
 dex = 13,
 con = 20,
 int = 2,
 wis = 14,
 cha = 8
]
\DndMonsterDetails[
saving-throws = {Str +12, Con +10, Wis +7},
skills = {Perception +12, Stealth +6},
damage-resistances = {fire, lightning},
senses = {darkvision 60 ft., passive perception 22},
challenge = 15,
]
\DndMonsterAction{Amphibious}
The froghemoth can breathe air and water.

\DndMonsterAction{Legendary Resistance (3/Day)}
If the froghemoth fails a saving throw, it can choose to succeed instead.

\DndMonsterAction{Shock Susceptibility}
If the froghemoth takes lightning damage, it suffers two effects until the end of its next turn: its speed is halved, and it has disadvantage on Dexterity saving throws.

\DndMonsterSection{Actions}
\DndMonsterAction{Multiattack}
The froghemoth makes one Bite attack and two Tentacle attacks, and it can use Tongue.

\DndMonsterAction{Bite}
\textsl{Melee Weapon Attack:} 12 to hit, reach 5 ft., one target. \textsl{Hit:} 23 (3d10 + 7) piercing damage, and the target is swallowed if it is a Medium or smaller creature. A swallowed creature has the blinded and restrained conditions, has total cover against attacks and other effects outside the froghemoth, and takes 10 (3d6) acid damage at the start of each of the froghemoth's turns.

The froghemoth's gullet can hold up to three creatures at a time. If the froghemoth takes 25 damage or more on a single turn from a creature inside it, the froghemoth must succeed on a DC 20 Constitution saving throw at the end of that turn or regurgitate all swallowed creatures, each of which lands in a space within 10 feet of the froghemoth and has the prone condition. If the froghemoth dies, any swallowed creatures are no longer restrained by it and can escape from the corpse using 10 feet of movement, exiting with the prone condition.

\DndMonsterAction{Tentacle}
\textsl{Melee Weapon Attack:} 12 to hit, reach 20 ft., one target. \textsl{Hit:} 20 (3d8 + 7) bludgeoning damage, and the target has the grappled condition (escape DC 20). Until this grapple ends, the froghemoth can't use this tentacle on another target. The froghemoth has four tentacles.

\DndMonsterAction{Tongue}
The froghemoth targets one Large or smaller creature that it can see within 25 feet of it. The target must make a DC 20 Strength saving throw. On a failed save, the target is pulled into an unoccupied space within 5 feet of the froghemoth.

\DndMonsterSection{Reactions}
\DndMonsterAction{Alien Gaze}
When a creature the froghemoth can see damages the froghemoth, the froghemoth swivels its eyestalk toward the creature and pierces the creature's mind with its otherworldly gaze. That creature must make a DC 18 Intelligence saving throw, taking 7 (2d6) psychic damage on a failed save or half as much damage on a successful one.

\DndMonsterAction{Leap}
Immediately after a creature the froghemoth can see ends its turn, the froghemoth jumps up to half its speed. Each creature within 5 feet of the froghemoth when it lands must succeed on a DC 20 Strength saving throw or have the prone condition.

\DndMonsterSection{Lair Actions}
On initiative count 20 (losing initiative ties), the froghemoth can take one of the following lair actions; the froghemoth can't take the same lair action two rounds in a row:


\begin{description}
\item[Croak.]
The froghemoth lets out an ear-splitting croak. Each creature within 30 feet of it must make a DC 18 Wisdom saving throw. On a failed save, the creature has the deafened and frightened conditions until the end of its next turn.
\item[Hypnotic Visions.]
One creature of the froghemoth's choice within 60 feet of it must make a DC 15 Wisdom saving throw. On a failed save, the target is mesmerized by otherworldly visions and has the charmed condition until the end of the target's next turn. While charmed in this way, the target has the incapacitated condition, and its speed is 0.
\end{description}

\DndMonsterSection{Regional Effects}
The region within 1 mile of a froghemoth elder's lair is altered by the froghemoth's unearthly presence, creating one or more of the following effects:


\begin{description}
\item[Alien Flora.]
Strange, invasive plants with bright colors and odd patterns sprout up around the lair.
\item[Mutations.]
Wildlife encountered near the lair often sport mutations, such as extra limbs or eyes, glowing secretions, or vestigial tentacles.
\end{description}

If the froghemoth elder dies, these effects fade over the course of 1d10 days.

\end{multicols}
\end{DndMonster}



\begin{DndMonster}[float=t]{Galeb Duhr}
\DndMonsterType{Medium elemental, neutral}
\DndMonsterBasics[
armor-class = {16 (natural armor)},
hit-points = {\DndDice{9d8 + 45}},
speed = {15 ft. (30 ft. when rolling, 60 ft. rolling downhill)}
]
\DndMonsterAbilityScores[
 str = 20,
 dex = 14,
 con = 20,
 int = 11,
 wis = 12,
 cha = 11
]
\DndMonsterDetails[
damage-resistances = {bludgeoning, piercing, slashing from nonmagical attacks},
damage-immunities = {poison},
condition-immunities = {exhaustion, paralyzed, petrified, poisoned},
senses = {darkvision 60 ft., tremorsense 60 ft., passive perception 11},
languages = {Terran},
challenge = 6,
]
\DndMonsterAction{False Appearance}
While the galeb duhr remains motionless, it is indistinguishable from a normal boulder.

\DndMonsterAction{Rolling Charge}
If the galeb duhr rolls at least 20 feet straight toward a target and then hits it with a slam attack on the same turn, the target takes an extra 7 (2d6) bludgeoning damage. If the target is a creature, it must succeed on a DC 16 Strength saving throw or be knocked prone.

\DndMonsterSection{Actions}
\DndMonsterAction{Slam}
\textsl{Melee Weapon Attack:} 8 to hit, reach 5 ft., one target. \textsl{Hit:} 12 (2d6 + 5) bludgeoning damage.

\DndMonsterAction{Animate Boulders (1/Day)}
The galeb duhr magically animates up to two boulders it can see within 60 feet of it. A boulder has statistics like those of a galeb duhr, except it has Intelligence 1 and Charisma 1, it can't be charmed or frightened, and it lacks this action option. A boulder remains animated as long as the galeb duhr maintains concentration, up to 1 minute (as if concentration on a spell).

\end{DndMonster}



\begin{DndMonster}[float=t]{Gargoyle}
\DndMonsterType{Medium elemental, chaotic evil}
\DndMonsterBasics[
armor-class = {15 (natural armor)},
hit-points = {\DndDice{7d8 + 21}},
speed = {30 ft., fly 60 ft.}
]
\DndMonsterAbilityScores[
 str = 15,
 dex = 11,
 con = 16,
 int = 6,
 wis = 11,
 cha = 7
]
\DndMonsterDetails[
damage-resistances = {bludgeoning, piercing, slashing from nonmagical attacks that aren't adamantine},
damage-immunities = {poison},
condition-immunities = {exhaustion, petrified, poisoned},
senses = {darkvision 60 ft., passive perception 10},
languages = {Terran},
challenge = 2,
]
\DndMonsterAction{False Appearance}
While the gargoyle remains motionless, it is indistinguishable from an inanimate statue.

\DndMonsterSection{Actions}
\DndMonsterAction{Multiattack}
The gargoyle makes two attacks: one with its bite and one with its claws.

\DndMonsterAction{Bite}
\textsl{Melee Weapon Attack:} 4 to hit, reach 5 ft., one target. \textsl{Hit:} 5 (1d6 + 2) piercing damage.

\DndMonsterAction{Claws}
\textsl{Melee Weapon Attack:} 4 to hit, reach 5 ft., one target. \textsl{Hit:} 5 (1d6 + 2) slashing damage.

\end{DndMonster}



\begin{DndMonster}[float=t]{Gas Spore}
\DndMonsterType{Large plant, unaligned}
\DndMonsterBasics[
armor-class = {5},
hit-points = {\DndDice{1d10 - 4}},
speed = {0 ft., fly 10 ft. \textit{(hover)}}
]
\DndMonsterAbilityScores[
 str = 5,
 dex = 1,
 con = 3,
 int = 1,
 wis = 1,
 cha = 1
]
\DndMonsterDetails[
damage-immunities = {poison},
condition-immunities = {blinded, deafened, frightened, paralyzed, poisoned, prone},
senses = {blindsight 30 ft. (blind beyond this radius), passive perception 5},
challenge = 1/2,
]
\DndMonsterAction{Death Burst}
The gas spore explodes when it drops to 0 hit points. Each creature within 20 feet of it must succeed on a DC 15 Constitution saving throw or take 10 (3d6) poison damage and become infected with a disease on a failed save. Creatures immune to the poisoned condition are immune to this disease.

Spores invade an infected creature's system, killing the creature in a number of hours equal to 1d12 + the creature's Constitution score, unless the disease is removed. In half that time, the creature becomes poisoned for the rest of the duration. After the creature dies, it sprouts 2d4 Tiny gas spores that grow to full size in 7 days.

\DndMonsterAction{Eerie Resemblance}
The gas spore resembles a beholder. A creature that can see the gas spore can discern its true nature with a successful DC 15 Intelligence (Nature) check.

\DndMonsterSection{Actions}
\DndMonsterAction{Touch}
\textsl{Melee Weapon Attack:} 0 to hit, reach 5 ft., one creature. \textsl{Hit:} 1 poison damage, and the creature must succeed on a DC 10 Constitution saving throw or become infected with the disease described in the Death Burst trait.

\end{DndMonster}



\begin{DndMonster}[float=t]{Gelatinous Cube}
\DndMonsterType{Large ooze, unaligned}
\DndMonsterBasics[
armor-class = {6},
hit-points = {\DndDice{8d10 + 40}},
speed = {15 ft.}
]
\DndMonsterAbilityScores[
 str = 14,
 dex = 3,
 con = 20,
 int = 1,
 wis = 6,
 cha = 1
]
\DndMonsterDetails[
condition-immunities = {blinded, charmed, deafened, exhaustion, frightened, prone},
senses = {blindsight 60 ft. (blind beyond this radius), passive perception 8},
challenge = 2,
]
\DndMonsterAction{Ooze Cube}
The cube takes up its entire space. Other creatures can enter the space, but a creature that does so is subjected to the cube's Engulf and has disadvantage on the saving throw.

Creatures inside the cube can be seen but have total cover.

A creature within 5 feet of the cube can take an action to pull a creature or object out of the cube. Doing so requires a successful DC 12 Strength check, and the creature making the attempt takes 10 (3d6) acid damage.

The cube can hold only one Large creature or up to four Medium or smaller creatures inside it at a time.

\DndMonsterAction{Transparent}
Even when the cube is in plain sight, it takes a successful DC 15 Wisdom (Perception) check to spot a cube that has neither moved nor attacked. A creature that tries to enter the cube's space while unaware of the cube is surprised by the cube.

\DndMonsterSection{Actions}
\DndMonsterAction{Pseudopod}
\textsl{Melee Weapon Attack:} 4 to hit, reach 5 ft., one creature. \textsl{Hit:} 10 (3d6) acid damage.

\DndMonsterAction{Engulf}
The cube moves up to its speed. While doing so, it can enter Large or smaller creatures' spaces. Whenever the cube enters a creature's space, the creature must make a DC 12 Dexterity saving throw.

On a successful save, the creature can choose to be pushed 5 feet back or to the side of the cube. A creature that chooses not to be pushed suffers the consequences of a failed saving throw.

On a failed save, the cube enters the creature's space, and the creature takes 10 (3d6) acid damage and is engulfed. The engulfed creature can't breathe, is restrained, and takes 21 (6d6) acid damage at the start of each of the cube's turns. When the cube moves, the engulfed creature moves with it.

An engulfed creature can try to escape by taking an action to make a DC 12 Strength check. On a success, the creature escapes and enters a space of its choice within 5 feet of the cube.

\end{DndMonster}



\begin{DndMonster}[float=t]{Ghast}
\DndMonsterType{Medium undead, chaotic evil}
\DndMonsterBasics[
armor-class = {13},
hit-points = {\DndDice{8d8}},
speed = {30 ft.}
]
\DndMonsterAbilityScores[
 str = 16,
 dex = 17,
 con = 10,
 int = 11,
 wis = 10,
 cha = 8
]
\DndMonsterDetails[
damage-resistances = {necrotic},
damage-immunities = {poison},
condition-immunities = {charmed, exhaustion, poisoned},
senses = {darkvision 60 ft., passive perception 10},
languages = {Common},
challenge = 2,
]
\DndMonsterAction{Stench}
Any creature that starts its turn within 5 feet of the ghast must succeed on a DC 10 Constitution saving throw or be poisoned until the start of its next turn. On a successful saving throw, the creature is immune to the ghast's Stench for 24 hours.

\DndMonsterAction{Turn Defiance}
The ghast and any ghouls within 30 feet of it have advantage on saving throws against effects that turn undead.

\DndMonsterSection{Actions}
\DndMonsterAction{Bite}
\textsl{Melee Weapon Attack:} 3 to hit, reach 5 ft., one creature. \textsl{Hit:} 12 (2d8 + 3) piercing damage.

\DndMonsterAction{Claws}
\textsl{Melee Weapon Attack:} 5 to hit, reach 5 ft., one target. \textsl{Hit:} 10 (2d6 + 3) slashing damage. If the target is a creature other than an undead, it must succeed on a DC 10 Constitution saving throw or be paralyzed for 1 minute. The target can repeat the saving throw at the end of each of its turns, ending the effect on itself on a success.

\end{DndMonster}



\begin{DndMonster}[float=t]{Ghost}
\DndMonsterType{Medium undead, any alignment}
\DndMonsterBasics[
armor-class = {11},
hit-points = {\DndDice{10d8}},
speed = {0 ft., fly 40 ft. \textit{(hover)}}
]
\DndMonsterAbilityScores[
 str = 7,
 dex = 13,
 con = 10,
 int = 10,
 wis = 12,
 cha = 17
]
\DndMonsterDetails[
damage-resistances = {acid; fire; lightning; thunder; bludgeoning, piercing, slashing from nonmagical attacks},
damage-immunities = {cold, necrotic, poison},
condition-immunities = {charmed, exhaustion, frightened, grappled, paralyzed, petrified, poisoned, prone, restrained},
senses = {darkvision 60 ft., passive perception 11},
languages = {any languages it knew in life},
challenge = 4,
]
\DndMonsterAction{Ethereal Sight}
The ghost can see 60 feet into the Ethereal Plane when it is on the Material Plane, and vice versa.

\DndMonsterAction{Incorporeal Movement}
The ghost can move through other creatures and objects as if they were difficult terrain. It takes 5 (1d10) force damage if it ends its turn inside an object.

\DndMonsterSection{Actions}
\DndMonsterAction{Withering Touch}
\textsl{Melee Weapon Attack:} 5 to hit, reach 5 ft., one target. \textsl{Hit:} 17 (4d6 + 3) necrotic damage.

\DndMonsterAction{Etherealness}
The ghost enters the Ethereal Plane from the Material Plane, or vice versa. It is visible on the Material Plane while it is in the Border Ethereal, and vice versa, yet it can't affect or be affected by anything on the other plane.

\DndMonsterAction{Horrifying Visage}
Each non-undead creature within 60 feet of the ghost that can see it must succeed on a DC 13 Wisdom saving throw or be frightened for 1 minute. If the save fails by 5 or more, the target also ages 1d4 × 10 years. A frightened target can repeat the saving throw at the end of each of its turns, ending the frightened condition on itself on a success. If a target's saving throw is successful or the effect ends for it, the target is immune to this ghost's Horrifying Visage for the next 24 hours. The aging effect can be reversed with a  \textit{greater restoration} spell, but only within 24 hours of it occurring.

\DndMonsterAction{Possession (Recharge 6)}
One humanoid that the ghost can see within 5 feet of it must succeed on a DC 13 Charisma saving throw or be possessed by the ghost; the ghost then disappears, and the target is incapacitated and loses control of its body. The ghost now controls the body but doesn't deprive the target of awareness. The ghost can't be targeted by any attack, spell, or other effect, except ones that turn undead, and it retains its alignment, Intelligence, Wisdom, Charisma, and immunity to being charmed and frightened. It otherwise uses the possessed target's statistics, but doesn't gain access to the target's knowledge, class features, or proficiencies.

The possession lasts until the body drops to 0 hit points, the ghost ends it as a bonus action, or the ghost is turned or forced out by an effect like the \textit{dispel evil and good} spell. When the possession ends, the ghost reappears in an unoccupied space within 5 feet of the body. The target is immune to this ghost's Possession for 24 hours after succeeding on the saving throw or after the possession ends.

\end{DndMonster}



\begin{DndMonster}[float=t]{Ghoul}
\DndMonsterType{Medium undead, chaotic evil}
\DndMonsterBasics[
armor-class = {12},
hit-points = {\DndDice{5d8}},
speed = {30 ft.}
]
\DndMonsterAbilityScores[
 str = 13,
 dex = 15,
 con = 10,
 int = 7,
 wis = 10,
 cha = 6
]
\DndMonsterDetails[
damage-immunities = {poison},
condition-immunities = {charmed, exhaustion, poisoned},
senses = {darkvision 60 ft., passive perception 10},
languages = {Common},
challenge = 1,
]
\DndMonsterSection{Actions}
\DndMonsterAction{Bite}
\textsl{Melee Weapon Attack:} 2 to hit, reach 5 ft., one creature. \textsl{Hit:} 9 (2d6 + 2) piercing damage.

\DndMonsterAction{Claws}
\textsl{Melee Weapon Attack:} 4 to hit, reach 5 ft., one target. \textsl{Hit:} 7 (2d4 + 2) slashing damage. If the target is a creature other than an elf or undead, it must succeed on a DC 10 Constitution saving throw or be paralyzed for 1 minute. The target can repeat the saving throw at the end of each of its turns, ending the effect on itself on a success.

\end{DndMonster}



\begin{DndMonster}[float=t]{Giant Badger}
\DndMonsterType{Medium beast, unaligned}
\DndMonsterBasics[
armor-class = {10},
hit-points = {\DndDice{2d8 + 4}},
speed = {30 ft., burrow 10 ft.}
]
\DndMonsterAbilityScores[
 str = 13,
 dex = 10,
 con = 15,
 int = 2,
 wis = 12,
 cha = 5
]
\DndMonsterDetails[
senses = {darkvision 30 ft., passive perception 11},
challenge = 1/4,
]
\DndMonsterAction{Keen Smell}
The badger has advantage on Wisdom (Perception) checks that rely on smell.

\DndMonsterSection{Actions}
\DndMonsterAction{Multiattack}
The badger makes two attacks: one with its bite and one with its claws.

\DndMonsterAction{Bite}
\textsl{Melee Weapon Attack:} 3 to hit, reach 5 ft., one target. \textsl{Hit:} 4 (1d6 + 1) piercing damage.

\DndMonsterAction{Claws}
\textsl{Melee Weapon Attack:} 3 to hit, reach 5 ft., one target. \textsl{Hit:} 6 (2d4 + 1) slashing damage.

\end{DndMonster}



\begin{DndMonster}[float=t]{Giant Bat}
\DndMonsterType{Large beast, unaligned}
\DndMonsterBasics[
armor-class = {13},
hit-points = {\DndDice{4d10}},
speed = {10 ft., fly 60 ft.}
]
\DndMonsterAbilityScores[
 str = 15,
 dex = 16,
 con = 11,
 int = 2,
 wis = 12,
 cha = 6
]
\DndMonsterDetails[
senses = {blindsight 60 ft., passive perception 11},
challenge = 1/4,
]
\DndMonsterAction{Echolocation}
The bat can't use its blindsight while deafened.

\DndMonsterAction{Keen Hearing}
The bat has advantage on Wisdom (Perception) checks that rely on hearing.

\DndMonsterSection{Actions}
\DndMonsterAction{Bite}
\textsl{Melee Weapon Attack:} 4 to hit, reach 5 ft., one creature. \textsl{Hit:} 5 (1d6 + 2) piercing damage.

\end{DndMonster}



\begin{DndMonster}[float=t]{Giant Centipede}
\DndMonsterType{Small beast, unaligned}
\DndMonsterBasics[
armor-class = {13 (natural armor)},
hit-points = {\DndDice{1d6 + 1}},
speed = {30 ft., climb 30 ft.}
]
\DndMonsterAbilityScores[
 str = 5,
 dex = 14,
 con = 12,
 int = 1,
 wis = 7,
 cha = 3
]
\DndMonsterDetails[
senses = {blindsight 30 ft., passive perception 8},
challenge = 1/4,
]
\DndMonsterSection{Actions}
\DndMonsterAction{Bite}
\textsl{Melee Weapon Attack:} 4 to hit, reach 5 ft., one creature. \textsl{Hit:} 4 (1d4 + 2) piercing damage, and the target must succeed on a DC 11 Constitution saving throw or take 10 (3d6) poison damage. If the poison damage reduces the target to 0 hit points, the target is stable but poisoned for 1 hour, even after regaining hit points, and is paralyzed while poisoned in this way.

\end{DndMonster}



\begin{DndMonster}[float=t]{Giant Constrictor Snake}
\DndMonsterType{Huge beast, unaligned}
\DndMonsterBasics[
armor-class = {12},
hit-points = {\DndDice{8d12 + 8}},
speed = {30 ft., swim 30 ft.}
]
\DndMonsterAbilityScores[
 str = 19,
 dex = 14,
 con = 12,
 int = 1,
 wis = 10,
 cha = 3
]
\DndMonsterDetails[
skills = {Perception +2},
senses = {blindsight 10 ft., passive perception 12},
challenge = 2,
]
\DndMonsterSection{Actions}
\DndMonsterAction{Bite}
\textsl{Melee Weapon Attack:} 6 to hit, reach 10 ft., one creature. \textsl{Hit:} 11 (2d6 + 4) piercing damage.

\DndMonsterAction{Constrict}
\textsl{Melee Weapon Attack:} 6 to hit, reach 5 ft., one creature. \textsl{Hit:} 13 (2d8 + 4) bludgeoning damage, and the target is grappled (escape DC 16). Until this grapple ends, the creature is restrained, and the snake can't constrict another target.

\end{DndMonster}


\clearpage

\begin{DndMonster}[float=t]{Giant Eagle}
\DndMonsterType{Large beast, neutral good}
\DndMonsterBasics[
armor-class = {13},
hit-points = {\DndDice{4d10 + 4}},
speed = {10 ft., fly 80 ft.}
]
\DndMonsterAbilityScores[
 str = 16,
 dex = 17,
 con = 13,
 int = 8,
 wis = 14,
 cha = 10
]
\DndMonsterDetails[
skills = {Perception +4},
languages = {Giant Eagle, understands Common and Auran but can't speak them},
challenge = 1,
]
\DndMonsterAction{Keen Sight}
The eagle has advantage on Wisdom (Perception) checks that rely on sight.

\DndMonsterSection{Actions}
\DndMonsterAction{Multiattack}
The eagle makes two attacks: one with its beak and one with its talons.

\DndMonsterAction{Beak}
\textsl{Melee Weapon Attack:} 5 to hit, reach 5 ft., one target. \textsl{Hit:} 6 (1d6 + 3) piercing damage.

\DndMonsterAction{Talons}
\textsl{Melee Weapon Attack:} 5 to hit, reach 5 ft., one target. \textsl{Hit:} 10 (2d6 + 3) slashing damage.

\end{DndMonster}



\begin{DndMonster}[float=t]{Giant Fire Beetle}
\DndMonsterType{Small beast, unaligned}
\DndMonsterBasics[
armor-class = {13 (natural armor)},
hit-points = {\DndDice{1d6 + 1}},
speed = {30 ft.}
]
\DndMonsterAbilityScores[
 str = 8,
 dex = 10,
 con = 12,
 int = 1,
 wis = 7,
 cha = 3
]
\DndMonsterDetails[
senses = {blindsight 30 ft., passive perception 8},
challenge = 0,
]
\DndMonsterAction{Illumination}
The beetle sheds bright light in a 10-foot radius and dim light for an additional 10 ft..

\DndMonsterSection{Actions}
\DndMonsterAction{Bite}
\textsl{Melee Weapon Attack:} 1 to hit, reach 5 ft., one target. \textsl{Hit:} 2 (1d6 - 1) slashing damage.

\end{DndMonster}



\begin{DndMonster}[float=t]{Giant Frog}
\DndMonsterType{Medium beast, unaligned}
\DndMonsterBasics[
armor-class = {11},
hit-points = {\DndDice{4d8}},
speed = {30 ft., swim 30 ft.}
]
\DndMonsterAbilityScores[
 str = 12,
 dex = 13,
 con = 11,
 int = 2,
 wis = 10,
 cha = 3
]
\DndMonsterDetails[
skills = {Perception +2, Stealth +3},
senses = {darkvision 30 ft., passive perception 12},
challenge = 1/4,
]
\DndMonsterAction{Amphibious}
The frog can breathe air and water.

\DndMonsterAction{Standing Leap}
The frog's long jump is up to 20 feet and its high jump is up to 10 feet, with or without a running start.

\DndMonsterSection{Actions}
\DndMonsterAction{Bite}
\textsl{Melee Weapon Attack:} 3 to hit, reach 5 ft., one target. \textsl{Hit:} 4 (1d6 + 1) piercing damage, and the target is grappled (escape DC 11). Until this grapple ends, the target is restrained, and the frog can't bite another target.

\DndMonsterAction{Swallow}
The frog makes one bite attack against a Small or smaller target it is grappling. If the attack hits, the target is swallowed, and the grapple ends. The swallowed target is blinded and restrained, it has total cover against attacks and other effects outside the frog, and it takes 5 (2d4) acid damage at the start of each of the frog's turns. The frog can have only one target swallowed at a time. If the frog dies, a swallowed creature is no longer restrained by it and can escape from the corpse using 5 feet of movement, exiting prone.

\end{DndMonster}



\begin{DndMonster}[float=t]{Giant Goat}
\DndMonsterType{Large beast, unaligned}
\DndMonsterBasics[
armor-class = {11 (natural armor)},
hit-points = {\DndDice{3d10 + 3}},
speed = {40 ft.}
]
\DndMonsterAbilityScores[
 str = 17,
 dex = 11,
 con = 12,
 int = 3,
 wis = 12,
 cha = 6
]
\DndMonsterDetails[
challenge = 1/2,
]
\DndMonsterAction{Charge}
If the goat moves at least 20 feet straight toward a target and then hits it with a ram attack on the same turn, the target takes an extra 5 (2d4) bludgeoning damage. If the target is a creature, it must succeed on a DC 13 Strength saving throw or be knocked prone.

\DndMonsterAction{Sure-Footed}
The goat has advantage on Strength and Dexterity saving throws made against effects that would knock it prone.

\DndMonsterSection{Actions}
\DndMonsterAction{Ram}
\textsl{Melee Weapon Attack:} 5 to hit, reach 5 ft., one target. \textsl{Hit:} 8 (2d4 + 3) bludgeoning damage.

\end{DndMonster}



\begin{DndMonster}[float=t]{Giant Lizard}
\DndMonsterType{Large beast, unaligned}
\DndMonsterBasics[
armor-class = {12 (natural armor)},
hit-points = {\DndDice{3d10 + 3}},
speed = {30 ft., climb 30 ft.}
]
\DndMonsterAbilityScores[
 str = 15,
 dex = 12,
 con = 13,
 int = 2,
 wis = 10,
 cha = 5
]
\DndMonsterDetails[
senses = {darkvision 30 ft., passive perception 10},
challenge = 1/4,
]
\DndMonsterSection{Actions}
\DndMonsterAction{Bite}
\textsl{Melee Weapon Attack:} 4 to hit, reach 5 ft., one target. \textsl{Hit:} 6 (1d8 + 2) piercing damage.

\end{DndMonster}



\begin{DndMonster}[float=t]{Giant Poisonous Snake}
\DndMonsterType{Medium beast, unaligned}
\DndMonsterBasics[
armor-class = {14},
hit-points = {\DndDice{2d8 + 2}},
speed = {30 ft., swim 30 ft.}
]
\DndMonsterAbilityScores[
 str = 10,
 dex = 18,
 con = 13,
 int = 2,
 wis = 10,
 cha = 3
]
\DndMonsterDetails[
skills = {Perception +2},
senses = {blindsight 10 ft., passive perception 12},
challenge = 1/4,
]
\DndMonsterSection{Actions}
\DndMonsterAction{Bite}
\textsl{Melee Weapon Attack:} 6 to hit, reach 10 ft., one target. \textsl{Hit:} 6 (1d4 + 4) piercing damage, and the target must make a DC 11 Constitution saving throw, taking 10 (3d6) poison damage on a failed save, or half as much damage on a successful one.

\end{DndMonster}



\begin{DndMonster}[float=t]{Giant Rat}
\DndMonsterType{Small beast, unaligned}
\DndMonsterBasics[
armor-class = {12},
hit-points = {\DndDice{2d6}},
speed = {30 ft.}
]
\DndMonsterAbilityScores[
 str = 7,
 dex = 15,
 con = 11,
 int = 2,
 wis = 10,
 cha = 4
]
\DndMonsterDetails[
senses = {darkvision 60 ft., passive perception 10},
challenge = 1/8,
]
\DndMonsterAction{Keen Smell}
The rat has advantage on Wisdom (Perception) checks that rely on smell.

\DndMonsterAction{Pack Tactics}
The rat has advantage on an attack roll against a creature if at least one of the rat's allies is within 5 feet of the creature and the ally isn't incapacitated.

\DndMonsterSection{Actions}
\DndMonsterAction{Bite}
\textsl{Melee Weapon Attack:} 4 to hit, reach 5 ft., one target. \textsl{Hit:} 4 (1d4 + 2) piercing damage.

\end{DndMonster}



\begin{DndMonster}[float=t]{Giant Spider}
\DndMonsterType{Large beast, unaligned}
\DndMonsterBasics[
armor-class = {14 (natural armor)},
hit-points = {\DndDice{4d10 + 4}},
speed = {30 ft., climb 30 ft.}
]
\DndMonsterAbilityScores[
 str = 14,
 dex = 16,
 con = 12,
 int = 2,
 wis = 11,
 cha = 4
]
\DndMonsterDetails[
skills = {Stealth +7},
senses = {blindsight 10 ft., darkvision 60 ft., passive perception 10},
challenge = 1,
]
\DndMonsterAction{Spider Climb}
The spider can climb difficult surfaces, including upside down on ceilings, without needing to make an ability check.

\DndMonsterAction{Web Sense}
While in contact with a web, the spider knows the exact location of any other creature in contact with the same web.

\DndMonsterAction{Web Walker}
The spider ignores movement restrictions caused by webbing.

\DndMonsterSection{Actions}
\DndMonsterAction{Bite}
\textsl{Melee Weapon Attack:} 5 to hit, reach 5 ft., one creature. \textsl{Hit:} 7 (1d8 + 3) piercing damage, and the target must make a DC 11 Constitution saving throw, taking 9 (2d8) poison damage on a failed save, or half as much damage on a successful one. If the poison damage reduces the target to 0 hit points, the target is stable but poisoned for 1 hour, even after regaining hit points, and is paralyzed while poisoned in this way.

\DndMonsterAction{Web (Recharge 5\textendash 6)}
\textsl{Ranged Weapon Attack:} 5 to hit, range 30/60 ft., one creature. \textsl{Hit:} The target is restrained by webbing. As an action, the restrained target can make a DC 12 Strength check, bursting the webbing on a success. The webbing can also be attacked and destroyed (AC 10; hp 5; vulnerability to fire damage; immunity to bludgeoning, poison, and psychic damage).

\end{DndMonster}



\begin{DndMonster}[float=t]{Giant Toad}
\DndMonsterType{Large beast, unaligned}
\DndMonsterBasics[
armor-class = {11},
hit-points = {\DndDice{6d10 + 6}},
speed = {20 ft., swim 40 ft.}
]
\DndMonsterAbilityScores[
 str = 15,
 dex = 13,
 con = 13,
 int = 2,
 wis = 10,
 cha = 3
]
\DndMonsterDetails[
senses = {darkvision 30 ft., passive perception 10},
challenge = 1,
]
\DndMonsterAction{Amphibious}
The toad can breathe air and water.

\DndMonsterAction{Standing Leap}
The toad's long jump is up to 20 feet and its high jump is up to 10 feet, with or without a running start.

\DndMonsterSection{Actions}
\DndMonsterAction{Bite}
\textsl{Melee Weapon Attack:} 4 to hit, reach 5 ft., one target. \textsl{Hit:} 7 (1d10 + 2) piercing damage plus 5 (1d10) poison damage, and the target is grappled (escape DC 13). Until this grapple ends, the target is restrained, and the toad can't bite another target.

\DndMonsterAction{Swallow}
The toad makes one bite attack against a Medium or smaller target it is grappling. If the attack hits, the target is swallowed, and the grapple ends. The swallowed target is blinded and restrained, it has total cover against attacks and other effects outside the toad, and it takes 10 (3d6) acid damage at the start of each of the toad's turns. The toad can have only one target swallowed at a time.

If the toad dies, a swallowed creature is no longer restrained by it and can escape from the corpse using 5 feet of movement, exiting prone.

\end{DndMonster}



\begin{DndMonster}[float=t]{Giant Weasel}
\DndMonsterType{Medium beast, unaligned}
\DndMonsterBasics[
armor-class = {13},
hit-points = {\DndDice{2d8}},
speed = {40 ft.}
]
\DndMonsterAbilityScores[
 str = 11,
 dex = 16,
 con = 10,
 int = 4,
 wis = 12,
 cha = 5
]
\DndMonsterDetails[
skills = {Perception +3, Stealth +5},
senses = {darkvision 60 ft., passive perception 13},
challenge = 1/8,
]
\DndMonsterAction{Keen Hearing and Smell}
The weasel has advantage on Wisdom (Perception) checks that rely on hearing or smell.

\DndMonsterSection{Actions}
\DndMonsterAction{Bite}
\textsl{Melee Weapon Attack:} 5 to hit, reach 5 ft., one target. \textsl{Hit:} 5 (1d4 + 3) piercing damage.

\end{DndMonster}



\begin{DndMonster}[float=t]{Gibbering Mouther}
\DndMonsterType{Medium aberration, neutral}
\DndMonsterBasics[
armor-class = {9},
hit-points = {\DndDice{9d8 + 27}},
speed = {10 ft., swim 10 ft.}
]
\DndMonsterAbilityScores[
 str = 10,
 dex = 8,
 con = 16,
 int = 3,
 wis = 10,
 cha = 6
]
\DndMonsterDetails[
condition-immunities = {prone},
senses = {darkvision 60 ft., passive perception 10},
challenge = 2,
]
\DndMonsterAction{Aberrant Ground}
The ground in a 10-foot radius around the mouther is doughlike difficult terrain. Each creature that starts its turn in that area must succeed on a DC 10 Strength saving throw or have its speed reduced to 0 until the start of its next turn.

\DndMonsterAction{Gibbering}
The mouther babbles incoherently while it can see any creature and isn't incapacitated. Each creature that starts its turn within 20 feet of the mouther and can hear the gibbering must succeed on a DC 10 Wisdom saving throw. On a failure, the creature can't take reactions until the start of its next turn and rolls a d8 to determine what it does during its turn. On a 1 to 4, the creature does nothing. On a 5 or 6, the creature takes no action or bonus action and uses all its movement to move in a randomly determined direction. On a 7 or 8, the creature makes a melee attack against a randomly determined creature within its reach or does nothing if it can't make such an attack.

\DndMonsterSection{Actions}
\DndMonsterAction{Multiattack}
The gibbering mouther makes one bite attack and, if it can, uses its Blinding Spittle.

\DndMonsterAction{Bites}
\textsl{Melee Weapon Attack:} 2 to hit, reach 5 ft., one creature. \textsl{Hit:} 17 (5d6) piercing damage. If the target is Medium or smaller, it must succeed on a DC 10 Strength saving throw or be knocked prone. If the target is killed by this damage, it is absorbed into the mouther.

\DndMonsterAction{Blinding Spittle (Recharge 5\textendash 6)}
The mouther spits a chemical glob at a point it can see within 15 feet of it. The glob explodes in a blinding flash of light on impact. Each creature within 5 feet of the flash must succeed on a DC 13 Dexterity saving throw or be blinded until the end of the mouther's next turn.

\end{DndMonster}



\begin{DndMonster}[float=t]{Gibberling}
\DndMonsterType{Small fiend (demon), chaotic evil}
\DndMonsterBasics[
armor-class = {12},
hit-points = {\DndDice{2d6}},
speed = {30 ft.}
]
\DndMonsterAbilityScores[
 str = 13,
 dex = 14,
 con = 11,
 int = 5,
 wis = 7,
 cha = 5
]
\DndMonsterDetails[
condition-immunities = {charmed},
senses = {darkvision 120 ft., passive perception 8},
languages = {Gibberling},
challenge = 0,
]
\DndMonsterAction{Aversion to Fire}
If the gibberling takes fire damage, it has disadvantage on attack rolls and ability checks until the end of its next turn.

\DndMonsterAction{Incessant Gibberish}
Any non-gibberling that is within 30 feet of the gibberling and doesn't have the deafened condition has disadvantage on Constitution saving throws to maintain concentration on spells and similar effects.

\DndMonsterSection{Actions}
\DndMonsterAction{Bite}
\textsl{Melee Weapon Attack:} 3 to hit, reach 5 ft., one target. \textsl{Hit:} 3 (1d4 + 1) piercing damage.

\end{DndMonster}



\begin{DndMonster}[float=t]{Githyanki Knight}
\DndMonsterType{Medium humanoid (gith), lawful evil}
\DndMonsterBasics[
armor-class = {18 (\textit{plate armor})},
hit-points = {\DndDice{14d8 + 28}},
speed = {30 ft.}
]
\DndMonsterAbilityScores[
 str = 16,
 dex = 14,
 con = 15,
 int = 14,
 wis = 14,
 cha = 15
]
\DndMonsterDetails[
saving-throws = {Con +5, Int +5, Wis +5},
languages = {Gith},
challenge = 8,
]
\DndMonsterAction{Innate Spellcasting (Psionics)}
The githyanki's innate spellcasting ability is Intelligence (spell save DC 13, 5 to hit with spell attacks). It can innately cast the following spells, requiring no components:
\begin{DndMonsterSpells}
  \DndInnateSpellLevel{mage hand}
  \DndInnateSpellLevel[3]{jump}
  \DndInnateSpellLevel[1]{plane shift}
\end{DndMonsterSpells}
\DndMonsterSection{Actions}
\DndMonsterAction{Multiattack}
The githyanki makes two silver greatsword attacks.

\DndMonsterAction{Silver Greatsword}
\textsl{Melee Weapon Attack:} 9 to hit, reach 5 ft., one target. \textsl{Hit:} 13 (2d6 + 6) slashing damage plus 10 (3d6) psychic damage. This is a magic weapon attack. On a critical hit against a target in an astral body (as with the \textit{astral projection} spell), the githyanki can cut the silvery cord that tethers the target to its material body, instead of dealing damage.

\end{DndMonster}



\begin{DndMonster}[float=t]{Githyanki Warrior}
\DndMonsterType{Medium humanoid (gith), lawful evil}
\DndMonsterBasics[
armor-class = {17 (\textit{half plate armor})},
hit-points = {\DndDice{9d8 + 9}},
speed = {30 ft.}
]
\DndMonsterAbilityScores[
 str = 15,
 dex = 14,
 con = 12,
 int = 13,
 wis = 13,
 cha = 10
]
\DndMonsterDetails[
saving-throws = {Con +3, Int +3, Wis +3},
languages = {Gith},
challenge = 3,
]
\DndMonsterAction{Innate Spellcasting (Psionics)}
The githyanki's innate spellcasting ability is Intelligence. It can innately cast the following spells, requiring no components:
\begin{DndMonsterSpells}
  \DndInnateSpellLevel{mage hand}
  \DndInnateSpellLevel[3]{jump}
\end{DndMonsterSpells}
\DndMonsterSection{Actions}
\DndMonsterAction{Multiattack}
The githyanki makes two greatsword attacks.

\DndMonsterAction{Greatsword}
\textsl{Melee Weapon Attack:} 4 to hit, reach 5 ft., one target. \textsl{Hit:} 9 (2d6 + 2) slashing damage plus 7 (2d6) psychic damage.

\end{DndMonster}



\begin{DndMonster}[float=t]{Gladiator}
\DndMonsterType{Medium humanoid (any race), any alignment}
\DndMonsterBasics[
armor-class = {16 (studded leather armor)},
hit-points = {\DndDice{15d8 + 45}},
speed = {30 ft.}
]
\DndMonsterAbilityScores[
 str = 18,
 dex = 15,
 con = 16,
 int = 10,
 wis = 12,
 cha = 15
]
\DndMonsterDetails[
saving-throws = {Str +7, Dex +5, Con +6},
skills = {Athletics +10, Intimidation +5},
languages = {any one language (usually Common)},
challenge = 5,
]
\DndMonsterAction{Brave}
The gladiator has advantage on saving throws against being frightened.

\DndMonsterAction{Brute}
A melee weapon deals one extra die of its damage when the gladiator hits with it (included in the attack).

\DndMonsterSection{Actions}
\DndMonsterAction{Multiattack}
The gladiator makes three melee attacks or two ranged attacks.

\DndMonsterAction{Spear}
\textsl{Melee or Ranged Weapon Attack:} 7 to hit, reach 5 ft. and range 20/60 ft., one target. \textsl{Hit:} 11 (2d6 + 4) piercing damage, or 13 (2d8 + 4) piercing damage if used with two hands to make a melee attack.

\DndMonsterAction{Shield Bash}
\textsl{Melee Weapon Attack:} 7 to hit, reach 5 ft., one creature. \textsl{Hit:} 9 (2d4 + 4) bludgeoning damage. If the target is a Medium or smaller creature, it must succeed on a DC 15 Strength saving throw or be knocked prone.

\DndMonsterSection{Reactions}
\DndMonsterAction{Parry}
The gladiator adds 3 to its AC against one melee attack that would hit it. To do so, the gladiator must see the attacker and be wielding a melee weapon.

\end{DndMonster}



\begin{DndMonster}[float=t]{Gorgon}
\DndMonsterType{Large monstrosity, unaligned}
\DndMonsterBasics[
armor-class = {19 (natural armor)},
hit-points = {\DndDice{12d10 + 48}},
speed = {40 ft.}
]
\DndMonsterAbilityScores[
 str = 20,
 dex = 11,
 con = 18,
 int = 2,
 wis = 12,
 cha = 7
]
\DndMonsterDetails[
skills = {Perception +4},
condition-immunities = {petrified},
senses = {darkvision 60 ft., passive perception 14},
challenge = 5,
]
\DndMonsterAction{Trampling Charge}
If the gorgon moves at least 20 feet straight toward a creature and then hits it with a gore attack on the same turn, that target must succeed on a DC 16 Strength saving throw or be knocked prone. If the target is prone, the gorgon can make one attack with its hooves against it as a bonus action.

\DndMonsterSection{Actions}
\DndMonsterAction{Gore}
\textsl{Melee Weapon Attack:} 8 to hit, reach 5 ft., one target. \textsl{Hit:} 18 (2d12 + 5) piercing damage.

\DndMonsterAction{Hooves}
\textsl{Melee Weapon Attack:} 8 to hit, reach 5 ft., one target. \textsl{Hit:} 16 (2d10 + 5) bludgeoning damage.

\DndMonsterAction{Petrifying Breath (Recharge 5\textendash 6)}
The gorgon exhales petrifying gas in a 30-foot cone. Each creature in that area must succeed on a DC 13 Constitution saving throw. On a failed save, a target begins to turn to stone and is restrained. The restrained target must repeat the saving throw at the end of its next turn. On a success, the effect ends on the target. On a failure, the target is petrified until freed by the  \textit{greater restoration} spell or other magic.

\end{DndMonster}



\begin{DndMonster}[float=t]{Goristro}
\DndMonsterType{Huge fiend (demon), chaotic evil}
\DndMonsterBasics[
armor-class = {19 (natural armor)},
hit-points = {\DndDice{23d12 + 161}},
speed = {40 ft.}
]
\DndMonsterAbilityScores[
 str = 25,
 dex = 11,
 con = 25,
 int = 6,
 wis = 13,
 cha = 14
]
\DndMonsterDetails[
saving-throws = {Str +13, Dex +6, Con +13, Wis +7},
skills = {Perception +7},
damage-resistances = {cold; fire; lightning; bludgeoning, piercing, slashing from nonmagical attacks},
damage-immunities = {poison},
condition-immunities = {poisoned},
senses = {darkvision 120 ft., passive perception 17},
languages = {Abyssal},
challenge = 17,
]
\DndMonsterAction{Charge}
If the goristro moves at least 15 feet straight toward a target and then hits it with a gore attack on the same turn, the target takes an extra 38 (7d10) piercing damage. If the target is a creature, it must succeed on a DC 21 Strength saving throw or be pushed up to 20 feet away and knocked prone.

\DndMonsterAction{Labyrinthine Recall}
The goristro can perfectly recall any path it has traveled.

\DndMonsterAction{Magic Resistance}
The goristro has advantage on saving throws against spells and other magical effects.

\DndMonsterAction{Siege Monster}
The goristro deals double damage to objects and structures.

\DndMonsterSection{Actions}
\DndMonsterAction{Multiattack}
The goristro makes three attacks: two with its fists and one with its hoof.

\DndMonsterAction{Fist}
\textsl{Melee Weapon Attack:} 13 to hit, reach 10 ft., one target. \textsl{Hit:} 20 (3d8 + 7) bludgeoning damage.

\DndMonsterAction{Hoof}
\textsl{Melee Weapon Attack:} 13 to hit, reach 5 ft., one target. \textsl{Hit:} 23 (3d10 + 7) bludgeoning damage. If the target is a creature, it must succeed on a DC 21 Strength saving throw or be knocked prone.

\DndMonsterAction{Gore}
\textsl{Melee Weapon Attack:} 13 to hit, reach 10 ft., one target. \textsl{Hit:} 45 (7d10 + 7) piercing damage.

\end{DndMonster}



\begin{DndMonster}[float=t]{Gray Slaad}
\DndMonsterType{Medium aberration (shapechanger), chaotic neutral}
\DndMonsterBasics[
armor-class = {18 (natural armor)},
hit-points = {\DndDice{17d8 + 51}},
speed = {30 ft.}
]
\DndMonsterAbilityScores[
 str = 17,
 dex = 17,
 con = 16,
 int = 13,
 wis = 8,
 cha = 14
]
\DndMonsterDetails[
skills = {Arcana +5, Perception +7},
damage-resistances = {acid, cold, fire, lightning, thunder},
senses = {blindsight 60 ft., darkvision 60 ft., passive perception 17},
languages = {Slaad, telepathy 60 ft.},
challenge = 9,
]
\DndMonsterAction{Shapechanger}
The slaad can use its action to polymorph into a Small or Medium humanoid, or back into its true form. Its statistics, other than its size, are the same in each form. Any equipment it is wearing or carrying isn't transformed. It reverts to its true form if it dies.

\DndMonsterAction{Magic Resistance}
The slaad has advantage on saving throws against spells and other magical effects.

\DndMonsterAction{Magic Weapons}
The slaad's weapon attacks are magical.

\DndMonsterAction{Regeneration}
The slaad regains 10 hit points at the start of its turn if it has at least 1 hit point.

\DndMonsterAction{Innate Spellcasting}
The slaad's innate spellcasting ability is Charisma (spell save DC 14). The slaad can innately cast the following spells, requiring no material components:
\begin{DndMonsterSpells}
  \DndInnateSpellLevel{detect magic}
  \DndInnateSpellLevel[]{plane shift}
  \DndInnateSpellLevel[2]{fear}
\end{DndMonsterSpells}
\DndMonsterSection{Actions}
\DndMonsterAction{Multiattack}
The slaad makes three attacks: one with its bite and two with its claws or greatsword.

\DndMonsterAction{Bite (Slaad Form Only)}
\textsl{Melee Weapon Attack:} 7 to hit, reach 5 ft., one target. \textsl{Hit:} 6 (1d6 + 3) piercing damage.

\DndMonsterAction{Claws (Slaad Form Only)}
\textsl{Melee Weapon Attack:} 7 to hit, reach 5 ft., one target. \textsl{Hit:} 8 (1d10 + 3) slashing damage.

\DndMonsterAction{Greatsword}
\textsl{Melee Weapon Attack:} 7 to hit, reach 5 ft., one target. \textsl{Hit:} 10 (2d6 + 3) slashing damage.

\end{DndMonster}



\begin{DndMonster}[float=t]{Grell}
\DndMonsterType{Medium aberration, neutral evil}
\DndMonsterBasics[
armor-class = {12},
hit-points = {\DndDice{10d8 + 10}},
speed = {10 ft., fly 30 ft. \textit{(hover)}}
]
\DndMonsterAbilityScores[
 str = 15,
 dex = 14,
 con = 13,
 int = 12,
 wis = 11,
 cha = 9
]
\DndMonsterDetails[
skills = {Perception +4, Stealth +6},
damage-immunities = {lightning},
condition-immunities = {blinded, prone},
senses = {blindsight 60 ft. (blind beyond this radius), passive perception 14},
languages = {Grell},
challenge = 3,
]
\DndMonsterSection{Actions}
\DndMonsterAction{Multiattack}
The grell makes two attacks: one with its tentacles and one with its beak.

\DndMonsterAction{Tentacles}
\textsl{Melee Weapon Attack:} 4 to hit, reach 10 ft., one creature. \textsl{Hit:} 7 (1d10 + 2) piercing damage, and the target must succeed on a DC 11 Constitution saving throw or be poisoned for 1 minute. The poisoned target is paralyzed, and it can repeat the saving throw at the end of each of its turns, ending the effect on a success.

The target is also grappled (escape DC 15). If the target is Medium or smaller, it is also restrained until this grapple ends. While grappling the target, the grell has advantage on attack rolls against it and can 't use this attack against other targets. When the grell moves, any Medium or smaller target it is grappling moves with it.

\DndMonsterAction{Beak}
\textsl{Melee Weapon Attack:} 4 to hit, reach 5 ft., one target. \textsl{Hit:} 7 (2d4 + 2) piercing damage.

\end{DndMonster}



\begin{DndMonster}[float=t]{Grick}
\DndMonsterType{Medium monstrosity, neutral}
\DndMonsterBasics[
armor-class = {14 (natural armor)},
hit-points = {\DndDice{6d8}},
speed = {30 ft., climb 30 ft.}
]
\DndMonsterAbilityScores[
 str = 14,
 dex = 14,
 con = 11,
 int = 3,
 wis = 14,
 cha = 5
]
\DndMonsterDetails[
damage-resistances = {bludgeoning, piercing, slashing from nonmagical attacks},
senses = {darkvision 60 ft., passive perception 12},
challenge = 2,
]
\DndMonsterAction{Stone Camouflage}
The grick has advantage on Dexterity (Stealth) checks made to hide in rocky terrain.

\DndMonsterSection{Actions}
\DndMonsterAction{Multiattack}
The grick makes one attack with its tentacles. If that attack hits, the grick can make one beak attack against the same target.

\DndMonsterAction{Tentacles}
\textsl{Melee Weapon Attack:} 4 to hit, reach 5 ft., one target. \textsl{Hit:} 9 (2d6 + 2) slashing damage.

\DndMonsterAction{Beak}
\textsl{Melee Weapon Attack:} 4 to hit, reach 5 ft., one target. \textsl{Hit:} 5 (1d6 + 2) piercing damage.

\end{DndMonster}


\clearpage

\begin{DndMonster}[float=t]{Griffon}
\DndMonsterType{Large monstrosity, unaligned}
\DndMonsterBasics[
armor-class = {12},
hit-points = {\DndDice{7d10 + 21}},
speed = {30 ft., fly 80 ft.}
]
\DndMonsterAbilityScores[
 str = 18,
 dex = 15,
 con = 16,
 int = 2,
 wis = 13,
 cha = 8
]
\DndMonsterDetails[
skills = {Perception +5},
senses = {darkvision 60 ft., passive perception 15},
challenge = 2,
]
\DndMonsterAction{Keen Sight}
The griffon has advantage on Wisdom (Perception) checks that rely on sight.

\DndMonsterSection{Actions}
\DndMonsterAction{Multiattack}
The griffon makes two attacks: one with its beak and one with its claws.

\DndMonsterAction{Beak}
\textsl{Melee Weapon Attack:} 6 to hit, reach 5 ft., one target. \textsl{Hit:} 8 (1d8 + 4) piercing damage.

\DndMonsterAction{Claws}
\textsl{Melee Weapon Attack:} 6 to hit, reach 5 ft., one target. \textsl{Hit:} 11 (2d6 + 4) slashing damage.

\end{DndMonster}



\begin{DndMonster}[float=t]{Guard}
\DndMonsterType{Medium humanoid (any race), any alignment}
\DndMonsterBasics[
armor-class = {16 (\textit{chain shirt})},
hit-points = {\DndDice{2d8 + 2}},
speed = {30 ft.}
]
\DndMonsterAbilityScores[
 str = 13,
 dex = 12,
 con = 12,
 int = 10,
 wis = 11,
 cha = 10
]
\DndMonsterDetails[
skills = {Perception +2},
languages = {any one language (usually Common)},
challenge = 1/8,
]
\DndMonsterSection{Actions}
\DndMonsterAction{Spear}
\textsl{Melee or Ranged Weapon Attack:} 3 to hit, reach 5 ft. or range 20/60 ft., one target. \textsl{Hit:} 4 (1d6 + 1) piercing damage, or 5 (1d8 + 1) piercing damage if used with two hands to make a melee attack.

\end{DndMonster}



\begin{DndMonster}[float=t]{Guardian of Gorm}
\DndMonsterType{Medium humanoid, lawful good}
\DndMonsterBasics[
armor-class = {16 (\textit{chain mail})},
hit-points = {\DndDice{2d8}},
speed = {30 ft.}
]
\DndMonsterAbilityScores[
 str = 15,
 dex = 10,
 con = 10,
 int = 10,
 wis = 13,
 cha = 12
]
\DndMonsterDetails[
skills = {Athletics +4, Insight +3, Religion +2},
languages = {Common},
challenge = 1/8,
]
\DndMonsterAction{Brave}
The guardian has advantage on saving throws against the frightened condition.

\DndMonsterSection{Actions}
\DndMonsterAction{Lightning Mace}
\textsl{Melee Weapon Attack:} 4 to hit, reach 5 ft., one target. \textsl{Hit:} 5 (1d6 + 2) piercing damage plus 2 (1d4) lightning damage.

\DndMonsterAction{Handaxe}
\textsl{Melee or Ranged Weapon Attack:} 4 to hit, reach 5 ft. or range 20/60 ft., one target. \textsl{Hit:} 5 (1d6 + 2) slashing damage.

\end{DndMonster}



\begin{DndMonster}[float*=b,width=\textwidth + 8pt]{Gynosphinx}
\begin{multicols}{2}
\DndMonsterType{Large monstrosity, lawful neutral}
\DndMonsterBasics[
armor-class = {17 (natural armor)},
hit-points = {\DndDice{16d10 + 48}},
speed = {40 ft., fly 60 ft.}
]
\DndMonsterAbilityScores[
 str = 18,
 dex = 15,
 con = 16,
 int = 18,
 wis = 18,
 cha = 18
]
\DndMonsterDetails[
skills = {Arcana +12, History +12, Perception +8, Religion +8},
damage-resistances = {bludgeoning, piercing, slashing from nonmagical attacks},
damage-immunities = {psychic},
condition-immunities = {charmed, frightened},
senses = {truesight 120 ft., passive perception 18},
languages = {Common, Sphinx},
challenge = 11,
]
\DndMonsterAction{Inscrutable}
The sphinx is immune to any effect that would sense its emotions or read its thoughts, as well as any divination spell that it refuses. Wisdom (Insight) checks made to ascertain the sphinx's intentions or sincerity have disadvantage.

\DndMonsterAction{Magic Weapons}
The sphinx's weapon attacks are magical.

\DndMonsterAction{Spellcasting}
The sphinx is a 9th-level spellcaster. Its spellcasting ability is Intelligence (spell save DC 16, 8 to hit with spell attacks). It requires no material components to cast its spells. The sphinx has the following wizard spells prepared:
\begin{DndMonsterSpells}
  \DndMonsterSpellLevel{mage hand}
   \DndMonsterSpellLevel[1][4]{detect magic}
   \DndMonsterSpellLevel[2][3]{darkness}
   \DndMonsterSpellLevel[3][3]{dispel magic}
   \DndMonsterSpellLevel[4][3]{banishment}
   \DndMonsterSpellLevel[5][1]{legend lore}
\end{DndMonsterSpells}
\DndMonsterSection{Actions}
\DndMonsterAction{Multiattack}
The sphinx makes two claw attacks.

\DndMonsterAction{Claw}
\textsl{Melee Weapon Attack:} 8 to hit, reach 5 ft., one target. \textsl{Hit:} 13 (2d8 + 4) slashing damage.

\DndMonsterSection{Legendary Actions}
Gynosphinx can take 3 legendary actions, choosing from the options below. Only one legendary action can be used at a time and only at the end of another creature's turn. Gynosphinx regains spent legendary actions at the start of its turn.

\begin{DndMonsterLegendaryActions}
\DndMonsterLegendaryAction{Claw Attack}{The sphinx makes one claw attack.
}
\DndMonsterLegendaryAction{Teleport (Costs 2 Actions)}{The sphinx magically teleports, along with any equipment it is wearing or carrying, up to 120 feet to an unoccupied space it can see.
}
\DndMonsterLegendaryAction{Cast a Spell (Costs 3 Actions)}{The sphinx casts a spell from its list of prepared spells, using a spell slot as normal.
}
\end{DndMonsterLegendaryActions}
\end{multicols}
\end{DndMonster}



\begin{DndMonster}[float=t]{Hawk}
\DndMonsterType{Tiny beast, unaligned}
\DndMonsterBasics[
armor-class = {13},
hit-points = {\DndDice{1d4 - 1}},
speed = {10 ft., fly 60 ft.}
]
\DndMonsterAbilityScores[
 str = 5,
 dex = 16,
 con = 8,
 int = 2,
 wis = 14,
 cha = 6
]
\DndMonsterDetails[
skills = {Perception +4},
challenge = 0,
]
\DndMonsterAction{Keen Sight}
The hawk has advantage on Wisdom (Perception) checks that rely on sight.

\DndMonsterSection{Actions}
\DndMonsterAction{Talons}
\textsl{Melee Weapon Attack:} 5 to hit, reach 5 ft., one target. \textsl{Hit:} 1 slashing damage.

\end{DndMonster}



\begin{DndMonster}[float=t]{Helmed Horror}
\DndMonsterType{Medium construct, unaligned}
\DndMonsterBasics[
armor-class = {20 (\textit{plate armor})},
hit-points = {\DndDice{8d8 + 24}},
speed = {30 ft., fly 30 ft.}
]
\DndMonsterAbilityScores[
 str = 18,
 dex = 13,
 con = 16,
 int = 10,
 wis = 10,
 cha = 10
]
\DndMonsterDetails[
skills = {Perception +4},
damage-resistances = {bludgeoning, piercing, slashing from nonmagical attacks that aren't adamantine},
damage-immunities = {force, necrotic, poison},
condition-immunities = {blinded, charmed, deafened, frightened, paralyzed, petrified, poisoned, stunned},
senses = {blindsight 60 ft. (blind beyond this radius), passive perception 14},
languages = {understands the languages of its creator but can't speak},
challenge = 4,
]
\DndMonsterAction{Magic Resistance}
The helmed horror has advantage on saving throws against spells and other magical effects.

\DndMonsterAction{Spell Immunity}
The helmed horror is immune to three spells chosen by its creator. Typical immunities include \textit{fireball}, \textit{heat metal}, and \textit{lightning bolt}.

\DndMonsterSection{Actions}
\DndMonsterAction{Multiattack}
The helmed horror makes two longsword attacks.

\DndMonsterAction{Longsword}
\textsl{Melee Weapon Attack:} 6 to hit, reach 5 ft., one target. \textsl{Hit:} 8 (1d8 + 4) slashing damage, or 9 (1d10 + 4) slashing damage if used with two hands.

\end{DndMonster}



\begin{DndMonster}[float=t]{Hezrou}
\DndMonsterType{Large fiend (demon), chaotic evil}
\DndMonsterBasics[
armor-class = {16 (natural armor)},
hit-points = {\DndDice{13d10 + 65}},
speed = {30 ft.}
]
\DndMonsterAbilityScores[
 str = 19,
 dex = 17,
 con = 20,
 int = 5,
 wis = 12,
 cha = 13
]
\DndMonsterDetails[
saving-throws = {Str +7, Con +8, Wis +4},
damage-resistances = {cold; fire; lightning; bludgeoning, piercing, slashing from nonmagical attacks},
damage-immunities = {poison},
condition-immunities = {poisoned},
senses = {darkvision 120 ft., passive perception 11},
languages = {Abyssal, telepathy 120 ft.},
challenge = 8,
]
\DndMonsterAction{Magic Resistance}
The hezrou has advantage on saving throws against spells and other magical effects.

\DndMonsterAction{Stench}
Any creature that starts its turn within 10 feet of the hezrou must succeed on a DC 14 Constitution saving throw or be poisoned until the start of its next turn. On a successful saving throw, the creature is immune to the hezrou's stench for 24 hours.

\DndMonsterSection{Actions}
\DndMonsterAction{Multiattack}
The hezrou makes three attacks: one with its bite and two with its claws.

\DndMonsterAction{Bite}
\textsl{Melee Weapon Attack:} 7 to hit, reach 5 ft., one target. \textsl{Hit:} 15 (2d10 + 4) piercing damage.

\DndMonsterAction{Claws}
\textsl{Melee Weapon Attack:} 7 to hit, reach 5 ft., one target. \textsl{Hit:} 11 (2d6 + 4) slashing damage.

\end{DndMonster}



\begin{DndMonster}[float=t]{Hill Giant}
\DndMonsterType{Huge giant, chaotic evil}
\DndMonsterBasics[
armor-class = {13 (natural armor)},
hit-points = {\DndDice{10d12 + 40}},
speed = {40 ft.}
]
\DndMonsterAbilityScores[
 str = 21,
 dex = 8,
 con = 19,
 int = 5,
 wis = 9,
 cha = 6
]
\DndMonsterDetails[
skills = {Perception +2},
languages = {Giant},
challenge = 5,
]
\DndMonsterSection{Actions}
\DndMonsterAction{Multiattack}
The giant makes two greatclub attacks.

\DndMonsterAction{Greatclub}
\textsl{Melee Weapon Attack:} 8 to hit, reach 10 ft., one target. \textsl{Hit:} 18 (3d8 + 5) bludgeoning damage.

\DndMonsterAction{Rock}
\textsl{Ranged Weapon Attack:} 8 to hit, range 60/240 ft., one target. \textsl{Hit:} 21 (3d10 + 5) bludgeoning damage.

\end{DndMonster}



\begin{DndMonster}[float=t]{Hobgoblin}
\DndMonsterType{Medium humanoid (goblinoid), lawful evil}
\DndMonsterBasics[
armor-class = {18 (\textit{chain mail})},
hit-points = {\DndDice{2d8 + 2}},
speed = {30 ft.}
]
\DndMonsterAbilityScores[
 str = 13,
 dex = 12,
 con = 12,
 int = 10,
 wis = 10,
 cha = 9
]
\DndMonsterDetails[
senses = {darkvision 60 ft., passive perception 10},
languages = {Common, Goblin},
challenge = 1/2,
]
\DndMonsterAction{Martial Advantage}
Once per turn, the hobgoblin can deal an extra 7 (2d6) damage to a creature it hits with a weapon attack if that creature is within 5 feet of an ally of the hobgoblin that isn't incapacitated.

\DndMonsterSection{Actions}
\DndMonsterAction{Longsword}
\textsl{Melee Weapon Attack:} 3 to hit, reach 5 ft., one target. \textsl{Hit:} 5 (1d8 + 1) slashing damage, or 6 (1d10 + 1) slashing damage if used with two hands.

\DndMonsterAction{Longbow}
\textsl{Ranged Weapon Attack:} 3 to hit, range 150/600 ft., one target. \textsl{Hit:} 5 (1d8 + 1) piercing damage.

\end{DndMonster}



\begin{DndMonster}[float=t]{Hobgoblin Warlord}
\DndMonsterType{Medium humanoid (goblinoid), lawful evil}
\DndMonsterBasics[
armor-class = {20 (\textit{plate armor})},
hit-points = {\DndDice{13d8 + 39}},
speed = {30 ft.}
]
\DndMonsterAbilityScores[
 str = 16,
 dex = 14,
 con = 16,
 int = 14,
 wis = 11,
 cha = 15
]
\DndMonsterDetails[
saving-throws = {Int +5, Wis +3, Cha +5},
senses = {darkvision 60 ft., passive perception 10},
languages = {Common, Goblin},
challenge = 6,
]
\DndMonsterAction{Martial Advantage}
Once per turn, the hobgoblin can deal an extra 14 (4d6) damage to a creature it hits with a weapon attack if that creature is within 5 feet of an ally of the hobgoblin that isn't incapacitated.

\DndMonsterSection{Actions}
\DndMonsterAction{Multiattack}
The hobgoblin makes three melee attacks. Alternatively, it can make two ranged attacks with its javelins.

\DndMonsterAction{Longsword}
\textsl{Melee Weapon Attack:} 9 to hit, reach 5 ft., one target. \textsl{Hit:} 7 (1d8 + 3) slashing damage, or 8 (1d10 + 3) slashing damage if used with two hands.

\DndMonsterAction{Shield Bash}
\textsl{Melee Weapon Attack:} 9 to hit, reach 5 ft., one creature. \textsl{Hit:} 5 (1d4 + 3) bludgeoning damage. If the target is Large or smaller, it must succeed on a DC 14 Strength saving throw or be knocked prone.

\DndMonsterAction{Javelin}
\textsl{Melee or Ranged Weapon Attack:} 9 to hit, reach 5 ft. or range 30/120 ft., one target. \textsl{Hit:} 6 (1d6 + 3) piercing damage.

\DndMonsterAction{Leadership (Recharges after a Short or Long Rest)}
For 1 minute, the hobgoblin can utter a special command or warning whenever a nonhostile creature that it can see within 30 feet of it makes an attack roll or a saving throw. The creature can add a d4 to its roll provided it can hear and understand the hobgoblin. A creature can benefit from only one Leadership die at a time. This effect ends if the hobgoblin is incapacitated.

\DndMonsterSection{Reactions}
\DndMonsterAction{Parry}
The hobgoblin adds 3 to its AC against one melee attack that would hit it. To do so, the hobgoblin must see the attacker and be wielding a melee weapon.

\end{DndMonster}



\begin{DndMonster}[float=t]{Hook Horror}
\DndMonsterType{Large monstrosity, neutral}
\DndMonsterBasics[
armor-class = {15 (natural armor)},
hit-points = {\DndDice{10d10 + 20}},
speed = {30 ft., climb 30 ft.}
]
\DndMonsterAbilityScores[
 str = 18,
 dex = 10,
 con = 15,
 int = 6,
 wis = 12,
 cha = 7
]
\DndMonsterDetails[
skills = {Perception +3},
senses = {blindsight 60 ft., darkvision 120 ft., passive perception 13},
languages = {Hook Horror},
challenge = 3,
]
\DndMonsterAction{Echolocation}
The hook horror can't use its blindsight while deafened.

\DndMonsterAction{Keen Hearing}
The hook horror has advantage on Wisdom (Perception) checks that rely on hearing.

\DndMonsterSection{Actions}
\DndMonsterAction{Multiattack}
The hook horror makes two hook attacks.

\DndMonsterAction{Hook}
\textsl{Melee Weapon Attack:} 6 to hit, reach 10 ft., one target. \textsl{Hit:} 11 (2d6 + 4) piercing damage.

\end{DndMonster}



\begin{DndMonster}[float=t]{Horrid Plant}
\DndMonsterType{Large plant, unaligned}
\DndMonsterBasics[
armor-class = {6},
hit-points = {\DndDice{5d10 + 15}},
speed = {5 ft.}
]
\DndMonsterAbilityScores[
 str = 18,
 dex = 3,
 con = 17,
 int = 1,
 wis = 10,
 cha = 1
]
\DndMonsterDetails[
condition-immunities = {blinded, deafened},
senses = {blindsight 60 ft. (can't see beyond this radius), passive perception 10},
challenge = 4,
]
\DndMonsterAction{Horrid Plant Varieties}
A horrid plant comes in one of three varieties (choose or roll a d6): 1-2, dew drinker; 3-4, purple blossom; or 5-6, snapper saw. This form determines certain traits in this stat block.

\DndMonsterAction{False Appearance}
If the horrid plant is motionless at the start of combat, it has advantage on its initiative roll. Moreover, if a creature hasn't observed the horrid plant move or act, that creature must succeed on a DC 18 Intelligence (Nature) check to discern that the horrid plant isn't an ordinary plant.

\DndMonsterSection{Actions}
\DndMonsterAction{Multiattack}
The horrid plant makes two Tendril attacks.

\DndMonsterAction{Tendril}
\textsl{Melee Weapon Attack:} 6 to hit, reach 30 ft., one target. \textsl{Hit:} 7 (2d6) bludgeoning damage. If the target is a Huge or smaller creature, it has the grappled condition (escape DC 14), and the horrid plant can pull the target up to 25 feet closer to itself. Until the grapple ends, the target has the restrained condition, and the horrid plant can't use the same tendril on another target. The horrid plant has two tendrils.

\DndMonsterSection{Bonus Actions}
\DndMonsterAction{Sap Squirt (Purple Blossom Only)}
The horrid plant targets one creature it can see within 15 feet of itself. The target must succeed on a DC 14 Dexterity saving throw or take 28 (8d6) acid damage and become covered in acidic sap. This sap lasts for 1 minute or until a creature uses its action to scrape the sap off itself or another creature it can reach. A creature covered in sap takes 14 (4d6) acid damage at the start of each of its turns.

\DndMonsterAction{Spiked Leaves (Snapper Saw Only)}
Creatures within 10 feet of the horrid plant and not behind total cover must make a DC 14 Dexterity saving throw, taking 21 (6d6) slashing damage on a failed save or half as much damage on a successful one.

\DndMonsterAction{Vampiric Tendril (Dew Drinker Only)}
One creature grappled by the horrid plant must make a DC 13 Constitution saving throw, taking 10 (3d6) necrotic damage on a failed save or half as much damage on a successful one. The target's hit point maximum is reduced by an amount equal to the damage taken, and the horrid plant regains hit points equal to that amount. This reduction lasts until the target finishes a long rest. The target dies if this effect reduces its hit point maximum to 0.

\end{DndMonster}



\begin{DndMonster}[float=t]{Imp}
\DndMonsterType{Tiny fiend (devil), lawful evil}
\DndMonsterBasics[
armor-class = {13},
hit-points = {\DndDice{3d4 + 3}},
speed = {20 ft., fly 40 ft.}
]
\DndMonsterAbilityScores[
 str = 6,
 dex = 17,
 con = 13,
 int = 11,
 wis = 12,
 cha = 14
]
\DndMonsterDetails[
skills = {Deception +4, Insight +3, Persuasion +4, Stealth +5},
damage-resistances = {cold; bludgeoning, piercing, slashing from nonmagical attacks not made with silvered weapons},
damage-immunities = {fire, poison},
condition-immunities = {poisoned},
senses = {darkvision 120 ft., passive perception 11},
languages = {Infernal, Common},
challenge = 1,
]
\DndMonsterAction{Shapechanger}
The imp can use its action to polymorph into a beast form that resembles a rat (speed 20 ft.), a raven (20 ft., fly 60 ft.), or a spider (20 ft., climb 20 ft.), or back into its true form. Its statistics are the same in each form, except for the speed changes noted. Any equipment it is wearing or carrying isn't transformed. It reverts to its true form if it dies.

\DndMonsterAction{Devil's Sight}
Magical darkness doesn't impede the imp's darkvision.

\DndMonsterAction{Magic Resistance}
The imp has advantage on saving throws against spells and other magical effects.

\DndMonsterSection{Actions}
\DndMonsterAction{Sting (Bite in Beast Form)}
\textsl{Melee Weapon Attack:} 5 to hit, reach 5 ft., one target. \textsl{Hit:} 5 (1d4 + 3) piercing damage, and the target must make a DC 11 Constitution saving throw, taking 10 (3d6) poison damage on a failed save, or half as much damage on a successful one.

\DndMonsterAction{Invisibility}
The imp magically turns invisible until it attacks, or until its concentration ends (as if concentration on a spell). Any equipment the imp wears or carries is invisible with it.

\end{DndMonster}



\begin{DndMonster}[float=t]{Intellect Devourer}
\DndMonsterType{Tiny aberration, lawful evil}
\DndMonsterBasics[
armor-class = {12},
hit-points = {\DndDice{6d4 + 6}},
speed = {40 ft.}
]
\DndMonsterAbilityScores[
 str = 6,
 dex = 14,
 con = 13,
 int = 12,
 wis = 11,
 cha = 10
]
\DndMonsterDetails[
skills = {Perception +2, Stealth +4},
damage-resistances = {bludgeoning, piercing, slashing from nonmagical attacks},
condition-immunities = {blinded},
senses = {blindsight 60 ft. (blind beyond this radius), passive perception 12},
languages = {understands Deep Speech but can't speak, telepathy 60 ft.},
challenge = 2,
]
\DndMonsterAction{Detect Sentience}
The intellect devourer can sense the presence and location of any creature within 300 feet of it that has an Intelligence of 3 or higher, regardless of interposing barriers, unless the creature is protected by a \textit{mind blank} spell.

\DndMonsterSection{Actions}
\DndMonsterAction{Multiattack}
The intellect devourer makes one attack with its claws and uses Devour Intellect.

\DndMonsterAction{Claws}
\textsl{Melee Weapon Attack:} 4 to hit, reach 5 ft., one target. \textsl{Hit:} 7 (2d4 + 2) slashing damage.

\DndMonsterAction{Devour Intellect}
The intellect devourer targets one creature it can see within 10 feet of it that has a brain. The target must succeed on a DC 12 Intelligence saving throw against this magic or take 11 (2d10) psychic damage. Also on a failure, roll 3d6: If the total equals or exceeds the target's Intelligence score, that score is reduced to 0. The target is stunned until it regains at least one point of Intelligence.

\DndMonsterAction{Body Thief}
The intellect devourer initiates an Intelligence contest with an incapacitated humanoid within 5 feet of it that isn't protected by \textit{protection from evil and good}. If it wins the contest, the intellect devourer magically consumes the target's brain, teleports into the target's skull, and takes control of the target's body. While inside a creature, the intellect devourer has total cover against attacks and other effects originating outside its host. The intellect devourer retains its Intelligence, Wisdom, and Charisma scores, as well as its understanding of Deep Speech, its telepathy, and its traits. It otherwise adopts the target's statistics. It knows everything the creature knew, including spells and languages.

If the host body dies, the intellect devourer must leave it. A \textit{protection from evil and good} spell cast on the body drives the intellect devourer out. The intellect devourer is also forced out if the target regains its devoured brain by means of a \textit{wish}. By spending 5 feet of its movement, the intellect devourer can voluntarily leave the body, teleporting to the nearest unoccupied space within 5 feet of it. The body then dies, unless its brain is restored within 1 round.

\end{DndMonster}



\begin{DndMonster}[float=t]{Invisible Stalker}
\DndMonsterType{Medium elemental, neutral}
\DndMonsterBasics[
armor-class = {14},
hit-points = {\DndDice{16d8 + 32}},
speed = {50 ft., fly 50 ft. \textit{(hover)}}
]
\DndMonsterAbilityScores[
 str = 16,
 dex = 19,
 con = 14,
 int = 10,
 wis = 15,
 cha = 11
]
\DndMonsterDetails[
skills = {Perception +8, Stealth +10},
damage-resistances = {bludgeoning, piercing, slashing from nonmagical attacks},
damage-immunities = {poison},
condition-immunities = {exhaustion, grappled, paralyzed, petrified, poisoned, prone, restrained, unconscious},
senses = {darkvision 60 ft., passive perception 18},
languages = {Auran, understands Common but doesn't speak it},
challenge = 6,
]
\DndMonsterAction{Invisibility}
The stalker is invisible.

\DndMonsterAction{Faultless Tracker}
The stalker is given a quarry by its summoner. The stalker knows the direction and distance to its quarry as long as the two of them are on the same plane of existence. The stalker also knows the location of its summoner.

\DndMonsterSection{Actions}
\DndMonsterAction{Multiattack}
The stalker makes two slam attacks.

\DndMonsterAction{Slam}
\textsl{Melee Weapon Attack:} 6 to hit, reach 5 ft., one target. \textsl{Hit:} 10 (2d6 + 3) bludgeoning damage.

\end{DndMonster}



\begin{DndMonster}[float=t]{Jackalwere}
\DndMonsterType{Medium humanoid (shapechanger), chaotic evil}
\DndMonsterBasics[
armor-class = {12},
hit-points = {\DndDice{4d8}},
speed = {40 ft.}
]
\DndMonsterAbilityScores[
 str = 11,
 dex = 15,
 con = 11,
 int = 13,
 wis = 11,
 cha = 10
]
\DndMonsterDetails[
skills = {Deception +4, Perception +2, Stealth +4},
damage-immunities = {bludgeoning, piercing, slashing from nonmagical attacks that aren't silvered},
languages = {Common (can't speak in jackal form)},
challenge = 1/2,
]
\DndMonsterAction{Shapechanger}
The jackalwere can use its action to polymorph into a specific Medium human or a jackal-humanoid hybrid, or back into its true form (that of a Small jackal). Other than its size, its statistics are the same in each form. Any equipment it is wearing or carrying isn't transformed. It reverts to its true form if it dies.

\DndMonsterAction{Keen Hearing and Smell}
The jackalwere has advantage on Wisdom (Perception) checks that rely on hearing or smell.

\DndMonsterAction{Pack Tactics}
The jackalwere has advantage on an attack roll against a creature if at least one of the jackalwere's allies is within 5 feet of the creature and the ally isn't incapacitated.

\DndMonsterSection{Actions}
\DndMonsterAction{Bite (Jackal or Hybrid Form Only)}
\textsl{Melee Weapon Attack:} 4 to hit, reach 5 ft., one target. \textsl{Hit:} 4 (1d4 + 2) piercing damage.

\DndMonsterAction{Scimitar (Human or Hybrid Form Only)}
\textsl{Melee Weapon Attack:} 4 to hit, reach 5 ft., one target. \textsl{Hit:} 5 (1d6 + 2) slashing damage.

\DndMonsterAction{Sleep Gaze}
The jackalwere gazes at one creature it can see within 30 feet of it. The target must make a DC 10 Wisdom saving throw. On a failed save, the target succumbs to a magical slumber, falling unconscious for 10 minutes or until someone uses an action to shake the target awake. A creature that successfully saves against the effect is immune to this jackalwere's gaze for the next 24 hours. Undead and creatures immune to being charmed aren't affected by it.

\end{DndMonster}



\begin{DndMonster}[float=t]{Knight}
\DndMonsterType{Medium humanoid (any race), any alignment}
\DndMonsterBasics[
armor-class = {18 (\textit{plate armor})},
hit-points = {\DndDice{8d8 + 16}},
speed = {30 ft.}
]
\DndMonsterAbilityScores[
 str = 16,
 dex = 11,
 con = 14,
 int = 11,
 wis = 11,
 cha = 15
]
\DndMonsterDetails[
saving-throws = {Con +4, Wis +2},
languages = {any one language (usually Common)},
challenge = 3,
]
\DndMonsterAction{Brave}
The knight has advantage on saving throws against being frightened.

\DndMonsterSection{Actions}
\DndMonsterAction{Multiattack}
The knight makes two melee attacks.

\DndMonsterAction{Greatsword}
\textsl{Melee Weapon Attack:} 5 to hit, reach 5 ft., one target. \textsl{Hit:} 10 (2d6 + 3) slashing damage.

\DndMonsterAction{Heavy Crossbow}
\textsl{Ranged Weapon Attack:} 2 to hit, range 100/400 ft., one target. \textsl{Hit:} 5 (1d10) piercing damage.

\DndMonsterAction{Leadership (Recharges after a Short or Long Rest)}
For 1 minute, the knight can utter a special command or warning whenever a nonhostile creature that it can see within 30 feet of it makes an attack roll or a saving throw. The creature can add a d4 to its roll provided it can hear and understand the knight. A creature can benefit from only one Leadership die at a time. This effect ends if the knight is incapacitated.

\DndMonsterSection{Reactions}
\DndMonsterAction{Parry}
The knight adds 2 to its AC against one melee attack that would hit it. To do so, the knight must see the attacker and be wielding a melee weapon.

\end{DndMonster}



\begin{DndMonster}[float=t]{Lamia}
\DndMonsterType{Large monstrosity, chaotic evil}
\DndMonsterBasics[
armor-class = {13 (natural armor)},
hit-points = {\DndDice{13d10 + 26}},
speed = {30 ft.}
]
\DndMonsterAbilityScores[
 str = 16,
 dex = 13,
 con = 15,
 int = 14,
 wis = 15,
 cha = 16
]
\DndMonsterDetails[
skills = {Deception +7, Insight +4, Stealth +3},
senses = {darkvision 60 ft., passive perception 12},
languages = {Abyssal, Common},
challenge = 4,
]
\DndMonsterAction{Innate Spellcasting}
The lamia's innate spellcasting ability is Charisma (spell save DC 13). It can innately cast the following spells, requiring no material components.
\begin{DndMonsterSpells}
  \DndInnateSpellLevel{disguise self}
  \DndInnateSpellLevel[]{geas}
  \DndInnateSpellLevel[3]{charm person}
\end{DndMonsterSpells}
\DndMonsterSection{Actions}
\DndMonsterAction{Multiattack}
The lamia makes two attacks: one with its claws and one with its dagger or Intoxicating Touch.

\DndMonsterAction{Claws}
\textsl{Melee Weapon Attack:} 5 to hit, reach 5 ft., one target. \textsl{Hit:} 14 (2d10 + 3) slashing damage.

\DndMonsterAction{Dagger}
\textsl{Melee Weapon Attack:} 5 to hit, reach 5 ft., one target. \textsl{Hit:} 5 (1d4 + 3) piercing damage.

\DndMonsterAction{Intoxicating Touch}
\textsl{Melee Spell Attack:} 5 to hit, reach 5 ft., one creature. \textsl{Hit:} The target is magically cursed for 1 hour. Until the curse ends, the target has disadvantage on Wisdom saving throws and all ability checks.

\end{DndMonster}



\begin{DndMonster}[float=t]{Leprechaun}
\DndMonsterType{Small fey, neutral}
\DndMonsterBasics[
armor-class = {13},
hit-points = {\DndDice{8d6 + 24}},
speed = {30 ft.}
]
\DndMonsterAbilityScores[
 str = 6,
 dex = 17,
 con = 16,
 int = 12,
 wis = 14,
 cha = 18
]
\DndMonsterDetails[
saving-throws = {Con +5, Wis +4},
skills = {Perception +4, Sleight Of Hand +7, Stealth +5},
senses = {darkvision 60 ft., passive perception 14},
languages = {Common, Sylvan},
challenge = 4,
]
\DndMonsterAction{Industrious}
The leprechaun is proficient with all artisan's tools and adds double its proficiency bonus to ability checks made with them.

\DndMonsterAction{Reluctant Refusal}
When a creature offers the leprechaun the chance to partake in merriment or revelry such as a song, a dance, or a good meal, the leprechaun must succeed on a DC 15 Wisdom saving throw or have the charmed condition for 24 hours. While charmed in this way, the leprechaun partakes of the offering, treats the creature as a trusted friend, and seeks to defend it from harm. The charmed condition ends if the creature or any of its allies damage the leprechaun, force the leprechaun to make a saving throw, or steal from the leprechaun.

\DndMonsterAction{Spellcasting}
The leprechaun casts one of the following spells, requiring no material components and using Charisma as the spellcasting ability (spell save DC 14):
\begin{DndMonsterSpells}
  \DndInnateSpellLevel{Mending}
  \DndInnateSpellLevel[2]{Invisibility}
  \DndInnateSpellLevel[1]{Fabricate}
\end{DndMonsterSpells}
\DndMonsterSection{Actions}
\DndMonsterAction{Multiattack}
The leprechaun makes two Cobbler's Hammer attacks and can use Spellcasting.

\DndMonsterAction{Cobbler's Hammer}
\textsl{Melee Weapon Attack:} 5 to hit, reach 5 ft., one target. \textsl{Hit:} 6 (1d6 + 3) bludgeoning damage plus 13 (3d8) force damage. If the target is a creature, its speed is halved until the start of the leprechaun's next turn.

\DndMonsterAction{Gift of Luck (1/Day)}
The leprechaun touches a creature and magically gifts the target a measure of luck. The creature gains the leprechaun's Astonishing Luck reaction. The creature can use the reaction three times, after which this gift goes away. The leprechaun can revoke this gift from a creature at any time (no action required). A creature can benefit from only one leprechaun's Gift of Luck at a time.

\DndMonsterSection{Bonus Actions}
\DndMonsterAction{Cunning Trick}
The leprechaun takes the Disengage or Hide action or makes a Dexterity (Sleight of Hand) check.

\DndMonsterSection{Reactions}
\DndMonsterAction{Astonishing Luck}
When the leprechaun fails an ability check, an attack roll, or a saving throw, it can roll a new d20 and choose which roll to use, potentially turning the failure into a success.

\end{DndMonster}



\begin{DndMonster}[float*=b,width=\textwidth + 8pt]{Lich}
\begin{multicols}{2}
\DndMonsterType{Medium undead, any evil alignment}
\DndMonsterBasics[
armor-class = {17 (natural armor)},
hit-points = {\DndDice{18d8 + 54}},
speed = {30 ft.}
]
\DndMonsterAbilityScores[
 str = 11,
 dex = 16,
 con = 16,
 int = 20,
 wis = 14,
 cha = 16
]
\DndMonsterDetails[
saving-throws = {Con +10, Int +12, Wis +9},
skills = {Arcana +19, History +12, Insight +9, Perception +9},
damage-resistances = {cold, lightning, necrotic},
damage-immunities = {poison; bludgeoning, piercing, slashing from nonmagical attacks},
condition-immunities = {charmed, exhaustion, frightened, paralyzed, poisoned},
senses = {truesight 120 ft., passive perception 19},
languages = {Common plus up to five other languages},
challenge = 21,
]
\DndMonsterAction{Legendary Resistance (3/Day)}
If the lich fails a saving throw, it can choose to succeed instead.

\DndMonsterAction{Rejuvenation}
If it has a phylactery, a destroyed lich gains a new body in 1d10 days, regaining all its hit points and becoming active again. The new body appears within 5 feet of the phylactery.

\DndMonsterAction{Turn Resistance}
The lich has advantage on saving throws against any effect that turns undead.

\DndMonsterAction{Spellcasting}
The lich is an 18th-level spellcaster. Its spellcasting ability is Intelligence (spell save DC 20, 12 to hit with spell attacks). The lich has the following wizard spells prepared:
\begin{DndMonsterSpells}
  \DndMonsterSpellLevel{mage hand}
   \DndMonsterSpellLevel[1][4]{detect magic}
   \DndMonsterSpellLevel[2][3]{detect thoughts}
   \DndMonsterSpellLevel[3][3]{animate dead}
   \DndMonsterSpellLevel[4][3]{blight}
   \DndMonsterSpellLevel[5][3]{cloudkill}
   \DndMonsterSpellLevel[6][1]{disintegrate}
   \DndMonsterSpellLevel[7][1]{finger of death}
   \DndMonsterSpellLevel[8][1]{dominate monster}
\end{DndMonsterSpells}
\DndMonsterSection{Actions}
\DndMonsterAction{Paralyzing Touch}
\textsl{Melee Spell Attack:} 12 to hit, reach 5 ft., one creature. \textsl{Hit:} 10 (3d6) cold damage. The target must succeed on a DC 18 Constitution saving throw or be paralyzed for 1 minute. The target can repeat the saving throw at the end of each of its turns, ending the effect on itself on a success.

\DndMonsterSection{Legendary Actions}
Lich can take 3 legendary actions, choosing from the options below. Only one legendary action can be used at a time and only at the end of another creature's turn. Lich regains spent legendary actions at the start of its turn.

\begin{DndMonsterLegendaryActions}
\DndMonsterLegendaryAction{Cantrip}{The lich casts a cantrip.
}
\DndMonsterLegendaryAction{Paralyzing Touch (Costs 2 Actions)}{The lich uses its Paralyzing Touch.
}
\DndMonsterLegendaryAction{Frightening Gaze (Costs 2 Actions)}{The lich fixes its gaze on one creature it can see within 10 feet of it. The target must succeed on a DC 18 Wisdom saving throw against this magic or become frightened for 1 minute. The frightened target can repeat the saving throw at the end of each of its turns, ending the effect on itself on a success. If a target's saving throw is successful or the effect ends for it, the target is immune to the lich's gaze for the next 24 hours.
}
\DndMonsterLegendaryAction{Disrupt Life (Costs 3 Actions)}{Each non-undead creature within 20 feet of the lich must make a DC 18 Constitution saving throw against this magic, taking 21 (6d6) necrotic damage on a failed save, or half as much damage on a successful one.
}
\end{DndMonsterLegendaryActions}
\DndMonsterSection{Lair Actions}
On initiative count 20 (losing initiative ties), the lich can take a lair action to cause one of the following magical effects; the lich can't use the same effect two rounds in a row:


\begin{itemize}
\item The lich rolls a d8 and regains a spell slot of that level or lower. If it has no spent spell slots of that level or lower, nothing happens.
\item The lich targets one creature it can see within 30 feet of it. A crackling cord of negative energy tethers the lich to the target. Whenever the lich takes damage, the target must make a DC 18 Constitution saving throw. On a failed save, the lich takes half the damage (rounded down), and the target takes the remaining damage. This tether lasts until initiative count 20 on the next round or until the lich or the target is no longer in the lich's lair.
\item The lich calls forth the spirits of creatures that died in its lair. These apparitions materialize and attack one creature that the lich can see within 60 feet of it. The target must succeed on a DC 18 Constitution saving throw, taking 52 (15d6) necrotic damage on a failed save, or half as much damage on a success. The apparitions then disappear.
\end{itemize}

\end{multicols}
\end{DndMonster}


\clearpage

\begin{DndMonster}[float=t]{Mage}
\DndMonsterType{Medium humanoid (any race), any alignment}
\DndMonsterBasics[
armor-class = {12 (15 with \textit{mage armor})},
hit-points = {\DndDice{9d8}},
speed = {30 ft.}
]
\DndMonsterAbilityScores[
 str = 9,
 dex = 14,
 con = 11,
 int = 17,
 wis = 12,
 cha = 11
]
\DndMonsterDetails[
saving-throws = {Int +6, Wis +4},
skills = {Arcana +6, History +6},
languages = {any four languages},
challenge = 6,
]
\DndMonsterAction{Spellcasting}
The mage is a 9th-level spellcaster. Its spellcasting ability is Intelligence (spell save DC 14, 6 to hit with spell attacks). The mage has the following wizard spells prepared:
\begin{DndMonsterSpells}
  \DndMonsterSpellLevel{fire bolt}
   \DndMonsterSpellLevel[1][4]{detect magic}
   \DndMonsterSpellLevel[2][3]{misty step}
   \DndMonsterSpellLevel[3][3]{counterspell}
   \DndMonsterSpellLevel[4][3]{greater invisibility}
   \DndMonsterSpellLevel[5][1]{cone of cold}
\end{DndMonsterSpells}
\DndMonsterSection{Actions}
\DndMonsterAction{Dagger}
\textsl{Melee or Ranged Weapon Attack:} 5 to hit, reach 5 ft. or range 20/60 ft., one target. \textsl{Hit:} 4 (1d4 + 2) piercing damage.

\end{DndMonster}



\begin{DndMonster}[float=t]{Mage of Usamigaras}
\DndMonsterType{Medium humanoid, chaotic neutral}
\DndMonsterBasics[
armor-class = {11},
hit-points = {\DndDice{2d8}},
speed = {30 ft.}
]
\DndMonsterAbilityScores[
 str = 10,
 dex = 12,
 con = 10,
 int = 15,
 wis = 10,
 cha = 13
]
\DndMonsterDetails[
skills = {Arcana +4, Deception +5, Sleight Of Hand +3},
languages = {Common},
challenge = 1/8,
]
\DndMonsterAction{Spellcasting}
The mage casts one of the following spells, using Intelligence as the spellcasting ability (spell save DC 12):
\begin{DndMonsterSpells}
  \DndInnateSpellLevel{Light}
  \DndInnateSpellLevel[1]{Charm Person}
\end{DndMonsterSpells}
\DndMonsterSection{Actions}
\DndMonsterAction{Arcane Burst}
\textsl{Melee or Ranged Spell Attack:} 4 to hit, reach 5 ft. or range 120 ft., one target. \textsl{Hit:} 7 (1d10 + 2) force damage.

\end{DndMonster}



\begin{DndMonster}[float=t]{Mammoth}
\DndMonsterType{Huge beast, unaligned}
\DndMonsterBasics[
armor-class = {13 (natural armor)},
hit-points = {\DndDice{11d12 + 55}},
speed = {40 ft.}
]
\DndMonsterAbilityScores[
 str = 24,
 dex = 9,
 con = 21,
 int = 3,
 wis = 11,
 cha = 6
]
\DndMonsterDetails[
challenge = 6,
]
\DndMonsterAction{Trampling Charge}
If the mammoth moves at least 20 feet straight toward a creature and then hits it with a gore attack on the same turn, that target must succeed on a DC 18 Strength saving throw or be knocked prone. If the target is prone, the mammoth can make one stomp attack against it as a bonus action.

\DndMonsterSection{Actions}
\DndMonsterAction{Gore}
\textsl{Melee Weapon Attack:} 10 to hit, reach 10 ft., one target. \textsl{Hit:} 25 (4d8 + 7) piercing damage.

\DndMonsterAction{Stomp}
\textsl{Melee Weapon Attack:} 10 to hit, reach 5 ft., one prone creature. \textsl{Hit:} 29 (4d10 + 7) bludgeoning damage.

\end{DndMonster}



\begin{DndMonster}[float=t]{Marid}
\DndMonsterType{Large elemental, chaotic neutral}
\DndMonsterBasics[
armor-class = {17 (natural armor)},
hit-points = {\DndDice{17d10 + 136}},
speed = {30 ft., fly 60 ft., swim 90 ft.}
]
\DndMonsterAbilityScores[
 str = 22,
 dex = 12,
 con = 26,
 int = 18,
 wis = 17,
 cha = 18
]
\DndMonsterDetails[
saving-throws = {Dex +5, Wis +7, Cha +8},
damage-resistances = {acid, cold, lightning},
senses = {blindsight 30 ft., darkvision 120 ft., passive perception 13},
languages = {Aquan},
challenge = 11,
]
\DndMonsterAction{Amphibious}
The marid can breathe air and water.

\DndMonsterAction{Elemental Demise}
If the marid dies, its body disintegrates into a burst of water and foam, leaving behind only equipment the marid was wearing or carrying.

\DndMonsterAction{Innate Spellcasting}
The marid's innate spellcasting ability is Charisma (spell save DC 16, 8 to hit with spell attacks). It can innately cast the following spells, requiring no material components:
\begin{DndMonsterSpells}
  \DndInnateSpellLevel{create or destroy water}
  \DndInnateSpellLevel[3]{tongues}
  \DndInnateSpellLevel[1]{conjure elemental}
\end{DndMonsterSpells}
\DndMonsterSection{Actions}
\DndMonsterAction{Multiattack}
The marid makes two trident attacks.

\DndMonsterAction{Trident}
\textsl{Melee or Ranged Weapon Attack:} 10 to hit, reach 5 ft. or range 20/60 ft., one target. \textsl{Hit:} 13 (2d6 + 6) piercing damage, or 15 (2d8 + 6) piercing damage if used with two hands to make a melee attack.

\DndMonsterAction{Water Jet}
The marid magically shoots water in a 60-foot line that is 5 feet wide. Each creature in that line must make a DC 16 Dexterity saving throw. On a failure, a target takes 21 (6d6) bludgeoning damage and, if it is Huge or smaller, is pushed up to 20 feet away from the marid and knocked prone. On a success, a target takes half the bludgeoning damage, but is neither pushed nor knocked prone.

\end{DndMonster}



\begin{DndMonster}[float=t]{Maschin-i-Bozorg}
\DndMonsterType{Large construct, unaligned}
\DndMonsterBasics[
armor-class = {17 (natural armor)},
hit-points = {\DndDice{9d10 + 45}},
speed = {20 ft.}
]
\DndMonsterAbilityScores[
 str = 18,
 dex = 16,
 con = 20,
 int = 2,
 wis = 10,
 cha = 1
]
\DndMonsterDetails[
damage-immunities = {poison, psychic},
condition-immunities = {charmed, deafened, exhaustion, frightened, paralyzed, poisoned, unconscious},
senses = {darkvision 120 ft., passive perception 10},
languages = {understands Gnomish but can't speak},
challenge = 8,
]
\DndMonsterAction{Overheat}
When the maschin-i-bozorg is reduced to 0 hit points, its power source overloads, briefly superheating its outer shell. Each creature within 10 feet of the maschin-i-bozorg must make a DC 16 Constitution saving throw, taking 7 (2d6) fire damage on a failed save or half as much damage on a successful one.

\DndMonsterSection{Actions}
\DndMonsterAction{Multiattack}
The maschin-i-bozorg makes two Poison Jab attacks.

\DndMonsterAction{Poison Jab}
\textsl{Melee or Ranged Weapon Attack:} 6 to hit, reach 5 ft. or range 30/120 ft., one target. \textsl{Hit:} 13 (4d4 + 3) piercing damage, and the target must succeed on a DC 16 Constitution saving throw or have the poisoned condition for 1 minute. The target can repeat the saving throw at the end of each of its turns, ending the effect on itself on a success.

\DndMonsterAction{Steam Jet (Recharge 5\textendash 6)}
The maschin-i-bozorg emits scalding steam in a 30-foot cone. Each creature in that area must make a DC 16 Constitution saving throw, taking 28 (8d6) fire damage on a failed save or half as much damage on a successful one.

\DndMonsterSection{Bonus Actions}
\DndMonsterAction{Crushing Stride}
The maschin-i-bozorg moves up to its speed in a straight line. During this movement, it can enter Medium and smaller creatures' spaces. A creature whose space the maschin-i-bozorg enters must make a DC 15 Dexterity saving throw. On a successful save, the creature is pushed to the nearest unoccupied space out of the maschin-i-bozorg's path. On a failed save, the creature takes 10 (3d6) bludgeoning damage and has the prone condition.

If the maschin-i-bozorg remains in the prone creature's space, the creature also has the restrained condition until it's no longer in the same space as the maschin-i-bozorg. While restrained in this way, the creature, or another creature within 5 feet of it, can use its action to make a DC 15 Strength (Athletics) check. On a successful check, the creature is shunted to an unoccupied space of its choice within 5 feet of the maschin-i-bozorg.

\end{DndMonster}



\begin{DndMonster}[float=t]{Mastiff}
\DndMonsterType{Medium beast, unaligned}
\DndMonsterBasics[
armor-class = {12},
hit-points = {\DndDice{1d8 + 1}},
speed = {40 ft.}
]
\DndMonsterAbilityScores[
 str = 13,
 dex = 14,
 con = 12,
 int = 3,
 wis = 12,
 cha = 7
]
\DndMonsterDetails[
skills = {Perception +3},
challenge = 1/8,
]
\DndMonsterAction{Keen Hearing and Smell}
The mastiff has advantage on Wisdom (Perception) checks that rely on hearing or smell.

\DndMonsterSection{Actions}
\DndMonsterAction{Bite}
\textsl{Melee Weapon Attack:} 3 to hit, reach 5 ft., one target. \textsl{Hit:} 4 (1d6 + 1) piercing damage. If the target is a creature, it must succeed on a DC 11 Strength saving throw or be knocked prone.

\end{DndMonster}



\begin{DndMonster}[float=t]{Memory Web}
\DndMonsterType{Large aberration, neutral evil}
\DndMonsterBasics[
armor-class = {14},
hit-points = {\DndDice{6d10 + 6}},
speed = {40 ft., climb 40 ft.}
]
\DndMonsterAbilityScores[
 str = 16,
 dex = 18,
 con = 13,
 int = 14,
 wis = 14,
 cha = 3
]
\DndMonsterDetails[
damage-resistances = {fire},
condition-immunities = {blinded, deafened},
senses = {blindsight 60 ft. (can't see beyond this radius), passive perception 12},
challenge = 4,
]
\DndMonsterAction{Damage Transfer}
While it is grappling a creature, the memory web takes only half the damage dealt to it, and the creature grappled by the web takes the other half.

\DndMonsterAction{False Appearance}
If the memory web is motionless at the start of combat, it has advantage on its initiative roll. Moreover, if a creature hasn't observed the memory web move or act, that creature must succeed on a DC 18 Intelligence (Investigation) check to discern that the memory web is animate.

\DndMonsterAction{Memory Flood}
When the memory web is reduced to 0 hit points, it discharges any memories it consumed over the past 24 hours in a telepathic deluge. Hazy, dreamlike visions of these discharged memories lodge in the minds of creatures up to 120 feet from the memory web.

\DndMonsterSection{Actions}
\DndMonsterAction{Ensnare}
\textsl{Melee Weapon Attack:} 5 to hit, reach 5 ft., one Large or smaller creature. \textsl{Hit:} The target has the grappled condition (escape DC 13). Until this grapple ends, the target has the restrained condition and takes 7 (1d8 + 3) bludgeoning damage at the start of each of its turns. The memory web can grapple only one creature at a time.

\DndMonsterSection{Bonus Actions}
\DndMonsterAction{Drain Memories}
The memory web targets one creature grappled by it. The target must make a DC 12 Intelligence saving throw. Constructs, Oozes, Plants, and Undead succeed on the save automatically. On a failed save, the target takes 5 (2d4) psychic damage and becomes memory drained until it finishes a long rest or the memory web is destroyed.

While memory drained, the target must roll a d4 each time it makes an ability check or attack roll, subtracting the d4 roll from it. Each time the target is memory drained beyond the first, the die size increases by one: the d4 becomes a d6, the d6 becomes a d8, and so on until the die becomes a d20, at which point the target has the unconscious condition for 1 hour. If a memory drained creature is the target of the Greater Restoration or \textit{Heal} spell, the memory drained effect ends on it. On a successful save, the target takes half as much damage only.

\end{DndMonster}



\begin{DndMonster}[float=t]{Merrow}
\DndMonsterType{Large monstrosity, chaotic evil}
\DndMonsterBasics[
armor-class = {13 (natural armor)},
hit-points = {\DndDice{6d10 + 12}},
speed = {10 ft., swim 40 ft.}
]
\DndMonsterAbilityScores[
 str = 18,
 dex = 10,
 con = 15,
 int = 8,
 wis = 10,
 cha = 9
]
\DndMonsterDetails[
senses = {darkvision 60 ft., passive perception 10},
languages = {Abyssal, Aquan},
challenge = 2,
]
\DndMonsterAction{Amphibious}
The merrow can breathe air and water.

\DndMonsterSection{Actions}
\DndMonsterAction{Multiattack}
The merrow makes two attacks: one with its bite and one with its claws or harpoon.

\DndMonsterAction{Bite}
\textsl{Melee Weapon Attack:} 6 to hit, reach 5 ft., one target. \textsl{Hit:} 8 (1d8 + 4) piercing damage.

\DndMonsterAction{Claws}
\textsl{Melee Weapon Attack:} 6 to hit, reach 5 ft., one target. \textsl{Hit:} 9 (2d4 + 4) slashing damage.

\DndMonsterAction{Harpoon}
\textsl{Melee or Ranged Weapon Attack:} 6 to hit, reach 5 ft. or range 20/60 ft., one target. \textsl{Hit:} 11 (2d6 + 4) piercing damage. If the target is a Huge or smaller creature, it must succeed on a Strength contest against the merrow or be pulled up to 20 feet toward the merrow.

\end{DndMonster}



\begin{DndMonster}[float=t]{Mimic}
\DndMonsterType{Medium monstrosity (shapechanger), neutral}
\DndMonsterBasics[
armor-class = {12 (natural armor)},
hit-points = {\DndDice{9d8 + 18}},
speed = {15 ft.}
]
\DndMonsterAbilityScores[
 str = 17,
 dex = 12,
 con = 15,
 int = 5,
 wis = 13,
 cha = 8
]
\DndMonsterDetails[
skills = {Stealth +5},
damage-immunities = {acid},
condition-immunities = {prone},
senses = {darkvision 60 ft., passive perception 11},
challenge = 2,
]
\DndMonsterAction{Shapechanger}
The mimic can use its action to polymorph into an object or back into its true, amorphous form. Its statistics are the same in each form. Any equipment it is wearing or carrying isn't transformed. It reverts to its true form if it dies.

\DndMonsterAction{Adhesive (Object Form Only)}
The mimic adheres to anything that touches it. A Huge or smaller creature adhered to the mimic is also grappled by it (escape DC 13). Ability checks made to escape this grapple have disadvantage.

\DndMonsterAction{False Appearance (Object Form Only)}
While the mimic remains motionless, it is indistinguishable from an ordinary object.

\DndMonsterAction{Grappler}
The mimic has advantage on attack rolls against any creature grappled by it.

\DndMonsterSection{Actions}
\DndMonsterAction{Pseudopod}
\textsl{Melee Weapon Attack:} 5 to hit, reach 5 ft., one target. \textsl{Hit:} 7 (1d8 + 3) bludgeoning damage. If the mimic is in object form, the target is subjected to its Adhesive trait.

\DndMonsterAction{Bite}
\textsl{Melee Weapon Attack:} 5 to hit, reach 5 ft., one target. \textsl{Hit:} 7 (1d8 + 3) piercing damage plus 4 (1d8) acid damage.

\end{DndMonster}



\begin{DndMonster}[float=t]{Mind Flayer}
\DndMonsterType{Medium aberration, lawful evil}
\DndMonsterBasics[
armor-class = {15 (\textit{breastplate})},
hit-points = {\DndDice{13d8 + 13}},
speed = {30 ft.}
]
\DndMonsterAbilityScores[
 str = 11,
 dex = 12,
 con = 12,
 int = 19,
 wis = 17,
 cha = 17
]
\DndMonsterDetails[
saving-throws = {Int +7, Wis +6, Cha +6},
skills = {Arcana +7, Deception +6, Insight +6, Perception +6, Persuasion +6, Stealth +4},
senses = {darkvision 120 ft., passive perception 16},
languages = {Deep Speech, Undercommon, telepathy 120 ft.},
challenge = 7,
]
\DndMonsterAction{Magic Resistance}
The mind flayer has advantage on saving throws against spells and other magical effects.

\DndMonsterAction{Innate Spellcasting (Psionics)}
The mind flayer's innate spellcasting ability is Intelligence (spell save DC 15). It can innately cast the following spells, requiring no components:
\begin{DndMonsterSpells}
  \DndInnateSpellLevel{detect thoughts}
  \DndInnateSpellLevel[1]{dominate monster}
\end{DndMonsterSpells}
\DndMonsterSection{Actions}
\DndMonsterAction{Tentacles}
\textsl{Melee Weapon Attack:} 7 to hit, reach 5 ft., one creature. \textsl{Hit:} 15 (2d10 + 4) psychic damage. If the target is Medium or smaller, it is grappled (escape DC 15) and must succeed on a DC 15 Intelligence saving throw or be stunned until this grapple ends.

\DndMonsterAction{Extract Brain}
\textsl{Melee Weapon Attack:} 7 to hit, reach 5 ft., one incapacitated humanoid grappled by the mind flayer. \textsl{Hit:} The target takes 55 (10d10) piercing damage. If this damage reduces the target to 0 hit points, the mind flayer kills the target by extracting and devouring its brain.

\DndMonsterAction{Mind Blast (Recharge 5\textendash 6)}
The mind flayer magically emits psychic energy in a 60-foot cone. Each creature in that area must succeed on a DC 15 Intelligence saving throw or take 22 (4d8 + 4) psychic damage and be stunned for 1 minute. A creature can repeat the saving throw at the end of each of its turns, ending the effect on itself on a success.

\end{DndMonster}



\begin{DndMonster}[float=t]{Minotaur}
\DndMonsterType{Large monstrosity, chaotic evil}
\DndMonsterBasics[
armor-class = {14 (natural armor)},
hit-points = {\DndDice{9d10 + 27}},
speed = {40 ft.}
]
\DndMonsterAbilityScores[
 str = 18,
 dex = 11,
 con = 16,
 int = 6,
 wis = 16,
 cha = 9
]
\DndMonsterDetails[
skills = {Perception +7},
senses = {darkvision 60 ft., passive perception 17},
languages = {Abyssal},
challenge = 3,
]
\DndMonsterAction{Charge}
If the minotaur moves at least 10 feet straight toward a target and then hits it with a gore attack on the same turn, the target takes an extra 9 (2d8) piercing damage. If the target is a creature, it must succeed on a DC 14 Strength saving throw or be pushed up to 10 feet away and knocked prone.

\DndMonsterAction{Labyrinthine Recall}
The minotaur can perfectly recall any path it has traveled.

\DndMonsterAction{Reckless}
At the start of its turn, the minotaur can gain advantage on all melee weapon attack rolls it makes during that turn, but attack rolls against it have advantage until the start of its next turn.

\DndMonsterSection{Actions}
\DndMonsterAction{Greataxe}
\textsl{Melee Weapon Attack:} 6 to hit, reach 5 ft., one target. \textsl{Hit:} 17 (2d12 + 4) slashing damage.

\DndMonsterAction{Gore}
\textsl{Melee Weapon Attack:} 6 to hit, reach 5 ft., one target. \textsl{Hit:} 13 (2d8 + 4) piercing damage.

\end{DndMonster}



\begin{DndMonster}[float=t]{Mummy}
\DndMonsterType{Medium undead, lawful evil}
\DndMonsterBasics[
armor-class = {11 (natural armor)},
hit-points = {\DndDice{9d8 + 18}},
speed = {20 ft.}
]
\DndMonsterAbilityScores[
 str = 16,
 dex = 8,
 con = 15,
 int = 6,
 wis = 10,
 cha = 12
]
\DndMonsterDetails[
saving-throws = {Wis +2},
damage-vulnerabilities = {fire},
damage-resistances = {bludgeoning, piercing, slashing from nonmagical attacks},
damage-immunities = {necrotic, poison},
condition-immunities = {charmed, exhaustion, frightened, paralyzed, poisoned},
senses = {darkvision 60 ft., passive perception 10},
languages = {the languages it knew in life},
challenge = 3,
]
\DndMonsterSection{Actions}
\DndMonsterAction{Multiattack}
The mummy can use its Dreadful Glare and makes one attack with its rotting fist.

\DndMonsterAction{Rotting Fist}
\textsl{Melee Weapon Attack:} 5 to hit, reach 5 ft., one target. \textsl{Hit:} 10 (2d6 + 3) bludgeoning damage plus 10 (3d6) necrotic damage. If the target is a creature, it must succeed on a DC 12 Constitution saving throw or be cursed with mummy rot. The cursed target can't regain hit points, and its hit point maximum decreases by 10 (3d6) for every 24 hours that elapse. If the curse reduces the target's hit point maximum to 0, the target dies, and its body turns to dust. The curse lasts until removed by the \textit{remove curse} spell or other magic.

\DndMonsterAction{Dreadful Glare}
The mummy targets one creature it can see within 60 feet of it. If the target can see the mummy, it must succeed on a DC 11 Wisdom saving throw against this magic or become frightened until the end of the mummy's next turn. If the target fails the saving throw by 5 or more, it is also paralyzed for the same duration. A target that succeeds on the saving throw is immune to the Dreadful Glare of all mummies (but not mummy lords) for the next 24 hours.

\end{DndMonster}



\begin{DndMonster}[float=t]{Myconid Sovereign}
\DndMonsterType{Large plant, lawful neutral}
\DndMonsterBasics[
armor-class = {13 (natural armor)},
hit-points = {\DndDice{8d10 + 16}},
speed = {30 ft.}
]
\DndMonsterAbilityScores[
 str = 12,
 dex = 10,
 con = 14,
 int = 13,
 wis = 15,
 cha = 10
]
\DndMonsterDetails[
senses = {darkvision 120 ft., passive perception 12},
challenge = 2,
]
\DndMonsterAction{Distress Spores}
When the myconid takes damage, all other myconids within 240 feet of it can sense its pain.

\DndMonsterAction{Sun Sickness}
While in sunlight, the myconid has disadvantage on ability checks, attack rolls, and saving throws. The myconid dies if it spends more than 1 hour in direct sunlight.

\DndMonsterSection{Actions}
\DndMonsterAction{Multiattack}
The myconid uses either its Hallucination Spores or its Pacifying Spores, then makes a fist attack.

\DndMonsterAction{Fist}
\textsl{Melee Weapon Attack:} 3 to hit, reach 5 ft., one target. \textsl{Hit:} 8 (3d4 + 1) bludgeoning damage plus 7 (3d4) poison damage.

\DndMonsterAction{Animating Spores (3/Day)}
The myconid targets one corpse of a humanoid or a Large or smaller beast within 5 feet of it and releases spores at the corpse. In 24 hours, the corpse rises as a spore servant. The corpse stays animated for 1d4 + 1 weeks or until destroyed, and it can't be animated again in this way.

\DndMonsterAction{Hallucination Spores}
The myconid ejects spores at one creature it can see within 5 feet of it. The target must succeed on a DC 12 Constitution saving throw or be poisoned for 1 minute. The poisoned target is incapacitated while it hallucinates. The target can repeat the saving throw at the end of each of its turns, ending the effect on itself on a success.

\DndMonsterAction{Pacifying Spores}
The myconid ejects spores at one creature it can see within 5 feet of it. The target must succeed on a DC 12 Constitution saving throw or be stunned for 1 minute. The target can repeat the saving throw at the end of each of its turns, ending the effect on itself on a success.

\DndMonsterAction{Rapport Spores}
A 30-foot radius of spores extends from the myconid. These spores can go around corners and affect only creatures with an Intelligence of 2 or higher that aren't undead, constructs, or elementals. Affected creatures can communicate telepathically with one another while they are within 30 feet of each other. The effect lasts for 1 hour.

\end{DndMonster}



\begin{DndMonster}[float*=b,width=\textwidth + 8pt]{Nafas}
\begin{multicols}{2}
\DndMonsterType{Large elemental, chaotic good}
\DndMonsterBasics[
armor-class = {19 (natural armor)},
hit-points = {\DndDice{28d10 + 196}},
speed = {30 ft., fly 90 ft. \textit{(hover)}}
]
\DndMonsterAbilityScores[
 str = 23,
 dex = 18,
 con = 24,
 int = 15,
 wis = 18,
 cha = 23
]
\DndMonsterDetails[
saving-throws = {Dex +11, Wis +11, Cha +13},
skills = {Perception +11, Performance +13},
damage-resistances = {poison; bludgeoning, piercing, slashing from nonmagical attacks},
damage-immunities = {lightning, thunder},
condition-immunities = {charmed, frightened, poisoned},
senses = {truesight 120 ft., passive perception 21},
languages = {Auran, Common},
challenge = 23,
]
\DndMonsterAction{Dimensionally Bound}
Nafas can't leave the Infinite Staircase or be trapped within a container (such as an Iron Flask). Attempts to transport Nafas to another plane are wasted.

\DndMonsterAction{Last Wish}
When Nafas drops to 0 hit points, his body disintegrates into a whirl of multiversal dust that surrounds one creature responsible for his demise. That creature then hears Nafas's last wish: for the creature to take his place.

If the creature accepts, it is transformed into a noble djinni. The creature's game statistics are replaced by those of Nafas (including this trait), though it retains its name, alignment, and personality. The creature also inherits Nafas's palace and all it contains.

If the creature refuses, Nafas gains a new body in 1d10 days, regaining all his hit points and appearing in a random safe location on the Infinite Staircase.

\DndMonsterAction{Legendary Resistance (5/Day)}
If Nafas fails a saving throw, he can choose to succeed instead.

\DndMonsterAction{Magic Resistance}
Nafas has advantage on saving throws against spells and other magical effects.

\DndMonsterAction{Noble Genie}
Nafas doesn't suffer any of the penalties that normally follow casting the \textit{Wish} spell to produce an effect other than duplicating another spell.

\DndMonsterAction{Spellcasting}
Nafas casts one of the following spells, requiring no spell components and using Charisma as the spellcasting ability (spell save DC 21):
\begin{DndMonsterSpells}
  \DndInnateSpellLevel{Detect Evil and Good}
  \DndInnateSpellLevel[3]{Create Food and Water}
  \DndInnateSpellLevel[1]{Gaseous Form}
\end{DndMonsterSpells}
\DndMonsterSection{Actions}
\DndMonsterAction{Multiattack}
Nafas makes three Storm Shamshir attacks and uses Create Vortex.

\DndMonsterAction{Storm Shamshir}
\textsl{Melee Weapon Attack:} 13 to hit, reach 5 ft., one target. \textsl{Hit:} 13 (2d6 + 6) slashing damage plus 14 (4d6) lightning or thunder damage (Nafas's choice).

\DndMonsterAction{Create Vortex}
A 10-foot-radius, 60-foot-tall cylinder of swirling cosmic dust forms on a point Nafas can see within 120 feet of him. The vortex lasts as long as Nafas maintains concentration (as if concentration on a spell). When the vortex appears, each creature other than Nafas in the vortex's area must make a DC 22 Strength saving throw. On a failed save, a creature takes 36 (8d8) force damage and has the restrained condition. On a successful save, a creature takes half as much damage only and moves to the nearest unoccupied space outside the vortex.

On subsequent turns, Nafas can use this action to move the vortex up to 60 feet. When the vortex enters a creature's space for the first time on a turn, the creature must make the same saving throw as when the vortex first appeared. Creatures restrained by the vortex move with it.

A creature restrained by the vortex, or another creature that can reach it, can use its action to make a DC 22 Strength check. On a successful check, the restrained creature is no longer restrained and moves to the nearest unoccupied space outside the vortex.

\DndMonsterSection{Reactions}
\DndMonsterAction{Blowback}
Immediately after a creature Nafas can see ends its turn, Nafas exhales forceful winds in a 30-foot cone. Large or smaller creatures in that area must succeed on a DC 22 Strength saving throw or be pushed up to 15 feet away from him.

\DndMonsterAction{Zephyr Step}
In response to being hit by an attack roll, Nafas moves up to half his flying speed without provoking opportunity attacks.

\DndMonsterSection{Lair Actions}
On initiative count 20 (losing initiative ties), Nafas can take one of the following lair actions; he can't take the same lair action two rounds in a row:


\begin{description}
\item[Downdraft.]
Nafas targets one creature he can see within 120 feet of himself. A downward current of air surrounds the target, reducing its flying speed (if any) to 0 feet until the end of its next turn, and the target must succeed on a DC 21 Strength saving throw or have the prone condition.
\item[Lightning Strike.]
A bolt of lightning strikes a point Nafas can see within 120 feet of him. Each creature within 10 feet of that point must succeed on a DC 21 Dexterity saving throw or take 16 (3d10) lightning damage.
\item[Stardust.]
Nafas creates a 20-foot-radius sphere of multiversal dust centered on a point anywhere in his lair. The sphere spreads around corners, and its area is heavily obscured. The dust lasts until Nafas disperses it (no action required) or uses this lair action again, and it can't be dispersed by wind.
\end{description}

\DndMonsterSection{Regional Effects}
The region containing Nafas's lair is warped by his magic, creating one or more of the following effects:


\begin{description}
\item[Mirages.]
Sapient creatures within 3 miles of the lair frequently see illusions of their hearts' desires, be they fame, fortune, or long-departed comrades. These mirages don't hold up to physical inspection and vanish after a brief interaction.
\item[Safe Descents.]
Within 3 miles of the lair, winds buoy creatures that fall, ushering them safely to the nearest staircase. Such creatures descend at a rate of 60 feet per round and take no damage from falling. Aberrations, Fiends, Undead, and creatures that fall due to Nafas's or his allies' actions don't gain this benefit, instead falling as normal.
\item[Wishful Winds.]
The wishes of creatures across the multiverse, including those within the Infinite Staircase, can be heard within 3 miles of the lair as fleeting voices carried on gentle breezes.
\end{description}

If Nafas is destroyed, all these effects end immediately, resuming when he either gains a new body or is replaced by the creature that slays him (see the Last Wish trait in Nafas's stat block).

\end{multicols}
\end{DndMonster}



\begin{DndMonster}[float=t]{Nafik}
\DndMonsterType{Medium undead, lawful evil}
\DndMonsterBasics[
armor-class = {13 (natural armor)},
hit-points = {\DndDice{11d8 + 33}},
speed = {20 ft.}
]
\DndMonsterAbilityScores[
 str = 18,
 dex = 10,
 con = 16,
 int = 10,
 wis = 16,
 cha = 14
]
\DndMonsterDetails[
saving-throws = {Int +3, Wis +6},
skills = {History +3, Religion +3},
damage-vulnerabilities = {fire},
damage-resistances = {bludgeoning, piercing, slashing from nonmagical attacks},
damage-immunities = {necrotic, poison},
condition-immunities = {charmed, exhaustion, frightened, paralyzed, poisoned},
senses = {darkvision 60 ft., passive perception 13},
languages = {Common},
challenge = 6,
]
\DndMonsterAction{Legendary Resistance (2/Day)}
If Nafik fails a saving throw, he can choose to succeed instead.

\DndMonsterAction{Magic Resistance}
Nafik has advantage on saving throws against spells and other magical effects.

\DndMonsterAction{Rejuvenation}
When he is destroyed, Nafik gains a new body in 24 hours if his heart is intact, regaining all his hit points. The new body appears in an unoccupied space within 5 feet of Nafik's heart.

\DndMonsterAction{Spellcasting}
Nafik casts one of the following spells, using Wisdom as the spellcasting ability (spell save DC 14):
\begin{DndMonsterSpells}
  \DndInnateSpellLevel{Thaumaturgy}
  \DndInnateSpellLevel[]{Insect Plague}
  \DndInnateSpellLevel[2]{Command}
\end{DndMonsterSpells}
\DndMonsterSection{Actions}
\DndMonsterAction{Multiattack}
Nafik can use his Dreadful Glare and makes one Rotting Fist or Unholy Beam attack.

\DndMonsterAction{Rotting Fist}
\textsl{Melee Weapon Attack:} 7 to hit, reach 5 ft., one target. \textsl{Hit:} 11 (2d6 + 4) bludgeoning damage plus 21 (6d6) necrotic damage. If the target is a creature, it must succeed on a DC 14 Constitution saving throw or be cursed with mummy rot. The cursed target can't regain hit points, and its hit point maximum decreases by 10 (3d6) for every 24 hours that elapse. If the curse reduces the target's hit point maximum to 0, the target dies and its body turns to dust. The curse lasts until removed by the \textit{Remove Curse} spell or other magic.

\DndMonsterAction{Unholy Beam}
\textsl{Ranged Spell Attack:} 6 to hit, range 120 ft., one target. \textsl{Hit:} 28 (8d6) necrotic damage, and the next attack roll made against this target before the end of Nafik's next turn has advantage.

\DndMonsterAction{Dreadful Glare}
Nafik targets one creature he can see within 60 feet of himself. The target must succeed on a DC 14 Wisdom saving throw or have the frightened condition until the end of Nafik's next turn. A target that succeeds on the saving throw is immune to Nafik's Dreadful Glare for the next 24 hours.

\end{DndMonster}



\begin{DndMonster}[float=t]{Nalfeshnee}
\DndMonsterType{Large fiend (demon), chaotic evil}
\DndMonsterBasics[
armor-class = {18 (natural armor)},
hit-points = {\DndDice{16d10 + 96}},
speed = {20 ft., fly 30 ft.}
]
\DndMonsterAbilityScores[
 str = 21,
 dex = 10,
 con = 22,
 int = 19,
 wis = 12,
 cha = 15
]
\DndMonsterDetails[
saving-throws = {Con +11, Int +9, Wis +6, Cha +7},
damage-resistances = {cold; fire; lightning; bludgeoning, piercing, slashing from nonmagical attacks},
damage-immunities = {poison},
condition-immunities = {poisoned},
senses = {truesight 120 ft., passive perception 11},
languages = {Abyssal, telepathy 120 ft.},
challenge = 13,
]
\DndMonsterAction{Magic Resistance}
The nalfeshnee has advantage on saving throws against spells and other magical effects.

\DndMonsterSection{Actions}
\DndMonsterAction{Multiattack}
The nalfeshnee uses Horror Nimbus if it can. It then makes three attacks: one with its bite and two with its claws.

\DndMonsterAction{Bite}
\textsl{Melee Weapon Attack:} 10 to hit, reach 5 ft., one target. \textsl{Hit:} 32 (5d10 + 5) piercing damage.

\DndMonsterAction{Claw}
\textsl{Melee Weapon Attack:} 10 to hit, reach 10 ft., one target. \textsl{Hit:} 15 (3d6 + 5) slashing damage.

\DndMonsterAction{Horror Nimbus (Recharge 5\textendash 6)}
The nalfeshnee magically emits scintillating, multicolored light. Each creature within 15 feet of the nalfeshnee that can see the light must succeed on a DC 15 Wisdom saving throw or be frightened for 1 minute. A creature can repeat the saving throw at the end of each of its turns, ending the effect on itself on a success. If a creature's saving throw is successful or the effect ends for it, the creature is immune to the nalfeshnee's Horror Nimbus for the next 24 hours.

\DndMonsterAction{Teleport}
The nalfeshnee magically teleports, along with any equipment it is wearing or carrying, up to 120 feet to an unoccupied space it can see.

\end{DndMonster}



\begin{DndMonster}[float=t]{Noble}
\DndMonsterType{Medium humanoid (any race), any alignment}
\DndMonsterBasics[
armor-class = {15 (\textit{breastplate})},
hit-points = {\DndDice{2d8}},
speed = {30 ft.}
]
\DndMonsterAbilityScores[
 str = 11,
 dex = 12,
 con = 11,
 int = 12,
 wis = 14,
 cha = 16
]
\DndMonsterDetails[
skills = {Deception +5, Insight +4, Persuasion +5},
languages = {any two languages},
challenge = 1/8,
]
\DndMonsterSection{Actions}
\DndMonsterAction{Rapier}
\textsl{Melee Weapon Attack:} 3 to hit, reach 5 ft., one target. \textsl{Hit:} 5 (1d8 + 1) piercing damage.

\DndMonsterSection{Reactions}
\DndMonsterAction{Parry}
The noble adds 2 to its AC against one melee attack that would hit it. To do so, the noble must see the attacker and be wielding a melee weapon.

\end{DndMonster}



\begin{DndMonster}[float=t]{Ochre Jelly}
\DndMonsterType{Large ooze, unaligned}
\DndMonsterBasics[
armor-class = {8},
hit-points = {\DndDice{6d10 + 12}},
speed = {10 ft., climb 10 ft.}
]
\DndMonsterAbilityScores[
 str = 15,
 dex = 6,
 con = 14,
 int = 2,
 wis = 6,
 cha = 1
]
\DndMonsterDetails[
damage-resistances = {acid},
damage-immunities = {lightning, slashing},
condition-immunities = {blinded, charmed, deafened, exhaustion, frightened, prone},
senses = {blindsight 60 ft. (blind beyond this radius), passive perception 8},
challenge = 2,
]
\DndMonsterAction{Amorphous}
The jelly can move through a space as narrow as 1 inch wide without squeezing.

\DndMonsterAction{Spider Climb}
The jelly can climb difficult surfaces, including upside down on ceilings, without needing to make an ability check.

\DndMonsterSection{Actions}
\DndMonsterAction{Pseudopod}
\textsl{Melee Weapon Attack:} 4 to hit, reach 5 ft., one target. \textsl{Hit:} 9 (2d6 + 2) bludgeoning damage plus 3 (1d6) acid damage.

\DndMonsterSection{Reactions}
\DndMonsterAction{Split}
When a jelly that is Medium or larger is subjected to lightning or slashing damage, it splits into two new jellies if it has at least 10 hit points. Each new jelly has hit points equal to half the original jelly's, rounded down. New jellies are one size smaller than the original jelly.

\end{DndMonster}



\begin{DndMonster}[float=t]{Ogre}
\DndMonsterType{Large giant, chaotic evil}
\DndMonsterBasics[
armor-class = {11 (\textit{hide armor})},
hit-points = {\DndDice{7d10 + 21}},
speed = {40 ft.}
]
\DndMonsterAbilityScores[
 str = 19,
 dex = 8,
 con = 16,
 int = 5,
 wis = 7,
 cha = 7
]
\DndMonsterDetails[
senses = {darkvision 60 ft., passive perception 8},
languages = {Common, Giant},
challenge = 2,
]
\DndMonsterSection{Actions}
\DndMonsterAction{Greatclub}
\textsl{Melee Weapon Attack:} 6 to hit, reach 5 ft., one target. \textsl{Hit:} 13 (2d8 + 4) bludgeoning damage.

\DndMonsterAction{Javelin}
\textsl{Melee or Ranged Weapon Attack:} 6 to hit, reach 5 ft. or range 30/120 ft., one target. \textsl{Hit:} 11 (2d6 + 4) piercing damage.

\end{DndMonster}



\begin{DndMonster}[float=t]{Owlbear}
\DndMonsterType{Large monstrosity, unaligned}
\DndMonsterBasics[
armor-class = {13 (natural armor)},
hit-points = {\DndDice{7d10 + 21}},
speed = {40 ft.}
]
\DndMonsterAbilityScores[
 str = 20,
 dex = 12,
 con = 17,
 int = 3,
 wis = 12,
 cha = 7
]
\DndMonsterDetails[
skills = {Perception +3},
senses = {darkvision 60 ft., passive perception 13},
challenge = 3,
]
\DndMonsterAction{Keen Sight and Smell}
The owlbear has advantage on Wisdom (Perception) checks that rely on sight or smell.

\DndMonsterSection{Actions}
\DndMonsterAction{Multiattack}
The owlbear makes two attacks: one with its beak and one with its claws.

\DndMonsterAction{Beak}
\textsl{Melee Weapon Attack:} 7 to hit, reach 5 ft., one creature. \textsl{Hit:} 10 (1d10 + 5) piercing damage.

\DndMonsterAction{Claws}
\textsl{Melee Weapon Attack:} 7 to hit, reach 5 ft., one target. \textsl{Hit:} 14 (2d8 + 5) slashing damage.

\end{DndMonster}


\clearpage

\begin{DndMonster}[float=t]{Pech}
\DndMonsterType{Small elemental, neutral good}
\DndMonsterBasics[
armor-class = {17 (natural armor)},
hit-points = {\DndDice{11d6 + 44}},
speed = {30 ft., burrow 20 ft.}
]
\DndMonsterAbilityScores[
 str = 19,
 dex = 11,
 con = 18,
 int = 11,
 wis = 14,
 cha = 10
]
\DndMonsterDetails[
saving-throws = {Con +6, Wis +4},
skills = {Athletics +6, Perception +4, Survival +4},
condition-immunities = {petrified},
senses = {darkvision 120 ft., tremorsense 120 ft., passive perception 14},
languages = {Common, Terran},
challenge = 4,
]
\DndMonsterAction{Magic Resistance}
The pech has advantage on saving throws against spells and other magical effects.

\DndMonsterAction{Sunlight Sensitivity}
While in sunlight, the pech has disadvantage on attack rolls.

\DndMonsterAction{Communal Spellcasting (2/Day)}
The pech works with three or more pechs to cast spells, requiring no spell components and using Wisdom as the spellcasting ability (save DC 12). If at least three other pechs are within 30 feet of it, the pech can cast \textit{Wall of Stone} . If at least seven other pechs are within 30 feet of it, it can cast \textit{Greater Restoration} . Each other pech involved in casting the spell can't have the incapacitated condition and must have at least one use of Communal Spellcasting remaining, which it must immediately expend to participate (no action required).
\begin{DndMonsterSpells}
  \DndInnateSpellLevel[]{Wall of Stone}
\end{DndMonsterSpells}
\DndMonsterSection{Actions}
\DndMonsterAction{Multiattack}
The pech makes two Fortified Pickaxe attacks. If it hits a Large or smaller creature with both attacks, the target must succeed on a DC 14 Strength saving throw or have the prone condition.

\DndMonsterAction{Fortified Pickaxe}
\textsl{Melee Weapon Attack:} 6 to hit, reach 5 ft., one target. \textsl{Hit:} 11 (2d6 + 4) force damage. If the target is a Construct or an object, the attack is automatically a critical hit.

\DndMonsterAction{Stone Shape (3/Day)}
The pech casts \textit{Stone Shape}, requiring no spell components and using Wisdom as the spellcasting ability.

\end{DndMonster}



\begin{DndMonster}[float=t]{Phase Spider}
\DndMonsterType{Large monstrosity, unaligned}
\DndMonsterBasics[
armor-class = {13 (natural armor)},
hit-points = {\DndDice{5d10 + 5}},
speed = {30 ft., climb 30 ft.}
]
\DndMonsterAbilityScores[
 str = 15,
 dex = 15,
 con = 12,
 int = 6,
 wis = 10,
 cha = 6
]
\DndMonsterDetails[
skills = {Stealth +6},
senses = {darkvision 60 ft., passive perception 10},
challenge = 3,
]
\DndMonsterAction{Ethereal Jaunt}
As a bonus action, the spider can magically shift from the Material Plane to the Ethereal Plane, or vice versa.

\DndMonsterAction{Spider Climb}
The spider can climb difficult surfaces, including upside down on ceilings, without needing to make an ability check.

\DndMonsterAction{Web Walker}
The spider ignores movement restrictions caused by webbing.

\DndMonsterSection{Actions}
\DndMonsterAction{Bite}
\textsl{Melee Weapon Attack:} 4 to hit, reach 5 ft., one creature. \textsl{Hit:} 7 (1d10 + 2) piercing damage, and the target must make a DC 11 Constitution saving throw, taking 18 (4d8) poison damage on a failed save, or half as much damage on a successful one. If the poison damage reduces the target to 0 hit points, the target is stable but poisoned for 1 hour, even after regaining hit points, and is paralyzed while poisoned in this way.

\end{DndMonster}



\begin{DndMonster}[float=t]{Pixie}
\DndMonsterType{Tiny fey, neutral good}
\DndMonsterBasics[
armor-class = {15},
hit-points = {\DndDice{1d4 - 1}},
speed = {10 ft., fly 30 ft.}
]
\DndMonsterAbilityScores[
 str = 2,
 dex = 20,
 con = 8,
 int = 10,
 wis = 14,
 cha = 15
]
\DndMonsterDetails[
skills = {Perception +4, Stealth +7},
languages = {Sylvan},
challenge = 1/4,
]
\DndMonsterAction{Magic Resistance}
The pixie has advantage on saving throws against spells and other magical effects.

\DndMonsterAction{Innate Spellcasting}
The pixie's innate spellcasting ability is Charisma (spell save DC 12). It can innately cast the following spells, requiring only its pixie dust as a component:
\begin{DndMonsterSpells}
  \DndInnateSpellLevel{druidcraft}
  \DndInnateSpellLevel[1]{confusion}
\end{DndMonsterSpells}
\DndMonsterSection{Actions}
\DndMonsterAction{Superior Invisibility}
The pixie magically turns invisible until its concentration ends (as if concentration on a spell). Any equipment the pixie wears or carries is invisible with it.

\end{DndMonster}



\begin{DndMonster}[float=t]{Poisonous Snake}
\DndMonsterType{Tiny beast, unaligned}
\DndMonsterBasics[
armor-class = {13},
hit-points = {\DndDice{1d4}},
speed = {30 ft., swim 30 ft.}
]
\DndMonsterAbilityScores[
 str = 2,
 dex = 16,
 con = 11,
 int = 1,
 wis = 10,
 cha = 3
]
\DndMonsterDetails[
senses = {blindsight 10 ft., passive perception 10},
challenge = 1/8,
]
\DndMonsterSection{Actions}
\DndMonsterAction{Bite}
\textsl{Melee Weapon Attack:} 5 to hit, reach 5 ft., one target. \textsl{Hit:} 1 piercing damage, and the target must make a DC 10 Constitution saving throw, taking 5 (2d4) poison damage on a failed save, or half as much damage on a successful one.

\end{DndMonster}



\begin{DndMonster}[float=t]{Poltergeist}
\DndMonsterType{Medium undead, chaotic evil}
\DndMonsterBasics[
armor-class = {12},
hit-points = {\DndDice{5d8}},
speed = {0 ft., fly 50 ft. \textit{(hover)}}
]
\DndMonsterAbilityScores[
 str = 1,
 dex = 14,
 con = 11,
 int = 10,
 wis = 10,
 cha = 11
]
\DndMonsterDetails[
damage-resistances = {acid; cold; fire; lightning; thunder; bludgeoning, piercing, slashing from nonmagical attacks},
damage-immunities = {necrotic, poison},
condition-immunities = {charmed, exhaustion, grappled, paralyzed, petrified, poisoned, prone, restrained, unconscious},
senses = {darkvision 60 ft., passive perception 10},
languages = {understands all languages it knew in life but can't speak},
challenge = 2,
]
\DndMonsterAction{Incorporeal Movement}
The poltergeist can move through other creatures and objects as if they were difficult terrain. It takes 5 (1d10) force damage if it ends its turn inside an object.

\DndMonsterAction{Sunlight Sensitivity}
While in sunlight, the poltergeist has disadvantage on attack rolls, as well as on Wisdom (Perception) checks that rely on sight.

\DndMonsterAction{Invisibility}
The poltergeist is invisible.

\DndMonsterSection{Actions}
\DndMonsterAction{Forceful Slam}
\textsl{Melee Weapon Attack:} 4 to hit, reach 5 ft., one creature. \textsl{Hit:} 10 (3d6) force damage.

\DndMonsterAction{Telekinetic Thrust}
The poltergeist targets a creature or unattended object within 30 feet of it. A creature must be Medium or smaller to be affected by this magic, and an object can weigh up to 150 pounds.

If the target is a creature, the poltergeist makes a Charisma check contested by the target's Strength check. If the poltergeist wins the contest, the poltergeist hurls the target up to 30 feet in any direction, including upward. If the target then comes into contact with a hard surface or heavy object, the target takes 1d6 damage per 10 feet moved.

If the target is an object that isn't being worn or carried, the poltergeist hurls it up to 30 feet in any direction. The poltergeist can use the object as a ranged weapon, attacking one creature along the object's path (4 to hit) and dealing 5 (2d4) bludgeoning damage on a hit.

\end{DndMonster}



\begin{DndMonster}[float=t]{Priest}
\DndMonsterType{Medium humanoid (any race), any alignment}
\DndMonsterBasics[
armor-class = {13 (\textit{chain shirt})},
hit-points = {\DndDice{5d8 + 5}},
speed = {30 ft.}
]
\DndMonsterAbilityScores[
 str = 10,
 dex = 10,
 con = 12,
 int = 13,
 wis = 16,
 cha = 13
]
\DndMonsterDetails[
skills = {Medicine +7, Persuasion +3, Religion +5},
languages = {any two languages},
challenge = 2,
]
\DndMonsterAction{Divine Eminence}
As a bonus action, the priest can expend a spell slot to cause its melee weapon attacks to magically deal an extra 10 (3d6) radiant damage to a target on a hit. This benefit lasts until the end of the turn. If the priest expends a spell slot of 2nd level or higher, the extra damage increases by 1d6 for each level above 1st.

\DndMonsterAction{Spellcasting}
The priest is a 5th-level spellcaster. Its spellcasting ability is Wisdom (spell save DC 13, 5 to hit with spell attacks). The priest has the following cleric spells prepared:
\begin{DndMonsterSpells}
  \DndMonsterSpellLevel{light}
   \DndMonsterSpellLevel[1][4]{cure wounds}
   \DndMonsterSpellLevel[2][3]{lesser restoration}
   \DndMonsterSpellLevel[3][2]{dispel magic}
\end{DndMonsterSpells}
\DndMonsterSection{Actions}
\DndMonsterAction{Mace}
\textsl{Melee Weapon Attack:} 2 to hit, reach 5 ft., one target. \textsl{Hit:} 3 (1d6) bludgeoning damage.

\end{DndMonster}



\begin{DndMonster}[float=t]{Rat}
\DndMonsterType{Tiny beast, unaligned}
\DndMonsterBasics[
armor-class = {10},
hit-points = {\DndDice{1d4 - 1}},
speed = {20 ft.}
]
\DndMonsterAbilityScores[
 str = 2,
 dex = 11,
 con = 9,
 int = 2,
 wis = 10,
 cha = 4
]
\DndMonsterDetails[
senses = {darkvision 30 ft., passive perception 10},
challenge = 0,
]
\DndMonsterAction{Keen Smell}
The rat has advantage on Wisdom (Perception) checks that rely on smell.

\DndMonsterSection{Actions}
\DndMonsterAction{Bite}
\textsl{Melee Weapon Attack:} 0 to hit, reach 5 ft., one target. \textsl{Hit:} 1 piercing damage.

\end{DndMonster}



\begin{DndMonster}[float=t]{Rhinoceros}
\DndMonsterType{Large beast, unaligned}
\DndMonsterBasics[
armor-class = {11 (natural armor)},
hit-points = {\DndDice{6d10 + 12}},
speed = {40 ft.}
]
\DndMonsterAbilityScores[
 str = 21,
 dex = 8,
 con = 15,
 int = 2,
 wis = 12,
 cha = 6
]
\DndMonsterDetails[
challenge = 2,
]
\DndMonsterAction{Charge}
If the rhinoceros moves at least 20 feet straight toward a target and then hits it with a gore attack on the same turn, the target takes an extra 9 (2d8) bludgeoning damage. If the target is a creature, it must succeed on a DC 15 Strength saving throw or be knocked prone.

\DndMonsterSection{Actions}
\DndMonsterAction{Gore}
\textsl{Melee Weapon Attack:} 7 to hit, reach 5 ft., one target. \textsl{Hit:} 14 (2d8 + 5) bludgeoning damage.

\end{DndMonster}



\begin{DndMonster}[float=t]{Riding Horse}
\DndMonsterType{Large beast, unaligned}
\DndMonsterBasics[
armor-class = {10},
hit-points = {\DndDice{2d10 + 2}},
speed = {60 ft.}
]
\DndMonsterAbilityScores[
 str = 16,
 dex = 10,
 con = 12,
 int = 2,
 wis = 11,
 cha = 7
]
\DndMonsterDetails[
challenge = 1/4,
]
\DndMonsterSection{Actions}
\DndMonsterAction{Hooves}
\textsl{Melee Weapon Attack:} 5 to hit, reach 5 ft., one target. \textsl{Hit:} 8 (2d4 + 3) bludgeoning damage.

\end{DndMonster}



\begin{DndMonster}[float=t]{Roper}
\DndMonsterType{Large monstrosity, neutral evil}
\DndMonsterBasics[
armor-class = {20 (natural armor)},
hit-points = {\DndDice{11d10 + 33}},
speed = {10 ft., climb 10 ft.}
]
\DndMonsterAbilityScores[
 str = 18,
 dex = 8,
 con = 17,
 int = 7,
 wis = 16,
 cha = 6
]
\DndMonsterDetails[
skills = {Perception +6, Stealth +5},
senses = {darkvision 60 ft., passive perception 16},
challenge = 5,
]
\DndMonsterAction{False Appearance}
While the roper remains motionless, it is indistinguishable from a normal cave formation, such as a stalagmite.

\DndMonsterAction{Grasping Tendrils}
The roper can have up to six tendrils at a time. Each tendril can be attacked (AC 20; 10 hit points; immunity to poison and psychic damage). Destroying a tendril deals no damage to the roper, which can extrude a replacement tendril on its next turn. A tendril can also be broken if a creature takes an action and succeeds on a DC 15 Strength check against it.

\DndMonsterAction{Spider Climb}
The roper can climb difficult surfaces, including upside down on ceilings, without needing to make an ability check.

\DndMonsterSection{Actions}
\DndMonsterAction{Multiattack}
The roper makes four attacks with its tendrils, uses Reel, and makes one attack with its bite.

\DndMonsterAction{Bite}
\textsl{Melee Weapon Attack:} 7 to hit, reach 5 ft., one target. \textsl{Hit:} 22 (4d8 + 4) piercing damage.

\DndMonsterAction{Tendril}
\textsl{Melee Weapon Attack:} 7 to hit, reach 50 ft., one creature. \textsl{Hit:} The target is grappled (escape DC 15). Until the grapple ends, the target is restrained and has disadvantage on Strength checks and Strength saving throws, and the roper can't use the same tendril on another target.

\DndMonsterAction{Reel}
The roper pulls each creature grappled by it up to 25 feet straight toward it.

\end{DndMonster}



\begin{DndMonster}[float=t]{Rug of Smothering}
\DndMonsterType{Large construct, unaligned}
\DndMonsterBasics[
armor-class = {12},
hit-points = {\DndDice{6d10}},
speed = {10 ft.}
]
\DndMonsterAbilityScores[
 str = 17,
 dex = 14,
 con = 10,
 int = 1,
 wis = 3,
 cha = 1
]
\DndMonsterDetails[
damage-immunities = {poison, psychic},
condition-immunities = {blinded, charmed, deafened, frightened, paralyzed, petrified, poisoned},
senses = {blindsight 60 ft. (blind beyond this radius), passive perception 6},
challenge = 2,
]
\DndMonsterAction{Antimagic Susceptibility}
The rug is incapacitated while in the area of an \textit{antimagic field}. If targeted by \textit{dispel magic}, the rug must succeed on a Constitution saving throw against the caster's spell save DC or fall unconscious for 1 minute.

\DndMonsterAction{Damage Transfer}
While it is grappling a creature, the rug takes only half the damage dealt to it, and the creature grappled by the rug takes the other half.

\DndMonsterAction{False Appearance}
While the rug remains motionless, it is indistinguishable from a normal rug.

\DndMonsterSection{Actions}
\DndMonsterAction{Smother}
\textsl{Melee Weapon Attack:} 5 to hit, reach 5 ft., one Medium or smaller creature. \textsl{Hit:} The creature is grappled (escape DC 13). Until this grapple ends, the target is restrained, blinded, and at risk of suffocating, and the rug can't smother another target. In addition, at the start of each of the target's turns, the target takes 10 (2d6 + 3) bludgeoning damage.

\end{DndMonster}



\begin{DndMonster}[float=t]{Rust Monster}
\DndMonsterType{Medium monstrosity, unaligned}
\DndMonsterBasics[
armor-class = {14 (natural armor)},
hit-points = {\DndDice{5d8 + 5}},
speed = {40 ft.}
]
\DndMonsterAbilityScores[
 str = 13,
 dex = 12,
 con = 13,
 int = 2,
 wis = 13,
 cha = 6
]
\DndMonsterDetails[
senses = {darkvision 60 ft., passive perception 11},
challenge = 1/2,
]
\DndMonsterAction{Iron Scent}
The rust monster can pinpoint, by scent, the location of ferrous metal within 30 feet of it.

\DndMonsterAction{Rust Metal}
Any nonmagical weapon made of metal that hits the rust monster corrodes. After dealing damage, the weapon takes a permanent and cumulative −1 penalty to damage rolls. If its penalty drops to −5, the weapon is destroyed. Non magical ammunition made of metal that hits the rust monster is destroyed after dealing damage.

\DndMonsterSection{Actions}
\DndMonsterAction{Bite}
\textsl{Melee Weapon Attack:} 3 to hit, reach 5 ft., one target. \textsl{Hit:} 5 (1d8 + 1) piercing damage.

\DndMonsterAction{Antennae}
The rust monster corrodes a nonmagical ferrous metal object it can see within 5 feet of it. If the object isn't being worn or carried, the touch destroys a 1-foot cube of it. If the object is being worn or carried by a creature, the creature can make a DC 11 Dexterity saving throw to avoid the rust monster's touch.

If the object touched is either metal armor or a metal shield being worn or carried, it takes a permanent and cumulative −1 penalty to the AC it offers. Armor reduced to an AC of 10 or a shield that drops to a +0 bonus is destroyed. If the object touched is a held metal weapon, it rusts as described in the Rust Metal trait.

\end{DndMonster}



\begin{DndMonster}[float=t]{Saber-Toothed Tiger}
\DndMonsterType{Large beast, unaligned}
\DndMonsterBasics[
armor-class = {12},
hit-points = {\DndDice{7d10 + 14}},
speed = {40 ft.}
]
\DndMonsterAbilityScores[
 str = 18,
 dex = 14,
 con = 15,
 int = 3,
 wis = 12,
 cha = 8
]
\DndMonsterDetails[
skills = {Perception +3, Stealth +6},
challenge = 2,
]
\DndMonsterAction{Keen Smell}
The tiger has advantage on Wisdom (Perception) checks that rely on smell.

\DndMonsterAction{Pounce}
If the tiger moves at least 20 feet straight toward a creature and then hits it with a claw attack on the same turn, that target must succeed on a DC 14 Strength saving throw or be knocked prone. If the target is prone, the tiger can make one bite attack against it as a bonus action.

\DndMonsterSection{Actions}
\DndMonsterAction{Bite}
\textsl{Melee Weapon Attack:} 6 to hit, reach 5 ft., one target. \textsl{Hit:} 10 (1d10 + 5) piercing damage.

\DndMonsterAction{Claw}
\textsl{Melee Weapon Attack:} 6 to hit, reach 5 ft., one target. \textsl{Hit:} 12 (2d6 + 5) slashing damage.

\end{DndMonster}



\begin{DndMonster}[float=t]{Satyr}
\DndMonsterType{Medium fey, chaotic neutral}
\DndMonsterBasics[
armor-class = {14 (\textit{leather armor})},
hit-points = {\DndDice{7d8}},
speed = {40 ft.}
]
\DndMonsterAbilityScores[
 str = 12,
 dex = 16,
 con = 11,
 int = 12,
 wis = 10,
 cha = 14
]
\DndMonsterDetails[
skills = {Perception +2, Performance +6, Stealth +5},
languages = {Common, Elvish, Sylvan},
challenge = 1/2,
]
\DndMonsterAction{Magic Resistance}
The satyr has advantage on saving throws against spells and other magical effects.

\DndMonsterSection{Actions}
\DndMonsterAction{Ram}
\textsl{Melee Weapon Attack:} 3 to hit, reach 5 ft., one target. \textsl{Hit:} 6 (2d4 + 1) bludgeoning damage.

\DndMonsterAction{Shortsword}
\textsl{Melee Weapon Attack:} 5 to hit, reach 5 ft., one target. \textsl{Hit:} 6 (1d6 + 3) piercing damage.

\DndMonsterAction{Shortbow}
\textsl{Ranged Weapon Attack:} 5 to hit, range 80/320 ft., one target. \textsl{Hit:} 6 (1d6 + 3) piercing damage.

\end{DndMonster}



\begin{DndMonster}[float=t]{Scout}
\DndMonsterType{Medium humanoid (any race), any alignment}
\DndMonsterBasics[
armor-class = {13 (\textit{leather armor})},
hit-points = {\DndDice{3d8 + 3}},
speed = {30 ft.}
]
\DndMonsterAbilityScores[
 str = 11,
 dex = 14,
 con = 12,
 int = 11,
 wis = 13,
 cha = 11
]
\DndMonsterDetails[
skills = {Nature +4, Perception +5, Stealth +6, Survival +5},
languages = {any one language (usually Common)},
challenge = 1/2,
]
\DndMonsterAction{Keen Hearing and Sight}
The scout has advantage on Wisdom (Perception) checks that rely on hearing or sight.

\DndMonsterSection{Actions}
\DndMonsterAction{Multiattack}
The scout makes two melee attacks or two ranged attacks.

\DndMonsterAction{Shortsword}
\textsl{Melee Weapon Attack:} 4 to hit, reach 5 ft., one target. \textsl{Hit:} 5 (1d6 + 2) piercing damage.

\DndMonsterAction{Longbow}
\textsl{Ranged Weapon Attack:} 4 to hit, ranged 150/600 ft., one target. \textsl{Hit:} 6 (1d8 + 2) piercing damage.

\end{DndMonster}



\begin{DndMonster}[float=t]{Shadow}
\DndMonsterType{Medium undead, chaotic evil}
\DndMonsterBasics[
armor-class = {12},
hit-points = {\DndDice{3d8 + 3}},
speed = {40 ft.}
]
\DndMonsterAbilityScores[
 str = 6,
 dex = 14,
 con = 13,
 int = 6,
 wis = 10,
 cha = 8
]
\DndMonsterDetails[
skills = {Stealth +4},
damage-vulnerabilities = {radiant},
damage-resistances = {acid; cold; fire; lightning; thunder; bludgeoning, piercing, slashing from nonmagical attacks},
damage-immunities = {necrotic, poison},
condition-immunities = {exhaustion, frightened, grappled, paralyzed, petrified, poisoned, prone, restrained},
senses = {darkvision 60 ft., passive perception 10},
challenge = 1/2,
]
\DndMonsterAction{Amorphous}
The shadow can move through a space as narrow as 1 inch wide without squeezing.

\DndMonsterAction{Shadow Stealth}
While in dim light or darkness, the shadow can take the Hide action as a bonus action. Its stealth bonus is also improved to +6.

\DndMonsterAction{Sunlight Weakness}
While in sunlight, the shadow has disadvantage on attack rolls, ability checks, and saving throws.

\DndMonsterSection{Actions}
\DndMonsterAction{Strength Drain}
\textsl{Melee Weapon Attack:} 4 to hit, reach 5 ft., one creature. \textsl{Hit:} 9 (2d6 + 2) necrotic damage, and the target's Strength score is reduced by 1d4. The target dies if this reduces its Strength to 0. Otherwise, the reduction lasts until the target finishes a short or long rest.

If a non-evil humanoid dies from this attack, a new shadow rises from the corpse 1d4 hours later.

\end{DndMonster}



\begin{DndMonster}[float=t]{Shambling Mound}
\DndMonsterType{Large plant, unaligned}
\DndMonsterBasics[
armor-class = {15 (natural armor)},
hit-points = {\DndDice{16d10 + 48}},
speed = {20 ft., swim 20 ft.}
]
\DndMonsterAbilityScores[
 str = 18,
 dex = 8,
 con = 16,
 int = 5,
 wis = 10,
 cha = 5
]
\DndMonsterDetails[
skills = {Stealth +2},
damage-resistances = {cold, fire},
damage-immunities = {lightning},
condition-immunities = {blinded, deafened, exhaustion},
senses = {blindsight 60 ft. (blind beyond this radius), passive perception 10},
challenge = 5,
]
\DndMonsterAction{Lightning Absorption}
Whenever the shambling mound is subjected to lightning damage, it takes no damage and regains a number of hit points equal to the lightning damage dealt.

\DndMonsterSection{Actions}
\DndMonsterAction{Multiattack}
The shambling mound makes two slam attacks. If both attacks hit a Medium or smaller target, the target is grappled (escape DC 14), and the shambling mound uses its Engulf on it.

\DndMonsterAction{Slam}
\textsl{Melee Weapon Attack:} 7 to hit, reach 5 ft., one target. \textsl{Hit:} 13 (2d8 + 4) bludgeoning damage.

\DndMonsterAction{Engulf}
The shambling mound engulfs a Medium or smaller creature grappled by it. The engulfed target is blinded, restrained, and unable to breathe, and it must succeed on a DC 14 Constitution saving throw at the start of each of the mound's turns or take 13 (2d8 + 4) bludgeoning damage. If the mound moves, the engulfed target moves with it. The mound can have only one creature engulfed at a time.

\end{DndMonster}



\begin{DndMonster}[float=t]{Shrieker}
\DndMonsterType{Medium plant, unaligned}
\DndMonsterBasics[
armor-class = {5},
hit-points = {\DndDice{3d8}},
speed = {0 ft.}
]
\DndMonsterAbilityScores[
 str = 1,
 dex = 1,
 con = 10,
 int = 1,
 wis = 3,
 cha = 1
]
\DndMonsterDetails[
condition-immunities = {blinded, deafened, frightened},
senses = {blindsight 30 ft. (blind beyond this radius), passive perception 6},
challenge = 0,
]
\DndMonsterAction{False Appearance}
While the shrieker remains motionless, it is indistinguishable from an ordinary fungus.

\DndMonsterSection{Reactions}
\DndMonsterAction{Shriek}
When bright light or a creature is within 30 feet of the shrieker, it emits a shriek audible within 300 feet of it. The shrieker continues to shriek until the disturbance moves out of range and for 1d4 of the shrieker's turns afterward.

\end{DndMonster}



\begin{DndMonster}[float=t]{Sion}
\DndMonsterType{Medium humanoid (human, sorcerer), neutral evil}
\DndMonsterBasics[
armor-class = {12},
hit-points = {\DndDice{9d8 + 9}},
speed = {30 ft.}
]
\DndMonsterAbilityScores[
 str = 16,
 dex = 14,
 con = 13,
 int = 12,
 wis = 14,
 cha = 16
]
\DndMonsterDetails[
saving-throws = {Con +3, Cha +5},
skills = {Arcana +3, Perception +4, Stealth +6},
senses = {darkvision 120 ft., passive perception 14},
languages = {Common},
challenge = 2,
]
\DndMonsterAction{Magic Resistance}
Sion has advantage on saving throws against spells and other magical effects if he is in dim light or darkness.

\DndMonsterAction{Shadesight}
Magical darkness doesn't impede Sion's darkvision.

\DndMonsterAction{Shadowy Demise}
If Sion dies, his body melts into shadow, leaving behind only equipment he was wearing or carrying.

\DndMonsterAction{Sunlight Weakness}
While in sunlight, Sion has disadvantage on attack rolls, ability checks, and saving throws.

\DndMonsterAction{Spellcasting}
Sion casts one of the following spells, using Charisma as the spellcasting ability (spell save DC 13):
\begin{DndMonsterSpells}
  \DndInnateSpellLevel{Mage Hand}
  \DndInnateSpellLevel[1]{Darkness}
\end{DndMonsterSpells}
\DndMonsterSection{Actions}
\DndMonsterAction{Multiattack}
Sion makes one Shadow Sword attack and uses Affix Shadow or Spellcasting.

\DndMonsterAction{Shadow Sword}
\textsl{Melee Weapon Attack:} 5 to hit, reach 5 ft., one target. \textsl{Hit:} 14 (2d10 + 3) necrotic damage.

\DndMonsterAction{Affix Shadow}
Sion targets a creature within 60 feet of himself that he can see. That target must make a DC 13 Charisma saving throw. On a failed save, the target has the grappled condition (escape DC 13) as its shadow wraps around it. Once the target escapes the grapple, its shadow returns to normal.

\DndMonsterSection{Bonus Actions}
\DndMonsterAction{Shadow Step}
While Sion is in dim light or darkness, he magically teleports up to 15 feet to an unoccupied space he can see that is also in dim light or darkness.

\end{DndMonster}



\begin{DndMonster}[float=t]{Skeleton}
\DndMonsterType{Medium undead, lawful evil}
\DndMonsterBasics[
armor-class = {13 (armor scraps)},
hit-points = {\DndDice{2d8 + 4}},
speed = {30 ft.}
]
\DndMonsterAbilityScores[
 str = 10,
 dex = 14,
 con = 15,
 int = 6,
 wis = 8,
 cha = 5
]
\DndMonsterDetails[
damage-vulnerabilities = {bludgeoning},
damage-immunities = {poison},
condition-immunities = {exhaustion, poisoned},
senses = {darkvision 60 ft., passive perception 9},
languages = {understands all languages it spoke in life but can't speak},
challenge = 1/4,
]
\DndMonsterSection{Actions}
\DndMonsterAction{Shortsword}
\textsl{Melee Weapon Attack:} 4 to hit, reach 5 ft., one target. \textsl{Hit:} 5 (1d6 + 2) piercing damage.

\DndMonsterAction{Shortbow}
\textsl{Ranged Weapon Attack:} 4 to hit, range 80/320 ft., one target. \textsl{Hit:} 5 (1d6 + 2) piercing damage.

\end{DndMonster}


\clearpage

\begin{DndMonster}[float=t]{Spectator}
\DndMonsterType{Medium aberration, lawful neutral}
\DndMonsterBasics[
armor-class = {14 (natural armor)},
hit-points = {\DndDice{6d8 + 12}},
speed = {0 ft., fly 30 ft. \textit{(hover)}}
]
\DndMonsterAbilityScores[
 str = 8,
 dex = 14,
 con = 14,
 int = 13,
 wis = 14,
 cha = 11
]
\DndMonsterDetails[
skills = {Perception +6},
condition-immunities = {prone},
senses = {darkvision 120 ft., passive perception 16},
languages = {Deep Speech, Undercommon, telepathy 120 ft.},
challenge = 3,
]
\DndMonsterSection{Actions}
\DndMonsterAction{Bite}
\textsl{Melee Weapon Attack:} 1 to hit, reach 5 ft., one target. \textsl{Hit:} 2 (1d6 - 1) piercing damage.

\DndMonsterAction{Eye Rays}
The spectator shoots up to two of the following magical eye rays at one or two creatures it can see within 90 feet of it. It can use each ray only once on a turn.


\begin{description}
\item[1. Confusion Ray.]
The target must succeed on a DC 13 Wisdom saving throw, or it can't take reactions until the end of its next turn. On its turn, the target can't move, and it uses its action to make a melee or ranged attack against a randomly determined creature within range. If the target can't attack, it does nothing on its turn.
\item[2. Paralyzing Ray.]
The target must succeed on a DC 13 Constitution saving throw or be paralyzed for 1 minute. The target can repeat the saving throw at the end of each of its turns, ending the effect on itself on a success.
\item[3. Fear Ray.]
The target must succeed on a DC 13 Wisdom saving throw or be frightened for 1 minute. The target can repeat the saving throw at the end of each of its turns, with disadvantage if the spectator is visible to the target, ending the effect on itself on a success.
\item[4. Wounding Ray.]
The target must make a DC 13 Constitution saving throw, taking 16 (3d10) necrotic damage on a failed save, or half as much damage on a successful one.
\end{description}

\DndMonsterAction{Create Food and Water}
The spectator magically creates enough food and water to sustain itself for 24 hours.

\DndMonsterSection{Reactions}
\DndMonsterAction{Spell Reflection}
If the spectator makes a successful saving throw against a spell, or a spell attack misses it, the spectator can choose another creature (including the spellcaster) it can see within 30 feet of it. The spell targets the chosen creature instead of the spectator. If the spell forced a saving throw, the chosen creature makes its own save. If the spell was an attack, the attack roll is rerolled against the chosen creature.

\end{DndMonster}



\begin{DndMonster}[float=t]{Specter}
\DndMonsterType{Medium undead, chaotic evil}
\DndMonsterBasics[
armor-class = {12},
hit-points = {\DndDice{5d8}},
speed = {0 ft., fly 50 ft. \textit{(hover)}}
]
\DndMonsterAbilityScores[
 str = 1,
 dex = 14,
 con = 11,
 int = 10,
 wis = 10,
 cha = 11
]
\DndMonsterDetails[
damage-resistances = {acid; cold; fire; lightning; thunder; bludgeoning, piercing, slashing from nonmagical attacks},
damage-immunities = {necrotic, poison},
condition-immunities = {charmed, exhaustion, grappled, paralyzed, petrified, poisoned, prone, restrained, unconscious},
senses = {darkvision 60 ft., passive perception 10},
languages = {understands all languages it knew in life but can't speak},
challenge = 1,
]
\DndMonsterAction{Incorporeal Movement}
The specter can move through other creatures and objects as if they were difficult terrain. It takes 5 (1d10) force damage if it ends its turn inside an object.

\DndMonsterAction{Sunlight Sensitivity}
While in sunlight, the specter has disadvantage on attack rolls, as well as on Wisdom (Perception) checks that rely on sight.

\DndMonsterSection{Actions}
\DndMonsterAction{Life Drain}
\textsl{Melee Spell Attack:} 4 to hit, reach 5 ft., one creature. \textsl{Hit:} 10 (3d6) necrotic damage. The target must succeed on a DC 10 Constitution saving throw or its hit point maximum is reduced by an amount equal to the damage taken. This reduction lasts until the creature finishes a long rest. The target dies if this effect reduces its hit point maximum to 0.

\end{DndMonster}



\begin{DndMonster}[float=t]{Sprite}
\DndMonsterType{Tiny fey, neutral good}
\DndMonsterBasics[
armor-class = {15 (\textit{leather armor})},
hit-points = {\DndDice{1d4}},
speed = {10 ft., fly 40 ft.}
]
\DndMonsterAbilityScores[
 str = 3,
 dex = 18,
 con = 10,
 int = 14,
 wis = 13,
 cha = 11
]
\DndMonsterDetails[
skills = {Perception +3, Stealth +8},
languages = {Common, Elvish, Sylvan},
challenge = 1/4,
]
\DndMonsterSection{Actions}
\DndMonsterAction{Longsword}
\textsl{Melee Weapon Attack:} 2 to hit, reach 5 ft., one target. \textsl{Hit:} 1 slashing damage.

\DndMonsterAction{Shortbow}
\textsl{Ranged Weapon Attack:} 6 to hit, range 40/160 ft., one target. \textsl{Hit:} 1 piercing damage, and the target must succeed on a DC 10 Constitution saving throw or become poisoned for 1 minute. If its saving throw result is 5 or lower, the poisoned target falls unconscious for the same duration, or until it takes damage or another creature takes an action to shake it awake.

\DndMonsterAction{Heart Sight}
The sprite touches a creature and magically knows the creature's current emotional state. If the target fails a DC 10 Charisma saving throw, the sprite also knows the creature's alignment. Celestials, fiends, and undead automatically fail the saving throw.

\DndMonsterAction{Invisibility}
The sprite magically turns invisible until it attacks or casts a spell, or until its concentration ends (as if concentration on a spell). Any equipment the sprite wears or carries is invisible with it.

\end{DndMonster}



\begin{DndMonster}[float=t]{Spy}
\DndMonsterType{Medium humanoid (any race), any alignment}
\DndMonsterBasics[
armor-class = {12},
hit-points = {\DndDice{6d8}},
speed = {30 ft.}
]
\DndMonsterAbilityScores[
 str = 10,
 dex = 15,
 con = 10,
 int = 12,
 wis = 14,
 cha = 16
]
\DndMonsterDetails[
skills = {Deception +5, Insight +4, Investigation +5, Perception +6, Persuasion +5, Sleight Of Hand +4, Stealth +4},
languages = {any two languages},
challenge = 1,
]
\DndMonsterAction{Cunning Action}
On each of its turns, the spy can use a bonus action to take the Dash, Disengage, or Hide action.

\DndMonsterAction{Sneak Attack (1/Turn)}
The spy deals an extra 7 (2d6) damage when it hits a target with a weapon attack and has advantage on the attack roll, or when the target is within 5 feet of an ally of the spy that isn't incapacitated and the spy doesn't have disadvantage on the attack roll.

\DndMonsterSection{Actions}
\DndMonsterAction{Multiattack}
The spy makes two melee attacks.

\DndMonsterAction{Shortsword}
\textsl{Melee Weapon Attack:} 4 to hit, reach 5 ft., one target. \textsl{Hit:} 5 (1d6 + 2) piercing damage.

\DndMonsterAction{Hand Crossbow}
\textsl{Ranged Weapon Attack:} 4 to hit, range 30/120 ft., one target. \textsl{Hit:} 5 (1d6 + 2) piercing damage.

\end{DndMonster}



\begin{DndMonster}[float=t]{Stirge}
\DndMonsterType{Tiny beast, unaligned}
\DndMonsterBasics[
armor-class = {14 (natural armor)},
hit-points = {\DndDice{1d4}},
speed = {10 ft., fly 40 ft.}
]
\DndMonsterAbilityScores[
 str = 4,
 dex = 16,
 con = 11,
 int = 2,
 wis = 8,
 cha = 6
]
\DndMonsterDetails[
senses = {darkvision 60 ft., passive perception 9},
challenge = 1/8,
]
\DndMonsterSection{Actions}
\DndMonsterAction{Blood Drain}
\textsl{Melee Weapon Attack:} 5 to hit, reach 5 ft., one creature. \textsl{Hit:} 5 (1d4 + 3) piercing damage, and the stirge attaches to the target. While attached, the stirge doesn't attack. Instead, at the start of each of the stirge's turns, the target loses 5 (1d4 + 3) hit points due to blood loss.

The stirge can detach itself by spending 5 feet of its movement. It does so after it drains 10 hit points of blood from the target or the target dies. A creature, including the target, can use its action to detach the stirge.

\end{DndMonster}



\begin{DndMonster}[float=t]{Stone Giant}
\DndMonsterType{Huge giant, neutral}
\DndMonsterBasics[
armor-class = {17 (natural armor)},
hit-points = {\DndDice{11d12 + 55}},
speed = {40 ft.}
]
\DndMonsterAbilityScores[
 str = 23,
 dex = 15,
 con = 20,
 int = 10,
 wis = 12,
 cha = 9
]
\DndMonsterDetails[
saving-throws = {Dex +5, Con +8, Wis +4},
skills = {Athletics +12, Perception +4},
senses = {darkvision 60 ft., passive perception 14},
languages = {Giant},
challenge = 7,
]
\DndMonsterAction{Stone Camouflage}
The giant has advantage on Dexterity (Stealth) checks made to hide in rocky terrain.

\DndMonsterSection{Actions}
\DndMonsterAction{Multiattack}
The giant makes two greatclub attacks.

\DndMonsterAction{Greatclub}
\textsl{Melee Weapon Attack:} 9 to hit, reach 15 ft., one target. \textsl{Hit:} 19 (3d8 + 6) bludgeoning damage.

\DndMonsterAction{Rock}
\textsl{Ranged Weapon Attack:} 9 to hit, range 60/240 ft., one target. \textsl{Hit:} 28 (4d10 + 6) bludgeoning damage. If the target is a creature, it must succeed on a DC 17 Strength saving throw or be knocked prone.

\DndMonsterSection{Reactions}
\DndMonsterAction{Rock Catching}
If a rock or similar object is hurled at the giant, the giant can, with a successful DC 10 Dexterity saving throw, catch the missile and take no bludgeoning damage from it.

\end{DndMonster}



\begin{DndMonster}[float=t]{Stone Golem}
\DndMonsterType{Large construct, unaligned}
\DndMonsterBasics[
armor-class = {17 (natural armor)},
hit-points = {\DndDice{17d10 + 85}},
speed = {30 ft.}
]
\DndMonsterAbilityScores[
 str = 22,
 dex = 9,
 con = 20,
 int = 3,
 wis = 11,
 cha = 1
]
\DndMonsterDetails[
damage-immunities = {poison; psychic; bludgeoning, piercing, slashing from nonmagical attacks that aren't adamantine},
condition-immunities = {charmed, exhaustion, frightened, paralyzed, petrified, poisoned},
senses = {darkvision 120 ft., passive perception 10},
languages = {understands the languages of its creator but can't speak},
challenge = 10,
]
\DndMonsterAction{Immutable Form}
The golem is immune to any spell or effect that would alter its form.

\DndMonsterAction{Magic Resistance}
The golem has advantage on saving throws against spells and other magical effects.

\DndMonsterAction{Magic Weapons}
The golem's weapon attacks are magical.

\DndMonsterSection{Actions}
\DndMonsterAction{Multiattack}
The golem makes two slam attacks.

\DndMonsterAction{Slam}
\textsl{Melee Weapon Attack:} 10 to hit, reach 5 ft., one target. \textsl{Hit:} 19 (3d8 + 6) bludgeoning damage.

\DndMonsterAction{Slow (Recharge 5\textendash 6)}
The golem targets one or more creatures it can see within 10 feet of it. Each target must make a DC 17 Wisdom saving throw against this magic. On a failed save, a target can't use reactions, its speed is halved, and it can't make more than one attack on its turn. In addition, the target can take either an action or a bonus action on its turn, not both. These effects last for 1 minute. A target can repeat the saving throw at the end of each of its turns, ending the effect on itself on a success.

\end{DndMonster}



\begin{DndMonster}[float=t]{Swarm of Bats}
\DndMonsterType{Medium beast, unaligned}
\DndMonsterBasics[
armor-class = {12},
hit-points = {\DndDice{5d8}},
speed = {0 ft., fly 30 ft.}
]
\DndMonsterAbilityScores[
 str = 5,
 dex = 15,
 con = 10,
 int = 2,
 wis = 12,
 cha = 4
]
\DndMonsterDetails[
damage-resistances = {bludgeoning, piercing, slashing},
condition-immunities = {charmed, frightened, grappled, paralyzed, petrified, prone, restrained, stunned},
senses = {blindsight 60 ft., passive perception 11},
challenge = 1/4,
]
\DndMonsterAction{Echolocation}
The swarm can't use its blindsight while deafened.

\DndMonsterAction{Keen Hearing}
The swarm has advantage on Wisdom (Perception) checks that rely on hearing.

\DndMonsterAction{Swarm}
The swarm can occupy another creature's space and vice versa, and the swarm can move through any opening large enough for a Tiny bat. The swarm can't regain hit points or gain temporary hit points.

\DndMonsterSection{Actions}
\DndMonsterAction{Bites}
\textsl{Melee Weapon Attack:} 4 to hit, reach 0 ft., one creature in the swarm's space. \textsl{Hit:} 5 (2d4) piercing damage, or 2 (1d4) piercing damage if the swarm has half of its hit points or fewer.

\end{DndMonster}



\begin{DndMonster}[float=t]{Swarm of Gibberlings}
\DndMonsterType{Large fiend, chaotic evil}
\DndMonsterBasics[
armor-class = {12},
hit-points = {\DndDice{7d10}},
speed = {30 ft.}
]
\DndMonsterAbilityScores[
 str = 18,
 dex = 14,
 con = 11,
 int = 5,
 wis = 7,
 cha = 5
]
\DndMonsterDetails[
damage-resistances = {bludgeoning, piercing, slashing},
condition-immunities = {charmed, frightened, grappled, paralyzed, petrified, prone, restrained, stunned},
senses = {darkvision 120 ft., passive perception 8},
languages = {Gibberling},
challenge = 3,
]
\DndMonsterAction{Aversion to Fire}
If the swarm takes fire damage, it has disadvantage on attack rolls and ability checks until the end of its next turn.

\DndMonsterAction{Incessant Gibberish}
Any non-gibberling that is within 60 feet of the swarm and doesn't have the deafened condition has disadvantage on Constitution saving throws to maintain concentration on spells and similar effects.

\DndMonsterAction{Swarm}
The swarm can occupy another creature's space and vice versa, and the swarm can move through any opening large enough to accommodate a Small gibberling. The swarm can't regain hit points or gain temporary hit points.

\DndMonsterSection{Actions}
\DndMonsterAction{Multiattack}
The swarm makes two Gnashing Teeth attacks.

\DndMonsterAction{Gnashing Teeth}
\textsl{Melee Weapon Attack:} 6 to hit, reach 0 ft., one target in the swarm's space. \textsl{Hit:} 14 (4d4 + 4) piercing damage, or 9 (2d4 + 4) piercing damage if the swarm has half of its hit points or fewer.

\end{DndMonster}



\begin{DndMonster}[float=t]{Swarm of Poisonous Snakes}
\DndMonsterType{Medium beast, unaligned}
\DndMonsterBasics[
armor-class = {14},
hit-points = {\DndDice{8d8}},
speed = {30 ft., swim 30 ft.}
]
\DndMonsterAbilityScores[
 str = 8,
 dex = 18,
 con = 11,
 int = 1,
 wis = 10,
 cha = 3
]
\DndMonsterDetails[
damage-resistances = {bludgeoning, piercing, slashing},
condition-immunities = {charmed, frightened, grappled, paralyzed, petrified, prone, restrained, stunned},
senses = {blindsight 10 ft., passive perception 10},
challenge = 2,
]
\DndMonsterAction{Swarm}
The swarm can occupy another creature's space and vice versa, and the swarm can move through any opening large enough for a Tiny snake. The swarm can't regain hit points or gain temporary hit points.

\DndMonsterSection{Actions}
\DndMonsterAction{Bites}
\textsl{Melee Weapon Attack:} 6 to hit, reach 0 ft., one creature in the swarm's space. \textsl{Hit:} 7 (2d6) piercing damage, or 3 (1d6) piercing damage if the swarm has half of its hit points or fewer. The target must make a DC 10 Constitution saving throw, taking 14 (4d6) poison damage on a failed save, or half as much damage on a successful one.

\end{DndMonster}



\begin{DndMonster}[float=t]{Swarm of Quippers}
\DndMonsterType{Medium beast, unaligned}
\DndMonsterBasics[
armor-class = {13},
hit-points = {\DndDice{8d8 - 8}},
speed = {0 ft., swim 40 ft.}
]
\DndMonsterAbilityScores[
 str = 13,
 dex = 16,
 con = 9,
 int = 1,
 wis = 7,
 cha = 2
]
\DndMonsterDetails[
damage-resistances = {bludgeoning, piercing, slashing},
condition-immunities = {charmed, frightened, grappled, paralyzed, petrified, prone, restrained, stunned},
senses = {darkvision 60 ft., passive perception 8},
challenge = 1,
]
\DndMonsterAction{Blood Frenzy}
The swarm has advantage on melee attack rolls against any creature that doesn't have all its hit points.

\DndMonsterAction{Swarm}
The swarm can occupy another creature's space and vice versa, and the swarm can move through any opening large enough for a Tiny quipper. The swarm can't regain hit points or gain temporary hit points.

\DndMonsterAction{Water Breathing}
The swarm can breathe only underwater.

\DndMonsterSection{Actions}
\DndMonsterAction{Bites}
\textsl{Melee Weapon Attack:} 5 to hit, reach 0 ft., one creature in the swarm's space. \textsl{Hit:} 14 (4d6) piercing damage, or 7 (2d6) piercing damage if the swarm has half of its hit points or fewer.

\end{DndMonster}



\begin{DndMonster}[float=t]{Swarm of Spiders}
\DndMonsterType{Medium beast, unaligned}
\DndMonsterBasics[
armor-class = {12 (natural armor)},
hit-points = {\DndDice{5d8}},
speed = {20 ft., climb 20 ft.}
]
\DndMonsterAbilityScores[
 str = 3,
 dex = 13,
 con = 10,
 int = 1,
 wis = 7,
 cha = 1
]
\DndMonsterDetails[
damage-resistances = {bludgeoning, piercing, slashing},
condition-immunities = {charmed, frightened, grappled, paralyzed, petrified, prone, restrained, stunned},
senses = {blindsight 10 ft., passive perception 8},
challenge = 1/2,
]
\DndMonsterAction{Swarm}
The swarm can occupy another creature's space and vice versa, and the swarm can move through any opening large enough for a Tiny insect. The swarm can't regain hit points or gain temporary hit points.

\DndMonsterAction{Spider Climb}
The swarm can climb difficult surfaces, including upside down on ceilings, without needing to make an ability check.

\DndMonsterAction{Web Sense}
While in contact with a web, the swarm knows the exact location of any other creature in contact with the same web.

\DndMonsterAction{Web Walker}
The swarm ignores movement restrictions caused by webbing.

\DndMonsterSection{Actions}
\DndMonsterAction{Bites}
\textsl{Melee Weapon Attack:} 3 to hit, reach 0 ft., one target in the swarm's space. \textsl{Hit:} 10 (4d4) piercing damage, or 5 (2d4) piercing damage if the swarm has half of its hit points or fewer.

\end{DndMonster}



\begin{DndMonster}[float=t]{Swarm of Wasps}
\DndMonsterType{Medium beast, unaligned}
\DndMonsterBasics[
armor-class = {12 (natural armor)},
hit-points = {\DndDice{5d8}},
speed = {5 ft., fly 30 ft.}
]
\DndMonsterAbilityScores[
 str = 3,
 dex = 13,
 con = 10,
 int = 1,
 wis = 7,
 cha = 1
]
\DndMonsterDetails[
damage-resistances = {bludgeoning, piercing, slashing},
condition-immunities = {charmed, frightened, grappled, paralyzed, petrified, prone, restrained, stunned},
senses = {blindsight 10 ft., passive perception 8},
challenge = 1/2,
]
\DndMonsterAction{Swarm}
The swarm can occupy another creature's space and vice versa, and the swarm can move through any opening large enough for a Tiny insect. The swarm can't regain hit points or gain temporary hit points.

\DndMonsterSection{Actions}
\DndMonsterAction{Bites}
\textsl{Melee Weapon Attack:} 3 to hit, reach 0 ft., one target in the swarm's space. \textsl{Hit:} 10 (4d4) piercing damage, or 5 (2d4) piercing damage if the swarm has half of its hit points or fewer.

\end{DndMonster}



\begin{DndMonster}[float*=b,width=\textwidth + 8pt]{The Gardener}
\begin{multicols}{2}
\DndMonsterType{Medium fey (archfey, druid), neutral good}
\DndMonsterBasics[
armor-class = {17 (natural armor)},
hit-points = {\DndDice{22d8 + 110}},
speed = {30 ft.}
]
\DndMonsterAbilityScores[
 str = 21,
 dex = 16,
 con = 20,
 int = 17,
 wis = 20,
 cha = 18
]
\DndMonsterDetails[
saving-throws = {Dex +7, Con +9, Wis +9, Cha +8},
skills = {Insight +9, Nature +11, Perception +9, Survival +13},
damage-resistances = {cold, fire},
damage-immunities = {poison},
condition-immunities = {charmed, frightened, poisoned},
senses = {truesight 120 ft., passive perception 19},
languages = {Common, Druidic, Elvish, Sylvan},
challenge = 12,
]
\DndMonsterAction{Fey Rebirth}
If the Gardener dies in the Eternal Garden, they revive with all their hit points 1d4 days later in a safe location within the garden.

\DndMonsterAction{Legendary Resistance (3/Day)}
If the Gardener fails a saving throw, they can choose to succeed instead.

\DndMonsterAction{Unfettered Steps}
The Gardener is unaffected by difficult terrain.

\DndMonsterAction{Spellcasting}
The Gardener casts one of the following spells, requiring no material components and using Wisdom as the spellcasting ability (spell save DC 17):
\begin{DndMonsterSpells}
  \DndInnateSpellLevel{Druidcraft}
  \DndInnateSpellLevel[2]{Awaken}
  \DndInnateSpellLevel[1]{Heroes' Feast}
\end{DndMonsterSpells}
\DndMonsterSection{Actions}
\DndMonsterAction{Multiattack}
The Gardener makes two Vine attacks and can use Breath of Tranquility if available.

\DndMonsterAction{Vine}
\textsl{Melee Weapon Attack:} 9 to hit, reach 10 ft., one target. \textsl{Hit:} 11 (1d12 + 5) bludgeoning damage plus 9 (2d8) psychic damage, and if the target is a Large or smaller creature, the vine wraps around the target, and the target has the grappled condition (escape DC 17). The vine vanishes when the target is no longer grappled, or when the Gardener wills it to (no action required). A creature reduced to 0 hit points by the vine has the unconscious condition but is stable instead of dying.

\DndMonsterAction{Breath of Tranquility (Recharge 5\textendash 6)}
The Gardener exhales soporific vapor in a 30-foot cone. Each creature in that area must succeed on a DC 17 Constitution saving throw or have the poisoned condition. A poisoned creature must repeat the saving throw at the end of its next turn. On a failed save, it has the unconscious condition, and on a successful save, the effect ends on it. An unconscious creature is no longer poisoned and remains unconscious for 1 hour, until it takes damage, or until a creature uses an action to shake it awake.

\DndMonsterSection{Reactions}
\DndMonsterAction{Pacification}
When a creature within 120 feet of the Gardener damages the Gardener, that creature takes 10 (3d6) psychic damage, and the Gardener teleports, along with anything they are wearing or carrying, to an unoccupied space they can see within 15 feet. A creature reduced to 0 hit points by the psychic damage has the unconscious condition and is stable instead of dying.

\DndMonsterAction{Spell Refuge}
When the Gardener or a creature within 30 feet of the Gardener takes damage from a spell, the Gardener chooses up to 5 creatures within 30 feet of themself. The Gardener and the chosen creatures have resistance to all damage from the triggering spell.

\DndMonsterSection{Lair Actions}
On initiative count 20 (losing initiative ties), the Gardener can take one of the following lair actions; they can't take the same lair action two rounds in a row:


\begin{description}
\item[Plant Walk.]
The Gardener moves up to their speed into the space of a Medium or larger plant within the lair, then teleports to an unoccupied space within 5 feet of any other Medium or larger plant within the lair.
\item[Vine Wall.]
A wall of grasping vines appears on the ground within 120 feet of the Gardener. The wall is up to 60 feet long, 10 feet high, and 5 feet thick, and it blocks line of sight. When the wall appears, each creature in its area must succeed on a DC 15 Dexterity saving throw or have the grappled condition (escape DC 15). A creature must make this saving throw when it starts its turn inside the wall or when it enters the wall for the first time on a turn. The wall's area is difficult terrain for a creature that isn't grappled. The wall sinks back into the ground, freeing all creatures grappled by it, when the Gardener uses this lair action again or when the Gardener dies.
\end{description}

\DndMonsterSection{Regional Effects}
The Eternal Garden has the following features:


\begin{description}
\item[Boundary.]
At the border of the domain is a lightly obscuring mist, tinged with a sweet blossom scent and the colors of sunset. No matter how far a creature travels into the mist, the creature finds itself back where it entered the border, returning to the garden.
\item[Eternal Summer.]
The climate in the garden is a perpetual, lovely summer, but weather conditions reflect the emotional tone of creatures in the garden. In times of danger, thunderstorms shake the sky. During times of melancholy or mourning, chill rain patters throughout the domain.
\item[Gardener's Vigil.]
The Gardener senses when anything dies in the garden. If a creature dies in the garden, roll a d6. On a 1, the Gardener appears on initiative count 20 of the next round (losing initiative ties) to investigate.
\item[Lighting.]
The sky above the garden is a perpetual twilight range of orange, red, pink, and yellow hues with no visible sun, moon, or stars. During daylight hours, it emits bright light, but at night, it never reaches full darkness, becoming dim light instead.
\item[Memory Loss.]
Unless otherwise stated, a creature that leaves the domain must make a DC 10 Wisdom saving throw; Fey creatures and creatures with the Fey Ancestry trait automatically succeed on the saving throw. On a failed save, the creature remembers nothing of its time in the garden or the Palace of Spires. On a successful save, the creature's memories remain intact, but they're hazy and dreamlike. A Remove Curse or Greater Restoration spell restores the creature's memories.
\item[Time.]
Time passes slower in the Eternal Garden relative to other planes. For each day spent in the garden, one year passes on the Material Plane.
\end{description}

\end{multicols}
\end{DndMonster}



\begin{DndMonster}[float=t]{Thri-kreen}
\DndMonsterType{Medium humanoid (thri-kreen), chaotic neutral}
\DndMonsterBasics[
armor-class = {15 (natural armor)},
hit-points = {\DndDice{6d8 + 6}},
speed = {40 ft.}
]
\DndMonsterAbilityScores[
 str = 12,
 dex = 15,
 con = 13,
 int = 8,
 wis = 12,
 cha = 7
]
\DndMonsterDetails[
skills = {Perception +3, Stealth +4, Survival +3},
senses = {darkvision 60 ft., passive perception 13},
languages = {Thri-kreen},
challenge = 1,
]
\DndMonsterAction{Chameleon Carapace}
The thri-kreen can change the color of its carapace to match the color and texture of its surroundings. As a result, it has advantage on Dexterity (Stealth) checks made to hide.

\DndMonsterAction{Standing Leap}
The thri-kreen's long jump is up to 30 feet and its high jump is up to 15 feet, with or without a running start.

\DndMonsterSection{Actions}
\DndMonsterAction{Multiattack}
The thri-kreen makes two attacks: one with its bite and one with its claws.

\DndMonsterAction{Bite}
\textsl{Melee Weapon Attack:} 3 to hit, reach 5 ft., one creature. \textsl{Hit:} 4 (1d6 + 1) piercing damage, and the target must succeed on a DC 11 Constitution saving throw or be poisoned for 1 minute. If the saving throw fails by 5 or more, the target is also paralyzed while poisoned in this way. The poisoned target can repeat the saving throw at the end of each of its turns, ending the effect on itself on a success.

\DndMonsterAction{Claws}
\textsl{Melee Weapon Attack:} 3 to hit, reach 5 ft., one target. \textsl{Hit:} 6 (2d4 + 1) slashing damage.

\end{DndMonster}



\begin{DndMonster}[float=t]{Tower Hand}
\DndMonsterType{Medium humanoid, any alignment}
\DndMonsterBasics[
armor-class = {14 (Unarmored Defense)},
hit-points = {\DndDice{4d8 + 4}},
speed = {40 ft.}
]
\DndMonsterAbilityScores[
 str = 10,
 dex = 14,
 con = 12,
 int = 12,
 wis = 14,
 cha = 9
]
\DndMonsterDetails[
skills = {Insight +4},
languages = {Common},
challenge = 1/2,
]
\DndMonsterAction{Unarmored Defense}
While the tower hand is wearing no armor and wielding no shield, its AC includes its Wisdom modifier.

\DndMonsterSection{Actions}
\DndMonsterAction{Multiattack}
The tower hand makes two Unarmed Strike attacks, two Dart attacks, or one of each.

\DndMonsterAction{Unarmed Strike}
\textsl{Melee Weapon Attack:} 4 to hit, reach 5 ft., one target. \textsl{Hit:} 5 (1d6 + 2) bludgeoning damage.

\DndMonsterAction{Dart}
\textsl{Ranged Weapon Attack:} 4 to hit, range 20/60 ft., one target. \textsl{Hit:} 4 (1d4 + 2) piercing damage.

\DndMonsterSection{Reactions}
\DndMonsterAction{Deflect Missile}
In response to being hit by a ranged weapon attack, the tower hand deflects the missile. The damage it takes from the attack is reduced by 7 (1d10 + 2).

\end{DndMonster}



\begin{DndMonster}[float=t]{Tower Sage}
\DndMonsterType{Medium humanoid, any alignment}
\DndMonsterBasics[
armor-class = {10 (13 with \textit{mage armor})},
hit-points = {\DndDice{5d8}},
speed = {30 ft.}
]
\DndMonsterAbilityScores[
 str = 10,
 dex = 10,
 con = 10,
 int = 16,
 wis = 14,
 cha = 16
]
\DndMonsterDetails[
skills = {Arcana +5, History +5, Insight +4},
languages = {Common plus any three languages},
challenge = 1,
]
\DndMonsterAction{Spellcasting}
The tower sage casts one of the following spells, using Intelligence as the spellcasting ability (spell save DC 13):
\begin{DndMonsterSpells}
  \DndInnateSpellLevel{Dancing Lights}
  \DndInnateSpellLevel[1]{Arcane Lock}
\end{DndMonsterSpells}
\DndMonsterSection{Actions}
\DndMonsterAction{Multiattack}
The tower sage makes two Arcane Burst attacks and can use Starry Radiance if available.

\DndMonsterAction{Arcane Burst}
\textsl{Melee or Ranged Spell Attack:} 5 to hit, reach 5 ft. or range 120 ft., one target. \textsl{Hit:} 8 (1d10 + 3) radiant damage.

\DndMonsterAction{Starry Radiance (Recharge 5\textendash 6)}
Dazzling light bursts from the tower sage's fingertips in a 15-foot cone. Each creature in that area must succeed on a DC 13 Constitution saving throw or have the blinded condition until the end of the tower sage's next turn.

\end{DndMonster}



\begin{DndMonster}[float=t]{Treant}
\DndMonsterType{Huge plant, chaotic good}
\DndMonsterBasics[
armor-class = {16 (natural armor)},
hit-points = {\DndDice{12d12 + 60}},
speed = {30 ft.}
]
\DndMonsterAbilityScores[
 str = 23,
 dex = 8,
 con = 21,
 int = 12,
 wis = 16,
 cha = 12
]
\DndMonsterDetails[
damage-vulnerabilities = {fire},
damage-resistances = {bludgeoning, piercing},
languages = {Common, Druidic, Elvish, Sylvan},
challenge = 9,
]
\DndMonsterAction{False Appearance}
While the treant remains motionless, it is indistinguishable from a normal tree.

\DndMonsterAction{Siege Monster}
The treant deals double damage to objects and structures.

\DndMonsterSection{Actions}
\DndMonsterAction{Multiattack}
The treant makes two slam attacks.

\DndMonsterAction{Slam}
\textsl{Melee Weapon Attack:} 10 to hit, reach 5 ft., one target. \textsl{Hit:} 16 (3d6 + 6) bludgeoning damage.

\DndMonsterAction{Rock}
\textsl{Ranged Weapon Attack:} 10 to hit, range 60/180 ft., one target. \textsl{Hit:} 28 (4d10 + 6) bludgeoning damage.

\DndMonsterAction{Animate Trees (1/Day)}
The treant magically animates one or two trees it can see within 60 feet of it. These trees have the same statistics as a \textbf{treant}, except they have Intelligence and Charisma scores of 1, they can't speak, and they have only the Slam action option. An animated tree acts as an ally of the treant. The tree remains animate for 1 day or until it dies; until the treant dies or is more than 120 feet from the tree; or until the treant takes a bonus action to turn it back into an inanimate tree. The tree then takes root if possible.

\end{DndMonster}



\begin{DndMonster}[float=t]{Tridrone}
\DndMonsterType{Medium construct, lawful neutral}
\DndMonsterBasics[
armor-class = {15 (natural armor)},
hit-points = {\DndDice{3d8 + 3}},
speed = {30 ft.}
]
\DndMonsterAbilityScores[
 str = 12,
 dex = 13,
 con = 12,
 int = 9,
 wis = 10,
 cha = 9
]
\DndMonsterDetails[
senses = {truesight 120 ft., passive perception 10},
languages = {Modron},
challenge = 1/2,
]
\DndMonsterAction{Axiomatic Mind}
The tridrone can't be compelled to act in a manner contrary to its nature or its instructions.

\DndMonsterAction{Disintegration}
If the tridrone dies, its body disintegrates into dust, leaving behind its weapons and anything else it was carrying.

\DndMonsterSection{Actions}
\DndMonsterAction{Multiattack}
The tridrone makes three fist attacks or three javelin attacks.

\DndMonsterAction{Fist}
\textsl{Melee Weapon Attack:} 3 to hit, reach 5 ft., one target. \textsl{Hit:} 3 (1d4 + 1) bludgeoning damage.

\DndMonsterAction{Javelin}
\textsl{Melee or Ranged Weapon Attack:} 3 to hit, reach 5 ft. or range 30/120 ft., one target. \textsl{Hit:} 4 (1d6 + 1) piercing damage.

\end{DndMonster}



\begin{DndMonster}[float=t]{Troglodyte}
\DndMonsterType{Medium humanoid (troglodyte), chaotic evil}
\DndMonsterBasics[
armor-class = {11 (natural armor)},
hit-points = {\DndDice{2d8 + 4}},
speed = {30 ft.}
]
\DndMonsterAbilityScores[
 str = 14,
 dex = 10,
 con = 14,
 int = 6,
 wis = 10,
 cha = 6
]
\DndMonsterDetails[
skills = {Stealth +2},
senses = {darkvision 60 ft., passive perception 10},
languages = {Troglodyte},
challenge = 1/4,
]
\DndMonsterAction{Chameleon Skin}
The troglodyte has advantage on Dexterity (Stealth) checks made to hide.

\DndMonsterAction{Stench}
Any creature other than a troglodyte that starts its turn within 5 feet of the troglodyte must succeed on a DC 12 Constitution saving throw or be poisoned until the start of the creature's next turn. On a successful saving throw, the creature is immune to the stench of all troglodytes for 1 hour.

\DndMonsterAction{Sunlight Sensitivity}
While in sunlight, the troglodyte has disadvantage on attack rolls, as well as on Wisdom (Perception) checks that rely on sight.

\DndMonsterSection{Actions}
\DndMonsterAction{Multiattack}
The troglodyte makes three attacks: one with its bite and two with its claws.

\DndMonsterAction{Bite}
\textsl{Melee Weapon Attack:} 4 to hit, reach 5 ft., one target. \textsl{Hit:} 4 (1d4 + 2) piercing damage.

\DndMonsterAction{Claw}
\textsl{Melee Weapon Attack:} 4 to hit, reach 5 ft., one target. \textsl{Hit:} 4 (1d4 + 2) slashing damage.

\end{DndMonster}


\clearpage

\begin{DndMonster}[float=t]{Troll}
\DndMonsterType{Large giant, chaotic evil}
\DndMonsterBasics[
armor-class = {15 (natural armor)},
hit-points = {\DndDice{8d10 + 40}},
speed = {30 ft.}
]
\DndMonsterAbilityScores[
 str = 18,
 dex = 13,
 con = 20,
 int = 7,
 wis = 9,
 cha = 7
]
\DndMonsterDetails[
skills = {Perception +2},
senses = {darkvision 60 ft., passive perception 12},
languages = {Giant},
challenge = 5,
]
\DndMonsterAction{Keen Smell}
The troll has advantage on Wisdom (Perception) checks that rely on smell.

\DndMonsterAction{Regeneration}
The troll regains 10 hit points at the start of its turn. If the troll takes acid or fire damage, this trait doesn't function at the start of the troll's next turn. The troll dies only if it starts its turn with 0 hit points and doesn't regenerate.

\DndMonsterSection{Actions}
\DndMonsterAction{Multiattack}
The troll makes three attacks: one with its bite and two with its claws.

\DndMonsterAction{Bite}
\textsl{Melee Weapon Attack:} 7 to hit, reach 5 ft., one target. \textsl{Hit:} 7 (1d6 + 4) piercing damage.

\DndMonsterAction{Claw}
\textsl{Melee Weapon Attack:} 7 to hit, reach 5 ft., one target. \textsl{Hit:} 11 (2d6 + 4) slashing damage.

\end{DndMonster}



\begin{DndMonster}[float=t]{Tyrannosaurus Rex}
\DndMonsterType{Huge beast, unaligned}
\DndMonsterBasics[
armor-class = {13 (natural armor)},
hit-points = {\DndDice{13d12 + 52}},
speed = {50 ft.}
]
\DndMonsterAbilityScores[
 str = 25,
 dex = 10,
 con = 19,
 int = 2,
 wis = 12,
 cha = 9
]
\DndMonsterDetails[
skills = {Perception +4},
challenge = 8,
]
\DndMonsterSection{Actions}
\DndMonsterAction{Multiattack}
The tyrannosaurus makes two attacks: one with its bite and one with its tail. It can't make both attacks against the same target.

\DndMonsterAction{Bite}
\textsl{Melee Weapon Attack:} 10 to hit, reach 10 ft., one target. \textsl{Hit:} 33 (4d12 + 7) piercing damage. If the target is a Medium or smaller creature, it is grappled (escape DC 17). Until this grapple ends, the target is restrained, and the tyrannosaurus can't bite another target.

\DndMonsterAction{Tail}
\textsl{Melee Weapon Attack:} 10 to hit, reach 10 ft., one target. \textsl{Hit:} 20 (3d8 + 7) bludgeoning damage.

\end{DndMonster}



\begin{DndMonster}[float=t]{Umber Hulk}
\DndMonsterType{Large monstrosity, chaotic evil}
\DndMonsterBasics[
armor-class = {18 (natural armor)},
hit-points = {\DndDice{11d10 + 33}},
speed = {30 ft., burrow 20 ft.}
]
\DndMonsterAbilityScores[
 str = 20,
 dex = 13,
 con = 16,
 int = 9,
 wis = 10,
 cha = 10
]
\DndMonsterDetails[
senses = {darkvision 120 ft., tremorsense 60 ft., passive perception 10},
languages = {Umber Hulk},
challenge = 5,
]
\DndMonsterAction{Confusing Gaze}
When a creature starts its turn within 30 feet of the umber hulk and is able to see the umber hulk's eyes, the umber hulk can magically force it to make a DC 15 Charisma saving throw, unless the umber hulk is incapacitated.

On a failed saving throw, the creature can't take reactions until the start of its next turn and rolls a d8 to determine what it does during that turn. On a 1 to 4, the creature does nothing. On a 5 or 6, the creature takes no action but uses all its movement to move in a random direction. On a 7 or 8, the creature makes one melee attack against a random creature, or it does nothing if no creature is within reach.

Unless surprised, a creature can avert its eyes to avoid the saving throw at the start of its turn. If the creature does so, it can't see the umber hulk until the start of its next turn, when it can avert its eyes again. If the creature looks at the umber hulk in the meantime, it must immediately make the save.

\DndMonsterAction{Tunneler}
The umber hulk can burrow through solid rock at half its burrowing speed and leaves a 5 foot-wide, 8-foot-high tunnel in its wake.

\DndMonsterSection{Actions}
\DndMonsterAction{Multiattack}
The umber hulk makes three attacks: two with its claws and one with its mandibles.

\DndMonsterAction{Claw}
\textsl{Melee Weapon Attack:} 8 to hit, reach 5 ft., one target. \textsl{Hit:} 9 (1d8 + 5) slashing damage.

\DndMonsterAction{Mandibles}
\textsl{Melee Weapon Attack:} 8 to hit, reach 5 ft., one target. \textsl{Hit:} 14 (2d8 + 5) slashing damage.

\end{DndMonster}



\begin{DndMonster}[float*=b,width=\textwidth + 8pt]{Unicorn}
\begin{multicols}{2}
\DndMonsterType{Large celestial, lawful good}
\DndMonsterBasics[
armor-class = {12},
hit-points = {\DndDice{9d10 + 18}},
speed = {50 ft.}
]
\DndMonsterAbilityScores[
 str = 18,
 dex = 14,
 con = 15,
 int = 11,
 wis = 17,
 cha = 16
]
\DndMonsterDetails[
damage-immunities = {poison},
condition-immunities = {charmed, paralyzed, poisoned},
senses = {darkvision 60 ft., passive perception 13},
languages = {Celestial, Elvish, Sylvan, telepathy 60 ft.},
challenge = 5,
]
\DndMonsterAction{Charge}
If the unicorn moves at least 20 feet straight toward a target and then hits it with a horn attack on the same turn, the target takes an extra 9 (2d8) piercing damage. If the target is a creature, it must succeed on a DC 15 Strength saving throw or be knocked prone.

\DndMonsterAction{Magic Resistance}
The unicorn has advantage on saving throws against spells and other magical effects.

\DndMonsterAction{Magic Weapons}
The unicorn's weapon attacks are magical.

\DndMonsterAction{Innate Spellcasting}
The unicorn's innate spellcasting ability is Charisma (spell save DC 14). The unicorn can innately cast the following spells, requiring no components:
\begin{DndMonsterSpells}
  \DndInnateSpellLevel{detect evil and good}
  \DndInnateSpellLevel[1]{calm emotions}
\end{DndMonsterSpells}
\DndMonsterSection{Actions}
\DndMonsterAction{Multiattack}
The unicorn makes two attacks: one with its hooves and one with its horn.

\DndMonsterAction{Hooves}
\textsl{Melee Weapon Attack:} 7 to hit, reach 5 ft., one target. \textsl{Hit:} 11 (2d6 + 4) bludgeoning damage.

\DndMonsterAction{Horn}
\textsl{Melee Weapon Attack:} 7 to hit, reach 5 ft., one target. \textsl{Hit:} 8 (1d8 + 4) piercing damage.

\DndMonsterAction{Healing Touch (3/Day)}
The unicorn touches another creature with its horn. The target magically regains 11 (2d8 + 2) hit points. In addition, the touch removes all diseases and neutralizes all poisons afflicting the target.

\DndMonsterAction{Teleport (1/Day)}
The unicorn magically teleports itself and up to three willing creatures it can see within 5 feet of it, along with any equipment they are wearing or carrying, to a location the unicorn is familiar with, up to 1 mile away.

\DndMonsterSection{Legendary Actions}
Unicorn can take 3 legendary actions, choosing from the options below. Only one legendary action can be used at a time and only at the end of another creature's turn. Unicorn regains spent legendary actions at the start of its turn.

\begin{DndMonsterLegendaryActions}
\DndMonsterLegendaryAction{Hooves}{The unicorn makes one attack with its hooves.
}
\DndMonsterLegendaryAction{Shimmering Shield (Costs 2 Actions)}{The unicorn creates a shimmering, magical field around itself or another creature it can see within 60 feet of it. The target gains a +2 bonus to AC until the end of the unicorn's next turn.
}
\DndMonsterLegendaryAction{Heal Self (Costs 3 Actions)}{The unicorn magically regains 11 (2d8 + 2) hit points.
}
\end{DndMonsterLegendaryActions}
\DndMonsterSection{Regional Effects}
Transformed by the creature's celestial presence, the domain of a unicorn might include any of the following magical effects:


\begin{itemize}
\item Open flames of a non magical nature are extinguished within the unicorn's domain. Torches and campfires refuse to burn, but closed lanterns are unaffected.
\item Creatures native to the unicorn's domain have an easier time hiding; they have advantage on all Dexterity (Stealth) checks made to hide.
\item When a good-aligned creature casts a spell or uses a magical effect that causes another good-aligned creature to regain hit points, the target regains the maximum number of hit points possible for the spell or effect.
\item Curses affecting any good-aligned creature are suppressed.
\end{itemize}

If the unicorn dies, these effects end immediately.

\end{multicols}
\end{DndMonster}



\begin{DndMonster}[float=t]{Vegepygmy Moldmaker}
\DndMonsterType{Small plant, neutral}
\DndMonsterBasics[
armor-class = {15 (natural armor)},
hit-points = {\DndDice{8d6 + 16}},
speed = {30 ft.}
]
\DndMonsterAbilityScores[
 str = 10,
 dex = 14,
 con = 14,
 int = 10,
 wis = 16,
 cha = 11
]
\DndMonsterDetails[
saving-throws = {Int +2, Wis +5},
skills = {Perception +5, Stealth +4},
damage-resistances = {lightning, piercing},
senses = {darkvision 60 ft., passive perception 15},
languages = {Vegepygmy},
challenge = 3,
]
\DndMonsterAction{Plant Camouflage}
The vegepygmy has advantage on Dexterity (Stealth) checks it makes in any terrain with ample obscuring vegetation.

\DndMonsterAction{Regeneration}
The vegepygmy regains 7 hit points at the start of its turn. If it takes cold, fire, or necrotic damage, this trait doesn't function at the start of the vegepygmy's next turn. The vegepygmy dies only if it starts its turn with 0 hit points and doesn't regenerate.

\DndMonsterSection{Actions}
\DndMonsterAction{Multiattack}
The vegepygmy makes two Claw attacks.

\DndMonsterAction{Claw}
\textsl{Melee Weapon Attack:} 4 to hit, reach 5 ft., one target. \textsl{Hit:} 12 (3d6 + 2) slashing damage.

\DndMonsterAction{Toxic Mold (2/Day)}
The vegepygmy targets a creature it can see within 60 feet of itself. If the target isn't a vegepygmy, it must make a DC 13 Constitution saving throw. On a failed save, the target takes 13 (3d8) poison damage and has the blinded and deafened conditions for 1 minute as it becomes covered in a thick layer of mold. On a successful save, the target takes half as much damage only.

The target can repeat the saving throw at the end of each of its turns, ending the effect on itself on a success.

\end{DndMonster}



\begin{DndMonster}[float=t]{Vegepygmy Scavenger}
\DndMonsterType{Small plant, neutral}
\DndMonsterBasics[
armor-class = {13 (natural armor)},
hit-points = {\DndDice{3d6 + 3}},
speed = {30 ft.}
]
\DndMonsterAbilityScores[
 str = 7,
 dex = 14,
 con = 13,
 int = 6,
 wis = 11,
 cha = 7
]
\DndMonsterDetails[
skills = {Perception +2, Sleight Of Hand +4, Stealth +4},
damage-resistances = {lightning, piercing},
senses = {darkvision 60 ft., passive perception 12},
languages = {Vegepygmy},
challenge = 1/4,
]
\DndMonsterAction{Plant Camouflage}
The vegepygmy has advantage on Dexterity (Stealth) checks it makes in any terrain with ample obscuring vegetation.

\DndMonsterAction{Regeneration}
The vegepygmy regains 3 hit points at the start of its turn. If it takes cold, fire, or necrotic damage, this trait doesn't function at the start of the vegepygmy's next turn. The vegepygmy dies only if it starts its turn with 0 hit points and doesn't regenerate.

\DndMonsterSection{Actions}
\DndMonsterAction{Claw}
\textsl{Melee Weapon Attack:} 4 to hit, reach 5 ft., one target. \textsl{Hit:} 5 (1d6 + 2) slashing damage.

\DndMonsterAction{Sling}
\textsl{Ranged Weapon Attack:} 4 to hit, range 30/120 ft., one target. \textsl{Hit:} 4 (1d4 + 2) bludgeoning damage.

\DndMonsterSection{Bonus Actions}
\DndMonsterAction{Nimble Escape}
The vegepygmy takes the Disengage or Hide action.

\end{DndMonster}



\begin{DndMonster}[float=t]{Vegepygmy Thorny Hunter}
\DndMonsterType{Medium plant, unaligned}
\DndMonsterBasics[
armor-class = {14 (natural armor)},
hit-points = {\DndDice{5d8 + 5}},
speed = {40 ft.}
]
\DndMonsterAbilityScores[
 str = 15,
 dex = 14,
 con = 13,
 int = 2,
 wis = 10,
 cha = 6
]
\DndMonsterDetails[
skills = {Perception +4, Stealth +4, Survival +4},
damage-resistances = {lightning, piercing},
senses = {darkvision 60 ft., passive perception 14},
challenge = 2,
]
\DndMonsterAction{Pack Tactics}
The vegepygmy has advantage on attack rolls against a creature if at least one of the vegepygmy's allies is within 5 feet of the creature and the ally doesn't have the incapacitated condition.

\DndMonsterAction{Plant Camouflage}
The vegepygmy has advantage on Dexterity (Stealth) checks it makes in any terrain with ample obscuring vegetation.

\DndMonsterAction{Regeneration}
The vegepygmy regains 5 hit points at the start of its turn. If it takes cold, fire, or necrotic damage, this trait doesn't function at the start of the vegepygmy's next turn. The vegepygmy dies only if it starts its turn with 0 hit points and doesn't regenerate.

\DndMonsterAction{Thorny Body}
At the start of its turn, the vegepygmy deals 2 (1d4) piercing damage to any creature grappling it.

\DndMonsterSection{Actions}
\DndMonsterAction{Bite}
\textsl{Melee Weapon Attack:} 4 to hit, reach 5 ft., one target. \textsl{Hit:} 11 (2d8 + 2) piercing damage. If the target is a creature, it must succeed on a DC 12 Strength saving throw or have the prone condition.

\end{DndMonster}



\begin{DndMonster}[float=t]{Veteran}
\DndMonsterType{Medium humanoid (any race), any alignment}
\DndMonsterBasics[
armor-class = {17 (\textit{splint armor})},
hit-points = {\DndDice{9d8 + 18}},
speed = {30 ft.}
]
\DndMonsterAbilityScores[
 str = 16,
 dex = 13,
 con = 14,
 int = 10,
 wis = 11,
 cha = 10
]
\DndMonsterDetails[
skills = {Athletics +5, Perception +2},
languages = {any one language (usually Common)},
challenge = 3,
]
\DndMonsterSection{Actions}
\DndMonsterAction{Multiattack}
The veteran makes two longsword attacks. If it has a shortsword drawn, it can also make a shortsword attack.

\DndMonsterAction{Longsword}
\textsl{Melee Weapon Attack:} 5 to hit, reach 5 ft., one target. \textsl{Hit:} 7 (1d8 + 3) slashing damage, or 8 (1d10 + 3) slashing damage if used with two hands.

\DndMonsterAction{Shortsword}
\textsl{Melee Weapon Attack:} 5 to hit, reach 5 ft., one target. \textsl{Hit:} 6 (1d6 + 3) piercing damage.

\DndMonsterAction{Heavy Crossbow}
\textsl{Ranged Weapon Attack:} 3 to hit, range 100/400 ft., one target. \textsl{Hit:} 6 (1d10 + 1) piercing damage.

\end{DndMonster}



\begin{DndMonster}[float=t]{Warrior of Madarua}
\DndMonsterType{Medium humanoid, chaotic good}
\DndMonsterBasics[
armor-class = {12 (\textit{leather armor})},
hit-points = {\DndDice{2d8 + 2}},
speed = {30 ft.}
]
\DndMonsterAbilityScores[
 str = 15,
 dex = 12,
 con = 13,
 int = 10,
 wis = 10,
 cha = 10
]
\DndMonsterDetails[
skills = {Acrobatics +3, Religion +2},
languages = {Common},
challenge = 1/8,
]
\DndMonsterSection{Actions}
\DndMonsterAction{Spear}
\textsl{Melee or Ranged Weapon Attack:} 4 to hit, reach 5 ft. or range 20/60 ft., one target. \textsl{Hit:} 5 (1d6 + 2) piercing damage, or 6 (1d8 + 2) piercing damage if used with two hands to make a melee attack.

\DndMonsterSection{Reactions}
\DndMonsterAction{Parry}
The warrior adds 2 to its AC against one melee attack that would hit it. To do so, the warrior must see the attacker and be wielding a melee weapon.

\end{DndMonster}



\begin{DndMonster}[float=t]{Water Elemental}
\DndMonsterType{Large elemental, neutral}
\DndMonsterBasics[
armor-class = {14 (natural armor)},
hit-points = {\DndDice{12d10 + 48}},
speed = {30 ft., swim 90 ft.}
]
\DndMonsterAbilityScores[
 str = 18,
 dex = 14,
 con = 18,
 int = 5,
 wis = 10,
 cha = 8
]
\DndMonsterDetails[
damage-resistances = {acid; bludgeoning, piercing, slashing from nonmagical attacks},
damage-immunities = {poison},
condition-immunities = {exhaustion, grappled, paralyzed, petrified, poisoned, prone, restrained, unconscious},
senses = {darkvision 60 ft., passive perception 10},
languages = {Aquan},
challenge = 5,
]
\DndMonsterAction{Water Form}
The elemental can enter a hostile creature's space and stop there. It can move through a space as narrow as 1 inch wide without squeezing.

\DndMonsterAction{Freeze}
If the elemental takes cold damage, it partially freezes; its speed is reduced by 20 feet until the end of its next turn.

\DndMonsterSection{Actions}
\DndMonsterAction{Multiattack}
The elemental makes two slam attacks.

\DndMonsterAction{Slam}
\textsl{Melee Weapon Attack:} 7 to hit, reach 5 ft., one target. \textsl{Hit:} 13 (2d8 + 4) bludgeoning damage.

\DndMonsterAction{Whelm (Recharge 4\textendash 6)}
Each creature in the elemental's space must make a DC 15 Strength saving throw. On a failure, a target takes 13 (2d8 + 4) bludgeoning damage. If it is Large or smaller, it is also grappled (escape DC 14). Until this grapple ends, the target is restrained and unable to breathe unless it can breathe water. If the saving throw is successful, the target is pushed out of the elemental's space.

The elemental can grapple one Large creature or up to two Medium or smaller creatures at one time. At the start of each of the elemental's turns, each target grappled by it takes 13 (2d8 + 4) bludgeoning damage. A creature within 5 feet of the elemental can pull a creature or object out of it by taking an action to make a DC 14 Strength check and succeeding.

\end{DndMonster}



\begin{DndMonster}[float=t]{Water Weird}
\DndMonsterType{Large elemental, neutral}
\DndMonsterBasics[
armor-class = {13},
hit-points = {\DndDice{9d10 + 9}},
speed = {0 ft., swim 60 ft.}
]
\DndMonsterAbilityScores[
 str = 17,
 dex = 16,
 con = 13,
 int = 11,
 wis = 10,
 cha = 10
]
\DndMonsterDetails[
damage-resistances = {fire; bludgeoning, piercing, slashing from nonmagical attacks},
damage-immunities = {poison},
condition-immunities = {exhaustion, grappled, paralyzed, poisoned, restrained, prone, unconscious},
senses = {blindsight 30 ft., passive perception 10},
languages = {understands Aquan but doesn't speak},
challenge = 3,
]
\DndMonsterAction{Invisible in Water}
The water weird is invisible while fully immersed in water.

\DndMonsterAction{Water Bound}
The water weird dies if it leaves the water to which it is bound or if that water is destroyed.

\DndMonsterSection{Actions}
\DndMonsterAction{Constrict}
\textsl{Melee Weapon Attack:} 5 to hit, reach 10 ft., one creature. \textsl{Hit:} 13 (3d6 + 3) bludgeoning damage. If the target is Medium or smaller, it is grappled (escape DC 13) and pulled 5 feet toward the water weird. Until this grapple ends, the target is restrained, the water weird tries to drown it, and the water weird can't constrict another target.

\end{DndMonster}



\begin{DndMonster}[float=t]{Weasel}
\DndMonsterType{Tiny beast, unaligned}
\DndMonsterBasics[
armor-class = {13},
hit-points = {\DndDice{1d4 - 1}},
speed = {30 ft.}
]
\DndMonsterAbilityScores[
 str = 3,
 dex = 16,
 con = 8,
 int = 2,
 wis = 12,
 cha = 3
]
\DndMonsterDetails[
skills = {Perception +3, Stealth +5},
challenge = 0,
]
\DndMonsterAction{Keen Hearing and Smell}
The weasel has advantage on Wisdom (Perception) checks that rely on hearing or smell.

\DndMonsterSection{Actions}
\DndMonsterAction{Bite}
\textsl{Melee Weapon Attack:} 5 to hit, reach 5 ft., one creature. \textsl{Hit:} 1 piercing damage.

\end{DndMonster}



\begin{DndMonster}[float=t]{Wererat}
\DndMonsterType{Medium humanoid (human, shapechanger), lawful evil}
\DndMonsterBasics[
armor-class = {12},
hit-points = {\DndDice{6d8 + 6}},
speed = {30 ft.}
]
\DndMonsterAbilityScores[
 str = 10,
 dex = 15,
 con = 12,
 int = 11,
 wis = 10,
 cha = 8
]
\DndMonsterDetails[
skills = {Perception +2, Stealth +4},
damage-immunities = {bludgeoning, piercing, slashing from nonmagical attacks that aren't silvered},
senses = {darkvision 60 ft. (rat form only), passive perception 12},
languages = {Common (can't speak in rat form)},
challenge = 2,
]
\DndMonsterAction{Shapechanger}
The wererat can use its action to polymorph into a rat-humanoid hybrid or into a giant rat, or back into its true form, which is humanoid. Its statistics, other than its size, are the same in each form. Any equipment it is wearing or carrying isn't transformed. It reverts to its true form if it dies.

\DndMonsterAction{Keen Smell}
The wererat has advantage on Wisdom (Perception) checks that rely on smell.

\DndMonsterSection{Actions}
\DndMonsterAction{Multiattack (Humanoid or Hybrid Form Only)}
The wererat makes two attacks, only one of which can be a bite.

\DndMonsterAction{Bite (Rat or Hybrid Form Only)}
\textsl{Melee Weapon Attack:} 4 to hit, reach 5 ft., one target. \textsl{Hit:} 4 (1d4 + 2) piercing damage. If the target is a humanoid, it must succeed on a DC 11 Constitution saving throw or be cursed with wererat lycanthropy.

\DndMonsterAction{Shortsword (Humanoid or Hybrid Form Only)}
\textsl{Melee Weapon Attack:} 4 to hit, reach 5 ft., one target. \textsl{Hit:} 5 (1d6 + 2) piercing damage.

\DndMonsterAction{Hand Crossbow (Humanoid or Hybrid Form Only)}
\textsl{Ranged Weapon Attack:} 4 to hit, range 30/120 ft., one target. \textsl{Hit:} 5 (1d6 + 2) piercing damage.

\end{DndMonster}



\begin{DndMonster}[float=t]{Wight}
\DndMonsterType{Medium undead, neutral evil}
\DndMonsterBasics[
armor-class = {14 (studded leather armor)},
hit-points = {\DndDice{6d8 + 18}},
speed = {30 ft.}
]
\DndMonsterAbilityScores[
 str = 15,
 dex = 14,
 con = 16,
 int = 10,
 wis = 13,
 cha = 15
]
\DndMonsterDetails[
skills = {Perception +3, Stealth +4},
damage-resistances = {necrotic; bludgeoning, piercing, slashing from nonmagical attacks that aren't silvered},
damage-immunities = {poison},
condition-immunities = {exhaustion, poisoned},
senses = {darkvision 60 ft., passive perception 13},
languages = {the languages it knew in life},
challenge = 3,
]
\DndMonsterAction{Sunlight Sensitivity}
While in sunlight, the wight has disadvantage on attack rolls, as well as on Wisdom (Perception) checks that rely on sight.

\DndMonsterSection{Actions}
\DndMonsterAction{Multiattack}
The wight makes two longsword attacks or two longbow attacks. It can use its Life Drain in place of one longsword attack.

\DndMonsterAction{Life Drain}
\textsl{Melee Weapon Attack:} 4 to hit, reach 5 ft., one creature. \textsl{Hit:} 5 (1d6 + 2) necrotic damage. The target must succeed on a DC 13 Constitution saving throw or its hit point maximum is reduced by an amount equal to the damage taken. This reduction lasts until the target finishes a long rest. The target dies if this effect reduces its hit point maximum to 0.

A humanoid slain by this attack rises 24 hours later as a \textbf{zombie} under the wight's control, unless the humanoid is restored to life or its body is destroyed. The wight can have no more than twelve zombies under its control at one time.

\DndMonsterAction{Longsword}
\textsl{Melee Weapon Attack:} 4 to hit, reach 5 ft., one target. \textsl{Hit:} 6 (1d8 + 2) slashing damage, or 7 (1d10 + 2) slashing damage if used with two hands.

\DndMonsterAction{Longbow}
\textsl{Ranged Weapon Attack:} 4 to hit, range 150/600 ft., one target. \textsl{Hit:} 6 (1d8 + 2) piercing damage.

\end{DndMonster}



\begin{DndMonster}[float=t]{Will-o'-Wisp}
\DndMonsterType{Tiny undead, chaotic evil}
\DndMonsterBasics[
armor-class = {19},
hit-points = {\DndDice{9d4}},
speed = {0 ft., fly 50 ft. \textit{(hover)}}
]
\DndMonsterAbilityScores[
 str = 1,
 dex = 28,
 con = 10,
 int = 13,
 wis = 14,
 cha = 11
]
\DndMonsterDetails[
damage-resistances = {acid; cold; fire; necrotic; thunder; bludgeoning, piercing, slashing from nonmagical attacks},
damage-immunities = {lightning, poison},
condition-immunities = {exhaustion, grappled, paralyzed, poisoned, prone, restrained, unconscious},
senses = {darkvision 120 ft., passive perception 12},
languages = {the languages it knew in life},
challenge = 2,
]
\DndMonsterAction{Consume Life}
As a bonus action, the will-o'-wisp can target one creature it can see within 5 feet of it that has 0 hit points and is still alive. The target must succeed on a DC 10 Constitution saving throw against this magic or die. If the target dies, the will-o'-wisp regains 10 (3d6) hit points.

\DndMonsterAction{Ephemeral}
The will-o'-wisp can't wear or carry anything.

\DndMonsterAction{Incorporeal Movement}
The will-o'-wisp can move through other creatures and objects as if they were difficult terrain. It takes 5 (1d10) force damage if it ends its turn inside an object.

\DndMonsterAction{Variable Illumination}
The will-o'-wisp sheds bright light in a 5 to 20-foot radius and dim light for an additional number of ft. equal to the chosen radius. The will-o'-wisp can alter the radius as a bonus action.

\DndMonsterSection{Actions}
\DndMonsterAction{Shock}
\textsl{Melee Spell Attack:} 4 to hit, reach 5 ft., one creature. \textsl{Hit:} 9 (2d8) lightning damage.

\DndMonsterAction{Invisibility}
The will-o'-wisp and its light magically become invisible until it attacks or uses its Consume Life, or until its concentration ends (as if concentration on a spell).

\end{DndMonster}



\begin{DndMonster}[float=t]{Wolf}
\DndMonsterType{Medium beast, unaligned}
\DndMonsterBasics[
armor-class = {13 (natural armor)},
hit-points = {\DndDice{2d8 + 2}},
speed = {40 ft.}
]
\DndMonsterAbilityScores[
 str = 12,
 dex = 15,
 con = 12,
 int = 3,
 wis = 12,
 cha = 6
]
\DndMonsterDetails[
skills = {Perception +3, Stealth +4},
challenge = 1/4,
]
\DndMonsterAction{Keen Hearing and Smell}
The wolf has advantage on Wisdom (Perception) checks that rely on hearing or smell.

\DndMonsterAction{Pack Tactics}
The wolf has advantage on an attack roll against a creature if at least one of the wolf's allies is within 5 feet of the creature and the ally isn't incapacitated.

\DndMonsterSection{Actions}
\DndMonsterAction{Bite}
\textsl{Melee Weapon Attack:} 4 to hit, reach 5 ft., one target. \textsl{Hit:} 7 (2d4 + 2) piercing damage. If the target is a creature, it must succeed on a DC 11 Strength saving throw or be knocked prone.

\end{DndMonster}



\begin{DndMonster}[float=t]{Wolf-in-Sheep's-Clothing}
\DndMonsterType{Medium plant, unaligned}
\DndMonsterBasics[
armor-class = {16 (natural armor)},
hit-points = {\DndDice{15d8 + 45}},
speed = {20 ft., climb 20 ft.}
]
\DndMonsterAbilityScores[
 str = 20,
 dex = 9,
 con = 16,
 int = 5,
 wis = 12,
 cha = 5
]
\DndMonsterDetails[
skills = {Perception +7, Stealth +5},
condition-immunities = {prone},
senses = {darkvision 60 ft., passive perception 17},
challenge = 7,
]
\DndMonsterAction{False Appearance}
If the wolf-in-sheep's-clothing is motionless (except for its lure) at the start of combat, it has advantage on its initiative roll. Moreover, if a creature hasn't observed the wolf-in-sheep's-clothing move or act (except for the lure), that creature must succeed on a DC 18 Intelligence (Investigation) check to discern that the wolf-in-sheep's-clothing is animate.

\DndMonsterSection{Actions}
\DndMonsterAction{Multiattack}
The wolf-in-sheep's-clothing makes one Bite attack and two Root Tentacle attacks.

\DndMonsterAction{Bite}
\textsl{Melee Weapon Attack:} 8 to hit, reach 5 ft., one target. \textsl{Hit:} 9 (1d8 + 5) piercing damage plus 7 (2d6) acid damage.

\DndMonsterAction{Root Tentacle}
\textsl{Melee Weapon Attack:} 8 to hit, reach 20 ft., one target. \textsl{Hit:} 8 (1d6 + 5) bludgeoning damage, and if the target is a Medium or smaller creature, it has the grappled condition (escape DC 16). While grappled in this way, the target has the restrained condition, and at the start of each of the wolf-in-sheep's-clothing's turns, the wolf-in-sheep's-clothing can pull the target up to 10 feet toward itself (no action required). The wolf-in-sheep's-clothing has four root tentacles, each of which can grapple one target.

\DndMonsterSection{Bonus Actions}
\DndMonsterAction{Completely Harmless Lure}
The wolf-in-sheep's-clothing can change the color, texture, and shape of its lure to resemble a Tiny Beast or Tiny object. It can move the lure to reinforce the resemblance (no action required), but the lure must remain within 15 feet of the wolf-in-sheep's-clothing, connected by nearly invisible filament-like tendrils.

\end{DndMonster}



\begin{DndMonster}[float=t]{Worker Robot}
\DndMonsterType{Medium construct, neutral}
\DndMonsterBasics[
armor-class = {13 (natural armor)},
hit-points = {\DndDice{10d8 + 30}},
speed = {30 ft., climb 30 ft.}
]
\DndMonsterAbilityScores[
 str = 20,
 dex = 9,
 con = 16,
 int = 10,
 wis = 12,
 cha = 10
]
\DndMonsterDetails[
skills = {Athletics +9},
damage-resistances = {acid, fire},
damage-immunities = {cold, poison},
condition-immunities = {exhaustion, poisoned},
senses = {darkvision 60 ft., passive perception 11},
languages = {Common plus the languages spoken by its creator},
challenge = 3,
]
\DndMonsterAction{Lightning Overload}
When the robot takes lightning damage, it must succeed on a DC 10 Constitution saving throw or have the stunned condition until the start of its next turn.

\DndMonsterSection{Actions}
\DndMonsterAction{Multiattack}
The robot makes two Cargo Tentacle attacks.

\DndMonsterAction{Cargo Tentacle}
\textsl{Melee Weapon Attack:} 7 to hit, reach 10 ft., one target. \textsl{Hit:} 9 (1d8 + 5) bludgeoning damage. If the target is a Medium or smaller creature, it has the grappled condition (escape DC 15). The robot has two cargo tentacles, each of which can grapple one target.

\DndMonsterAction{Tractor Beam}
The robot casts \textit{Telekinesis}, targeting only creatures with the incapacitated condition or objects. It requires no spell components and uses Wisdom as the spellcasting ability.

\end{DndMonster}



\begin{DndMonster}[float=t]{Wraith}
\DndMonsterType{Medium undead, neutral evil}
\DndMonsterBasics[
armor-class = {13},
hit-points = {\DndDice{9d8 + 27}},
speed = {0 ft., fly 60 ft. \textit{(hover)}}
]
\DndMonsterAbilityScores[
 str = 6,
 dex = 16,
 con = 16,
 int = 12,
 wis = 14,
 cha = 15
]
\DndMonsterDetails[
damage-resistances = {acid; cold; fire; lightning; thunder; bludgeoning, piercing, slashing from nonmagical attacks that aren't silvered},
damage-immunities = {necrotic, poison},
condition-immunities = {charmed, exhaustion, grappled, paralyzed, petrified, poisoned, prone, restrained},
senses = {darkvision 60 ft., passive perception 12},
languages = {the languages it knew in life},
challenge = 5,
]
\DndMonsterAction{Incorporeal Movement}
The wraith can move through other creatures and objects as if they were difficult terrain. It takes 5 (1d10) force damage if it ends its turn inside an object.

\DndMonsterAction{Sunlight Sensitivity}
While in sunlight, the wraith has disadvantage on attack rolls, as well as on Wisdom (Perception) checks that rely on sight.

\DndMonsterSection{Actions}
\DndMonsterAction{Life Drain}
\textsl{Melee Weapon Attack:} 6 to hit, reach 5 ft., one creature. \textsl{Hit:} 21 (4d8 + 3) necrotic damage. The target must succeed on a DC 14 Constitution saving throw or its hit point maximum is reduced by an amount equal to the damage taken. This reduction lasts until the target finishes a long rest. The target dies if this effect reduces its hit point maximum to 0.

\DndMonsterAction{Create Specter}
The wraith targets a humanoid within 10 feet of it that has been dead for no longer than 1 minute and died violently. The target's spirit rises as a \textbf{specter} in the space of its corpse or in the nearest unoccupied space. The \textbf{specter} is under the wraith's control. The wraith can have no more than seven specters under its control at one time.

\end{DndMonster}



\begin{DndMonster}[float=t]{Wyvern}
\DndMonsterType{Large dragon, unaligned}
\DndMonsterBasics[
armor-class = {13 (natural armor)},
hit-points = {\DndDice{13d10 + 39}},
speed = {20 ft., fly 80 ft.}
]
\DndMonsterAbilityScores[
 str = 19,
 dex = 10,
 con = 16,
 int = 5,
 wis = 12,
 cha = 6
]
\DndMonsterDetails[
skills = {Perception +4},
senses = {darkvision 60 ft., passive perception 14},
challenge = 6,
]
\DndMonsterSection{Actions}
\DndMonsterAction{Multiattack}
The wyvern makes two attacks: one with its bite and one with its stinger. While flying, it can use its claws in place of one other attack.

\DndMonsterAction{Bite}
\textsl{Melee Weapon Attack:} 7 to hit, reach 10 ft., one creature. \textsl{Hit:} 11 (2d6 + 4) piercing damage.

\DndMonsterAction{Claws}
\textsl{Melee Weapon Attack:} 7 to hit, reach 5 ft., one target. \textsl{Hit:} 13 (2d8 + 4) slashing damage.

\DndMonsterAction{Stinger}
\textsl{Melee Weapon Attack:} 7 to hit, reach 10 ft., one creature. \textsl{Hit:} 11 (2d6 + 4) piercing damage. The target must make a DC 15 Constitution saving throw, taking 24 (7d6) poison damage on a failed save, or half as much damage on a successful one.

\end{DndMonster}


\clearpage

\begin{DndMonster}[float=t]{Xorn}
\DndMonsterType{Medium elemental, neutral}
\DndMonsterBasics[
armor-class = {19 (natural armor)},
hit-points = {\DndDice{7d8 + 42}},
speed = {20 ft., burrow 20 ft.}
]
\DndMonsterAbilityScores[
 str = 17,
 dex = 10,
 con = 22,
 int = 11,
 wis = 10,
 cha = 11
]
\DndMonsterDetails[
skills = {Perception +6, Stealth +3},
damage-resistances = {piercing, slashing from nonmagical attacks that aren't adamantine},
senses = {darkvision 60 ft., tremorsense 60 ft., passive perception 16},
languages = {Terran},
challenge = 5,
]
\DndMonsterAction{Earth Glide}
The xorn can burrow through nonmagical, unworked earth and stone. While doing so, the xorn doesn't disturb the material it moves through.

\DndMonsterAction{Stone Camouflage}
The xorn has advantage on Dexterity (Stealth) checks made to hide in rocky terrain.

\DndMonsterAction{Treasure Sense}
The xorn can pinpoint, by scent, the location of precious metals and stones, such as coins and gems, within 60 feet of it.

\DndMonsterSection{Actions}
\DndMonsterAction{Multiattack}
The xorn makes three claw attacks and one bite attack.

\DndMonsterAction{Bite}
\textsl{Melee Weapon Attack:} 6 to hit, reach 5 ft., one target. \textsl{Hit:} 13 (3d6 + 3) piercing damage.

\DndMonsterAction{Claw}
\textsl{Melee Weapon Attack:} 6 to hit, reach 5 ft., one target. \textsl{Hit:} 6 (1d6 + 3) slashing damage.

\end{DndMonster}



\begin{DndMonster}[float=t]{Young Blue Dragon}
\DndMonsterType{Large dragon, lawful evil}
\DndMonsterBasics[
armor-class = {18 (natural armor)},
hit-points = {\DndDice{16d10 + 64}},
speed = {40 ft., burrow 20 ft., fly 80 ft.}
]
\DndMonsterAbilityScores[
 str = 21,
 dex = 10,
 con = 19,
 int = 14,
 wis = 13,
 cha = 17
]
\DndMonsterDetails[
saving-throws = {Dex +4, Con +8, Wis +5, Cha +7},
skills = {Perception +9, Stealth +4},
damage-immunities = {lightning},
senses = {blindsight 30 ft., darkvision 120 ft., passive perception 19},
languages = {Common, Draconic},
challenge = 9,
]
\DndMonsterSection{Actions}
\DndMonsterAction{Multiattack}
The dragon makes three attacks: one with its bite and two with its claws.

\DndMonsterAction{Bite}
\textsl{Melee Weapon Attack:} 9 to hit, reach 10 ft., one target. \textsl{Hit:} 16 (2d10 + 5) piercing damage plus 5 (1d10) lightning damage.

\DndMonsterAction{Claw}
\textsl{Melee Weapon Attack:} 9 to hit, reach 5 ft., one target. \textsl{Hit:} 12 (2d6 + 5) slashing damage.

\DndMonsterAction{Lightning Breath (Recharge 5\textendash 6)}
The dragon exhales lightning in a 60-foot line that is 5 feet wide. Each creature in that line must make a DC 16 Dexterity saving throw, taking 55 (10d10) lightning damage on a failed save, or half as much damage on a successful one.

\end{DndMonster}



\begin{DndMonster}[float=t]{Young Bronze Dragon}
\DndMonsterType{Large dragon, lawful good}
\DndMonsterBasics[
armor-class = {18 (natural armor)},
hit-points = {\DndDice{15d10 + 60}},
speed = {40 ft., fly 80 ft., swim 40 ft.}
]
\DndMonsterAbilityScores[
 str = 21,
 dex = 10,
 con = 19,
 int = 14,
 wis = 13,
 cha = 17
]
\DndMonsterDetails[
saving-throws = {Dex +3, Con +7, Wis +4, Cha +6},
skills = {Insight +4, Perception +7, Stealth +3},
damage-immunities = {lightning},
senses = {blindsight 30 ft., darkvision 120 ft., passive perception 17},
languages = {Common, Draconic},
challenge = 8,
]
\DndMonsterAction{Amphibious}
The dragon can breathe air and water.

\DndMonsterSection{Actions}
\DndMonsterAction{Multiattack}
The dragon makes three attacks: one with its bite and two with its claws.

\DndMonsterAction{Bite}
\textsl{Melee Weapon Attack:} 8 to hit, reach 10 ft., one target. \textsl{Hit:} 16 (2d10 + 5) piercing damage.

\DndMonsterAction{Claw}
\textsl{Melee Weapon Attack:} 8 to hit, reach 5 ft., one target. \textsl{Hit:} 12 (2d6 + 5) slashing damage.

\DndMonsterAction{Breath Weapons (Recharge 5\textendash 6)}
The dragon uses one of the following breath weapons.


\begin{description}
\item[Lightning Breath.]
The dragon exhales lightning in a 60-foot line that is 5 feet wide. Each creature in that line must make a DC 15 Dexterity saving throw, taking 55 (10d10) lightning damage on a failed save, or half as much damage on a successful one.
\item[Repulsion Breath.]
The dragon exhales repulsion energy in a 30-foot cone. Each creature in that area must succeed on a DC 15 Strength saving throw. On a failed save, the creature is pushed 40 feet away from the dragon.
\end{description}

\end{DndMonster}



\begin{DndMonster}[float=t]{Young Red Dragon}
\DndMonsterType{Large dragon, chaotic evil}
\DndMonsterBasics[
armor-class = {18 (natural armor)},
hit-points = {\DndDice{17d10 + 85}},
speed = {40 ft., climb 40 ft., fly 80 ft.}
]
\DndMonsterAbilityScores[
 str = 23,
 dex = 10,
 con = 21,
 int = 14,
 wis = 11,
 cha = 19
]
\DndMonsterDetails[
saving-throws = {Dex +4, Con +9, Wis +4, Cha +8},
skills = {Perception +8, Stealth +4},
damage-immunities = {fire},
senses = {blindsight 30 ft., darkvision 120 ft., passive perception 18},
languages = {Common, Draconic},
challenge = 10,
]
\DndMonsterSection{Actions}
\DndMonsterAction{Multiattack}
The dragon makes three attacks: one with its bite and two with its claws.

\DndMonsterAction{Bite}
\textsl{Melee Weapon Attack:} 10 to hit, reach 10 ft., one target. \textsl{Hit:} 17 (2d10 + 6) piercing damage plus 3 (1d6) fire damage.

\DndMonsterAction{Claw}
\textsl{Melee Weapon Attack:} 10 to hit, reach 5 ft., one target. \textsl{Hit:} 13 (2d6 + 6) slashing damage.

\DndMonsterAction{Fire Breath (Recharge 5\textendash 6)}
The dragon exhales fire in a 30-foot cone. Each creature in that area must make a DC 17 Dexterity saving throw, taking 56 (16d6) fire damage on a failed save, or half as much damage on a successful one.

\end{DndMonster}



\begin{DndMonster}[float=t]{Young Red Shadow Dragon}
\DndMonsterType{Large dragon, chaotic evil}
\DndMonsterBasics[
armor-class = {18 (natural armor)},
hit-points = {\DndDice{17d10 + 85}},
speed = {40 ft., climb 40 ft., fly 80 ft.}
]
\DndMonsterAbilityScores[
 str = 23,
 dex = 10,
 con = 21,
 int = 14,
 wis = 11,
 cha = 19
]
\DndMonsterDetails[
saving-throws = {Dex +5, Con +10, Wis +5, Cha +9},
skills = {Perception +10, Stealth +10},
damage-resistances = {necrotic},
damage-immunities = {fire},
senses = {blindsight 30 ft., darkvision 120 ft., passive perception 18},
languages = {Common, Draconic},
challenge = 13,
]
\DndMonsterAction{Living Shadow}
While in dim light or darkness, the dragon has resistance to damage that isn't force, psychic, or radiant.

\DndMonsterAction{Shadow Stealth}
While in dim light or darkness, the dragon can take the Hide action as a bonus action.

\DndMonsterAction{Sunlight Sensitivity}
While in sunlight, the dragon has disadvantage on attack rolls, as well as on Wisdom (Perception) checks that rely on sight.

\DndMonsterSection{Actions}
\DndMonsterAction{Multiattack}
The dragon makes three attacks: one with its bite and two with its claws.

\DndMonsterAction{Bite}
\textsl{Melee Weapon Attack:} 11 to hit, reach 10 ft., one target. \textsl{Hit:} 17 (2d10 + 6) piercing damage plus 3 (1d6) necrotic damage.

\DndMonsterAction{Claw}
\textsl{Melee Weapon Attack:} 11 to hit, reach 5 ft., one target. \textsl{Hit:} 13 (2d6 + 6) slashing damage.

\DndMonsterAction{Shadow Breath (Recharge 5\textendash 6)}
The dragon exhales shadowy fire in a 30-foot cone. Each creature in that area must make a DC 18 Dexterity saving throw, taking 56 (16d6) necrotic damage on a failed save, or half as much damage on a successful one. A humanoid reduced to 0 hit points by this damage dies, and an undead \textbf{shadow} rises from its corpse and acts immediately after the dragon in the initiative count. The shadow is under the dragon's control.

\end{DndMonster}



\begin{DndMonster}[float*=b,width=\textwidth + 8pt]{Zargon the Returner}
\begin{multicols}{2}
\DndMonsterType{Huge aberration, lawful evil}
\DndMonsterBasics[
armor-class = {18 (natural armor)},
hit-points = {\DndDice{22d12 + 110}},
speed = {40 ft., swim 80 ft.}
]
\DndMonsterAbilityScores[
 str = 22,
 dex = 10,
 con = 20,
 int = 14,
 wis = 18,
 cha = 18
]
\DndMonsterDetails[
saving-throws = {Dex +6, Cha +10},
skills = {History +8, Perception +10},
damage-resistances = {cold; fire; bludgeoning, piercing, slashing from nonmagical attacks},
damage-immunities = {acid, poison},
condition-immunities = {exhaustion, poisoned},
senses = {truesight 120 ft., passive perception 20},
languages = {all, telepathy 120 ft.},
challenge = 17,
]
\DndMonsterAction{Legendary Resistance (4/Day)}
If Zargon fails a saving throw, it can choose to succeed instead.

\DndMonsterAction{Regeneration}
Zargon regains 20 hit points at the start of each of its turns. If Zargon takes cold or fire damage, this trait doesn't function at the start of Zargon's next turn. Zargon dies only if it starts its turn with 0 hit points and doesn't regenerate.

\DndMonsterAction{Shrouded Being}
Zargon can't be targeted by divination magic or perceived through magical scrying sensors.

\DndMonsterAction{Slimy Demise}
When Zargon dies, its body dissolves into foul slime, leaving only its horn behind. Zargon re-forms in 1d10 days, regrowing from the horn. The horn is immune to all damage and can be destroyed only by submerging it in a cleansing waterfall on one of the Upper Planes for 101 days. While the horn is submerged in this way, Zargon doesn't re-form, and the horn slowly dissolves, sending corrupting slime downriver that permanently fouls the water for 10 miles from the place where the horn dissolved. The fouled water is unfit to drink, chokes aquatic wildlife, and withers plants.

\DndMonsterSection{Actions}
\DndMonsterAction{Multiattack}
Zargon makes two Barbed Tentacle attacks, one Bite attack, and one Gore attack.

\DndMonsterAction{Barbed Tentacle}
\textsl{Melee Weapon Attack:} 12 to hit, reach 20 ft., one target. \textsl{Hit:} 15 (2d8 + 6) piercing damage. If the target is a Large or smaller creature, it has the grappled condition (escape DC 20), and Zargon can pull the creature up to 20 feet straight toward itself. Zargon has six tentacles, each of which can grapple one creature. Zargon can move at its full speed while dragging creatures it is grappling.

\DndMonsterAction{Bite}
\textsl{Melee Weapon Attack:} 12 to hit, reach 5 ft., one target. \textsl{Hit:} 19 (2d12 + 6) piercing damage plus 7 (2d6) acid damage.

\DndMonsterAction{Gore}
\textsl{Melee Weapon Attack:} 12 to hit, reach 10 ft., one target. \textsl{Hit:} 15 (2d8 + 6) force damage, and a 10-foot-radius invisible sphere of antimagic, like that created by an \textit{Antimagic Field} spell, surrounds the target. The sphere is centered on the target, moves with the target, and lasts until the end of Zargon's next turn.

\DndMonsterAction{Slime Wave (Recharge 5\textendash 6)}
Zargon spews slime in a 60-foot cone. Each creature in that area that isn't an Aberration or Ooze must make a DC 19 Constitution saving throw. On a failed save, the creature takes 38 (7d10) acid damage and has the poisoned condition for 1 minute. On a successful save, the creature takes half as much damage only. A poisoned creature can repeat the saving throw at the end of each of its turns, ending the effect on itself on a success.

A creature reduced to 0 hit points by the acid damage dies and dissolves into a puddle of slime that rises as a gibbering mouther at the start of Zargon's next turn. The creature obeys Zargon's commands and takes its turn immediately after Zargon's. Only a \textit{Wish} spell can reverse this transformation and restore the creature to life.

\DndMonsterSection{Reactions}
\DndMonsterAction{Defiant Essence}
When a creature casts a spell that targets Zargon or would deal damage to it, Zargon attempts to absorb the magic into its horn. The creature must make a DC 19 Charisma saving throw. On a failed save, the creature takes 6 (1d12) force damage, and the spell it cast fails and is wasted.

\DndMonsterAction{Slime Spray}
When a creature ends its turn within 30 feet of Zargon, Zargon sprays toxic slime at the creature. The target must make a DC 19 Dexterity saving throw (with disadvantage if it has the poisoned condition). On a failed save, the creature takes 7 (2d6) poison damage. On a successful save, it takes half as much damage.

\DndMonsterSection{Lair Actions}
On initiative count 20 (losing initiative ties), Zargon can take one of the following lair actions; it can't take the same lair action two rounds in a row:


\begin{description}
\item[Dissolution.]
Zargon dissolves into a glob of slime, seeps to an unoccupied space anywhere in the lair without provoking opportunity attacks, then re-forms.
\item[Spell Breaker.]
Zargon targets up to two creatures it can see within the lair. All spells of 5th level or lower affecting the targets end.
\item[Sucking Slime.]
Zargon creates a 20-foot-radius puddle of slime on the ground centered on a point within the lair. The slime is difficult terrain, and each creature in that area when the slime appears has the grappled condition (escape DC 15). Aberrations and Oozes are unaffected by the slime. The slime lasts until Zargon dies or until it uses this lair action again.
\end{description}

\DndMonsterSection{Regional Effects}
The region containing Zargon's lair is warped by Zargon's unnatural presence, creating one or more of the following effects:


\begin{description}
\item[Arcane Enervation.]
Enemies of Zargon within 1 mile of the lair have disadvantage on saving throws made to maintain concentration.
\item[Corrupted Water.]
Rain and water sources within 1 mile of the lair are tainted. A creature that drinks the water must succeed on a DC 15 Constitution saving throw or have the poisoned condition until it finishes a long rest.
\item[Eerie Weather.]
The area within 10 miles of the lair experiences strange weather patterns that defy the region's typical climate.
\end{description}

If Zargon dies, these effects fade over the course of 1d10 days.

\end{multicols}
\end{DndMonster}



\begin{DndMonster}[float=t]{Zombie}
\DndMonsterType{Medium undead, neutral evil}
\DndMonsterBasics[
armor-class = {8},
hit-points = {\DndDice{3d8 + 9}},
speed = {20 ft.}
]
\DndMonsterAbilityScores[
 str = 13,
 dex = 6,
 con = 16,
 int = 3,
 wis = 6,
 cha = 5
]
\DndMonsterDetails[
saving-throws = {Wis +0},
damage-immunities = {poison},
condition-immunities = {poisoned},
senses = {darkvision 60 ft., passive perception 8},
languages = {understands all languages it spoke in life but can't speak},
challenge = 1/4,
]
\DndMonsterAction{Undead Fortitude}
If damage reduces the zombie to 0 hit points, it must make a Constitution saving throw with a DC of 5 + the damage taken, unless the damage is radiant or from a critical hit. On a success, the zombie drops to 1 hit point instead.

\DndMonsterSection{Actions}
\DndMonsterAction{Slam}
\textsl{Melee Weapon Attack:} 3 to hit, reach 5 ft., one target. \textsl{Hit:} 4 (1d6 + 1) bludgeoning damage.

\end{DndMonster}


\end{document}
